\section*{Introduction}
\begin{center}\Large\textbf{MANIFESTE RBK 2.0}\end{center}
\vspace{0.5em}
\vspace{0.5em}
\vspace{0.5em}
\vspace{0.5em}
EL\\
Le Paradigme « Senior-by-Design »\\
\vspace{0.5em}
\vspace{0.5em}
\vspace{0.5em}
\vspace{0.5em}
TI\\
EN\\
Partenariat exclusif RBK – Nexus Réussite\\
\vspace{0.5em}
\vspace{0.5em}
\vspace{0.5em}
\vspace{0.5em}
D\\
FI\\
N\\
Alaeddine BEN RHOUMA\\
CO\\
En partenariat avec Nexus Réussite – Maître d’œuvre pédagogique\\
\vspace{0.5em}
\vspace{0.5em}
RBK 2.0 x Nexus Réussite\\
Utilisant la plateforme Venture Engine de Money Factory AI\\
\vspace{0.5em}
Décembre 2025\\
\vspace{0.5em}
Ce programme est opéré par RBK et exécuté pédagogiquement par Nexus Réussite, utilisant la plateforme Venture Engine de Money Factory AI.\\
RBK 2.0 — Whitepaper Architecte Web3\\
\vspace{0.5em}
\vspace{0.5em}
\vspace{0.5em}
TABLE DES MATIÈRES\\
\vspace{0.5em}
Guide de Lecture . . . . . . . . . . . . . . . . . . . . . . . . . . . . . . . . . . . . . . . . . . . . . . . . . . . . . . . 18\\
Liste des Acronymes . . . . . . . . . . . . . . . . . . . . . . . . . . . . . . . . . . . . . . . . . . . . . . . . . . . . . 21\\
Factsheet RBK 2.0 . . . . . . . . . . . . . . . . . . . . . . . . . . . . . . . . . . . . . . . . . . . . . . . . . . . . . . . 22\\
Executive Summary . . . . . . . . . . . . . . . . . . . . . . . . . . . . . . . . . . . . . . . . . . . . . . . . . . . . . . 27\\
1     Vision \textbackslash{}& Manifeste . . . . . . . . . . . . . . . . . . . . . . . . . . . . . . . . . . . . . . . . . . . . . . . . . . . 30\\
\vspace{0.5em}
\vspace{0.5em}
\vspace{0.5em}
\vspace{0.5em}
EL\\
\subsection{1.1     La Thèse Centrale : Former des Architectes, pas des Codeurs . . . . . . . . . . . . . . . . . . . . . . . . . . . . . . . . . . . . . . . . . . . . . . . . . . . . . . . . . . . . . . . . . . . . . 30}
\subsection{1.1.0.1           Définition opérationnelle d’un Architecte Web3 . . . . . . . . . . . . . . . . . . . . . . . . . . . . . . .                                                                                                               31}
\vspace{0.5em}
\vspace{0.5em}
\vspace{0.5em}
\vspace{0.5em}
TI\\
\subsection{1.1.0.2           Pourquoi le ”code basique” ne suffit plus à l’ère des LLM . . . . . . . . . . . . . . . . . . . . . . . . . .                                                                                                              32}
\subsection{1.1.0.3           Le Mécanisme Senior-by-Design . . . . . . . . . . . . . . . . . . . . . . . . . . . . . . . . . . . . . . .                                                                                                                32}
\vspace{0.5em}
\vspace{0.5em}
\vspace{0.5em}
\vspace{0.5em}
EN\\
\subsection{1.1.0.4           Métriques de Succès et Méthode de Mesure . . . . . . . . . . . . . . . . . . . . . . . . . . . . . . . . .                                                                                                                 33}
\vspace{0.5em}
\vspace{0.5em}
\vspace{0.5em}
\vspace{0.5em}
D\\
\subsection{1.2     Pourquoi RBK 2.0 ? . . . . . . . . . . . . . . . . . . . . . . . . . . . . . . . . . . . . . . . . . . . . . . . . . . . . . . . . . . . . . . . . . . . . . . . . . . . . . . . . . . . . . . . . . . . . . . . . . . . . . . . . . . . . . . 34}
\vspace{0.5em}
\vspace{0.5em}
\vspace{0.5em}
\vspace{0.5em}
FI\\
\subsection{1.2.0.1 Diagnostic : L’Écart de Compétence (Skills Gap) . . . . . . . . . . . . . . . . . . . . . . . . . . . . . . .                                                                                                                        34}
\subsection{1.2.0.2 Les Différenciateurs RBK 2.0 . . . . . . . . . . . . . . . . . . . . . . . . . . . . . . . . . . . . . . . . .                                                                                                                       34}
N\\
\subsection{1.2.0.3 Ce que RBK 2.0 n’est pas . . . . . . . . . . . . . . . . . . . . . . . . . . . . . . . . . . . . . . . . . . .                                                                                                                       38}
CO\\
\subsection{1.2.0.4 Changement de Paradigme . . . . . . . . . . . . . . . . . . . . . . . . . . . . . . . . . . . . . . . . . .                                                                                                                          39}
2     Stratégie : Note de Cadrage . . . . . . . . . . . . . . . . . . . . . . . . . . . . . . . . . . . . . . . . . . . . . 40\\
\subsection{2.1     Contexte du projet . . . . . . . . . . . . . . . . . . . . . . . . . . . . . . . . . . . . . . . . . . . . . . . . . . . . . . . . . . . . . . . . . . . . . . . . . . . . . . . . . . . . . . . . . . . . . . . . . . . . . . . . . . . . . . . 40}
\subsection{2.1.1               Situation actuelle . . . . . . . . . . . . . . . . . . . . . . . . . . . . . . . . . . . . . . . . . . . . . . . . . . . . . . . . . . . . . . . . . . . . . . . . . . . . . . . . . . . . . . . . . . . . . . . . . . . . . . 40}
\subsection{2.1.2               Motivations et enjeux . . . . . . . . . . . . . . . . . . . . . . . . . . . . . . . . . . . . . . . . . . . . . . . . . . . . . . . . . . . . . . . . . . . . . . . . . . . . . . . . . . . . . . . . . . . . . . . . . . . 41}
\subsection{2.2     Enjeux (Stratégiques, Opérationnels, Réglementaires) . . . . . . . . . . . . . . . . . . . . . . . . . . . . . . . . . . . . . . . . . . . . . . . . . . . . . . . . . . . . . . . . . . . . . . . . . . . . 41}
\subsection{2.3     Matrice SWOT (Synthèse) . . . . . . . . . . . . . . . . . . . . . . . . . . . . . . . . . . . . . . . . . . . . . . . . . . . . . . . . . . . . . . . . . . . . . . . . . . . . . . . . . . . . . . . . . . . . . . . . . . . . . . . . 42}
\subsection{2.4     Priorisation MoSCoW (Fonctionnalités Clés) . . . . . . . . . . . . . . . . . . . . . . . . . . . . . . . . . . . . . . . . . . . . . . . . . . . . . . . . . . . . . . . . . . . . . . . . . . . . . . . . . . . . . . 43}
\vspace{0.5em}
\vspace{0.5em}
\vspace{0.5em}
\vspace{0.5em}
\begin{confidential}CONFIDENTIEL — PROJET RBK 2.0                                                                              Page 1 sur 316\end{confidential}
RBK 2.0 — Whitepaper Architecte Web3\\
\vspace{0.5em}
\vspace{0.5em}
\subsection{2.5     Objectifs SMART . . . . . . . . . . . . . . . . . . . . . . . . . . . . . . . . . . . . . . . . . . . . . . . . . . . . . . . . . . . . . . . . . . . . . . . . . . . . . . . . . . . . . . . . . . . . . . . . . . . . . . . . . . . . . . . . . 44}
\subsection{2.6     Analyse des besoins . . . . . . . . . . . . . . . . . . . . . . . . . . . . . . . . . . . . . . . . . . . . . . . . . . . . . . . . . . . . . . . . . . . . . . . . . . . . . . . . . . . . . . . . . . . . . . . . . . . . . . . . . . . . . . 44}
\subsection{2.6.1               Audit de l’existant . . . . . . . . . . . . . . . . . . . . . . . . . . . . . . . . . . . . . . . . . . . . . . . . . . . . . . . . . . . . . . . . . . . . . . . . . . . . . . . . . . . . . . . . . . . . . . . . . . . . . . 44}
\subsection{2.6.2               Besoins Utilisateurs et Parties Prenantes . . . . . . . . . . . . . . . . . . . . . . . . . . . . . . . . . . . . . . . . . . . . . . . . . . . . . . . . . . . . . . . . . . . . . . . . . . . . . . . . . . 45}
\subsection{2.7     Solutions techniques . . . . . . . . . . . . . . . . . . . . . . . . . . . . . . . . . . . . . . . . . . . . . . . . . . . . . . . . . . . . . . . . . . . . . . . . . . . . . . . . . . . . . . . . . . . . . . . . . . . . . . . . . . . . . 45}
\subsection{2.7.1               Justification du Dual-Track (EVM + Solana) . . . . . . . . . . . . . . . . . . . . . . . . . . . . . . . . . . . . . . . . . . . . . . . . . . . . . . . . . . . . . . . . . . . . . . . . . . . . . . 45}
\subsection{2.7.2               Architecture Technique du Studio . . . . . . . . . . . . . . . . . . . . . . . . . . . . . . . . . . . . . . . . . . . . . . . . . . . . . . . . . . . . . . . . . . . . . . . . . . . . . . . . . . . . . . . . 46}
\subsection{2.8     Évaluation des risques . . . . . . . . . . . . . . . . . . . . . . . . . . . . . . . . . . . . . . . . . . . . . . . . . . . . . . . . . . . . . . . . . . . . . . . . . . . . . . . . . . . . . . . . . . . . . . . . . . . . . . . . . . . 46}
\subsection{2.9     Gouvernance \textbackslash{}& pilotage . . . . . . . . . . . . . . . . . . . . . . . . . . . . . . . . . . . . . . . . . . . . . . . . . . . . . . . . . . . . . . . . . . . . . . . . . . . . . . . . . . . . . . . . . . . . . . . . . . . . . . . . . 46}
\subsection{2.9.1               Méthodologie . . . . . . . . . . . . . . . . . . . . . . . . . . . . . . . . . . . . . . . . . . . . . . . . . . . . . . . . . . . . . . . . . . . . . . . . . . . . . . . . . . . . . . . . . . . . . . . . . . . . . . . . . . 46}
\vspace{0.5em}
\vspace{0.5em}
\vspace{0.5em}
\vspace{0.5em}
EL\\
\subsection{2.9.2               Addendum — Gouvernance contractuelle . . . . . . . . . . . . . . . . . . . . . . . . . . . . . . . . . . . . . . . . . . . . . . . . . . . . . . . . . . . . . . . . . . . . . . . . . . . . . . . . . 47}
\vspace{0.5em}
\vspace{0.5em}
\vspace{0.5em}
\vspace{0.5em}
TI\\
\subsection{2.10    Matrice RACI (Responsabilités Macro) . . . . . . . . . . . . . . . . . . . . . . . . . . . . . . . . . . . . . . . . . . . . . . . . . . . . . . . . . . . . . . . . . . . . . . . . . . . . . . . . . . . . . . . . . . . 47}
\subsection{2.11    Planification, budget, indicateurs . . . . . . . . . . . . . . . . . . . . . . . . . . . . . . . . . . . . . . . . . . . . . . . . . . . . . . . . . . . . . . . . . . . . . . . . . . . . . . . . . . . . . . . . . . . . . . . . 47}
\vspace{0.5em}
\vspace{0.5em}
\vspace{0.5em}
\vspace{0.5em}
EN\\
\subsection{2.11.1              Rétroplanning des Grands Jalons . . . . . . . . . . . . . . . . . . . . . . . . . . . . . . . . . . . . . . . . . . . . . . . . . . . . . . . . . . . . . . . . . . . . . . . . . . . . . . . . . . . . . . . . 48}
\subsection{2.11.2              Budget et Coûts . . . . . . . . . . . . . . . . . . . . . . . . . . . . . . . . . . . . . . . . . . . . . . . . . . . . . . . . . . . . . . . . . . . . . . . . . . . . . . . . . . . . . . . . . . . . . . . . . . . . . . . . 48}
\vspace{0.5em}
\vspace{0.5em}
\vspace{0.5em}
\vspace{0.5em}
D\\
\subsection{2.12    Conduite du changement . . . . . . . . . . . . . . . . . . . . . . . . . . . . . . . . . . . . . . . . . . . . . . . . . . . . . . . . . . . . . . . . . . . . . . . . . . . . . . . . . . . . . . . . . . . . . . . . . . . . . . . . 48}
\vspace{0.5em}
\vspace{0.5em}
\vspace{0.5em}
\vspace{0.5em}
FI\\
\subsection{2.13    Conclusion de la Note . . . . . . . . . . . . . . . . . . . . . . . . . . . . . . . . . . . . . . . . . . . . . . . . . . . . . . . . . . . . . . . . . . . . . . . . . . . . . . . . . . . . . . . . . . . . . . . . . . . . . . . . . . . . 48}
\vspace{0.5em}
3                                                                                  N\\
Gouvernance, Partenariat et RACI . . . . . . . . . . . . . . . . . . . . . . . . . . . . . . . . . . . . . . . . . 49\\
CO\\
\subsection{3.1     Cadre général du partenariat . . . . . . . . . . . . . . . . . . . . . . . . . . . . . . . . . . . . . . . . . . . . . . . . . . . . . . . . . . . . . . . . . . . . . . . . . . . . . . . . . . . . . . . . . . . . . . . . . . . . . 49}
\subsection{3.2     Architecture détaillée du partenariat . . . . . . . . . . . . . . . . . . . . . . . . . . . . . . . . . . . . . . . . . . . . . . . . . . . . . . . . . . . . . . . . . . . . . . . . . . . . . . . . . . . . . . . . . . . . . 49}
\subsection{3.3     Matrice RACI (Répartition des responsabilités) . . . . . . . . . . . . . . . . . . . . . . . . . . . . . . . . . . . . . . . . . . . . . . . . . . . . . . . . . . . . . . . . . . . . . . . . . . . . . . . . . . . 50}
\vspace{0.5em}
4     ANALYSE DU CONTEXTE . . . . . . . . . . . . . . . . . . . . . . . . . . . . . . . . . . . . . . . . . . . . . . . 52\\
\subsection{4.1     L’Opportunité Web3 \textbackslash{}& Solana . . . . . . . . . . . . . . . . . . . . . . . . . . . . . . . . . . . . . . . . . . . . . . . . . . . . . . . . . . . . . . . . . . . . . . . . . . . . . . . . . . . . . . . . . . . . . . . . . . . . 52}
\subsection{4.1.0.1           Définitions Minimales (Lexique Opérationnel) . . . . . . . . . . . . . . . . . . . . . . . . . . . . . . . .                                                                                                                 52}
\subsection{4.1.0.2           Segmentation de la Demande . . . . . . . . . . . . . . . . . . . . . . . . . . . . . . . . . . . . . . . . .                                                                                                                 52}
\subsection{4.1.0.3           Indicateurs 2025 actualisés . . . . . . . . . . . . . . . . . . . . . . . . . . . . . . . . . . . . . . . . . .                                                                                                              53}
\vspace{0.5em}
\vspace{0.5em}
\vspace{0.5em}
\begin{confidential}CONFIDENTIEL — PROJET RBK 2.0                                                                               Page 2 sur 316\end{confidential}
RBK 2.0 — Whitepaper Architecte Web3\\
\vspace{0.5em}
\vspace{0.5em}
\subsection{4.1.0.4           Pourquoi Solana est un Accélérateur d’Employabilité                                                         . . . . . . . . . . . . . . . . . . . . . . . . . . . .                                                    53}
\subsection{4.2     Dynamique Salariale . . . . . . . . . . . . . . . . . . . . . . . . . . . . . . . . . . . . . . . . . . . . . . . . . . . . . . . . . . . . . . . . . . . . . . . . . . . . . . . . . . . . . . . . . . . . . . . . . . . . . . . . . . . . . 53}
\subsection{4.2.0.1 Indicateurs 2025 actualisés . . . . . . . . . . . . . . . . . . . . . . . . . . . . . . . . . . . . . . . . . .                                                                                                                  54}
\subsection{4.2.0.2 Hypothèses de Lecture (TND vs USD) . . . . . . . . . . . . . . . . . . . . . . . . . . . . . . . . . . . . .                                                                                                                     58}
\subsection{4.2.0.3 Grille de Rémunération Standard . . . . . . . . . . . . . . . . . . . . . . . . . . . . . . . . . . . . . . .                                                                                                                    59}
\subsection{4.2.0.4 Modèle ROI Candidat (Simulation 1 an) . . . . . . . . . . . . . . . . . . . . . . . . . . . . . . . . . . .                                                                                                                      59}
\subsection{4.3     Croissance du Marché . . . . . . . . . . . . . . . . . . . . . . . . . . . . . . . . . . . . . . . . . . . . . . . . . . . . . . . . . . . . . . . . . . . . . . . . . . . . . . . . . . . . . . . . . . . . . . . . . . . . . . . . . . . . 59}
\subsection{4.3.0.1           Définition de l’Index . . . . . . . . . . . . . . . . . . . . . . . . . . . . . . . . . . . . . . . . . . . . .                                                                                                        60}
\subsection{4.3.0.2           Lecture Stratégique . . . . . . . . . . . . . . . . . . . . . . . . . . . . . . . . . . . . . . . . . . . . . .                                                                                                        60}
5       Architecture de Partenariat . . . . . . . . . . . . . . . . . . . . . . . . . . . . . . . . . . . . . . . . . . . . . 62\\
\vspace{0.5em}
\vspace{0.5em}
\vspace{0.5em}
\vspace{0.5em}
EL\\
\subsection{5.1     Le partenariat structurant . . . . . . . . . . . . . . . . . . . . . . . . . . . . . . . . . . . . . . . . . . . . . . . . . . . . . . . . . . . . . . . . . . . . . . . . . . . . . . . . . . . . . . . . . . . . . . . . . . . . . . . . 62}
\vspace{0.5em}
\vspace{0.5em}
\vspace{0.5em}
\vspace{0.5em}
TI\\
\subsection{5.1.1              RBK — l’opérateur et la face publique . . . . . . . . . . . . . . . . . . . . . . . . . . . . . . . . . . . . . . . . . . . . . . . . . . . . . . . . . . . . . . . . . . . . . . . . . . . . . . . . . . . . 62}
\subsection{5.1.2              Nexus Réussite — le Maître d’œuvre pédagogique1 . . . . . . . . . . . . . . . . . . . . . . . . . . . . . . . . . . . . . . . . . . . . . . . . . . . . . . . . . . . . . . . . . . . . . . . . 63}
\vspace{0.5em}
\vspace{0.5em}
\vspace{0.5em}
\vspace{0.5em}
EN\\
\subsection{5.1.3              Gouvernance conjointe et comités . . . . . . . . . . . . . . . . . . . . . . . . . . . . . . . . . . . . . . . . . . . . . . . . . . . . . . . . . . . . . . . . . . . . . . . . . . . . . . . . . . . . . . . 63}
\subsection{5.2     Rôle de Money Factory AI : fournisseur technologique . . . . . . . . . . . . . . . . . . . . . . . . . . . . . . . . . . . . . . . . . . . . . . . . . . . . . . . . . . . . . . . . . . . . . . . . . . . 63}
\vspace{0.5em}
\vspace{0.5em}
\vspace{0.5em}
\vspace{0.5em}
D\\
Flux contractuels et financiers . . . . . . . . . . . . . . . . . . . . . . . . . . . . . . . . . . . . . . . . . . . . . . . . . . . . . . . . . . . . . . . . . . . . . . . . . . . . . . . . . . . . . . . . . . . . . . . . . . . 64\\
\vspace{0.5em}
\vspace{0.5em}
\vspace{0.5em}
\vspace{0.5em}
FI\\
\subsection{5.3}
\subsection{5.4     Sécurisation des intérêts de RBK . . . . . . . . . . . . . . . . . . . . . . . . . . . . . . . . . . . . . . . . . . . . . . . . . . . . . . . . . . . . . . . . . . . . . . . . . . . . . . . . . . . . . . . . . . . . . . . . . 64}
\subsection{5.5                                                                         N}
Matrice RACI macro révisée . . . . . . . . . . . . . . . . . . . . . . . . . . . . . . . . . . . . . . . . . . . . . . . . . . . . . . . . . . . . . . . . . . . . . . . . . . . . . . . . . . . . . . . . . . . . . . . . . . . . . 68\\
CO\\
6       MÉTHODOLOGIE CYBORG 2.0 . . . . . . . . . . . . . . . . . . . . . . . . . . . . . . . . . . . . . . . . . . . . 69\\
\subsection{6.1     Philosophie Pédagogique : Intégration du Bien-être . . . . . . . . . . . . . . . . . . . . . . . . . . . . . . . . . . . . . . . . . . . . . . . . . . . . . . . . . . . . . . . . . . . . . . . . . . . . . . 69}
\subsection{6.2     Modèle 3 Niveaux RBK 3.0 . . . . . . . . . . . . . . . . . . . . . . . . . . . . . . . . . . . . . . . . . . . . . . . . . . . . . . . . . . . . . . . . . . . . . . . . . . . . . . . . . . . . . . . . . . . . . . . . . . . . . . 70}
\subsection{6.2.1              Cycle Skill Mirror \textbackslash{}& Block Check . . . . . . . . . . . . . . . . . . . . . . . . . . . . . . . . . . . . . . . . . . . . . . . . . . . . . . . . . . . . . . . . . . . . . . . . . . . . . . . . . . . . . . . . 70}
\subsection{6.3     Genesis Pool : Bootcamp Sélectif . . . . . . . . . . . . . . . . . . . . . . . . . . . . . . . . . . . . . . . . . . . . . . . . . . . . . . . . . . . . . . . . . . . . . . . . . . . . . . . . . . . . . . . . . . . . . . . . . 71}
\subsection{6.3.1              Structure et Évaluation . . . . . . . . . . . . . . . . . . . . . . . . . . . . . . . . . . . . . . . . . . . . . . . . . . . . . . . . . . . . . . . . . . . . . . . . . . . . . . . . . . . . . . . . . . . . . . . . . 72}
1\\
Gouvernance — Partenaire responsable contractuellement de la conception, de l’exécution et de la qualité pédagogique au nom de RBK, avec obligations de\\
moyens renforcées et reporting dédié.\\
\vspace{0.5em}
\vspace{0.5em}
\vspace{0.5em}
\vspace{0.5em}
\begin{confidential}CONFIDENTIEL — PROJET RBK 2.0                                                                               Page 3 sur 316\end{confidential}
RBK 2.0 — Whitepaper Architecte Web3\\
\vspace{0.5em}
\vspace{0.5em}
\subsection{6.4     Gamification Web3 et Système de Réputation . . . . . . . . . . . . . . . . . . . . . . . . . . . . . . . . . . . . . . . . . . . . . . . . . . . . . . . . . . . . . . . . . . . . . . . . . . . . . . . . . . . . 72}
\subsection{6.5     Protocole Anti-Burnout . . . . . . . . . . . . . . . . . . . . . . . . . . . . . . . . . . . . . . . . . . . . . . . . . . . . . . . . . . . . . . . . . . . . . . . . . . . . . . . . . . . . . . . . . . . . . . . . . . . . . . . . . . 73}
\subsection{6.5.1              Organigramme de Monitoring . . . . . . . . . . . . . . . . . . . . . . . . . . . . . . . . . . . . . . . . . . . . . . . . . . . . . . . . . . . . . . . . . . . . . . . . . . . . . . . . . . . . . . . . . . . 74}
\subsection{6.5.2              Tableau de bord hebdomadaire . . . . . . . . . . . . . . . . . . . . . . . . . . . . . . . . . . . . . . . . . . . . . . . . . . . . . . . . . . . . . . . . . . . . . . . . . . . . . . . . . . . . . . . . . . 74}
\vspace{0.5em}
7     STRUCTURE DU CURSUS AMÉLIORÉE . . . . . . . . . . . . . . . . . . . . . . . . . . . . . . . . . . . . . . . 76\\
\subsection{7.1     Architecture pédagogique — 48 semaines de Genesis vers Launch Proof . . . . . . . . . . . . . . . . . . . . . . . . . . . . . . . . . . . . . . . . . . . . . . . . . . . . . . . . . . 76}
\subsection{7.1.1              Diagramme de Structure du Curriculum (48 Semaines) . . . . . . . . . . . . . . . . . . . . . . . . . . . . . . . . . . . . . . . . . . . . . . . . . . . . . . . . . . . . . . . . . . . . 77}
\subsection{7.1.2              Timeline pédagogique – Semaines 1 à 48 . . . . . . . . . . . . . . . . . . . . . . . . . . . . . . . . . . . . . . . . . . . . . . . . . . . . . . . . . . . . . . . . . . . . . . . . . . . . . . . . . 77}
\subsection{7.2     Nouveau Track C : Web3 Product \textbackslash{}& Ecosystem Strategy . . . . . . . . . . . . . . . . . . . . . . . . . . . . . . . . . . . . . . . . . . . . . . . . . . . . . . . . . . . . . . . . . . . . . . . . . . 77}
Modules du Track C (16 semaines) . . . . . . . . . . . . . . . . . . . . . . . . . . . . . . . . . . . . . . . . . . . . . . . . . . . . . . . . . . . . . . . . . . . . . . . . . . . . . . . . . . . . . . . 78\\
\vspace{0.5em}
\vspace{0.5em}
\vspace{0.5em}
\vspace{0.5em}
EL\\
\subsection{7.2.1}
\subsection{7.2.2              Compétences Visées du Track C . . . . . . . . . . . . . . . . . . . . . . . . . . . . . . . . . . . . . . . . . . . . . . . . . . . . . . . . . . . . . . . . . . . . . . . . . . . . . . . . . . . . . . . . . 78}
\vspace{0.5em}
\vspace{0.5em}
\vspace{0.5em}
\vspace{0.5em}
TI\\
\subsection{7.3     Timeline pédagogique – Semaines 1 à 48 . . . . . . . . . . . . . . . . . . . . . . . . . . . . . . . . . . . . . . . . . . . . . . . . . . . . . . . . . . . . . . . . . . . . . . . . . . . . . . . . . . . . . . . . . 79}
\vspace{0.5em}
\vspace{0.5em}
\vspace{0.5em}
\vspace{0.5em}
EN\\
8     SYLLABUS TECHNIQUE . . . . . . . . . . . . . . . . . . . . . . . . . . . . . . . . . . . . . . . . . . . . . . . . 80\\
\subsection{8.1     Calendrier Pédagogique Global . . . . . . . . . . . . . . . . . . . . . . . . . . . . . . . . . . . . . . . . . . . . . . . . . . . . . . . . . . . . . . . . . . . . . . . . . . . . . . . . . . . . . . . . . . . . . . . . . . 80}
\vspace{0.5em}
\vspace{0.5em}
\vspace{0.5em}
\vspace{0.5em}
D\\
\subsection{8.2     Genesis Pool : Sélection Intensives (S1-S4) . . . . . . . . . . . . . . . . . . . . . . . . . . . . . . . . . . . . . . . . . . . . . . . . . . . . . . . . . . . . . . . . . . . . . . . . . . . . . . . . . . . . . . . 81}
\vspace{0.5em}
\vspace{0.5em}
\vspace{0.5em}
\vspace{0.5em}
FI\\
\subsection{8.2.1              Plan d’Exécution Explorer (S05-S16) . . . . . . . . . . . . . . . . . . . . . . . . . . . . . . . . . . . . . . . . . . . . . . . . . . . . . . . . . . . . . . . . . . . . . . . . . . . . . . . . . . . . . 81}
\subsection{8.3}
N\\
Tracks Bâtisseur A / B / C (S17-S32) . . . . . . . . . . . . . . . . . . . . . . . . . . . . . . . . . . . . . . . . . . . . . . . . . . . . . . . . . . . . . . . . . . . . . . . . . . . . . . . . . . . . . . . . . . . . . 83\\
CO\\
\subsection{8.4     STUDIO DE PRODUCTION (S29-S48) . . . . . . . . . . . . . . . . . . . . . . . . . . . . . . . . . . . . . . . . . . . . . . . . . . . . . . . . . . . . . . . . . . . . . . . . . . . . . . . . . . . . . . . . . . . . 86}
\subsection{8.5     Modules de Diversification (Electifs) . . . . . . . . . . . . . . . . . . . . . . . . . . . . . . . . . . . . . . . . . . . . . . . . . . . . . . . . . . . . . . . . . . . . . . . . . . . . . . . . . . . . . . . . . . . . . 88}
\subsection{8.5.1              Module ZK : Zero-Knowledge Proofs (8 semaines) . . . . . . . . . . . . . . . . . . . . . . . . . . . . . . . . . . . . . . . . . . . . . . . . . . . . . . . . . . . . . . . . . . . . . . . . . 88}
\subsection{8.5.2              Module DePIN : Decentralized Physical Infra (6 semaines) . . . . . . . . . . . . . . . . . . . . . . . . . . . . . . . . . . . . . . . . . . . . . . . . . . . . . . . . . . . . . . . . . 88}
\subsection{8.5.3              Module Cross-Chain \textbackslash{}& Interop (4 semaines) . . . . . . . . . . . . . . . . . . . . . . . . . . . . . . . . . . . . . . . . . . . . . . . . . . . . . . . . . . . . . . . . . . . . . . . . . . . . . . 88}
\vspace{0.5em}
9     Track A : Solana Engineer . . . . . . . . . . . . . . . . . . . . . . . . . . . . . . . . . . . . . . . . . . . . . . 89\\
\subsection{9.1     Philosophie du Track : L’Excellence par Rust . . . . . . . . . . . . . . . . . . . . . . . . . . . . . . . . . . . . . . . . . . . . . . . . . . . . . . . . . . . . . . . . . . . . . . . . . . . . . . . . . . . . . 89}
\subsection{9.2     Structure Pédagogique : De l’Architecture au Produit (16 Semaines) . . . . . . . . . . . . . . . . . . . . . . . . . . . . . . . . . . . . . . . . . . . . . . . . . . . . . . . . . . . . . . 91}
\subsection{9.2.1              MODULE 1 : Le Modèle Solana \textbackslash{}& Rust Natif (Semaines 13-16) . . . . . . . . . . . . . . . . . . . . . . . . . . . . . . . . . . . . . . . . . . . . . . . . . . . . . . . . . . . . . . 91}
\vspace{0.5em}
\vspace{0.5em}
\vspace{0.5em}
\begin{confidential}CONFIDENTIEL — PROJET RBK 2.0                                                                            Page 4 sur 316\end{confidential}
RBK 2.0 — Whitepaper Architecte Web3\\
\vspace{0.5em}
\vspace{0.5em}
\subsection{9.2.2               MODULE 2 : Maîtrise du Framework Anchor (Semaines 17-20) . . . . . . . . . . . . . . . . . . . . . . . . . . . . . . . . . . . . . . . . . . . . . . . . . . . . . . . . . . . . . 92}
\subsection{9.2.3               MODULE 3 : Architectures Avancées \textbackslash{}& Innovation (Semaines 21-24) . . . . . . . . . . . . . . . . . . . . . . . . . . . . . . . . . . . . . . . . . . . . . . . . . . . . . . . . 93}
\subsection{9.2.4               MODULE 4 : Production Hardening \textbackslash{}& UX Performance (Semaines 25-28) . . . . . . . . . . . . . . . . . . . . . . . . . . . . . . . . . . . . . . . . . . . . . . . . . . . . 93}
\subsection{9.3     Stack Technique Spécifique . . . . . . . . . . . . . . . . . . . . . . . . . . . . . . . . . . . . . . . . . . . . . . . . . . . . . . . . . . . . . . . . . . . . . . . . . . . . . . . . . . . . . . . . . . . . . . . . . . . . . . 94}
\subsection{9.4     RBK Solana Validator Track . . . . . . . . . . . . . . . . . . . . . . . . . . . . . . . . . . . . . . . . . . . . . . . . . . . . . . . . . . . . . . . . . . . . . . . . . . . . . . . . . . . . . . . . . . . . . . . . . . . . . . 94}
\subsection{9.5     Profil de Sortie : Le « Guardian » . . . . . . . . . . . . . . . . . . . . . . . . . . . . . . . . . . . . . . . . . . . . . . . . . . . . . . . . . . . . . . . . . . . . . . . . . . . . . . . . . . . . . . . . . . . . . . . . . 96}
\vspace{0.5em}
10    Track B : EVM Engineer . . . . . . . . . . . . . . . . . . . . . . . . . . . . . . . . . . . . . . . . . . . . . . . . 97\\
\subsection{10.1    Philosophie du Track : La Maîtrise du Standard Industriel . . . . . . . . . . . . . . . . . . . . . . . . . . . . . . . . . . . . . . . . . . . . . . . . . . . . . . . . . . . . . . . . . . . . . . . . 97}
\subsection{10.2    Structure Pédagogique : De la Logique au Durcissement (16 Semaines) . . . . . . . . . . . . . . . . . . . . . . . . . . . . . . . . . . . . . . . . . . . . . . . . . . . . . . . . . . . 98}
\vspace{0.5em}
\vspace{0.5em}
\vspace{0.5em}
\vspace{0.5em}
EL\\
\subsection{10.2.1              MODULE 1 : Smart Contract Basics \textbackslash{}& Solidity Deep Dive (Semaines 13-15) . . . . . . . . . . . . . . . . . . . . . . . . . . . . . . . . . . . . . . . . . . . . . . . . . . 98}
\subsection{10.2.2              MODULE 2 : Environnement de Développement Pro (Semaines 16-18) . . . . . . . . . . . . . . . . . . . . . . . . . . . . . . . . . . . . . . . . . . . . . . . . . . . . . . 99}
\vspace{0.5em}
\vspace{0.5em}
\vspace{0.5em}
\vspace{0.5em}
TI\\
\subsection{10.2.3              MODULE 3 : Token Standards \textbackslash{}& Composabilité (Semaines 19-21) . . . . . . . . . . . . . . . . . . . . . . . . . . . . . . . . . . . . . . . . . . . . . . . . . . . . . . . . . . . 99}
\vspace{0.5em}
\vspace{0.5em}
\vspace{0.5em}
\vspace{0.5em}
EN\\
\subsection{10.2.4              MODULE 4 : dApp Development \textbackslash{}& Web3 Integration (Semaines 22-24) . . . . . . . . . . . . . . . . . . . . . . . . . . . . . . . . . . . . . . . . . . . . . . . . . . . . . 99}
\subsection{10.2.5              MODULE 5 : L2 Scaling \textbackslash{}& Advanced Patterns (Semaines 25-26) . . . . . . . . . . . . . . . . . . . . . . . . . . . . . . . . . . . . . . . . . . . . . . . . . . . . . . . . . . . . 99}
\vspace{0.5em}
\vspace{0.5em}
\vspace{0.5em}
\vspace{0.5em}
D\\
\subsection{10.2.6              MODULE 6 : Production Hardening \textbackslash{}& Security (Semaines 27-28) . . . . . . . . . . . . . . . . . . . . . . . . . . . . . . . . . . . . . . . . . . . . . . . . . . . . . . . . . . . 99}
\vspace{0.5em}
\vspace{0.5em}
\vspace{0.5em}
\vspace{0.5em}
FI\\
\subsection{10.3    Stack Technique Spécifique . . . . . . . . . . . . . . . . . . . . . . . . . . . . . . . . . . . . . . . . . . . . . . . . . . . . . . . . . . . . . . . . . . . . . . . . . . . . . . . . . . . . . . . . . . . . . . . . . . . . . . 100}
Profil de Sortie : L’Ingénieur d’Infrastructure EVM . . . . . . . . . . . . . . . . . . . . . . . . . . . . . . . . . . . . . . . . . . . . . . . . . . . . . . . . . . . . . . . . . . . . . . . . . . . . . . . 100\\
N\\
\subsection{10.4}
CO\\
11    Track C : Product Engineer . . . . . . . . . . . . . . . . . . . . . . . . . . . . . . . . . . . . . . . . . . . . . . 102\\
\subsection{11.1    Philosophie du Track : Le ”Product Builder” Complet . . . . . . . . . . . . . . . . . . . . . . . . . . . . . . . . . . . . . . . . . . . . . . . . . . . . . . . . . . . . . . . . . . . . . . . . . . . . 102}
\subsection{11.2    Structure Pédagogique : Fondations communes, spécialisations binômes (16 Semaines) . . . . . . . . . . . . . . . . . . . . . . . . . . . . . . . . . . . . . . . . . . . 103}
\subsection{11.2.1              Bloc Fondations (Semaines 13-18) . . . . . . . . . . . . . . . . . . . . . . . . . . . . . . . . . . . . . . . . . . . . . . . . . . . . . . . . . . . . . . . . . . . . . . . . . . . . . . . . . . . . . . . 103}
\subsection{11.2.2              Parcours de spécialisation (Semaines 19-28) . . . . . . . . . . . . . . . . . . . . . . . . . . . . . . . . . . . . . . . . . . . . . . . . . . . . . . . . . . . . . . . . . . . . . . . . . . . . . . 104}
\subsection{11.3    Profil de Sortie . . . . . . . . . . . . . . . . . . . . . . . . . . . . . . . . . . . . . . . . . . . . . . . . . . . . . . . . . . . . . . . . . . . . . . . . . . . . . . . . . . . . . . . . . . . . . . . . . . . . . . . . . . . . . . . . . . 105}
\vspace{0.5em}
12    MODULE SOFT SKILLS \textbackslash{}& PROFESSIONNALISATION . . . . . . . . . . . . . . . . . . . . . . . . . . . . . . 106\\
\subsection{12.1    Structure du Module (4 semaines) . . . . . . . . . . . . . . . . . . . . . . . . . . . . . . . . . . . . . . . . . . . . . . . . . . . . . . . . . . . . . . . . . . . . . . . . . . . . . . . . . . . . . . . . . . . . . . . 106}
\subsection{12.2    Rubrique d’Évaluation . . . . . . . . . . . . . . . . . . . . . . . . . . . . . . . . . . . . . . . . . . . . . . . . . . . . . . . . . . . . . . . . . . . . . . . . . . . . . . . . . . . . . . . . . . . . . . . . . . . . . . . . . . . 108}
\vspace{0.5em}
\vspace{0.5em}
\vspace{0.5em}
\begin{confidential}CONFIDENTIEL — PROJET RBK 2.0                                                                                Page 5 sur 316\end{confidential}
RBK 2.0 — Whitepaper Architecte Web3\\
\vspace{0.5em}
\vspace{0.5em}
13    CAPSTONES (PROJETS SIGNATURES) . . . . . . . . . . . . . . . . . . . . . . . . . . . . . . . . . . . . . . . 110\\
\subsection{13.1    Philosophie du Capstone : Le Standard « Studio » . . . . . . . . . . . . . . . . . . . . . . . . . . . . . . . . . . . . . . . . . . . . . . . . . . . . . . . . . . . . . . . . . . . . . . . . . . . . . . . . . 110}
\subsection{13.2    Les 3 Projets Signatures (Cahier des Charges Industriel) . . . . . . . . . . . . . . . . . . . . . . . . . . . . . . . . . . . . . . . . . . . . . . . . . . . . . . . . . . . . . . . . . . . . . . . . . . 110}
\subsection{13.2.1            Capstone 1 — Wallet \textbackslash{}& Transaction Reliability Pack (Frontend/UX) . . . . . . . . . . . . . . . . . . . . . . . . . . . . . . . . . . . . . . . . . . . . . . . . . . . . . . . . . 110}
\subsection{13.2.1.1 Spécifications Fonctionnelles (User Stories) . . . . . . . . . . . . . . . . . . . . . . . . . . . . . . . . . 111}
\subsection{13.2.1.2 Spécifications Techniques . . . . . . . . . . . . . . . . . . . . . . . . . . . . . . . . . . . . . . . . . . . 111}
\subsection{13.2.1.3 Checklist Sécurité \textbackslash{}& Qualité . . . . . . . . . . . . . . . . . . . . . . . . . . . . . . . . . . . . . . . . . 111}
\subsection{13.2.1.4 Critères d’Acceptation \textbackslash{}& CI . . . . . . . . . . . . . . . . . . . . . . . . . . . . . . . . . . . . . . . . . . 111}
\subsection{13.2.2    Capstone 2 — Tokenization \textbackslash{}& Admin Control Center (RWA/DeFi) . . . . . . . . . . . . . . . . . . . . . . . . . . . . . . . . . . . . . . . . . . . . . . . . . . . . . . . . . . 112}
\subsection{13.2.2.1 Spécifications Fonctionnelles . . . . . . . . . . . . . . . . . . . . . . . . . . . . . . . . . . . . . . . . . 112}
\vspace{0.5em}
\vspace{0.5em}
\vspace{0.5em}
\vspace{0.5em}
EL\\
\subsection{13.2.2.2 Spécifications Techniques . . . . . . . . . . . . . . . . . . . . . . . . . . . . . . . . . . . . . . . . . . . 112}
\subsection{13.2.2.3 Checklist Sécurité Dédiée . . . . . . . . . . . . . . . . . . . . . . . . . . . . . . . . . . . . . . . . . . . 112}
\vspace{0.5em}
\vspace{0.5em}
\vspace{0.5em}
\vspace{0.5em}
TI\\
\subsection{13.2.3    Capstone 3 — Digital Assets \textbackslash{}& Utility Ecosystem (NFT/Gaming) . . . . . . . . . . . . . . . . . . . . . . . . . . . . . . . . . . . . . . . . . . . . . . . . . . . . . . . . . . . . 112}
\vspace{0.5em}
\vspace{0.5em}
\vspace{0.5em}
\vspace{0.5em}
EN\\
\subsection{13.2.3.1 Spécifications Fonctionnelles . . . . . . . . . . . . . . . . . . . . . . . . . . . . . . . . . . . . . . . . . 113}
\subsection{13.2.3.2 Spécifications Techniques . . . . . . . . . . . . . . . . . . . . . . . . . . . . . . . . . . . . . . . . . . . 113}
\vspace{0.5em}
\vspace{0.5em}
\vspace{0.5em}
\vspace{0.5em}
D\\
\subsection{13.2.3.3 Checklist Sécurité Dédiée . . . . . . . . . . . . . . . . . . . . . . . . . . . . . . . . . . . . . . . . . . . 113}
\vspace{0.5em}
\vspace{0.5em}
\vspace{0.5em}
\vspace{0.5em}
FI\\
\subsection{13.3    Matrice d’Évaluation (Rubric Studio) . . . . . . . . . . . . . . . . . . . . . . . . . . . . . . . . . . . . . . . . . . . . . . . . . . . . . . . . . . . . . . . . . . . . . . . . . . . . . . . . . . . . . . . . . . . . . 113}
\subsection{13.4}
N\\
Délivrables de Sortie (Le ”Package”) . . . . . . . . . . . . . . . . . . . . . . . . . . . . . . . . . . . . . . . . . . . . . . . . . . . . . . . . . . . . . . . . . . . . . . . . . . . . . . . . . . . . . . . . . . . . . 114\\
CO\\
14    FICHES MÉTIERS \textbackslash{}& ÉCONOMIE DU DIPLÔMÉ . . . . . . . . . . . . . . . . . . . . . . . . . . . . . . . . . . 115\\
\subsection{14.1    Fiche Métier 1 : Smart Contract Engineer \textbackslash{}& Auditor (Le « Guardian ») . . . . . . . . . . . . . . . . . . . . . . . . . . . . . . . . . . . . . . . . . . . . . . . . . . . . . . . . . . . . . 115}
\subsection{14.2    Fiche Métier 2 : Protocol \textbackslash{}& Ecosystem Strategist (Le « Visionnaire ») . . . . . . . . . . . . . . . . . . . . . . . . . . . . . . . . . . . . . . . . . . . . . . . . . . . . . . . . . . . . . . 117}
\subsection{14.3    Fiche Métier 3 : Web3 Product Builder / Entrepreneur (Le « Builder ») . . . . . . . . . . . . . . . . . . . . . . . . . . . . . . . . . . . . . . . . . . . . . . . . . . . . . . . . . . . . 120}
\subsection{14.4    Fiche Métier 4 : Solana dApp Engineer (Front Web3) . . . . . . . . . . . . . . . . . . . . . . . . . . . . . . . . . . . . . . . . . . . . . . . . . . . . . . . . . . . . . . . . . . . . . . . . . . . . . 122}
\subsection{14.5    Fiche Métier 5 : Tokenization \textbackslash{}& DePIN Architect . . . . . . . . . . . . . . . . . . . . . . . . . . . . . . . . . . . . . . . . . . . . . . . . . . . . . . . . . . . . . . . . . . . . . . . . . . . . . . . . . 124}
\subsection{14.6    Fiche Métier 6 : Web3 QA \textbackslash{}& Test Automation Engineer . . . . . . . . . . . . . . . . . . . . . . . . . . . . . . . . . . . . . . . . . . . . . . . . . . . . . . . . . . . . . . . . . . . . . . . . . . 126}
\subsection{14.7    Fiche Métier 7 : Developer Advocate \textbackslash{}& Technical Writer . . . . . . . . . . . . . . . . . . . . . . . . . . . . . . . . . . . . . . . . . . . . . . . . . . . . . . . . . . . . . . . . . . . . . . . . . 127}
\subsection{14.8    Perspectives Économiques \textbackslash{}& Carrière . . . . . . . . . . . . . . . . . . . . . . . . . . . . . . . . . . . . . . . . . . . . . . . . . . . . . . . . . . . . . . . . . . . . . . . . . . . . . . . . . . . . . . . . . . . . 129}
\vspace{0.5em}
\vspace{0.5em}
\begin{confidential}CONFIDENTIEL — PROJET RBK 2.0                                                                       Page 6 sur 316\end{confidential}
RBK 2.0 — Whitepaper Architecte Web3\\
\vspace{0.5em}
\vspace{0.5em}
\subsection{14.8.1              Revenus Annuels Cibles 2025 . . . . . . . . . . . . . . . . . . . . . . . . . . . . . . . . . . . . . . . . . . . . . . . . . . . . . . . . . . . . . . . . . . . . . . . . . . . . . . . . . . . . . . . . . . . 129}
\subsection{14.8.2              Ce qui compose réellement la rémunération (important) . . . . . . . . . . . . . . . . . . . . . . . . . . . . . . . . . . . . . . . . . . . . . . . . . . . . . . . . . . . . . . . . . . . 130}
\subsection{14.8.3              Comment atteindre le palier . . . . . . . . . . . . . . . . . . . . . . . . . . . . . . . . . . . . . . . . . . . . . . . . . . . . . . . . . . . . . . . . . . . . . . . . . . . . . . . . . . . . . . . . . . . . . 130}
\vspace{0.5em}
15    Business Plan . . . . . . . . . . . . . . . . . . . . . . . . . . . . . . . . . . . . . . . . . . . . . . . . . . . . . . 132\\
\subsection{15.1    Modèle Économique Hybride . . . . . . . . . . . . . . . . . . . . . . . . . . . . . . . . . . . . . . . . . . . . . . . . . . . . . . . . . . . . . . . . . . . . . . . . . . . . . . . . . . . . . . . . . . . . . . . . . . . . 132}
\subsection{15.1.1              Modèle économique détaillé — Waterfall \textbackslash{}& caps mentors . . . . . . . . . . . . . . . . . . . . . . . . . . . . . . . . . . . . . . . . . . . . . . . . . . . . . . . . . . . . . . . . . . 133}
\subsection{15.2    Hypothèses \textbackslash{}& Sources du Modèle . . . . . . . . . . . . . . . . . . . . . . . . . . . . . . . . . . . . . . . . . . . . . . . . . . . . . . . . . . . . . . . . . . . . . . . . . . . . . . . . . . . . . . . . . . . . . . . . 133}
\subsection{15.2.1              Hypothèses Clés . . . . . . . . . . . . . . . . . . . . . . . . . . . . . . . . . . . . . . . . . . . . . . . . . . . . . . . . . . . . . . . . . . . . . . . . . . . . . . . . . . . . . . . . . . . . . . . . . . . . . . . 133}
\subsection{15.2.2              Structure des Coûts Directs . . . . . . . . . . . . . . . . . . . . . . . . . . . . . . . . . . . . . . . . . . . . . . . . . . . . . . . . . . . . . . . . . . . . . . . . . . . . . . . . . . . . . . . . . . . . . . 134}
\subsection{15.3    Funnel d’Acquisition \textbackslash{}& Sourcing . . . . . . . . . . . . . . . . . . . . . . . . . . . . . . . . . . . . . . . . . . . . . . . . . . . . . . . . . . . . . . . . . . . . . . . . . . . . . . . . . . . . . . . . . . . . . . . . . 135}
\vspace{0.5em}
\vspace{0.5em}
\vspace{0.5em}
\vspace{0.5em}
EL\\
\subsection{15.4    Le Pilier B2B : Corporate Upskilling . . . . . . . . . . . . . . . . . . . . . . . . . . . . . . . . . . . . . . . . . . . . . . . . . . . . . . . . . . . . . . . . . . . . . . . . . . . . . . . . . . . . . . . . . . . . . . 135}
\vspace{0.5em}
\vspace{0.5em}
\vspace{0.5em}
\vspace{0.5em}
TI\\
\subsection{15.5    Trajectoire Financière (36 Mois) . . . . . . . . . . . . . . . . . . . . . . . . . . . . . . . . . . . . . . . . . . . . . . . . . . . . . . . . . . . . . . . . . . . . . . . . . . . . . . . . . . . . . . . . . . . . . . . . . 136}
\vspace{0.5em}
\vspace{0.5em}
\vspace{0.5em}
\vspace{0.5em}
EN\\
\subsection{15.6    Analyse de Sensibilité . . . . . . . . . . . . . . . . . . . . . . . . . . . . . . . . . . . . . . . . . . . . . . . . . . . . . . . . . . . . . . . . . . . . . . . . . . . . . . . . . . . . . . . . . . . . . . . . . . . . . . . . . . . . 138}
\subsection{15.6.1              Gestion du Risque Crédit ISA . . . . . . . . . . . . . . . . . . . . . . . . . . . . . . . . . . . . . . . . . . . . . . . . . . . . . . . . . . . . . . . . . . . . . . . . . . . . . . . . . . . . . . . . . . . . 138}
\vspace{0.5em}
\vspace{0.5em}
\vspace{0.5em}
\vspace{0.5em}
D\\
\subsection{15.7    Financements et Partenariats Stratégiques . . . . . . . . . . . . . . . . . . . . . . . . . . . . . . . . . . . . . . . . . . . . . . . . . . . . . . . . . . . . . . . . . . . . . . . . . . . . . . . . . . . . . . . . 139}
\vspace{0.5em}
\vspace{0.5em}
\vspace{0.5em}
\vspace{0.5em}
FI\\
\subsection{15.7.1             1. Écosystème Web3 (Grants) . . . . . . . . . . . . . . . . . . . . . . . . . . . . . . . . . . . . . . . . . . . . . . . . . . . . . . . . . . . . . . . . . . . . . . . . . . . . . . . . . . . . . . . . . . . . 139}
2. Bailleurs de Fonds Institutionnels . . . . . . . . . . . . . . . . . . . . . . . . . . . . . . . . . . . . . . . . . . . . . . . . . . . . . . . . . . . . . . . . . . . . . . . . . . . . . . . . . . . . . . 139\\
N\\
\subsection{15.7.2}
\subsection{15.7.3              3. Modèle de Franchise (Scale Africa) . . . . . . . . . . . . . . . . . . . . . . . . . . . . . . . . . . . . . . . . . . . . . . . . . . . . . . . . . . . . . . . . . . . . . . . . . . . . . . . . . . . . 139}
CO\\
16    Stratégie marketing \textbackslash{}& acquisition . . . . . . . . . . . . . . . . . . . . . . . . . . . . . . . . . . . . . . . . . 142\\
\subsection{16.1    Programme ”Building in Public” . . . . . . . . . . . . . . . . . . . . . . . . . . . . . . . . . . . . . . . . . . . . . . . . . . . . . . . . . . . . . . . . . . . . . . . . . . . . . . . . . . . . . . . . . . . . . . . . . 142}
\subsection{16.2    Assets visuels Sprint 0 (Maquettes à livrer) . . . . . . . . . . . . . . . . . . . . . . . . . . . . . . . . . . . . . . . . . . . . . . . . . . . . . . . . . . . . . . . . . . . . . . . . . . . . . . . . . . . . . . . 143}
\subsection{16.3    Simulateur de ROI Interactif . . . . . . . . . . . . . . . . . . . . . . . . . . . . . . . . . . . . . . . . . . . . . . . . . . . . . . . . . . . . . . . . . . . . . . . . . . . . . . . . . . . . . . . . . . . . . . . . . . . . . 145}
\subsection{16.4    Stratégie Multi-Canaux . . . . . . . . . . . . . . . . . . . . . . . . . . . . . . . . . . . . . . . . . . . . . . . . . . . . . . . . . . . . . . . . . . . . . . . . . . . . . . . . . . . . . . . . . . . . . . . . . . . . . . . . . . 147}
\subsection{16.5    Budget d’Acquisition \textbackslash{}& CAC Cible . . . . . . . . . . . . . . . . . . . . . . . . . . . . . . . . . . . . . . . . . . . . . . . . . . . . . . . . . . . . . . . . . . . . . . . . . . . . . . . . . . . . . . . . . . . . . . . . 147}
\subsection{16.5.1              Discord : Le cœur communautaire de la pré-sélection . . . . . . . . . . . . . . . . . . . . . . . . . . . . . . . . . . . . . . . . . . . . . . . . . . . . . . . . . . . . . . . . . . . . . . 148}
\subsection{16.6    Programme de Référence \textbackslash{}& Bounties . . . . . . . . . . . . . . . . . . . . . . . . . . . . . . . . . . . . . . . . . . . . . . . . . . . . . . . . . . . . . . . . . . . . . . . . . . . . . . . . . . . . . . . . . . . . . 149}
\vspace{0.5em}
\vspace{0.5em}
\vspace{0.5em}
\begin{confidential}CONFIDENTIEL — PROJET RBK 2.0                                                                              Page 7 sur 316\end{confidential}
RBK 2.0 — Whitepaper Architecte Web3\\
\vspace{0.5em}
\vspace{0.5em}
17    ANALYSE DES RISQUES \textbackslash{}& MODÈLE DE RÉSILIENCE . . . . . . . . . . . . . . . . . . . . . . . . . . . . . . 151\\
\subsection{17.1    Risques Réglementaires et Conformité . . . . . . . . . . . . . . . . . . . . . . . . . . . . . . . . . . . . . . . . . . . . . . . . . . . . . . . . . . . . . . . . . . . . . . . . . . . . . . . . . . . . . . . . . . . 151}
\subsection{17.1.1            Loi des Changes et Crypto-Actifs (Tunisie) . . . . . . . . . . . . . . . . . . . . . . . . . . . . . . . . . . . . . . . . . . . . . . . . . . . . . . . . . . . . . . . . . . . . . . . . . . . . . . . . 151}
\subsection{17.1.2            GDPR et Données Étudiantes On-Chain . . . . . . . . . . . . . . . . . . . . . . . . . . . . . . . . . . . . . . . . . . . . . . . . . . . . . . . . . . . . . . . . . . . . . . . . . . . . . . . . . . . 152}
\subsection{17.1.3            Cadre Légal des ISA (Income Share Agreements) . . . . . . . . . . . . . . . . . . . . . . . . . . . . . . . . . . . . . . . . . . . . . . . . . . . . . . . . . . . . . . . . . . . . . . . . . . 152}
\subsection{17.2    Matrice de Risques Dynamique . . . . . . . . . . . . . . . . . . . . . . . . . . . . . . . . . . . . . . . . . . . . . . . . . . . . . . . . . . . . . . . . . . . . . . . . . . . . . . . . . . . . . . . . . . . . . . . . . . 152}
\subsection{17.3    Plan de Continuité Pédagogique (PCP) . . . . . . . . . . . . . . . . . . . . . . . . . . . . . . . . . . . . . . . . . . . . . . . . . . . . . . . . . . . . . . . . . . . . . . . . . . . . . . . . . . . . . . . . . . . 154}
\subsection{17.3.1            Kit de survie pédagogique . . . . . . . . . . . . . . . . . . . . . . . . . . . . . . . . . . . . . . . . . . . . . . . . . . . . . . . . . . . . . . . . . . . . . . . . . . . . . . . . . . . . . . . . . . . . . . . 154}
\subsection{17.3.2            Tableau de réponse rapide . . . . . . . . . . . . . . . . . . . . . . . . . . . . . . . . . . . . . . . . . . . . . . . . . . . . . . . . . . . . . . . . . . . . . . . . . . . . . . . . . . . . . . . . . . . . . . 155}
\subsection{17.3.3            Protocole de tests et amélioration . . . . . . . . . . . . . . . . . . . . . . . . . . . . . . . . . . . . . . . . . . . . . . . . . . . . . . . . . . . . . . . . . . . . . . . . . . . . . . . . . . . . . . . . 155}
\vspace{0.5em}
\vspace{0.5em}
\vspace{0.5em}
\vspace{0.5em}
EL\\
\subsection{17.4    Plan de Réponse aux Incidents Crypto (”Black Swan”) . . . . . . . . . . . . . . . . . . . . . . . . . . . . . . . . . . . . . . . . . . . . . . . . . . . . . . . . . . . . . . . . . . . . . . . . . . . . 156}
Scénario A : Effondrement de l’Écosystème Solana . . . . . . . . . . . . . . . . . . . . . . . . . . . . . . . . . . . . . . . . . . . . . . . . . . . . . . . . . . . . . . . . . . . . . . . . . 156\\
\vspace{0.5em}
\vspace{0.5em}
\vspace{0.5em}
\vspace{0.5em}
TI\\
\subsection{17.4.1}
\subsection{17.4.2            Scénario B : Hack d’un Bridge / Protocole Partenaire . . . . . . . . . . . . . . . . . . . . . . . . . . . . . . . . . . . . . . . . . . . . . . . . . . . . . . . . . . . . . . . . . . . . . . 156}
\vspace{0.5em}
\vspace{0.5em}
\vspace{0.5em}
\vspace{0.5em}
EN\\
\subsection{17.5    Tableau de Bord des Risques Critiques . . . . . . . . . . . . . . . . . . . . . . . . . . . . . . . . . . . . . . . . . . . . . . . . . . . . . . . . . . . . . . . . . . . . . . . . . . . . . . . . . . . . . . . . . . . 157}
\vspace{0.5em}
18    Guide Compliance Web3 . . . . . . . . . . . . . . . . . . . . . . . . . . . . . . . . . . . . . . . . . . . . . . . 158\\
\vspace{0.5em}
\vspace{0.5em}
\vspace{0.5em}
\vspace{0.5em}
D\\
FI\\
\subsection{18.1    Operating Model Compliant : Scénarios pour la Tunisie . . . . . . . . . . . . . . . . . . . . . . . . . . . . . . . . . . . . . . . . . . . . . . . . . . . . . . . . . . . . . . . . . . . . . . . . . . 159}
Scénario A : Exportateur de Services Logiciels (Le Standard) . . . . . . . . . . . . . . . . . . . . . . . . . . . . . . . . . . . . . . . . . . . . . . . . . . . . . . . . . . . . . . . . 159\\
\vspace{0.5em}
N\\
\subsection{18.1.1}
\subsection{18.1.2            Scénario B : Filiale Offshore (Le Scale-Up) . . . . . . . . . . . . . . . . . . . . . . . . . . . . . . . . . . . . . . . . . . . . . . . . . . . . . . . . . . . . . . . . . . . . . . . . . . . . . . . . 159}
CO\\
\subsection{18.1.3            Scénario C : Freelance ”Portage Salarial” (Le Simple) . . . . . . . . . . . . . . . . . . . . . . . . . . . . . . . . . . . . . . . . . . . . . . . . . . . . . . . . . . . . . . . . . . . . . . 160}
\subsection{18.2    Matrice des Risques de Conformité . . . . . . . . . . . . . . . . . . . . . . . . . . . . . . . . . . . . . . . . . . . . . . . . . . . . . . . . . . . . . . . . . . . . . . . . . . . . . . . . . . . . . . . . . . . . . . . 161}
\subsection{18.3    Politique de Communication (Communication Policy) . . . . . . . . . . . . . . . . . . . . . . . . . . . . . . . . . . . . . . . . . . . . . . . . . . . . . . . . . . . . . . . . . . . . . . . . . . . . 161}
\subsection{18.3.1            Interdits Absolus (Red Lines) . . . . . . . . . . . . . . . . . . . . . . . . . . . . . . . . . . . . . . . . . . . . . . . . . . . . . . . . . . . . . . . . . . . . . . . . . . . . . . . . . . . . . . . . . . . . 161}
\subsection{18.3.2            Obligations (Green Lines) . . . . . . . . . . . . . . . . . . . . . . . . . . . . . . . . . . . . . . . . . . . . . . . . . . . . . . . . . . . . . . . . . . . . . . . . . . . . . . . . . . . . . . . . . . . . . . . 162}
\subsection{18.4    KYC/AML Décentralisé – La Conformité par la Technologie . . . . . . . . . . . . . . . . . . . . . . . . . . . . . . . . . . . . . . . . . . . . . . . . . . . . . . . . . . . . . . . . . . . . . . 162}
\subsection{18.4.1            Philosophie du ”Privacy by Design” . . . . . . . . . . . . . . . . . . . . . . . . . . . . . . . . . . . . . . . . . . . . . . . . . . . . . . . . . . . . . . . . . . . . . . . . . . . . . . . . . . . . . . 162}
\subsection{18.4.2            Architecture Technique . . . . . . . . . . . . . . . . . . . . . . . . . . . . . . . . . . . . . . . . . . . . . . . . . . . . . . . . . . . . . . . . . . . . . . . . . . . . . . . . . . . . . . . . . . . . . . . . . 162}
\subsection{18.4.3            Stack Pratique Enseignée . . . . . . . . . . . . . . . . . . . . . . . . . . . . . . . . . . . . . . . . . . . . . . . . . . . . . . . . . . . . . . . . . . . . . . . . . . . . . . . . . . . . . . . . . . . . . . . 163}
\vspace{0.5em}
\vspace{0.5em}
\vspace{0.5em}
\begin{confidential}CONFIDENTIEL — PROJET RBK 2.0                                                                           Page 8 sur 316\end{confidential}
RBK 2.0 — Whitepaper Architecte Web3\\
\vspace{0.5em}
\vspace{0.5em}
\subsection{18.5    GDPR \textbackslash{}& Données On-Chain . . . . . . . . . . . . . . . . . . . . . . . . . . . . . . . . . . . . . . . . . . . . . . . . . . . . . . . . . . . . . . . . . . . . . . . . . . . . . . . . . . . . . . . . . . . . . . . . . . . . . . 163}
\subsection{18.5.1             Le Conflit Immuabilité vs Droit à l’Oubli . . . . . . . . . . . . . . . . . . . . . . . . . . . . . . . . . . . . . . . . . . . . . . . . . . . . . . . . . . . . . . . . . . . . . . . . . . . . . . . . . 163}
\subsection{18.5.2             Patterns Architecturaux . . . . . . . . . . . . . . . . . . . . . . . . . . . . . . . . . . . . . . . . . . . . . . . . . . . . . . . . . . . . . . . . . . . . . . . . . . . . . . . . . . . . . . . . . . . . . . . . . 163}
\subsection{18.6    Fiscalité Crypto \textbackslash{}& Statut ETE . . . . . . . . . . . . . . . . . . . . . . . . . . . . . . . . . . . . . . . . . . . . . . . . . . . . . . . . . . . . . . . . . . . . . . . . . . . . . . . . . . . . . . . . . . . . . . . . . . . . 163}
\subsection{18.6.1             Le Guide de l’Ingénieur-Exportateur . . . . . . . . . . . . . . . . . . . . . . . . . . . . . . . . . . . . . . . . . . . . . . . . . . . . . . . . . . . . . . . . . . . . . . . . . . . . . . . . . . . . . 164}
\subsection{18.6.2             Flux Financier Recommandé . . . . . . . . . . . . . . . . . . . . . . . . . . . . . . . . . . . . . . . . . . . . . . . . . . . . . . . . . . . . . . . . . . . . . . . . . . . . . . . . . . . . . . . . . . . . 164}
\vspace{0.5em}
19    Gouvernance \textbackslash{}& transparence . . . . . . . . . . . . . . . . . . . . . . . . . . . . . . . . . . . . . . . . . . . . 165\\
\subsection{19.1    Comité Éthique \textbackslash{}& Pédagogique (CEP) . . . . . . . . . . . . . . . . . . . . . . . . . . . . . . . . . . . . . . . . . . . . . . . . . . . . . . . . . . . . . . . . . . . . . . . . . . . . . . . . . . . . . . . . . . . . 165}
\subsection{19.1.1             Composition (5 Membres) . . . . . . . . . . . . . . . . . . . . . . . . . . . . . . . . . . . . . . . . . . . . . . . . . . . . . . . . . . . . . . . . . . . . . . . . . . . . . . . . . . . . . . . . . . . . . . . 165}
\subsection{19.1.2             Mandat . . . . . . . . . . . . . . . . . . . . . . . . . . . . . . . . . . . . . . . . . . . . . . . . . . . . . . . . . . . . . . . . . . . . . . . . . . . . . . . . . . . . . . . . . . . . . . . . . . . . . . . . . . . . . . . 165}
\vspace{0.5em}
\vspace{0.5em}
\vspace{0.5em}
\vspace{0.5em}
EL\\
\subsection{19.2    Transparence Radicale (Open Metrics) . . . . . . . . . . . . . . . . . . . . . . . . . . . . . . . . . . . . . . . . . . . . . . . . . . . . . . . . . . . . . . . . . . . . . . . . . . . . . . . . . . . . . . . . . . . 166}
\vspace{0.5em}
\vspace{0.5em}
\vspace{0.5em}
\vspace{0.5em}
TI\\
\subsection{19.3    Charte de Déontologie . . . . . . . . . . . . . . . . . . . . . . . . . . . . . . . . . . . . . . . . . . . . . . . . . . . . . . . . . . . . . . . . . . . . . . . . . . . . . . . . . . . . . . . . . . . . . . . . . . . . . . . . . . . 166}
\vspace{0.5em}
\vspace{0.5em}
\vspace{0.5em}
\vspace{0.5em}
EN\\
\subsection{19.4    Structure Juridique et Rôles (Branding) . . . . . . . . . . . . . . . . . . . . . . . . . . . . . . . . . . . . . . . . . . . . . . . . . . . . . . . . . . . . . . . . . . . . . . . . . . . . . . . . . . . . . . . . . . 167}
\vspace{0.5em}
20    IMPACT SOCIAL \textbackslash{}& ALIGNEMENT ODD . . . . . . . . . . . . . . . . . . . . . . . . . . . . . . . . . . . . . . . 168\\
\vspace{0.5em}
\vspace{0.5em}
\vspace{0.5em}
\vspace{0.5em}
D\\
\subsection{20.1    Contribution aux Objectifs de Développement Durable (ONU) . . . . . . . . . . . . . . . . . . . . . . . . . . . . . . . . . . . . . . . . . . . . . . . . . . . . . . . . . . . . . . . . . . . . 168}
\vspace{0.5em}
\vspace{0.5em}
\vspace{0.5em}
\vspace{0.5em}
FI\\
\subsection{20.2    Indicateurs de Performance Sociale . . . . . . . . . . . . . . . . . . . . . . . . . . . . . . . . . . . . . . . . . . . . . . . . . . . . . . . . . . . . . . . . . . . . . . . . . . . . . . . . . . . . . . . . . . . . . . 168}
\subsection{20.2.1}
N\\
1. Inclusion des Femmes dans la Tech . . . . . . . . . . . . . . . . . . . . . . . . . . . . . . . . . . . . . . . . . . . . . . . . . . . . . . . . . . . . . . . . . . . . . . . . . . . . . . . . . . . . 168\\
CO\\
\subsection{20.2.2             2. Décentralisation Régionale . . . . . . . . . . . . . . . . . . . . . . . . . . . . . . . . . . . . . . . . . . . . . . . . . . . . . . . . . . . . . . . . . . . . . . . . . . . . . . . . . . . . . . . . . . . 169}
\subsection{20.2.3             3. Empreinte Carbone et Compensation . . . . . . . . . . . . . . . . . . . . . . . . . . . . . . . . . . . . . . . . . . . . . . . . . . . . . . . . . . . . . . . . . . . . . . . . . . . . . . . . . . 169}
\vspace{0.5em}
21    INFRASTRUCTURE SBT \textbackslash{}& CERTIFICATION . . . . . . . . . . . . . . . . . . . . . . . . . . . . . . . . . . . . 170\\
\subsection{21.1    Philosophie : ”Don’t Trust, Verify” . . . . . . . . . . . . . . . . . . . . . . . . . . . . . . . . . . . . . . . . . . . . . . . . . . . . . . . . . . . . . . . . . . . . . . . . . . . . . . . . . . . . . . . . . . . . . . . 170}
\subsection{21.2    Stack Technique SBT . . . . . . . . . . . . . . . . . . . . . . . . . . . . . . . . . . . . . . . . . . . . . . . . . . . . . . . . . . . . . . . . . . . . . . . . . . . . . . . . . . . . . . . . . . . . . . . . . . . . . . . . . . . . 170}
\subsection{21.2.1             Choix du Standard . . . . . . . . . . . . . . . . . . . . . . . . . . . . . . . . . . . . . . . . . . . . . . . . . . . . . . . . . . . . . . . . . . . . . . . . . . . . . . . . . . . . . . . . . . . . . . . . . . . . . 170}
\subsection{21.3    Cycle de Vie de la Certification . . . . . . . . . . . . . . . . . . . . . . . . . . . . . . . . . . . . . . . . . . . . . . . . . . . . . . . . . . . . . . . . . . . . . . . . . . . . . . . . . . . . . . . . . . . . . . . . . . 171}
\subsection{21.4    Conformité RGPD \textbackslash{}& Privacy . . . . . . . . . . . . . . . . . . . . . . . . . . . . . . . . . . . . . . . . . . . . . . . . . . . . . . . . . . . . . . . . . . . . . . . . . . . . . . . . . . . . . . . . . . . . . . . . . . . . . 172}
\subsection{21.5    Cas d’Usage : Le Recrutement Instantané . . . . . . . . . . . . . . . . . . . . . . . . . . . . . . . . . . . . . . . . . . . . . . . . . . . . . . . . . . . . . . . . . . . . . . . . . . . . . . . . . . . . . . . . . 172}
\vspace{0.5em}
\vspace{0.5em}
\vspace{0.5em}
\begin{confidential}CONFIDENTIEL — PROJET RBK 2.0                                                                                Page 9 sur 316\end{confidential}
RBK 2.0 — Whitepaper Architecte Web3\\
\vspace{0.5em}
\vspace{0.5em}
22    FEUILLE DE ROUTE 120 JOURS . . . . . . . . . . . . . . . . . . . . . . . . . . . . . . . . . . . . . . . . . . . 173\\
\subsection{22.1    Timeline des Opérations . . . . . . . . . . . . . . . . . . . . . . . . . . . . . . . . . . . . . . . . . . . . . . . . . . . . . . . . . . . . . . . . . . . . . . . . . . . . . . . . . . . . . . . . . . . . . . . . . . . . . . . . . 173}
\subsection{22.2    Jalons Clés \textbackslash{}& Actions . . . . . . . . . . . . . . . . . . . . . . . . . . . . . . . . . . . . . . . . . . . . . . . . . . . . . . . . . . . . . . . . . . . . . . . . . . . . . . . . . . . . . . . . . . . . . . . . . . . . . . . . . . . . 175}
\subsection{22.3    Diagramme de Gantt Macro . . . . . . . . . . . . . . . . . . . . . . . . . . . . . . . . . . . . . . . . . . . . . . . . . . . . . . . . . . . . . . . . . . . . . . . . . . . . . . . . . . . . . . . . . . . . . . . . . . . . . . 177}
\vspace{0.5em}
23    FEUILLE DE ROUTE : LE PLAN DE LANCEMENT (90 JOURS) . . . . . . . . . . . . . . . . . . . . . . . . . 179\\
\subsection{23.1    Mois 1 . . . . . . . . . . . . . . . . . . . . . . . . . . . . . . . . . . . . . . . . . . . . . . . . . . . . . . . . . . . . . . . . . . . . . . . . . . . . . . . . . . . . . . . . . . . . . . . . . . . . . . . . . . . . . . . . . . . . . . . . . . . 179}
\subsection{23.1.1               Validation \textbackslash{}& Cadrage Stratégique . . . . . . . . . . . . . . . . . . . . . . . . . . . . . . . . . . . . . . . . . . . . . . . . . . . . . . . . . . . . . . . . . . . . . . . . . . . . . . . . . . . . . . . . 179}
\subsection{23.1.2               Constitution de l’Alliance Écosystémique . . . . . . . . . . . . . . . . . . . . . . . . . . . . . . . . . . . . . . . . . . . . . . . . . . . . . . . . . . . . . . . . . . . . . . . . . . . . . . . . . 179}
\subsection{23.1.3               Recrutement de l’Équipe Pilote . . . . . . . . . . . . . . . . . . . . . . . . . . . . . . . . . . . . . . . . . . . . . . . . . . . . . . . . . . . . . . . . . . . . . . . . . . . . . . . . . . . . . . . . . . 180}
\vspace{0.5em}
\vspace{0.5em}
\vspace{0.5em}
\vspace{0.5em}
EL\\
\subsection{23.2    Mois 2 . . . . . . . . . . . . . . . . . . . . . . . . . . . . . . . . . . . . . . . . . . . . . . . . . . . . . . . . . . . . . . . . . . . . . . . . . . . . . . . . . . . . . . . . . . . . . . . . . . . . . . . . . . . . . . . . . . . . . . . . . . . 180}
\subsection{23.2.1               Ingénierie Pédagogique (Les « Golden Templates ») . . . . . . . . . . . . . . . . . . . . . . . . . . . . . . . . . . . . . . . . . . . . . . . . . . . . . . . . . . . . . . . . . . . . . . . . 180}
\vspace{0.5em}
\vspace{0.5em}
\vspace{0.5em}
\vspace{0.5em}
TI\\
\subsection{23.2.2               Mise en place du Cockpit Technique . . . . . . . . . . . . . . . . . . . . . . . . . . . . . . . . . . . . . . . . . . . . . . . . . . . . . . . . . . . . . . . . . . . . . . . . . . . . . . . . . . . . . 180}
\vspace{0.5em}
\vspace{0.5em}
\vspace{0.5em}
\vspace{0.5em}
EN\\
\subsection{23.2.3               Lancement Commercial \textbackslash{}& Marketing . . . . . . . . . . . . . . . . . . . . . . . . . . . . . . . . . . . . . . . . . . . . . . . . . . . . . . . . . . . . . . . . . . . . . . . . . . . . . . . . . . . . . 180}
\subsection{23.3    Mois 3 . . . . . . . . . . . . . . . . . . . . . . . . . . . . . . . . . . . . . . . . . . . . . . . . . . . . . . . . . . . . . . . . . . . . . . . . . . . . . . . . . . . . . . . . . . . . . . . . . . . . . . . . . . . . . . . . . . . . . . . . . . . 181}
\vspace{0.5em}
\vspace{0.5em}
\vspace{0.5em}
\vspace{0.5em}
D\\
\subsection{23.3.1               Processus de Sélection d’Élite . . . . . . . . . . . . . . . . . . . . . . . . . . . . . . . . . . . . . . . . . . . . . . . . . . . . . . . . . . . . . . . . . . . . . . . . . . . . . . . . . . . . . . . . . . . 181}
\vspace{0.5em}
\vspace{0.5em}
\vspace{0.5em}
\vspace{0.5em}
FI\\
\subsection{23.3.2               Finalisation de la Cohorte . . . . . . . . . . . . . . . . . . . . . . . . . . . . . . . . . . . . . . . . . . . . . . . . . . . . . . . . . . . . . . . . . . . . . . . . . . . . . . . . . . . . . . . . . . . . . . . 181}
\subsection{23.3.3               Kick-off Opérationnel . . . . . . . . . . . . . . . . . . . . . . . . . . . . . . . . . . . . . . . . . . . . . . . . . . . . . . . . . . . . . . . . . . . . . . . . . . . . . . . . . . . . . . . . . . . . . . . . . . . 181}
\subsection{23.4                                                                            N}
Jalons clés . . . . . . . . . . . . . . . . . . . . . . . . . . . . . . . . . . . . . . . . . . . . . . . . . . . . . . . . . . . . . . . . . . . . . . . . . . . . . . . . . . . . . . . . . . . . . . . . . . . . . . . . . . . . . . . . . . . . . . . 182\\
CO\\
24    TOKEN DE RÉPUTATION \textbackslash{}& ALUMNI PROGRAM . . . . . . . . . . . . . . . . . . . . . . . . . . . . . . . . . 183\\
\subsection{24.1    RBK Soulbound Tokens (SBTs) . . . . . . . . . . . . . . . . . . . . . . . . . . . . . . . . . . . . . . . . . . . . . . . . . . . . . . . . . . . . . . . . . . . . . . . . . . . . . . . . . . . . . . . . . . . . . . . . . . . 183}
\subsection{24.2    Usages des SBT . . . . . . . . . . . . . . . . . . . . . . . . . . . . . . . . . . . . . . . . . . . . . . . . . . . . . . . . . . . . . . . . . . . . . . . . . . . . . . . . . . . . . . . . . . . . . . . . . . . . . . . . . . . . . . . . . . 185}
\subsection{24.3    Alumni Program Structuré . . . . . . . . . . . . . . . . . . . . . . . . . . . . . . . . . . . . . . . . . . . . . . . . . . . . . . . . . . . . . . . . . . . . . . . . . . . . . . . . . . . . . . . . . . . . . . . . . . . . . . . 186}
\vspace{0.5em}
25    ÉLÉMENTS DE DIFFÉRENCIATION . . . . . . . . . . . . . . . . . . . . . . . . . . . . . . . . . . . . . . . . . 188\\
\subsection{25.1    Le Paradigme « Senior-by-Design » . . . . . . . . . . . . . . . . . . . . . . . . . . . . . . . . . . . . . . . . . . . . . . . . . . . . . . . . . . . . . . . . . . . . . . . . . . . . . . . . . . . . . . . . . . . . . . . 188}
\subsection{25.2    Approche « Cyborg » : IA-Augmented Engineering . . . . . . . . . . . . . . . . . . . . . . . . . . . . . . . . . . . . . . . . . . . . . . . . . . . . . . . . . . . . . . . . . . . . . . . . . . . . . . . . 189}
\subsection{25.3    Dual Track Solana/EVM : Flexibilité Stratégique . . . . . . . . . . . . . . . . . . . . . . . . . . . . . . . . . . . . . . . . . . . . . . . . . . . . . . . . . . . . . . . . . . . . . . . . . . . . . . . . . 190}
\vspace{0.5em}
\vspace{0.5em}
\begin{confidential}CONFIDENTIEL — PROJET RBK 2.0                                                                                  Page 10 sur 316\end{confidential}
RBK 2.0 — Whitepaper Architecte Web3\\
\vspace{0.5em}
\vspace{0.5em}
\subsection{25.4    Intégration Superteam : Opportunités Directes . . . . . . . . . . . . . . . . . . . . . . . . . . . . . . . . . . . . . . . . . . . . . . . . . . . . . . . . . . . . . . . . . . . . . . . . . . . . . . . . . . . 190}
\subsection{25.5    « On-Chain Resume » : Preuve de Travail Public . . . . . . . . . . . . . . . . . . . . . . . . . . . . . . . . . . . . . . . . . . . . . . . . . . . . . . . . . . . . . . . . . . . . . . . . . . . . . . . . . . 191}
\subsection{25.6    Ancrage Tunisie + Export : Software Factory Future . . . . . . . . . . . . . . . . . . . . . . . . . . . . . . . . . . . . . . . . . . . . . . . . . . . . . . . . . . . . . . . . . . . . . . . . . . . . . 191}
\subsection{25.6.1              Comparatif RBK 2.0 vs Bootcamps Classiques . . . . . . . . . . . . . . . . . . . . . . . . . . . . . . . . . . . . . . . . . . . . . . . . . . . . . . . . . . . . . . . . . . . . . . . . . . . . . 192}
\vspace{0.5em}
26    CONCLUSION \textbackslash{}& FEUILLE DE ROUTE . . . . . . . . . . . . . . . . . . . . . . . . . . . . . . . . . . . . . . . . 193\\
\subsection{26.1    Priorités Immédiates (Semaine 1–4) . . . . . . . . . . . . . . . . . . . . . . . . . . . . . . . . . . . . . . . . . . . . . . . . . . . . . . . . . . . . . . . . . . . . . . . . . . . . . . . . . . . . . . . . . . . . . . 193}
\subsection{26.2    KPI de Succès . . . . . . . . . . . . . . . . . . . . . . . . . . . . . . . . . . . . . . . . . . . . . . . . . . . . . . . . . . . . . . . . . . . . . . . . . . . . . . . . . . . . . . . . . . . . . . . . . . . . . . . . . . . . . . . . . . . . 193}
\subsection{26.3    Engagement Qualité Formel . . . . . . . . . . . . . . . . . . . . . . . . . . . . . . . . . . . . . . . . . . . . . . . . . . . . . . . . . . . . . . . . . . . . . . . . . . . . . . . . . . . . . . . . . . . . . . . . . . . . . 194}
\subsection{26.4    Forge de l’Élite Africaine . . . . . . . . . . . . . . . . . . . . . . . . . . . . . . . . . . . . . . . . . . . . . . . . . . . . . . . . . . . . . . . . . . . . . . . . . . . . . . . . . . . . . . . . . . . . . . . . . . . . . . . . . 194}
\vspace{0.5em}
\vspace{0.5em}
\vspace{0.5em}
\vspace{0.5em}
EL\\
\subsection{26.5    Synthèse Valeur Stratégique . . . . . . . . . . . . . . . . . . . . . . . . . . . . . . . . . . . . . . . . . . . . . . . . . . . . . . . . . . . . . . . . . . . . . . . . . . . . . . . . . . . . . . . . . . . . . . . . . . . . . 194}
\subsection{26.6    Synthèse Investisseur – Différenciateurs Nexus . . . . . . . . . . . . . . . . . . . . . . . . . . . . . . . . . . . . . . . . . . . . . . . . . . . . . . . . . . . . . . . . . . . . . . . . . . . . . . . . . . . 195}
\vspace{0.5em}
\vspace{0.5em}
\vspace{0.5em}
\vspace{0.5em}
TI\\
\subsection{26.7    Appel à l’Action . . . . . . . . . . . . . . . . . . . . . . . . . . . . . . . . . . . . . . . . . . . . . . . . . . . . . . . . . . . . . . . . . . . . . . . . . . . . . . . . . . . . . . . . . . . . . . . . . . . . . . . . . . . . . . . . . 195}
\vspace{0.5em}
\vspace{0.5em}
\vspace{0.5em}
\vspace{0.5em}
EN\\
\subsection{26.8    Message Final au CEO . . . . . . . . . . . . . . . . . . . . . . . . . . . . . . . . . . . . . . . . . . . . . . . . . . . . . . . . . . . . . . . . . . . . . . . . . . . . . . . . . . . . . . . . . . . . . . . . . . . . . . . . . . . 196}
\subsection{26.9    Profil de Sortie . . . . . . . . . . . . . . . . . . . . . . . . . . . . . . . . . . . . . . . . . . . . . . . . . . . . . . . . . . . . . . . . . . . . . . . . . . . . . . . . . . . . . . . . . . . . . . . . . . . . . . . . . . . . . . . . . . 196}
\vspace{0.5em}
\vspace{0.5em}
\vspace{0.5em}
\vspace{0.5em}
D\\
FI\\
Liste des figures . . . . . . . . . . . . . . . . . . . . . . . . . . . . . . . . . . . . . . . . . . . . . . . . . . . . . . . . 197\\
\vspace{0.5em}
N\\
Liste des tableaux . . . . . . . . . . . . . . . . . . . . . . . . . . . . . . . . . . . . . . . . . . . . . . . . . . . . . . . 199\\
CO\\
ANNEXES                                                                                                                                                                                                                                                      203\\
A     Gabarits opérationnels \textbackslash{}& stratégiques . . . . . . . . . . . . . . . . . . . . . . . . . . . . . . . . . . . . . . . 203\\
A.1     Matrice SWOT (Forces, Faiblesses, Opportunités, Menaces) . . . . . . . . . . . . . . . . . . . . . . . . . . . . . . . . . . . . . . . . . . . . . . . . . . . . . . . . . . . . . . . . . . . . . . . 204\\
A.2     Priorisation MoSCoW (Must, Should,\\
Could, Won’t) . . . . . . . . . . . . . . . . . . . . . . . . . . . . . . . . . . . . . . . . . . . . . . . . . . . . . . . . . . . . . . . . . . . . . . . . . . . . . . . . . . . . . . . . . . . . . . . . . . . . . . . . . . . . . . . . . . . . 205\\
A.3     Matrice RACI (Responsable, Accountable,\\
Consulted, Informed) . . . . . . . . . . . . . . . . . . . . . . . . . . . . . . . . . . . . . . . . . . . . . . . . . . . . . . . . . . . . . . . . . . . . . . . . . . . . . . . . . . . . . . . . . . . . . . . . . . . . . . . . . . . . 207\\
\vspace{0.5em}
\vspace{0.5em}
\vspace{0.5em}
\vspace{0.5em}
\begin{confidential}CONFIDENTIEL — PROJET RBK 2.0                                                                               Page 11 sur 316\end{confidential}
RBK 2.0 — Whitepaper Architecte Web3\\
\vspace{0.5em}
\vspace{0.5em}
A.4     Gabarits Studio (Ops \textbackslash{}& Technique) . . . . . . . . . . . . . . . . . . . . . . . . . . . . . . . . . . . . . . . . . . . . . . . . . . . . . . . . . . . . . . . . . . . . . . . . . . . . . . . . . . . . . . . . . . . . . . . 209\\
A.4.1              Template : Incident Drill Postmortem . . . . . . . . . . . . . . . . . . . . . . . . . . . . . . . . . . . . . . . . . . . . . . . . . . . . . . . . . . . . . . . . . . . . . . . . . . . . . . . . . . . . 209\\
A.5     Références Standards (Voir Annexe A) . . . . . . . . . . . . . . . . . . . . . . . . . . . . . . . . . . . . . . . . . . . . . . . . . . . . . . . . . . . . . . . . . . . . . . . . . . . . . . . . . . . . . . . . . . . 209\\
\vspace{0.5em}
B     SYLLABUS TECHNIQUE\\
DÉTAILLÉ ET STANDARDS . . . . . . . . . . . . . . . . . . . . . . . . . . . . . . . . . . . . . . . . . . . . . . 210\\
B.1     Gabarit de Dépôt Obligatoire . . . . . . . . . . . . . . . . . . . . . . . . . . . . . . . . . . . . . . . . . . . . . . . . . . . . . . . . . . . . . . . . . . . . . . . . . . . . . . . . . . . . . . . . . . . . . . . . . . . . 210\\
B.1.1              Exigences de Reproductibilité (Scripts Obligatoires) . . . . . . . . . . . . . . . . . . . . . . . . . . . . . . . . . . . . . . . . . . . . . . . . . . . . . . . . . . . . . . . . . . . . . . . 212\\
B.2     Rubrique d’Évaluation \textbackslash{}& Processus . . . . . . . . . . . . . . . . . . . . . . . . . . . . . . . . . . . . . . . . . . . . . . . . . . . . . . . . . . . . . . . . . . . . . . . . . . . . . . . . . . . . . . . . . . . . . . . 213\\
B.2.1              Revues Hebdomadaires (Sprints) . . . . . . . . . . . . . . . . . . . . . . . . . . . . . . . . . . . . . . . . . . . . . . . . . . . . . . . . . . . . . . . . . . . . . . . . . . . . . . . . . . . . . . . . 213\\
Format de Jury (Gate) . . . . . . . . . . . . . . . . . . . . . . . . . . . . . . . . . . . . . . . . . . . . . . . . . . . . . . . . . . . . . . . . . . . . . . . . . . . . . . . . . . . . . . . . . . . . . . . . . . 213\\
\vspace{0.5em}
\vspace{0.5em}
\vspace{0.5em}
\vspace{0.5em}
EL\\
B.2.2\\
B.3     Standards de Documentation et Templates . . . . . . . . . . . . . . . . . . . . . . . . . . . . . . . . . . . . . . . . . . . . . . . . . . . . . . . . . . . . . . . . . . . . . . . . . . . . . . . . . . . . . . . 213\\
\vspace{0.5em}
\vspace{0.5em}
\vspace{0.5em}
\vspace{0.5em}
TI\\
B.3.1              Règles ADR (Architecture Decision Records) . . . . . . . . . . . . . . . . . . . . . . . . . . . . . . . . . . . . . . . . . . . . . . . . . . . . . . . . . . . . . . . . . . . . . . . . . . . . . . 213\\
Structure PROOFS.md (Journal des Preuves) . . . . . . . . . . . . . . . . . . . . . . . . . . . . . . . . . . . . . . . . . . . . . . . . . . . . . . . . . . . . . . . . . . . . . . . . . . . . . . 214\\
\vspace{0.5em}
\vspace{0.5em}
\vspace{0.5em}
\vspace{0.5em}
EN\\
B.3.2\\
B.3.3              Template : PROOFS.md . . . . . . . . . . . . . . . . . . . . . . . . . . . . . . . . . . . . . . . . . . . . . . . . . . . . . . . . . . . . . . . . . . . . . . . . . . . . . . . . . . . . . . . . . . . . . . . . . 214\\
Template : THREAT\textbackslash{}_MODEL.md . . . . . . . . . . . . . . . . . . . . . . . . . . . . . . . . . . . . . . . . . . . . . . . . . . . . . . . . . . . . . . . . . . . . . . . . . . . . . . . . . . . . . . . . . 215\\
\vspace{0.5em}
\vspace{0.5em}
\vspace{0.5em}
\vspace{0.5em}
D\\
B.3.4\\
B.3.5              Template : adr/NNNN-template.md . . . . . . . . . . . . . . . . . . . . . . . . . . . . . . . . . . . . . . . . . . . . . . . . . . . . . . . . . . . . . . . . . . . . . . . . . . . . . . . . . . . . . . 216\\
\vspace{0.5em}
\vspace{0.5em}
\vspace{0.5em}
\vspace{0.5em}
FI\\
B.4     Grille de Notation Standardisée (/100) et Seuils . . . . . . . . . . . . . . . . . . . . . . . . . . . . . . . . . . . . . . . . . . . . . . . . . . . . . . . . . . . . . . . . . . . . . . . . . . . . . . . . . . 217\\
B.4.1\\
N\\
Détail de la Notation (/100) . . . . . . . . . . . . . . . . . . . . . . . . . . . . . . . . . . . . . . . . . . . . . . . . . . . . . . . . . . . . . . . . . . . . . . . . . . . . . . . . . . . . . . . . . . . . . 217\\
CO\\
B.4.2              Seuils de Décision (GO / WARN / NO-GO) . . . . . . . . . . . . . . . . . . . . . . . . . . . . . . . . . . . . . . . . . . . . . . . . . . . . . . . . . . . . . . . . . . . . . . . . . . . . . . . . 218\\
B.5     Format de Jury Gate et Procès-Verbal . . . . . . . . . . . . . . . . . . . . . . . . . . . . . . . . . . . . . . . . . . . . . . . . . . . . . . . . . . . . . . . . . . . . . . . . . . . . . . . . . . . . . . . . . . . . 219\\
B.5.1              Déroulé Standard (45–60 min) . . . . . . . . . . . . . . . . . . . . . . . . . . . . . . . . . . . . . . . . . . . . . . . . . . . . . . . . . . . . . . . . . . . . . . . . . . . . . . . . . . . . . . . . . . 219\\
B.5.2              Checklist de Questions Fixes . . . . . . . . . . . . . . . . . . . . . . . . . . . . . . . . . . . . . . . . . . . . . . . . . . . . . . . . . . . . . . . . . . . . . . . . . . . . . . . . . . . . . . . . . . . . 220\\
B.5.3              Barème de Décision (GO / WARN / NO-GO) . . . . . . . . . . . . . . . . . . . . . . . . . . . . . . . . . . . . . . . . . . . . . . . . . . . . . . . . . . . . . . . . . . . . . . . . . . . . . . 221\\
B.5.4              Modèle de Procès-Verbal de Jury . . . . . . . . . . . . . . . . . . . . . . . . . . . . . . . . . . . . . . . . . . . . . . . . . . . . . . . . . . . . . . . . . . . . . . . . . . . . . . . . . . . . . . . . 222\\
B.6     Contrat de Livrable (1 page) . . . . . . . . . . . . . . . . . . . . . . . . . . . . . . . . . . . . . . . . . . . . . . . . . . . . . . . . . . . . . . . . . . . . . . . . . . . . . . . . . . . . . . . . . . . . . . . . . . . . . 224\\
B.7     Syllabus Opérationnel (48 lignes / 48 semaines) . . . . . . . . . . . . . . . . . . . . . . . . . . . . . . . . . . . . . . . . . . . . . . . . . . . . . . . . . . . . . . . . . . . . . . . . . . . . . . . . . 225\\
\vspace{0.5em}
C     Annexe B — Modèle financier . . . . . . . . . . . . . . . . . . . . . . . . . . . . . . . . . . . . . . . . . . . . 229\\
\vspace{0.5em}
\vspace{0.5em}
\begin{confidential}CONFIDENTIEL — PROJET RBK 2.0                                                                           Page 12 sur 316\end{confidential}
RBK 2.0 — Whitepaper Architecte Web3\\
\vspace{0.5em}
\vspace{0.5em}
C.1     Référentiel financier contractuel . . . . . . . . . . . . . . . . . . . . . . . . . . . . . . . . . . . . . . . . . . . . . . . . . . . . . . . . . . . . . . . . . . . . . . . . . . . . . . . . . . . . . . . . . . . . . . . . . 229\\
C.1.1               Définitions comptables . . . . . . . . . . . . . . . . . . . . . . . . . . . . . . . . . . . . . . . . . . . . . . . . . . . . . . . . . . . . . . . . . . . . . . . . . . . . . . . . . . . . . . . . . . . . . . . . . 229\\
C.1.2               Synthèse des montants contractuels (TND) . . . . . . . . . . . . . . . . . . . . . . . . . . . . . . . . . . . . . . . . . . . . . . . . . . . . . . . . . . . . . . . . . . . . . . . . . . . . . . . . 230\\
C.2     Cohorte 1 — Paramètres imposés + Minimum garanti . . . . . . . . . . . . . . . . . . . . . . . . . . . . . . . . . . . . . . . . . . . . . . . . . . . . . . . . . . . . . . . . . . . . . . . . . . . 230\\
C.3     Rémunération de Nexus — Structure à 4 composantes . . . . . . . . . . . . . . . . . . . . . . . . . . . . . . . . . . . . . . . . . . . . . . . . . . . . . . . . . . . . . . . . . . . . . . . . . . . 230\\
C.4     Waterfall (ordre de distribution) — Clause contractuelle . . . . . . . . . . . . . . . . . . . . . . . . . . . . . . . . . . . . . . . . . . . . . . . . . . . . . . . . . . . . . . . . . . . . . . . . . 231\\
C.5     Table de projections (OBLIGATOIRE) — 3 scénarios fixes . . . . . . . . . . . . . . . . . . . . . . . . . . . . . . . . . . . . . . . . . . . . . . . . . . . . . . . . . . . . . . . . . . . . . . . . 233\\
C.6     Clauses comptables \textbackslash{}& fiscales (sans chiffres hasardeux) . . . . . . . . . . . . . . . . . . . . . . . . . . . . . . . . . . . . . . . . . . . . . . . . . . . . . . . . . . . . . . . . . . . . . . . . . . 234\\
C.7     Garanties \textbackslash{}& protections contractuelles . . . . . . . . . . . . . . . . . . . . . . . . . . . . . . . . . . . . . . . . . . . . . . . . . . . . . . . . . . . . . . . . . . . . . . . . . . . . . . . . . . . . . . . . . . . 234\\
\vspace{0.5em}
\vspace{0.5em}
\vspace{0.5em}
\vspace{0.5em}
EL\\
D     Addendum juridique . . . . . . . . . . . . . . . . . . . . . . . . . . . . . . . . . . . . . . . . . . . . . . . . . . 235\\
D.1     Périmètre contractuel RBK–Nexus . . . . . . . . . . . . . . . . . . . . . . . . . . . . . . . . . . . . . . . . . . . . . . . . . . . . . . . . . . . . . . . . . . . . . . . . . . . . . . . . . . . . . . . . . . . . . . . . 235\\
\vspace{0.5em}
\vspace{0.5em}
\vspace{0.5em}
\vspace{0.5em}
TI\\
D.1.1               Livrables contractuels par phase . . . . . . . . . . . . . . . . . . . . . . . . . . . . . . . . . . . . . . . . . . . . . . . . . . . . . . . . . . . . . . . . . . . . . . . . . . . . . . . . . . . . . . . . . 235\\
\vspace{0.5em}
\vspace{0.5em}
\vspace{0.5em}
\vspace{0.5em}
EN\\
D.1.2               Domaines transverses . . . . . . . . . . . . . . . . . . . . . . . . . . . . . . . . . . . . . . . . . . . . . . . . . . . . . . . . . . . . . . . . . . . . . . . . . . . . . . . . . . . . . . . . . . . . . . . . . . . 236\\
D.2     KPIs contractuels et bonus/malus . . . . . . . . . . . . . . . . . . . . . . . . . . . . . . . . . . . . . . . . . . . . . . . . . . . . . . . . . . . . . . . . . . . . . . . . . . . . . . . . . . . . . . . . . . . . . . . . 237\\
\vspace{0.5em}
\vspace{0.5em}
\vspace{0.5em}
\vspace{0.5em}
D\\
D.3     Clauses de propriété intellectuelle \textbackslash{}& licences . . . . . . . . . . . . . . . . . . . . . . . . . . . . . . . . . . . . . . . . . . . . . . . . . . . . . . . . . . . . . . . . . . . . . . . . . . . . . . . . . . . . . 238\\
\vspace{0.5em}
\vspace{0.5em}
\vspace{0.5em}
\vspace{0.5em}
FI\\
D.4     Processus de pilotage et résolution des litiges . . . . . . . . . . . . . . . . . . . . . . . . . . . . . . . . . . . . . . . . . . . . . . . . . . . . . . . . . . . . . . . . . . . . . . . . . . . . . . . . . . . . 238\\
D.5\\
N\\
Référentiel financier . . . . . . . . . . . . . . . . . . . . . . . . . . . . . . . . . . . . . . . . . . . . . . . . . . . . . . . . . . . . . . . . . . . . . . . . . . . . . . . . . . . . . . . . . . . . . . . . . . . . . . . . . . . . . 238\\
CO\\
D.6     Modalités de paiement . . . . . . . . . . . . . . . . . . . . . . . . . . . . . . . . . . . . . . . . . . . . . . . . . . . . . . . . . . . . . . . . . . . . . . . . . . . . . . . . . . . . . . . . . . . . . . . . . . . . . . . . . . . 238\\
D.7     Clause de ”Gross-Up” . . . . . . . . . . . . . . . . . . . . . . . . . . . . . . . . . . . . . . . . . . . . . . . . . . . . . . . . . . . . . . . . . . . . . . . . . . . . . . . . . . . . . . . . . . . . . . . . . . . . . . . . . . . . 239\\
D.8     Droit d’audit . . . . . . . . . . . . . . . . . . . . . . . . . . . . . . . . . . . . . . . . . . . . . . . . . . . . . . . . . . . . . . . . . . . . . . . . . . . . . . . . . . . . . . . . . . . . . . . . . . . . . . . . . . . . . . . . . . . . . 239\\
D.9     Responsabilités contractuelles . . . . . . . . . . . . . . . . . . . . . . . . . . . . . . . . . . . . . . . . . . . . . . . . . . . . . . . . . . . . . . . . . . . . . . . . . . . . . . . . . . . . . . . . . . . . . . . . . . . 239\\
D.10    Force majeure . . . . . . . . . . . . . . . . . . . . . . . . . . . . . . . . . . . . . . . . . . . . . . . . . . . . . . . . . . . . . . . . . . . . . . . . . . . . . . . . . . . . . . . . . . . . . . . . . . . . . . . . . . . . . . . . . . . 239\\
D.11    Résiliation anticipée . . . . . . . . . . . . . . . . . . . . . . . . . . . . . . . . . . . . . . . . . . . . . . . . . . . . . . . . . . . . . . . . . . . . . . . . . . . . . . . . . . . . . . . . . . . . . . . . . . . . . . . . . . . . . 239\\
\vspace{0.5em}
E     TEMPLATE DE RAPPORT D’AUDIT DE SÉCURITÉ . . . . . . . . . . . . . . . . . . . . . . . . . . . . . . . . 241\\
E.1     Structure du Rapport . . . . . . . . . . . . . . . . . . . . . . . . . . . . . . . . . . . . . . . . . . . . . . . . . . . . . . . . . . . . . . . . . . . . . . . . . . . . . . . . . . . . . . . . . . . . . . . . . . . . . . . . . . . . 241\\
\vspace{0.5em}
\vspace{0.5em}
\vspace{0.5em}
\begin{confidential}CONFIDENTIEL — PROJET RBK 2.0                                                                                Page 13 sur 316\end{confidential}
RBK 2.0 — Whitepaper Architecte Web3\\
\vspace{0.5em}
\vspace{0.5em}
E.2     Classification des Risques . . . . . . . . . . . . . . . . . . . . . . . . . . . . . . . . . . . . . . . . . . . . . . . . . . . . . . . . . . . . . . . . . . . . . . . . . . . . . . . . . . . . . . . . . . . . . . . . . . . . . . . . 241\\
E.3     Fiche Finding Type . . . . . . . . . . . . . . . . . . . . . . . . . . . . . . . . . . . . . . . . . . . . . . . . . . . . . . . . . . . . . . . . . . . . . . . . . . . . . . . . . . . . . . . . . . . . . . . . . . . . . . . . . . . . . . 242\\
\vspace{0.5em}
F     LE COCKPIT DE L’ARCHITECTE . . . . . . . . . . . . . . . . . . . . . . . . . . . . . . . . . . . . . . . . . . . 243\\
F.1     Stack Outillage Minimal . . . . . . . . . . . . . . . . . . . . . . . . . . . . . . . . . . . . . . . . . . . . . . . . . . . . . . . . . . . . . . . . . . . . . . . . . . . . . . . . . . . . . . . . . . . . . . . . . . . . . . . . . 243\\
F.2     Journée Type (Productivité) . . . . . . . . . . . . . . . . . . . . . . . . . . . . . . . . . . . . . . . . . . . . . . . . . . . . . . . . . . . . . . . . . . . . . . . . . . . . . . . . . . . . . . . . . . . . . . . . . . . . . 243\\
\vspace{0.5em}
G     Modèle ISA (Income Share Agreement) . . . . . . . . . . . . . . . . . . . . . . . . . . . . . . . . . . . . . . 244\\
G.1     Objet et Principes . . . . . . . . . . . . . . . . . . . . . . . . . . . . . . . . . . . . . . . . . . . . . . . . . . . . . . . . . . . . . . . . . . . . . . . . . . . . . . . . . . . . . . . . . . . . . . . . . . . . . . . . . . . . . . . . 244\\
G.2     Éligibilité (Gating) . . . . . . . . . . . . . . . . . . . . . . . . . . . . . . . . . . . . . . . . . . . . . . . . . . . . . . . . . . . . . . . . . . . . . . . . . . . . . . . . . . . . . . . . . . . . . . . . . . . . . . . . . . . . . . . 244\\
G.3     Paramètres ISA unifiés . . . . . . . . . . . . . . . . . . . . . . . . . . . . . . . . . . . . . . . . . . . . . . . . . . . . . . . . . . . . . . . . . . . . . . . . . . . . . . . . . . . . . . . . . . . . . . . . . . . . . . . . . . . 244\\
\vspace{0.5em}
\vspace{0.5em}
\vspace{0.5em}
\vspace{0.5em}
EL\\
G.4     Règles de Pause, Chômage, Variabilité . . . . . . . . . . . . . . . . . . . . . . . . . . . . . . . . . . . . . . . . . . . . . . . . . . . . . . . . . . . . . . . . . . . . . . . . . . . . . . . . . . . . . . . . . . . 245\\
\vspace{0.5em}
\vspace{0.5em}
\vspace{0.5em}
\vspace{0.5em}
TI\\
G.5     Cas Limites (Edge Cases) . . . . . . . . . . . . . . . . . . . . . . . . . . . . . . . . . . . . . . . . . . . . . . . . . . . . . . . . . . . . . . . . . . . . . . . . . . . . . . . . . . . . . . . . . . . . . . . . . . . . . . . . . 245\\
\vspace{0.5em}
\vspace{0.5em}
\vspace{0.5em}
\vspace{0.5em}
EN\\
G.6     Conformité Éthique (Musharaka) . . . . . . . . . . . . . . . . . . . . . . . . . . . . . . . . . . . . . . . . . . . . . . . . . . . . . . . . . . . . . . . . . . . . . . . . . . . . . . . . . . . . . . . . . . . . . . . . . 245\\
G.7     Exemples Chiffrés (Seuil 2500, Taux 15\textbackslash{}%) . . . . . . . . . . . . . . . . . . . . . . . . . . . . . . . . . . . . . . . . . . . . . . . . . . . . . . . . . . . . . . . . . . . . . . . . . . . . . . . . . . . . . . . 246\\
\vspace{0.5em}
\vspace{0.5em}
\vspace{0.5em}
\vspace{0.5em}
D\\
H     GUIDE DE SÉLECTION \textbackslash{}& SCORING « PISCINE RUST » . . . . . . . . . . . . . . . . . . . . . . . . . . . . . . 247\\
\vspace{0.5em}
\vspace{0.5em}
\vspace{0.5em}
FI\\
H.1     Grille de Scoring . . . . . . . . . . . . . . . . . . . . . . . . . . . . . . . . . . . . . . . . . . . . . . . . . . . . . . . . . . . . . . . . . . . . . . . . . . . . . . . . . . . . . . . . . . . . . . . . . . . . . . . . . . . . . . . . . 247\\
H.1.1\\
N\\
Test d’Entrée Explorer (Bypass Genesis Pool) . . . . . . . . . . . . . . . . . . . . . . . . . . . . . . . . . . . . . . . . . . . . . . . . . . . . . . . . . . . . . . . . . . . . . . . . . . . . . . 247\\
CO\\
H.1.2               Test d’Entrée Bâtisseur/Architecte (Bypass Explorer) . . . . . . . . . . . . . . . . . . . . . . . . . . . . . . . . . . . . . . . . . . . . . . . . . . . . . . . . . . . . . . . . . . . . . . 248\\
\vspace{0.5em}
I     SPÉCIFICATIONS TECHNIQUES SBT . . . . . . . . . . . . . . . . . . . . . . . . . . . . . . . . . . . . . . . . 249\\
I.1     Schéma de Métadonnées (JSON) . . . . . . . . . . . . . . . . . . . . . . . . . . . . . . . . . . . . . . . . . . . . . . . . . . . . . . . . . . . . . . . . . . . . . . . . . . . . . . . . . . . . . . . . . . . . . . . . . 249\\
I.2     Processus de Vérification . . . . . . . . . . . . . . . . . . . . . . . . . . . . . . . . . . . . . . . . . . . . . . . . . . . . . . . . . . . . . . . . . . . . . . . . . . . . . . . . . . . . . . . . . . . . . . . . . . . . . . . . 250\\
\vspace{0.5em}
J     DASHBOARD DE SUIVI PROMO . . . . . . . . . . . . . . . . . . . . . . . . . . . . . . . . . . . . . . . . . . . 251\\
J.1     Indicateurs Hebdomadaires (KPI) . . . . . . . . . . . . . . . . . . . . . . . . . . . . . . . . . . . . . . . . . . . . . . . . . . . . . . . . . . . . . . . . . . . . . . . . . . . . . . . . . . . . . . . . . . . . . . . . 251\\
J.2     Questionnaire Bien-être Minimal . . . . . . . . . . . . . . . . . . . . . . . . . . . . . . . . . . . . . . . . . . . . . . . . . . . . . . . . . . . . . . . . . . . . . . . . . . . . . . . . . . . . . . . . . . . . . . . . . 251\\
\vspace{0.5em}
K     OFFRE COMMERCIALE \textbackslash{}& MODALITÉS . . . . . . . . . . . . . . . . . . . . . . . . . . . . . . . . . . . . . . . 252\\
\vspace{0.5em}
\vspace{0.5em}
\begin{confidential}CONFIDENTIEL — PROJET RBK 2.0                                                                              Page 14 sur 316\end{confidential}
RBK 2.0 — Whitepaper Architecte Web3\\
\vspace{0.5em}
\vspace{0.5em}
K.1     Le Pack RBK 2.0 . . . . . . . . . . . . . . . . . . . . . . . . . . . . . . . . . . . . . . . . . . . . . . . . . . . . . . . . . . . . . . . . . . . . . . . . . . . . . . . . . . . . . . . . . . . . . . . . . . . . . . . . . . . . . . . . . 252\\
K.2     Pricing \textbackslash{}& Conditions (Value Ladder) . . . . . . . . . . . . . . . . . . . . . . . . . . . . . . . . . . . . . . . . . . . . . . . . . . . . . . . . . . . . . . . . . . . . . . . . . . . . . . . . . . . . . . . . . . . . . 252\\
K.2.1               Mécanisme d’Incitation (Upgrade) . . . . . . . . . . . . . . . . . . . . . . . . . . . . . . . . . . . . . . . . . . . . . . . . . . . . . . . . . . . . . . . . . . . . . . . . . . . . . . . . . . . . . . . 254\\
K.2.2               Admission Directe (Passerelles) . . . . . . . . . . . . . . . . . . . . . . . . . . . . . . . . . . . . . . . . . . . . . . . . . . . . . . . . . . . . . . . . . . . . . . . . . . . . . . . . . . . . . . . . . . 254\\
K.2.3               Offre ISA (Income Share Agreement) . . . . . . . . . . . . . . . . . . . . . . . . . . . . . . . . . . . . . . . . . . . . . . . . . . . . . . . . . . . . . . . . . . . . . . . . . . . . . . . . . . . . . 254\\
K.3     Objections \textbackslash{}& Réponses . . . . . . . . . . . . . . . . . . . . . . . . . . . . . . . . . . . . . . . . . . . . . . . . . . . . . . . . . . . . . . . . . . . . . . . . . . . . . . . . . . . . . . . . . . . . . . . . . . . . . . . . . . . 255\\
K.4     Politique de Remboursement et Report . . . . . . . . . . . . . . . . . . . . . . . . . . . . . . . . . . . . . . . . . . . . . . . . . . . . . . . . . . . . . . . . . . . . . . . . . . . . . . . . . . . . . . . . . . . 255\\
\vspace{0.5em}
L     Stratégie mentorat . . . . . . . . . . . . . . . . . . . . . . . . . . . . . . . . . . . . . . . . . . . . . . . . . . . 256\\
L.1     Programme ”Mentor-in-a-Box” . . . . . . . . . . . . . . . . . . . . . . . . . . . . . . . . . . . . . . . . . . . . . . . . . . . . . . . . . . . . . . . . . . . . . . . . . . . . . . . . . . . . . . . . . . . . . . . . . . . 256\\
\vspace{0.5em}
\vspace{0.5em}
\vspace{0.5em}
\vspace{0.5em}
EL\\
L.2     Le Pipeline ”Train the Trainer” . . . . . . . . . . . . . . . . . . . . . . . . . . . . . . . . . . . . . . . . . . . . . . . . . . . . . . . . . . . . . . . . . . . . . . . . . . . . . . . . . . . . . . . . . . . . . . . . . . 256\\
L.3     Modèle de Rémunération Incitatif . . . . . . . . . . . . . . . . . . . . . . . . . . . . . . . . . . . . . . . . . . . . . . . . . . . . . . . . . . . . . . . . . . . . . . . . . . . . . . . . . . . . . . . . . . . . . . . . 257\\
\vspace{0.5em}
\vspace{0.5em}
\vspace{0.5em}
\vspace{0.5em}
TI\\
L.4     Contrat type Mentor Nexus . . . . . . . . . . . . . . . . . . . . . . . . . . . . . . . . . . . . . . . . . . . . . . . . . . . . . . . . . . . . . . . . . . . . . . . . . . . . . . . . . . . . . . . . . . . . . . . . . . . . . . 257\\
\vspace{0.5em}
\vspace{0.5em}
\vspace{0.5em}
\vspace{0.5em}
EN\\
L.5     Plan de Relève et Continuité . . . . . . . . . . . . . . . . . . . . . . . . . . . . . . . . . . . . . . . . . . . . . . . . . . . . . . . . . . . . . . . . . . . . . . . . . . . . . . . . . . . . . . . . . . . . . . . . . . . . . 258\\
\vspace{0.5em}
M     Offre partenariat B2B . . . . . . . . . . . . . . . . . . . . . . . . . . . . . . . . . . . . . . . . . . . . . . . . . 259\\
\vspace{0.5em}
\vspace{0.5em}
\vspace{0.5em}
\vspace{0.5em}
D\\
Modèle d’Offre Corporate . . . . . . . . . . . . . . . . . . . . . . . . . . . . . . . . . . . . . . . . . . . . . . . . . . . . . . . . . . . . . . . . . . . . . . . . . . . . . . . . . . . . . . . . . . . . . . . . . . . . . . . . 259\\
\vspace{0.5em}
\vspace{0.5em}
\vspace{0.5em}
\vspace{0.5em}
FI\\
M.1\\
M.1.1               Les Packs Entreprise . . . . . . . . . . . . . . . . . . . . . . . . . . . . . . . . . . . . . . . . . . . . . . . . . . . . . . . . . . . . . . . . . . . . . . . . . . . . . . . . . . . . . . . . . . . . . . . . . . . . 259\\
M.2                                                                           N\\
Conditions Particulières . . . . . . . . . . . . . . . . . . . . . . . . . . . . . . . . . . . . . . . . . . . . . . . . . . . . . . . . . . . . . . . . . . . . . . . . . . . . . . . . . . . . . . . . . . . . . . . . . . . . . . . . . . 259\\
CO\\
N     Outillage \textbackslash{}& stack technique . . . . . . . . . . . . . . . . . . . . . . . . . . . . . . . . . . . . . . . . . . . . . 261\\
N.1     Stack de Développement (Cyborg-Ready) . . . . . . . . . . . . . . . . . . . . . . . . . . . . . . . . . . . . . . . . . . . . . . . . . . . . . . . . . . . . . . . . . . . . . . . . . . . . . . . . . . . . . . . . 261\\
N.1.1               Environnement Local . . . . . . . . . . . . . . . . . . . . . . . . . . . . . . . . . . . . . . . . . . . . . . . . . . . . . . . . . . . . . . . . . . . . . . . . . . . . . . . . . . . . . . . . . . . . . . . . . . . 261\\
N.1.2               Chain Stack . . . . . . . . . . . . . . . . . . . . . . . . . . . . . . . . . . . . . . . . . . . . . . . . . . . . . . . . . . . . . . . . . . . . . . . . . . . . . . . . . . . . . . . . . . . . . . . . . . . . . . . . . . . . 261\\
N.2     Outils de Productivité \textbackslash{}& IA . . . . . . . . . . . . . . . . . . . . . . . . . . . . . . . . . . . . . . . . . . . . . . . . . . . . . . . . . . . . . . . . . . . . . . . . . . . . . . . . . . . . . . . . . . . . . . . . . . . . . . 262\\
N.3     Infrastructure CI/CD (Github Actions) . . . . . . . . . . . . . . . . . . . . . . . . . . . . . . . . . . . . . . . . . . . . . . . . . . . . . . . . . . . . . . . . . . . . . . . . . . . . . . . . . . . . . . . . . . . 262\\
\vspace{0.5em}
O     RÉFÉRENTIEL DE COMPÉTENCES . . . . . . . . . . . . . . . . . . . . . . . . . . . . . . . . . . . . . . . . . . 263\\
O.1     Matrice de Compétences . . . . . . . . . . . . . . . . . . . . . . . . . . . . . . . . . . . . . . . . . . . . . . . . . . . . . . . . . . . . . . . . . . . . . . . . . . . . . . . . . . . . . . . . . . . . . . . . . . . . . . . . . 263\\
\vspace{0.5em}
\vspace{0.5em}
\vspace{0.5em}
\begin{confidential}CONFIDENTIEL — PROJET RBK 2.0                                                                               Page 15 sur 316\end{confidential}
RBK 2.0 — Whitepaper Architecte Web3\\
\vspace{0.5em}
\vspace{0.5em}
P     Charte de qualité \textbackslash{}& règles d’or . . . . . . . . . . . . . . . . . . . . . . . . . . . . . . . . . . . . . . . . . . . 264\\
P.1     Les 4 Commandements de l’Ingénieur RBK . . . . . . . . . . . . . . . . . . . . . . . . . . . . . . . . . . . . . . . . . . . . . . . . . . . . . . . . . . . . . . . . . . . . . . . . . . . . . . . . . . . . . . . 264\\
P.2     Matrice de Conformité (Sanctions) . . . . . . . . . . . . . . . . . . . . . . . . . . . . . . . . . . . . . . . . . . . . . . . . . . . . . . . . . . . . . . . . . . . . . . . . . . . . . . . . . . . . . . . . . . . . . . . 264\\
P.3     Processus de Validation Qualité . . . . . . . . . . . . . . . . . . . . . . . . . . . . . . . . . . . . . . . . . . . . . . . . . . . . . . . . . . . . . . . . . . . . . . . . . . . . . . . . . . . . . . . . . . . . . . . . . . 265\\
\vspace{0.5em}
Q     CONTRAT DE PARTENARIAT PÉDAGOGIQUE ET DE PARTAGE DE SUCCÈS (CPPS) . . . . . . . . . . . 266\\
Q.1     Objet du Contrat . . . . . . . . . . . . . . . . . . . . . . . . . . . . . . . . . . . . . . . . . . . . . . . . . . . . . . . . . . . . . . . . . . . . . . . . . . . . . . . . . . . . . . . . . . . . . . . . . . . . . . . . . . . . . . . . . 266\\
Q.2     Périmètre des revenus éligibles . . . . . . . . . . . . . . . . . . . . . . . . . . . . . . . . . . . . . . . . . . . . . . . . . . . . . . . . . . . . . . . . . . . . . . . . . . . . . . . . . . . . . . . . . . . . . . . . . . 266\\
Q.3     Définitions clés . . . . . . . . . . . . . . . . . . . . . . . . . . . . . . . . . . . . . . . . . . . . . . . . . . . . . . . . . . . . . . . . . . . . . . . . . . . . . . . . . . . . . . . . . . . . . . . . . . . . . . . . . . . . . . . . . . 267\\
Q.4     Obligations de RBK . . . . . . . . . . . . . . . . . . . . . . . . . . . . . . . . . . . . . . . . . . . . . . . . . . . . . . . . . . . . . . . . . . . . . . . . . . . . . . . . . . . . . . . . . . . . . . . . . . . . . . . . . . . . . . 267\\
\vspace{0.5em}
\vspace{0.5em}
\vspace{0.5em}
\vspace{0.5em}
EL\\
Q.5     Obligations du bénéficiaire . . . . . . . . . . . . . . . . . . . . . . . . . . . . . . . . . . . . . . . . . . . . . . . . . . . . . . . . . . . . . . . . . . . . . . . . . . . . . . . . . . . . . . . . . . . . . . . . . . . . . . 267\\
\vspace{0.5em}
\vspace{0.5em}
\vspace{0.5em}
\vspace{0.5em}
TI\\
Q.6     Collecte et validation des revenus . . . . . . . . . . . . . . . . . . . . . . . . . . . . . . . . . . . . . . . . . . . . . . . . . . . . . . . . . . . . . . . . . . . . . . . . . . . . . . . . . . . . . . . . . . . . . . . . 267\\
Q.7     Modalités de paiement . . . . . . . . . . . . . . . . . . . . . . . . . . . . . . . . . . . . . . . . . . . . . . . . . . . . . . . . . . . . . . . . . . . . . . . . . . . . . . . . . . . . . . . . . . . . . . . . . . . . . . . . . . . 268\\
\vspace{0.5em}
\vspace{0.5em}
\vspace{0.5em}
\vspace{0.5em}
EN\\
Q.8     Gestion des litiges et cas de force majeure . . . . . . . . . . . . . . . . . . . . . . . . . . . . . . . . . . . . . . . . . . . . . . . . . . . . . . . . . . . . . . . . . . . . . . . . . . . . . . . . . . . . . . . . 268\\
\vspace{0.5em}
\vspace{0.5em}
\vspace{0.5em}
\vspace{0.5em}
D\\
Q.9     Recouvrement et protection du bénéficiaire . . . . . . . . . . . . . . . . . . . . . . . . . . . . . . . . . . . . . . . . . . . . . . . . . . . . . . . . . . . . . . . . . . . . . . . . . . . . . . . . . . . . . . 268\\
\vspace{0.5em}
\vspace{0.5em}
\vspace{0.5em}
\vspace{0.5em}
FI\\
R     MODÈLE DE PARTENARIAT B2B (HIRING) . . . . . . . . . . . . . . . . . . . . . . . . . . . . . . . . . . . . 269\\
R.1\\
N\\
Offre ”Hire Train Deploy” . . . . . . . . . . . . . . . . . . . . . . . . . . . . . . . . . . . . . . . . . . . . . . . . . . . . . . . . . . . . . . . . . . . . . . . . . . . . . . . . . . . . . . . . . . . . . . . . . . . . . . . . 269\\
CO\\
R.2     Offre ”Corporate Upskilling” . . . . . . . . . . . . . . . . . . . . . . . . . . . . . . . . . . . . . . . . . . . . . . . . . . . . . . . . . . . . . . . . . . . . . . . . . . . . . . . . . . . . . . . . . . . . . . . . . . . . . 269\\
\vspace{0.5em}
S     Kit de survie juridique . . . . . . . . . . . . . . . . . . . . . . . . . . . . . . . . . . . . . . . . . . . . . . . . 270\\
S.1     Modèle de Contrat de Prestation Freelance (Extraits) . . . . . . . . . . . . . . . . . . . . . . . . . . . . . . . . . . . . . . . . . . . . . . . . . . . . . . . . . . . . . . . . . . . . . . . . . . . . . 270\\
S.2     Checklist : Créer sa Micro-Entreprise Exportatrice . . . . . . . . . . . . . . . . . . . . . . . . . . . . . . . . . . . . . . . . . . . . . . . . . . . . . . . . . . . . . . . . . . . . . . . . . . . . . . . . 270\\
S.3     Guide Visuel : Recevoir un Salaire en Crypto . . . . . . . . . . . . . . . . . . . . . . . . . . . . . . . . . . . . . . . . . . . . . . . . . . . . . . . . . . . . . . . . . . . . . . . . . . . . . . . . . . . . . 271\\
S.4     Red Flags (Vigilance) . . . . . . . . . . . . . . . . . . . . . . . . . . . . . . . . . . . . . . . . . . . . . . . . . . . . . . . . . . . . . . . . . . . . . . . . . . . . . . . . . . . . . . . . . . . . . . . . . . . . . . . . . . . . 271\\
\vspace{0.5em}
T     Annexe : Modèle Piscine \textbackslash{}& Staffing . . . . . . . . . . . . . . . . . . . . . . . . . . . . . . . . . . . . . . . . 272\\
T.1     Synthèse du Modèle Piscine et Staffing . . . . . . . . . . . . . . . . . . . . . . . . . . . . . . . . . . . . . . . . . . . . . . . . . . . . . . . . . . . . . . . . . . . . . . . . . . . . . . . . . . . . . . . . . . . 272\\
\vspace{0.5em}
\vspace{0.5em}
\vspace{0.5em}
\begin{confidential}CONFIDENTIEL — PROJET RBK 2.0                                                                               Page 16 sur 316\end{confidential}
RBK 2.0 — Whitepaper Architecte Web3\\
\vspace{0.5em}
\vspace{0.5em}
U     Registre des risques opérationnels . . . . . . . . . . . . . . . . . . . . . . . . . . . . . . . . . . . . . . . . . 275\\
V     Plan d’actions risques \textbackslash{}& tableau de bord KPI (v1) . . . . . . . . . . . . . . . . . . . . . . . . . . . . . . . 281\\
V.1     Annexe X — Traceability Matrix (SSOT \textbackslash{}& Release Gate) . . . . . . . . . . . . . . . . . . . . . . . . . . . . . . . . . . . . . . . . . . . . . . . . . . . . . . . . . . . . . . . . . . . . . . . . . . 301\\
V.2     Annexe Y — Change Control (SSOT \textbackslash{}& Release Gate) . . . . . . . . . . . . . . . . . . . . . . . . . . . . . . . . . . . . . . . . . . . . . . . . . . . . . . . . . . . . . . . . . . . . . . . . . . . . . 304\\
\vspace{0.5em}
W     Glossaire des termes, acronymes et jargon . . . . . . . . . . . . . . . . . . . . . . . . . . . . . . . . . . . . 306\\
Journal des modifications v5.0 . . . . . . . . . . . . . . . . . . . . . . . . . . . . . . . . . . . . . . . . . . . . . . 312\\
References \textbackslash{}& Bibliographie . . . . . . . . . . . . . . . . . . . . . . . . . . . . . . . . . . . . . . . . . . . . . . . . . 315\\
W.1     Documentation Technique . . . . . . . . . . . . . . . . . . . . . . . . . . . . . . . . . . . . . . . . . . . . . . . . . . . . . . . . . . . . . . . . . . . . . . . . . . . . . . . . . . . . . . . . . . . . . . . . . . . . . . . 315\\
\vspace{0.5em}
\vspace{0.5em}
\vspace{0.5em}
\vspace{0.5em}
EL\\
W.2     Rapports de l’Industrie . . . . . . . . . . . . . . . . . . . . . . . . . . . . . . . . . . . . . . . . . . . . . . . . . . . . . . . . . . . . . . . . . . . . . . . . . . . . . . . . . . . . . . . . . . . . . . . . . . . . . . . . . . . 315\\
W.3     Rapports de Marché . . . . . . . . . . . . . . . . . . . . . . . . . . . . . . . . . . . . . . . . . . . . . . . . . . . . . . . . . . . . . . . . . . . . . . . . . . . . . . . . . . . . . . . . . . . . . . . . . . . . . . . . . . . . . 315\\
\vspace{0.5em}
\vspace{0.5em}
\vspace{0.5em}
\vspace{0.5em}
TI\\
W.4     Outils Cités . . . . . . . . . . . . . . . . . . . . . . . . . . . . . . . . . . . . . . . . . . . . . . . . . . . . . . . . . . . . . . . . . . . . . . . . . . . . . . . . . . . . . . . . . . . . . . . . . . . . . . . . . . . . . . . . . . . . . . 316\\
\vspace{0.5em}
\vspace{0.5em}
\vspace{0.5em}
\vspace{0.5em}
D            EN\\
FI\\
N\\
CO\\
\vspace{0.5em}
\vspace{0.5em}
\vspace{0.5em}
\vspace{0.5em}
\begin{confidential}CONFIDENTIEL — PROJET RBK 2.0                                                                                Page 17 sur 316\end{confidential}
GUIDE DE LECTURE\\
\vspace{0.5em}
\vspace{0.5em}
\vspace{0.5em}
\vspace{0.5em}
EL\\
Ce document est conçu pour servir plusieurs audiences. Voici les parcours de lecture recommandés pour naviguer efficacement :\\
\vspace{0.5em}
\vspace{0.5em}
\vspace{0.5em}
\vspace{0.5em}
TI\\
þ Pour les Investisseurs : Concentrez-vous sur la validation du modèle économique, la scalabilité et la gestion des risques.\\
\vspace{0.5em}
\vspace{0.5em}
\vspace{0.5em}
\vspace{0.5em}
EN\\
• Chapitre 1 (Vision) : La thèse d’investissement (”Senior-by-Design”).\\
• Chapitre 8 (Business Plan) : Le modèle hybride, les Unit Economics et le P\textbackslash{}&L.\\
\vspace{0.5em}
\vspace{0.5em}
\vspace{0.5em}
\vspace{0.5em}
D\\
• Chapitre 10 (Risques) : La matrice de risques et la mitigation (notamment ISA).\\
\vspace{0.5em}
\vspace{0.5em}
\vspace{0.5em}
FI\\
• Annexe B (Finance) \textbackslash{}& F (ISA) : Les détails techniques des hypothèses financières.\\
N\\
CO\\
ŵ Pour les Candidats (Étudiants) : Comprenez l’intensité du programme et les pré-requis pour réussir.\\
\vspace{0.5em}
• Chapitre 1 (Vision) : Pourquoi RBK n’est pas une ”école” classique.\\
• Chapitre 5 (Structure) \textbackslash{}& 6 (Syllabus) : Le rythme, les phases et les livrables.\\
• Annexe J (Offre) : Les tarifs, le fonctionnement des Niveaux et du Pack.\\
• Annexe G (Sélection) : Comment se préparer aux tests d’entrée.\\
\vspace{0.5em}
ǹ Pour les Partenaires B2B : Découvrez comment intégrer vos technologies ou recruter nos talents d’élite.\\
\vspace{0.5em}
• Chapitre 8 (Business Plan) : L’offre Corporate et le Hiring.\\
\vspace{0.5em}
\vspace{0.5em}
\vspace{0.5em}
18\\
RBK 2.0 — Whitepaper Architecte Web3\\
\vspace{0.5em}
\vspace{0.5em}
• Chapitre 4 (Méthodologie) : La rigueur de notre process ”Cyborg”.\\
• Annexe M (Partenariats) : Les modalités de collaboration (Sponsoring, Recrutement).\\
\vspace{0.5em}
\vspace{0.5em}
\vspace{0.5em}
\vspace{0.5em}
EL\\
TI\\
D   EN\\
FI\\
N\\
CO\\
\vspace{0.5em}
\vspace{0.5em}
\vspace{0.5em}
\vspace{0.5em}
\begin{confidential}CONFIDENTIEL — PROJET RBK 2.0                             Page 19 sur 316\end{confidential}
Liste des Acronymes\\
\vspace{0.5em}
\vspace{0.5em}
\vspace{0.5em}
\vspace{0.5em}
EL\\
API Application Programming Interface\\
\vspace{0.5em}
\vspace{0.5em}
\vspace{0.5em}
\vspace{0.5em}
TI\\
CAGR Compound Annual Growth Rate (Taux de croissance annuel moyen)\\
\vspace{0.5em}
\vspace{0.5em}
\vspace{0.5em}
\vspace{0.5em}
EN\\
CI/CD Continuous Integration / Continuous Deployment\\
\vspace{0.5em}
CLI Command Line Interface\\
\vspace{0.5em}
\vspace{0.5em}
\vspace{0.5em}
\vspace{0.5em}
D\\
FI\\
DAO Decentralized Autonomous Organization\\
\vspace{0.5em}
DApp Decentralized Application\\
N\\
CO\\
DeFi Decentralized Finance\\
\vspace{0.5em}
DePIN Decentralized Physical Infrastructure Networks\\
\vspace{0.5em}
DoD Definition of Done\\
\vspace{0.5em}
EVM Ethereum Virtual Machine\\
\vspace{0.5em}
ISA Income Share Agreement\\
\vspace{0.5em}
KPI Key Performance Indicator\\
\vspace{0.5em}
\vspace{0.5em}
\vspace{0.5em}
20\\
RBK 2.0 — Whitepaper Architecte Web3\\
\vspace{0.5em}
\vspace{0.5em}
L1/L2 Layer 1 (Blockchain de base) / Layer 2 (Couche de mise à l’échelle)\\
\vspace{0.5em}
MVP Minimum Viable Product\\
\vspace{0.5em}
PDA Program Derived Address (Solana)\\
\vspace{0.5em}
PoS Proof of Stake\\
\vspace{0.5em}
PR Pull Request\\
\vspace{0.5em}
ROI Return on Investment\\
\vspace{0.5em}
RPC Remote Procedure Call\\
\vspace{0.5em}
\vspace{0.5em}
\vspace{0.5em}
\vspace{0.5em}
EL\\
SBT Soulbound Token\\
\vspace{0.5em}
\vspace{0.5em}
\vspace{0.5em}
\vspace{0.5em}
TI\\
SVM Solana Virtual Machine\\
\vspace{0.5em}
\vspace{0.5em}
\vspace{0.5em}
\vspace{0.5em}
EN\\
TVL Total Value Locked\\
\vspace{0.5em}
\vspace{0.5em}
\vspace{0.5em}
\vspace{0.5em}
D\\
UX/UI User Experience / User Interface\\
\vspace{0.5em}
\vspace{0.5em}
\vspace{0.5em}
\vspace{0.5em}
FI\\
N\\
CO\\
\vspace{0.5em}
\vspace{0.5em}
\vspace{0.5em}
\vspace{0.5em}
\begin{confidential}CONFIDENTIEL — PROJET RBK 2.0                              Page 21 sur 316\end{confidential}
Factsheet RBK 2.0\\
\vspace{0.5em}
\vspace{0.5em}
\vspace{0.5em}
\vspace{0.5em}
EL\\
Ď Résumé Opérationnel\\
\vspace{0.5em}
\vspace{0.5em}
\vspace{0.5em}
\vspace{0.5em}
TI\\
Ce document synthétise les paramètres clés du programme RBK (ReBoot Kamp) pour l’année 2026. Il fait foi pour les partenaires,\\
\vspace{0.5em}
\vspace{0.5em}
\vspace{0.5em}
\vspace{0.5em}
EN\\
investisseurs et candidats.\\
\vspace{0.5em}
\vspace{0.5em}
\vspace{0.5em}
\vspace{0.5em}
D\\
FI\\
Paramètres de la Cohorte\\
\vspace{0.5em}
N\\
CO\\
\vspace{0.5em}
\vspace{0.5em}
\vspace{0.5em}
\vspace{0.5em}
22\\
RBK 2.0 — Whitepaper Architecte Web3\\
\vspace{0.5em}
\vspace{0.5em}
\vspace{0.5em}
\vspace{0.5em}
Paramètre                   Valeur Validée\\
\vspace{0.5em}
Cadence annuelle            48 semaines (Genesis Pool 4 + Explorer 12 + Bâtisseur 16 + Architecte 16) complétées par un Launchpad\\
optionnel de 4 semaines pour les projets incubés.\\
\vspace{0.5em}
Format                      Campus Tunis (présentiel augmenté) et Remote synchrone, cohorte baptisée “Genesis Cohort” avec mixité\\
internationale.\\
\vspace{0.5em}
Rythme                      Full-time (9h00 – 18h00) + sessions de veille technique planifiées ; aucune exigence nocturne, protocole\\
\vspace{0.5em}
\vspace{0.5em}
\vspace{0.5em}
\vspace{0.5em}
EL\\
anti-burnout appliqué.\\
\vspace{0.5em}
\vspace{0.5em}
\vspace{0.5em}
\vspace{0.5em}
TI\\
Phase d’entrée              Genesis Pool intensive de 4 semaines (bootcamp sélectif, 120h, 80\textbackslash{}% pair-programming).\\
\vspace{0.5em}
\vspace{0.5em}
\vspace{0.5em}
\vspace{0.5em}
EN\\
Niveaux pédagogiques        Explorer (12 sem.), Bâtisseur (16 sem.), Architecte (16 sem.). Chaque niveau s’achève par un Block Check\\
certifiant.\\
\vspace{0.5em}
\vspace{0.5em}
\vspace{0.5em}
\vspace{0.5em}
D\\
Taille Cohorte              25 – 30 apprenants par Genesis Cohort, ratio mentor référent 1 :25 sur Explorer.\\
\vspace{0.5em}
\vspace{0.5em}
\vspace{0.5em}
\vspace{0.5em}
FI\\
Tracks                      Solana Smart Contracts, EVM Engineering, Product Strategy, RBK Solana Validator Track (nouveau).\\
N\\
CO\\
Langue                      Anglais (documentation, code) / Français (animation, soutenances).\\
\vspace{0.5em}
\vspace{0.5em}
\vspace{0.5em}
Conditions d’Accès \textbackslash{}& Tarifs\\
• Prérequis :\\
\vspace{0.5em}
– Algorithmique solide (Test technique éliminatoire).\\
– Anglais technique lu/écrit (Niveau B2 minimum).\\
– Engagement full-time indispensable.\\
\vspace{0.5em}
\vspace{0.5em}
\begin{confidential}CONFIDENTIEL — PROJET RBK 2.0                                 Page 23 sur 316\end{confidential}
RBK 2.0 — Whitepaper Architecte Web3\\
\vspace{0.5em}
\vspace{0.5em}
• Modèle Économique :\\
\vspace{0.5em}
– Frais annuels : 15 900 TND TTC par apprenant pour le cycle complet Genesis Pool + Explorer + Bâtisseur + Architecte.\\
\vspace{0.5em}
\vspace{0.5em}
– Forfait opérateur : Nexus Réussite, maître d’œuvre pédagogique exclusif, facture 3 300 TND par apprenant (design, mentors,\\
supervision et usage de la plateforme Venture Engine de Money Factory AI).\\
\vspace{0.5em}
\vspace{0.5em}
– Partage incubateur : Nexus Réussite perçoit 30\textbackslash{}% des revenus générés par les projets incubés via le Launchpad.\\
\vspace{0.5em}
\vspace{0.5em}
– ISA : Seuil : 2500 TND brut ; Partage : 15\textbackslash{}% du brut ; Cap : 20000 TND ; Durée : 36 mensualités payées max. (réf. Annexe K) ;\\
\vspace{0.5em}
\vspace{0.5em}
\vspace{0.5em}
\vspace{0.5em}
EL\\
contrat RBK – Apprenant géré intégralement par RBK via Fonds de Garantie ISA2 dédié.\\
\vspace{0.5em}
\vspace{0.5em}
\vspace{0.5em}
\vspace{0.5em}
TI\\
– Entreprise : Parcours Corporate “Dev DAO” (B2B) sur devis, aligné sur la grille Nexus.\\
\vspace{0.5em}
\vspace{0.5em}
\vspace{0.5em}
\vspace{0.5em}
EN\\
ISA3 .\\
\vspace{0.5em}
\vspace{0.5em}
\vspace{0.5em}
\vspace{0.5em}
D\\
FI\\
Paramètres ISA — version unifiée (Annexe K)\\
\vspace{0.5em}
Cadre juridique : Contrat de Partenariat Pédagogiquea et de Partage de Succès (CPPS) signé entre RBK et l’apprenant. Seuil : 2 500\\
N\\
TND brut/mois. Taux : 15\textbackslash{}% du brut mensuel. Cap : 20 000 TND. Durée : 36 mensualités payées max. Pause automatique si\\
CO\\
revenu \textless{} seuil. Référence : Annexe K (Q).\\
a\\
Finance/Juridique/Compliance — Accord conditionnant le remboursement à l’obtention d’un revenu dans le périmètre Web3, encadrant droits, obligations\\
et procédures de recours pour les apprenants RBK.\\
\vspace{0.5em}
\vspace{0.5em}
\vspace{0.5em}
\vspace{0.5em}
Mentorat de la première cohorte\\
2\\
Finance/Juridique/Compliance — Compte séquestré capitalisant les contributions RBK/ISA pour couvrir défauts temporaires, avances de bourses et incidents de\\
paiement, gouverné conjointement avec Nexus.\\
3\\
Finance/Juridique/Compliance — Income Share Agreement (ISA) : contrat où l’étudiant rembourse une part de revenu selon des paramètres centralisés dans la\\
SSOT.\\
\vspace{0.5em}
\vspace{0.5em}
\vspace{0.5em}
\begin{confidential}CONFIDENTIEL — PROJET RBK 2.0                                          Page 24 sur 316\end{confidential}
RBK 2.0 — Whitepaper Architecte Web3\\
\vspace{0.5em}
\vspace{0.5em}
\vspace{0.5em}
• Explorer (Sem. 1-12) : 2 formateurs référents Nexus (Senior Solana/EVM) assurent 20\textbackslash{}% d’animation synchrone, le reste via Skill\\
Mirror et NFT Skills.\\
\vspace{0.5em}
• Bâtisseur (Sem. 13-28) : Passage à un modèle hybride, mentors internes (Architecte N-1) validés par NFT Skills “Mentor Lead”.\\
\vspace{0.5em}
• Architecte (Sem. 29-44) : Encadrement par les Architectes sortants (Dev DAO) + coach produit, monitoring qualité par Block Check\\
hebdomadaire.\\
\vspace{0.5em}
• Launchpad (Option) : Mentorat business par Nexus Réussite, alignment incubateur et suivi des revenus partagés.\\
\vspace{0.5em}
\vspace{0.5em}
\vspace{0.5em}
\vspace{0.5em}
EL\\
Objectifs de Sortie (KPIs)\\
\vspace{0.5em}
\vspace{0.5em}
\vspace{0.5em}
\vspace{0.5em}
TI\\
Les KPI4 ci-dessous structurent le pilotage cohorte et la viabilité du modèle.\\
\vspace{0.5em}
\vspace{0.5em}
\vspace{0.5em}
\vspace{0.5em}
EN\\
KPI \textbackslash{}& ROI\\
\vspace{0.5em}
\vspace{0.5em}
\vspace{0.5em}
\vspace{0.5em}
D\\
Indicateur                       Définition                     Cible            Mesure / Source\\
\vspace{0.5em}
\vspace{0.5em}
\vspace{0.5em}
\vspace{0.5em}
FI\\
Taux de complétion               Graduation                     90\textbackslash{}%              Output Capstone\\
Placement (post-parcours)\\
N  CDI / Freelance / Grant        80\textbackslash{}%              Suivi Carrière\\
CO\\
Salaire moyen                    Premier emploi                 \textgreater{} 2.5k TND       Contrats signés\\
Satisfaction                     NPS                            \textgreater{} 70             Enquêtes anonymes\\
Break-even                       Rentabilité Promo              M+12             Cashflow ISA\\
\vspace{0.5em}
\vspace{0.5em}
\vspace{0.5em}
Livrables Apprenant (Portfolio)\\
4\\
Projet/Management/Qualité — Key Performance Indicator (KPI) : indicateur clé mesurant l’atteinte d’un objectif stratégique ou opérationnel.\\
\vspace{0.5em}
\vspace{0.5em}
\vspace{0.5em}
\vspace{0.5em}
\begin{confidential}CONFIDENTIEL — PROJET RBK 2.0                                            Page 25 sur 316\end{confidential}
RBK 2.0 — Whitepaper Architecte Web3\\
\vspace{0.5em}
\vspace{0.5em}
\vspace{0.5em}
À la fin du cursus, chaque ”Survivor” possède :\\
\vspace{0.5em}
1. 3 Projets Capstone complets (Code, Tests, Docs) sur GitHub.\\
\vspace{0.5em}
2. 1 Audit Report (Sécurité) sur un protocole tiers.\\
\vspace{0.5em}
3. 1 Contribution Open Source validée (PR).\\
\vspace{0.5em}
4. NFT Skills “Launch Proof” attestant de la réussite Block Check final, couplé au SBT “Audit-Ready”.\\
\vspace{0.5em}
Disclaimer\\
\vspace{0.5em}
\vspace{0.5em}
\vspace{0.5em}
\vspace{0.5em}
EL\\
RBK est un programme d’ingénierie logicielle. Ce n’est pas une académie de trading, d’investissement ou de conseil financier. Nous\\
formons des constructeurs, pas des spéculateurs. Tout contenu relatif aux tokens est abordé sous l’angle technique et technologique.\\
\vspace{0.5em}
\vspace{0.5em}
\vspace{0.5em}
\vspace{0.5em}
TI\\
D   EN\\
FI\\
N\\
CO\\
\vspace{0.5em}
\vspace{0.5em}
\vspace{0.5em}
\vspace{0.5em}
\begin{confidential}CONFIDENTIEL — PROJET RBK 2.0                                Page 26 sur 316\end{confidential}
EXECUTIVE SUMMARY\\
\vspace{0.5em}
\vspace{0.5em}
\vspace{0.5em}
\vspace{0.5em}
EL\\
Le Constat : La Fin du ”Junior” et l’Urgence Web3\\
\vspace{0.5em}
\vspace{0.5em}
\vspace{0.5em}
\vspace{0.5em}
TI\\
L’industrie technologique traverse une mutation violente. L’intelligence artificielle générative (LLMs) a commoditisé la production de\\
\vspace{0.5em}
\vspace{0.5em}
\vspace{0.5em}
\vspace{0.5em}
EN\\
code simple, rendant le profil de ”développeur junior” économiquement obsolèteIllustration Concrete5 .\\
Parallèlement, l’économie décentralisée (Web36 ) connaît une croissance institutionnelle sans précédent (+25\textbackslash{}% CAGR), créant une\\
\vspace{0.5em}
\vspace{0.5em}
\vspace{0.5em}
\vspace{0.5em}
D\\
pénurie mondiale de talents capables de concevoir des architectures sécurisées et complexes.\\
Le marché ne cherche plus des exécutants ; il cherche des Architectes Web37 .\\
\vspace{0.5em}
\vspace{0.5em}
\vspace{0.5em}
FI\\
N\\
CO\\
La Solution : RBK 3.0 ”Senior-by-Design”\\
RBK Web3 Studio n’est pas une simple école de code. C’est un changement de paradigme éducatif opéré par RBK, entité légale et\\
commerciale, et exécuté pédagogiquement par Nexus Réussite, son maître d’œuvre pédagogique exclusif. Ce binôme industrialise la\\
montée en compétence Web3 : notre modèle RBK 3.0 décline l’inspiration 42 en quatre strates : Genesis Pool sélective, niveaux Explorer,\\
Bâtisseur, Architecte et un Launchpad optionnel. Chaque jalon est validé par un Skill Mirror pair-à-pair puis un Block Check certifiant\\
générant des NFT Skills et documentés dans Venture Engine (Money Factory AI).\\
5\\
Pédagogie — Exemple concret illustratif utilisé pour faciliter la compréhension d’un concept.\\
6\\
Web3 — Web 3.0 (Web3) : génération d’Internet décentralisée basée sur la propriété numérique via la blockchain, redonnant le contrôle aux utilisateurs.\\
7\\
Carrière — Ingénieur capable de concevoir des systèmes décentralisés complets (Smart Contracts + Frontend + Indexing), en maîtrisant les enjeux de sécurité,\\
de coût (Gas) et de performance.\\
\vspace{0.5em}
\vspace{0.5em}
\vspace{0.5em}
27\\
RBK 2.0 — Whitepaper Architecte Web3\\
\vspace{0.5em}
\vspace{0.5em}
Notre Promesse : Former en 48 semaines (Genesis Pool 4 + Explorer 12 + Bâtisseur 16 + Architecte 16) des ingénieurs possédant la\\
maturité technique d’un profil de 3 ans d’expérience (”Senior-by-Design8 ”), audités, certifiés on-chain via les NFT Skills “Launch Proof”,\\
et prêts à déployer de la valeur dès le jour 1.\\
\vspace{0.5em}
• Genesis Pool (4 sem.) : bootcamp intensif sans IA, culture Dev DAO, sélection des 30 meilleurs profils.\\
\vspace{0.5em}
• Explorer (12 sem.) : fondamentaux Web3 + pipeline produit, validation Block Check “Explorer”.\\
\vspace{0.5em}
• Bâtisseur (16 sem.) : projets pairés, spécialisation Solana/EVM/Product, préparation RBK Solana Validator Track.\\
\vspace{0.5em}
• Architecte (16 sem.) : leadership technique, mentorat des niveaux inférieurs, production d’audits.\\
\vspace{0.5em}
\vspace{0.5em}
\vspace{0.5em}
\vspace{0.5em}
EL\\
• Launchpad (4 sem., sur dossier) : incubation, revenue share 30\textbackslash{}% Nexus / 70\textbackslash{}% RBK ou startup.\\
\vspace{0.5em}
\vspace{0.5em}
\vspace{0.5em}
\vspace{0.5em}
TI\\
Cibles et Personas\\
\vspace{0.5em}
\vspace{0.5em}
\vspace{0.5em}
\vspace{0.5em}
EN\\
• Le Junior Ambitieux : Diplômé CS ou autodidacte talentueux bloqué par le ”plafond de verre” du marché local.\\
\vspace{0.5em}
\vspace{0.5em}
\vspace{0.5em}
\vspace{0.5em}
D\\
• Le Tech Switcher : Ingénieur Web2 confirmé (Java/JS) cherchant à pivoter vers la blockchain et le remote international.\\
\vspace{0.5em}
\vspace{0.5em}
\vspace{0.5em}
\vspace{0.5em}
FI\\
N\\
• Le Validator Architect : Profil infra/SRE visant la certification RBK Solana Validator Track pour opérer des nœuds institutionnels.\\
CO\\
• Le Stratège (Track C) : Profil business/finance souhaitant maîtriser la Tokenomics et la gouvernance DAO.\\
\vspace{0.5em}
(Voir Annexe G pour les critères détaillés d’admission.)\\
\vspace{0.5em}
\vspace{0.5em}
Chiffres Clés \textbackslash{}& Objectifs 2026\\
8\\
Organisation — Principe de conception du programme visant à inculquer dès la formation les réflexes et standards d’un profil senior.\\
\vspace{0.5em}
\vspace{0.5em}
\vspace{0.5em}
\vspace{0.5em}
\begin{confidential}CONFIDENTIEL — PROJET RBK 2.0                                             Page 28 sur 316\end{confidential}
RBK 2.0 — Whitepaper Architecte Web3\\
\vspace{0.5em}
\vspace{0.5em}
\vspace{0.5em}
\vspace{0.5em}
Métrique                     Objectif Alpha 2026\\
Durée du cursus              48 semaines (Genesis Pool 4 + Explorer 12 + Bâtisseur 16 + Architecte 16) + Launchpad optionnel\\
Validation                   Skill Mirror hebdo + Block Check “Launch Proof” (NFT Skills)\\
Modèle économique            15 900 TND TTC/an + forfait Nexus 3 300 TND (prestations pédagogiques complètes) + ISA (15\textbackslash{}%)\\
Taux de placement            80\textbackslash{}% (international, remote, grants structurés)\\
(cible)\\
Salaire moyen de             3 500 TND brut / mois (équivalent \textbackslash{}$)\\
sortie (cible)\\
\vspace{0.5em}
\vspace{0.5em}
\vspace{0.5em}
\vspace{0.5em}
EL\\
Opérateur technique          Nexus Réussite (Maître d’œuvre pédagogique : design, mentors, évaluations) utilisant la plateforme Venture Engine de\\
Money Factory AI (fournisseur technologique)\\
\vspace{0.5em}
\vspace{0.5em}
\vspace{0.5em}
\vspace{0.5em}
TI\\
EN\\
* Les chiffres de placement et salaire sont des objectifs cibles et ne constituent pas une garantie contractuelle.\\
\vspace{0.5em}
L’ISA est un Contrat de Partenariat Pédagogique et de Partage de Succès conclu uniquement entre l’apprenant et RBK ; RBK assure\\
\vspace{0.5em}
\vspace{0.5em}
\vspace{0.5em}
\vspace{0.5em}
D\\
le recouvrement via un fonds de garantie dédié et un prestataire spécialisé (cf. Annexe K).\\
\vspace{0.5em}
\vspace{0.5em}
\vspace{0.5em}
\vspace{0.5em}
FI\\
Appel à l’Action                                                N\\
CO\\
La Genesis Cohort Alpha (30 sièges) ouvre ses tests de sélection en Mars 2026. Droits d’entrée : acompte 25\textbackslash{}% sur 15 900 TND, solde\\
financé via ISA ou subventions publiques. RBK offre ici l’opportunité unique de rejoindre l’élite technologique africaine grâce à l’exécution\\
opérationnelle de Nexus Réussite et à la collaboration stratégique avec la Superteam.\\
\vspace{0.5em}
\vspace{0.5em}
\vspace{0.5em}
\vspace{0.5em}
RBK 2.0 : De Codeur à Architecte. De Local à Global.\\
\vspace{0.5em}
\vspace{0.5em}
\vspace{0.5em}
\vspace{0.5em}
\begin{confidential}CONFIDENTIEL — PROJET RBK 2.0                                                   Page 29 sur 316\end{confidential}
\section{VISION \textbackslash{}& MANIFESTE}
\vspace{0.5em}
\vspace{0.5em}
\vspace{0.5em}
\vspace{0.5em}
EL\\
\subsection{1.1 La Thèse Centrale : Former des Architectes, pas des Codeurs}
Le marché n’a plus besoin de ”développeurs exécutants”. L’IA le fait mieux, plus vite, et moins cher. Ce qui manque cruellement, ce\\
\vspace{0.5em}
\vspace{0.5em}
\vspace{0.5em}
\vspace{0.5em}
TI\\
sont des Architectes de Systèmes Distribués.\\
\vspace{0.5em}
\vspace{0.5em}
\vspace{0.5em}
\vspace{0.5em}
EN\\
Ď Le Manifeste RBK 2.0\\
\vspace{0.5em}
\vspace{0.5em}
\vspace{0.5em}
D\\
Manifeste : RBK 2.0 forge des Architectes Web3 immédiatement opérationnels, capables de concevoir, auditer et sécuriser des\\
\vspace{0.5em}
\vspace{0.5em}
\vspace{0.5em}
\vspace{0.5em}
FI\\
systèmes décentralisés dès leur sortie. Notre promesse : un diplômé RBK possède la rigueur d’un ingénieur senior et la productivité\\
d’une équipe junior assistée par l’IA.\\
N\\
CO\\
Indicateurs Clés de Performance (KPIs)\\
• KPI 1 : Taux de Placement (TP) : 80\textbackslash{}% des diplômés en CDI/CDD de 6+ mois dans le domaine (Web3/Solana/EVM) dans\\
les 6 mois suivant la graduation.\\
\vspace{0.5em}
• KPI 2 : Salaire Moyen de Sortie (SM) : 2500 TND brut/mois (seuil ISA) pour les profils juniors, 20000 TND+ mensuel\\
pour les profils seniors-by-design.\\
\vspace{0.5em}
• KPI 3 : Temps Moyen de Premier Emploi (TPE) : \textless{} 3 mois pour les 80\textbackslash{}% des diplômés.\\
\vspace{0.5em}
\vspace{0.5em}
\vspace{0.5em}
\vspace{0.5em}
30\\
RBK 2.0 — Whitepaper Architecte Web3                                                                         Chapitre 1. Vision \textbackslash{}& Manifeste\\
\vspace{0.5em}
\vspace{0.5em}
\vspace{0.5em}
• KPI 4 : Satisfaction Étudiante (NPS) : \textgreater{} 60 points (échelle : -100 à +100).\\
\vspace{0.5em}
• KPI 5 : Score Technique de Diplôme : \textgreater{} 85/100 basé sur évaluation finale objectivée.\\
\vspace{0.5em}
\vspace{0.5em}
\subsection{1.1.0.1 Définition opérationnelle d’un Architecte Web3}
\vspace{0.5em}
Un Architecte Web3 ne se contente pas d’écrire des smart contracts ; il conçoit des systèmes financiers inarrêtables. Sa responsabilité\\
principale est la gestion du risque. Contrairement au développeur Web2 qui optimise pour la vitesse de livraison, l’architecte Web3\\
optimise pour la sécurité et la résilience (Trust Minimization).\\
Concrètement, un architecte RBK maîtrise :\\
\vspace{0.5em}
\vspace{0.5em}
\vspace{0.5em}
\vspace{0.5em}
EL\\
• Le Design de Protocoles : Définition des invariants économiques et des surfaces d’attaque (Threat Modeling).\\
\vspace{0.5em}
\vspace{0.5em}
\vspace{0.5em}
\vspace{0.5em}
TI\\
• L’Optimisation Bas-Niveau : Gestion fine des Compute Units et du stockage on-chain (PDA Seeds, Merkle Trees).\\
\vspace{0.5em}
\vspace{0.5em}
\vspace{0.5em}
\vspace{0.5em}
EN\\
• Les Patterns de Sécurité : Protection contre les attaques classiques (Re-entrancy, CPI hijacking, Sybil attacks).\\
\vspace{0.5em}
\vspace{0.5em}
\vspace{0.5em}
\vspace{0.5em}
D\\
• L’Observabilité : Capacité à monitorer l’état du système en temps réel (Indexing, RPCs).\\
\vspace{0.5em}
\vspace{0.5em}
\vspace{0.5em}
\vspace{0.5em}
FI\\
N\\
Livrables attendus d’un Architecte RBK :\\
CO\\
• Diagrammes d’architecture (C4 Model) et de flux de données.\\
\vspace{0.5em}
• Rapport de Threat Modeling identifiant les vecteurs d’attaque.\\
\vspace{0.5em}
• Suite de tests exhaustive (Unitaires + Fuzzing + Invariants).\\
\vspace{0.5em}
• Code audité et documenté (NatSpec / RustDoc).\\
\vspace{0.5em}
• Runbook d’incident (Procédure de pause/fixation d’urgence).\\
\vspace{0.5em}
\vspace{0.5em}
\vspace{0.5em}
\vspace{0.5em}
\begin{confidential}CONFIDENTIEL — PROJET RBK 2.0                                Page 31 sur 316\end{confidential}
RBK 2.0 — Whitepaper Architecte Web3                                                                                        Chapitre 1. Vision \textbackslash{}& Manifeste\\
\vspace{0.5em}
\vspace{0.5em}
\subsection{1.1.0.2 Pourquoi le ”code basique” ne suffit plus à l’ère des LLM}
\vspace{0.5em}
L’avènement des LLMs (GPT-4, Claude 3.5 Sonnet) a commoditisé la production de code syntaxique. Générer un ERC-201 ou un\\
programme Anchor standard prend désormais 30 secondes et coûte 0.01\textbackslash{}$. La valeur ajoutée du ”codeur” qui traduit une spec en fonctions\\
s’effondre.\\
Cependant, l’IA ne sait pas raisonner sur l’intention. Elle peut générer un code qui compile parfaitement mais qui contient des failles\\
logiques dévastatrices.\\
\vspace{0.5em}
Le Risque des ”Failles Invisibles” (IA-Generated)\\
1. Hypothèses Non Vérifiées : L’IA suppose que l’utilisateur est honnête, omettant les contrôles d’accès (Missing Access Control).\\
Impact : Vol de fonds.\\
\vspace{0.5em}
\vspace{0.5em}
\vspace{0.5em}
\vspace{0.5em}
EL\\
2. Invariants Économiques : L’IA ne vérifie pas si ‘total\textbackslash{}_minted \textless{}= max\textbackslash{}_supply‘ après un calcul complexe. Impact : Inflation\\
\vspace{0.5em}
\vspace{0.5em}
\vspace{0.5em}
\vspace{0.5em}
TI\\
infinie.\\
\vspace{0.5em}
\vspace{0.5em}
\vspace{0.5em}
\vspace{0.5em}
N\\
3. Edge Cases : L’IA oublie les cas limites (division par zéro, overflow, array vide). Impact : Blocage du protocole (DoS).\\
\vspace{0.5em}
\vspace{0.5em}
\vspace{0.5em}
\vspace{0.5em}
DE\\
C’est pourquoi RBK 2.0 adopte l’approche ”Learning by Auditing”. Nous formons les étudiants à considérer tout code (humain ou IA)\\
\vspace{0.5em}
\vspace{0.5em}
\vspace{0.5em}
\vspace{0.5em}
FI\\
comme potentiellement hostile jusqu’à preuve du contraire.\\
\vspace{0.5em}
\vspace{0.5em}
N\\
CO\\
\subsection{1.1.0.3 Le Mécanisme Senior-by-Design}
\vspace{0.5em}
Comment transformer un profil junior en architecte senior en 48 semaines ? Par un conditionnement intensif en quatre étapes séquen-\\
tielles :\\
1. Genesis Pool (4 semaines) : Bootcamp de sélection où l’apprenant prouve sa résilience, maîtrise les fondamentaux Rust et adopte les\\
rituels Skill Mirror. 2. Explorer (12 semaines) : Consolidation technique et produit ; développement sur Solana/EVM, sécurité applicative,\\
Block Check “Explorer”. 3. Bâtisseur \textbackslash{}& Architecte (16 + 16 semaines) : Spécialisation par tracks, audits croisés, leadership technique\\
et préparation au Launch Proof. 4. Launchpad (optionnel) : Validation de marché via le parcours Nexus (forfait 3 300 TND + partage\\
30\textbackslash{}%).\\
1\\
Web3/Blockchain — Ethereum Request for Comments 20 (ERC-20) : standard définissant l’interface minimale d’un jeton fongible (transfert, approbation, total\\
supply).\\
\vspace{0.5em}
\vspace{0.5em}
\vspace{0.5em}
\begin{confidential}CONFIDENTIEL — PROJET RBK 2.0                                         Page 32 sur 316\end{confidential}
RBK 2.0 — Whitepaper Architecte Web3                                                                           Chapitre 1. Vision \textbackslash{}& Manifeste\\
\vspace{0.5em}
\vspace{0.5em}
\vspace{0.5em}
Mécanisme                 Habitude Créée                     Preuve Tangible\\
Code Review Obligatoire   ”Mon code sera lu par un humain”   Qualité des PRs, Commentaires\\
Fuzzing Systématique      ”Le happy-path ne suffit pas”      Rapports de couverture \textgreater{} 90\textbackslash{}%\\
Threat Modeling           ”Penser comme un attaquant”        Documents d’architecture défensive\\
Démo Publique             ”Je dois défendre mes choix”       Vidéos de pitch, README pro\\
\vspace{0.5em}
\vspace{0.5em}
\vspace{0.5em}
\vspace{0.5em}
\subsection{1.1.0.4 Métriques de Succès et Méthode de Mesure}
\vspace{0.5em}
\vspace{0.5em}
\vspace{0.5em}
\vspace{0.5em}
EL\\
Nous ne vendons pas du rêve, nous vendons des résultats mesurables.\\
\vspace{0.5em}
\vspace{0.5em}
\vspace{0.5em}
\vspace{0.5em}
TI\\
• Taux de placement (3 mois) : Pourcentage des diplômés ayant signé un contrat (CDI, Freelance \textgreater{} 3 mois, ou Grant \textgreater{} 5k\textbackslash{}$) 90\\
jours après la fin du cursus.\\
\vspace{0.5em}
\vspace{0.5em}
\vspace{0.5em}
\vspace{0.5em}
EN\\
• Salaire Moyen de Sortie : Moyenne des rémunérations annualisées (converties en TND), hors equity/tokens non-liquides.\\
\vspace{0.5em}
\vspace{0.5em}
\vspace{0.5em}
\vspace{0.5em}
D\\
• Time-to-First-Revenue : Délai moyen entre le début de la Phase 3 et le premier dollar gagné (souvent via un Bounty Superteam).\\
\vspace{0.5em}
\vspace{0.5em}
\vspace{0.5em}
\vspace{0.5em}
FI\\
N\\
\begin{center}\begin{tabular}TAB. 1.2 : Métriques de Succès RBK 2.0\end{tabular}\end{center}
CO\\
Indicateur     Définition                         Cible     Méthode                         Preuve\\
Placement      Contrat signé ou facture émise     90\textbackslash{}%       Suivi Alumni J+90               Contrats, Relevés\\
Salaire        Revenu net mensuel équivalent     \textgreater{}3k TND    Déclaration sur l’honneur       Fiches de paie\\
Satisfaction   NPS (Net Promoter Score)           \textgreater{}70       Enquête anonyme fin de cursus   Typeform Export\\
Niveau Tech    Score aux tests finaux           \textgreater{}850/1000   Platforme d’examen (LMS)        Certificat On-chain\\
\vspace{0.5em}
\vspace{0.5em}
\vspace{0.5em}
\vspace{0.5em}
\begin{confidential}CONFIDENTIEL — PROJET RBK 2.0                                  Page 33 sur 316\end{confidential}
RBK 2.0 — Whitepaper Architecte Web3                                                                                    Chapitre 1. Vision \textbackslash{}& Manifeste\\
\vspace{0.5em}
\vspace{0.5em}
Entrée (Bases)      IA (x10)        Cyborg 2.0       Réseau (Superteam)     Architecte\\
Méthode \textbackslash{}& Rigueur                    Preuves : Portfolio, Audit\\
\vspace{0.5em}
FIG. 1.1 : La Chaîne de Valeur RBK 2.0\\
\vspace{0.5em}
\vspace{0.5em}
\vspace{0.5em}
\subsection{1.2 Pourquoi RBK 2.0 ?}
RBK 2.0 répond à un décalage structurel entre l’offre éducative classique et les exigences d’architectes Web3 seniors : nous cadrons ici\\
le ”pourquoi” avant de détailler les sous-piliers.\\
\vspace{0.5em}
\subsection{1.2.0.1 Diagnostic : L’Écart de Compétence (Skills Gap)}
\vspace{0.5em}
\vspace{0.5em}
\vspace{0.5em}
\vspace{0.5em}
EL\\
Le fossé entre l’offre de formation classique et la demande du marché Web3 est béant.\\
\vspace{0.5em}
\vspace{0.5em}
\vspace{0.5em}
\vspace{0.5em}
TI\\
1. Évaluation obsolète : Les écoles notent la mémorisation ; le marché paie la résolution de problèmes inconnus.\\
\vspace{0.5em}
\vspace{0.5em}
\vspace{0.5em}
\vspace{0.5em}
EN\\
2. Absence de Sécurité : La sécurité est souvent une option ou un module théorique. En Web3, c’est le prérequis absolu.\\
\vspace{0.5em}
\vspace{0.5em}
\vspace{0.5em}
\vspace{0.5em}
D\\
3. Pas de Production Réelle : Les projets d’école finissent dans un dossier ”brouillon”. Un profil senior doit montrer un code en\\
\vspace{0.5em}
\vspace{0.5em}
\vspace{0.5em}
\vspace{0.5em}
FI\\
production.\\
\vspace{0.5em}
\vspace{0.5em}
N\\
4. Signaux Marché Faibles : Un diplôme papier ne prouve rien à une DAO internationale. Seul le code (GitHub) et la réputation\\
CO\\
(On-chain) comptent.\\
\vspace{0.5em}
\subsection{1.2.0.2 Les Différenciateurs RBK 2.0}
\vspace{0.5em}
• Méthodologie Cyborg 2.0 (voir Chap. 4) : Nous intégrons l’IA comme outil de base, pas comme aide à la triche.\\
\vspace{0.5em}
• Intensité 48 semaines (voir Chap. 7) : Une immersion totale Genesis →Explorer →Bâtisseur →Architecte pour ancrer le mindset\\
“senior-by-design”.\\
\vspace{0.5em}
• Preuve de Travail (Proof of Work) (voir Chap. 11) : Chaque ligne de code contribue à un portfolio public auditable.\\
\vspace{0.5em}
• Réseau Global : Connexion directe avec la Superteam et les opportunités internationales.\\
\vspace{0.5em}
\vspace{0.5em}
\begin{confidential}CONFIDENTIEL — PROJET RBK 2.0                                 Page 34 sur 316\end{confidential}
RBK 2.0 — Whitepaper Architecte Web3                           Chapitre 1. Vision \textbackslash{}& Manifeste\\
\vspace{0.5em}
\vspace{0.5em}
RBK 1.0 vs RBK 2.0 : Impact Business\\
\vspace{0.5em}
\vspace{0.5em}
\vspace{0.5em}
\vspace{0.5em}
EL\\
TI\\
EN\\
D\\
FI\\
N\\
CO\\
\vspace{0.5em}
\vspace{0.5em}
\vspace{0.5em}
\vspace{0.5em}
\begin{confidential}CONFIDENTIEL — PROJET RBK 2.0                Page 35 sur 316\end{confidential}
RBK 2.0 — Whitepaper Architecte Web3                                                                      Chapitre 1. Vision \textbackslash{}& Manifeste\\
\vspace{0.5em}
\vspace{0.5em}
\vspace{0.5em}
Indicateur clé       RBK 1.0 (historique)                RBK 2.0 (modèle Nexus)              Signal pour RBK\\
CA moyen par         Parcours courts modulaires (ticket Offre architecte premium facturée Revenus récurrents structurés sur 12\\
apprenant            moyen \textless{} 2500 TND, forte volatili- 15 900 TND avec jalons Genesis/Ex- mois\\
té bourses)                        plorer/Architecte\\
Taux de placement Empiriquement \textless{} 50\textbackslash{}% sur les co- Cible 65\textbackslash{}% sécurisée par accords Justifie ISA et bonus/malus Nexus\\
6 mois            hortes Web2 (absence d’offres par- Superteam/MFAI et Launchpad op-\\
tenaires dédiées)                  tionnel\\
Marge                Capex lourd internalisé, dépen- Nexus absorbe le coût pédagogique Marge nette positive Cohorte 2 (voir\\
opérationnelle       dance subventions ponctuelles   (R sur delivery), RBK concentre sa Chap. 15)\\
\vspace{0.5em}
\vspace{0.5em}
\vspace{0.5em}
\vspace{0.5em}
EL\\
marge sur commercial et ISA\\
\vspace{0.5em}
\vspace{0.5em}
\vspace{0.5em}
\vspace{0.5em}
TI\\
Coût d’acquisition   Campagnes génériques, CAC non Playbook    performance dédié Visibilité CAC/LTV dans les dash-\\
suivi, dérive \textgreater{} 800 TND       (GROWTH SQUAD, 500 TND cible) boards\\
\vspace{0.5em}
\vspace{0.5em}
\vspace{0.5em}
\vspace{0.5em}
EN\\
+ budget 3 000 TND\\
\vspace{0.5em}
\vspace{0.5em}
\vspace{0.5em}
\vspace{0.5em}
D\\
Protection de        Qualité hétérogène selon mentors, SLA pédagogiques, audits conjoints RBK garde la signature et valide\\
\vspace{0.5em}
\vspace{0.5em}
\vspace{0.5em}
\vspace{0.5em}
FI\\
marque               absence de normes QA              et protocole anti-burnout          chaque livrable\\
Positionnement\\
N\\
Bootcamp généraliste perçu « full- École d’Architectes Web3 senior-by- Renforce le capital réputation RBK\\
CO\\
marché               stack junior »                     design avec preuves on-chain\\
Relation             Relation opportuniste,    peu   de Account management structuré (Ta- Fidélise les partenaires corporate RBK\\
employeurs           comptes stratégiques               lent Partners, Career Studio) + re-\\
porting mensuel\\
Contrôle financier   Flux mentors internalisés peu tra- RBK encaisse 100\textbackslash{}% et redistribue Traçabilité et auditabilité renforcées\\
çables, ISA ad hoc                 selon jalons, fonds ISA séquestré\\
150 000 TND\\
\vspace{0.5em}
\vspace{0.5em}
\vspace{0.5em}
\vspace{0.5em}
\begin{confidential}CONFIDENTIEL — PROJET RBK 2.0                               Page 36 sur 316\end{confidential}
RBK 2.0 — Whitepaper Architecte Web3                                                                       Chapitre 1. Vision \textbackslash{}& Manifeste\\
\vspace{0.5em}
\vspace{0.5em}
\vspace{0.5em}
Gestion du risque    Pas de clauses de sortie formalisées Clauses go/no-go 4 semaines, Risque opérationnel borné contrac-\\
droit de suspension Nexus, comité tuellement\\
risques\\
Expérience           Variation forte selon batch, pas de NPS piloté (\textgreater{}60), Block Checks, Assure cohérence expérience RBK\\
apprenants           mesure systématique                 Skill Mirror, support MFAI\\
Notoriété            Visibilité locale principalement    Intégration Superteam, showcases Élargit pipeline candidats et investis-\\
internationale                                           Solana, partenariats DePIN       seurs\\
Monétisation long    Peu de services post-formation      Launchpad, Talent Services, Corpo- Nouveaux profit centers pour RBK\\
terme                                                    rate Seats 15 900 TND\\
\vspace{0.5em}
\vspace{0.5em}
\vspace{0.5em}
\vspace{0.5em}
EL\\
Pilotage             Processus informels, décisions réac- Comités CEC/Qualité, RACI clarifié, RBK arbitre, Nexus exécute sous\\
\vspace{0.5em}
\vspace{0.5em}
\vspace{0.5em}
\vspace{0.5em}
TI\\
gouvernance          tives                                reporting KPI partagé               contrôle\\
\vspace{0.5em}
\vspace{0.5em}
\vspace{0.5em}
\vspace{0.5em}
EN\\
Preuves              Dossiers internes non mutualisés    Dashboards MFAI, repos Git véri- Démonstration objective pour le board\\
d’exécution                                              fiés, certificats on-chain       RBK\\
\vspace{0.5em}
\vspace{0.5em}
\vspace{0.5em}
\vspace{0.5em}
D\\
Alignement           Offre 1.0 centrée hardware/networ- Vision RBK 2.0 alignée mission d’ex- Narratif cohérent pour les dirigeants\\
\vspace{0.5em}
\vspace{0.5em}
\vspace{0.5em}
\vspace{0.5em}
FI\\
stratégique          king                               cellence tunisienne + expansion RBK\\
\vspace{0.5em}
N         Web3\\
CO\\
Continuité service   Risque rupture si équipe part       Plan de succession mentors, docu- Assure résilience opérationnelle\\
mentation, shadowing Nexus/RBK\\
\vspace{0.5em}
\vspace{0.5em}
Cette approche protège RBK contre toute dilution de marque tout en maximisant la valeur vie apprenant : la qualité reste sous gouver-\\
nance RBK (validation comités), la promesse marché est renforcée par des preuves auditées, et chaque clause contractuelle prévoit une\\
sortie maîtrisée.\\
\vspace{0.5em}
\vspace{0.5em}
\vspace{0.5em}
\vspace{0.5em}
\begin{confidential}CONFIDENTIEL — PROJET RBK 2.0                               Page 37 sur 316\end{confidential}
RBK 2.0 — Whitepaper Architecte Web3                                                                        Chapitre 1. Vision \textbackslash{}& Manifeste\\
\vspace{0.5em}
\vspace{0.5em}
Bouclier de Marque RBK\\
\vspace{0.5em}
Ownership relation apprenant : RBK conserve la signature contractuelle, l’encaissement et la communication officielle ; Nexus\\
opère en marque blanche sous SLA qualité. Standards éditoriaux unifiés : tout contenu pédagogique ou marketing passe par un\\
comité RBK, garantissant cohérence de ton et conformité compliance (voir Chap. 18). Traçabilité des décisions : procès-verbaux\\
des comités, audits trimestriels et Skill Mirror créent un historique opposable en cas de litige ou de due diligence investisseurs.\\
\vspace{0.5em}
\vspace{0.5em}
\subsection{1.2.0.3 Ce que RBK 2.0 n’est pas}
\vspace{0.5em}
Il est crucial d’aligner les attentes. RBK 2.0 n’est :\\
\vspace{0.5em}
\vspace{0.5em}
\vspace{0.5em}
\vspace{0.5em}
EL\\
• Pas un cours vidéo passif : L’apprentissage se fait par la pratique douloureuse et gratifiante (Hard Fun).\\
\vspace{0.5em}
• Pas un bootcamp JavaScript : Nous formons des ingénieurs système (Rust/Solidity), pas des développeurs frontend React (bien\\
\vspace{0.5em}
\vspace{0.5em}
\vspace{0.5em}
\vspace{0.5em}
TI\\
que ce soit une compétence annexe).\\
\vspace{0.5em}
\vspace{0.5em}
\vspace{0.5em}
\vspace{0.5em}
EN\\
• Pas une promesse magique : L’ISA et le placement dépendent à 100\textbackslash{}% de l’engagement de l’étudiant.\\
\vspace{0.5em}
\vspace{0.5em}
\vspace{0.5em}
\vspace{0.5em}
D\\
Positionnement Stratégique\\
\vspace{0.5em}
\vspace{0.5em}
\vspace{0.5em}
\vspace{0.5em}
FI\\
RBK 2.0 est une ”School of Engineering” accélérée, positionnée entre le bootcamp d’élite (type 42) et l’incubateur de startups Web3.\\
\vspace{0.5em}
N\\
CO\\
\vspace{0.5em}
\vspace{0.5em}
\vspace{0.5em}
\vspace{0.5em}
\begin{confidential}CONFIDENTIEL — PROJET RBK 2.0                               Page 38 sur 316\end{confidential}
RBK 2.0 — Whitepaper Architecte Web3                                                                               Chapitre 1. Vision \textbackslash{}& Manifeste\\
\vspace{0.5em}
\vspace{0.5em}
\subsection{1.2.0.4 Changement de Paradigme}
\vspace{0.5em}
\vspace{0.5em}
\begin{center}\begin{tabular}TAB. 1.5 : Le Changement de Paradigme RBK 2.0 (Détaillé)\end{tabular}\end{center}
\vspace{0.5em}
\vspace{0.5em}
Dimension     Ancien Monde (Univ/Bootcamps)      RBK 2.0 (Senior-by-Design)         Signal Recruteur\\
\vspace{0.5em}
Objectif      Valider des modules                Livrer de la valeur                GitHub Activity\\
Outils        Interdits (Pas d’IA)               Obligatoires (Cursor, Copilot)     Vitesse d’exécution\\
Rythme        Linéaire, théorique                Cyclique, intense, pratique        Résilience\\
\vspace{0.5em}
\vspace{0.5em}
\vspace{0.5em}
\vspace{0.5em}
EL\\
Sécurité      Optionnelle / Théorique            By Design (Audit Flow)             Portfolio d’audits\\
Santé         Ignorée                            Gérée (Protocole Anti-Burnout)     Stabilité émotionnelle\\
\vspace{0.5em}
\vspace{0.5em}
\vspace{0.5em}
\vspace{0.5em}
TI\\
Sortie        Stage sous-payé                    Consultance / CDI Senior / Grant   Contrats signés\\
\vspace{0.5em}
\vspace{0.5em}
\vspace{0.5em}
\vspace{0.5em}
D   EN\\
FI\\
N\\
CO\\
\vspace{0.5em}
\vspace{0.5em}
\vspace{0.5em}
\vspace{0.5em}
\begin{confidential}CONFIDENTIEL — PROJET RBK 2.0                                  Page 39 sur 316\end{confidential}
\section{Note de cadrage — RBK 2.0 (Web3 Studio Dual-Track}
EVM \textbackslash{}& Solana)\\
\vspace{0.5em}
\vspace{0.5em}
\vspace{0.5em}
\vspace{0.5em}
EL\\
TI\\
Avant-propos stratégique\\
\vspace{0.5em}
Le passage à RBK 2.0 marque une rupture avec les modèles éducatifs traditionnels. En pivotant vers une architecture de Web3 Build\\
\vspace{0.5em}
\vspace{0.5em}
\vspace{0.5em}
\vspace{0.5em}
EN\\
Studio hybride (EVM \textbackslash{}& Solana), nous ne formons plus de simples développeurs, mais des architectes de la souveraineté numérique.\\
Cette note de cadrage formalise l’ambition, les moyens et le pilotage de cette transformation structurante.\\
\vspace{0.5em}
\vspace{0.5em}
\vspace{0.5em}
\vspace{0.5em}
D\\
FI\\
\subsection{2.1 Contexte du projet                                    N}
CO\\
Cette section situe le pivot RBK 2.0 dans son environnement actuel avant d’entrer dans le détail des enjeux et des rôles.\\
\vspace{0.5em}
\vspace{0.5em}
\subsection{2.1.1 Situation actuelle}
RBK (ReBootKamp) s’est imposé comme un acteur clé de la formation intensive en développement logiciel. Toutefois, l’écosystème\\
technologique mondial subit une mutation profonde avec l’avènement du Web3, de la finance décentralisée (DeFi1 ) et des architectures\\
distribuées. La version 1.0 du modèle a prouvé son efficacité pédagogique, mais doit désormais intégrer une dimension industrielle et\\
1\\
Web3/Blockchain — Decentralized Finance (DeFi) : finance décentralisée construite sur des smart contracts, offrant prêts, échanges et rendements sans intermé-\\
diaire bancaire traditionnel.\\
\vspace{0.5em}
\vspace{0.5em}
\vspace{0.5em}
\vspace{0.5em}
40\\
RBK 2.0 — Whitepaper Architecte Web3                                                                                  Chapitre 2. Stratégie : Note de Cadrage\\
\vspace{0.5em}
\vspace{0.5em}
commerciale directe. L’environnement actuel se caractérise par une forte demande de profils ”Seniors-by-Design” capables d’opérer sur\\
des environnements de production critiques (Solana, Ethereum), face à une pénurie structurelle de talents vérifiés.\\
\vspace{0.5em}
\vspace{0.5em}
\subsection{2.1.2 Motivations et enjeux}
La motivation principale est de repositionner RBK non plus comme une école, mais comme un Talent/Venture Studio. L’enjeu est triple :\\
\vspace{0.5em}
1. Technologique : Maîtriser le double standard du marché (Rust/Solana pour la performance, Solidity/EVM pour l’interopérabilité).\\
\vspace{0.5em}
2. Économique : Sécuriser le modèle via les ISA et la production de valeur tangible (MVP2 , audits) durant la formation.\\
\vspace{0.5em}
3. Humain : Offrir une employabilité internationale immédiate aux talents tunisiens et africains.\\
\vspace{0.5em}
\vspace{0.5em}
\vspace{0.5em}
\vspace{0.5em}
EL\\
TI\\
\subsection{2.2 Enjeux (Stratégiques, Opérationnels, Réglementaires)}
\vspace{0.5em}
\vspace{0.5em}
\vspace{0.5em}
\vspace{0.5em}
EN\\
Légende des rôles (RACI3 ).\\
\vspace{0.5em}
\vspace{0.5em}
\vspace{0.5em}
\vspace{0.5em}
D\\
• CEO : Direction RBK (décisions stratégiques).\\
\vspace{0.5em}
\vspace{0.5em}
\vspace{0.5em}
\vspace{0.5em}
FI\\
• HoP : Head of Program (qualité, curriculum).\\
N\\
CO\\
• TechLeads : Référents techniques (EVM/Solana).\\
\vspace{0.5em}
• SecLead : Référent sécurité/audit.\\
\vspace{0.5em}
• Ops : Opérations \textbackslash{}& Career : Carrière.\\
2\\
Produit/API/Architecture — Minimum Viable Product (MVP) : produit minimum viable permettant de valider une hypothèse avec l’effort minimal avant itérations.\\
3\\
Web3/Blockchain — Decentralized Finance (DeFi) : finance décentralisée construite sur des smart contracts, offrant prêts, échanges et rendements sans intermé-\\
diaire bancaire traditionnel.\\
\vspace{0.5em}
\vspace{0.5em}
\vspace{0.5em}
\vspace{0.5em}
\begin{confidential}CONFIDENTIEL — PROJET RBK 2.0                                           Page 41 sur 316\end{confidential}
RBK 2.0 — Whitepaper Architecte Web3                                                                 Chapitre 2. Stratégie : Note de Cadrage\\
\vspace{0.5em}
\vspace{0.5em}
\vspace{0.5em}
\subsection{2.3 Matrice SWOT (Synthèse)}
SWOT — Analyse Stratégique RBK 2.0\\
\vspace{0.5em}
\vspace{0.5em}
Forces                                                           Faiblesses\\
• Positionnement premium : “Senior-by-Design”.                    • Risque de sur-ambition (contenu dense).\\
• DualTrack EVM/Solana.                                           • Dépendance à des experts rares.\\
• Méthodologie Studio (PR reviews, CI).                           • Exigence élevée (sélectivité).\\
• Production d’un portfolio vérifiable.                           • Nécessité d’une cohérence chiffrée stricte.\\
• Réseau mentors (diaspora, builders).                            • Sensibilité réglementaire locale.\\
\vspace{0.5em}
\vspace{0.5em}
\vspace{0.5em}
\vspace{0.5em}
EL\\
Opportunités                                                     Menaces\\
\vspace{0.5em}
\vspace{0.5em}
\vspace{0.5em}
\vspace{0.5em}
TI\\
• Marché remote Web3 global.                                      • Volatilité Web3 (cycles marché).\\
\vspace{0.5em}
\vspace{0.5em}
\vspace{0.5em}
\vspace{0.5em}
EN\\
• Demande profils sécurité / audit.                               • Risque réputationnel (confusion crypto/trading).\\
• Partenariats protocols/infra.                                   • Risques sécurité (mauvaises pratiques).\\
\vspace{0.5em}
\vspace{0.5em}
\vspace{0.5em}
\vspace{0.5em}
D\\
• Hub régional (Afrique du Nord).                                 • Concurrence internationale.\\
\vspace{0.5em}
\vspace{0.5em}
\vspace{0.5em}
\vspace{0.5em}
FI\\
• Opportunités B2B \textbackslash{}& Studio.                                      • Incertitudes réglementaires.\\
\vspace{0.5em}
\vspace{0.5em}
N\\
CO\\
SWOT dynamique et trajectoire d’atténuation\\
\vspace{0.5em}
Tableau dynamique SWOT (90 jours \textbackslash{}& 12 mois)\\
\vspace{0.5em}
\vspace{0.5em}
Axe              Situation Q1 2025                     Cap 90 jours                           Cap 12 mois\\
Forces           DualTrack industrialisé sur 3 modules Finaliser ”Golden Repos” \textbackslash{}& lancer QA   Transition RBKops →RBK autonome\\
Genesis/Explorer, mentors seniors ali- labs (actions A-010/011)              (shadowing Nexus, handover mentors)\\
gnés\\
\vspace{0.5em}
\vspace{0.5em}
\vspace{0.5em}
\vspace{0.5em}
\begin{confidential}CONFIDENTIEL — PROJET RBK 2.0                               Page 42 sur 316\end{confidential}
RBK 2.0 — Whitepaper Architecte Web3                                                                                   Chapitre 2. Stratégie : Note de Cadrage\\
\vspace{0.5em}
\vspace{0.5em}
\vspace{0.5em}
Faiblesses          Dépendance experts critiques et sylla- Lancer pool remplaçants \textbackslash{}& bridge skills              Internaliser 60\textbackslash{}% des tech leads via par-\\
bus dense                              (actions A-013/016)                                  cours ”TL Academy” \textbackslash{}& déployer moni-\\
toring burn-out\\
Opportunités        Pipeline Superteam/MFAI ouvert (9           Convertir 2 LOI en contrats cadres \textbackslash{}& ou- Atteindre 20 placements remote, 3 caps-\\
partenaires actifs)                         vrir 1 corporate seat pilote             tones DePIN opérés \textbackslash{}& lancer Launchpad\\
V1\\
Menaces             Sensibilité réputationnelle (crypto/tra- Charte marketing + disclaimers signés              Assurance RC pro active, fonds ISA sé-\\
ding) \textbackslash{}& choc marché                      (A-027) \textbackslash{}& comité risques mensuel                   questré \textbackslash{}& PRA complet testé (annexe\\
PCA)\\
\vspace{0.5em}
Indicateurs de suivi : revue hebdomadaire des actions A-010, A-013, A-027 (Annexe V), NPS cohorte \textgreater{} 60, backlog Block Checks \textless{} 5\textbackslash{}%, et avancement handover\\
\vspace{0.5em}
\vspace{0.5em}
\vspace{0.5em}
\vspace{0.5em}
EL\\
RBK ≥ 50\textbackslash{}% à M+12.\\
\vspace{0.5em}
\vspace{0.5em}
\vspace{0.5em}
\vspace{0.5em}
TI\\
N\\
DE\\
\subsection{2.4 Priorisation MoSCoW (Fonctionnalités Clés)}
\vspace{0.5em}
\vspace{0.5em}
\vspace{0.5em}
\vspace{0.5em}
FI\\
Priorisation MoSCoW — Fonctionnalités \textbackslash{}& Exigences\\
\vspace{0.5em}
\vspace{0.5em}
\vspace{0.5em}
\vspace{0.5em}
N\\
Must have                             Should have\\
CO                          Could have                              Won’t have (for now)\\
• Programme DualTrack.                • Micro-certifications (Badges).       • Incident Drills hebdos.               • Trading / Spéculation.\\
• Méthode Studio obligatoire.         • Réseau Mentors structuré.            • Module Security Avancé.               • Contournement fiscal.\\
• Capstones (3) + Projet Final.       • Outillage standardisé (Templates).   • Offre B2B Corporate.                  • Contenu ”Fluff” sans pratique.\\
• Note de cadrage remplie.            • Listes automatiques (Figures/Acro-   • Incubation légère (post-demo).        • Mainnet non audité obligatoire.\\
• Factsheet unique.                   nymes).                                • Scénarios de scalabilité.\\
• Conformité/Éthique (Disclaimers).   • Bibliographie sourcée.\\
• Employabilité (Portfolio).\\
\vspace{0.5em}
\vspace{0.5em}
\vspace{0.5em}
\vspace{0.5em}
\subsection{2.5 Objectifs SMART}
\vspace{0.5em}
\vspace{0.5em}
\begin{confidential}CONFIDENTIEL — PROJET RBK 2.0                                          Page 43 sur 316\end{confidential}
RBK 2.0 — Whitepaper Architecte Web3                                                                                          Chapitre 2. Stratégie : Note de Cadrage\\
\vspace{0.5em}
\vspace{0.5em}
\vspace{0.5em}
Le projet RBK 2.0 s’articule autour d’objectifs précis garantissant la qualité, la maîtrise des coûts et le respect des délais.\\
\vspace{0.5em}
Tableau de bord - Objectifs SMART\\
\vspace{0.5em}
\vspace{0.5em}
\vspace{0.5em}
Indicateur                                Définition                                Cible                                    Mesure / Source\\
Objectif                                  Définition SMART                          Cible                                    KPI / Mesure\\
Excellence Tech                           Former des profils capables de dé-        100\textbackslash{}% de réussite aux Capstones           Taux de déploiement Mainnet\\
ployer en mainnet\\
Employabilité                             Placement des talents sous 90 jours       \textgreater{} 90\textbackslash{}%                                    Taux d’insertion Pro\\
post-graduation\\
Rentabilité                               ROI positif des cohortes via ISA et       Break-even à M+12                        ARR / Cohorte\\
\vspace{0.5em}
\vspace{0.5em}
\vspace{0.5em}
\vspace{0.5em}
EL\\
Factory\\
Vélocité                                  Durée de formation optimisée sans         48 semaines (Genesis Pool 4 + Ex-        Time-to-Skills\\
\vspace{0.5em}
\vspace{0.5em}
\vspace{0.5em}
\vspace{0.5em}
TI\\
perte de qualité                          plorer 12 + Bâtisseur 16 + Archi-\\
\vspace{0.5em}
\vspace{0.5em}
\vspace{0.5em}
\vspace{0.5em}
N\\
tecte 16)\\
\vspace{0.5em}
\vspace{0.5em}
\vspace{0.5em}
\vspace{0.5em}
DE\\
FI\\
ROI et Métriques Financières\\
\vspace{0.5em}
\vspace{0.5em}
N\\
CO\\
L’hypothèse de ROI repose sur une valorisation moyenne des profils sortants à \textbackslash{}textasciitilde{}150 k TND/an (marché international) et sur le recouvrement ISA (net des pauses).\\
Le modèle ”Factory” (production de MVP pour tiers) génère un revenu complémentaire estimé à 15\textbackslash{}% du CA global.\\
\vspace{0.5em}
\vspace{0.5em}
\vspace{0.5em}
\subsection{2.6 Analyse des besoins}
Nous clarifions ici ce qui manque aujourd’hui (tech, outils, expérience utilisateur) pour rendre le studio pleinement opérationnel.\\
\vspace{0.5em}
\vspace{0.5em}
\subsection{2.6.1 Audit de l’existant}
L’infrastructure actuelle (locaux, connexion, serveurs) est robuste pour du développement Web2 classique. Le passage au Web3 nécessite une mise à niveau :\\
• Infrastructure Node : Nécessité de nœuds RPC privés ou dédiés pour les tests de charge.\\
\vspace{0.5em}
\vspace{0.5em}
\vspace{0.5em}
\begin{confidential}CONFIDENTIEL — PROJET RBK 2.0                                                Page 44 sur 316\end{confidential}
RBK 2.0 — Whitepaper Architecte Web3                                                                                         Chapitre 2. Stratégie : Note de Cadrage\\
\vspace{0.5em}
\vspace{0.5em}
• Sécurité : Environnements isolés (Sandbox) pour les manipulations de smart contracts4 .\\
\vspace{0.5em}
\vspace{0.5em}
\subsection{2.6.2 Besoins Utilisateurs et Parties Prenantes}
• Apprenants (Talents) : Recherchent une expertise rare, un mentorat de haut niveau et une insertion rapide.\\
• Partenaires (Hiring Partners) : Recherchent des profils ”Plug \textbackslash{}& Play”, auditables via leur code sur GitHub.\\
• Instruction Team : A besoin d’outils de suivi automatisé (CI/CD5 pédagogique) et de supports à jour.\\
\vspace{0.5em}
Priorisation des Fonctionnalités - MVP Studio\\
\vspace{0.5em}
\vspace{0.5em}
\vspace{0.5em}
Must have                               Should have                             Could have                               Won’t have (for now)\\
• Cursus Dual-Track complet             • Module Audit Sécurité Avancé          • Hackathons Internationaux              • Token de gouvernance DAO (V2)\\
\vspace{0.5em}
\vspace{0.5em}
\vspace{0.5em}
\vspace{0.5em}
EL\\
• Plateforme d’évaluation auto.         • Certification Soulbound (SBT)         • Incubateur physique dédié              • Expansion multi-sites\\
\vspace{0.5em}
\vspace{0.5em}
\vspace{0.5em}
\vspace{0.5em}
TI\\
• Noeuds RPC Testnet\\
\vspace{0.5em}
\vspace{0.5em}
\vspace{0.5em}
\vspace{0.5em}
N\\
DE\\
\subsection{2.7 Solutions techniques}
\vspace{0.5em}
\vspace{0.5em}
\vspace{0.5em}
FI\\
N\\
Les choix technologiques proposés ci-dessous concrétisent le dual-track et sécurisent la montée en échelle.\\
CO\\
\subsection{2.7.1 Justification du Dual-Track (EVM + Solana)}
Le choix de couvrir à la fois l’EVM (Ethereum Virtual Machine) et Solana répond à une logique de couverture de marché totale.\\
• EVM (Solidity/Foundry) : Standard industriel, essentiel pour la DeFi institutionnelle et l’interopérabilité (L2s).\\
• Solana (Rust/Anchor) : Performance extrême, essentiel pour les applications grand public (DePIN, Payments) et l’innovation haute fréquence.\\
4\\
Web3/Blockchain — Smart contract : programme déployé sur blockchain exécutant automatiquement des règles vérifiables, immuables et auditées par le réseau.\\
5\\
Produit/API/Architecture — Continuous Integration / Continuous Delivery (CI/CD) : chaîne d’intégration et de livraison continues automatisant tests, builds,\\
audits et déploiements.\\
\vspace{0.5em}
\vspace{0.5em}
\vspace{0.5em}
\vspace{0.5em}
\begin{confidential}CONFIDENTIEL — PROJET RBK 2.0                                              Page 45 sur 316\end{confidential}
RBK 2.0 — Whitepaper Architecte Web3                                                                                      Chapitre 2. Stratégie : Note de Cadrage\\
\vspace{0.5em}
\vspace{0.5em}
\subsection{2.7.2 Architecture Technique du Studio}
L’architecture du studio repose sur des principes de ”DevOps-first”. Chaque apprenant opère dans un conteneurisé, avec des pipelines CI/CD imposant des standards\\
de qualité (linting, tests unitaires, couverture). L’observabilité est assurée par un dashboard centralisé suivant la progression des compétences (Skill Tree).\\
\vspace{0.5em}
\vspace{0.5em}
\vspace{0.5em}
\subsection{2.8 Évaluation des risques}
La gestion des risques est intégrée dès la conception du programme (”Risk-by-Design”).\\
\vspace{0.5em}
Registre des Risques Prioritaires\\
\vspace{0.5em}
\vspace{0.5em}
\vspace{0.5em}
Risque                                 P   I   Mesures d’atténuation                             Owner\\
\vspace{0.5em}
\vspace{0.5em}
\vspace{0.5em}
\vspace{0.5em}
EL\\
Saturation cognitive des apprenants    4   5   Coaching mental, Pauses actives, Suivi psy        Chief Happiness\\
\vspace{0.5em}
\vspace{0.5em}
\vspace{0.5em}
\vspace{0.5em}
TI\\
Obsolescence technique rapide          5   4   Veille hebdo, Mises à jour syllabus en continu    Tech Lead\\
Défaut de paiement ISA                 3   4   Sélection rigoureuse, Cadre juridique fort        Legal / Finance\\
\vspace{0.5em}
\vspace{0.5em}
\vspace{0.5em}
\vspace{0.5em}
N\\
DE\\
Bugs critiques en prod (Factory)       2   5   Audits croisés, Bounties, Assurance               QA Lead\\
\vspace{0.5em}
\vspace{0.5em}
\vspace{0.5em}
\vspace{0.5em}
FI\\
Légende : P = Probabilité (1-5), I = Impact (1-5).\\
\vspace{0.5em}
\vspace{0.5em}
\vspace{0.5em}
N\\
CO\\
\subsection{2.9 Gouvernance \textbackslash{}& pilotage}
Nous posons les mécanismes de pilotage et de coordination nécessaires pour exécuter le plan sans dette organisationnelle.\\
\vspace{0.5em}
\vspace{0.5em}
\subsection{2.9.1 Méthodologie}
Le pilotage suit une approche Agile/Scrum adaptée à la pédagogie. Des sprints de 2 semaines rythment l’apprentissage et la production.\\
\vspace{0.5em}
\vspace{0.5em}
\vspace{0.5em}
\vspace{0.5em}
\begin{confidential}CONFIDENTIEL — PROJET RBK 2.0                                            Page 46 sur 316\end{confidential}
RBK 2.0 — Whitepaper Architecte Web3                                                                                      Chapitre 2. Stratégie : Note de Cadrage\\
\vspace{0.5em}
\vspace{0.5em}
\subsection{2.9.2 Addendum — Gouvernance contractuelle}
\vspace{0.5em}
Note sur la gouvernance étendue\\
\vspace{0.5em}
La gouvernance de l’exécution et du partenariat est détaillée de manière contractuelle dans les documents suivants :\\
• Le chapitre 3 sur la Gouvernance, Partenariat et RACI.\\
• L’annexe T sur le Modèle Piscine \textbackslash{}& Staffing.\\
• Les addenda aux annexes financières (Annexe C) et juridiques (Annexe D).\\
\vspace{0.5em}
\vspace{0.5em}
\vspace{0.5em}
\vspace{0.5em}
\subsection{2.10 Matrice RACI (Responsabilités Macro)}
\vspace{0.5em}
\vspace{0.5em}
\vspace{0.5em}
\vspace{0.5em}
EL\\
Matrice RACI — Rôles \textbackslash{}& Responsabilités\\
\vspace{0.5em}
\vspace{0.5em}
\vspace{0.5em}
\vspace{0.5em}
TI\\
N\\
Activité                               R             A      C                    I\\
\vspace{0.5em}
\vspace{0.5em}
\vspace{0.5em}
\vspace{0.5em}
DE\\
Définir la vision/positionnement       CEO           CEO    HoP, Legal           Ops\\
\vspace{0.5em}
\vspace{0.5em}
\vspace{0.5em}
\vspace{0.5em}
FI\\
Figer la “Factsheet”                   HoP           CEO    Ops, Legal           Students\\
\vspace{0.5em}
\vspace{0.5em}
\vspace{0.5em}
\vspace{0.5em}
N\\
Concevoir les Tracks (EVM/Solana)      TechLeads     HoP    SecLead              CEO\\
CO\\
Définir rubrics \textbackslash{}& Capstones            SecLead       HoP    TechLeads            CEO\\
Admissions \textbackslash{}& Sélection                 Mkt           CEO    HoP, TechLeads       Ops\\
Encadrement hebdo (Sprints)            HoP           HoP    TechLeads, Mentors   Ops\\
Conformité \textbackslash{}& Légal                     Legal         CEO    HoP                  Students\\
Demo Day \textbackslash{}& Carrière                    Career        CEO    HoP, Mentors         Ops\\
\vspace{0.5em}
\vspace{0.5em}
\vspace{0.5em}
\vspace{0.5em}
\subsection{2.11 Planification, budget, indicateurs}
Ce bloc synthétise la trajectoire temporelle, les ressources et les KPI pour suivre l’exécution du cadrage.\\
\vspace{0.5em}
\vspace{0.5em}
\vspace{0.5em}
\vspace{0.5em}
\begin{confidential}CONFIDENTIEL — PROJET RBK 2.0                                             Page 47 sur 316\end{confidential}
RBK 2.0 — Whitepaper Architecte Web3                                                                                    Chapitre 2. Stratégie : Note de Cadrage\\
\vspace{0.5em}
\vspace{0.5em}
\subsection{2.11.1 Rétroplanning des Grands Jalons}
• Mois 1-2 : Finalisation Ingénierie Pédagogique \textbackslash{}& Recrutement Staff.\\
• Mois 3 : Lancement Campagne Candidats (Marketing).\\
• Mois 4 : Sélection \textbackslash{}& Bootcamps pré-rentrée.\\
• Mois 5 : KICK-OFF Cohorte \textbackslash{}#1 (S0).\\
• Mois 11 : Demo Day \textbackslash{}& Graduation.\\
\vspace{0.5em}
\vspace{0.5em}
\subsection{2.11.2 Budget et Coûts}
Le budget prévisionnel distingue les CAPEX (Infrastructure matériel, Contenu propriétaire) des OPEX (Salaires staff, Marketing, Cloud). Une provision pour risque\\
\vspace{0.5em}
\vspace{0.5em}
\vspace{0.5em}
\vspace{0.5em}
EL\\
de 10\textbackslash{}% est intégrée.\\
\vspace{0.5em}
\vspace{0.5em}
\vspace{0.5em}
\vspace{0.5em}
TI\\
N\\
\subsection{2.12 Conduite du changement}
\vspace{0.5em}
\vspace{0.5em}
\vspace{0.5em}
\vspace{0.5em}
DE\\
La transformation vers RBK 2.0 demande un accompagnement soutenu :\\
\vspace{0.5em}
\vspace{0.5em}
\vspace{0.5em}
\vspace{0.5em}
FI\\
• Formation des formateurs : Montée en compétence obligatoire sur Rust et Solidity Avancé.\\
\vspace{0.5em}
\vspace{0.5em}
\vspace{0.5em}
N\\
• Communication : Clarifier le passage d’une ”école de code” à un ”Centre d’Excellence Web3”.\\
CO\\
• Adhésion : Impliquer les anciens (Alumni) comme mentors pour faciliter la transition culturelle.\\
\vspace{0.5em}
\vspace{0.5em}
\vspace{0.5em}
\subsection{2.13 Conclusion de la Note}
Cette note de cadrage valide la faisabilité et la pertinence du pivot RBK 2.0. En alignant l’excellence technique sur la réalité du marché Web3, RBK se dote d’un\\
avantage concurrentiel décisif. La structure Dual-Track, soutenue par une gouvernance rigoureuse et une gestion des risques proactive, assure la pérennité du modèle.\\
\vspace{0.5em}
w Recommandation\\
Il est recommandé de VALIDER ce cadrage et de lancer immédiatement la phase d’exécution (Recrutement Staff Technique \textbackslash{}& Préparation Infrastructure).\\
\vspace{0.5em}
\vspace{0.5em}
\vspace{0.5em}
\vspace{0.5em}
\begin{confidential}CONFIDENTIEL — PROJET RBK 2.0                                            Page 48 sur 316\end{confidential}
\section{Gouvernance, Partenariat et RACI}
\vspace{0.5em}
\vspace{0.5em}
Ce chapitre détaille l’architecture du partenariat, les rôles des entités impliquées et la répartition des responsabilités (RACI) pour l’exécution du programme RBK\\
\subsection{2.0. Il constitue un addendum contractuel au présent manifeste.}
\vspace{0.5em}
\vspace{0.5em}
\vspace{0.5em}
\vspace{0.5em}
EL\\
TI\\
\subsection{3.1 Cadre général du partenariat}
\vspace{0.5em}
\vspace{0.5em}
\vspace{0.5em}
\vspace{0.5em}
N\\
DE\\
\subsection{3.2 Architecture détaillée du partenariat}
\vspace{0.5em}
\vspace{0.5em}
\vspace{0.5em}
\vspace{0.5em}
FI\\
Architecture Détaillée du Partenariat\\
\vspace{0.5em}
\vspace{0.5em}
\vspace{0.5em}
N\\
CO\\
Entité            Apports Clés                                                           Documents Contractuels Clés\\
\vspace{0.5em}
RBK               Marque, conformité, relation client, encaissements, contrats ap- Contrat-cadre RBK ↔ Nexus (MSA) : « Prestations d’ingénie-\\
(l’enseigne et    prenants (y compris ISA). Apports : locaux, logistique, admi- rie pédagogique \textbackslash{}& pilotage ». Bon de commande / SOW (RBK\\
l’infrastruc-     nistration, CRM, budget acquisition 3 000 TND/mois, abonne- ↔ Nexus) : périmètre cohorte 1 (durée, livrables, KPIs, staf-\\
ture)             ments logiciels “tiers”.                                         fing, plafonds d’heures).\\
\vspace{0.5em}
\vspace{0.5em}
\vspace{0.5em}
\vspace{0.5em}
49\\
RBK 2.0 — Whitepaper Architecte Web3                                                                            Chapitre 3. Gouvernance, Partenariat et RACI\\
\vspace{0.5em}
\vspace{0.5em}
\vspace{0.5em}
\vspace{0.5em}
Nexus             Unique responsable de la livraison pédagogique de A à Z : de- Contrat-cadre RBK ↔ Nexus (MSA) (voir ci-dessus). Bon de\\
Réussite          sign du programme, sélection/contrats/paiement des mentors, commande / SOW (RBK ↔ Nexus) (voir ci-dessus). Contrats\\
(maître           standards qualité, examens, jurys, traçabilité, preuves d’exécu- mentors ↔ Nexus : sous-traitance, confidentialité, cession/u-\\
d’œuvre           tion. Interface contractuelle : Nexus facture RBK, Nexus sous- sage limité des supports, non-sollicitation. Politique données\\
pédagogique       traite les mentors.                                              (DPA) : RBK ↔ Nexus ↔ MFAI.\\
et opérateur)\\
\vspace{0.5em}
Money             Détient l’IP et l’exploitation de journey.mfai.app. Consent Licence MFAI ↔ Nexus : usage journey.mfai.app, IP, SLA,\\
Factory AI        à Nexus une licence d’usage (gratuite cohorte pilote, au- réversibilité, données. Politique données (DPA) (voir ci-\\
(propriétaire     delà, accord contractuel prédéfini (redevance par cohorte ou dessus).\\
techno / IP)      MFAI ne facture pas RBK : RBK contracte avec Nexus ; Nexus\\
contracte avec MFAI.\\
\vspace{0.5em}
\vspace{0.5em}
\vspace{0.5em}
\vspace{0.5em}
EL\\
TI\\
N\\
DE\\
\subsection{3.3 Matrice RACI (Répartition des responsabilités)}
Le tableau ci-dessous présente une macro-répartition des tâches clés. Une version détaillée sera annexée au contrat-cadre.\\
\vspace{0.5em}
\vspace{0.5em}
\vspace{0.5em}
\vspace{0.5em}
FI\\
N\\
Matrice RACI — Macro-Processus RBK 2.0\\
CO\\
Activité / Tâche                                          RBK                 Nexus Réussite           Money Factory AI\\
Phase Commerciale\\
Marketing \textbackslash{}& Acquisition Apprenants                          A                        C                        I\\
Inscriptions \textbackslash{}& Contrats Apprenants (ISA inclus)             A                        R                        I\\
Exécution Pédagogique\\
Design du programme \textbackslash{}& Syllabus                              C                        R                        A\\
Sélection \textbackslash{}& contractualisation des mentors                  C                        A                        I\\
Paiement des mentors                                        I                        A                        I\\
\vspace{0.5em}
\vspace{0.5em}
\vspace{0.5em}
\vspace{0.5em}
\begin{confidential}CONFIDENTIEL — PROJET RBK 2.0                                           Page 50 sur 316\end{confidential}
RBK 2.0 — Whitepaper Architecte Web3                                                             Chapitre 3. Gouvernance, Partenariat et RACI\\
\vspace{0.5em}
\vspace{0.5em}
\vspace{0.5em}
Organisation des jurys \textbackslash{}& examens                    C                     A                       R\\
Suivi qualité \textbackslash{}& conformité                          R                     A                       C\\
Technologie \textbackslash{}& Plateforme\\
Fourniture plateforme journey.mfai.app              I                     R                       A\\
Support technique plateforme (SLA)                  I                     R                       A\\
Gestion des données apprenants (DPA)                A                     R                       C\\
Finance \textbackslash{}& Reporting\\
Facturation Apprenants \textbackslash{}& Recouvrement ISA           A                     C                       I\\
Facturation RBK → Nexus                             A                     R                       I\\
Reporting d’exécution (KPIs)                        R                     A                       I\\
\vspace{0.5em}
\vspace{0.5em}
\vspace{0.5em}
\vspace{0.5em}
EL\\
Légende : RResponsable (Responsible), AApprobateur (Accountable), CConsulté (Consulted), IInformé (Informed).\\
\vspace{0.5em}
\vspace{0.5em}
\vspace{0.5em}
\vspace{0.5em}
TI\\
N\\
DE\\
FI\\
N\\
CO\\
\vspace{0.5em}
\vspace{0.5em}
\vspace{0.5em}
\vspace{0.5em}
\begin{confidential}CONFIDENTIEL — PROJET RBK 2.0                                 Page 51 sur 316\end{confidential}
\section{ANALYSE DU CONTEXTE}
\vspace{0.5em}
\vspace{0.5em}
\subsection{4.1 L’Opportunité Web3 \textbackslash{}& Solana}
\vspace{0.5em}
\vspace{0.5em}
\vspace{0.5em}
\vspace{0.5em}
EL\\
Avant de comparer les verticales, nous posons ici le cadre de lecture du marché Web3 et de la position singulière de Solana.\\
\vspace{0.5em}
\vspace{0.5em}
\vspace{0.5em}
\vspace{0.5em}
TI\\
N\\
\subsection{4.1.0.1 Définitions Minimales (Lexique Opérationnel)}
\vspace{0.5em}
\vspace{0.5em}
\vspace{0.5em}
\vspace{0.5em}
DE\\
Pour comprendre l’arbitrage RBK, il faut maîtriser le vocabulaire du marché :\\
\vspace{0.5em}
\vspace{0.5em}
\vspace{0.5em}
\vspace{0.5em}
FI\\
• Web3 : Un internet où les utilisateurs possèdent leurs données et leurs actifs, sécurisé par des réseaux décentralisés (Blockchains).\\
\vspace{0.5em}
\vspace{0.5em}
\vspace{0.5em}
\vspace{0.5em}
N\\
• Solana (SVM) : La blockchain la plus performante à ce jour (65k TPS théoriques), optimisée pour des applications grand public (Payments, Gaming, DePIN).\\
CO\\
• DeFi (Decentralized Finance) : Services financiers (prêt, échange) sans intermédiaire bancaire.\\
• DePIN (Decentralized Physical Infrastructure) : Réseaux physiques (Wifi, GPU) gérés par des incitations crypto.\\
• Bounty : Mission à la tâche rémunérée en stablecoins1 (USDC), souvent premier revenu d’un étudiant.\\
\vspace{0.5em}
\vspace{0.5em}
\subsection{4.1.0.2 Segmentation de la Demande}
Le marché ne cherche pas ”un dev blockchain”, mais des spécialistes par verticale.\\
1\\
Web3/Blockchain — Stablecoin : actif numérique cherchant à maintenir une parité stable via collatéral, algorithme ou réserves hors chaîne auditées.\\
\vspace{0.5em}
\vspace{0.5em}
\vspace{0.5em}
\vspace{0.5em}
52\\
RBK 2.0 — Whitepaper Architecte Web3                                                                                           Chapitre 4. ANALYSE DU CONTEXTE\\
\vspace{0.5em}
\vspace{0.5em}
\begin{center}\begin{tabular}TAB. 4.1 : Segmentation des Rôles Web3 (2025)\end{tabular}\end{center}
\vspace{0.5em}
\vspace{0.5em}
\subsection{4.1.0.3 Indicateurs                                                                      2025                                                            actualisés}
\vspace{0.5em}
\vspace{0.5em}
Segment        Rôles Clés                Livrables Concrets                     Compétence Dominante\\
DeFi           Smart Contract Eng.       AMM, Lending Protocol, Vaults          Mathématiques \textbackslash{}& Sécurité\\
DePIN          Rust Embedded Eng.        Drivers IoT, Proof-of-Coverage         Optimisation Bas-niveau\\
Infra          DevOps / RPC Eng.         Indexers, Validators, Nodes            Linux, Docker, Rust\\
Consumer       Mobile dApp Dev.          Wallet UI, Payment SDK                 UX/UI, React Native\\
\vspace{0.5em}
\vspace{0.5em}
\vspace{0.5em}
\vspace{0.5em}
EL\\
\subsection{4.1.0.4 Pourquoi Solana est un Accélérateur d’Employabilité}
\vspace{0.5em}
\vspace{0.5em}
\vspace{0.5em}
\vspace{0.5em}
TI\\
Contrairement à Ethereum (EVM) qui est saturé et fragmenté (L2s), Solana offre un écosystème unifié et en hyper-croissance (+500\textbackslash{}% d’adresses actives en 2024).\\
\vspace{0.5em}
\vspace{0.5em}
\vspace{0.5em}
\vspace{0.5em}
N\\
Pour un junior, la courbe d’apprentissage est plus raide (Rust), mais la concurrence est moindre et les primes sont plus élevées. La Superteam offre un pipeline direct\\
\vspace{0.5em}
\vspace{0.5em}
\vspace{0.5em}
\vspace{0.5em}
DE\\
vers l’emploi via Earn.\\
\vspace{0.5em}
\vspace{0.5em}
\vspace{0.5em}
\vspace{0.5em}
FI\\
Statistiques Marché 2025\\
\vspace{0.5em}
\vspace{0.5em}
\vspace{0.5em}
\vspace{0.5em}
N\\
1. Postes ouverts : 15 000+ offres actives en Remote GlobalWeb3.careera .\\
CO\\
2. Pénurie : 58\textbackslash{}% des Lead Techs citent le recrutement d'ingénieurs Rust seniors comme leur blocage n°1.\\
\vspace{0.5em}
3. Développeurs Actifs : \textless{} 25 000 développeurs crypto mensuels vs 25M devs Web2. L'opportunité d'arbitrage est de x1000Electric\\
Capitalb .\\
a\\
Source --- Données agrégées par Web3.career \textbackslash{}& TrueUp Tech Jobs Report, Q4 2024.\\
b\\
Source --- Source : Electric Capital Developer Report 2023.\\
\vspace{0.5em}
\vspace{0.5em}
\vspace{0.5em}
\vspace{0.5em}
\subsection{4.2 Dynamique Salariale}
Cette section éclaire l’arbitrage salarial et fiscal afin d’ancrer les promesses de revenu dans des ordres de grandeur réalistes.\\
\vspace{0.5em}
\vspace{0.5em}
\vspace{0.5em}
\vspace{0.5em}
\begin{confidential}CONFIDENTIEL — PROJET RBK 2.0                                               Page 53 sur 316\end{confidential}
RBK 2.0 — Whitepaper Architecte Web3                                                                                        Chapitre 4. ANALYSE DU CONTEXTE\\
\vspace{0.5em}
\vspace{0.5em}
\subsection{4.2.0.1 Indicateurs 2025 actualisés}
Le marché confirme l’accélération détectée par notre audit : l’offre de talents reste inférieure à la traction produits, ouvrant une fenêtre 18–24 mois pour capter\\
la demande.\\
\vspace{0.5em}
KPIs Web3 Solana 2025\\
\vspace{0.5em}
\vspace{0.5em}
\vspace{0.5em}
\vspace{0.5em}
EL\\
TI\\
N\\
DE\\
FI\\
N\\
CO\\
\vspace{0.5em}
\vspace{0.5em}
\vspace{0.5em}
\vspace{0.5em}
\begin{confidential}CONFIDENTIEL — PROJET RBK 2.0                                             Page 54 sur 316\end{confidential}
RBK 2.0 — Whitepaper Architecte Web3                                                                                  Chapitre 4. ANALYSE DU CONTEXTE\\
\vspace{0.5em}
\vspace{0.5em}
\vspace{0.5em}
Indicateur                    Lecture 2025                      Source\\
Développeurs actifs           ∼1 900 profils (+35\textbackslash{}% YoY)         Electric Capital Developer Report 2024 (addendum Q1 2025)\\
mensuels Solana               concentrés sur les couches core\\
et validator\\
TVL Solana                    4,6 Md\textbackslash{}$ (sept. 2024) → 6,5 Md\textbackslash{}$ Messari State of Solana Q3 2024, Artemis dashboard\\
(projection Q1 2025)\\
Volumes DEX Solana            30 Md\textbackslash{}$/mois au S4 2024 (x7 vs     CoinGecko Derivatives Monitor, 15/12/2025\\
2023)\\
Investissements MENA          310 M\textbackslash{}$ levés en 2024 (+22\textbackslash{}%        RootData ”MENA Web3 Funding” 2025 Outlook\\
Web3                          YoY)\\
Talents Rust/Web3             \textless{} 350 profils référencés          Nexus Talent Intelligence, déc. 2025\\
\vspace{0.5em}
\vspace{0.5em}
\vspace{0.5em}
\vspace{0.5em}
EL\\
disponibles (MENA)            (LinkedIn, Superteam Earn)\\
\vspace{0.5em}
\vspace{0.5em}
\vspace{0.5em}
\vspace{0.5em}
TI\\
Pipeline employeurs RBK       9 partenaires (3 Solana, 3 EVM,   MoU/LOI Nexus Réussite Superteam (𝑛 = 9)\\
3 DePIN) couvrant 120 postes\\
\vspace{0.5em}
\vspace{0.5em}
\vspace{0.5em}
\vspace{0.5em}
N\\
2025\\
\vspace{0.5em}
\vspace{0.5em}
\vspace{0.5em}
\vspace{0.5em}
DE\\
ISA-friendly roles (Remote)   42\textbackslash{}% des offres Solana proposées   Web3.career Salary Index, nov. 2025\\
en remote total avec packages\\
\vspace{0.5em}
\vspace{0.5em}
\vspace{0.5em}
\vspace{0.5em}
FI\\
\textgreater{} 60k\textbackslash{}$/an\\
\vspace{0.5em}
\vspace{0.5em}
N\\
Durée moyenne                 86 jours (contre 42 jours Web2    Lightspeed Crypto Talent Pulse 2025\\
CO\\
recrutement Rust Senior       senior)\\
Solana validator count        3 400+ noeuds (doublement vs      Solana Foundation Validator Health Report 2025\\
2023) 33\textbackslash{}% en régions\\
émergentes\\
Activité bounties Superteam 2 7 M\textbackslash{}$ redistribués en 2024,        Superteam Earn Transparency Report 2024\\
ticket médian 1 2k\textbackslash{}$\\
Temps médian premier          48 jours post-onboarding pour     Nexus Earn Analytics, cohorte pilote 2024\\
revenu                        un builder Genesis\\
Coût moyen bug bounty L1      75k\textbackslash{}$ (x1,5 vs 2023)               Immunefi Crypto Losses Bug Bounties 2025\\
\vspace{0.5em}
\vspace{0.5em}
\vspace{0.5em}
\vspace{0.5em}
\begin{confidential}CONFIDENTIEL — PROJET RBK 2.0                                        Page 55 sur 316\end{confidential}
RBK 2.0 — Whitepaper Architecte Web3                                                                                     Chapitre 4. ANALYSE DU CONTEXTE\\
\vspace{0.5em}
\vspace{0.5em}
\vspace{0.5em}
Croissance DePIN              +180\textbackslash{}% utilisateurs actifs          Messari DePIN Landscape 2025\\
Solana/EVM                    (Helium, Hivemapper, io.net)\\
Taux attrition talents Web3   12\textbackslash{}% (vs 22\textbackslash{}% Web2)                  Terminal.dev Talent Survey 2025\\
Cadres réglementaires         3 programmes incitatifs (Startup   RBK Legal Watch, déc. 2025\\
Tunisie                       Act, Export IT, Sandbox BCT)\\
Budget formation corporate    4 3 M\textbackslash{}$ engagés par les top         Solana Foundation Training Grants Memo 2025\\
Solana                        protocols (Anza, Helius, Jump)\\
Taux de conversion            27\textbackslash{}% (parcours Superteam)           Superteam Outcomes Report 2025\\
bounties → CDI\\
Ratio offres Senior vs Junior 3,1 :1 en faveur des profils       Electric Capital, LinkedIn Talent Insights 2025\\
Senior/Architecte\\
\vspace{0.5em}
\vspace{0.5em}
\vspace{0.5em}
\vspace{0.5em}
EL\\
Programmes financiers         6 dispositifs (FOPRODI, PEE,       APII Ministère TIC 2025\\
\vspace{0.5em}
\vspace{0.5em}
\vspace{0.5em}
\vspace{0.5em}
TI\\
publics mobilisables          FODEC, Caisse des Dépôts,\\
Startup Act)\\
\vspace{0.5em}
\vspace{0.5em}
\vspace{0.5em}
\vspace{0.5em}
N\\
DE\\
Synthèse Nexus Réussite consolidant Electric Capital, Messari, RootData, Superteam Earn, Artemis, Galaxy Research (rapports 2024-2025).\\
\vspace{0.5em}
\vspace{0.5em}
\vspace{0.5em}
\vspace{0.5em}
FI\\
N\\
CO\\
\vspace{0.5em}
\vspace{0.5em}
\vspace{0.5em}
\vspace{0.5em}
\begin{confidential}CONFIDENTIEL — PROJET RBK 2.0                                            Page 56 sur 316\end{confidential}
RBK 2.0 — Whitepaper Architecte Web3                                                                                                            Chapitre 4. ANALYSE DU CONTEXTE\\
\vspace{0.5em}
\vspace{0.5em}
\vspace{0.5em}
3,000                                                                                Lecture :\\
• Talent gap persistant\\
• TVL ↑ traction produits\\
• Offres emploi x4 en 4 ans\\
2,500\\
Développeurs actifs Solana\\
\vspace{0.5em}
\vspace{0.5em}
\vspace{0.5em}
\vspace{0.5em}
2,000\\
\vspace{0.5em}
\vspace{0.5em}
\vspace{0.5em}
\vspace{0.5em}
EL\\
1,000\\
\vspace{0.5em}
\vspace{0.5em}
\vspace{0.5em}
\vspace{0.5em}
TI\\
N\\
500\\
\vspace{0.5em}
\vspace{0.5em}
\vspace{0.5em}
\vspace{0.5em}
DE\\
FI\\
0\\
2,021          2,022                               2,023                           2,024                             2,025              2,026\\
\vspace{0.5em}
\vspace{0.5em}
N\\
Année\\
CO\\
\vspace{0.5em}
Développeurs actifs (Electric Capital)   TVL Solana (M\textbackslash{}$)       Offres emploi Web3 (index base 2021)\\
\vspace{0.5em}
\vspace{0.5em}
FIG. 4.1 : Corrélation talents Solana vs traction financière (2021-2026)\\
\vspace{0.5em}
Benchmark formations Web3 (2025)\\
\vspace{0.5em}
\vspace{0.5em}
\vspace{0.5em}
\vspace{0.5em}
\begin{confidential}CONFIDENTIEL — PROJET RBK 2.0                                                            Page 57 sur 316\end{confidential}
RBK 2.0 — Whitepaper Architecte Web3                                                                                      Chapitre 4. ANALYSE DU CONTEXTE\\
\vspace{0.5em}
\vspace{0.5em}
\vspace{0.5em}
Programme             Durée/Format           Modèle                 Gaps vs RBK 2.0\\
économique\\
Encode Club Solana    12 semaines, remote    Gratuit (ISA 17\textbackslash{}%       Parcours court, charte santé absente, pas de transfert gouvernance pour partenaires lo-\\
Fellowship            intensif               sur 24 mois si         caux\\
placement)\\
01Founders x          9 mois, hybride        Stipend +              Capacité limitée (40 places), absence d’ISA MENA, pas de pipeline DePIN\\
Solana Foundation     Londres/Remote         recrutement ∼70\textbackslash{}%\\
ChainShot             16 semaines,           Tuition 3 800\textbackslash{}$ +       Scope smart contract EVM uniquement, pas de gouvernance partagée, santé non adressée\\
(Alchemy) relaunch    bootcamp contract      hiring network\\
Holberton x Smart     12 mois, présentiel    ISA 17\textbackslash{}% (contrat       Curriculum Web2 généraliste, peu d’ancrage Solana, pas d’audit sécurité avancé\\
Tunisia                                      local)\\
\vspace{0.5em}
\vspace{0.5em}
\vspace{0.5em}
\vspace{0.5em}
EL\\
Simplon Tunisia       4 mois, bootcamp       Subvention             Focus JS/React, pas de portfolio on-chain, pas d’offre architecte senior\\
Web3                  intensif               publique + bourses\\
\vspace{0.5em}
\vspace{0.5em}
\vspace{0.5em}
\vspace{0.5em}
TI\\
Platzi Web3 School    6 mois,                Subscription           Faible sélectivité, parcours self-paced, pas de garanties placement\\
\vspace{0.5em}
\vspace{0.5em}
\vspace{0.5em}
\vspace{0.5em}
N\\
asynchronous           (29\textbackslash{}$/mois)\\
\vspace{0.5em}
\vspace{0.5em}
\vspace{0.5em}
\vspace{0.5em}
DE\\
Superteam Africa      8 semaines, remote     Bounties + grants      Excellent tremplin sortie, nécessite pipeline formation amont (complément RBK 2.0)\\
Residency             cohort\\
\vspace{0.5em}
\vspace{0.5em}
\vspace{0.5em}
\vspace{0.5em}
FI\\
RBK 2.0 (RBK ×        48 semaines, hybride Tuition 15 900 TND       Gouvernance co-pilotée, SLA qualité, transfert savoir-faire, plan continuité, pipeline em-\\
\vspace{0.5em}
\vspace{0.5em}
\vspace{0.5em}
N\\
Nexus)                présentiel/remote    + ISA + forfait\\
CO               ployeurs signé RBK\\
Nexus\\
\vspace{0.5em}
Sources : Encode Club 2025 Impact Report, Solana Foundation Talent Update 2025, Alchemy ChainShot relaunch note 2024, Holberton Tunisia ISA deck 2024,\\
Simplon Tunisie Web3 brochure 2024, Platzi Investor Deck 2025, Superteam Africa Residency recap 2024.\\
\vspace{0.5em}
Lecture pour la direction RBK\\
\vspace{0.5em}
RBK 2.0 occupe la niche ”Architecte senior-by-design” en combinant durée longue, charte santé, ISA robuste et co-gouvernance. Aucun concurrent ne propose\\
simultanément transfert de savoir-faire, plan de continuité et pipeline employeurs sécurisés en Afrique du Nord.\\
\vspace{0.5em}
\vspace{0.5em}
\vspace{0.5em}
\subsection{4.2.0.2 Hypothèses de Lecture (TND vs USD)}
Les chiffres présentés ci-dessous sont exprimés en USD brut annuel. Pour un talent tunisien en remote :\\
\vspace{0.5em}
\vspace{0.5em}
\vspace{0.5em}
\vspace{0.5em}
\begin{confidential}CONFIDENTIEL — PROJET RBK 2.0                                           Page 58 sur 316\end{confidential}
RBK 2.0 — Whitepaper Architecte Web3                                                                                            Chapitre 4. ANALYSE DU CONTEXTE\\
\vspace{0.5em}
\vspace{0.5em}
• Conversion : 1 USD ≈ 3.1 TND.\\
• Fiscalité : En statut ”Exportateur de Services” (entreprise totalement exportatrice), l’imposition est avantageuse, maximisant le net.\\
• Réalité Marché : Le salaire ”Junior” Web3 (60k\textbackslash{}$) correspond souvent à un salaire ”VP Engineering” sur le marché local.\\
\vspace{0.5em}
\vspace{0.5em}
\subsection{4.2.0.3 Grille de Rémunération Standard}
\vspace{0.5em}
\vspace{0.5em}
\begin{center}\begin{tabular}TAB. 4.5 : Grille Salariale Web3 (Remote Global) vs Local\end{tabular}\end{center}
\vspace{0.5em}
\vspace{0.5em}
Rôle                          Junior (0-2 ans)        Senior (3+ ans)          Pré-requis\\
Solana Rust Engineer            60k\textbackslash{}$ - 90k\textbackslash{}$            140k\textbackslash{}$ - 220k\textbackslash{}$           Portfolio GitHub Solide\\
\vspace{0.5em}
\vspace{0.5em}
\vspace{0.5em}
\vspace{0.5em}
EL\\
Security Auditor                80k\textbackslash{}$ - 120k\textbackslash{}$              250k\textbackslash{}$+               Track Record de vulnérabilités trouvées\\
Fullstack dApp                  50k\textbackslash{}$ - 80k\textbackslash{}$            110k\textbackslash{}$ - 160k\textbackslash{}$           Portfolio React + Anchor\\
\vspace{0.5em}
\vspace{0.5em}
\vspace{0.5em}
\vspace{0.5em}
TI\\
Dev Web2 (Tunisie)             15k - 25k TND           40k - 60k TND           Diplôme Ingénieur\\
\vspace{0.5em}
\vspace{0.5em}
\vspace{0.5em}
\vspace{0.5em}
N\\
DE\\
Sources : Web3.career, Pantera Capital Salary Survey 2024. Note : Les montants Web3 sont en Brut Global. En Tunisie, grâce au statut exportateur (off-shore/startup act), le\\
Net est maximisé (charges allégées), rendant le pouvoir d’achat x3 supérieur au local.\\
\vspace{0.5em}
\vspace{0.5em}
\vspace{0.5em}
\vspace{0.5em}
FI\\
N\\
\subsection{4.2.0.4 Modèle ROI Candidat (Simulation 1 an)}
CO\\
Scénario     Revenu Cible        Time-to-Revenue               Risques\\
Prudent      1 500 \textbackslash{}$/mois        4 mois post-cursus            Marché Bear, Anglais moyen\\
Médian       2 500 \textbackslash{}$/mois        2 mois post-cursus            Concurrence, Portfolio standard\\
Top Gun      5 000 \textbackslash{}$/mois        Pendant le cursus (S20)       Burnout, Gestion charge travail\\
\vspace{0.5em}
\vspace{0.5em}
\vspace{0.5em}
\subsection{4.3 Croissance du Marché}
Nous analysons la traction emploi/produits pour valider la soutenabilité du modèle RBK au-delà des cycles crypto.\\
\vspace{0.5em}
\vspace{0.5em}
\vspace{0.5em}
\begin{confidential}CONFIDENTIEL — PROJET RBK 2.0                                               Page 59 sur 316\end{confidential}
RBK 2.0 — Whitepaper Architecte Web3                                                                                         Chapitre 4. ANALYSE DU CONTEXTE\\
\vspace{0.5em}
\vspace{0.5em}
\subsection{4.3.0.1 Définition de l’Index}
Le graphique ci-dessous agrège le volume d’offres d’emploi techniques (Engineering, Product, Design) postées sur les 5 principaux job boards crypto, normalisé\\
sur une base 100 en Janvier 2021.\\
\vspace{0.5em}
\vspace{0.5em}
\subsection{4.3.0.2 Lecture Stratégique}
La corrélation avec le prix des actifs (BTC/SOL) diminue : les entreprises construisent (Build) même en bear market. Cela signifie que l’embauche se professionnalise\\
et devient moins volatile. Pour RBK, cela valide la stratégie de ”formation longue” (7 mois) qui lisse les cycles de court terme.\\
\vspace{0.5em}
\vspace{0.5em}
\vspace{0.5em}
\vspace{0.5em}
EL\\
TI\\
N\\
DE\\
FI\\
N\\
CO\\
\vspace{0.5em}
\vspace{0.5em}
\vspace{0.5em}
\vspace{0.5em}
\begin{confidential}CONFIDENTIEL — PROJET RBK 2.0                                              Page 60 sur 316\end{confidential}
RBK 2.0 — Whitepaper Architecte Web3                                                                                                                                            Chapitre 4. ANALYSE DU CONTEXTE\\
\vspace{0.5em}
\vspace{0.5em}
Index Offres d’Emploi Web3 vs Web2 (2021-2026)\\
450\\
Index du Volume d’Offres (Base 100 = Jan 2021)           Lecture du Graphique :\\
100 = Volume de référence (2021)\\
400     200 = x2 (Doublement du volume)\\
300 = x3 (Triplement)\\
\vspace{0.5em}
350\\
\vspace{0.5em}
Adoption\\
300\\
Institutionnelle\\
\vspace{0.5em}
250\\
\vspace{0.5em}
\vspace{0.5em}
\vspace{0.5em}
\vspace{0.5em}
EL\\
200\\
\vspace{0.5em}
\vspace{0.5em}
\vspace{0.5em}
\vspace{0.5em}
TI\\
150\\
\vspace{0.5em}
\vspace{0.5em}
\vspace{0.5em}
\vspace{0.5em}
N\\
DE\\
100\\
\vspace{0.5em}
\vspace{0.5em}
\vspace{0.5em}
\vspace{0.5em}
FI\\
2021                            2022                    2023                                              2024                2025                  2026\\
\vspace{0.5em}
\vspace{0.5em}
\vspace{0.5em}
N\\
Année\\
CO\\
Web3 (Rust/Solidity) - Hypothèse Croissance\\
Web2 (Marché IT Classique)\\
\vspace{0.5em}
\vspace{0.5em}
\vspace{0.5em}
\vspace{0.5em}
Source : Projection interne basée sur Electric Capital Reports \textbackslash{}& LinkedIn Data.\\
\vspace{0.5em}
\vspace{0.5em}
\vspace{0.5em}
\vspace{0.5em}
\begin{confidential}CONFIDENTIEL — PROJET RBK 2.0                                                                                         Page 61 sur 316\end{confidential}
\section{ARCHITECTURE DE PARTENARIAT \textbackslash{}& GOUVERNANCE OPÉ-}
RATIONNELLE (RBK ↔ Nexus Réussite)\\
\vspace{0.5em}
\vspace{0.5em}
\vspace{0.5em}
\vspace{0.5em}
EL\\
\subsection{5.1 Le partenariat structurant}
\vspace{0.5em}
\vspace{0.5em}
\vspace{0.5em}
\vspace{0.5em}
TI\\
N\\
Ce chapitre formalise le cadre contractuel liant RBK, porteur officiel de l’offre de formation, et Nexus Réussite, maître d’œuvre pédagogique exclusif. Cette\\
\vspace{0.5em}
\vspace{0.5em}
\vspace{0.5em}
\vspace{0.5em}
DE\\
architecture bipartite garantit la traçabilité des responsabilités, l’alignement des incentives et le recours systématique à des indicateurs vérifiables.\\
\vspace{0.5em}
Ď RBK et Nexus : un binôme d’exécution indissociable\\
\vspace{0.5em}
FI\\
N\\
Le modèle opérationnel repose sur deux entités liées par un contrat-cadre de services : RBK agit comme opérateur légal, commercial et responsable final\\
CO\\
vis-à-vis des apprenants ; Nexus Réussite en est le Maître d’œuvre pédagogique exclusif, responsable contractuellement de la conception et de l’exécution.\\
Chaque promesse faite au marché est adossée à une responsabilité documentée (voir Annexe D).\\
\vspace{0.5em}
\vspace{0.5em}
\vspace{0.5em}
\subsection{5.1.1 RBK — l’opérateur et la face publique}
• Déploie et finance l’infrastructure locale (campus, logistique, assurances) et porte la conformité administrative tunisienne.\\
• Commercialise l’offre, encaisse les frais initiaux, signe et administre les contrats apprenants et le ISA.\\
• Supervise la relation client et la qualité de service perçue, s’assurant que les engagements contractuels de Nexus sont respectés.\\
\vspace{0.5em}
\vspace{0.5em}
\vspace{0.5em}
\vspace{0.5em}
62\\
RBK 2.0 — Whitepaper Architecte Web3                                                                                          Chapitre 5. Architecture de Partenariat\\
\vspace{0.5em}
\vspace{0.5em}
\subsection{5.1.2 Nexus Réussite — le Maître d’œuvre pédagogique}
• Conçoit le curriculum intégral, pilote les Block Checks, anime les jurys et produit toutes les preuves d’exécution (reports, dashboards, procès-verbaux).\\
• Recrute, forme, contractualise et paie l’équipe pédagogique (mentors, responsables de tracks, intervenants externes).\\
• Garantit la qualité d’exécution : respect du Definition of Done, contrôle continu, audits internes, traçabilité des décisions.\\
\vspace{0.5em}
\vspace{0.5em}
\subsection{5.1.3 Gouvernance conjointe et comités}
• Comité exécutif conjoint (CEC) : RBK et Nexus se réunissent chaque semaine (pilotage opérationnel) et chaque mois (revue stratégique) pour arbitrer les priorités,\\
approuver les plans correctifs et valider les budgets ; Money Factory AI peut être convié à titre consultatif sur les sujets plateforme.\\
• Comité éthique et pédagogique : instance Nexus présidée par RBK qui statue sur les cas d’escalade (burnout, litiges académiques, recours ISA).\\
\vspace{0.5em}
\vspace{0.5em}
\vspace{0.5em}
\vspace{0.5em}
EL\\
• Comité qualité et risques : suit les indicateurs critiques (voir Chapitre 17) et déclenche les plans de continuité.\\
\vspace{0.5em}
\vspace{0.5em}
\vspace{0.5em}
\vspace{0.5em}
TI\\
N\\
\subsection{5.2 Rôle de Money Factory AI : fournisseur technologique}
\vspace{0.5em}
\vspace{0.5em}
\vspace{0.5em}
\vspace{0.5em}
DE\\
Licence Venture Engine et responsabilités\\
\vspace{0.5em}
\vspace{0.5em}
\vspace{0.5em}
\vspace{0.5em}
FI\\
N\\
Élément                      Description contractuelleCO\\
Licence logicielle           Nexus bénéficie d’une licence d’usage de la plateforme Venture Engine (journey.mfai.app) incluant maintenance, mises à jour et\\
support SLA 99,5\textbackslash{}%. RBK n’est pas redevable directement envers MFAI.\\
Réversibilité des            MFAI garantit la portabilité complète des données apprenants vers RBK en cas de sortie de Nexus (format ouvert, délai maximal\\
données                      30 jours).\\
Support et sécurité          MFAI fournit un canal d’escalade 24/7, applique une politique ”privacy-by-design” et notifie RBK et Nexus de tout incident de\\
sécurité en moins de 12 heures.\\
\vspace{0.5em}
\vspace{0.5em}
\vspace{0.5em}
\vspace{0.5em}
\subsection{5.3 Flux contractuels et financiers}
\vspace{0.5em}
\vspace{0.5em}
\vspace{0.5em}
\begin{confidential}CONFIDENTIEL — PROJET RBK 2.0                                               Page 63 sur 316\end{confidential}
RBK 2.0 — Whitepaper Architecte Web3                                                                                       Chapitre 5. Architecture de Partenariat\\
\vspace{0.5em}
\vspace{0.5em}
\vspace{0.5em}
1. Contrat-cadre de services (RBK ↔ Nexus) : définit le périmètre pédagogique, les KPIs, les clauses de pénalité, la propriété intellectuelle partagée et la gouvernance\\
(CEC).\\
2. Bon de commande / SOW par cohorte : précise les livrables phase par phase, le staffing validé, les plafonds d’heures mentors et les seuils de performance\\
déclenchant bonus/malus.\\
3. Facturation : Nexus facture RBK en trois jalons (Genesis, Bâtisseur, Architecte) ; RBK facture les apprenants (frais initiaux, échéanciers ISA) et constitue le Fonds\\
de Garantie ISA (voir §15.6.1).\\
4. Flux ISA : RBK conserve l’entier recouvrement, Nexus transmet les preuves de graduation et de placement nécessaires aux déclenchements.\\
5. Gestion de la trésorerie : les comptes séquestres (ISA, bourses) sont opérés par RBK avec droit de regard Nexus ; un plancher de trésorerie minimal est fixé\\
contractuellement (cf. Chapitre 15).\\
\vspace{0.5em}
\vspace{0.5em}
\vspace{0.5em}
\vspace{0.5em}
EL\\
\subsection{5.4 Sécurisation des intérêts de RBK}
\vspace{0.5em}
\vspace{0.5em}
\vspace{0.5em}
\vspace{0.5em}
TI\\
Cette approche protège RBK contre les risques de dérive d’exécution tout en maximisant la valeur contractuelle : RBK garde la main sur les décisions structurantes,\\
\vspace{0.5em}
\vspace{0.5em}
\vspace{0.5em}
\vspace{0.5em}
N\\
bénéficie d’un droit de suspension à clauses multiples et peut internaliser l’ensemble des actifs en moins de 60 jours.\\
\vspace{0.5em}
\vspace{0.5em}
\vspace{0.5em}
\vspace{0.5em}
DE\\
Clauses de réversibilité go/no-go\\
\vspace{0.5em}
\vspace{0.5em}
\vspace{0.5em}
\vspace{0.5em}
FI\\
N\\
Déclencheur           Mécanisme contractuel          CO                                                                                  Décision RBK\\
constaté\\
Marqueurs             Comité Qualité convoqué sous 5 jours, plan correctif écrit, jalon go/no-go à +4 semaines                           Validation RBK\\
qualité : NPS \textless{} 60                                                                                                                       obligatoire ; échec ⇒\\
sur 2 comités                                                                                                                            suspension jalon\\
consécutifs ou                                                                                                                           suivant activation\\
Learning Velocity                                                                                                                        plan de substitution\\
en alerte rouge                                                                                                                          mentors\\
\vspace{0.5em}
\vspace{0.5em}
\vspace{0.5em}
\vspace{0.5em}
\begin{confidential}CONFIDENTIEL — PROJET RBK 2.0                                             Page 64 sur 316\end{confidential}
RBK 2.0 — Whitepaper Architecte Web3                                                                                   Chapitre 5. Architecture de Partenariat\\
\vspace{0.5em}
\vspace{0.5em}
\vspace{0.5em}
Sous-performance      Pénalité financière automatique (malus) et revue CEC extraordinaire ; Nexus doit présenter plan de rattrapage RBK peut imposer gel\\
opérationnelle :      + staffing de réserve                                                                                         des nouvelles\\
SLA mentors non                                                                                                                     cohortes et activer\\
tenus (fill rate                                                                                                                    recours prestataires\\
\textless{} 95\textbackslash{}%) ou budget                                                                                                                    tiers listés en annexe\\
acquisition                                                                                                                         D\\
dépassement\\
\textgreater{} 15\textbackslash{}%\\
Incident critique :   Clause de suspension immédiate (45 jours) + transfert des accès clés (LMS, data) ; audit externe déclenché par RBK arbitre maintien\\
manquement            RBK                                                                                                            ou résiliation ; Nexus\\
conformité, fuite                                                                                                                    doit assister à la\\
données, litige                                                                                                                      transition complète\\
\vspace{0.5em}
\vspace{0.5em}
\vspace{0.5em}
\vspace{0.5em}
EL\\
juridique majeur\\
Divergence            Obligation de passage en Comité Stratégique (CEC mensuel) avec vote RBK double voix                            RBK conserve veto et\\
\vspace{0.5em}
\vspace{0.5em}
\vspace{0.5em}
\vspace{0.5em}
TI\\
stratégique :                                                                                                                        peut rebasculer la\\
\vspace{0.5em}
\vspace{0.5em}
\vspace{0.5em}
\vspace{0.5em}
N\\
changement de                                                                                                                        priorisation produit\\
\vspace{0.5em}
\vspace{0.5em}
\vspace{0.5em}
\vspace{0.5em}
DE\\
roadmap produit\\
imposé par Nexus\\
\vspace{0.5em}
\vspace{0.5em}
\vspace{0.5em}
\vspace{0.5em}
FI\\
sans accord RBK\\
\vspace{0.5em}
\vspace{0.5em}
\vspace{0.5em}
N\\
CO\\
Ď Transfert progressif de savoir-faire (24 mois)\\
• Binômage systématique : chaque fonction critique Nexus (Head of Curriculum, Lead Mentor, Ops ISA) possède un homologue RBK shadowé 2\\
jours/semaine ; rotation des binômes tous les 6 mois documentée.\\
• Playbooks propriétaires : production trimestrielle de dossiers HOW WE RUN (curriculum, QA, comités) validés par RBK et stockés dans le référentiel\\
interne (voir Annexe D).\\
• Certification mentors RBK : objectif 60\textbackslash{}% des mentors labellisés Certification mentor RBKa en 18 mois, permettant à RBK de redéployer l’offre sans\\
Nexus.\\
• Handover planifié : jalons T0, T+12, T+24 mois définissant le passage progressif de la responsabilité pédagogique (R 𝑖𝑔ℎ𝑡𝑎𝑟𝑟𝑜𝑤 A) vers RBK si\\
décision d’internalisation.\\
\vspace{0.5em}
\vspace{0.5em}
\vspace{0.5em}
\vspace{0.5em}
\begin{confidential}CONFIDENTIEL — PROJET RBK 2.0                                          Page 65 sur 316\end{confidential}
RBK 2.0 — Whitepaper Architecte Web3                                                                                    Chapitre 5. Architecture de Partenariat\\
\vspace{0.5em}
\vspace{0.5em}
\vspace{0.5em}
a\\
Pédagogie — Processus d’évaluation et de validation des mentors RBK/Nexus couvrant expertise technique, soft skills et conformité aux standards de\\
mentoring Senior-by-Design.\\
\vspace{0.5em}
\vspace{0.5em}
Audits et continuité pédagogique\\
\vspace{0.5em}
Audits croisés : RBK mandate un cabinet tiers au moins une fois par an sur trois volets (académique, sécurité, finance) ; Nexus supporte la préparation\\
documentaire et l’accès aux environnements. Plan de continuité : un registre de ressources critiques (mentors, outils, données) est tenu par le Comité Risques ;\\
des remplaçants RBK pré-qualifiés sont contractualisés. Escalade structurée : incidents classés en 3 niveaux (Majeur, Critique, Catastrophique) avec délais de\\
réponse (4h, 12h, 24h) et processus de communication apprenant/employeur.\\
\vspace{0.5em}
\vspace{0.5em}
\vspace{0.5em}
\vspace{0.5em}
EL\\
TI\\
N\\
DE\\
FI\\
N\\
CO\\
\vspace{0.5em}
\vspace{0.5em}
\vspace{0.5em}
\vspace{0.5em}
\begin{confidential}CONFIDENTIEL — PROJET RBK 2.0                                           Page 66 sur 316\end{confidential}
RBK 2.0 — Whitepaper Architecte Web3                                                                                                Chapitre 5. Architecture de Partenariat\\
\vspace{0.5em}
\vspace{0.5em}
Monitoring quotidien\\
(Dashboards\\
MFAI, NPS)\\
\vspace{0.5em}
\vspace{0.5em}
\vspace{0.5em}
\vspace{0.5em}
Comité CEC\\
(RBK + Nexus)                           Non\\
Alerte KPIou Incident ?\\
Actions correc-\\
tives mineures\\
\vspace{0.5em}
\vspace{0.5em}
\vspace{0.5em}
\vspace{0.5em}
EL\\
Oui\\
\vspace{0.5em}
\vspace{0.5em}
\vspace{0.5em}
\vspace{0.5em}
TI\\
N\\
Comité Risques RBK\\
\vspace{0.5em}
\vspace{0.5em}
\vspace{0.5em}
\vspace{0.5em}
DE\\
(Qualité, Sécu-                                                    Audit ciblé\\
rité, Finance)                                                         (externe ou interne)\\
\vspace{0.5em}
\vspace{0.5em}
\vspace{0.5em}
\vspace{0.5em}
FI\\
N\\
CO\\
\vspace{0.5em}
Décision RBK                                              Plan d’action     Activation plan continuité\\
Go / Pause / Résilia-                                                    (mentors réserve, backup\\
tion conditionnelle                                                        LMS, communication)\\
\vspace{0.5em}
\vspace{0.5em}
\vspace{0.5em}
\vspace{0.5em}
\begin{confidential}CONFIDENTIEL — PROJET RBK 2.0          Principes directeurs : RBK conservePage\end{confidential}
le droit67  sur stratégique,\\
de veto  316        les jalons go/no-go sont\\
actés par procès-verbal, et chaque décision est tracée dans Skill Mirror Governance.\\
RBK 2.0 — Whitepaper Architecte Web3                                                                                Chapitre 5. Architecture de Partenariat\\
\vspace{0.5em}
\vspace{0.5em}
\vspace{0.5em}
\subsection{5.5 Matrice RACI macro révisée}
Matrice RACI — Partenariat RBK ↔ Nexus\\
\vspace{0.5em}
\vspace{0.5em}
\vspace{0.5em}
Activité clé                            RBK (Opérateur)        Nexus Réussite       Money Factory AI\\
(Maître d’œuvre        (Fournisseur)\\
pédagogique)\\
Marketing \textbackslash{}& acquisition                        AR                     C                      I\\
Contrats apprenants \textbackslash{}& ISA                      AR                      I                     I\\
Conception syllabus \textbackslash{}& standards                 C                    AR                      C\\
Recrutement \textbackslash{}& paiement des mentors               I                   AR                      I\\
\vspace{0.5em}
\vspace{0.5em}
\vspace{0.5em}
\vspace{0.5em}
EL\\
Animation cours \textbackslash{}& Block Checks                   I                   AR                      I\\
\vspace{0.5em}
\vspace{0.5em}
\vspace{0.5em}
\vspace{0.5em}
TI\\
Qualité pédagogique \textbackslash{}& audits internes           A                     R                      I\\
\vspace{0.5em}
\vspace{0.5em}
\vspace{0.5em}
\vspace{0.5em}
N\\
Plateforme LMS \textbackslash{}& support                         I                    A                     R\\
\vspace{0.5em}
\vspace{0.5em}
\vspace{0.5em}
\vspace{0.5em}
DE\\
Gestion des données \textbackslash{}& conformité DPA           AR                     R                      C\\
Recouvrement ISA \textbackslash{}& bonus anticipé              AR                     C                      I\\
\vspace{0.5em}
\vspace{0.5em}
\vspace{0.5em}
\vspace{0.5em}
FI\\
Reporting KPI \textbackslash{}& comités                         A                     R                      I\\
\vspace{0.5em}
\vspace{0.5em}
N\\
CO\\
Plan de continuité pédagogique                  A                     R                      C\\
Gestion incidents technologiques                 I                    C                     AR\\
\vspace{0.5em}
\vspace{0.5em}
Articulation contractuelle\\
\vspace{0.5em}
Chaque activité listée est adossée à un livrable contractuel explicitement décrit dans l’Annexe D. Les indicateurs de performance (NPS, taux de placement,\\
respect des SLA) conditionnent les bonus/malus appliqués sur la facturation Nexus.\\
\vspace{0.5em}
\vspace{0.5em}
\vspace{0.5em}
\vspace{0.5em}
\begin{confidential}CONFIDENTIEL — PROJET RBK 2.0                                         Page 68 sur 316\end{confidential}
\section{MÉTHODOLOGIE CYBORG 2.0}
\vspace{0.5em}
\vspace{0.5em}
\subsection{6.1 Philosophie Pédagogique : Intégration du Bien-être}
\vspace{0.5em}
\vspace{0.5em}
\vspace{0.5em}
\vspace{0.5em}
EL\\
Ď Approche Holistique du Développement Technique\\
\vspace{0.5em}
\vspace{0.5em}
\vspace{0.5em}
\vspace{0.5em}
TI\\
Principe Fondateur : L’excellence technique ne doit pas se faire au détriment du bien-être mental. Nos méthodes intègrent des mécanismes de prévention du\\
\vspace{0.5em}
\vspace{0.5em}
\vspace{0.5em}
\vspace{0.5em}
N\\
burnout qui maintiennent la charge cognitive optimale (zone de « productive struggle ») sans basculer dans l’overwhelm.\\
\vspace{0.5em}
\vspace{0.5em}
\vspace{0.5em}
\vspace{0.5em}
DE\\
Approche : Alternance rythmée d’efforts cognitifs intenses (pair programming, code reviews) et de moments de consolidation (pause active, veille technolo-\\
gique, mentorat doublé).\\
\vspace{0.5em}
\vspace{0.5em}
\vspace{0.5em}
\vspace{0.5em}
FI\\
N\\
CO\\
\subsection{6.2 Modèle 3 Niveaux RBK 3.0}
\vspace{0.5em}
\vspace{0.5em}
\vspace{0.5em}
\vspace{0.5em}
69\\
RBK 2.0 — Whitepaper Architecte Web3                                                                   Chapitre 6. MÉTHODOLOGIE CYBORG 2.0\\
\vspace{0.5em}
\vspace{0.5em}
\vspace{0.5em}
\vspace{0.5em}
\begin{center}\begin{tabular}TAB. 6.1 : Architecture pédagogique alignée sur le modèle 42 adapté RBK\end{tabular}\end{center}
\vspace{0.5em}
\vspace{0.5em}
Phase                Durée     Objectifs opérationnels                                   Validation (Block Check)\\
Genesis Pool         4 sem.    Sélection intensive, culture Dev DAO, maîtrise Rust de    Skill Mirror x3 + Block Check “Genesis” (NFT Skills “Ac-\\
base sans assistance IA.                                  cess”).\\
Explorer             12 sem.   Fondamentaux Web3, pipelines Git, culture protocolaire,   Block Check “Explorer” (NFT Skills “Core”).\\
autonomie sur mini-projets.\\
Bâtisseur            16 sem.   Projets pairés, spécialisation Solana/EVM/Product,        Block Check “Build” + audit croisé Skill Mirror.\\
contribution Dev DAO.\\
Architecte           16 sem.   Leadership technique, mentorat des Explorer/Bâtisseur,    Block Check “Architecte” (NFT Skills “Launch Proof”).\\
\vspace{0.5em}
\vspace{0.5em}
\vspace{0.5em}
\vspace{0.5em}
EL\\
industrialisation.\\
\vspace{0.5em}
\vspace{0.5em}
\vspace{0.5em}
\vspace{0.5em}
TI\\
Launchpad (option)   4 sem.    Incubation projet et levée de fonds, accompagnement       Launch Proof pitch + accord revenue share (30\textbackslash{}%\\
\vspace{0.5em}
\vspace{0.5em}
\vspace{0.5em}
\vspace{0.5em}
N\\
Nexus.                                                    Nexus).\\
\vspace{0.5em}
\vspace{0.5em}
\vspace{0.5em}
\vspace{0.5em}
DE\\
FI\\
\subsection{6.2.1 Cycle Skill Mirror \textbackslash{}& Block Check}
\vspace{0.5em}
\vspace{0.5em}
N\\
CO\\
Mercredi\\
Lundi        Mardi                           Jeudi         Vendredi\\
Dev DAO\\
PO Sync     Dev Deep                       Skill Mirror   Block Check\\
Sprint\\
\vspace{0.5em}
\vspace{0.5em}
\vspace{0.5em}
Vendredi 2\\
Lundi 2      Mardi 2       Mercredi 2        Jeudi 2\\
Retro Launch\\
Mentor Sync     Lab Work       Office Hour     Skill Mirror\\
Proof\\
\vspace{0.5em}
\vspace{0.5em}
\vspace{0.5em}
\vspace{0.5em}
\begin{confidential}CONFIDENTIEL — PROJET RBK 2.0                                   Page 70 sur 316\end{confidential}
RBK 2.0 — Whitepaper Architecte Web3                                                                                Chapitre 6. MÉTHODOLOGIE CYBORG 2.0\\
\vspace{0.5em}
\vspace{0.5em}
\begin{center}\begin{tabular}TAB. 6.3 : Mesures Anti-Burnout Intégrées dans le Cycle Hebdomadaire\end{tabular}\end{center}
\vspace{0.5em}
\vspace{0.5em}
Type de Mesure         Mise en Œuvre                                                             Indicateur de Suivi\\
Monitoring Mental      Check-ins hebdo (1-1 mentor-mentoré) + mood tracker                       Taux de complétion des check-ins, score de bien-être (1-10)\\
Pause Cognitive        2h/semaine ”Pause Active” (sport, jeu, art) + 1h veille douce             Temps dédié aux pauses, taux de participation\\
Limites Claires        Max 6h de deep-work par jour, blocage week-end, seuil de PR par jour      Respect des quotas hebdo, heures de connexion\\
Support Immédiat       Ligne directe ”Well-being” + réponse \textless{} 24h                                Temps de réponse moyen, nb tickets\\
\vspace{0.5em}
\vspace{0.5em}
\vspace{0.5em}
\vspace{0.5em}
\subsection{6.3 Genesis Pool : Bootcamp Sélectif}
\vspace{0.5em}
\vspace{0.5em}
\vspace{0.5em}
\vspace{0.5em}
EL\\
͆ Structure Genesis Pool — Sélection Niveau Zéro\\
\vspace{0.5em}
\vspace{0.5em}
\vspace{0.5em}
\vspace{0.5em}
TI\\
La Genesis Pool est une phase intensive de 4 semaines conçue pour évaluer la capacité cognitive des candidats et leur propension à résoudre des problèmes\\
complexes en Rust sans assistance IA (Copilot désactivé). Cette phase filtre les candidats selon leur potentiel d’apprentissage et leur capacité à s’intégrer au\\
\vspace{0.5em}
\vspace{0.5em}
\vspace{0.5em}
\vspace{0.5em}
N\\
dispositif Dev DAO.\\
\vspace{0.5em}
\vspace{0.5em}
\vspace{0.5em}
\vspace{0.5em}
DE\\
Objectifs :\\
\vspace{0.5em}
\vspace{0.5em}
\vspace{0.5em}
\vspace{0.5em}
FI\\
• Mesurer la capacité d’apprentissage en situation de stress cognitif\\
\vspace{0.5em}
\vspace{0.5em}
\vspace{0.5em}
\vspace{0.5em}
N\\
• Évaluer les concepts fondamentaux de Rust (ownership, borrowing, lifetimes)\\
CO\\
• Tester la persistance face à des défis techniques non triviaux\\
• Identifier les profils qui réussissent à apprendre sans assistance IA\\
\vspace{0.5em}
\vspace{0.5em}
\vspace{0.5em}
\vspace{0.5em}
\begin{confidential}CONFIDENTIEL — PROJET RBK 2.0                                              Page 71 sur 316\end{confidential}
RBK 2.0 — Whitepaper Architecte Web3                                                                                 Chapitre 6. MÉTHODOLOGIE CYBORG 2.0\\
\vspace{0.5em}
\vspace{0.5em}
\subsection{6.3.1 Structure et Évaluation}
\vspace{0.5em}
\vspace{0.5em}
\begin{center}\begin{tabular}TAB. 6.5 : Compétences et évaluations ciblées durant la Genesis Pool\end{tabular}\end{center}
\vspace{0.5em}
\vspace{0.5em}
Semaine       Compétence Cible                 Évaluation Spécifique\\
Semaine 1     Concepts Rust de base            Ownership, borrowing, types, erreurs (Result/Option)\\
Semaine 2     Structures et généricité         Enums, structs, traits, impls, patterns matching\\
Semaine 3     Programmation fonctionnelle      Closures, iterators, modules, tests unitaires\\
Semaine 4     Projet complet                   CLI tool Rust + parsing binaire + Ed25519 signing + première Skill Mirror\\
\vspace{0.5em}
\vspace{0.5em}
\vspace{0.5em}
\vspace{0.5em}
EL\\
\subsection{6.4 Gamification Web3 et Système de Réputation}
\vspace{0.5em}
\vspace{0.5em}
\vspace{0.5em}
\vspace{0.5em}
TI\\
N\\
̪ Mécanismes Genesis Pool → Launch Proof\\
\vspace{0.5em}
\vspace{0.5em}
\vspace{0.5em}
\vspace{0.5em}
DE\\
• Skill Mirror : à chaque jalon, les pairs évaluent le livrable via une grille Rubik (technique, documentation, sécurité). Deux validations positives débloquent\\
l’accès au Block Check.\\
\vspace{0.5em}
\vspace{0.5em}
\vspace{0.5em}
\vspace{0.5em}
FI\\
N\\
• NFT Skills : NFT non-transférables frappés sur Devnet Solana à chaque Block Check (Genesis, Explorer, Build, Architecte). Ils encodent le niveau, les badges\\
“Mentor Lead” et les heures de mentoring réalisées.\\
CO\\
• Launch Proof : certificat final combinant NFT Skills Architecte et SBT “Audit-Ready”. Nécessaire pour accéder au Launchpad et au partage de revenus.\\
• Dev DAO : les projets sont menés en “Dev DAO” (escouades auto-gérées). Chaque escouade gère son trésor fictif et soumet un rapport on-chain.\\
• Block Check : rituels hebdomadaires (vendredi) permettant de clôturer les sprints, générer des feedbacks horodatés et déclencher l’émission automatique\\
des points de réputation.\\
\vspace{0.5em}
\vspace{0.5em}
\vspace{0.5em}
\vspace{0.5em}
\begin{confidential}CONFIDENTIEL — PROJET RBK 2.0                                            Page 72 sur 316\end{confidential}
RBK 2.0 — Whitepaper Architecte Web3                                                                                Chapitre 6. MÉTHODOLOGIE CYBORG 2.0\\
\vspace{0.5em}
\vspace{0.5em}
\begin{center}\begin{tabular}TAB. 6.7 : Correspondance NFT Skills et responsabilités pédagogiques\end{tabular}\end{center}
\vspace{0.5em}
\vspace{0.5em}
NFT Skills                   Pré-requis                                   Responsabilités activées\\
Genesis Access               Skill Mirror validés en Genesis Pool         Accès Explorer, participation Dev DAO.\\
Explorer Core                Block Check Explorer réussi                  Pairing obligatoire, relecture de Skill Mirror niveau 1.\\
Build Master                 Block Check Bâtisseur + audit sécurisé       Mentorat Explorer, animation Skill Mirror, accès labs Validator.\\
Architecte Launch Proof      Block Check final + SBT Audit-Ready          Mentorat Bâtisseur, jury Block Check, éligibilité Launchpad.\\
Mentor Lead                  30h mentorat + NPS \textgreater{} 4.5                     Suivi cohortes Genesis, modération Skill Mirror.\\
\vspace{0.5em}
\vspace{0.5em}
\vspace{0.5em}
\vspace{0.5em}
EL\\
\subsection{6.5 Protocole Anti-Burnout}
\vspace{0.5em}
\vspace{0.5em}
\vspace{0.5em}
\vspace{0.5em}
TI\\
Ʋ Gouvernance Bien-être\\
(Responsable Bien-êtrea \textbackslash{}& Résilience)\\
\vspace{0.5em}
\vspace{0.5em}
\vspace{0.5em}
\vspace{0.5em}
N\\
DE\\
a\\
Organisation — Référent Nexus habilité à ajuster charge, deadlines et protocoles de repos pour maintenir la santé mentale et physique des apprenants.\\
\vspace{0.5em}
\vspace{0.5em}
\vspace{0.5em}
\vspace{0.5em}
FI\\
Objectif : déclencher des temps de récupération avant la rupture et protéger la qualité pédagogique.\\
Pilotage opérationnel :\\
\vspace{0.5em}
\vspace{0.5em}
N\\
CO\\
• Responsable Bien-être \textbackslash{}& Résilience (Nexus) : mandat adossé au contrat RBK ↔ Nexus, doté d’un pouvoir de veto sur les deadlines. Peut imposer\\
jusqu’à 48 heures de pause pédagogique ou redistribuer les livrables d’une squad.\\
• Boucle d’escalade : mentor référent → Responsable Bien-être → HoP Nexus → RBK Ops pour arbitrage budgétaire si besoin.\\
• Traçabilité : chaque décision est journalisée dans Venture Engine (ticket ”Health”) et partagée au Comité éthique (hebdomadaire).\\
Seuils quantitatifs déclencheurs :\\
1. Score Fatigue \textgreater{} 7/10 pendant deux semaines consécutives (questionnaire anonyme).\\
2. Sommeil \textless{} 5 h sur trois nuits consécutives (auto-tracking via formulaire).\\
3. Productivité : baisse \textgreater{} 30\textbackslash{}% du throughput squad (tickets ”Done”) ou absence à deux rituels critiques sur une semaine.\\
4. Signal pair : deux alertes ”burn” déclenchées par les camarades sur Skill Mirror → repos obligatoire.\\
\vspace{0.5em}
\vspace{0.5em}
\vspace{0.5em}
\vspace{0.5em}
\begin{confidential}CONFIDENTIEL — PROJET RBK 2.0                                            Page 73 sur 316\end{confidential}
RBK 2.0 — Whitepaper Architecte Web3                                                                             Chapitre 6. MÉTHODOLOGIE CYBORG 2.0\\
\vspace{0.5em}
\vspace{0.5em}
\subsection{6.5.1 Organigramme de Monitoring}
\vspace{0.5em}
Comité exécutif\\
conjoint (RBK/-\\
Nexus/MFAI)\\
\vspace{0.5em}
\vspace{0.5em}
\vspace{0.5em}
\vspace{0.5em}
Responsable Bien-être\\
Escalade vers cellule psychologique partenaire si besoin\\
\textbackslash{}& Résilience (Nexus)\\
\vspace{0.5em}
\vspace{0.5em}
\vspace{0.5em}
\vspace{0.5em}
EL\\
TI\\
Mentors référents                        RBK Ops \textbackslash{}& Carrière\\
\vspace{0.5em}
\vspace{0.5em}
\vspace{0.5em}
\vspace{0.5em}
N\\
(daily check)                           (soutien logistique)\\
\vspace{0.5em}
\vspace{0.5em}
\vspace{0.5em}
\vspace{0.5em}
DE\\
Squads apprenants\\
\textbackslash{}& délégués santé\\
\vspace{0.5em}
\vspace{0.5em}
\vspace{0.5em}
\vspace{0.5em}
FI\\
N\\
\subsection{6.5.2 Tableau de bord hebdomadaire}
CO\\
Monitoring santé cohorte\\
\vspace{0.5em}
\vspace{0.5em}
\vspace{0.5em}
Indicateur                    Seuil GO         Seuil Alerte    Action Responsable Bien-être\\
Score bien-être moyen         ≤ 5/10             ≥ 7/10        Pause 24h, coaching individuel, ajustement backlog.\\
Heures de sommeil               ≥ 6h           ≤ 5h trois      Mise en ”mode sécurité” : pas de livrable critique sous 48h.\\
auto-reportées                                    nuits\\
Présence rituels               ≥ 90\textbackslash{}%             ≤ 80\textbackslash{}%         Entretien croisé mentor + RBK Ops, adaptation planning.\\
obligatoires\\
\vspace{0.5em}
\vspace{0.5em}
\vspace{0.5em}
\vspace{0.5em}
\begin{confidential}CONFIDENTIEL — PROJET RBK 2.0                                        Page 74 sur 316\end{confidential}
RBK 2.0 — Whitepaper Architecte Web3                                                                 Chapitre 6. MÉTHODOLOGIE CYBORG 2.0\\
\vspace{0.5em}
\vspace{0.5em}
\vspace{0.5em}
Charge hebdomadaire    +/-10\textbackslash{}% cible   \textgreater{} 30\textbackslash{}% écart   Repriorisation via HoP, possibilité de décaler Block Check.\\
(tickets)\\
Alertes Skill Mirror        0            ≥2         Activation cellule support + suivi quotidien.\\
”burn”\\
\vspace{0.5em}
\vspace{0.5em}
\vspace{0.5em}
\vspace{0.5em}
EL\\
TI\\
N\\
DE\\
FI\\
N\\
CO\\
\vspace{0.5em}
\vspace{0.5em}
\vspace{0.5em}
\vspace{0.5em}
\begin{confidential}CONFIDENTIEL — PROJET RBK 2.0                               Page 75 sur 316\end{confidential}
\section{STRUCTURE DU CURSUS AMÉLIORÉE}
\vspace{0.5em}
\vspace{0.5em}
\subsection{7.1 Architecture pédagogique — 48 semaines de Genesis vers Launch Proof}
\vspace{0.5em}
\vspace{0.5em}
\vspace{0.5em}
\vspace{0.5em}
EL\\
Le cursus RBK 3.0 repose sur un modèle hybride inspiré de 42 et adapté Web3. Il est structuré en quatre blocs obligatoires et un Launchpad optionnel :\\
\vspace{0.5em}
\vspace{0.5em}
\vspace{0.5em}
\vspace{0.5em}
TI\\
• Genesis Pool : 4 semaines de bootcamp sélectif (120h, Skill Mirror intensif).\\
\vspace{0.5em}
\vspace{0.5em}
\vspace{0.5em}
\vspace{0.5em}
N\\
• Explorer : 12 semaines de consolidations fondamentales (Rust, Web3, GitOps) avec Block Check “Explorer”.\\
\vspace{0.5em}
\vspace{0.5em}
\vspace{0.5em}
\vspace{0.5em}
DE\\
• Bâtisseur : 16 semaines de spécialisation et de projets Dev DAO pairés.\\
\vspace{0.5em}
\vspace{0.5em}
\vspace{0.5em}
\vspace{0.5em}
FI\\
• Architecte : 16 semaines de leadership technique, mentorat et audit readiness.\\
\vspace{0.5em}
\vspace{0.5em}
\vspace{0.5em}
N\\
• Launchpad Nexus (optionnel) : 4 semaines pour incuber les projets et activer le partage de revenus 30\textbackslash{}% / 70\textbackslash{}%.\\
CO\\
* Note : Le Launchpad est optionnel et nécessite l’approbation de Nexus.\\
\vspace{0.5em}
ˌ Structure du Curriculum RBK 3.0\\
\vspace{0.5em}
Objectif : Former des constructeurs Web3 capables de livrer, auditer et opérer des protocoles Solana/EVM en moins d’un an, tout en préparant\\
l’incubation Launchpad.\\
Réorganisation : Passage de l’ancien découpage en niveaux au pipeline Genesis Pool →Explorer →Bâtisseur →Architecte, chaque Block Check\\
déclenchant un NFT Skills non transférable.\\
Logique : Les apprenants progressent de la sélection intensive (Genesis Pool) vers la consolidation (Explorer), l’accélération projet (Bâtisseur) et\\
la posture d’Architecte mentor/auditeur. Le Launchpad sécurise la dimension business et la répartition des revenus.\\
\vspace{0.5em}
\vspace{0.5em}
\vspace{0.5em}
\vspace{0.5em}
76\\
RBK 2.0 — Whitepaper Architecte Web3                                                            Chapitre 7. STRUCTURE DU CURSUS AMÉLIORÉE\\
\vspace{0.5em}
\vspace{0.5em}
\subsection{7.1.1 Diagramme de Structure du Curriculum (48 Semaines)}
\vspace{0.5em}
Cursus\\
48 Sem.\\
\vspace{0.5em}
\vspace{0.5em}
Genesis Pool                        Explorer                          Bâtisseur                         Architecte\\
\vspace{0.5em}
\vspace{0.5em}
S1-S4                            S5-S16                            S17-S32                            S33-S48\\
\vspace{0.5em}
\vspace{0.5em}
\vspace{0.5em}
\vspace{0.5em}
EL\\
\subsection{7.1.2 Timeline pédagogique – Semaines 1 à 48}
\vspace{0.5em}
Genesis Pool        Explorer          Bâtisseur         Architecte         Launchpad\\
\vspace{0.5em}
\vspace{0.5em}
\vspace{0.5em}
\vspace{0.5em}
TI\\
S1-S4             S5-S16            S17-S32           S33-S48            S49-S52*\\
\vspace{0.5em}
\vspace{0.5em}
\vspace{0.5em}
\vspace{0.5em}
N\\
DE\\
\subsection{7.2 Nouveau Track C : Web3 Product \textbackslash{}& Ecosystem Strategy}
\vspace{0.5em}
\vspace{0.5em}
\vspace{0.5em}
FI\\
N\\
Ż Spécialisation en Stratégie Produit et Écosystème Web3\\
CO\\
Objectif : Former des experts capables de concevoir, lancer et scaler des produits Web3 avec une compréhension approfondie de l’écosystème,\\
des communautés et des mécanismes économiques.\\
Public Cible : Profils business/finance avec intérêt pour la technologie, ou ingénieurs souhaitant évoluer vers des rôles de produit/solution\\
architect dans l’écosystème Web3.\\
Durée : 16 semaines, avec un accent sur la création de valeur et la compréhension des dynamiques de marché.\\
\vspace{0.5em}
\vspace{0.5em}
\vspace{0.5em}
\vspace{0.5em}
\begin{confidential}CONFIDENTIEL — PROJET RBK 2.0                                   Page 77 sur 316\end{confidential}
RBK 2.0 — Whitepaper Architecte Web3                                                                    Chapitre 7. STRUCTURE DU CURSUS AMÉLIORÉE\\
\vspace{0.5em}
\vspace{0.5em}
\subsection{7.2.1 Modules du Track C (16 semaines)}
\vspace{0.5em}
\vspace{0.5em}
\begin{center}\begin{tabular}TAB. 7.1 : Modules Détaillés du Track C - Web3 Product \textbackslash{}& Ecosystem Strategy\end{tabular}\end{center}
\vspace{0.5em}
\vspace{0.5em}
Module                                Contenu                                   Livrables                               Semaines\\
M1 : Tokenomics \textbackslash{}& Governance          Modèles économiques, jetons utili-        Whitepaper Tokenomics, Simulation       4 semaines\\
taires, DAOs, mécanismes de gouver-       Economique\\
nance\\
M2 : Écosystème \textbackslash{}& GTM                 Analyse des protocoles, stratégies de     Analyse de marché, Stratégie d’entrée   4 semaines\\
go-to-market, partnerships\\
M3 : Produit \textbackslash{}& UX                     Conception produit, expérience utilisa-   Figma UI/UX, Prototype interactif       4 semaines\\
\vspace{0.5em}
\vspace{0.5em}
\vspace{0.5em}
\vspace{0.5em}
EL\\
teur, interfaces Web3, onboarding\\
M4 : Lancement \textbackslash{}& Scale                Lancement de produit, croissance, ges-    Plan de lancement, Campagnes marke-     4 semaines\\
\vspace{0.5em}
\vspace{0.5em}
\vspace{0.5em}
\vspace{0.5em}
TI\\
tion de communauté, métriques             ting\\
\vspace{0.5em}
\vspace{0.5em}
\vspace{0.5em}
\vspace{0.5em}
N\\
DE\\
\subsection{7.2.2 Compétences Visées du Track C}
\vspace{0.5em}
\vspace{0.5em}
\vspace{0.5em}
\vspace{0.5em}
FI\\
N\\
Tokenomics              Gouvernance\\
CO\\
\vspace{0.5em}
Produit                  Marché\\
\vspace{0.5em}
\vspace{0.5em}
\vspace{0.5em}
\vspace{0.5em}
\subsection{7.3 Timeline pédagogique – Semaines 1 à 48}
\vspace{0.5em}
\vspace{0.5em}
\vspace{0.5em}
\vspace{0.5em}
\begin{confidential}CONFIDENTIEL — PROJET RBK 2.0                                         Page 78 sur 316\end{confidential}
RBK 2.0 — Whitepaper Architecte Web3                                                          Chapitre 7. STRUCTURE DU CURSUS AMÉLIORÉE\\
\vspace{0.5em}
\vspace{0.5em}
\vspace{0.5em}
\vspace{0.5em}
\begin{center}\begin{tabular}TAB. 7.3 : Timeline pédagogique détaillée (Genesis Pool →Launchpad)\end{tabular}\end{center}
\vspace{0.5em}
\vspace{0.5em}
Semaine      Mois     Jalons pédagogiques                                  Évaluations \textbackslash{}& livrables\\
1–4        M1       Genesis Pool – Sélection Rust, culture Dev DAO       3 Skill Mirror + Block Check “Genesis” (NFT Access)\\
5–8        M2       Explorer – Fondamentaux Rust/Web3                    Skill Mirror Fundamentals, mini-projets pairés\\
9–12       M3       Explorer – Protocoles, sécurité de base              Block Check “Explorer” + Portfolio pack\\
13–16       M4       Explorer – Delivery produit, UX wallet               Block Check “Explorer” consolidation\\
17–20       M5       Bâtisseur – Solana/EVM track, Dev DAO sprint         Skill Mirror avançés, Block Check “Build” 1\\
21–24       M6       Bâtisseur – Architecture \textbackslash{}& audits croisés            Audit croisé, Block Check “Build” 2\\
\vspace{0.5em}
\vspace{0.5em}
\vspace{0.5em}
\vspace{0.5em}
EL\\
25–28       M7       Bâtisseur – Indexation, monitoring, DevOps           Release candidate + NFT Skills “Build Master”\\
\vspace{0.5em}
\vspace{0.5em}
\vspace{0.5em}
\vspace{0.5em}
TI\\
29–32       M8       Bâtisseur – Capstone Track (Validator/Product)       Jury Dev DAO, pitch technique\\
33–36       M9       Architecte – Leadership, mentorat Explorer           Skill Mirror modéré, Block Check “Architecte” 1\\
\vspace{0.5em}
\vspace{0.5em}
\vspace{0.5em}
\vspace{0.5em}
N\\
37–40       M10      Architecte – Audit, sécurité, opérations             Audit pack + NFT Skills “Mentor Lead”\\
\vspace{0.5em}
\vspace{0.5em}
\vspace{0.5em}
\vspace{0.5em}
DE\\
41–44       M11      Architecte – Industrialisation, runbooks             Block Check “Architecte” final (Launch Proof)\\
\vspace{0.5em}
\vspace{0.5em}
\vspace{0.5em}
\vspace{0.5em}
FI\\
45–48       M12      Architecte – Déploiement marché, alliances           Soutenance Launch Proof, signature SBT Audit-Ready\\
49–52*      M13      Launchpad (option) – Incubation, revenue model       Pitch investisseur, pacte Nexus (30\textbackslash{}% / 70\textbackslash{}%)\\
\vspace{0.5em}
N\\
CO\\
* Activation Launchpad : dossier validé par Nexus.\\
\vspace{0.5em}
\vspace{0.5em}
\vspace{0.5em}
\vspace{0.5em}
\begin{confidential}CONFIDENTIEL — PROJET RBK 2.0                                   Page 79 sur 316\end{confidential}
\section{SYLLABUS TECHNIQUE COMPLET}
(48 SEMAINES)\\
\vspace{0.5em}
\vspace{0.5em}
\vspace{0.5em}
\vspace{0.5em}
EL\\
Note : Ce chapitre détaille l’exécution technique. L’intégration des Soft Skills (S45-S48) est traitée au Chapitre 7.\\
\vspace{0.5em}
\vspace{0.5em}
\vspace{0.5em}
\vspace{0.5em}
TI\\
N\\
\subsection{8.1 Calendrier Pédagogique Global}
\vspace{0.5em}
\vspace{0.5em}
\vspace{0.5em}
\vspace{0.5em}
DE\\
FI\\
N\\
Genesis Pool               Explorer                Bâtisseur            Architecte               Launchpad*\\
(S1-S4)                  (S5-S16)                (S17-S32)            (S33-S48)                 (S49-S52)\\
CO\\
Sélection, Skill Mirror   Rust, Web3, GitOps     Solana/EVM, Dev DAO   Mentorat, Audits, Ops   Incubation, revenue share\\
\vspace{0.5em}
\vspace{0.5em}
\vspace{0.5em}
FIG. 8.1 : Timeline Macro du Cursus\\
\vspace{0.5em}
* Activation Launchpad conditionnée à la validation du dossier Nexus (pacte 30\textbackslash{}% / 70\textbackslash{}%).\\
\vspace{0.5em}
GENESIS POOL (S1-S4)\\
\vspace{0.5em}
\subsection{8.2 Genesis Pool : Sélection Intensives (S1-S4)}
\vspace{0.5em}
\vspace{0.5em}
80\\
RBK 2.0 — Whitepaper Architecte Web3                                                                                   Chapitre 8. SYLLABUS TECHNIQUE\\
\vspace{0.5em}
\vspace{0.5em}
\vspace{0.5em}
\vspace{0.5em}
\begin{center}\begin{tabular}TAB. 8.1 : Genesis Pool – Synthèse des 4 semaines\end{tabular}\end{center}
\vspace{0.5em}
\vspace{0.5em}
Sem.     Focus Technique        Livrable Pivot                Skill Mirror / Block Check\\
S1       OS \textbackslash{}& Git Internals     Réplique ‘ls -la‘ en Rust     Skill Mirror “Core Shell”\\
S2       Memory Safety          Custom Allocator              Skill Mirror “Memory”\\
S3       Concurrence            HTTP Server multi-thread      Skill Mirror “Concurrency”\\
S4       Cryptographie          CLI Wallet (Ed25519)          Block Check “Genesis” + NFT Access\\
\vspace{0.5em}
\vspace{0.5em}
La Stack du Vainqueur\\
\vspace{0.5em}
\vspace{0.5em}
\vspace{0.5em}
\vspace{0.5em}
EL\\
• Rust: Langage système, sécurité mémoire garantie sans Garbage Collector.\\
\vspace{0.5em}
\vspace{0.5em}
\vspace{0.5em}
\vspace{0.5em}
TI\\
• Anchor: Framework de développement Solana qui sécurise et accélère le code.\\
\vspace{0.5em}
\vspace{0.5em}
\vspace{0.5em}
\vspace{0.5em}
N\\
DE\\
• Solidity: Langage historique des Smart Contracts (EVM).\\
\vspace{0.5em}
\vspace{0.5em}
\vspace{0.5em}
\vspace{0.5em}
FI\\
• Foundry: Outil de test et déploiement Ethereum écrit en Rust.\\
\vspace{0.5em}
\vspace{0.5em}
\vspace{0.5em}
\vspace{0.5em}
N\\
CO\\
\subsection{8.2.1 Plan d’Exécution Explorer (S05-S16)}
GATE1 structure les jalons.\\
Les contrôles CD ancrent les exigences de qualité dès la formation.\\
\vspace{0.5em}
\vspace{0.5em}
\vspace{0.5em}
\vspace{0.5em}
1\\
Projet/Management/Qualité — Point de décision formalisé (GO/NO-GO) après vérification build, scripts et absence de hardcodes ; bloque la promotion en cas\\
d’échec.\\
\vspace{0.5em}
\vspace{0.5em}
\begin{confidential}CONFIDENTIEL — PROJET RBK 2.0                                               Page 81 sur 316\end{confidential}
RBK 2.0 — Whitepaper Architecte Web3                                                                                     Chapitre 8. SYLLABUS TECHNIQUE\\
\vspace{0.5em}
\vspace{0.5em}
Sem     Cours (Objectifs)                    Labs (TP Guidés)                   Livrable (Preuve)        Critères (Skill Mirror\\
/ DoD)\\
S05     Crypto 2 : signatures, keys,         Sign/verify, key mgmt              Module signature         Skill Mirror “Crypto”\\
threats                                                                                          validée, gestion clés OK\\
S06     Blockchain 1 : tx, RPC2 , finality   Client RPC, parsing tx             ”explorer” CLI           Résilience réseau, Skill\\
Mirror “Explorer”\\
S07     API Design : REST, cache, limits     API endpoints, OpenAPI3            Micro-API + Spec         OpenAPI validée,\\
latence \textless{} 200ms\\
S08     Indexation : events, storage local   Indexer minimal, backfill          Indexer + Wallet API     Données cohérentes,\\
Skill Mirror “Data”\\
S09     UX4 Wallet : signature, pending,     Flow connect/sign/read             Mini front/CLI UX        Tests UI automatisés,\\
errors                                                                                           satisfaction pair \textgreater{} 4/5\\
\vspace{0.5em}
\vspace{0.5em}
\vspace{0.5em}
\vspace{0.5em}
EL\\
S10     Tokenomics 1 : supply,               Simu simple, stress tests          ”tokenomics memo”        Hypothèses explicites,\\
incentives                                                                                       revue croisée validée\\
\vspace{0.5em}
\vspace{0.5em}
\vspace{0.5em}
\vspace{0.5em}
TI\\
S11     Intégration E2E : vertical slice     1 user story complète              Mini-dApp E2E            Démo fonctionnelle,\\
tests E2E OK\\
\vspace{0.5em}
\vspace{0.5em}
\vspace{0.5em}
\vspace{0.5em}
EN\\
S12     Observabilité de base, métriques     Logs structurés, dashboards        Dashboard Grafana        Alertes configurées,\\
Skill Mirror “Observ”\\
\vspace{0.5em}
\vspace{0.5em}
\vspace{0.5em}
\vspace{0.5em}
D\\
S13     Sécurité applicative (OWASP          Threat modeling, lint sécurité     Threat Model v1          Validation mentor, zéro\\
Web3)                                                                                            faille critique\\
\vspace{0.5em}
\vspace{0.5em}
\vspace{0.5em}
\vspace{0.5em}
FI\\
S14     Product Delivery cadence             Retro Dev DAO, velocity tracking   Rapport Dev DAO          Velocity stable,\\
\vspace{0.5em}
N                                                          feedback pair \textgreater{} 4/5\\
CO\\
S15     Préparation Block Check              Hardening, documentation           Release Candidate        Tous tests verts, doc\\
Explorer                                                                Explorer                 audit prête\\
S16     Block Check “Explorer”               Soutenance, panel pairs/mentors    NFT Skills Explorer      Score ≥ 88/100 (tech\\
Core                     + soft), vote mentor\\
\vspace{0.5em}
\vspace{0.5em}
\vspace{0.5em}
\vspace{0.5em}
2\\
Projet/Management/Qualité — Point de décision formalisé (GO/NO-GO) après vérification build, scripts et absence de hardcodes ; bloque la promotion en cas d’échec.\\
3\\
Projet/Management/Qualité — Point de décision formalisé (GO/NO-GO) après vérification build, scripts et absence de hardcodes ; bloque la promotion en cas d’échec.\\
4\\
Projet/Management/Qualité — Point de décision formalisé (GO/NO-GO) après vérification build, scripts et absence de hardcodes ; bloque la promotion en cas d’échec.\\
\vspace{0.5em}
\vspace{0.5em}
\begin{confidential}CONFIDENTIEL — PROJET RBK 2.0                                          Page 82 sur 316\end{confidential}
RBK 2.0 — Whitepaper Architecte Web3                                                                     Chapitre 8. SYLLABUS TECHNIQUE\\
\vspace{0.5em}
\vspace{0.5em}
\vspace{0.5em}
BÂTISSEUR (S17-S32)\\
\vspace{0.5em}
\subsection{8.3 Tracks Bâtisseur A / B / C (S17-S32)}
Tracks : A (Solana), B (EVM), C (Product). Les semaines sont synchronisées et alignées sur la cadence Dev DAO.\\
\vspace{0.5em}
\vspace{0.5em}
\vspace{0.5em}
\vspace{0.5em}
EL\\
TI\\
D    EN\\
FI\\
N\\
CO\\
\vspace{0.5em}
\vspace{0.5em}
\vspace{0.5em}
\vspace{0.5em}
\begin{confidential}CONFIDENTIEL — PROJET RBK 2.0                                 Page 83 sur 316\end{confidential}
RBK 2.0 — Whitepaper Architecte Web3                                                                                      Chapitre 8. SYLLABUS TECHNIQUE\\
\vspace{0.5em}
\vspace{0.5em}
Sem Cours (A/B/C)                       Labs (A/B/C)                    Livrable                  DoD / Skill Mirror\\
S17 A : Modèle Solana / B :             A : Prog natif / B : ERC5 min / Repo S17 (Tracké)         Tests + README6 ,\\
Solidity bases / C : Discovery      C : User Stories                                          Skill Mirror\\
“Kick-off”\\
S18 A : Sérialisation, Budget / B : A : Constraints / B : Events / PR + ADR \textbackslash{}#1                   ADR clair,\\
Gas, Erreurs / C : Wireframes C : Proto.                                                     reviewable\\
S19 A : Sécurité Auth. / B : Access A : Auth. / B : Roles / C :        ”Module pack” v1          Checklist sécu, Skill\\
Ctrl / C : KPIs                  Metrics Plan                                                Mirror “Security”\\
S20 A : Tests (Locaux) / B :         A : Tests état / B : Fuzz intro / PRD/Tool Gate             Pipeline test OK\\
Foundry Intro / C : PRD Final C : Roadmap\\
S21 A : Anchor IDL / B : Foundry A : Refactor Anchor / B : Fuzz Repo S21                          Tests exec, doc\\
\vspace{0.5em}
\vspace{0.5em}
\vspace{0.5em}
\vspace{0.5em}
EL\\
Mastery / C : Tokenomics         / C : Incentives\\
S22 A : Anchor Constraints / B :     A : Seeds/PDA / B : ERCs / C : ”Simu pack”                   Résultats repro., Skill\\
\vspace{0.5em}
\vspace{0.5em}
\vspace{0.5em}
\vspace{0.5em}
TI\\
ERC Stds / C : Simu.             Scénarios                                                    Mirror “Architecture”\\
\vspace{0.5em}
\vspace{0.5em}
\vspace{0.5em}
\vspace{0.5em}
N\\
S23 A : CPI / B : Architecture / C : A : CPI / B : Timelocks / C :     ADR \textbackslash{}#2 + PR                Menaces listées\\
\vspace{0.5em}
\vspace{0.5em}
\vspace{0.5em}
\vspace{0.5em}
DE\\
Risques                          Mitigations\\
S24 A : Hardening / B : Invariants A : Edge cases / B : Coverage / Module 2 Clôture               Audit notes\\
\vspace{0.5em}
\vspace{0.5em}
\vspace{0.5em}
\vspace{0.5em}
FI\\
/ C : Paper                      C : Paper v1\\
\vspace{0.5em}
\vspace{0.5em}
N\\
S25 A : DeFi Primitives / B : DApp A : Vault / B : Sign Flow / C : Repo S25                       Logs propres, Skill\\
Intég. / C : Analytics           Tracking                                                     Mirror “Analytics”\\
CO\\
S26 A : Pricing / B : Oracles / C : A : Calculs / B : Oracle mock / Dash/Indexer v1               Métriques lisibles\\
Dashboards                       C : Dash v1\\
S27 A : Indexation / B : Events /    A : Events / B : Endpoints / C : ”Data pack”                 Schéma doc, Skill\\
C : Cohortes                     Funnels                                                      Mirror “Data Ops”\\
S28 A : DeFi Pack / B : L2 Scale / A : Stable / B : Archi. L2 / C : Module 3 Clôture              Démo 10 min, vote\\
C : Analytics Final              Final Instr.                                                 mentors\\
\vspace{0.5em}
\vspace{0.5em}
\vspace{0.5em}
\vspace{0.5em}
5\\
Projet/Management/Qualité — Point de décision formalisé (GO/NO-GO) après vérification build, scripts et absence de hardcodes ; bloque la promotion en cas d’échec.\\
6\\
Projet/Management/Qualité — Point de décision formalisé (GO/NO-GO) après vérification build, scripts et absence de hardcodes ; bloque la promotion en cas d’échec.\\
\vspace{0.5em}
\vspace{0.5em}
\begin{confidential}CONFIDENTIEL — PROJET RBK 2.0                                            Page 84 sur 316\end{confidential}
RBK 2.0 — Whitepaper Architecte Web3                                                                                      Chapitre 8. SYLLABUS TECHNIQUE\\
\vspace{0.5em}
\vspace{0.5em}
Sem Cours (A/B/C)                 Labs (A/B/C)                     Livrable                        DoD / Skill Mirror\\
S29 A : Prod Perf / B : Sec       A : Bench / B : Threat Model7 Repo S29 + TM                      TM présent, CI verte\\
Hardening / C : Gouv. Design / C : Rules\\
S30 A : Monitoring / B :          A : Metrics / B : Fix findings / ”Ops pack”                      Runbook8 incidents\\
Correction \textbackslash{}& Preuve / C : GTM C : Launch plan\\
S31 A : Wallet Intég. / B : Audit A : Adapters / B : Audit pack / Release Cand.                    Tag RC, Changelog\\
Notes / C : Playbook          C : Metrics\\
S32 Block Check “Build”           Stabilisation + Docs             PACK Bâtisseur                 GO/WARN/NO-GO\\
+ NFT “Build\\
Master”\\
\vspace{0.5em}
\vspace{0.5em}
\vspace{0.5em}
\vspace{0.5em}
EL\\
TI\\
D     EN\\
FI\\
N\\
CO\\
\vspace{0.5em}
\vspace{0.5em}
\vspace{0.5em}
\vspace{0.5em}
7\\
Projet/Management/Qualité — Point de décision formalisé (GO/NO-GO) après vérification build, scripts et absence de hardcodes ; bloque la promotion en cas d’échec.\\
8\\
Projet/Management/Qualité — Point de décision formalisé (GO/NO-GO) après vérification build, scripts et absence de hardcodes ; bloque la promotion en cas d’échec.\\
\vspace{0.5em}
\vspace{0.5em}
\begin{confidential}CONFIDENTIEL — PROJET RBK 2.0                                            Page 85 sur 316\end{confidential}
RBK 2.0 — Whitepaper Architecte Web3                                                Chapitre 8. SYLLABUS TECHNIQUE\\
\vspace{0.5em}
\vspace{0.5em}
\vspace{0.5em}
NIVEAU 3 : STUDIO \textbackslash{}& PROFESSIONNALISATION (S29-S48)\\
\vspace{0.5em}
\subsection{8.4 STUDIO DE PRODUCTION (S29-S48)}
Objectif : Livrer des artefacts employables (Capstones) et préparer l’audit.\\
\vspace{0.5em}
\vspace{0.5em}
\vspace{0.5em}
\vspace{0.5em}
EL\\
TI\\
D    EN\\
FI\\
N\\
CO\\
\vspace{0.5em}
\vspace{0.5em}
\vspace{0.5em}
\vspace{0.5em}
\begin{confidential}CONFIDENTIEL — PROJET RBK 2.0                                     Page 86 sur 316\end{confidential}
RBK 2.0 — Whitepaper Architecte Web3                                                                                      Chapitre 8. SYLLABUS TECHNIQUE\\
\vspace{0.5em}
\vspace{0.5em}
Sem Objectifs (Sprint)                   Labs (Production)                Livrable                 DoD\\
S29 Capstone 1 : Reliability, ADR,       ADR \textbackslash{}#3, Threat Model, Spec       Spec + ADR + TM          Menaces mitigées\\
TM\\
S30 Sessions, RPC Fallback               Impl. fallback, retries          Module Session           Pas de blocage\\
S31 Tx Builder, Simulation               Builder, Idempotence             Tx Pipeline v1           Scénarios test\\
S32 Observabilité, Support               Metrics, Traces, Playbook        Obs. Pack                Dashboards\\
S33 GATE CAPSTONE 1                      Tests E2E, Démo filmée           Reliability Pack         Démo stable\\
S34 Capstone 2 : Protocole               Invariants, Cas limites          Spec + Invariants        Stratégie tests\\
On-chain\\
S35 Implémentation Core                  Core v1, Unit tests              Core v1                  CI verte\\
S36 Sécurité                             Simu attaques, Mitigations       Security Checklist       Preuves fixes\\
\vspace{0.5em}
\vspace{0.5em}
\vspace{0.5em}
\vspace{0.5em}
EL\\
S37 Audit Readiness                      Scripts repro, Diagrammes        Audit Pack v1            Repro 1 cmd\\
\vspace{0.5em}
\vspace{0.5em}
\vspace{0.5em}
\vspace{0.5em}
TI\\
S38 GATE CAPSTONE 2                      Hardening, Perf                  RC + Audit Pack          RC Tag\\
S39 Capstone 3 : Ops, SLO                SLO, Alerting                    SLO + Dashboards         Alertes utiles\\
\vspace{0.5em}
\vspace{0.5em}
\vspace{0.5em}
\vspace{0.5em}
EN\\
S40 Incident Simu                        Chaos monkey, Drill              Postmortem Simulé        Actions corr.\\
S41 Deploy Prod-like                     Pipeline, Migration              Ops Playbook v1          Rollback doc\\
\vspace{0.5em}
\vspace{0.5em}
\vspace{0.5em}
\vspace{0.5em}
D\\
S42 Perf \textbackslash{}& Cout                          Profiling, Optim                 Rapport Perf             Mesures\\
\vspace{0.5em}
\vspace{0.5em}
\vspace{0.5em}
\vspace{0.5em}
FI\\
S43 GATE CAPSTONE 3                      Go-live simulé                   Go-Live Pack             Runbook\\
S44 Placement 1 : GitHub9                Nettoyage, README                Portfolio Pack           Proof links\\
S45 Placement 2 : CV/LinkedIn            N\\
”Project one-liners”             CV + LinkedIn            Claims vérifiables\\
CO\\
S46 Placement 3 : Demos                  Script, Vidéo, Slides            Demo Kit                 Timing respecté\\
S47 Placement 4 : Interviews             Mock, System Design              Interview Notes          Réponses rigoureuses\\
S48 GATE FINAL Architecte                Soutenance                       PACK FINAL               GO MARCHÉ\\
\vspace{0.5em}
\vspace{0.5em}
\vspace{0.5em}
\vspace{0.5em}
9\\
Projet/Management/Qualité — Point de décision formalisé (GO/NO-GO) après vérification build, scripts et absence de hardcodes ; bloque la promotion en cas d’échec.\\
\vspace{0.5em}
\vspace{0.5em}
\begin{confidential}CONFIDENTIEL — PROJET RBK 2.0                                            Page 87 sur 316\end{confidential}
RBK 2.0 — Whitepaper Architecte Web3                                                                  Chapitre 8. SYLLABUS TECHNIQUE\\
\vspace{0.5em}
\vspace{0.5em}
\vspace{0.5em}
EXTENSIONS \textbackslash{}& OPTIONS\\
\vspace{0.5em}
\subsection{8.5 Modules de Diversification (Electifs)}
Pour les étudiants souhaitant élargir leur spectre technique, RBK 2.0 propose des modules intensifs accessibles en parallèle ou post-\\
cursus.\\
\vspace{0.5em}
\vspace{0.5em}
\subsection{8.5.1 Module ZK : Zero-Knowledge Proofs (8 semaines)}
• Contenu : Arithmétisation, R1CS, Plonk, Langage Noir et Circom.\\
\vspace{0.5em}
\vspace{0.5em}
\vspace{0.5em}
\vspace{0.5em}
EL\\
• Projet : Concevoir un mélangeur de tokens (Mixer) compliant (Privacy Pools).\\
\vspace{0.5em}
\vspace{0.5em}
\vspace{0.5em}
\vspace{0.5em}
TI\\
• Pré-requis : Niveau Mathématique A+ (Algèbre linéaire).\\
\vspace{0.5em}
\vspace{0.5em}
\vspace{0.5em}
\vspace{0.5em}
EN\\
\subsection{8.5.2 Module DePIN : Decentralized Physical Infra (6 semaines)}
\vspace{0.5em}
\vspace{0.5em}
\vspace{0.5em}
\vspace{0.5em}
D\\
FI\\
• Contenu : Helium Network, Filecoin, IoT Integration, Proof of Coverage.\\
\vspace{0.5em}
\vspace{0.5em}
N\\
• Projet : Déployer un réseau de capteurs LoRaWAN incentivé par token.\\
CO\\
\subsection{8.5.3 Module Cross-Chain \textbackslash{}& Interop (4 semaines)}
• Contenu : Wormhole, LayerZero, Axelar. Design de messages asynchrones.\\
\vspace{0.5em}
• Projet : Bridge NFT Solana ↔ Ethereum.\\
\vspace{0.5em}
\vspace{0.5em}
\vspace{0.5em}
\vspace{0.5em}
\begin{confidential}CONFIDENTIEL — PROJET RBK 2.0                               Page 88 sur 316\end{confidential}
\section{TRACK A : SOLANA SMART CONTRACT ENGINEER}
(RUST/ANCHOR)\\
\vspace{0.5em}
\vspace{0.5em}
\vspace{0.5em}
\vspace{0.5em}
EL\\
TI\\
\subsection{9.1 Philosophie du Track : L’Excellence par Rust}
\vspace{0.5em}
\vspace{0.5em}
\vspace{0.5em}
\vspace{0.5em}
EN\\
Solana n’est pas une simple blockchain rapide ; c’est un système d’exploitation décentralisé massivement parallèle (Sealevel). Pour y\\
développer, comprendre la syntaxe ne suffit pas. Il faut maîtriser la gestion mémoire, la concurrence d’accès aux données (Account Model)\\
\vspace{0.5em}
\vspace{0.5em}
\vspace{0.5em}
\vspace{0.5em}
D\\
et l’optimisation des cycles CPU (Compute Units). Le choix de Rust n’est pas anodin : il impose une rigueur absolue (Safety) qui, combinée\\
\vspace{0.5em}
\vspace{0.5em}
\vspace{0.5em}
\vspace{0.5em}
FI\\
aux contraintes de Solana, forme des ingénieurs d’élite. Notre objectif est de former des ”Guardians” : des développeurs obsédés par la\\
sécurité des fonds, la performance du code et la résilience de l’architecture.\\
N\\
CO\\
Positionnement ”Guardian” Un Guardian ne se contente pas de coder une feature. Il pense ”adversarial”. Il sait comment une transaction\\
peut échouer, comment un attaquant peut manipuler une instruction, et comment le réseau va réagir sous charge. C’est un profil hybride\\
entre Architecte Système et Auditeur de Sécurité.\\
\vspace{0.5em}
\vspace{0.5em}
\vspace{0.5em}
\vspace{0.5em}
89\\
RBK 2.0 — Whitepaper Architecte Web3                                                                    Chapitre 9. Track A : Solana Engineer\\
\vspace{0.5em}
\vspace{0.5em}
\vspace{0.5em}
NON NÉGOCIABLE : LE STANDARD QUALITÉ\\
• Tests Systematiques : Pas de PR sans tests (Unit + E2E). Coverage \textgreater{} 80\textbackslash{}%.\\
\vspace{0.5em}
• Reproductibilité : Builds déterministes (Verifiable Builds).\\
\vspace{0.5em}
• Audit-Readiness : Code commenté, Documentation d’architecture jour 1, Threat Model explicite.\\
\vspace{0.5em}
\vspace{0.5em}
\vspace{0.5em}
\vspace{0.5em}
EL\\
Rust                      SVM\\
Safety                     Perf\\
\vspace{0.5em}
\vspace{0.5em}
\vspace{0.5em}
\vspace{0.5em}
TI\\
The Guardian\\
\vspace{0.5em}
\vspace{0.5em}
\vspace{0.5em}
\vspace{0.5em}
EN\\
Security\\
\vspace{0.5em}
\vspace{0.5em}
\vspace{0.5em}
\vspace{0.5em}
D\\
FI\\
FIG. 9.1 : Pourquoi Solana est un track d’excellence\\
N\\
O\\
\begin{center}\begin{tabular}TAB. 9.1 : Compétences Cibles vs Preuves\end{tabular}\end{center}
C\\
\vspace{0.5em}
\vspace{0.5em}
Domaine      Compétence     Preuve Attendue                 Standard\\
Architecture Gestion PDAs   Diagramme de dérivations        Pas de collision\\
Sécurité     Signer Checks  Tests d’invocation malveillante 100\textbackslash{}% checked\\
Performance CU Optimization Rapport profilage transaction \textless{} 200k CU\\
\vspace{0.5em}
\vspace{0.5em}
\vspace{0.5em}
\vspace{0.5em}
\begin{confidential}CONFIDENTIEL — PROJET RBK 2.0                                   Page 90 sur 316\end{confidential}
RBK 2.0 — Whitepaper Architecte Web3                                                                 Chapitre 9. Track A : Solana Engineer\\
\vspace{0.5em}
\vspace{0.5em}
\vspace{0.5em}
\subsection{9.2 Structure Pédagogique : De l’Architecture au Produit (16 Semaines)}
Le cursus s’articule autour de la montée en puissance technique. (Voir le détail semaine par semaine dans le Syllabus Opérationnel en\\
Annexe B.7).\\
Le parcours est découpé en 4 modules progressifs. On commence par ”souffrir” avec Rust natif pour comprendre la mécanique interne,\\
puis on accélère avec Anchor, avant de plonger dans l’architecture complexe (CPI) et le durcissement pour la production. Chaque module\\
se solde par un ”Livrable Portfolio” qui prouve l’acquisition de la compétence.\\
\vspace{0.5em}
Syllabus Opérationnel\\
\vspace{0.5em}
Le syllabus opérationnel détaillé est présenté dans l’Annexe B.7.\\
\vspace{0.5em}
\vspace{0.5em}
\vspace{0.5em}
\vspace{0.5em}
EL\\
\begin{center}\begin{tabular}TAB. 9.3 : Carte des Modules (Résumé Exécutif)\end{tabular}\end{center}
\vspace{0.5em}
\vspace{0.5em}
\vspace{0.5em}
\vspace{0.5em}
TI\\
EN\\
Module Sem Objectif                        Lab Principal    Portfolio\\
1. Native 13-16 Comprendre l’Account Model Mini-Vault Natif Repo ”Raw”\\
\vspace{0.5em}
\vspace{0.5em}
\vspace{0.5em}
\vspace{0.5em}
D\\
2. Anchor 17-20 Productivité               Marketplace NFT Program IDL\\
\vspace{0.5em}
\vspace{0.5em}
\vspace{0.5em}
\vspace{0.5em}
FI\\
3. Arch   21-24 Composabilité              CPI Orchestrator Diagramme Arch\\
4. Prod\\
N\\
25-28 Hardening                  Full dApp        Audit Report\\
CO\\
\vspace{0.5em}
\subsection{9.2.1 MODULE 1 : Le Modèle Solana \textbackslash{}& Rust Natif (Semaines 13-16)}
Objectifs Opérationnels\\
\vspace{0.5em}
• Maîtriser la dé-sérialisation manuelle des données (Borsh).\\
\vspace{0.5em}
• Gérer le ”Rent” et l’allocation d’espace (realloc).\\
\vspace{0.5em}
• Comprendre le cycle de vie d’une transaction (Signer, Writable).\\
\vspace{0.5em}
\vspace{0.5em}
\vspace{0.5em}
\begin{confidential}CONFIDENTIEL — PROJET RBK 2.0                                Page 91 sur 316\end{confidential}
RBK 2.0 — Whitepaper Architecte Web3                                                                                   Chapitre 9. Track A : Solana Engineer\\
\vspace{0.5em}
\vspace{0.5em}
Labs \textbackslash{}& Livrables Lab A (Messagerie On-chain1 ) : Créer un programme qui permet à des utilisateurs de poster des messages stockés\\
dans des comptes dédiés. Lab B (Mini-Escrow) : Un contrat qui bloque des fonds jusqu’à validation par un tiers. Livrable : Repo GitHub\\
structuré avec tests TS (Mocha/Chai) interagissant avec ‘solana-test-validator‘.\\
\vspace{0.5em}
\begin{center}\begin{tabular}TAB. 9.5 : Checklist Sécurité Module 1\end{tabular}\end{center}
\vspace{0.5em}
\vspace{0.5em}
Contrôle     Vérification                                 Fail Typique\\
Owner Check Vérifier que account\textbackslash{}_info.owner == program\textbackslash{}_id Injection de données\\
Signer Check Vérifier account\textbackslash{}_info.is\textbackslash{}_signer              Usurpation\\
Rent Exempt Le compte est-il assez fondé ?                Compte purgé\\
\vspace{0.5em}
\vspace{0.5em}
\vspace{0.5em}
\vspace{0.5em}
EL\\
TI\\
\subsection{9.2.2 MODULE 2 : Maîtrise du Framework Anchor (Semaines 17-20)}
\vspace{0.5em}
\vspace{0.5em}
\vspace{0.5em}
\vspace{0.5em}
EN\\
Objectifs Opérationnels\\
\vspace{0.5em}
• Utiliser les macros Anchor pour sécuriser le code (\textbackslash{}#[account(...)]).\\
\vspace{0.5em}
\vspace{0.5em}
\vspace{0.5em}
\vspace{0.5em}
D\\
FI\\
• Gérer les PDAs (Program Derived Addresses) de manière déterministe.\\
\vspace{0.5em}
• Émettre des Events pour l’indexation.\\
N\\
CO\\
Labs \textbackslash{}& Livrables Lab A (Counter PDA) : Un compteur global et des compteurs user-specific utilisant des seeds. Lab B (Staking Vault) :\\
Un utilisateur dépose des tokens, le programme tracking le solde et le temps. Livrable : Code Anchor propre, Tests TypeScript étendus,\\
IDL publié.\\
1\\
Web3/Blockchain — Donnée ou logique exécutée et stockée directement sur la blockchain, vérifiée par le consensus du réseau.\\
\vspace{0.5em}
\vspace{0.5em}
\vspace{0.5em}
\vspace{0.5em}
\begin{confidential}CONFIDENTIEL — PROJET RBK 2.0                                            Page 92 sur 316\end{confidential}
RBK 2.0 — Whitepaper Architecte Web3                                                                                 Chapitre 9. Track A : Solana Engineer\\
\vspace{0.5em}
\vspace{0.5em}
Client (TS)          IDL        Program (Rust)\\
\vspace{0.5em}
\vspace{0.5em}
\vspace{0.5em}
Accounts (Data)\\
\vspace{0.5em}
FIG. 9.2 : Flux Anchor\\
\vspace{0.5em}
\vspace{0.5em}
\subsection{9.2.3 MODULE 3 : Architectures Avancées \textbackslash{}& Innovation (Semaines 21-24)}
Objectifs Opérationnels\\
\vspace{0.5em}
\vspace{0.5em}
\vspace{0.5em}
\vspace{0.5em}
EL\\
• CPI (Cross-Program Invocations) : Appeler un programme depuis un autre (ex : Transfert SPL Token).\\
\vspace{0.5em}
\vspace{0.5em}
\vspace{0.5em}
\vspace{0.5em}
TI\\
• Token Extensions (Token-20222 ) : Metadata, Transfer Hooks.\\
\vspace{0.5em}
\vspace{0.5em}
\vspace{0.5em}
\vspace{0.5em}
EN\\
• Architecture Modulaire : Séparer la logique métier du stockage.\\
\vspace{0.5em}
\vspace{0.5em}
\vspace{0.5em}
\vspace{0.5em}
D\\
Labs \textbackslash{}& Livrables Lab A (CPI Challenge) : Un programme ”Master” qui contrôle un programme ”Slave” via CPI signée (PDA Signer).\\
\vspace{0.5em}
\vspace{0.5em}
\vspace{0.5em}
\vspace{0.5em}
FI\\
Livrable : Architecture complexe documentée (C4 Model) et tests d’intégration multi-programmes.\\
\vspace{0.5em}
N\\
CO\\
\subsection{9.2.4 MODULE 4 : Production Hardening \textbackslash{}& UX Performance (Semaines 25-28)}
Objectifs Opérationnels\\
\vspace{0.5em}
• Optimisation des Compute Units (CU) pour réduire les coûts et la latence.\\
\vspace{0.5em}
• Gestion des erreurs custom et logs structurés.\\
\vspace{0.5em}
• Préparation à l’audit (Freeze Authority, Upgradeability).\\
2\\
Web3/Blockchain — SPL Token-2022 : standard Solana ajoutant des extensions (règles, hooks, métadonnées étendues) pour des jetons programmables.\\
\vspace{0.5em}
\vspace{0.5em}
\vspace{0.5em}
\vspace{0.5em}
\begin{confidential}CONFIDENTIEL — PROJET RBK 2.0                                          Page 93 sur 316\end{confidential}
RBK 2.0 — Whitepaper Architecte Web3                                                                                          Chapitre 9. Track A : Solana Engineer\\
\vspace{0.5em}
\vspace{0.5em}
Projet Final Une dApp complète (ex : AMM simplifié ou DAO Voting) déployée sur Devnet, avec une UI fonctionnelle, une suite de tests\\
CI/CD, et un rapport d’auto-audit. La préparation inclut des contrôles Ops3 pour sécuriser la mise en production.\\
\vspace{0.5em}
\begin{center}\begin{tabular}TAB. 9.7 : Production Readiness Review (PRR)\end{tabular}\end{center}
\vspace{0.5em}
\vspace{0.5em}
Domaine Critère           Preuve                          Statut\\
Sécurité Fuzzing Tests    Corps de cas limites testés     Obligatoire\\
Ops      Multisig Upgrade Clés gérées par Squads/Multisig Obligatoire\\
Doc      Architecture     Diagramme Mermaid à jour        Obligatoire\\
\vspace{0.5em}
\vspace{0.5em}
\vspace{0.5em}
\vspace{0.5em}
EL\\
\subsection{9.3 Stack Technique Spécifique}
\vspace{0.5em}
\vspace{0.5em}
\vspace{0.5em}
\vspace{0.5em}
TI\\
La stack Solana évolue vite. Nous imposons une version LTS (Long Term Support) et des outils standards.\\
\vspace{0.5em}
\vspace{0.5em}
\vspace{0.5em}
\vspace{0.5em}
EN\\
\begin{center}\begin{tabular}TAB. 9.9 : Stack Track A (Standard)\end{tabular}\end{center}
\vspace{0.5em}
\vspace{0.5em}
\vspace{0.5em}
\vspace{0.5em}
D\\
FI\\
Catégorie Outils                         Usage\\
\vspace{0.5em}
\vspace{0.5em}
N\\
Core      Rust, Solana CLI, Anchor       Dev, Deploy, Test\\
Client    TypeScript, web3.js, Anchor.ts Intégration Front/Tests\\
CO\\
Security  Trident (Fuzzing), Soteria     Audit auto\\
DevOps    GitHub Actions, Solana Verify CI/CD\\
\vspace{0.5em}
\vspace{0.5em}
\vspace{0.5em}
\vspace{0.5em}
\subsection{9.4 RBK Solana Validator Track}
Le Validator Track prolonge le niveau Bâtisseur vers l’Architecte : les apprenants qui réussissent le Block Check “Build” peuvent intégrer\\
un parcours opérationnel focalisé sur l’exploitation d’un validateur mainnet-beta Solana. L’objectif est double : maîtriser les exigences SRE\\
3\\
Ops/SRE — « Ops » désigne les pratiques et activités d’exploitation (monitoring, incidents, runbooks, disponibilité).\\
\vspace{0.5em}
\vspace{0.5em}
\vspace{0.5em}
\begin{confidential}CONFIDENTIEL — PROJET RBK 2.0                                               Page 94 sur 316\end{confidential}
RBK 2.0 — Whitepaper Architecte Web3                                                                   Chapitre 9. Track A : Solana Engineer\\
\vspace{0.5em}
\vspace{0.5em}
du réseau et préparer l’offre Launchpad Nexus (forfait 3 300 TND + 30\textbackslash{}% de partage sur les protocoles accompagnés).\\
\vspace{0.5em}
Objectifs clés\\
\vspace{0.5em}
• Garantir un uptime \textgreater{} 99,5\textbackslash{}% et une latence de signature stable en conditions de production.\\
\vspace{0.5em}
• Automatiser la rotation des clés vote et staker avec un Runbook reproductible.\\
\vspace{0.5em}
• Exposer des métriques GraphQL/Prometheus et alertes Discord intégrées au cockpit RBK.\\
\vspace{0.5em}
• Déployer un plan de réponse incident (rollback snapshot, vote account recovery) validé par les mentors SRE.\\
\vspace{0.5em}
\vspace{0.5em}
\vspace{0.5em}
\vspace{0.5em}
EL\\
\begin{center}\begin{tabular}TAB. 9.11 : Itinéraire Validator (6 semaines)\end{tabular}\end{center}
\vspace{0.5em}
\vspace{0.5em}
\vspace{0.5em}
\vspace{0.5em}
TI\\
Semaine Focus                                Livrable                                         NFT Skills\\
\vspace{0.5em}
\vspace{0.5em}
\vspace{0.5em}
\vspace{0.5em}
EN\\
1       Provision bare-metal / Cloud hybride Terraform + Script Ansible validés (Stage CI/CD) Genesis Access (validator)\\
2       Performance \textbackslash{}& Observabilité          Dash Grafana + Alertes PagerDuty                 Skill Mirror “Ops” \textbackslash{}#1\\
\vspace{0.5em}
\vspace{0.5em}
\vspace{0.5em}
\vspace{0.5em}
D\\
3       Sécurité Clés \textbackslash{}& Gouvernance          Procédure rotation clés, Multisig Squads         Badge “Guardian Ops”\\
\vspace{0.5em}
\vspace{0.5em}
\vspace{0.5em}
\vspace{0.5em}
FI\\
4       Participation Réseau                 Vote account sur mainnet-beta, Score RPC stable  Skill Mirror “Reliability”\\
5       Monétisation Launchpad\\
N    Modèle de rémunération, prévision ISA/Launchpad NFT “Validator Lead”\\
CO\\
6       Incident Game Day                    Postmortem complet + plan remédiation            Contribution Launch Proof\\
\vspace{0.5em}
\vspace{0.5em}
\vspace{0.5em}
Livrables attendus\\
\vspace{0.5em}
• Validator Runbook : procédures cold/hot restart, gestion snapshots, mapping dépendances réseau.\\
\vspace{0.5em}
• Observability Pack : dashboards partagés avec l’équipe Launchpad et intégrations alerting.\\
\vspace{0.5em}
• NFT Skills Validator Lead : conditionnée à ≥ 2 Skill Mirror “Ops” validés et audit mentor.\\
\vspace{0.5em}
• Dossier Nexus : fiche projet incluant modèle économique (tuition 15 900 TND, forfait 3 300 TND, partage 30\textbackslash{}%).\\
\vspace{0.5em}
\vspace{0.5em}
\begin{confidential}CONFIDENTIEL — PROJET RBK 2.0                                 Page 95 sur 316\end{confidential}
RBK 2.0 — Whitepaper Architecte Web3                                                                      Chapitre 9. Track A : Solana Engineer\\
\vspace{0.5em}
\vspace{0.5em}
Les diplômés du Validator Track deviennent des points focaux pour les cohortes suivantes : ils co-animent les Block Check “Build”,\\
fournissent des audits PRR pour les projets Launchpad et siègent au jury Gate GO/NO-GO. Cette boucle d’excellence garantit la résilience\\
de la communauté RBK et renforce la proposition de valeur auprès des partenaires institutionnels.\\
\vspace{0.5em}
\vspace{0.5em}
\subsection{9.5 Profil de Sortie : Le « Guardian »}
Le Guardian est un ingénieur rare. Il ne ”bricole” pas des scripts. Il construit des infrastructures financières immuables. Il est capable\\
de livrer un protocole DeFi sécurisé, documenté et performant en autonomie. Son employabilité est maximale car il maîtrise la chaîne de\\
valeur complète : du bas niveau (Rust/BPF) au haut niveau (Architecture/Produit).\\
\vspace{0.5em}
\vspace{0.5em}
\vspace{0.5em}
\vspace{0.5em}
EL\\
Missions Types en Entreprise\\
\vspace{0.5em}
\vspace{0.5em}
\vspace{0.5em}
\vspace{0.5em}
TI\\
• Construire un DEX (Decentralized Exchange) à haute fréquence.\\
\vspace{0.5em}
\vspace{0.5em}
\vspace{0.5em}
\vspace{0.5em}
EN\\
• Auditer un protocole de Lending pour détecter les failles de liquidité.\\
\vspace{0.5em}
\vspace{0.5em}
\vspace{0.5em}
\vspace{0.5em}
D\\
• Optimiser les coûts de gas d’un programme NFT à fort volume (Compression).\\
\vspace{0.5em}
\vspace{0.5em}
\vspace{0.5em}
\vspace{0.5em}
FI\\
\begin{center}\begin{tabular}TAB. 9.13 : Checklist Portfolio Guardian\end{tabular}\end{center}
N\\
CO\\
Artefact       Contenu                              Critère\\
3 Repos GitHub Code Rust clean, Tests, CI           Green CI Badge\\
Audit Report   Analyse d’un protocole tiers         3 failles identifiées\\
Demo Live      Vidéo Loom (3 min) expliquant l’arch Clarté orale\\
\vspace{0.5em}
\vspace{0.5em}
\vspace{0.5em}
\vspace{0.5em}
\begin{confidential}CONFIDENTIEL — PROJET RBK 2.0                                  Page 96 sur 316\end{confidential}
\section{TRACK B : EVM ENGINEER (SOLIDITY/FOUNDRY)}
\vspace{0.5em}
\vspace{0.5em}
\vspace{0.5em}
\vspace{0.5em}
EL\\
\subsection{10.1 Philosophie du Track : La Maîtrise du Standard Industriel}
L’EVM1 (Ethereum Virtual Machine) est le standard mondial des smart contracts. Maîtriser Solidity et Foundry, c’est s’ouvrir les portes\\
\vspace{0.5em}
\vspace{0.5em}
\vspace{0.5em}
\vspace{0.5em}
TI\\
de l’écosystème le plus riche (Ethereum, Arbitrum, Optimism, Base, Polygon). Notre approche est ”Infrastructure-First”. Nous ne formons\\
\vspace{0.5em}
\vspace{0.5em}
\vspace{0.5em}
\vspace{0.5em}
EN\\
pas des développeurs qui copient-collent du code OpenZeppelin, mais des ingénieurs capables de comprendre le stockage bas niveau,\\
l’optimisation du Gas et les subtilités des upgrades (Proxies).\\
\vspace{0.5em}
\vspace{0.5em}
\vspace{0.5em}
\vspace{0.5em}
D\\
FI\\
Positionnement ”Infra Engineer” L’ingénieur EVM RBK est un bâtisseur de protocoles. Il maîtrise la chaîne DevOps (Foundry, CI,\\
\vspace{0.5em}
N\\
Verification), la sécurité offensive (Fuzzing, Invariants) et les patterns de composabilité (DeFi Lego).\\
CO\\
NON NÉGOCIABLE : AUDIT-READINESS\\
• Test-First : TDD strict avec Foundry. Fuzzing obligatoire.\\
\vspace{0.5em}
• Gas Optimization : Chaque Opcode compte (Assembly si nécessaire).\\
\vspace{0.5em}
• Security Mindset : ”Don’t trust, verify”. Protection Reentrancy, Overflow, Access Control.\\
\vspace{0.5em}
1\\
Web3/Blockchain — Ethereum Virtual Machine (EVM) : moteur d’exécution des smart contracts bytecode compatibles sur Ethereum et chaînes EVM, assurant\\
portabilité des dApps et outils.\\
\vspace{0.5em}
\vspace{0.5em}
\vspace{0.5em}
\vspace{0.5em}
97\\
RBK 2.0 — Whitepaper Architecte Web3                                                                Chapitre 10. Track B : EVM Engineer\\
\vspace{0.5em}
\vspace{0.5em}
Solidity         Tests (Foundry)        Security (Slither)\\
\vspace{0.5em}
\vspace{0.5em}
\vspace{0.5em}
Deploy (Scripts)\\
\vspace{0.5em}
FIG. 10.1 : Chaîne de Valeur EVM\\
\vspace{0.5em}
\vspace{0.5em}
\vspace{0.5em}
\subsection{10.2 Structure Pédagogique : De la Logique au Durcissement (16 Semaines)}
Un parcours intensif qui commence par les fondations (Storage Layout) pour aller jusqu’au déploiement multi-chain complexe.\\
\vspace{0.5em}
\vspace{0.5em}
\vspace{0.5em}
\vspace{0.5em}
EL\\
\begin{center}\begin{tabular}TAB. 10.1 : Carte des Modules Track B\end{tabular}\end{center}
\vspace{0.5em}
\vspace{0.5em}
\vspace{0.5em}
\vspace{0.5em}
TI\\
Module       Sem Objectif          Livrable\\
\vspace{0.5em}
\vspace{0.5em}
\vspace{0.5em}
\vspace{0.5em}
EN\\
1. Basics    13-15 EVM Internals   Vault Natif\\
2. Pro Env   16-18 Foundry Mastery CI Pipeline\\
\vspace{0.5em}
\vspace{0.5em}
\vspace{0.5em}
\vspace{0.5em}
D\\
3. Standards 19-21 ERC20/721       Token System\\
\vspace{0.5em}
\vspace{0.5em}
\vspace{0.5em}
\vspace{0.5em}
FI\\
4. dApp      22-24 Intégration     Full Stack dApp\\
\vspace{0.5em}
N\\
5. Scaling   25-26 L2 \textbackslash{}& Upgrades   UUPS Proxy\\
CO\\
6. Hardening 27-28 Sécurité        Audit Report\\
\vspace{0.5em}
\vspace{0.5em}
\vspace{0.5em}
\vspace{0.5em}
\subsection{10.2.1 MODULE 1 : Smart Contract Basics \textbackslash{}& Solidity Deep Dive (Semaines 13-15)}
Objectifs Comprendre comment l’EVM stocke les données (Slots), la différence Memory/CallData, et les structures de contrôle de base.\\
Lab A (Vault Sécurisé) : Un contrat de dépôt/retrait avec gestion des rôles (Ownable). Critères : Tests de cas nominaux et d’erreurs\\
(Revert).\\
\vspace{0.5em}
\vspace{0.5em}
\vspace{0.5em}
\vspace{0.5em}
\begin{confidential}CONFIDENTIEL — PROJET RBK 2.0                                Page 98 sur 316\end{confidential}
RBK 2.0 — Whitepaper Architecte Web3                                                                                      Chapitre 10. Track B : EVM Engineer\\
\vspace{0.5em}
\vspace{0.5em}
\subsection{10.2.2 MODULE 2 : Environnement de Développement Pro (Semaines 16-18)}
\vspace{0.5em}
Objectifs Passer de Remix à Foundry. Maîtriser ‘forge test‘, ‘cast‘, et le Fuzzing. Lab (Test Suite) : Écrire une suite de tests exhaustive\\
(Unit + Fuzz) pour un protocole existant (ex : Uniswap V2 Pair simplifié). Livrable : Repo avec GitHub Actions qui lance les tests à chaque\\
PR.\\
\vspace{0.5em}
\vspace{0.5em}
\subsection{10.2.3 MODULE 3 : Token Standards \textbackslash{}& Composabilité (Semaines 19-21)}
Objectifs Implémenter ERC20, ERC721, ERC1155. Comprendre ‘approve‘, ‘transferFrom‘ et les risques associés. Lab (DeFi Lego) : Un\\
contrat qui ”wrap” un token pour ajouter du rendement (Staking2 ). Critères : Interopérabilité vérifiée avec les standards.\\
\vspace{0.5em}
\vspace{0.5em}
\vspace{0.5em}
\vspace{0.5em}
EL\\
\subsection{10.2.4 MODULE 4 : dApp Development \textbackslash{}& Web3 Integration (Semaines 22-24)}
\vspace{0.5em}
\vspace{0.5em}
\vspace{0.5em}
\vspace{0.5em}
TI\\
Objectifs Connecter un Front (React/Next) au contrat via Wagmi/Viem. Gérer le cycle de vie de la transaction (Pending, Confirmed,\\
\vspace{0.5em}
\vspace{0.5em}
\vspace{0.5em}
\vspace{0.5em}
EN\\
Failed). Lab (Mini-DEX UI) : Interface pour swapper des tokens (simulation). Checklist : Gestion des erreurs RPC, Feedback utilisateur.\\
\vspace{0.5em}
\vspace{0.5em}
\vspace{0.5em}
\vspace{0.5em}
D\\
FI\\
\subsection{10.2.5 MODULE 5 : L2 Scaling \textbackslash{}& Advanced Patterns (Semaines 25-26)}
Objectifs Comprendre les Rollups (Optimistic/ZK). Déployer sur Arbitrum/Base. Gérer l’upgradeabilité (Proxies). Lab (UUPS Upgrade) :\\
N\\
Déployer une V1, puis upgrader vers une V2 sans perdre l’état (Storage).\\
CO\\
\subsection{10.2.6 MODULE 6 : Production Hardening \textbackslash{}& Security (Semaines 27-28)}
Objectifs Sécurisation finale. Audit interne. Projet Final : Déploiement d’un protocole complet sur Testnet (Sepolia/Goerli) avec scripts\\
de vérification Etherscan automatisés.\\
2\\
Web3/Blockchain — Staking : blocage de tokens pour sécuriser un réseau Proof-of-Stake et recevoir des récompenses en échange de la participation au consensus.\\
\vspace{0.5em}
\vspace{0.5em}
\vspace{0.5em}
\vspace{0.5em}
\begin{confidential}CONFIDENTIEL — PROJET RBK 2.0                                             Page 99 sur 316\end{confidential}
RBK 2.0 — Whitepaper Architecte Web3                                                                   Chapitre 10. Track B : EVM Engineer\\
\vspace{0.5em}
\vspace{0.5em}
\begin{center}\begin{tabular}TAB. 10.3 : Security Checklist EVM\end{tabular}\end{center}
\vspace{0.5em}
\vspace{0.5em}
Vulnérabilité Contrôle                          Outil\\
Reentrancy     Checks-Effects-Interactions      Slither\\
Access Control Modifiers corrects               Manual Review\\
Arithmetic     Overflow (Solidity \textless{}0.8 checked) Fuzzing\\
\vspace{0.5em}
\vspace{0.5em}
\vspace{0.5em}
\vspace{0.5em}
\subsection{10.3 Stack Technique Spécifique}
\vspace{0.5em}
\vspace{0.5em}
\vspace{0.5em}
\vspace{0.5em}
EL\\
On privilégie la stack moderne (Rust-based) pour sa rapidité.\\
\vspace{0.5em}
\vspace{0.5em}
\vspace{0.5em}
\vspace{0.5em}
TI\\
\begin{center}\begin{tabular}TAB. 10.5 : Stack Track B (Foundry)\end{tabular}\end{center}
\vspace{0.5em}
\vspace{0.5em}
\vspace{0.5em}
\vspace{0.5em}
EN\\
Catégorie Outils                       Usage\\
\vspace{0.5em}
\vspace{0.5em}
\vspace{0.5em}
\vspace{0.5em}
D\\
Framework Foundry (Forge, Cast, Anvil) Tout-en-un\\
\vspace{0.5em}
\vspace{0.5em}
\vspace{0.5em}
\vspace{0.5em}
FI\\
Libs      OpenZeppelin Contracts       Standards sécu\\
Client    Viem, Wagmi                  Front-end\\
\vspace{0.5em}
N\\
Analysis  Slither, Aderyn              Static Analysis\\
CO\\
\vspace{0.5em}
\subsection{10.4 Profil de Sortie : L’Ingénieur d’Infrastructure EVM}
L’Ingénieur EVM RBK est prêt pour intégrer une équipe Core Protocol ou une start-up DeFi. Il sait écrire du code qui gère des millions\\
de dollars.\\
\vspace{0.5em}
Missions Types en Entreprise\\
\vspace{0.5em}
• Développer une stratégie de Loop Staking sur un protocole de Lending (ex : Aave).\\
\vspace{0.5em}
\vspace{0.5em}
\begin{confidential}CONFIDENTIEL — PROJET RBK 2.0                                Page 100 sur 316\end{confidential}
RBK 2.0 — Whitepaper Architecte Web3                                                                                       Chapitre 10. Track B : EVM Engineer\\
\vspace{0.5em}
\vspace{0.5em}
• Migrer un token Gouvernance d’un Layer 1 vers un Layer 2 (Bridge3 Architecture).\\
\vspace{0.5em}
• Écrire les scripts de déploiement automatisés pour une collection NFT de 10k pièces avec whitelist Merkle Tree.\\
\vspace{0.5em}
La trajectoire comprend un volet Ops pour fiabiliser les déploiements et vérifier les contrats.\\
\vspace{0.5em}
\begin{center}\begin{tabular}TAB. 10.7 : Matrice Compétences Infra EVM\end{tabular}\end{center}
\vspace{0.5em}
\vspace{0.5em}
Domaine Attendu                       Preuve\\
Solidity Expert (Storage, Assembly) Gas Golfing Repo\\
Testing  Expert (Fuzzing, Invariants) Coverage \textgreater{} 95\textbackslash{}%\\
\vspace{0.5em}
\vspace{0.5em}
\vspace{0.5em}
\vspace{0.5em}
EL\\
Ops      Autonome (Scripts, Verify) Déploiements vérifiés\\
\vspace{0.5em}
\vspace{0.5em}
\vspace{0.5em}
\vspace{0.5em}
TI\\
D     EN\\
FI\\
N\\
CO\\
\vspace{0.5em}
\vspace{0.5em}
\vspace{0.5em}
\vspace{0.5em}
3\\
Web3/Blockchain — Mécanisme de transfert ou de verrouillage d’actifs entre blockchains distinctes, basé sur des preuves ou des validateurs fédérés.\\
\vspace{0.5em}
\vspace{0.5em}
\begin{confidential}CONFIDENTIEL — PROJET RBK 2.0                                            Page 101 sur 316\end{confidential}
\section{TRACK C : PRODUCT \textbackslash{}& GROWTH ENGINEER (FULL}
STACK + ANALYTICS)\\
\vspace{0.5em}
\vspace{0.5em}
\vspace{0.5em}
\vspace{0.5em}
EL\\
TI\\
\subsection{11.1 Philosophie du Track : Le ”Product Builder” Complet}
\vspace{0.5em}
\vspace{0.5em}
\vspace{0.5em}
\vspace{0.5em}
EN\\
Le monde Web3 regorge de protocoles techniquement brillants mais inutilisables. Le Track C forme le chaînon manquant : l’ingénieur\\
capable de construire une dApp de bout en bout (Front, Indexing, Analytics) et de piloter sa croissance. Ce n’est pas un track ”No-Code”.\\
\vspace{0.5em}
\vspace{0.5em}
\vspace{0.5em}
\vspace{0.5em}
D\\
C’est un track ”Full-Code” orienté utilisateur. L’étudiant apprend à connecter les smart contracts au monde réel, à visualiser la data\\
On-chain, et à itérer sur le produit en fonction des métriques.\\
\vspace{0.5em}
\vspace{0.5em}
FI\\
N\\
Positionnement ”Product Engineer”\\
Il maîtrise la stack Next.js/React, l’indexation (Subgraphs/Squid), et l’analyse de données (Dune/SQL). Il sait que le code n’est qu’un\\
O\\
\vspace{0.5em}
moyen de livrer de la valeur.\\
C\\
\vspace{0.5em}
\vspace{0.5em}
NON NÉGOCIABLE : USER-CENTRICITY\\
• Zero-Friction : Gérer les erreurs RPC, les wallets et l’UX dégradée sans perdre l’utilisateur.\\
\vspace{0.5em}
• Data-Driven : Pas de feature sans tracking. Dashboard analytique jour 1.\\
\vspace{0.5em}
• Shipping : Déploiement continu (Vercel/Fleek) et tests E2E (Playwright).\\
\vspace{0.5em}
\vspace{0.5em}
\vspace{0.5em}
102\\
RBK 2.0 — Whitepaper Architecte Web3                                                                Chapitre 11. Track C : Product Engineer\\
\vspace{0.5em}
\vspace{0.5em}
\vspace{0.5em}
\subsection{11.2 Structure Pédagogique : Fondations communes, spécialisations binômes (16}
Semaines)\\
Le Track C se compose de six semaines de fondations techniques partagées, puis de deux parcours complémentaires : Web3 Product\\
Manager et Web3 Full-Stack Engineer. Les apprenants avancent en binômes mixtes, garantissant un transfert croisé des compétences\\
avant le Capstone final.\\
\vspace{0.5em}
\vspace{0.5em}
\subsection{11.2.1 Bloc Fondations (Semaines 13-18)}
\vspace{0.5em}
Fondations communes Track C\\
\vspace{0.5em}
\vspace{0.5em}
\vspace{0.5em}
\vspace{0.5em}
EL\\
TI\\
Semaines Module                     Objectifs clefs                                Livrable\\
13-15     Web3 Connectivity Wallet adapters (RainbowKit/Solana), gestion Universal   Profile              Viewer      multi-chain\\
\vspace{0.5em}
\vspace{0.5em}
\vspace{0.5em}
\vspace{0.5em}
EN\\
\textbackslash{}& State           d’état asynchrone (TanStack Query), résilience (EVM+SVM)\\
Management        UX\\
\vspace{0.5em}
\vspace{0.5em}
\vspace{0.5em}
\vspace{0.5em}
D\\
16-17     Indexing \textbackslash{}& Data         Pipelines The Graph/Goldsky, sync Postgres, stra- Subgraph performant (\textless{}100ms) sur un DEX exis-\\
\vspace{0.5em}
\vspace{0.5em}
\vspace{0.5em}
\vspace{0.5em}
FI\\
Layer                   tégies reorg                                      tant\\
18      Analytics Temps\\
N\\
Dune/Flipside, SQL analytique, instrumentation Tableau ”Whale Watcher” avec alerting Slack\\
CO\\
Réel                    produit\\
\vspace{0.5em}
\vspace{0.5em}
\vspace{0.5em}
\vspace{0.5em}
\begin{confidential}CONFIDENTIEL — PROJET RBK 2.0                                  Page 103 sur 316\end{confidential}
RBK 2.0 — Whitepaper Architecte Web3                                                                Chapitre 11. Track C : Product Engineer\\
\vspace{0.5em}
\vspace{0.5em}
\subsection{11.2.2 Parcours de spécialisation (Semaines 19-28)}
\vspace{0.5em}
Choix de spécialisation\\
\vspace{0.5em}
\vspace{0.5em}
Axe                                     Modules \textbackslash{}& objectifs                           Livrables contractuels\\
Option A — Web3 Product Manager          • Tokenomics \textbackslash{}& Incentive Design : modélisa- • Playbook GTM complet (Pricing, messa-\\
tion cash-flow, gouvernance.                ging, funnels).\\
\vspace{0.5em}
• Growth Engineering : campagnes referral • Smart sheet ”Growth Experiments” avec\\
on-chain, segmentation CRM.               KPIs.\\
\vspace{0.5em}
\vspace{0.5em}
\vspace{0.5em}
\vspace{0.5em}
EL\\
• GTM \textbackslash{}& Compliance : pricing, réglementa- • Dossier Contrat de Partenariat Pédago-\\
\vspace{0.5em}
\vspace{0.5em}
\vspace{0.5em}
\vspace{0.5em}
TI\\
tion marketplaces, storyboards UX.        gique aligné produit.\\
Option B — Web3 Full-Stack Engineer • Automation \textbackslash{}& Bots : Chainlink Automa- • Bot opérationnel (sniping/monitoring) au-\\
\vspace{0.5em}
\vspace{0.5em}
\vspace{0.5em}
\vspace{0.5em}
EN\\
tion, bots Discord/TG.                    dité sécurité.\\
\vspace{0.5em}
\vspace{0.5em}
\vspace{0.5em}
\vspace{0.5em}
D\\
• Observabilité \textbackslash{}& Sécurité : OpenTelemetry, • Tableau SRE avec budgets d’erreur.\\
\vspace{0.5em}
\vspace{0.5em}
\vspace{0.5em}
\vspace{0.5em}
FI\\
SLO, alertes Threat Intel.\\
• Infra codée (Terraform/Fly.io) + runbook\\
\vspace{0.5em}
N\\
• Production \textbackslash{}& DevOps : pipelines CI/CD       incident.\\
CO\\
(GitHub Actions), IaC pour déploiements\\
Web3.\\
\vspace{0.5em}
\vspace{0.5em}
Travail en binôme et Capstone conjoints Chaque binôme associe un Product Manager et un Full-Stack. Ils co-conçoivent un SaaS Web3 :\\
la composante Produit définit la roadmap, la métrologie et les boucles de croissance ; la composante Full-Stack implémente l’architecture\\
technique, l’automatisation et les opérations. Le Capstone final est évalué sur la cohérence GTM, la robustesse technique et la capacité à\\
itérer sur des métriques réelles.\\
\vspace{0.5em}
\vspace{0.5em}
\subsection{11.3 Profil de Sortie}
\vspace{0.5em}
\begin{confidential}CONFIDENTIEL — PROJET RBK 2.0                               Page 104 sur 316\end{confidential}
RBK 2.0 — Whitepaper Architecte Web3                                                                 Chapitre 11. Track C : Product Engineer\\
\vspace{0.5em}
\vspace{0.5em}
\vspace{0.5em}
Un ”Full Stack Web3” capable de lancer sa propre startup ou de rejoindre une équipe Growth dans un grand protocole.\\
\vspace{0.5em}
\begin{center}\begin{tabular}TAB. 11.3 : Matrice Compétences Product\end{tabular}\end{center}
\vspace{0.5em}
\vspace{0.5em}
Domaine Attendu                          Preuve\\
Front-end Expert React/Next.js           App fluide 60fps\\
Data      Capable d’écrire des subgraphs API GraphQL\\
Growth    Comprend les funnels           Analytics intégrés\\
\vspace{0.5em}
\vspace{0.5em}
\vspace{0.5em}
\vspace{0.5em}
EL\\
TI\\
D   EN\\
FI\\
N\\
CO\\
\vspace{0.5em}
\vspace{0.5em}
\vspace{0.5em}
\vspace{0.5em}
\begin{confidential}CONFIDENTIEL — PROJET RBK 2.0                               Page 105 sur 316\end{confidential}
\section{MODULE SOFT SKILLS \textbackslash{}& PROFESSIONNALISATION}
\vspace{0.5em}
\vspace{0.5em}
\vspace{0.5em}
\vspace{0.5em}
EL\\
\subsection{12.1 Structure du Module (4 semaines)}
Ce module de 4 semaines (Phase 3) est le pont critique entre l’étudiant et le professionnel. Dans un marché où la compétence technique\\
\vspace{0.5em}
\vspace{0.5em}
\vspace{0.5em}
\vspace{0.5em}
TI\\
est un pré-requis, c’est la ”Seniorité Attitude” qui déclenche l’embauche. Nous ne formons pas seulement des codeurs, mais des ingénieurs\\
\vspace{0.5em}
\vspace{0.5em}
\vspace{0.5em}
\vspace{0.5em}
EN\\
capables de gérer des incidents, de communiquer avec des stakeholders non-techniques, et de vendre leur valeur. C’est l’étape finale de\\
transformation ”Senior-by-Design”.\\
\vspace{0.5em}
\vspace{0.5em}
\vspace{0.5em}
\vspace{0.5em}
D\\
FI\\
Livrables Finaux du Module (Obligatoires pour Certification)        Pour valider cette phase, l’étudiant doit produire et faire valider :\\
\vspace{0.5em}
N\\
1. Rapport d’Audit Professionnel : Basé sur le template Code4rena/OtterSec, analysant un protocole réel.\\
CO\\
2. Documentation Technique (GitBook) : Une documentation utilisateur et développeur complète pour leur Capstone.\\
\vspace{0.5em}
3. Package Freelance : Une proposition commerciale (SOW) type, une grille tarifaire (TJM) et un contrat de service.\\
\vspace{0.5em}
4. Board Projet (Jira/Notion) : L’historique des sprints, user stories, et une rétrospective écrite post-mortem.\\
\vspace{0.5em}
5. Pitch Deck (10 slides) \textbackslash{}& Démo : Une présentation vidéo (3-5 min) et un pitch deck investisseur.\\
\vspace{0.5em}
6. Profil Public : GitHub (Green dots, Readme profil) et LinkedIn (Headline, About, Featured) optimisés.\\
\vspace{0.5em}
\vspace{0.5em}
\vspace{0.5em}
\vspace{0.5em}
106\\
RBK 2.0 — Whitepaper Architecte Web3                                           Chapitre 12. MODULE SOFT SKILLS \textbackslash{}& PROFESSIONNALISATION\\
\vspace{0.5em}
\vspace{0.5em}
\vspace{0.5em}
\vspace{0.5em}
S25 : Comms                  S26 : Business            S27 : Gestion                S28 : Leadership             DEMO DAY\\
FIG. 12.1 : Timeline 4 semaines — Soft Skills \textbackslash{}& Pro\\
\vspace{0.5em}
\vspace{0.5em}
\begin{center}\begin{tabular}TAB. 12.1 : Vue d’ensemble du module (4 semaines)\end{tabular}\end{center}
\vspace{0.5em}
\vspace{0.5em}
Sem. Thème              Livrable Principal          Évaluation\\
S25 Communication Tech Audit Report \textbackslash{}& Documentation Revue par Pairs + Mentor\\
S26 Négociation \textbackslash{}& Biz   Simulation Freelance (SOW)  Roleplay Client/Vendeur\\
\vspace{0.5em}
\vspace{0.5em}
\vspace{0.5em}
\vspace{0.5em}
EL\\
S27 Gestion Projet Web3 Board Notion \textbackslash{}& Retro        Audit de Process\\
\vspace{0.5em}
\vspace{0.5em}
\vspace{0.5em}
\vspace{0.5em}
TI\\
S28 Leadership          Pitch Deck \textbackslash{}& Démo           Jury Final (Investisseurs)\\
\vspace{0.5em}
\vspace{0.5em}
\vspace{0.5em}
\vspace{0.5em}
EN\\
Détail Semaine 25 : Communication Technique Objectifs : Savoir vulgariser sans simplifier à l’excès. Rédiger pour être lu. Ateliers :\\
\vspace{0.5em}
\vspace{0.5em}
\vspace{0.5em}
\vspace{0.5em}
D\\
”Writing for Developers” (Docs), ”Audit Reporting Standards”. Exercice : Réécrire le README d’un projet open-source complexe pour le\\
\vspace{0.5em}
\vspace{0.5em}
\vspace{0.5em}
\vspace{0.5em}
FI\\
rendre accessible. Livrable : Rapport d’incident (Post-Mortem) fictif sur un hack historique. Critères : Clarté, Précision technique, Ton\\
professionnel, Anglais technique impeccable.\\
N\\
CO\\
Détail Semaine 26 : Business \textbackslash{}& Négociation Objectifs : Se vendre, chiffrer, contractualiser. Ateliers : ”Pricing your TJM”, ”Mock\\
Négociation Client”, ”Structuring a DAO Proposal”. Exercice : Répondre à un appel d’offre réel (Upwork/Bounties) ou simulé. Livrable :\\
Proposition Commerciale (Statement of Work) complète. Critères : Réalisme du chiffrage, couverture des risques (clauses), force de\\
conviction.\\
\vspace{0.5em}
Détail Semaine 27 : Gestion de Projet Agile/Web3 Objectifs : Délivrer de la valeur en continu, gérer le chaos. Ateliers : ”Scrum for\\
Web3”, ”Async Communication Rules”, ”Github Flow”. Exercice : Organiser le sprint final du Capstone. Livrable : Board Projet propre +\\
Rétrospective Sincère (Start/Stop/Continue). Critères : Transparence, granularité des tickets, gestion des bloquants.\\
\vspace{0.5em}
\vspace{0.5em}
\vspace{0.5em}
\begin{confidential}CONFIDENTIEL — PROJET RBK 2.0                               Page 107 sur 316\end{confidential}
RBK 2.0 — Whitepaper Architecte Web3                                             Chapitre 12. MODULE SOFT SKILLS \textbackslash{}& PROFESSIONNALISATION\\
\vspace{0.5em}
\vspace{0.5em}
Détail Semaine 28 : Leadership \textbackslash{}& Pitch Objectifs : Inspirer la confiance, présenter une vision. Ateliers : ”Public Speaking”, ”Pitch\\
Deck Design”, ”Demo Day Rehearsal”. Exercice : Crash-test du pitch devant des ”commis d’office” hostiles. Livrable : Pitch Deck Final +\\
Vidéo Démo. Critères : Storytelling, Body Language, Gestion du Q\textbackslash{}&A, Qualité visuelle.\\
\vspace{0.5em}
\vspace{0.5em}
\subsection{12.2 Rubrique d’Évaluation}
L’évaluation des Soft Skills chez RBK n’est pas une ”note de participation”. C’est une évaluation professionnelle basée sur des preuves\\
tangibles (artefacts). Nous utilisons une grille stricte pour objectiver la progression. Le barème est conçu pour protéger l’étudiant : on ne\\
juge pas la personnalité, mais les comportements professionnels et les livrables.\\
\vspace{0.5em}
\vspace{0.5em}
\vspace{0.5em}
\vspace{0.5em}
EL\\
Axes d’Évaluation et Pondération\\
\vspace{0.5em}
\vspace{0.5em}
\vspace{0.5em}
\vspace{0.5em}
TI\\
• Communication Technique (30\textbackslash{}%) : Capacité à transmettre de l’information complexe (écrit/oral).\\
\vspace{0.5em}
\vspace{0.5em}
\vspace{0.5em}
\vspace{0.5em}
EN\\
• Collaboration \textbackslash{}& Leadership (30\textbackslash{}%) : Capacité à travailler en équipe, gérer les conflits et driver le projet.\\
\vspace{0.5em}
\vspace{0.5em}
\vspace{0.5em}
\vspace{0.5em}
D\\
• Professionnalisme (40\textbackslash{}%) : Fiabilité, ponctualité, rigueur, gestion du temps, ”Doer” attitude.\\
\vspace{0.5em}
\vspace{0.5em}
\vspace{0.5em}
\vspace{0.5em}
FI\\
Échelle de Notation\\
N\\
CO\\
• Insuffisant (0-9) : Bloquant pour l’emploi. Attitude passive ou toxique. Livrables bâclés.\\
\vspace{0.5em}
• En Progrès (10-13) : Junior standard. Fait le job mais nécessite un management serré.\\
\vspace{0.5em}
• Pro (14-17) : L’objectif RBK. Autonome, fiable, communique proactivement. ”Fire and Forget”.\\
\vspace{0.5em}
• Excellent (18-20) : Top Gun. Tire l’équipe vers le haut, anticipe les problèmes, livre au-delà des attentes.\\
\vspace{0.5em}
\vspace{0.5em}
\vspace{0.5em}
\vspace{0.5em}
\begin{confidential}CONFIDENTIEL — PROJET RBK 2.0                                 Page 108 sur 316\end{confidential}
RBK 2.0 — Whitepaper Architecte Web3                                                  Chapitre 12. MODULE SOFT SKILLS \textbackslash{}& PROFESSIONNALISATION\\
\vspace{0.5em}
\vspace{0.5em}
\begin{center}\begin{tabular}TAB. 12.3 : Rubrique d’Évaluation des Soft Skills\end{tabular}\end{center}
\vspace{0.5em}
\vspace{0.5em}
Axe        Poids Preuves de niveau ”Pro” (14-17)                                       Preuves attendues (Artefacts)\\
Comm. Tech 30\textbackslash{}% Documentation claire, PR descriptions détaillées, sait expliquer le     GitBook du Capstone, Historique de PRs, Rapport d’Audit.\\
”pourquoi” technique.\\
Collab.    30\textbackslash{}% Débloque les autres, ne blâme pas, feedback constructif, utilise les    Commentaires Code Review, Activity Log Discord/Jira.\\
outils async correctement.\\
Pro.       40\textbackslash{}% Respect absolu des deadlines, communication immédiate en cas de         Ponctualité rendus, Qualité finition (typos, UX), Suivi planning.\\
retard, proactivité sur les problèmes.\\
\vspace{0.5em}
\vspace{0.5em}
\vspace{0.5em}
\vspace{0.5em}
EL\\
TI\\
D    EN\\
FI\\
N\\
CO\\
\vspace{0.5em}
\vspace{0.5em}
\vspace{0.5em}
\vspace{0.5em}
\begin{confidential}CONFIDENTIEL — PROJET RBK 2.0                                      Page 109 sur 316\end{confidential}
\section{CAPSTONES (PROJETS SIGNATURES)}
\vspace{0.5em}
\vspace{0.5em}
\vspace{0.5em}
\vspace{0.5em}
EL\\
\subsection{13.1 Philosophie du Capstone : Le Standard « Studio »}
Chez RBK, un Capstone n’est pas un projet d’école. C’est un produit ”Mainnet-Ready” qui respecte les standards d’un studio de déve-\\
\vspace{0.5em}
\vspace{0.5em}
\vspace{0.5em}
\vspace{0.5em}
TI\\
loppement professionnel. Il ne s’agit pas de prouver que ”ça marche”, mais que ”ça ne peut pas casser”.\\
\vspace{0.5em}
\vspace{0.5em}
\vspace{0.5em}
\vspace{0.5em}
EN\\
\begin{center}\begin{tabular}TAB. 13.1 : Studio-Grade Checklist (Non-Négociable)\end{tabular}\end{center}
\vspace{0.5em}
\vspace{0.5em}
\vspace{0.5em}
\vspace{0.5em}
D\\
Catégorie Exigence                                      Preuve\\
\vspace{0.5em}
\vspace{0.5em}
\vspace{0.5em}
\vspace{0.5em}
FI\\
Sécurité  Threat Model formalisé AVANT le code          Doc STRIDE\\
\vspace{0.5em}
\vspace{0.5em}
N\\
Qualité   Zéro warning linter, coverage \textgreater{} 80\textbackslash{}%           Rapport CI\\
Ops       Déploiement scripté et reproductible          Makefile / Taskfile\\
CO\\
\vspace{0.5em}
\subsection{13.2 Les 3 Projets Signatures (Cahier des Charges Industriel)}
Chaque Capstone simule un livrable client. Le niveau d’exigence est ”Audit Ready”.\\
\vspace{0.5em}
\vspace{0.5em}
\subsection{13.2.1 Capstone 1 — Wallet \textbackslash{}& Transaction Reliability Pack (Frontend/UX)}
Focus : Expérience Utilisateur \textbackslash{}& Fiabilité RPC.\\
\vspace{0.5em}
\vspace{0.5em}
\vspace{0.5em}
\vspace{0.5em}
110\\
RBK 2.0 — Whitepaper Architecte Web3                                                     Chapitre 13. CAPSTONES (PROJETS SIGNATURES)\\
\vspace{0.5em}
\vspace{0.5em}
\subsection{13.2.1.1 Spécifications Fonctionnelles (User Stories)}
\vspace{0.5em}
1. US.1 (Feedback) : ”En tant qu’utilisateur, je veux voir un feedback précis (Spinner, Toast) pour chaque état de transaction (Signing,\\
Sending, Confirming, Finalized).”\\
\vspace{0.5em}
2. US.2 (Retry) : ”En tant qu’utilisateur, si le RPC échoue, je veux que l’app réessaie automatiquement (Backoff exponentiel) sans que\\
je reclique.”\\
\vspace{0.5em}
3. US.3 (history) : ”En tant que dev, je veux un log persistant des transactions échouées pour debugger.”\\
\vspace{0.5em}
\subsection{13.2.1.2 Spécifications Techniques}
\vspace{0.5em}
\vspace{0.5em}
\vspace{0.5em}
\vspace{0.5em}
EL\\
• Stack : React (Next.js), TanStack Query, Viem (EVM) ou Solana Web3.js.\\
\vspace{0.5em}
\vspace{0.5em}
\vspace{0.5em}
\vspace{0.5em}
TI\\
• State Machine : Implémentation stricte (Idle → Signing → Broadcasting → Confirming).\\
\vspace{0.5em}
\vspace{0.5em}
\vspace{0.5em}
\vspace{0.5em}
EN\\
• Failover : Configuration d’au moins 2 RPC (Primaire + Fallback).\\
\vspace{0.5em}
\vspace{0.5em}
\vspace{0.5em}
\vspace{0.5em}
D\\
\subsection{13.2.1.3 Checklist Sécurité \textbackslash{}& Qualité}
\vspace{0.5em}
\vspace{0.5em}
\vspace{0.5em}
\vspace{0.5em}
FI\\
Pas de Private Key stockée (Wallet Adapter Only).\\
N\\
CO\\
Sanitize des inputs utilisateurs (Adresses, Montants).\\
\vspace{0.5em}
Gestion du ”Slippage” (Tolérance affichée).\\
\vspace{0.5em}
\subsection{13.2.1.4 Critères d’Acceptation \textbackslash{}& CI}
\vspace{0.5em}
• Test : Simulation de coupure réseau (Mock RPC) → le Retry fonctionne.\\
\vspace{0.5em}
• CI : Linting strict (ESLint), Prettier, Build success.\\
\vspace{0.5em}
• Doc : README avec diagramme de séquence (Mermaid).\\
\vspace{0.5em}
\vspace{0.5em}
\vspace{0.5em}
\begin{confidential}CONFIDENTIEL — PROJET RBK 2.0                                 Page 111 sur 316\end{confidential}
RBK 2.0 — Whitepaper Architecte Web3                                                                        Chapitre 13. CAPSTONES (PROJETS SIGNATURES)\\
\vspace{0.5em}
\vspace{0.5em}
\subsection{13.2.2 Capstone 2 — Tokenization \textbackslash{}& Admin Control Center (RWA/DeFi)}
\vspace{0.5em}
Focus : Gestion des droits (RBAC1 ) et Logique Métier.\\
\vspace{0.5em}
\subsection{13.2.2.1 Spécifications Fonctionnelles}
\vspace{0.5em}
1. US.1 (Role) : ”En tant que SuperAdmin, je peux nommer ou révoquer un ’Operator’.”\\
\vspace{0.5em}
2. US.2 (Freeze) : ”En tant qu’Operator, je peux geler un compte utilisateur suspect (Compliance).”\\
\vspace{0.5em}
3. US.3 (Audit) : ”En tant qu’Auditeur, je peux consulter l’historique on-chain de toutes les actions Admin.”\\
\vspace{0.5em}
\vspace{0.5em}
\vspace{0.5em}
\vspace{0.5em}
EL\\
\subsection{13.2.2.2 Spécifications Techniques}
\vspace{0.5em}
\vspace{0.5em}
\vspace{0.5em}
\vspace{0.5em}
TI\\
• Stack : Solidity (OZ AccessControl) ou Anchor (PDA Authorities).\\
\vspace{0.5em}
\vspace{0.5em}
\vspace{0.5em}
\vspace{0.5em}
EN\\
• Invariants : TotalSupply doit toujours égaler la somme des balances (sauf mint/burn explicite).\\
\vspace{0.5em}
\vspace{0.5em}
\vspace{0.5em}
\vspace{0.5em}
D\\
• Events : Émission d’événement pour CHAQUE changement d’état critique.\\
\vspace{0.5em}
\vspace{0.5em}
\vspace{0.5em}
\vspace{0.5em}
FI\\
\subsection{13.2.2.3 Checklist Sécurité Dédiée}
N\\
CO\\
Access Control : Vérifier onlyRole sur toutes les fonctions sensibles.\\
\vspace{0.5em}
Two-Step Transfer : Changement d’admin en 2 étapes (Propose + Accept).\\
\vspace{0.5em}
Checks-Effects-Interactions : Respect strict du pattern anti-reentrancy.\\
\vspace{0.5em}
\vspace{0.5em}
\subsection{13.2.3 Capstone 3 — Digital Assets \textbackslash{}& Utility Ecosystem (NFT/Gaming)}
Focus : Performance \textbackslash{}& Vérification de propriété.\\
1\\
Sécurité/Infra/Ops — Role-Based Access Control (RBAC) : modèle limitant les privilèges par rôles prédéfinis pour réduire la surface d’attaque.\\
\vspace{0.5em}
\vspace{0.5em}
\vspace{0.5em}
\vspace{0.5em}
\begin{confidential}CONFIDENTIEL — PROJET RBK 2.0                                             Page 112 sur 316\end{confidential}
RBK 2.0 — Whitepaper Architecte Web3                                                     Chapitre 13. CAPSTONES (PROJETS SIGNATURES)\\
\vspace{0.5em}
\vspace{0.5em}
\subsection{13.2.3.1 Spécifications Fonctionnelles}
\vspace{0.5em}
1. US.1 (Gating) : ”En tant qu’utilisateur, je ne peux accéder au contenu Beta que si je possède le NFT ’Early Bird’.”\\
\vspace{0.5em}
2. US.2 (Claim) : ”En tant que joueur, je peux clamer mon reward une seule fois par jour.”\\
\vspace{0.5em}
\subsection{13.2.3.2 Spécifications Techniques}
\vspace{0.5em}
• Mécanisme : Merkle Tree (Allowlist) ou Signature verifier backend.\\
\vspace{0.5em}
• Optimisation : Utilisation de la compression (cNFT Solana) ou ERC-721A (EVM) pour réduire les coûts de mint.\\
\vspace{0.5em}
\vspace{0.5em}
\vspace{0.5em}
\vspace{0.5em}
EL\\
\subsection{13.2.3.3 Checklist Sécurité Dédiée}
\vspace{0.5em}
\vspace{0.5em}
\vspace{0.5em}
\vspace{0.5em}
TI\\
Bot Protection : Captcha ou signature nonce pour empêcher le mint massif par bots.\\
\vspace{0.5em}
\vspace{0.5em}
\vspace{0.5em}
\vspace{0.5em}
EN\\
Replay Attack : Vérification que la signature n’a pas déjà été utilisée.\\
\vspace{0.5em}
\vspace{0.5em}
\vspace{0.5em}
\vspace{0.5em}
D\\
Metadata : Fichiers hébergés sur IPFS (pas de liens HTTP centralisés).\\
\vspace{0.5em}
\vspace{0.5em}
\vspace{0.5em}
\vspace{0.5em}
FI\\
\subsection{13.3 Matrice d’Évaluation (Rubric Studio)}
N\\
CO\\
Chaque projet est noté sur 100 points selon une grille industrielle stricte.\\
\vspace{0.5em}
\vspace{0.5em}
\vspace{0.5em}
\vspace{0.5em}
\begin{confidential}CONFIDENTIEL — PROJET RBK 2.0                                 Page 113 sur 316\end{confidential}
RBK 2.0 — Whitepaper Architecte Web3                                                          Chapitre 13. CAPSTONES (PROJETS SIGNATURES)\\
\vspace{0.5em}
\vspace{0.5em}
\begin{center}\begin{tabular}TAB. 13.2 : Rubric d’Évaluation Studio (100 pts)\end{tabular}\end{center}
\vspace{0.5em}
\vspace{0.5em}
Pilier      Pts Critères Évalués                                                                    Fail Condition\\
Engineering 40 Architecture propre, Code modulaire, Gestion d’erreurs (Option/Result), Choix de Code monolithique, ”Unwrap” sauvages.\\
structures de données.\\
Sécurité    30 Respect du Threat Model, Validation des inputs, Access Control, Tests de cas limites Vulnérabilité critique, Clé privée dans repo.\\
(Fuzzing basic).\\
Delivery    20 CI/CD fonctionnel, Git Flow propre (Branches, Commits), README complet, Dé- Pas de doc, Build cassé.\\
ploiement reproductible.\\
Comms       10 Vidéo démo claire, Pitch deck concis, Capacité à expliquer les choix techniques Démo illisible, Incapable de justifier un choix.\\
(Trade-offs).\\
\vspace{0.5em}
\vspace{0.5em}
\vspace{0.5em}
\vspace{0.5em}
EL\\
TI\\
EN\\
\subsection{13.4 Délivrables de Sortie (Le ”Package”)}
Chaque étudiant doit remettre un ”Package” zippé (ou repo public) contenant : 1. Le Code : Clean, commenté, testé. 2. La Documen-\\
\vspace{0.5em}
\vspace{0.5em}
\vspace{0.5em}
\vspace{0.5em}
D\\
tation : Architecture, Setup, API. 3. Le Rapport d’Auto-Audit : ”J’ai cherché à me hacker, voici ce que j’ai trouvé”. 4. La Vidéo Démo :\\
\vspace{0.5em}
\vspace{0.5em}
\vspace{0.5em}
\vspace{0.5em}
FI\\
3 minutes max, scénarisée.\\
\vspace{0.5em}
N\\
Ce package est votre passeport pour l’emploi. Il remplace le CV.\\
CO\\
Code         Test       Audit Doc          Release        Demo\\
\vspace{0.5em}
FIG. 13.1 : Pipeline Packaging\\
\vspace{0.5em}
\vspace{0.5em}
\vspace{0.5em}
\vspace{0.5em}
\begin{confidential}CONFIDENTIEL — PROJET RBK 2.0                                   Page 114 sur 316\end{confidential}
\section{FICHES MÉTIERS \textbackslash{}& ÉCONOMIE DU DIPLÔMÉ}
\vspace{0.5em}
\vspace{0.5em}
Ce chapitre détaille les 7 profils de sortie du cursus RBK. Chaque fiche est un standard industriel définissant les attentes, responsabilités\\
\vspace{0.5em}
\vspace{0.5em}
\vspace{0.5em}
\vspace{0.5em}
EL\\
et preuves exigées. L’objectif n’est pas de ”décrire un poste”, mais de rendre le niveau observable : un diplômé RBK doit être capable de\\
produire des livrables auditables, rejouables, et comparables aux standards d’un studio Web3.\\
\vspace{0.5em}
\vspace{0.5em}
\vspace{0.5em}
\vspace{0.5em}
TI\\
EN\\
\subsection{14.1 Fiche Métier 1 : Smart Contract Engineer \textbackslash{}& Auditor (Le « Guardian »)}
\vspace{0.5em}
\vspace{0.5em}
\vspace{0.5em}
\vspace{0.5em}
D\\
Résumé métier Le Guardian est le profil le plus critique : il construit des protocoles qui manipulent de la valeur et il sait les attaquer\\
\vspace{0.5em}
\vspace{0.5em}
\vspace{0.5em}
\vspace{0.5em}
FI\\
mentalement avant que d’autres ne le fassent. Il délivre du code audit-ready, documenté, testé, instrumenté, et il sait gérer le post-prod\\
\vspace{0.5em}
N\\
(incident, patch, gouvernance, communication). Un Guardian qui ne sait pas écrire des tests négatifs et formaliser des invariants n’est pas\\
”junior” : il est dangereux.\\
CO\\
Mission et périmètre\\
\vspace{0.5em}
• Concevoir et livrer des smart contracts (Solana/Anchor ou EVM selon track) avec garanties : invariants, contrôles d’accès, intégrité\\
économique.\\
\vspace{0.5em}
• Réaliser des audits internes et externes : revue ”impitoyable”, modèle de menaces, classification des findings, correctifs, preuves.\\
\vspace{0.5em}
• Organiser la sécurité opérationnelle : Runbooks, Multisig ops, war room, surveillance.\\
\vspace{0.5em}
\vspace{0.5em}
\vspace{0.5em}
\vspace{0.5em}
115\\
RBK 2.0 — Whitepaper Architecte Web3                                                               Chapitre 14. FICHES MÉTIERS \textbackslash{}& ÉCONOMIE DU DIPLÔMÉ\\
\vspace{0.5em}
\vspace{0.5em}
Responsabilités cœur (opérationnelles)\\
\vspace{0.5em}
1. Threat Modeling (avant le code) : Définir actifs critiques, surfaces d’attaque, hypothèses. Produire un threat model exploitable\\
(STRIDE simplifié).\\
\vspace{0.5em}
2. High-Assurance Coding : Formaliser les invariants (soldes, monotonicité). Machine à états explicite.\\
\vspace{0.5em}
3. Review \textbackslash{}& Audit : Revue structurée (checklist auth, CPI, oracles). Rédiger findings et findings.\\
\vspace{0.5em}
4. Hardening (pré-prod) : Tests négatifs, fuzz/invariants, quality gates bloquants.\\
\vspace{0.5em}
5. Emergency Response : War room, diagnostic, patch, rollback, communication.\\
\vspace{0.5em}
\vspace{0.5em}
\vspace{0.5em}
\vspace{0.5em}
EL\\
Trajectoire Carrière \textbackslash{}& Mission\\
\vspace{0.5em}
\vspace{0.5em}
\vspace{0.5em}
\vspace{0.5em}
TI\\
• Junior (0-2 ans) : Écrit des tests, corrige des bugs, audite des modules isolés. Mission : Implémenter le Staking Reward d’un protocole.\\
\vspace{0.5em}
\vspace{0.5em}
\vspace{0.5em}
\vspace{0.5em}
N\\
DE\\
• Senior (3+ ans) : Lead l’architecture, gère la War Room, valide les upgrades critiques. Mission : Sécuriser un Bridge cross-chain à\\
500M\textbackslash{}$ TVL1 .\\
\vspace{0.5em}
\vspace{0.5em}
\vspace{0.5em}
\vspace{0.5em}
FI\\
N\\
Interactions Builder/PM (cadrage), dApp Engineer (API/erreurs), QA Engineer (stratégie tests), Visionnaire (hypothèses éco).\\
CO\\
Livrables standard          docs/threat-model.md, AUDIT\textbackslash{}_REPORT.md (10 pages), Architecture doc, Tests suite (unit, int, fuzz), RUNBOOK.md.\\
\vspace{0.5em}
\vspace{0.5em}
KPIs Taux de findings critiques avant audit externe (”zéro surprise”), MTTR (correction vuln), Couverture utile, Fréquence incidents.\\
\vspace{0.5em}
Table : Matrice Compétence → Preuve → Outil\\
\vspace{0.5em}
Compétence      Preuve attendue                 Outil\\
Threat modeling Threat model (actifs, surfaces) STRIDE\\
\vspace{0.5em}
\vspace{0.5em}
1\\
Web3/Blockchain — Total Value Locked (TVL) : valeur totale des actifs déposés ou immobilisés dans des protocoles DeFi, indicateur de liquidité et de traction.\\
\vspace{0.5em}
\vspace{0.5em}
\begin{confidential}CONFIDENTIEL — PROJET RBK 2.0                                            Page 116 sur 316\end{confidential}
RBK 2.0 — Whitepaper Architecte Web3                                                         Chapitre 14. FICHES MÉTIERS \textbackslash{}& ÉCONOMIE DU DIPLÔMÉ\\
\vspace{0.5em}
\vspace{0.5em}
\vspace{0.5em}
Sécurité          Tests négatifs auth/PDA/CPI       CI Checklist\\
Audit mindset     Audit report protocole tiers      Markdown\\
Fuzz / Invariants Campagne fuzz + rapport           Trident/Foundry\\
Hardening         Quality gates bloquants           CI/CD\\
Perf budget       Profiling compute/gas             CLI / Gas report\\
Post-prod ops     Runbook + war-room drill          Simulation\\
\vspace{0.5em}
\vspace{0.5em}
\vspace{0.5em}
Spec             Threat Model              Code\\
\vspace{0.5em}
\vspace{0.5em}
\vspace{0.5em}
\vspace{0.5em}
EL\\
TI\\
Deploy             Review             Tests (Fuzz)\\
\vspace{0.5em}
\vspace{0.5em}
\vspace{0.5em}
\vspace{0.5em}
D    EN\\
FI\\
Monitor             Incident Drill\\
\vspace{0.5em}
\vspace{0.5em}
N\\
FIG. 14.1 : Boucle Guardian (SecDevOps)\\
CO\\
\vspace{0.5em}
\subsection{14.2 Fiche Métier 2 : Protocol \textbackslash{}& Ecosystem Strategist (Le « Visionnaire »)}
Résumé métier Le Visionnaire transforme une idée en système incitatif robuste. Il conçoit la Tokenomics2 , les règles de gouvernance\\
(DAO), les mécanismes d’incitation, et encadre les risques (techniques, économiques, réglementaires, réputationnels). Il ne code pas le\\
”comment” : il rend le ”quoi” mesurable, testable et défendable.\\
2\\
Web3/Blockchain — Token economics (tokenomics) : conception économique d’un token couvrant émission, distribution, utilité, incitations et mécanismes de\\
capture de valeur.\\
\vspace{0.5em}
\vspace{0.5em}
\vspace{0.5em}
\vspace{0.5em}
\begin{confidential}CONFIDENTIEL — PROJET RBK 2.0                                       Page 117 sur 316\end{confidential}
RBK 2.0 — Whitepaper Architecte Web3                                              Chapitre 14. FICHES MÉTIERS \textbackslash{}& ÉCONOMIE DU DIPLÔMÉ\\
\vspace{0.5em}
\vspace{0.5em}
Mission et périmètre\\
\vspace{0.5em}
• Définir le design économique complet d’un protocole : émissions, distributions, sinks/sources, capture de valeur, mécanismes anti-\\
manipulation.\\
\vspace{0.5em}
• Construire des modèles paramétriques (sheets / notebooks) et des scénarios (stress tests, bank-run, depeg, oracle failure).\\
\vspace{0.5em}
• Spécifier la gouvernance : rôles, pouvoirs d’urgence, timelocks, quorum, garde-fous, processus de décision.\\
\vspace{0.5em}
• Fournir le cadre Go-to-market : bootstrap liquidité, incentives, grants, programmes builders, partenariats.\\
\vspace{0.5em}
Responsabilités cœur (opérationnelles)\\
\vspace{0.5em}
\vspace{0.5em}
\vspace{0.5em}
\vspace{0.5em}
EL\\
1. Tokenomics Design : émission (inflation/désinflation), distribution, vesting, allocations, lockups, mécanismes de burn/buyback.\\
\vspace{0.5em}
\vspace{0.5em}
\vspace{0.5em}
\vspace{0.5em}
TI\\
2. Incentive Modeling : boucles positives vs boucles toxiques (mercenaires), design anti-sybil, anti-wash, anti-farming.\\
\vspace{0.5em}
\vspace{0.5em}
\vspace{0.5em}
\vspace{0.5em}
EN\\
3. Market Microstructure : liquidité (AMM/CLMM), slippage, profondeur, LTV, impact du liquidity mining.\\
\vspace{0.5em}
\vspace{0.5em}
\vspace{0.5em}
\vspace{0.5em}
D\\
4. Governance \textbackslash{}& Control Plane : quorum, délégation, timelocks, emergency powers, procédures d’upgrade.\\
\vspace{0.5em}
\vspace{0.5em}
\vspace{0.5em}
\vspace{0.5em}
FI\\
N\\
5. Risk Framing : depeg, bank-run, oracle manipulation, MEV, bridge risk, dépendances offchain.\\
CO\\
6. GTM \textbackslash{}& Ecosystem : incentives calendar, budget, KPIs d’adoption, programmes devs/communauté.\\
\vspace{0.5em}
Trajectoire Carrière \textbackslash{}& Mission\\
\vspace{0.5em}
• Junior (0–2 ans) : produit des modèles, supporte le PM/Builder, rédige un litepaper clair. Mission : modéliser un staking + emissions\\
schedule + stress tests.\\
\vspace{0.5em}
• Senior (3+ ans) : arbitre la cohérence globale, définit la gouvernance, pilote crises économiques. Mission : redesign tokenomics après\\
attaque d’incentives / mercenary liquidity.\\
\vspace{0.5em}
\vspace{0.5em}
\vspace{0.5em}
\vspace{0.5em}
\begin{confidential}CONFIDENTIEL — PROJET RBK 2.0                               Page 118 sur 316\end{confidential}
RBK 2.0 — Whitepaper Architecte Web3                                                              Chapitre 14. FICHES MÉTIERS \textbackslash{}& ÉCONOMIE DU DIPLÔMÉ\\
\vspace{0.5em}
\vspace{0.5em}
Interactions Guardian (risques/exploitabilité), Builder/PM (scope, contraintes produit), dApp Engineer (UX incentives), Legal/Com-\\
pliance (si tokenisation/RWA), Growth/Community (programmes).\\
\vspace{0.5em}
Livrables  standard LITEPAPER.pdf, TOKENOMICS\textbackslash{}_SPEC.md, SIMULATION\textbackslash{}_MODEL.xlsx                                        (ou   notebook),   RISK\textbackslash{}_REGISTER.md,\\
GOVERNANCE\textbackslash{}_SPEC.md, INCENTIVE\textbackslash{}_PLAN.md (budget, durée, KPI, antifraude).\\
\vspace{0.5em}
\vspace{0.5em}
KPIs Stabilité (volatilité, drawdown), Rétention3 vs mercenaires, croissance organique des intégrations, ratio incentives / usage réel,\\
incidents de gouvernance, ”surprises” économiques évitées.\\
\vspace{0.5em}
Table : Matrice Compétence → Preuve → Outil (Visionnaire)\\
\vspace{0.5em}
\vspace{0.5em}
\vspace{0.5em}
\vspace{0.5em}
EL\\
Compétence Preuve attendue                                      Outil\\
Tokenomics     Spec complète + paramètres + justifications      Markdown / Docs\\
\vspace{0.5em}
\vspace{0.5em}
\vspace{0.5em}
\vspace{0.5em}
TI\\
Simulation     Modèle rejouable + scénarios stress              Sheets / Notebook\\
\vspace{0.5em}
\vspace{0.5em}
\vspace{0.5em}
\vspace{0.5em}
EN\\
Risk           Risk register prob/impact + mitigations          Matrice / Checklist\\
Governance     Règles testables (quorum, timelock, rôles)       DAO spec\\
\vspace{0.5em}
\vspace{0.5em}
\vspace{0.5em}
\vspace{0.5em}
D\\
Microstructure Analyse liquidité / slippage / impact incentives DEX analytics\\
\vspace{0.5em}
\vspace{0.5em}
\vspace{0.5em}
\vspace{0.5em}
FI\\
GTM            Incentive plan + budget + KPIs                   Roadmap / Notion\\
\vspace{0.5em}
\vspace{0.5em}
N\\
CO\\
3\\
Finance/Juridique/Compliance — Capacité à conserver clients ou utilisateurs sur la durée, inverse du churn.\\
\vspace{0.5em}
\vspace{0.5em}
\vspace{0.5em}
\vspace{0.5em}
\begin{confidential}CONFIDENTIEL — PROJET RBK 2.0                                            Page 119 sur 316\end{confidential}
RBK 2.0 — Whitepaper Architecte Web3                                                  Chapitre 14. FICHES MÉTIERS \textbackslash{}& ÉCONOMIE DU DIPLÔMÉ\\
\vspace{0.5em}
\vspace{0.5em}
\vspace{0.5em}
Idea / Thesis           Tokenomics Spec               Simulation\\
\vspace{0.5em}
\vspace{0.5em}
\vspace{0.5em}
\vspace{0.5em}
Governance             Risk Register          Stress Tests\\
\vspace{0.5em}
\vspace{0.5em}
\vspace{0.5em}
\vspace{0.5em}
GTM Plan\\
\vspace{0.5em}
FIG. 14.2 : Boucle Visionnaire (Design → Simulation → Risques → Déploiement)\\
\vspace{0.5em}
\vspace{0.5em}
\vspace{0.5em}
\vspace{0.5em}
EL\\
TI\\
\subsection{14.3 Fiche Métier 3 : Web3 Product Builder / Entrepreneur (Le « Builder »)}
\vspace{0.5em}
\vspace{0.5em}
\vspace{0.5em}
\vspace{0.5em}
EN\\
Résumé métier Le Builder est obsédé par la livraison utilisable. Il transforme une hypothèse en produit, cadre le MVP, orchestre le\\
\vspace{0.5em}
\vspace{0.5em}
\vspace{0.5em}
\vspace{0.5em}
D\\
delivery multi-métiers (smart contracts, front, infra, ops, growth) et garantit la qualité. C’est le rôle qui empêche le projet de devenir un\\
\vspace{0.5em}
\vspace{0.5em}
\vspace{0.5em}
\vspace{0.5em}
FI\\
”prototype éternel”.\\
\vspace{0.5em}
Mission et périmètre                                N\\
CO\\
• Conduire le cycle Discovery → MVP → Release → Iteration.\\
\vspace{0.5em}
• Produire une spécification exécutable : PRD, user stories, acceptance criteria, Definition of Done.\\
\vspace{0.5em}
• Orchestrer sécurité/qualité : checklists, audits internes, bug bounty, observabilité.\\
\vspace{0.5em}
• Tenir la boucle business : activation, rétention, pricing (si SaaS), partenariats.\\
\vspace{0.5em}
Responsabilités cœur (opérationnelles)\\
\vspace{0.5em}
1. Product Discovery : définir personas, pain points, jobs-to-be-done, hypothèses testables.\\
\vspace{0.5em}
\vspace{0.5em}
\begin{confidential}CONFIDENTIEL — PROJET RBK 2.0                                  Page 120 sur 316\end{confidential}
RBK 2.0 — Whitepaper Architecte Web3                                                     Chapitre 14. FICHES MÉTIERS \textbackslash{}& ÉCONOMIE DU DIPLÔMÉ\\
\vspace{0.5em}
\vspace{0.5em}
2. Spec \textbackslash{}& Scope : rédiger PRD, user stories, risques, contraintes chain (fees, latence, finalité).\\
\vspace{0.5em}
3. Delivery Coordination : roadmap, sprints, arbitrages, dépendances (indexer, RPC, wallets).\\
\vspace{0.5em}
4. Quality \textbackslash{}& Release : DoD robuste (tests, docs, rollback), release tagging, changelog.\\
\vspace{0.5em}
5. Business Loop : instrumentation (events), KPIs (activation, retention, churn), itération.\\
\vspace{0.5em}
Trajectoire Carrière \textbackslash{}& Mission\\
\vspace{0.5em}
• Junior : cadre un MVP, gère backlog + release, produit un PRD solide. Mission : lancer un MVP staking avec dashboard et support\\
utilisateur.\\
\vspace{0.5em}
\vspace{0.5em}
\vspace{0.5em}
\vspace{0.5em}
EL\\
• Senior : pilote multi-produits, gère incidents de prod, scale équipes, priorise sécurité/UX. Mission : orchestrer le lancement public d’un\\
\vspace{0.5em}
\vspace{0.5em}
\vspace{0.5em}
\vspace{0.5em}
TI\\
protocole avec bug bounty et plan de rollback.\\
\vspace{0.5em}
\vspace{0.5em}
\vspace{0.5em}
\vspace{0.5em}
N\\
Interactions Visionnaire (hypothèses éco), Guardian (risques et gates), dApp Engineer (UX), QA (strategy), DevRel (docs/adoption),\\
\vspace{0.5em}
\vspace{0.5em}
\vspace{0.5em}
\vspace{0.5em}
DE\\
Ops/SRE (monitoring).\\
\vspace{0.5em}
\vspace{0.5em}
\vspace{0.5em}
\vspace{0.5em}
FI\\
N\\
Livrables standard    PRD (PRD.md), backlog priorisé, roadmap, RELEASE\textbackslash{}_PLAN.md, CHANGELOG.md, ROLLBACK\textbackslash{}_PLAN.md, dashboards (KPIs\\
CO\\
produit + erreurs).\\
\vspace{0.5em}
KPIs Time-to-Value (TTFV), activation (wallet connect → action), rétention, churn, NPS/support load, taux d’échec transactions, cadence\\
release sans régressions.\\
\vspace{0.5em}
Table : Definition of Done (Builder)\\
\vspace{0.5em}
Axe      Critère DoD (minimum studio-grade)\\
Sécurité Audit interne + checklist + restrictions roles + plan incident\\
Qualité Tests unit/int + scénario critique rejouable + CI verte\\
UX       Messages d’erreur clairs + retry/failover + latence sous seuil\\
\vspace{0.5em}
\vspace{0.5em}
\vspace{0.5em}
\begin{confidential}CONFIDENTIEL — PROJET RBK 2.0                                    Page 121 sur 316\end{confidential}
RBK 2.0 — Whitepaper Architecte Web3                                                 Chapitre 14. FICHES MÉTIERS \textbackslash{}& ÉCONOMIE DU DIPLÔMÉ\\
\vspace{0.5em}
\vspace{0.5em}
\vspace{0.5em}
Obs     Dashboard KPIs + erreurs RPC + alerting + logs corrélés\\
Docs    README fresh clone + user guide + runbook support\\
Release Tag + changelog + rollback plan + release notes\\
\vspace{0.5em}
\vspace{0.5em}
\vspace{0.5em}
Discovery          PRD / Scope           Build\\
\vspace{0.5em}
\vspace{0.5em}
\vspace{0.5em}
\vspace{0.5em}
Measure (KPIs)           Release          QA / Gates\\
\vspace{0.5em}
\vspace{0.5em}
\vspace{0.5em}
\vspace{0.5em}
EL\\
FIG. 14.3 : Boucle Builder (Produit → Release → Mesure → Itération)\\
\vspace{0.5em}
\vspace{0.5em}
\vspace{0.5em}
\vspace{0.5em}
TI\\
EN\\
\subsection{14.4 Fiche Métier 4 : Solana dApp Engineer (Front Web3)}
\vspace{0.5em}
\vspace{0.5em}
\vspace{0.5em}
\vspace{0.5em}
D\\
Résumé métier Le dApp Engineer rend une blockchain instable utilisable humainement. Il maîtrise le cycle transactionnel (simula-\\
\vspace{0.5em}
\vspace{0.5em}
\vspace{0.5em}
\vspace{0.5em}
FI\\
tion, signature, envoi, confirmation, finalité), les erreurs RPC, l’UX wallet, et la cohérence des données (indexer vs onchain). C’est le rôle\\
qui transforme des primitives chain en expérience produit fiable.\\
N\\
CO\\
Mission et périmètre\\
\vspace{0.5em}
• Concevoir le front web3 (React/Next/Vite) et l’intégration wallet (Solana Wallet Adapter / EVM wallets selon track).\\
\vspace{0.5em}
• Implémenter une machine à états transactionnelle robuste, avec retries, fallback RPC, et messages d’erreur actionnables.\\
\vspace{0.5em}
• Gérer la couche data : caching, invalidation, indexers, re-sync onchain, pagination, perf.\\
\vspace{0.5em}
• Mettre en place l’observabilité : logs client, métriques UX, erreurs chain, corrélations.\\
\vspace{0.5em}
\vspace{0.5em}
\vspace{0.5em}
\vspace{0.5em}
\begin{confidential}CONFIDENTIEL — PROJET RBK 2.0                                   Page 122 sur 316\end{confidential}
RBK 2.0 — Whitepaper Architecte Web3                                              Chapitre 14. FICHES MÉTIERS \textbackslash{}& ÉCONOMIE DU DIPLÔMÉ\\
\vspace{0.5em}
\vspace{0.5em}
Responsabilités cœur (opérationnelles)\\
\vspace{0.5em}
1. Transaction Lifecycle UI : simulation, préparation comptes, signature, submit, confirm, finality, affichage receipt.\\
\vspace{0.5em}
2. RPC Management : rate limit, failover multi-RPC, backoff, circuit breaker, timeouts.\\
\vspace{0.5em}
3. Wallet UX : états wallet (locked, missing), chain mismatch, multi-wallet, deep links mobile, signature prompts.\\
\vspace{0.5em}
4. Data Layer : indexer lag, fallback onchain, cache (SWR/React Query), revalidation.\\
\vspace{0.5em}
5. Observabilité : tracking erreurs, funnel, latency, replay des transactions échouées.\\
\vspace{0.5em}
Trajectoire Carrière \textbackslash{}& Mission\\
\vspace{0.5em}
\vspace{0.5em}
\vspace{0.5em}
\vspace{0.5em}
EL\\
• Junior : implémente flows tx simples, gère erreurs classiques, met en place cache. Mission : mint NFT + listing + transaction receipts\\
\vspace{0.5em}
\vspace{0.5em}
\vspace{0.5em}
\vspace{0.5em}
TI\\
propres.\\
\vspace{0.5em}
\vspace{0.5em}
\vspace{0.5em}
\vspace{0.5em}
N\\
DE\\
• Senior : architecture data + multi-RPC + indexer strategy, UX de crise, monitoring. Mission : construire une dApp à fort trafic sans\\
”tx stuck”.\\
\vspace{0.5em}
\vspace{0.5em}
\vspace{0.5em}
\vspace{0.5em}
FI\\
N\\
Interactions Guardian (préconditions, erreurs, events), QA (scénarios), Builder (DoD, funnel), DevRel (docs intégration), Ops (RPC\\
CO\\
providers, status pages).\\
\vspace{0.5em}
Livrables standard TX\textbackslash{}_STATE\textbackslash{}_MACHINE.md, ERROR\textbackslash{}_TAXONOMY.md, RPC\textbackslash{}_FAILOVER.md, dashboards erreurs + latences, tests e2e (Play-\\
wright/Cypress) sur flows web3.\\
\vspace{0.5em}
KPIs Tx success rate, temps moyen confirmation, taux ”blockhash expired / nonce issues”, support tickets wallet, TTFV, latence RPC,\\
erreurs indexer.\\
\vspace{0.5em}
Table : Taxonomie erreurs (Extrait)\\
\vspace{0.5em}
\vspace{0.5em}
\vspace{0.5em}
\vspace{0.5em}
\begin{confidential}CONFIDENTIEL — PROJET RBK 2.0                               Page 123 sur 316\end{confidential}
RBK 2.0 — Whitepaper Architecte Web3                                                Chapitre 14. FICHES MÉTIERS \textbackslash{}& ÉCONOMIE DU DIPLÔMÉ\\
\vspace{0.5em}
\vspace{0.5em}
\vspace{0.5em}
Cause                   Mitigation Standard\\
RPC Rate Limit          Backoff exponentiel + failover + circuit breaker\\
Simulation Failed       Précondition claire + CTA ”Fix” + lien doc\\
Blockhash Expired       Auto-refresh + re-sign guidé + gestion timeout\\
Stale Indexer           Fallback onchain + UI ”syncing” + invalidation cache\\
Wallet Locked / Missing Détection + guide utilisateur + deep-link\\
\vspace{0.5em}
\vspace{0.5em}
\vspace{0.5em}
Prepare / Simulate         Sign          Send        Confirm\\
\vspace{0.5em}
\vspace{0.5em}
\vspace{0.5em}
\vspace{0.5em}
EL\\
Retry / Failover        Receipt + UI Update           Finality\\
\vspace{0.5em}
\vspace{0.5em}
\vspace{0.5em}
\vspace{0.5em}
TI\\
EN\\
FIG. 14.4 : Machine à états transactionnelle (dApp Engineer)\\
\vspace{0.5em}
\vspace{0.5em}
\vspace{0.5em}
\vspace{0.5em}
D\\
FI\\
\subsection{14.5 Fiche Métier 5 : Tokenization \textbackslash{}& DePIN Architect}
\vspace{0.5em}
N\\
Résumé métier Le Tokenization \textbackslash{}& DePIN Architect relie le monde réel à la blockchain : actifs, droits, identités, conformité, opérations.\\
CO\\
Il pense cycle de vie (Mint → Transfer → Freeze → Update Policy → Burn) et contrôle (RBAC, audits, traçabilité), avec une lecture\\
”risque + régulateur”.\\
\vspace{0.5em}
Mission et périmètre\\
\vspace{0.5em}
• Concevoir l’architecture de tokenisation (RWA / droits / attestations / DePIN).\\
\vspace{0.5em}
• Définir les rôles et permissions (RBAC), les contrôles d’accès, les garde-fous (timelocks, multi-approval).\\
\vspace{0.5em}
• Mettre en place la conformité opérationnelle : KYC/AML (si requis), audit trail, policy updates.\\
\vspace{0.5em}
• Organiser l’ops : émission, registres, procédures de gel/dégel, recovery, incident.\\
\vspace{0.5em}
\vspace{0.5em}
\begin{confidential}CONFIDENTIEL — PROJET RBK 2.0                                 Page 124 sur 316\end{confidential}
RBK 2.0 — Whitepaper Architecte Web3                                                 Chapitre 14. FICHES MÉTIERS \textbackslash{}& ÉCONOMIE DU DIPLÔMÉ\\
\vspace{0.5em}
\vspace{0.5em}
Responsabilités cœur (opérationnelles)\\
\vspace{0.5em}
1. Asset Modeling : définir l’actif, le droit représenté, le lien juridique, et les événements de cycle de vie.\\
\vspace{0.5em}
2. RBAC \textbackslash{}& Policy : rôles (issuer, admin, compliance, operator), séparation des pouvoirs, journalisation.\\
\vspace{0.5em}
3. Compliance by Design : règles (transfers allowed/blocked), listes (allow/deny), preuves, audits internes.\\
\vspace{0.5em}
4. Lifecycle Ops : procédures mint/burn/freeze, délais, approbations, runbooks.\\
\vspace{0.5em}
5. DePIN Considerations : data provenance, attestations capteurs, anti-fraude, slashing/incentives.\\
\vspace{0.5em}
Trajectoire Carrière \textbackslash{}& Mission\\
\vspace{0.5em}
\vspace{0.5em}
\vspace{0.5em}
\vspace{0.5em}
EL\\
• Junior : modélise un actif simple, implémente RBAC + audit trail. Mission : tokeniser des ”credits” internes avec plafonds et logs.\\
\vspace{0.5em}
\vspace{0.5em}
\vspace{0.5em}
\vspace{0.5em}
TI\\
N\\
• Senior : conçoit une architecture RWA/DePIN complète + procédures conformité. Mission : mettre en production un système de\\
\vspace{0.5em}
\vspace{0.5em}
\vspace{0.5em}
\vspace{0.5em}
DE\\
tokenisation avec freeze, recovery et gouvernance.\\
\vspace{0.5em}
\vspace{0.5em}
\vspace{0.5em}
\vspace{0.5em}
FI\\
Interactions Visionnaire (incitations / économie), Guardian (contrôles d’accès, upgradeability), Legal/Compliance, Ops, dApp Engineer\\
\vspace{0.5em}
\vspace{0.5em}
N\\
(UX transfert/gel).\\
CO\\
Livrables standard    ASSET\textbackslash{}_MODEL.md, RBAC\textbackslash{}_SPEC.md, POLICY\textbackslash{}_ENGINE.md, COMPLIANCE\textbackslash{}_RUNBOOK.md, AUDIT\textbackslash{}_TRAIL\textbackslash{}_SPEC.md, procédures\\
freeze/burn.\\
\vspace{0.5em}
KPIs Incidents conformité, traçabilité complète (audit trail), taux d’erreurs opérationnelles (mint/freeze), temps de réponse incident,\\
alignement policy vs implémentation.\\
\vspace{0.5em}
Table : Matrice RBAC (Extrait)\\
\vspace{0.5em}
\vspace{0.5em}
\vspace{0.5em}
\vspace{0.5em}
\begin{confidential}CONFIDENTIEL — PROJET RBK 2.0                                 Page 125 sur 316\end{confidential}
RBK 2.0 — Whitepaper Architecte Web3                                                  Chapitre 14. FICHES MÉTIERS \textbackslash{}& ÉCONOMIE DU DIPLÔMÉ\\
\vspace{0.5em}
\vspace{0.5em}
Permission   Admin Issuer User Garde-fou Standard\\
Mint          Oui    Oui Non Plafond + logs + approbation\\
Freeze        Oui   Non Non Timelock + justification + audit\\
Transfer      Non   Non Oui Règles policy + limites + recovery\\
Update Policy Oui   Non Non Review + multi-approval + versioning\\
Burn          Oui    Oui Non Motif + audit trail\\
\vspace{0.5em}
\vspace{0.5em}
\vspace{0.5em}
Mint           Transfer          Freeze\\
\vspace{0.5em}
\vspace{0.5em}
\vspace{0.5em}
\vspace{0.5em}
EL\\
Audit Trail          Burn           Update Policy\\
\vspace{0.5em}
\vspace{0.5em}
\vspace{0.5em}
\vspace{0.5em}
TI\\
EN\\
FIG. 14.5 : Cycle de vie tokenisé + Audit Trail (Tokenization \textbackslash{}& DePIN Architect)\\
\vspace{0.5em}
\vspace{0.5em}
\vspace{0.5em}
\vspace{0.5em}
D\\
FI\\
\subsection{14.6 Fiche Métier 6 : Web3 QA \textbackslash{}& Test Automation Engineer}
\vspace{0.5em}
N\\
Résumé métier Le QA Web3 écrit du code qui teste le code. C’est un rôle de sécurité et de fiabilité : il conçoit une stratégie de tests\\
CO\\
multi-niveaux (unit, integration, fork, fuzz, invariants), automatise les gates CI et empêche les régressions. Un QA Web3 qui ne sait pas\\
faire du fork testing et du fuzzing n’est pas ”incomplet” : il est aveugle.\\
\vspace{0.5em}
Mission et périmètre\\
\vspace{0.5em}
• Définir la stratégie qualité (risques, scénarios critiques, niveaux de tests).\\
\vspace{0.5em}
• Implémenter l’automatisation CI : lint, unit, integration, fork, fuzz, coverage, artefacts de preuve.\\
\vspace{0.5em}
• Tester le comportement sous conditions chain : latence, reorg (EVM), congestion, RPC rate limit, comptes manquants.\\
\vspace{0.5em}
• Maintenir une discipline de non-régression : suites rejouables, test data, snapshots.\\
\vspace{0.5em}
\vspace{0.5em}
\begin{confidential}CONFIDENTIEL — PROJET RBK 2.0                                  Page 126 sur 316\end{confidential}
RBK 2.0 — Whitepaper Architecte Web3                                                  Chapitre 14. FICHES MÉTIERS \textbackslash{}& ÉCONOMIE DU DIPLÔMÉ\\
\vspace{0.5em}
\vspace{0.5em}
Responsabilités cœur (opérationnelles)\\
1. Test Strategy : cartographie risques → tests, priorisation ”money paths”.\\
\vspace{0.5em}
2. Automation (CI) : gates bloquants, rapports, artifacts, règles de merge.\\
\vspace{0.5em}
3. Fork / Simulation : tests sur fork mainnet/devnet, scénarios réalistes.\\
\vspace{0.5em}
4. Fuzz / Invariants : propriétés, campagnes fuzz, rapports d’échec minimisés.\\
\vspace{0.5em}
5. Regression Discipline : tests e2e, snapshots, versions, changelog test.\\
\vspace{0.5em}
Trajectoire Carrière \textbackslash{}& Mission\\
\vspace{0.5em}
\vspace{0.5em}
\vspace{0.5em}
\vspace{0.5em}
EL\\
• Junior : écrit tests unit/int, maintient CI, reproduit bugs. Mission : écrire une suite d’intégration pour staking + claim + withdraw.\\
\vspace{0.5em}
\vspace{0.5em}
\vspace{0.5em}
\vspace{0.5em}
TI\\
• Senior : construit frameworks internes, fuzz/invariants systématiques, fork pipelines. Mission : empêcher une régression critique avant\\
audit externe.\\
\vspace{0.5em}
\vspace{0.5em}
\vspace{0.5em}
\vspace{0.5em}
N\\
DE\\
Interactions Guardian (propriétés/invariants), dApp Engineer (e2e), Builder (DoD), Ops (environnements), DevRel (repro guides).\\
\vspace{0.5em}
\vspace{0.5em}
\vspace{0.5em}
\vspace{0.5em}
FI\\
Livrables standard     TEST\textbackslash{}_STRATEGY.md, CI\textbackslash{}_GATES.md, suites unit/int/fork/fuzz, rapports coverage, REGRESSION\textbackslash{}_PLAYBOOK.md, bugs\\
repro steps.\\
N\\
CO\\
KPIs Taux de régression, couverture utile (pas ”ligne”), temps de reproduction bug, taux de flakiness, findings évités, qualité des artefacts\\
CI.\\
Table : Pipeline Qualité\\
Étape           Gate Bloquant (minimum)\\
Lint/Format     Échec si style/format non conforme\\
Unit Tests      Échec si logique locale incorrecte\\
Integration     Échec si scénario critique cassé\\
Fork Tests      Échec si comportement diverge sur fork\\
Fuzz/Invariants Échec si invariant violé ou crash\\
\vspace{0.5em}
\vspace{0.5em}
\begin{confidential}CONFIDENTIEL — PROJET RBK 2.0                                  Page 127 sur 316\end{confidential}
RBK 2.0 — Whitepaper Architecte Web3                                                Chapitre 14. FICHES MÉTIERS \textbackslash{}& ÉCONOMIE DU DIPLÔMÉ\\
\vspace{0.5em}
\vspace{0.5em}
Table : Matrice Compétence → Preuve → Outil (QA Web3)\\
\vspace{0.5em}
Compétence       Preuve attendue                  Outil\\
Unit/Integration Suite rejouable + coverage       Foundry/Anchor tests\\
Fork testing     Scénarios mainnet-like           Anvil/Hardhat fork\\
Fuzz/Invariants Campagne + rapport + minimisation Foundry fuzz/Trident\\
CI Gates         Pipeline bloquante + artefacts   CI/CD\\
E2E dApp         Flows wallet + tx receipts       Playwright/Cypress\\
\vspace{0.5em}
\vspace{0.5em}
\vspace{0.5em}
\subsection{14.7 Fiche Métier 7 : Developer Advocate \textbackslash{}& Technical Writer}
\vspace{0.5em}
\vspace{0.5em}
\vspace{0.5em}
\vspace{0.5em}
EL\\
Résumé métier Le DevRel/Technical Writer est la voix technique du protocole. Il rend le produit adoptable via documentation, SDKs,\\
\vspace{0.5em}
\vspace{0.5em}
\vspace{0.5em}
\vspace{0.5em}
TI\\
exemples, support d’intégration, et boucle de feedback vers l’équipe core. C’est un multiplicateur de croissance : il transforme une tech\\
”bonne” en tech ”utilisée”.\\
\vspace{0.5em}
\vspace{0.5em}
\vspace{0.5em}
\vspace{0.5em}
EN\\
Mission et périmètre\\
\vspace{0.5em}
\vspace{0.5em}
\vspace{0.5em}
\vspace{0.5em}
D\\
FI\\
• Produire une documentation complète : Quickstart, guides, API, architecture, runbooks intégration.\\
\vspace{0.5em}
N\\
• Maintenir des SDKs, examples, starters et templates.\\
CO\\
• Assurer le support technique (issues, Discord, forums), et formaliser les patterns de résolution.\\
\vspace{0.5em}
• Organiser le feedback produit : bugs, frictions d’intégration, demandes API, priorisation.\\
\vspace{0.5em}
Responsabilités cœur (opérationnelles)\\
\vspace{0.5em}
1. Docs Architecture : structure doc (learning path), style guide, versioning, dépréciations.\\
\vspace{0.5em}
2. SDKs \textbackslash{}& Examples : starters maintenus, samples ”copy-paste safe”, tests des snippets.\\
\vspace{0.5em}
3. Integration Support : playbooks d’intégration, debug guides, taxonomie d’erreurs.\\
\vspace{0.5em}
\vspace{0.5em}
\begin{confidential}CONFIDENTIEL — PROJET RBK 2.0                                 Page 127 sur 316\end{confidential}
RBK 2.0 — Whitepaper Architecte Web3                                              Chapitre 14. FICHES MÉTIERS \textbackslash{}& ÉCONOMIE DU DIPLÔMÉ\\
\vspace{0.5em}
\vspace{0.5em}
4. Developer Experience : friction mapping, time-to-first-success, simplification flows.\\
\vspace{0.5em}
5. Content \textbackslash{}& Adoption : posts techniques, workshops, templates, bounties doc.\\
\vspace{0.5em}
Trajectoire Carrière \textbackslash{}& Mission\\
\vspace{0.5em}
• Junior : écrit docs, maintient examples, répond support niveau 1. Mission : produire un Quickstart intégration wallet + transaction +\\
indexer.\\
\vspace{0.5em}
• Senior : pilote DX, versioning docs/SDK, influence roadmap produit. Mission : réduire de 50\textbackslash{}% le temps d’intégration moyen via docs +\\
SDK.\\
\vspace{0.5em}
\vspace{0.5em}
\vspace{0.5em}
\vspace{0.5em}
EL\\
Interactions dApp Engineer (DX et erreurs), Guardian (sécurité des patterns), Builder (priorités), QA (snippets testés), Community/-\\
\vspace{0.5em}
\vspace{0.5em}
\vspace{0.5em}
\vspace{0.5em}
TI\\
Growth.\\
\vspace{0.5em}
\vspace{0.5em}
\vspace{0.5em}
\vspace{0.5em}
EN\\
Livrables standard  Doc set complet (Quickstart, Concepts, API, Guides), STARTER\textbackslash{}_KIT maintenu, INTEGRATION\textbackslash{}_PLAYBOOK.md,\\
\vspace{0.5em}
\vspace{0.5em}
\vspace{0.5em}
\vspace{0.5em}
D\\
TROUBLESHOOTING.md, templates (CI, env, config), FAQ.\\
\vspace{0.5em}
\vspace{0.5em}
\vspace{0.5em}
\vspace{0.5em}
FI\\
KPIs Time-to-first-success, taux d’intégrations réussies, support load, résolution issues, adoption SDK, taux de docs à jour, satisfaction\\
devs.                                            N\\
CO\\
Table : Matrice Compétence → Preuve → Outil (DevRel)\\
\vspace{0.5em}
Compétence       Preuve attendue                  Outil\\
Docs system      Doc set versionné + style guide  Docusaurus/MkDocs\\
Snippets fiables Snippets testés en CI            CI + examples\\
SDK \textbackslash{}& starters Starter kit maintenu + releases    NPM/Cargo\\
Support          Playbook + base de connaissances Issues/Discord\\
DX analytics     Mesure TFFS + friction points    Analytics\\
\vspace{0.5em}
\vspace{0.5em}
\vspace{0.5em}
\vspace{0.5em}
\begin{confidential}CONFIDENTIEL — PROJET RBK 2.0                               Page 128 sur 316\end{confidential}
RBK 2.0 — Whitepaper Architecte Web3                                                   Chapitre 14. FICHES MÉTIERS \textbackslash{}& ÉCONOMIE DU DIPLÔMÉ\\
\vspace{0.5em}
\vspace{0.5em}
\vspace{0.5em}
\subsection{14.8 Perspectives Économiques \textbackslash{}& Carrière}
Nous cadrons ici les fourchettes de revenus et les conditions pour atteindre chaque palier, afin de relier compétences et réalité marché.\\
Principe RBK : le salaire est corrélé à la réduction de risque que vous apportez (sécurité, fiabilité, delivery, adoption). Les profils capables\\
de produire des preuves studio-grade atteignent les hauts de fourchettes.\\
\vspace{0.5em}
\vspace{0.5em}
\subsection{14.8.1 Revenus Annuels Cibles 2025}
Ce tableau présente des ordres de grandeur. Le haut de fourchette n’est accessible qu’avec des preuves de compétence ”Studio-Grade”\\
(portfolio, livrables, track record, recommandations). Note : Le différentiel apparent ”Tunisie vs Remote” doit être pondéré par le coût de la\\
vie et la fiscalité. À compétences égales, le plafond ”remote global” dépend aussi du niveau d’anglais, de la capacité à travailler async, et de la\\
\vspace{0.5em}
\vspace{0.5em}
\vspace{0.5em}
\vspace{0.5em}
EL\\
crédibilité du portfolio.\\
\vspace{0.5em}
\vspace{0.5em}
\vspace{0.5em}
\vspace{0.5em}
TI\\
\begin{center}\begin{tabular}TAB. 14.7 : Revenus Indicatifs (2025)\end{tabular}\end{center}
\vspace{0.5em}
\vspace{0.5em}
\vspace{0.5em}
\vspace{0.5em}
EN\\
Métier          Remote Global Tunisie          Condition Top Tier (preuve)\\
\vspace{0.5em}
\vspace{0.5em}
\vspace{0.5em}
\vspace{0.5em}
D\\
FI\\
Guardian        \textbackslash{}$80k – \textbackslash{}$150k    5k – 10k TND 3 repos + Audit Report + fuzz/invariants + runbook\\
\vspace{0.5em}
Auditor (Elite) \textbackslash{}$120k – \textbackslash{}$250k N/A\\
N           Track record de findings critiques + rapports publics\\
CO\\
Strategist      \textbackslash{}$90k – \textbackslash{}$160k    Consultant     Modèles rejouables + risk register + décisions défendables\\
\vspace{0.5em}
Builder         \textbackslash{}$70k – \textbackslash{}$140k    4k – 9k TND    PRD + releases régulières + obs + incidents gérés proprement\\
\vspace{0.5em}
dApp Eng.       \textbackslash{}$60k – \textbackslash{}$110k    3k – 6k TND    UX tx irréprochable + failover RPC + e2e + métriques\\
\vspace{0.5em}
Token Arch.     \textbackslash{}$90k – \textbackslash{}$180k    Consultant     RBAC/policy + compliance ops + audit trail + lifecycle complet\\
\vspace{0.5em}
QA Eng.         \textbackslash{}$60k – \textbackslash{}$120k    3k – 7k TND    CI gates + fork suite + fuzz + faible flakiness\\
\vspace{0.5em}
DevRel          \textbackslash{}$50k – \textbackslash{}$110k    3k – 6k TND    Doc system + SDK + starters + baisse friction intégration\\
\vspace{0.5em}
\vspace{0.5em}
\vspace{0.5em}
\vspace{0.5em}
\begin{confidential}CONFIDENTIEL — PROJET RBK 2.0                                   Page 129 sur 316\end{confidential}
RBK 2.0 — Whitepaper Architecte Web3                                                   Chapitre 14. FICHES MÉTIERS \textbackslash{}& ÉCONOMIE DU DIPLÔMÉ\\
\vspace{0.5em}
\vspace{0.5em}
\subsection{14.8.2 Ce qui compose réellement la rémunération (important)}
• Salaire fixe (base).\\
\vspace{0.5em}
• Tokens/Equity (souvent vestés, risqués) : valeur dépendante du marché.\\
\vspace{0.5em}
• Bonus / bounties (audit findings, bug bounty, delivery).\\
\vspace{0.5em}
• Freelance/consulting : facturation au livrable, souvent plus lucrative si portfolio solide.\\
\vspace{0.5em}
\vspace{0.5em}
\subsection{14.8.3 Comment atteindre le palier}
\vspace{0.5em}
\vspace{0.5em}
\vspace{0.5em}
\vspace{0.5em}
EL\\
Palier commun ”RBK Ready”\\
\vspace{0.5em}
\vspace{0.5em}
\vspace{0.5em}
\vspace{0.5em}
TI\\
• Portfolio GitHub : 3 repos studio-grade (tests, docs, CI, releases).\\
\vspace{0.5em}
\vspace{0.5em}
\vspace{0.5em}
\vspace{0.5em}
EN\\
• 1 démo rejouable + 1 runbook (post-prod).\\
\vspace{0.5em}
• Communication pro : README, changelog, versioning, tickets/PR propres.\\
\vspace{0.5em}
\vspace{0.5em}
\vspace{0.5em}
\vspace{0.5em}
D\\
FI\\
• Preuves : pas ”j’ai fait”, mais ”voici les artefacts, reproductibles, audités”.\\
\vspace{0.5em}
N\\
CO\\
Preuves par métier\\
\vspace{0.5em}
• Guardian : threat model + audit report + tests négatifs + fuzz/invariants.\\
\vspace{0.5em}
• Visionnaire : simulation paramétrique + stress tests + risk register + governance spec.\\
\vspace{0.5em}
• Builder : PRD + backlog + releases + métriques + rollback plan.\\
\vspace{0.5em}
• dApp : tx state machine + error taxonomy + failover RPC + e2e.\\
\vspace{0.5em}
• Token Arch : asset model + RBAC/policy + audit trail + runbooks compliance.\\
\vspace{0.5em}
• QA : CI gates + fork suite + fuzz/invariants + anti-flakiness.\\
\vspace{0.5em}
\vspace{0.5em}
\begin{confidential}CONFIDENTIEL — PROJET RBK 2.0                                  Page 130 sur 316\end{confidential}
RBK 2.0 — Whitepaper Architecte Web3                                                Chapitre 14. FICHES MÉTIERS \textbackslash{}& ÉCONOMIE DU DIPLÔMÉ\\
\vspace{0.5em}
\vspace{0.5em}
• DevRel : doc system versionné + starter kits + SDK + playbooks intégration.\\
\vspace{0.5em}
Disqualifiants (marché)\\
\vspace{0.5em}
Pas de tests négatifs, documentation inexistante, CI rouge, absence de specs, ”projet non rejouable” (setup flou), incapacité à produire\\
un runbook incident.\\
\vspace{0.5em}
Accélérateurs (marché)\\
1 rapport d’audit public solide, 1 démo utilisée par des tiers, 1 starter kit populaire, contributions open-source, capacité à gérer un\\
incident et à communiquer proprement (post-mortem).\\
\vspace{0.5em}
\vspace{0.5em}
\vspace{0.5em}
\vspace{0.5em}
EL\\
TI\\
D    EN\\
FI\\
N\\
CO\\
\vspace{0.5em}
\vspace{0.5em}
\vspace{0.5em}
\vspace{0.5em}
\begin{confidential}CONFIDENTIEL — PROJET RBK 2.0                                 Page 131 sur 316\end{confidential}
\section{BUSINESS PLAN \textbackslash{}& STRATÉGIE DE CROISSANCE}
\vspace{0.5em}
\vspace{0.5em}
\vspace{0.5em}
\vspace{0.5em}
EL\\
\subsection{15.1 Modèle Économique Hybride}
RBK 2.0 repose sur une diversification des sources de revenus pour garantir sa pérennité indépendamment des cycles du marché crypto.\\
\vspace{0.5em}
\vspace{0.5em}
\vspace{0.5em}
\vspace{0.5em}
TI\\
D     EN\\
B2C (70\textbackslash{}%)          ISA (20\textbackslash{}%)      B2B (10\textbackslash{}%)\\
\vspace{0.5em}
\vspace{0.5em}
\vspace{0.5em}
\vspace{0.5em}
FI\\
Corporate\\
Formation Initiale   Success Fees\\
\vspace{0.5em}
N\\
CO\\
FIG. 15.1 : Mix Revenus Cible (Année 3)\\
\vspace{0.5em}
\vspace{0.5em}
\vspace{0.5em}
\vspace{0.5em}
132\\
RBK 2.0 — Whitepaper Architecte Web3                                                                           Chapitre 15. Business Plan\\
\vspace{0.5em}
\vspace{0.5em}
\subsection{15.1.1 Modèle économique détaillé — Waterfall \textbackslash{}& caps mentors}
\vspace{0.5em}
Référence contractuelle\\
Le modèle économique et les flux financiers entre les entités partenaires sont précisément décrits dans les annexes. Ces documents\\
définissent les mécanismes de partage des revenus, de plafonnement des coûts et de protection des marges.\\
\vspace{0.5em}
• Pour le détail des flux financiers (Management Fee, Mentor Pool), voir l’addendum financier en Annexe, section C.\\
\vspace{0.5em}
• Pour le plafonnement des heures des mentors (caps), voir l’Annexe T.\\
\vspace{0.5em}
\vspace{0.5em}
\vspace{0.5em}
\vspace{0.5em}
EL\\
\subsection{15.2 Hypothèses \textbackslash{}& Sources du Modèle}
\vspace{0.5em}
\vspace{0.5em}
\vspace{0.5em}
\vspace{0.5em}
TI\\
Ce modèle repose sur des données de marché comparables (Bootcamps Web3, Market Research 2024-2025) et sur le modèle financier\\
contractuel défini dans l’Annexe B — Modèle financier \textbackslash{}& Waterfall contractuel (cf. C).\\
\vspace{0.5em}
\vspace{0.5em}
\vspace{0.5em}
\vspace{0.5em}
D   EN\\
\subsection{15.2.1 Hypothèses Clés}
\vspace{0.5em}
\vspace{0.5em}
\vspace{0.5em}
\vspace{0.5em}
FI\\
Scénarios macro-économiques utilisés pour la modélisation\\
\vspace{0.5em}
N\\
CO\\
Hypothèse              Scénario      Scénario     Scénario\\
prudent        réaliste   ambitieux\\
Coût d’acquisition     800 TND       500 TND      300 TND\\
(CAC) par étudiant\\
converti\\
Taux de placement à      40\textbackslash{}%             65\textbackslash{}%        85\textbackslash{}%\\
6 mois\\
\vspace{0.5em}
\vspace{0.5em}
\vspace{0.5em}
\vspace{0.5em}
\begin{confidential}CONFIDENTIEL — PROJET RBK 2.0                                Page 133 sur 316\end{confidential}
RBK 2.0 — Whitepaper Architecte Web3                                                                          Chapitre 15. Business Plan\\
\vspace{0.5em}
\vspace{0.5em}
\vspace{0.5em}
Salaire brut de sortie   2 000 TND     3 000 TND   4 000 TND\\
(médiane)\\
Part d’apprenants           50\textbackslash{}%             40\textbackslash{}%      30\textbackslash{}%\\
finançant via ISA\\
Volume annuel                24             30        36\\
d’admis (Genesis)\\
Taux de graduation          65\textbackslash{}%             80\textbackslash{}%      90\textbackslash{}%\\
Genesis → Architecte\\
\vspace{0.5em}
\vspace{0.5em}
\vspace{0.5em}
\vspace{0.5em}
EL\\
Part de revenus              8\textbackslash{}%             12\textbackslash{}%      18\textbackslash{}%\\
B2B/Corporate\\
\vspace{0.5em}
\vspace{0.5em}
\vspace{0.5em}
\vspace{0.5em}
TI\\
Budget marketing du      75 000 TND 75 000 TND 75 000 TND\\
\vspace{0.5em}
\vspace{0.5em}
\vspace{0.5em}
\vspace{0.5em}
EN\\
Sprint 0 (90 jours)\\
Dotation initiale du     150 000 TND (compte séquestré RBK)\\
\vspace{0.5em}
\vspace{0.5em}
\vspace{0.5em}
\vspace{0.5em}
D\\
fonds de garantie ISA\\
\vspace{0.5em}
\vspace{0.5em}
\vspace{0.5em}
\vspace{0.5em}
FI\\
N\\
Le scénario réaliste est retenu comme base pour les projections financières présentées ci-dessous. Les scénarios prudent et ambitieux\\
CO\\
servent aux tests de robustesse (voir Analyse de sensibilité).\\
\vspace{0.5em}
\vspace{0.5em}
\subsection{15.2.2 Structure des Coûts Directs}
• Forfait Nexus Réussite : 3 300 TND par étudiant et par cohorte couvrant la conception pédagogique, la plateforme Venture Engine\\
et le pilotage qualité.\\
\vspace{0.5em}
• Pool mentors : variable indexée sur le volume d’heures effectives (20\textbackslash{}% du chiffre d’affaires formation dans le scénario réaliste,\\
plafonné par SOW cohorte).\\
\vspace{0.5em}
• Opérations RBK : masse salariale (opérations, carrière, juridique), loyers, assurances et coûts de conformité estimés à 240 000 TND\\
\vspace{0.5em}
\vspace{0.5em}
\begin{confidential}CONFIDENTIEL — PROJET RBK 2.0                                   Page 134 sur 316\end{confidential}
RBK 2.0 — Whitepaper Architecte Web3                                                                                           Chapitre 15. Business Plan\\
\vspace{0.5em}
\vspace{0.5em}
la première année (+30 000 TND/an).\\
\vspace{0.5em}
• Marketing acquisition : Sprint 0 à 75 000 TND, puis CAC conforme au scénario (500 TND réaliste).\\
\vspace{0.5em}
• Dotation fonds de garantie ISA : 150 000 TND affectés en année 1 dans un compte séquestré dédié, maintenu ensuite par des\\
abondements annuels ciblés de 30 000 TND.\\
\vspace{0.5em}
\vspace{0.5em}
\subsection{15.3 Funnel d’Acquisition \textbackslash{}& Sourcing}
Le modèle de rentabilité dépend de la sélectivité en entrée (qualité des profils = placement garanti).\\
\vspace{0.5em}
\vspace{0.5em}
\vspace{0.5em}
\vspace{0.5em}
EL\\
1. Top of Funnel1 (Leads) : 1000 Inscrits Webinaire / Downloads Livre Blanc.\\
\vspace{0.5em}
\vspace{0.5em}
\vspace{0.5em}
\vspace{0.5em}
TI\\
2. Middle of Funnel (Applicants) : 200 Candidats passent le test technique Python/JS.\\
\vspace{0.5em}
\vspace{0.5em}
\vspace{0.5em}
\vspace{0.5em}
EN\\
3. Bottom of Funnel (Selection) : 50 Admissibles après entretien.\\
\vspace{0.5em}
4. Taux de conversion2 (Cohorte) : 30 Inscrits payants (ou boursiers validés).\\
\vspace{0.5em}
\vspace{0.5em}
\vspace{0.5em}
\vspace{0.5em}
D\\
FI\\
Le ratio cible est de 3\textbackslash{}% de conversion Lead → Client, assurant une densité de talents élevée.\\
\vspace{0.5em}
N\\
CO\\
\subsection{15.4 Le Pilier B2B : Corporate Upskilling}
Pour réduire la dépendance aux frais de scolarité individuels, nous lançons une offre dédiée aux entreprises (Banques, ESN, Telcos)\\
souhaitant monter une ”Blockchain Factory” interne.\\
Offre ”Corporate Cohort” :\\
\vspace{0.5em}
• Principe : Une entreprise réserve un bloc de 3 à 5 sièges dans une cohorte pour ses employés.\\
\vspace{0.5em}
• Tarif : 15 900 TND / siège (Premium Pricing).\\
1\\
Projet/Management/Qualité — Parcours d’acquisition décomposé en étapes (awareness → consideration → conversion) avec mesures à chaque étape.\\
2\\
Projet/Management/Qualité — Proportion d’utilisateurs qui passent à l’étape suivante du funnel ou réalisent l’action cible.\\
\vspace{0.5em}
\vspace{0.5em}
\begin{confidential}CONFIDENTIEL — PROJET RBK 2.0                                        Page 135 sur 316\end{confidential}
RBK 2.0 — Whitepaper Architecte Web3                                                                            Chapitre 15. Business Plan\\
\vspace{0.5em}
\vspace{0.5em}
• Avantages : Suivi RH dédié, Capstone orienté sur un Use-Case de l’entreprise, Clause de confidentialité.\\
\vspace{0.5em}
• Objectif : 30\textbackslash{}% du CA total en Année 3.\\
\vspace{0.5em}
\vspace{0.5em}
\subsection{15.5 Trajectoire Financière (36 Mois)}
La trésorerie est le nerf de la guerre. Notre modèle prend en compte le ”Cash Drag” (décalage) des ISA.\\
Projection du Volume Étudiant (Funnel) — scénario réaliste :\\
\vspace{0.5em}
Palier               Année 1 Année 2 Année 3\\
Genesis Pool           30      60      90\\
\vspace{0.5em}
\vspace{0.5em}
\vspace{0.5em}
\vspace{0.5em}
EL\\
Explorer               27      54      81\\
Bâtisseur/Architecte   24      48      72\\
\vspace{0.5em}
\vspace{0.5em}
\vspace{0.5em}
\vspace{0.5em}
TI\\
D   EN\\
FI\\
N\\
CO\\
\vspace{0.5em}
\vspace{0.5em}
\vspace{0.5em}
\vspace{0.5em}
\begin{confidential}CONFIDENTIEL — PROJET RBK 2.0                               Page 136 sur 316\end{confidential}
RBK 2.0 — Whitepaper Architecte Web3                                                                                Chapitre 15. Business Plan\\
\vspace{0.5em}
\vspace{0.5em}
\begin{center}\begin{tabular}TAB. 15.2 : Projection financière — scénario réaliste (TND)\end{tabular}\end{center}
\vspace{0.5em}
\vspace{0.5em}
Indicateur                                 Année 1          Année 2          Année 3\\
Effectif Genesis annuel                          30              60               90\\
Diplômés (80\textbackslash{}% de rétention)                      24              48               72\\
CA formation facturé (B2C)                  477 000          954 000        1 431 000\\
Revenus B2B / Corporate                      60 000          120 000          180 000\\
CA total (formation + B2B)                  537 000        1 074 000        1 611 000\\
Encaissements ISA de l’année                         0        90 000          216 000\\
\vspace{0.5em}
\vspace{0.5em}
\vspace{0.5em}
\vspace{0.5em}
EL\\
Forfait Nexus Réussite                       99 000          198 000          297 000\\
Pool mentors (20\textbackslash{}% du CA formation)           95 400          190 800          286 200\\
\vspace{0.5em}
\vspace{0.5em}
\vspace{0.5em}
\vspace{0.5em}
TI\\
Charges opérations RBK                 240 000 TND       270 000 TND   320 000 TND\\
\vspace{0.5em}
\vspace{0.5em}
\vspace{0.5em}
\vspace{0.5em}
EN\\
Marketing (CAC scénario)                75 000 TND        30 000 TND       45 000 TND\\
Dotation fonds de garantie ISA         150 000 TND        30 000 TND       30 000 TND\\
\vspace{0.5em}
\vspace{0.5em}
\vspace{0.5em}
\vspace{0.5em}
D\\
Résultat opérationnel après ISA           −122 400           445 200          848 800\\
\vspace{0.5em}
\vspace{0.5em}
\vspace{0.5em}
\vspace{0.5em}
FI\\
N\\
CO\\
Lecture clé\\
• Année 1 volontairement déficitaire : la dotation du fonds de garantie et le budget marketing Sprint 0 génèrent un déficit de\\
−122 400 TND, pré-financé par l’apport initial et/ou une ligne de crédit court terme.\\
\vspace{0.5em}
• Point mort : l’activité devient cash-flow positive dès l’année 2 avec un excédent opérationnel de 445 200 TND, permettant de\\
reconstituer les réserves et de financer la montée en puissance.\\
\vspace{0.5em}
• Levier ISA : les encaissements différés dépassent 200 000 TND en année 3 ; ils sont systématiquement fléchés vers le fonds de\\
garantie et le renforcement du service carrière.\\
\vspace{0.5em}
Analyse de Trésorerie : L’Année 1 est financée par le mix Upfront. L’effet de levier ISA commence à impacter significativement la\\
\vspace{0.5em}
\vspace{0.5em}
\begin{confidential}CONFIDENTIEL — PROJET RBK 2.0                                 Page 137 sur 316\end{confidential}
RBK 2.0 — Whitepaper Architecte Web3                                                                             Chapitre 15. Business Plan\\
\vspace{0.5em}
\vspace{0.5em}
trésorerie au milieu de l’Année 2, créant un ”Fond de Roulement” naturel pour l’expansion.\\
\vspace{0.5em}
\vspace{0.5em}
\subsection{15.6 Analyse de Sensibilité}
Nous stress-testons le modèle selon 3 variables critiques.\\
\vspace{0.5em}
\begin{center}\begin{tabular}TAB. 15.4 : Analyse de sensibilité sur le scénario réaliste\end{tabular}\end{center}
\vspace{0.5em}
Variable testée                                 Fourchette Effet observé\\
Taux de placement (prudent 40\textbackslash{}% / ambitieux 85\textbackslash{}%) 40\textbackslash{}% → 85\textbackslash{}% EBITDA A3 passe de 848 800 TND à 312 000 TND (prudent) ou 1 042 000 TND (ambitieux).\\
CAC (800 TND / 300 TND)                         ±200 TND Le besoin marketing annuel varie de +90 000 TND (prudent) à −54 000 TND (ambitieux).\\
\vspace{0.5em}
\vspace{0.5em}
\vspace{0.5em}
\vspace{0.5em}
EL\\
Part d’apprenants en ISA (50\textbackslash{}% / 30\textbackslash{}%)            30\textbackslash{}% → 50\textbackslash{}% Le cash différé à 12 mois augmente de 180 000 TND en scénario prudent, exigeant une ligne\\
\vspace{0.5em}
\vspace{0.5em}
\vspace{0.5em}
\vspace{0.5em}
TI\\
Le scénario prudent impose un besoin de trésorerie supplémentaire de 180 000 TND (fonds de roulement ISA + marketing), ce qui justifie la\\
\vspace{0.5em}
\vspace{0.5em}
\vspace{0.5em}
\vspace{0.5em}
EN\\
sécurisation d’une ligne de crédit confirmée avant lancement. À l’inverse, le scénario ambitieux permet d’atteindre la marge opérationnelle\\
cible (\textgreater{}45\textbackslash{}%) dès l’année 2 sous réserve d’une capacité d’absorption mentors certifiée par Nexus.\\
\vspace{0.5em}
\vspace{0.5em}
\vspace{0.5em}
\vspace{0.5em}
D\\
FI\\
\subsection{15.6.1 Gestion du Risque Crédit ISA}
\vspace{0.5em}
N\\
L’Income Share Agreement est traité comme un portefeuille de créances à risque. La gouvernance financière s’appuie sur trois garde-fous\\
CO\\
complémentaires :\\
\vspace{0.5em}
1. Fonds de garantie séquestré : 150 000 TND déposés dès l’année 1 dans un compte bloqué RBK dédié aux défauts, alimentés par\\
20\textbackslash{}% des encaissements upfront (acompte + premières mensualités ISA), puis complétés chaque année (30 000 TND) pour refléter\\
la montée en charge des cohortes.\\
\vspace{0.5em}
2. Recouvrement professionnel : mandat confié à un cabinet spécialisé (ex. TCM Tunisia pour le national, Intrum pour l’international)\\
avec lettre de mission signée avant ouverture de la cohorte. RBK déclenche le recouvrement seulement après avis du Comité éthique\\
(Chapitre 5).\\
\vspace{0.5em}
3. Incitations alignées :\\
\vspace{0.5em}
\vspace{0.5em}
\begin{confidential}CONFIDENTIEL — PROJET RBK 2.0                                   Page 138 sur 316\end{confidential}
RBK 2.0 — Whitepaper Architecte Web3                                                                                               Chapitre 15. Business Plan\\
\vspace{0.5em}
\vspace{0.5em}
• Bonus de paiement anticipé : remise contractuelle de 20\textbackslash{}% du solde si l’apprenant rembourse intégralement dans les 6 premiers\\
mois d’emploi.\\
• Signal réputationnel contrôlé : le SBT3 ”Launch Proof” est marqué ”en anomalie” après validation contradictoire, sans effacer\\
la preuve de compétences.\\
• Assurance conformité : toutes les données de revenus sont collectées via un mandat SEPA / attestations employeur, puis\\
archivées dans Venture Engine pour audit.\\
\vspace{0.5em}
Ce dispositif maintient la responsabilité légale chez RBK tout en sécurisant les flux pour Nexus et pour les financeurs tiers.\\
\vspace{0.5em}
\vspace{0.5em}
\vspace{0.5em}
\vspace{0.5em}
EL\\
\subsection{15.7 Financements et Partenariats Stratégiques}
Pour accélérer sans diluer le capital, RBK active les leviers non-dilutifs :\\
\vspace{0.5em}
\vspace{0.5em}
\vspace{0.5em}
\vspace{0.5em}
TI\\
EN\\
\subsection{15.7.1 1. Écosystème Web3 (Grants)}
\vspace{0.5em}
\vspace{0.5em}
\vspace{0.5em}
\vspace{0.5em}
D\\
• Solana Foundation : Demande de grant ”Education” pour financer les serveurs et les bourses (Target : 50k\textbackslash{}$).\\
\vspace{0.5em}
\vspace{0.5em}
\vspace{0.5em}
\vspace{0.5em}
FI\\
• Superteam : Sponsoring des Hackathons de fin de cohorte (Prize pool).\\
\vspace{0.5em}
N\\
CO\\
\subsection{15.7.2 2. Bailleurs de Fonds Institutionnels}
• Union Européenne (Erasmus+ / Horizon Europe) : Projets de mobilité des talents numériques Afrique-Europe.\\
\vspace{0.5em}
• Banque Africaine de Développement (BAD) : Programme ”Coding for Jobs”.\\
\vspace{0.5em}
\vspace{0.5em}
\subsection{15.7.3 3. Modèle de Franchise (Scale Africa)}
Dès l’Année 3, le modèle ”RBK in a Box” (LMS + Programme + Brand) sera proposé en franchise à des hubs technologiques au Sénégal\\
et en Côte d’Ivoire.\\
3\\
Web3/Blockchain — Soulbound Token non transférable, attaché à une identité ou un compte pour prouver un statut, une attestation ou une réussite.\\
\vspace{0.5em}
\vspace{0.5em}
\begin{confidential}CONFIDENTIEL — PROJET RBK 2.0                                           Page 139 sur 316\end{confidential}
RBK 2.0 — Whitepaper Architecte Web3                                                                           Chapitre 15. Business Plan\\
\vspace{0.5em}
\vspace{0.5em}
• Modèle : Revenue Share (20\textbackslash{}% du CA Franchise).\\
\vspace{0.5em}
• Apport : RBK fournit la plateforme et la certification SBT. Le partenaire gère le local et le sourcing.\\
\vspace{0.5em}
\vspace{0.5em}
Étudiants\\
Payants\\
\vspace{0.5em}
Fonds de                    Revenus\\
Garantie ISA               ISA Collectés\\
\vspace{0.5em}
\vspace{0.5em}
\vspace{0.5em}
\vspace{0.5em}
EL\\
Demandes\\
de\\
\vspace{0.5em}
\vspace{0.5em}
\vspace{0.5em}
\vspace{0.5em}
TI\\
Garantie\\
\vspace{0.5em}
\vspace{0.5em}
\vspace{0.5em}
\vspace{0.5em}
EN\\
(échec)\\
\vspace{0.5em}
FIG. 15.2 : Fonds de Garantie ISA (Cash Collateral)\\
\vspace{0.5em}
\vspace{0.5em}
\vspace{0.5em}
\vspace{0.5em}
D\\
FI\\
N\\
CO\\
\vspace{0.5em}
\vspace{0.5em}
\vspace{0.5em}
\vspace{0.5em}
\begin{confidential}CONFIDENTIEL — PROJET RBK 2.0                                  Page 140 sur 316\end{confidential}
RBK 2.0 — Whitepaper Architecte Web3                                                 Chapitre 15. Business Plan\\
\vspace{0.5em}
\vspace{0.5em}
B2C (70\textbackslash{}%)\\
Formation\\
Initiale\\
\vspace{0.5em}
\vspace{0.5em}
ISA (20\textbackslash{}%)\\
Partage\\
de Revenu\\
\vspace{0.5em}
\vspace{0.5em}
B2B (10\textbackslash{}%)         Revenue Streams\\
Formations         Total : 100\textbackslash{}%\\
\vspace{0.5em}
\vspace{0.5em}
\vspace{0.5em}
\vspace{0.5em}
EL\\
Sur-Mesure\\
\vspace{0.5em}
\vspace{0.5em}
\vspace{0.5em}
\vspace{0.5em}
TI\\
FIG. 15.3 : Modèle de Revenus Multi-Couches\\
\vspace{0.5em}
\vspace{0.5em}
\vspace{0.5em}
\vspace{0.5em}
D   EN\\
FI\\
N\\
CO\\
\vspace{0.5em}
\vspace{0.5em}
\vspace{0.5em}
\vspace{0.5em}
\begin{confidential}CONFIDENTIEL — PROJET RBK 2.0                       Page 141 sur 316\end{confidential}
\section{STRATÉGIE MARKETING \textbackslash{}& ACQUISITION RENFORCÉE}
\vspace{0.5em}
\vspace{0.5em}
\vspace{0.5em}
\vspace{0.5em}
EL\\
\subsection{16.1 Programme ”Building in Public”}
RBK 2.0 ne fait pas de publicité, elle produit de la preuve. Notre stratégie d’acquisition repose sur le ”Building in Public”. Nous\\
\vspace{0.5em}
\vspace{0.5em}
\vspace{0.5em}
\vspace{0.5em}
TI\\
documentons publiquement nos succès, nos échecs, nos audits et nos outils. Cette transparence radicale a trois objectifs : 1. Crédibilité :\\
\vspace{0.5em}
\vspace{0.5em}
\vspace{0.5em}
\vspace{0.5em}
EN\\
Montrer le niveau technique réel avant même l’inscription. 2. Confiance : Rassurer les candidats (et leurs parents) sur le sérieux de la\\
pédagogie. 3. Communauté : Attirer des mentors et des entreprises qui partagent nos valeurs.\\
\vspace{0.5em}
\vspace{0.5em}
\vspace{0.5em}
\vspace{0.5em}
D\\
FI\\
Les 3 Piliers de Contenu\\
\vspace{0.5em}
N\\
1. Technique (The Code) : Partage de snippets Rust, analyses de hacks récents, tutoriels Solana. Cible : Développeurs, CTOs.\\
CO\\
2. Pédagogie (The Journey) : Avant/Après des étudiants, rediffusion de Code Reviews, partage de ressources (Cheat Sheets). Cible :\\
Candidats.\\
\vspace{0.5em}
3. Success Stories (The Result) : Interviews d’Alumni, montants des bounties gagnés, projets lancés. Cible : Grand public.\\
\vspace{0.5em}
\vspace{0.5em}
\subsection{16.2 Assets visuels Sprint 0 (Maquettes à livrer)}
\vspace{0.5em}
\vspace{0.5em}
\vspace{0.5em}
\vspace{0.5em}
142\\
RBK 2.0 — Whitepaper Architecte Web3                                                            Chapitre 16. Stratégie marketing \textbackslash{}& acquisition\\
\vspace{0.5em}
\vspace{0.5em}
\vspace{0.5em}
\vspace{0.5em}
LinkedIn Carousel — ”Architecte vs Codeur” (6 slides).\\
Slide 1 : promesse RBK 2.0.\\
Slide 2 : comparaison RBK 1.0 / RBK 2.0.\\
Slide 3 : chiffres Indicateurs 2025.\\
Slide 4 : témoignage alumni.\\
Slide 5 : call-to-action ”Genesis Pool”.\\
Slide 6 : coordonnées RBK Discord.\\
Statut actuel : story-board validé, design final en production (cf. Chapitre 23).\\
\vspace{0.5em}
FIG. 16.1 : Maquette LinkedIn (placeholder visuel) — livrable D-21 avant lancement\\
\vspace{0.5em}
\vspace{0.5em}
\vspace{0.5em}
\vspace{0.5em}
EL\\
TI\\
Landing Page Hero — RBK Web3 Studio.\\
\vspace{0.5em}
\vspace{0.5em}
\vspace{0.5em}
\vspace{0.5em}
EN\\
Hero : visuel Validator + claim ”De Codeur à Architecte Web3”.\\
Sections : preuves (NFT Skills), pipeline employeurs, simulateur ROI.\\
\vspace{0.5em}
\vspace{0.5em}
\vspace{0.5em}
\vspace{0.5em}
D\\
Bloc CTA : formulaire double (Genesis Pool / Corporate Seat).\\
\vspace{0.5em}
\vspace{0.5em}
\vspace{0.5em}
\vspace{0.5em}
FI\\
Statut actuel : wireframe Figma prêt, intégration prévue Sprint 0 semaine 3.\\
\vspace{0.5em}
\vspace{0.5em}
N\\
FIG. 16.2 : Wireframe landing page (placeholder) — référentiel product marketing\\
CO\\
Script vidéo 60s — ”RBK 2.0 in Action”.\\
Hook (0-10s) : question ”Où se forment les architectes Web3 tunisiens ?”\\
Proof (10-35s) : extraits audits, dashboards, mentors.\\
CTA (35-60s) : ouverture Genesis Cohort Alpha.\\
Statut actuel : storyboard validé, tournage planifié Sprint 0 semaine 5.\\
\vspace{0.5em}
FIG. 16.3 : Storyboard vidéo (placeholder) — diffusion multi-canale\\
\vspace{0.5em}
\vspace{0.5em}
\vspace{0.5em}
\vspace{0.5em}
\begin{confidential}CONFIDENTIEL — PROJET RBK 2.0                                         Page 143 sur 316\end{confidential}
RBK 2.0 — Whitepaper Architecte Web3                                                        Chapitre 16. Stratégie marketing \textbackslash{}& acquisition\\
\vspace{0.5em}
\vspace{0.5em}
\vspace{0.5em}
\vspace{0.5em}
Contenu           Trust\\
\vspace{0.5em}
\vspace{0.5em}
\vspace{0.5em}
\vspace{0.5em}
Résultats      Candidats\\
\vspace{0.5em}
\vspace{0.5em}
FIG. 16.4 : Flywheel Building in Public\\
\vspace{0.5em}
\vspace{0.5em}
\vspace{0.5em}
\vspace{0.5em}
EL\\
TI\\
\begin{center}\begin{tabular}TAB. 16.1 : Calendrier Éditorial Type (Cycle 12 Semaines)\end{tabular}\end{center}
\vspace{0.5em}
\vspace{0.5em}
\vspace{0.5em}
\vspace{0.5em}
EN\\
Semaine Thème               Canal            KPI Cible\\
\vspace{0.5em}
\vspace{0.5em}
\vspace{0.5em}
\vspace{0.5em}
D\\
S1-S4   Rust Tips \textbackslash{}& Tricks Twitter/X         10k Impressions\\
S5-S8   Démo Projets Élèves YouTube/LinkedIn 50 Leads (Inscrits Webinar)\\
\vspace{0.5em}
\vspace{0.5em}
\vspace{0.5em}
\vspace{0.5em}
FI\\
S9-S12  Audit \textbackslash{}& Sécurité    Blog/Medium      5 Partenariats Entreprise\\
\vspace{0.5em}
N\\
CO\\
\vspace{0.5em}
\vspace{0.5em}
\vspace{0.5em}
\vspace{0.5em}
\begin{confidential}CONFIDENTIEL — PROJET RBK 2.0                               Page 144 sur 316\end{confidential}
RBK 2.0 — Whitepaper Architecte Web3                                                                                    Chapitre 16. Stratégie marketing \textbackslash{}& acquisition\\
\vspace{0.5em}
\vspace{0.5em}
\begin{center}\begin{tabular}TAB. 16.3 : Roadmap éditoriale Sprint 0 (90 jours détaillés)\end{tabular}\end{center}
\vspace{0.5em}
Semaine    Livrable              Owner principal      KPI de validation\\
S1         Article    ”Pourquoi CEO \textbackslash{}& Content Lead    3k vues blog, 20 partages X\\
RBK 2.0” + thread\\
X\\
S2         Live Discord ”Ask Me Head of Program \textbackslash{}& Ca- 80 participants, 30 questions qualifiées\\
Anything”             reer Studio\\
S3         Landing page MVP Growth Squad              Taux de conversion ≥ 5\textbackslash{}%\\
mise en ligne\\
S4         Carousel     LinkedIn Design Studio        150 leads MQL, 5 commentaires partenaires\\
(maquette ci-dessus)\\
S5         Vidéo 60s (story- Production Squad         Script validé, plan de tournage verrouillé\\
board finalisé)\\
\vspace{0.5em}
\vspace{0.5em}
\vspace{0.5em}
\vspace{0.5em}
EL\\
S6         Webinar ”Security Security Lead \textbackslash{}& Nexus 100 inscrits, 40 restream YouTube\\
Sprint”               Réussite\\
S7         Campagne         paid Growth Squad         CAC ≤ objectif prudent (Annexe F)\\
\vspace{0.5em}
\vspace{0.5em}
\vspace{0.5em}
\vspace{0.5em}
TI\\
”Builders”\\
S8         Témoignages alumni Career Studio           3 témoignages finalisés, satisfaction ≥ 4.5/5\\
\vspace{0.5em}
\vspace{0.5em}
\vspace{0.5em}
\vspace{0.5em}
EN\\
(capsules vidéo)\\
S9         Rapport KPI public Ops \textbackslash{}& Data              Dashboard publié, 200 vues uniques\\
(GitHub)\\
\vspace{0.5em}
\vspace{0.5em}
\vspace{0.5em}
\vspace{0.5em}
D\\
S10        Étude de cas DePIN Product Track           3 téléchargements corporate, 1 demande rendez-vous\\
\vspace{0.5em}
\vspace{0.5em}
\vspace{0.5em}
\vspace{0.5em}
FI\\
(PDF)\\
S11        Sprint bounty com- Career Studio           30 participations, 10 bounties clos\\
\vspace{0.5em}
\vspace{0.5em}
\vspace{0.5em}
N\\
munautaire\\
S12        Demo Day simulé Nexus Réussite             5 employeurs présents, 2 LOI signées\\
CO\\
(live)\\
\vspace{0.5em}
\vspace{0.5em}
\vspace{0.5em}
\vspace{0.5em}
\subsection{16.3 Simulateur de ROI Interactif}
Pour contrer l’objection du prix, nous proposons un outil permettant de comparer 3 scénarios d’investissement. L’objectif est de montrer\\
que le risque est nul si on commence petit et de rendre le ROI1 explicite dès le premier contact.\\
\vspace{0.5em}
Les 3 Options du Simulateur\\
1\\
Finance/Juridique/Compliance — Return on Investment (ROI) : rapport entre gain généré et investissement engagé.\\
\vspace{0.5em}
\vspace{0.5em}
\vspace{0.5em}
\begin{confidential}CONFIDENTIEL — PROJET RBK 2.0                                              Page 145 sur 316\end{confidential}
RBK 2.0 — Whitepaper Architecte Web3                                                              Chapitre 16. Stratégie marketing \textbackslash{}& acquisition\\
\vspace{0.5em}
\vspace{0.5em}
1. Option Genesis (Découverte) : Accès à la Genesis Pool (4 semaines) pour valider son appétence technique avec un investissement\\
initial limité et 100\textbackslash{}% flexible.\\
\vspace{0.5em}
2. Option Explorer (Immersion) : Parcours Genesis →Explorer (4 + 12 semaines) financé par mensualités alignées sur la scolarité\\
(tuition annuelle 15 900 TND).\\
\vspace{0.5em}
3. Option Nexus Launch (Complet) : Parcours intégral 48 semaines (Genesis Pool 4 + Explorer 12 + Bâtisseur 16 + Architecte 16)\\
suivi du Launchpad (forfait 3 300 TND + partage 30\textbackslash{}%), idéal pour les profils entrepreneuriaux.\\
\vspace{0.5em}
Cas d’Usage : Le “Smart Start”\\
Profil : Étudiant hésitant avec side-job freelance.\\
Action : Rejoint la Genesis Pool, valide les Skill Mirror d’origine, signe ensuite le plan Explorer (paiement mensuel indexé\\
\vspace{0.5em}
\vspace{0.5em}
\vspace{0.5em}
\vspace{0.5em}
EL\\
sur 15 900 TND).\\
\vspace{0.5em}
\vspace{0.5em}
\vspace{0.5em}
\vspace{0.5em}
TI\\
Résultat : Remporte un bounty communautaire de 600\textbackslash{}$ et couvre ses mensualités Genesis + Explorer.\\
Décision : Intègre le Launchpad pour préparer son projet et accepte le forfait 3 300 TND financé par le bounty.\\
\vspace{0.5em}
\vspace{0.5em}
\vspace{0.5em}
\vspace{0.5em}
EN\\
ROI : Il autofinance la transition Genesis →Explorer et abaisse le risque perçu à zéro.\\
\vspace{0.5em}
\vspace{0.5em}
\vspace{0.5em}
\vspace{0.5em}
D\\
FI\\
\begin{center}\begin{tabular}TAB. 16.5 : ROI Comparatif par Option (Sortie Junior : 2 500 TND brut/mois)\end{tabular}\end{center}
\vspace{0.5em}
N\\
CO\\
Option            Cash-out Initial                                            Point d’Équilibre        Avantage Principal\\
Genesis Pool      Frais modulaires, sans engagement long terme                \textless{} 3 mois (revenus bounties)\\
Découverte sécurisée, Skill Mirror\\
“Genesis Access”\\
Explorer Track  Mensualités alignées sur 15 900 TND                       ≈ 6 mois (mission freelance) Portfolio opérationnel, accès Block\\
Check Explorer\\
Nexus Launchpad Tuition complète + 3 300 TND (financé par ISA ou projets) \textless{} 12 mois (contrat CDI/CTO) Co-accompagnement marché + par-\\
tage 30\textbackslash{}% win-win\\
\vspace{0.5em}
\vspace{0.5em}
\vspace{0.5em}
\vspace{0.5em}
\subsection{16.4 Stratégie Multi-Canaux}
\vspace{0.5em}
\begin{confidential}CONFIDENTIEL — PROJET RBK 2.0                                    Page 146 sur 316\end{confidential}
RBK 2.0 — Whitepaper Architecte Web3                                                             Chapitre 16. Stratégie marketing \textbackslash{}& acquisition\\
\vspace{0.5em}
\vspace{0.5em}
\vspace{0.5em}
Nous ne cherchons pas à être partout, mais à dominer 5 canaux spécifiques où se trouve notre cible ”Elite”.\\
\vspace{0.5em}
Les 5 Canaux Prioritaires\\
\vspace{0.5em}
1. LinkedIn (La Vitrine) : Pour les parents, les recruteurs et les partenariats corporatifs. Cadence : 2 posts/semaine.\\
\vspace{0.5em}
2. X / Twitter (L’Arène) : Pour la crédibilité technique crypto, les news Rust, et l’engagement communautaire. Cadence : Quotidien.\\
\vspace{0.5em}
3. YouTube (La Preuve) : Replays de workshops, Démos de Capstones, Témoignages. Cadence : 2 vidéos/mois.\\
\vspace{0.5em}
4. GitHub (Le CV) : C’est notre canal d’acquisition ”silencieux”. Des repos propres et étoilés attirent les curieux techniques.\\
\vspace{0.5em}
\vspace{0.5em}
\vspace{0.5em}
\vspace{0.5em}
EL\\
5. Discord (Le Salon) : Conversion des leads chauds, support, Q\textbackslash{}&A avant inscription.\\
\vspace{0.5em}
\vspace{0.5em}
\vspace{0.5em}
\vspace{0.5em}
TI\\
EN\\
\subsection{16.5 Budget d’Acquisition \textbackslash{}& CAC Cible}
Le plan marketing respecte la doctrine Sprint 0 : investir 75 000 TND sur 12 semaines pour installer le funnel Discord et documenter\\
\vspace{0.5em}
\vspace{0.5em}
\vspace{0.5em}
\vspace{0.5em}
D\\
chaque victoire. Les années suivantes appliquent un run rate maîtrisé (30 000 TND puis 45 000 TND) afin de rester sous le plafond CAC ≤\\
\vspace{0.5em}
\vspace{0.5em}
\vspace{0.5em}
\vspace{0.5em}
FI\\
500 TND.\\
\vspace{0.5em}
N\\
\begin{center}\begin{tabular}TAB. 16.7 : Allocation budgétaire et objectifs CAC\end{tabular}\end{center}
CO\\
Phase Budget        Focales marketing                                                      Garde-fou CAC\\
Sprint 0 75 000 TND Paid social ciblé, production ”Building in Public”, onboarding Discord 800 TND\\
Année 1 30 000 TND Webinaires mensuels, nurturing e-mail, nurturing alumni                 500 TND\\
Année 2 45 000 TND Programme ambassadeurs + partenariats HR, événements hybrides           300 TND\\
\vspace{0.5em}
\vspace{0.5em}
\vspace{0.5em}
\vspace{0.5em}
\begin{confidential}CONFIDENTIEL — PROJET RBK 2.0                                Page 147 sur 316\end{confidential}
RBK 2.0 — Whitepaper Architecte Web3                                                           Chapitre 16. Stratégie marketing \textbackslash{}& acquisition\\
\vspace{0.5em}
\vspace{0.5em}
Gouvernance budgétaire Chaque dépense est tracée dans le cockpit marketing (Google Sheet de référence), relié aux playbooks détaillés\\
dans le document Marketing Strategy - Nexus x RBK 2024. Les arbitrages sont pilotés mensuellement en comité Funnel avec un reporting\\
CAC vs. placement signé par Nexus Réussite.\\
\vspace{0.5em}
\vspace{0.5em}
\subsection{16.5.1 Discord : Le cœur communautaire de la pré-sélection}
Discord concentre 80\textbackslash{}% de notre énergie Sprint 0. Il assure la qualification des leads, la montée en confiance et la détection des talents\\
”builders”.\\
\vspace{0.5em}
Cadence opérationnelle\\
\vspace{0.5em}
\vspace{0.5em}
\vspace{0.5em}
\vspace{0.5em}
EL\\
• Weekly Office Hours animées par Nexus Réussite pour répondre aux questions admissions et expliquer la trajectoire Genesis →\\
Explorer.\\
\vspace{0.5em}
\vspace{0.5em}
\vspace{0.5em}
\vspace{0.5em}
TI\\
EN\\
• Challenge ”Ship or Die” bi-hebdomadaire : mini-sprints 48h avec jury RBK, publiés dans le canal \textbackslash{}#build-in-public.\\
\vspace{0.5em}
• Canal \textbackslash{}#ask-mentor modéré par les mentors Launchpad afin de créer le réflexe pair → pair dès la pré-inscription.\\
\vspace{0.5em}
\vspace{0.5em}
\vspace{0.5em}
\vspace{0.5em}
D\\
FI\\
• Onboarding automatisé (Notion Checklist + Skill Mirror) pour orienter immédiatement vers les ressources Genesis Pool.\\
\vspace{0.5em}
KPIs suivis                                       N\\
CO\\
• Taux de conversion lead → candidature complète (≥ 40\textbackslash{}%).\\
\vspace{0.5em}
• Nombre de builders actifs (au moins 3 commits ou prototypes/mois).\\
\vspace{0.5em}
• Temps moyen de réponse mentor \textless{} 12 heures.\\
\vspace{0.5em}
Les apprentissages issus de Discord sont intégrés dans les campagnes LinkedIn/X la semaine suivante afin de garder le storytelling\\
aligné sur les objections réelles du terrain.\\
\vspace{0.5em}
\vspace{0.5em}
\vspace{0.5em}
\vspace{0.5em}
\begin{confidential}CONFIDENTIEL — PROJET RBK 2.0                                 Page 148 sur 316\end{confidential}
RBK 2.0 — Whitepaper Architecte Web3                                                         Chapitre 16. Stratégie marketing \textbackslash{}& acquisition\\
\vspace{0.5em}
\vspace{0.5em}
Awareness (X/LinkedIn)\\
\vspace{0.5em}
Considération (YouTube/Webinar)\\
\vspace{0.5em}
Intention (Discord/Simulateur)\\
\vspace{0.5em}
Action (Candidature)\\
\vspace{0.5em}
FIG. 16.5 : Funnel d’Acquisition Simplifié\\
\vspace{0.5em}
\vspace{0.5em}
\vspace{0.5em}
\subsection{16.6 Programme de Référence \textbackslash{}& Bounties}
\vspace{0.5em}
\vspace{0.5em}
\vspace{0.5em}
\vspace{0.5em}
EL\\
Le ”Word of Mouth” est notre canal le plus rentable (CAC ≈ 0). Nous l’industrialisons.\\
\vspace{0.5em}
\vspace{0.5em}
\vspace{0.5em}
\vspace{0.5em}
TI\\
EN\\
Système de Parrainage (Referral)       Tout Alumni ou Étudiant validé peut parrainer un candidat.\\
\vspace{0.5em}
• Pour le Parrain : 500 TND (Cash ou déduction ISA) versés APRÈS la validation de la Période d’Essai du filleul (Anti-fraude).\\
\vspace{0.5em}
\vspace{0.5em}
\vspace{0.5em}
\vspace{0.5em}
D\\
FI\\
• Pour le Filleul : 5\textbackslash{}% de réduction immédiate sur les frais Upfront.\\
\vspace{0.5em}
N\\
CO\\
Programme ”Bug Bounties” Pédagogiques           Nous payons (en crédits ou token réputation) pour l’amélioration du cursus.\\
\vspace{0.5em}
• Typo majeure dans le cours : 10 pts.\\
\vspace{0.5em}
• Optimisation d’un exercice de code : 50 pts.\\
\vspace{0.5em}
• Fix de sécurité sur l’infra école : 500 TND.\\
\vspace{0.5em}
Cela crée une culture de ”Contribution” dès le premier jour.\\
\vspace{0.5em}
\vspace{0.5em}
\vspace{0.5em}
\vspace{0.5em}
\begin{confidential}CONFIDENTIEL — PROJET RBK 2.0                                  Page 149 sur 316\end{confidential}
RBK 2.0 — Whitepaper Architecte Web3                                                           Chapitre 16. Stratégie marketing \textbackslash{}& acquisition\\
\vspace{0.5em}
\vspace{0.5em}
\begin{center}\begin{tabular}TAB. 16.9 : Catalogue des Incentives\end{tabular}\end{center}
\vspace{0.5em}
\vspace{0.5em}
Mécanisme Bénéficiaire Récompense      Condition Anti-Fraude\\
Parrainage  Alumni      500 TND        Filleul valide le SPRINT 1 (pas juste inscrit).\\
Ambassadeur Influenceur 10\textbackslash{}% Commission Lien tracké + KYC obligatoire.\\
Bounty Code Étudiant    Goodies / Cash Pull Request validée par le Lead Tech.\\
\vspace{0.5em}
\vspace{0.5em}
\vspace{0.5em}
\vspace{0.5em}
EL\\
TI\\
D    EN\\
FI\\
N\\
CO\\
\vspace{0.5em}
\vspace{0.5em}
\vspace{0.5em}
\vspace{0.5em}
\begin{confidential}CONFIDENTIEL — PROJET RBK 2.0                                Page 150 sur 316\end{confidential}
\section{Á ANALYSE DES RISQUES \textbackslash{}& MODÈLE DE RÉSILIENCE}
\vspace{0.5em}
\vspace{0.5em}
RBK 2.0 opère à l’intersection de deux secteurs volatils : l’éducation technologique et les actifs numériques. Cette position exige une\\
\vspace{0.5em}
\vspace{0.5em}
\vspace{0.5em}
\vspace{0.5em}
EL\\
gestion des risques de niveau institutionnel.\\
\vspace{0.5em}
\vspace{0.5em}
\vspace{0.5em}
\vspace{0.5em}
TI\\
EN\\
\subsection{17.1 Risques Réglementaires et Conformité}
La pérennité de RBK repose sur une veille juridique proactive, particulièrement en Tunisie (siège opérationnel) et en Europe (marché\\
\vspace{0.5em}
\vspace{0.5em}
\vspace{0.5em}
\vspace{0.5em}
D\\
cible).\\
\vspace{0.5em}
\vspace{0.5em}
\vspace{0.5em}
\vspace{0.5em}
FI\\
\subsection{17.1.1 Loi des Changes et Crypto-Actifs (Tunisie)}
N\\
CO\\
Risque : La détention de crypto-actifs reste une zone grise. Une interdiction stricte pourrait bloquer les paiements en Stablecoins.\\
Mitigation :\\
\vspace{0.5em}
• Structure Off-shore : RBK facture via une entité non-résidente (ou partenaire) pour les flux internationaux, en conformité totale\\
avec le code des changes.\\
\vspace{0.5em}
• Modèle Opérationnel Détaillé : Voir Chapitre 16 pour l’application des Scénarios A (Export), B (Offshore) ou C (Portage) selon la\\
maturité de l’apprenant.\\
\vspace{0.5em}
• Flux Fiat Prioritaire : 100\textbackslash{}% des frais de scolarité locaux sont encaissés en TND via virement bancaire classique. La crypto n’est\\
qu’un rail technologique optionnel pour les bourses étrangères.\\
\vspace{0.5em}
\vspace{0.5em}
151\\
RBK 2.0 — Whitepaper Architecte Web3                                        Chapitre 17. ANALYSE DES RISQUES \textbackslash{}& MODÈLE DE RÉSILIENCE\\
\vspace{0.5em}
\vspace{0.5em}
• Lobbying Actif : RBK participe aux groupes de travail de la BCT (Banque Centrale) pour encadrer le statut de ”Service Exporter”\\
blockchain.\\
\vspace{0.5em}
\vspace{0.5em}
\subsection{17.1.2 GDPR et Données Étudiantes On-Chain}
Risque : Les SBT (Soulbound Tokens) sont immuables. Si des données personnelles y sont inscrites, le ”Droit à l’oubli” est impossible.\\
Mitigation :\\
\vspace{0.5em}
• Architecture Privacy-First : Aucun nom, email ou IP n’est stocké sur la blockchain. Le SBT contient uniquement un Hash Cryptogra-\\
phique (ex : Keccak256(Diplome\textbackslash{}_PDF)).\\
\vspace{0.5em}
\vspace{0.5em}
\vspace{0.5em}
\vspace{0.5em}
EL\\
• Consentement Explicite : L’étudiant signe une décharge explicite pour le minting de ses résultats.\\
\vspace{0.5em}
• Droit à la Révocation : Le contrat intelligent permet à l’admin (sur demande de l’étudiant) de ”brûler” un token, rompant le lien\\
\vspace{0.5em}
\vspace{0.5em}
\vspace{0.5em}
\vspace{0.5em}
TI\\
public.\\
\vspace{0.5em}
\vspace{0.5em}
\vspace{0.5em}
\vspace{0.5em}
D    EN\\
\subsection{17.1.3 Cadre Légal des ISA (Income Share Agreements)}
\vspace{0.5em}
\vspace{0.5em}
\vspace{0.5em}
\vspace{0.5em}
FI\\
Risque : Requalification du contrat ISA en crédit à la consommation déguisé ou clause abusive. Mitigation :\\
\vspace{0.5em}
N\\
• Juridiction Compétente : Contrats régis par le droit commercial (prestation de service avec paiement différé) et non le droit de la\\
CO\\
consommation.\\
\vspace{0.5em}
• Clauses Protectrices : Plafond de remboursement (Cap) strict et durée limitée pour éviter toute notion de ”servitude”.\\
\vspace{0.5em}
• Enforceability : Partenariat avec des cabinets de recouvrement locaux en Tunisie, Maroc et Côte d’Ivoire.\\
\vspace{0.5em}
\vspace{0.5em}
\subsection{17.2 Matrice de Risques Dynamique}
Nous évaluons la résilience du modèle selon trois scénarios de marché.\\
\vspace{0.5em}
\vspace{0.5em}
\vspace{0.5em}
\vspace{0.5em}
\begin{confidential}CONFIDENTIEL — PROJET RBK 2.0                               Page 152 sur 316\end{confidential}
RBK 2.0 — Whitepaper Architecte Web3                                        Chapitre 17. ANALYSE DES RISQUES \textbackslash{}& MODÈLE DE RÉSILIENCE\\
\vspace{0.5em}
\vspace{0.5em}
\begin{center}\begin{tabular}TAB. 17.1 : Matrice des Risques — Référentiel Commun (RBK – Nexus Réussite)\end{tabular}\end{center}
\vspace{0.5em}
\vspace{0.5em}
Risque                                               Impact Probabilité Score Mitigation\\
Technique\\
Mauvaise estimation de la complexité Solana           Élevé     Moyen         12   Approche itérative, prototypage, veille technique ac-\\
tive\\
Mismatch d’API entre versions                        Modéré      Élevé        15   Ancrage sur versions LTS, tests d’intégration, CI auto-\\
matique\\
Atteinte de performance (QoS)                         Élevé     Faible        6    Benchmarks, profiling, optimisation continue\\
Pédagogique\\
\vspace{0.5em}
\vspace{0.5em}
\vspace{0.5em}
\vspace{0.5em}
EL\\
Taux de chute élevé (\textgreater{}40\textbackslash{}%)                            Élevé     Faible        6    Suivi individualisé, soutien psychologique, adapta-\\
tion du rythme\\
\vspace{0.5em}
\vspace{0.5em}
\vspace{0.5em}
\vspace{0.5em}
TI\\
Non-alignement avec le marché                         Élevé     Modéré        12   Veille métier, feedbacks industriels, itération conti-\\
\vspace{0.5em}
\vspace{0.5em}
\vspace{0.5em}
\vspace{0.5em}
EN\\
nue\\
Manque de diversité pédagogique                      Modéré     Modéré        9    Mix pédagogique (pairs, projets, mentors), rotation\\
\vspace{0.5em}
\vspace{0.5em}
\vspace{0.5em}
\vspace{0.5em}
D\\
du staff\\
\vspace{0.5em}
\vspace{0.5em}
\vspace{0.5em}
\vspace{0.5em}
FI\\
Commercial\\
Absence de débouchés (placement \textless{}70\textbackslash{}%)                 Élevé     Faible        6    Partenariats entreprises, réseau alumni, événements\\
N                                  networking\\
CO\\
Mauvaise perception du marché                        Modéré     Faible        3    Communication proactive, résultats transparents,\\
preuves sociales\\
Financier\\
Retard de la mise en place ISA                        Élevé     Faible        3    Processus juridique en amont, test de faisabilité\\
Insolvabilité d’un grand partenaire                   Élevé   Très faible     1    Diversification des partenaires, contrats avec garantie\\
Réglementaire\\
Changements législatifs (crypto)                     Modéré     Faible        2    Veille juridique, conseils externes, conformité proac-\\
tive\\
Problèmes de conformité ISA                          Modéré     Faible        2    Suivi juridique, audit externe, mise à jour régulière\\
\vspace{0.5em}
\vspace{0.5em}
\vspace{0.5em}
\begin{confidential}CONFIDENTIEL — PROJET RBK 2.0                              Page 153 sur 316\end{confidential}
RBK 2.0 — Whitepaper Architecte Web3                                        Chapitre 17. ANALYSE DES RISQUES \textbackslash{}& MODÈLE DE RÉSILIENCE\\
\vspace{0.5em}
\vspace{0.5em}
\begin{center}\begin{tabular}TAB. 17.3 : Nouvelle Matrice des Risques (Version 2026)\end{tabular}\end{center}
\vspace{0.5em}
\vspace{0.5em}
Risque Probabilité Impact Mesures d’atténuation Responsable\\
R-XX   Stratégique 3      5                     15\\
\vspace{0.5em}
\vspace{0.5em}
\vspace{0.5em}
\vspace{0.5em}
\subsection{17.3 Plan de Continuité Pédagogique (PCP)}
Le PCP garantit la poursuite des apprentissages malgré une indisponibilité de la plateforme, de la blockchain cible ou des ressources\\
humaines critiques. Il est testé chaque trimestre sous la supervision conjointe du Responsable Bien-être et du CTO Nexus.\\
\vspace{0.5em}
\vspace{0.5em}
\vspace{0.5em}
\vspace{0.5em}
EL\\
TI\\
\subsection{17.3.1 Kit de survie pédagogique}
\vspace{0.5em}
\vspace{0.5em}
\vspace{0.5em}
\vspace{0.5em}
EN\\
• Contenus Offline : miroirs Git des Golden Templates, labs exportés en PDF/Markdown, vidéos masterclass stockées sur NAS RBK.\\
\vspace{0.5em}
• Simulateurs Locaux : images Docker prêtes à l’emploi (Solana test validator, Foundry/Anvil, PostgreSQL seed), scripts de provisioning\\
\vspace{0.5em}
\vspace{0.5em}
\vspace{0.5em}
\vspace{0.5em}
D\\
automatisés.\\
\vspace{0.5em}
\vspace{0.5em}
\vspace{0.5em}
\vspace{0.5em}
FI\\
• Playbooks de Substitution : plans jour par jour pour basculer vers des activités ”Rust Core” ou ”Security Drills” sans dépendance\\
on-chain.                                      N\\
CO\\
• Canaux de Secours : serveur Discord privé Continuity, numéros d’astreinte, accès à la FAQ offline (Notion export).\\
\vspace{0.5em}
\vspace{0.5em}
\vspace{0.5em}
\vspace{0.5em}
\begin{confidential}CONFIDENTIEL — PROJET RBK 2.0                                Page 154 sur 316\end{confidential}
RBK 2.0 — Whitepaper Architecte Web3                                            Chapitre 17. ANALYSE DES RISQUES \textbackslash{}& MODÈLE DE RÉSILIENCE\\
\vspace{0.5em}
\vspace{0.5em}
\subsection{17.3.2 Tableau de réponse rapide}
\vspace{0.5em}
\begin{center}\begin{tabular}TAB. 17.5 : Déclencheurs et plans de continuité\end{tabular}\end{center}
\vspace{0.5em}
\vspace{0.5em}
Scénario           Déclencheur         Actions H+4 → H+24                                                                            Responsable\\
Outage Venture     SLA \textless{} 99\textbackslash{}%,          Bascule vers LMS backup (Moodle RBK), diffusion labs offline, collecte livrables via Git mir- Nexus CTO\\
Engine             plateforme          ror, information cohorte                                                                      (R), RBK Ops\\
\vspace{0.5em}
\vspace{0.5em}
\vspace{0.5em}
\vspace{0.5em}
EL\\
indisponible \textgreater{} 2h                                                                                                 (A)\\
Panne blockchain Solana/RPC            Suspension travaux on-chain, lancement modules Rust Fundamentals, distribution simula- Lead Track\\
majeure            injoignable,        teurs locaux, briefing sécurité                                                               Solana (R),\\
\vspace{0.5em}
\vspace{0.5em}
\vspace{0.5em}
\vspace{0.5em}
TI\\
incident sécurité                                                                                                 Responsable\\
critique                                                                                                          Bien-être (C)\\
\vspace{0.5em}
\vspace{0.5em}
\vspace{0.5em}
\vspace{0.5em}
EN\\
Rupture mentor clé Indisponibilité \textgreater{}   Activation Mentor-in-a-Box, rotation bench alumni, replanification séances, communication Nexus People\\
48h d’un mentor     cohorte                                                                                       Lead (R), RBK\\
Core                                                                                                              Program\\
\vspace{0.5em}
\vspace{0.5em}
\vspace{0.5em}
\vspace{0.5em}
D\\
Director (A)\\
Fermeture campus Crise sanitaire,      Passage full-remote, allocation matériel (kits 4G, laptops), mise en place plages Wellbeing RBK COO (A),\\
\vspace{0.5em}
\vspace{0.5em}
\vspace{0.5em}
FI\\
sinistre locaux       renforcées                                                                                    Nexus\\
Responsable\\
N\\
Bien-être (R)\\
O\\
C\\
\vspace{0.5em}
\subsection{17.3.3 Protocole de tests et amélioration}
• Exercices de simulation trimestriels (table-top) documentés et revus en comité qualité.\\
\vspace{0.5em}
• Post-mortem systématique après chaque activation du PCP, intégrant feedbacks apprenants et mentors.\\
\vspace{0.5em}
• Mise à jour des indicateurs de fatigue et des seuils d’alerte dans le cockpit central.\\
\vspace{0.5em}
• Audit externe annuel pour valider la conformité du PCP aux normes ISO 22301 (Business Continuity) adaptées.\\
\vspace{0.5em}
\vspace{0.5em}
\begin{confidential}CONFIDENTIEL — PROJET RBK 2.0                                  Page 155 sur 316\end{confidential}
RBK 2.0 — Whitepaper Architecte Web3                                         Chapitre 17. ANALYSE DES RISQUES \textbackslash{}& MODÈLE DE RÉSILIENCE\\
\vspace{0.5em}
\vspace{0.5em}
\vspace{0.5em}
\subsection{17.4 Plan de Réponse aux Incidents Crypto (”Black Swan”)}
Face à la volatilité intrinsèque du secteur, nous déployons un plan de continuité ”Grade Militaire”.\\
\vspace{0.5em}
\vspace{0.5em}
\subsection{17.4.1 Scénario A : Effondrement de l’Écosystème Solana}
Déclencheur : Panne du réseau \textgreater{} 72 h ou chute du token SOL \textless{} 10\textbackslash{}$.\\
\vspace{0.5em}
1. Immédiat (H+1) : Communication de crise rassurance (”Nous formons des ingénieurs, pas des spéculateurs”).\\
\vspace{0.5em}
2. Pivot Pédagogique (H+24) : Bascule des modules ”Solana Specific” vers ”Rust Générique” (valable pour Polkadot, Near, ou Backend\\
Web2).\\
\vspace{0.5em}
\vspace{0.5em}
\vspace{0.5em}
\vspace{0.5em}
EL\\
3. Trésorerie : Conversion automatique de tous les assets crypto en Fiat/Stablecoin dès que la volatilité dépasse un seuil d’alerte\\
\vspace{0.5em}
\vspace{0.5em}
\vspace{0.5em}
\vspace{0.5em}
TI\\
(Stop-Loss).\\
\vspace{0.5em}
\vspace{0.5em}
\vspace{0.5em}
\vspace{0.5em}
D    EN\\
\subsection{17.4.2 Scénario B : Hack d’un Bridge / Protocole Partenaire}
\vspace{0.5em}
\vspace{0.5em}
\vspace{0.5em}
\vspace{0.5em}
FI\\
Déclencheur : Un outil utilisé dans le cours (ex : Wormhole) est compromis. Ce scénario couvre la compromission d’un Bridge tiers.\\
\vspace{0.5em}
\vspace{0.5em}
N\\
1. Arrêt des Nodes : Les étudiants déconnectent leurs environnements de dev locaux.\\
CO\\
2. Learning Moment : L’incident devient un cas d’étude ”Live”. Analyse on-chain du hack en cours de sécurité.\\
\vspace{0.5em}
3. Fonds de Sécurité : Si des bourses étudiantes étaient bloquées, le Fonds de Garantie ISA avance la liquidité.\\
\vspace{0.5em}
\vspace{0.5em}
\subsection{17.5 Tableau de Bord des Risques Critiques}
\vspace{0.5em}
\vspace{0.5em}
\vspace{0.5em}
\vspace{0.5em}
\begin{confidential}CONFIDENTIEL — PROJET RBK 2.0                                Page 156 sur 316\end{confidential}
RBK 2.0 — Whitepaper Architecte Web3                                        Chapitre 17. ANALYSE DES RISQUES \textbackslash{}& MODÈLE DE RÉSILIENCE\\
\vspace{0.5em}
\vspace{0.5em}
\vspace{0.5em}
\vspace{0.5em}
\begin{center}\begin{tabular}TAB. 17.7 : Top 5 Risques et Mitigations (2026)\end{tabular}\end{center}
\vspace{0.5em}
\vspace{0.5em}
Risque              Prob. Imp. Plan de Mitigation\\
Défaut Paiement ISA 3/5 5/5 Sélection stricte (Top 30\textbackslash{}%), Fonds de Garantie ISA (120k TND), Assurance.\\
Obsolescence Tech    4/5 3/5 Comité Pédagogique trimestriel, Track Agnostique (Focus Fundamentals).\\
Fuite des Mentors    2/5 4/5 Programme ”Train the Trainer”, Satisfaction Index, Bonus Performance.\\
Cyber-attaque École 3/5 4/5 Infra isolée, 2FA Hardware (Yubikey) pour staff, Audit annuel.\\
Perte Réputation     2/5 5/5 Transparence totale (Building in Public), Charte Éthique stricte.\\
\vspace{0.5em}
\vspace{0.5em}
\vspace{0.5em}
\vspace{0.5em}
EL\\
TI\\
D    EN\\
FI\\
N\\
CO\\
\vspace{0.5em}
\vspace{0.5em}
\vspace{0.5em}
\vspace{0.5em}
\begin{confidential}CONFIDENTIEL — PROJET RBK 2.0                               Page 157 sur 316\end{confidential}
\section{COMPLIANCE \textbackslash{}& RÉGULATION WEB3 – GUIDE PRATIQUE}
\vspace{0.5em}
\vspace{0.5em}
\vspace{0.5em}
\vspace{0.5em}
EL\\
Ce chapitre transforme la contrainte réglementaire en avantage compétitif. Dans le Web3, la conformité n’est pas un frein, mais une\\
\vspace{0.5em}
\vspace{0.5em}
\vspace{0.5em}
\vspace{0.5em}
TI\\
fonctionnalité architecturale.\\
\vspace{0.5em}
Ď Résumé Exécutif : La Conformité comme Avantage Compétitif\\
\vspace{0.5em}
\vspace{0.5em}
\vspace{0.5em}
\vspace{0.5em}
EN\\
Le Web3 n’est pas une zone de non-droit. RBK 2.0 forme des ingénieurs capables de naviguer dans la complexité réglementaire\\
(RGPDa , MiCA, ETE). Ce chapitre détaille les protocoles techniques pour garantir une conformité ”By Design” sans sacrifier la\\
\vspace{0.5em}
\vspace{0.5em}
\vspace{0.5em}
\vspace{0.5em}
D\\
décentralisation.\\
\vspace{0.5em}
\vspace{0.5em}
\vspace{0.5em}
FI\\
a\\
Finance/Juridique/Compliance — General Data Protection Regulation (GDPR / RGPD) : règlement encadrant collecte, traitement et droits des per-\\
sonnes en Europe.\\
N\\
Disclaimer Légal\\
O\\
\vspace{0.5em}
RBK est un programme d’ingénierie logicielle.\\
C\\
\vspace{0.5em}
\vspace{0.5em}
• Pas de Conseil Financier : Nous n’enseignons pas le trading, l’analyse technique ou l’investissement.\\
\vspace{0.5em}
• Neutralité Technologique : L’étude des protocoles (DeFi, DAO) est purement technique (smart contracts, sécurité).\\
\vspace{0.5em}
• Gains : Aucune promesse de gains passifs ou de rendements n’est faite aux apprenants.\\
\vspace{0.5em}
\vspace{0.5em}
\vspace{0.5em}
\vspace{0.5em}
158\\
RBK 2.0 — Whitepaper Architecte Web3                                                                 Chapitre 18. Guide Compliance Web3\\
\vspace{0.5em}
\vspace{0.5em}
\vspace{0.5em}
\subsection{18.1 Operating Model Compliant : Scénarios pour la Tunisie}
AVERTISSEMENT\\
Les schémas décrits ci-dessous sont des modèles types. Chaque apprenant, alumni ou startup issue de RBK doit impérativement\\
solliciter un expert-comptable et un avocat fiscaliste tunisien avant toute mise en œuvre. RBK et Nexus Réussite recommandent la\\
signature préalable d’un mandat avec un cabinet spécialisé pour valider la stratégie individuelle.\\
\vspace{0.5em}
Pour opérer légalement depuis la Tunisie tout en servant un marché mondial Web3, nous structurons l’activité selon trois scénarios\\
validés. Chaque scénario correspond à un stade de maturité de l’apprenant/entrepreneur.\\
\vspace{0.5em}
\vspace{0.5em}
\vspace{0.5em}
\vspace{0.5em}
EL\\
\subsection{18.1.1 Scénario A : Exportateur de Services Logiciels (Le Standard)}
\vspace{0.5em}
\vspace{0.5em}
\vspace{0.5em}
\vspace{0.5em}
TI\\
Cible : Freelance individuel ou petite équipe (1-3 pers).\\
\vspace{0.5em}
\vspace{0.5em}
\vspace{0.5em}
\vspace{0.5em}
EN\\
Dimension      Détails Opérationnels\\
\vspace{0.5em}
\vspace{0.5em}
\vspace{0.5em}
\vspace{0.5em}
D\\
Périmètre      Développement, Audit, Consulting Technique. Pas de détention d’actifs pour compte de tiers.\\
\vspace{0.5em}
\vspace{0.5em}
\vspace{0.5em}
\vspace{0.5em}
FI\\
Statut         Personne Physique (Patente Catégorie Services Informatiques) ou SUARL Exportatrice.\\
\vspace{0.5em}
\vspace{0.5em}
N\\
Flux Financier Client (Devises) → Facture (EUR/USD) → Virement SWIFT → Compte Pro Devises TND. Crypto acceptée uniquement via processeur\\
de paiement (Bitwage, Request Finance) convertissant instantanément.\\
CO\\
Risques        Volatilité de change (si délai de conversion), Blocage virement (Compliance bancaire).\\
Limite         Interdiction de faire du ”Custody” (Garde de fonds) ou de l’Exchange.\\
\vspace{0.5em}
\vspace{0.5em}
\vspace{0.5em}
\subsection{18.1.2 Scénario B : Filiale Offshore (Le Scale-Up)}
Cible : Startup Web3 lancée par des Alumni RBK.\\
\vspace{0.5em}
\vspace{0.5em}
\vspace{0.5em}
\vspace{0.5em}
\begin{confidential}CONFIDENTIEL — PROJET RBK 2.0                                  Page 159 sur 316\end{confidential}
RBK 2.0 — Whitepaper Architecte Web3                                                                             Chapitre 18. Guide Compliance Web3\\
\vspace{0.5em}
\vspace{0.5em}
\vspace{0.5em}
Dimension      Détails Opérationnels\\
Périmètre      Édition de logiciel SaaS/DApp, Tokenisation, Revenus protocolaires.\\
Structure      Holding (Estonie/Delaware/Dubaï) qui encaisse la Crypto + Filiale Tunisie (Cost Center) qui produit.\\
Flux Financier La Holding facture les clients, encaisse les tokens, gère la trésorerie. La Filiale refacture ses coûts (Salaires + Marge) à la Holding en\\
Fiat (EUR).\\
Risques        Prix de transfert (doit être justifié), Substance fiscale (Holding doit avoir une gestion réelle).\\
\vspace{0.5em}
\vspace{0.5em}
\vspace{0.5em}
\subsection{18.1.3 Scénario C : Freelance ”Portage Salarial” (Le Simple)}
Cible : Débutant ou mission courte durée.\\
\vspace{0.5em}
\vspace{0.5em}
\vspace{0.5em}
\vspace{0.5em}
EL\\
• Mecanisme : Utilisation d’un EoR (Employer of Record) comme Deel ou Remote.\\
\vspace{0.5em}
\vspace{0.5em}
\vspace{0.5em}
\vspace{0.5em}
TI\\
• Avantage : Zéro gestion administrative. Protection sociale complète (CNSS).\\
\vspace{0.5em}
\vspace{0.5em}
\vspace{0.5em}
\vspace{0.5em}
EN\\
• Inconvénient : Coût (Frais 5-10\textbackslash{}%) + Moins d’optimisation fiscale que le statut ETE.\\
\vspace{0.5em}
\vspace{0.5em}
\vspace{0.5em}
\vspace{0.5em}
D\\
FI\\
\subsection{18.2 Matrice des Risques de Conformité}
N\\
CO\\
\vspace{0.5em}
\vspace{0.5em}
\vspace{0.5em}
\vspace{0.5em}
\begin{confidential}CONFIDENTIEL — PROJET RBK 2.0                                      Page 160 sur 316\end{confidential}
RBK 2.0 — Whitepaper Architecte Web3                                                                         Chapitre 18. Guide Compliance Web3\\
\vspace{0.5em}
\vspace{0.5em}
\vspace{0.5em}
\vspace{0.5em}
\begin{center}\begin{tabular}TAB. 18.3 : Risques Compliance \textbackslash{}& Atténuation\end{tabular}\end{center}
\vspace{0.5em}
\vspace{0.5em}
Domaine       Risque Identifié                                                 Niveau Atténuation (Contrôle)\\
Légal         Qualification en ”Vente de titres financiers” (Security Token)    Élevé   Disclaimer strict, Utility\\
Token uniquement, Avis\\
juridique préalable.\\
Fiscal        Redressement pour revenus non déclarés (Crypto)                  Moyen    Conversion Fiat systématique,\\
Déclaration BCT trimestrielle.\\
\vspace{0.5em}
\vspace{0.5em}
\vspace{0.5em}
\vspace{0.5em}
EL\\
Réputation    Association à des projets ”Scam” ou ”Ponzi”                      Critique Due Diligence des partenaires,\\
Refus de tout projet\\
\vspace{0.5em}
\vspace{0.5em}
\vspace{0.5em}
\vspace{0.5em}
TI\\
High-Yield non audité.\\
\vspace{0.5em}
\vspace{0.5em}
\vspace{0.5em}
\vspace{0.5em}
EN\\
Technique     Sanctions internationales (Tornado Cash, OFAC)                    Élevé   Filtrage des adresses\\
(Chainalysis), Refus des\\
\vspace{0.5em}
\vspace{0.5em}
\vspace{0.5em}
\vspace{0.5em}
D\\
interactions mixeurs.\\
\vspace{0.5em}
\vspace{0.5em}
\vspace{0.5em}
\vspace{0.5em}
FI\\
N\\
CO\\
\subsection{18.3 Politique de Communication (Communication Policy)}
Pour protéger l’école et ses membres, RBK impose une charte de communication stricte.\\
\vspace{0.5em}
\vspace{0.5em}
\subsection{18.3.1 Interdits Absolus (Red Lines)}
1. Promesse de Gains : Interdiction totale d’utiliser des termes comme ”ROI”, ”Lambo”, ”To the moon”, ”Revenu passif garanti”.\\
\vspace{0.5em}
2. Conseil Financier : Aucun conseil d’achat/vente de token. L’analyse est toujours technique (Smart Contract) ou fondamentale\\
(Tokenomics), jamais spéculative.\\
\vspace{0.5em}
3. Incitation à l’Évasion : Interdiction de promouvoir des schémas de contournement fiscal ou de change.\\
\vspace{0.5em}
\vspace{0.5em}
\begin{confidential}CONFIDENTIEL — PROJET RBK 2.0                                  Page 161 sur 316\end{confidential}
RBK 2.0 — Whitepaper Architecte Web3                                                                                    Chapitre 18. Guide Compliance Web3\\
\vspace{0.5em}
\vspace{0.5em}
\subsection{18.3.2 Obligations (Green Lines)}
• Disclaimer Systématique : ”Ceci n’est pas un conseil financier” sur chaque contenu.\\
\vspace{0.5em}
• Pédagogie du Risque : Rappeler que ”Code is Law” implique aussi ”Bugs are loss”.\\
\vspace{0.5em}
• Testnet First : Toutes les démos se font sur Testnet (argent fictif) sauf nécessité absolue de Mainnet (avec petits montants).\\
\vspace{0.5em}
Avertissement\\
Ce document ne constitue pas un avis juridique. Les scénarios ci-dessus sont des modèles types donnés à titre informatif. Chaque\\
situation entrepreneuriale est unique et nécessite la consultation d’un expert-comptable et d’un avocat spécialisé en droit des affaires\\
tunisien et international.\\
\vspace{0.5em}
\vspace{0.5em}
\vspace{0.5em}
\vspace{0.5em}
EL\\
TI\\
\subsection{18.4 KYC/AML Décentralisé – La Conformité par la Technologie}
\vspace{0.5em}
\vspace{0.5em}
\vspace{0.5em}
\vspace{0.5em}
EN\\
Cette section introduit l’approche privacy-by-design appliquée au KYC1 /AML2 pour concilier conformité et protection des données.\\
\vspace{0.5em}
\vspace{0.5em}
\vspace{0.5em}
\vspace{0.5em}
D\\
FI\\
\subsection{18.4.1 Philosophie du ”Privacy by Design”}
\vspace{0.5em}
N\\
Nous enseignons à passer du KYC centralisé (documents stockés sur serveur vulnérable) à l’Identité Auto-Souveraine (SSI) et aux Preuves\\
CO\\
à Divulgation Nulle de Connaissance (ZK-Proofs).\\
\vspace{0.5em}
\vspace{0.5em}
\subsection{18.4.2 Architecture Technique}
1. Vérification (Claim) : L’utilisateur se vérifie une fois auprès d’un Issuer (ex : Civic) et reçoit une ”Verifiable Credential” (VC).\\
\vspace{0.5em}
2. Stockage (Wallet) : La VC est stockée localement dans le portefeuille de l’utilisateur.\\
1\\
Sécurité/Infra/Ops — Know Your Customer (KYC) : processus de connaissance client vérifiant l’identité avant d’accéder à un service financier ou réglementé.\\
2\\
Sécurité/Infra/Ops — Anti-Money Laundering (AML) : lutte contre le blanchiment d’argent couvrant surveillance des flux, filtrage des sanctions et déclarations\\
aux autorités.\\
\vspace{0.5em}
\vspace{0.5em}
\vspace{0.5em}
\begin{confidential}CONFIDENTIEL — PROJET RBK 2.0                                          Page 162 sur 316\end{confidential}
RBK 2.0 — Whitepaper Architecte Web3                                                                Chapitre 18. Guide Compliance Web3\\
\vspace{0.5em}
\vspace{0.5em}
3. Preuve (Proof) : Pour accéder à un protocole, le portefeuille génère une preuve ZK : ”J’ai +18 ans et je ne suis pas résident US”,\\
sans révéler l’identité réelle.\\
\vspace{0.5em}
\vspace{0.5em}
\subsection{18.4.3 Stack Pratique Enseignée}
• Polygon ID : Création de contrats de staking avec gating géographique via VC.\\
\vspace{0.5em}
• Civic Pass : Protection anti-sybil pour les airdrops.\\
\vspace{0.5em}
• Sismo : Badges ZK pour la réputation (gouvernance DAO).\\
\vspace{0.5em}
\vspace{0.5em}
\vspace{0.5em}
\vspace{0.5em}
EL\\
\subsection{18.5 GDPR \textbackslash{}& Données On-Chain}
\vspace{0.5em}
\vspace{0.5em}
\vspace{0.5em}
\vspace{0.5em}
TI\\
Nous clarifions ici comment respecter le RGPD dans un contexte d’immuabilité on-chain en privilégiant des patterns sûrs.\\
\vspace{0.5em}
\vspace{0.5em}
\vspace{0.5em}
\vspace{0.5em}
EN\\
\subsection{18.5.1 Le Conflit Immuabilité vs Droit à l’Oubli}
\vspace{0.5em}
\vspace{0.5em}
\vspace{0.5em}
\vspace{0.5em}
D\\
FI\\
Règle d’or : Jamais de PII (Personally Identifiable Information) on-chain.\\
\vspace{0.5em}
\vspace{0.5em}
N\\
CO\\
\subsection{18.5.2 Patterns Architecturaux}
• Hash-Only : Stocker keccak256(data) on-chain. La donnée réelle est off-chain avec contrôle d’accès.\\
\vspace{0.5em}
• Chiffrement Asymétrique : Données chiffrées avec la clé publique du destinataire.\\
\vspace{0.5em}
• Pointeurs IPFS : Stocker uniquement le CID (Content ID) sur la blockchain.\\
\vspace{0.5em}
\vspace{0.5em}
\subsection{18.6 Fiscalité Crypto \textbackslash{}& Statut ETE}
Ce segment résume le chemin financier recommandé pour rester conforme tout en optimisant le statut exportateur.\\
\vspace{0.5em}
\vspace{0.5em}
\vspace{0.5em}
\begin{confidential}CONFIDENTIEL — PROJET RBK 2.0                                Page 163 sur 316\end{confidential}
RBK 2.0 — Whitepaper Architecte Web3                                                                   Chapitre 18. Guide Compliance Web3\\
\vspace{0.5em}
\vspace{0.5em}
\subsection{18.6.1 Le Guide de l’Ingénieur-Exportateur}
Les revenus en crypto sont des revenus en devises étrangères. Le statut ETE (Entreprise Totalement Exportatrice) est la clé de l’optimi-\\
sation légale.\\
\vspace{0.5em}
\vspace{0.5em}
\subsection{18.6.2 Flux Financier Recommandé}
1. Réception : Client → Passerelle (Grey.co/Bitwage) → Virement TND.\\
\vspace{0.5em}
2. Comptabilité : Enregistrement au taux du jour BCT.\\
\vspace{0.5em}
3. Déclaration : Trimestrielle auprès de la banque centrale.\\
\vspace{0.5em}
\vspace{0.5em}
\vspace{0.5em}
\vspace{0.5em}
EL\\
TI\\
D    EN\\
FI\\
N\\
CO\\
\vspace{0.5em}
\vspace{0.5em}
\vspace{0.5em}
\vspace{0.5em}
\begin{confidential}CONFIDENTIEL — PROJET RBK 2.0                                Page 164 sur 316\end{confidential}
\section{GOUVERNANCE, ÉTHIQUE \textbackslash{}& TRANSPARENCE}
\vspace{0.5em}
\vspace{0.5em}
RBK 2.0 aspire à devenir une institution de confiance. Cela exige une gouvernance partagée et une transparence radicale sur nos\\
\vspace{0.5em}
\vspace{0.5em}
\vspace{0.5em}
\vspace{0.5em}
EL\\
résultats.\\
\vspace{0.5em}
\vspace{0.5em}
\vspace{0.5em}
\vspace{0.5em}
TI\\
EN\\
\subsection{19.1 Comité Éthique \textbackslash{}& Pédagogique (CEP)}
Le CEP est l’organe de contre-pouvoir indépendant qui garantit que l’école reste fidèle à sa mission.\\
\vspace{0.5em}
\vspace{0.5em}
\vspace{0.5em}
\vspace{0.5em}
D\\
FI\\
\subsection{19.1.1 Composition (5 Membres)}
\vspace{0.5em}
N\\
1. Président : Une figure de la Tech en Tunisie (ex : CTO d’une Startup à succès, non lié à RBK).\\
CO\\
2. Représentant Alumni : Élu par la DAO des anciens.\\
\vspace{0.5em}
3. Représentant Étudiants : Délégué de la promo en cours.\\
\vspace{0.5em}
4. Expert Éducation : Un pédagogue ou universitaire.\\
\vspace{0.5em}
5. CEO RBK : Voix consultative (ne vote pas sur sa propre rémunération ou sa révocation).\\
\vspace{0.5em}
\vspace{0.5em}
\subsection{19.1.2 Mandat}
Le CEP se réunit trimestriellement pour :\\
\vspace{0.5em}
\vspace{0.5em}
165\\
RBK 2.0 — Whitepaper Architecte Web3                                                             Chapitre 19. Gouvernance \textbackslash{}& transparence\\
\vspace{0.5em}
\vspace{0.5em}
• Valider les changements majeurs de curriculum.\\
\vspace{0.5em}
• Arbitrer les contentieux ISA complexes (ex : demande de grâce pour cas de force majeure).\\
\vspace{0.5em}
• Auditer les taux de placement déclarés.\\
\vspace{0.5em}
\vspace{0.5em}
\subsection{19.2 Transparence Radicale (Open Metrics)}
Contrairement aux écoles opaques, RBK publie ses KPI en temps réel sur une page publique ”Status” (et On-chain).\\
\vspace{0.5em}
\vspace{0.5em}
\vspace{0.5em}
\vspace{0.5em}
EL\\
Metriques Publiques (Dashboard)\\
\vspace{0.5em}
\vspace{0.5em}
\vspace{0.5em}
\vspace{0.5em}
TI\\
• Taux de Placement Réel : Calculé à J+180 (CDI/Freelance).\\
\vspace{0.5em}
\vspace{0.5em}
\vspace{0.5em}
\vspace{0.5em}
EN\\
• Salaire Médian de Sortie : Basé sur les fiches de paie anonymisées.\\
\vspace{0.5em}
• Taux de Remboursement ISA : \textbackslash{}% de recouvrement (indicateur de santé financière).\\
\vspace{0.5em}
\vspace{0.5em}
\vspace{0.5em}
\vspace{0.5em}
D\\
• Diversité : Ratio Homme/Femme et Répartition Géographique (Hors Grand Tunis).\\
\vspace{0.5em}
\vspace{0.5em}
\vspace{0.5em}
\subsection{19.3 Charte de Déontologie}
FI\\
N\\
RBK s’engage formellement sur les points suivants :\\
O\\
\vspace{0.5em}
1. Pas de Diplôme de Complaisance : Un étudiant qui paie Upfront mais échoue aux examens techniques ne reçoit PAS de certification.\\
C\\
\vspace{0.5em}
\vspace{0.5em}
Le niveau ne s’achète pas.\\
\vspace{0.5em}
2. Consentement Éclairé ISA : Chaque candidat reçoit une simulation ”Worst Case” (Salaire élevé = Paiement max) avant de signer.\\
\vspace{0.5em}
3. Neutralité Technologique : Bien que financés par des écosystèmes (ex : Solana), nous enseignons l’ingénierie fondamentale, pas le\\
dogmatisme. Nous critiquons les faiblesses de chaque chaine.\\
\vspace{0.5em}
4. Protection des Données : Refus de monétiser les données étudiants auprès de recruteurs tiers sans ”Opt-in” spécifique.\\
\vspace{0.5em}
\vspace{0.5em}
\begin{confidential}CONFIDENTIEL — PROJET RBK 2.0                              Page 166 sur 316\end{confidential}
RBK 2.0 — Whitepaper Architecte Web3                                                                              Chapitre 19. Gouvernance \textbackslash{}& transparence\\
\vspace{0.5em}
\vspace{0.5em}
\vspace{0.5em}
\subsection{19.4 Structure Juridique et Rôles (Branding)}
Pour assurer une clarté totale vis-à-vis des étudiants et partenaires, nous distinguons les entités comme suit :\\
\vspace{0.5em}
ReBootKamp (RBK) Tunisie : L’entité légale opératrice historique. Elle porte l’agrément de formation, gère les locaux (Ariana), les\\
contrats étudiants (ISA inclus via véhicule dédié) et le staff administratif. C’est le garant de la conformité locale.\\
\vspace{0.5em}
Nexus Réussite : Le maître d’œuvre pédagogique exclusif mandaté par RBK. Nexus conçoit et met en oeuvre le curriculum, contracte et\\
pilote les mentors, valide les Block Checks et fournit les reporting exigés par la gouvernance.\\
\vspace{0.5em}
Money Factory AI : Le fournisseur technologique référencé par Nexus. Basé à Dubai/Singapour, Money Factory délivre la plateforme\\
Venture Engine, maintient l’infrastructure de certification On-chain (SBT) et assure le support technique selon les SLA1 contractuels.\\
\vspace{0.5em}
\vspace{0.5em}
\vspace{0.5em}
\vspace{0.5em}
EL\\
Le Programme ”RBK 3.0” : Est une Offre RBK opérée avec Nexus Réussite et adossée aux briques technologiques de Money Factory AI,\\
\vspace{0.5em}
\vspace{0.5em}
\vspace{0.5em}
\vspace{0.5em}
TI\\
garantissant une chaîne de valeur clairement contractuelle et maîtrisée.\\
\vspace{0.5em}
\vspace{0.5em}
\vspace{0.5em}
\vspace{0.5em}
D     EN\\
FI\\
N\\
CO\\
\vspace{0.5em}
\vspace{0.5em}
\vspace{0.5em}
\vspace{0.5em}
1\\
Finance/Juridique/Compliance — Service Level Agreement contractuel décrivant engagements de service, métriques, pénalités et modalités de mesure.\\
\vspace{0.5em}
\vspace{0.5em}
\begin{confidential}CONFIDENTIEL — PROJET RBK 2.0                                          Page 167 sur 316\end{confidential}
\section{IMPACT SOCIAL \textbackslash{}& ALIGNEMENT ODD}
\vspace{0.5em}
\vspace{0.5em}
RBK 2.0 est une entreprise à mission. Notre but est de transférer de la richesse du PIB mondial (Web3) vers l’économie locale tunisienne\\
\vspace{0.5em}
\vspace{0.5em}
\vspace{0.5em}
\vspace{0.5em}
EL\\
et africaine.\\
\vspace{0.5em}
\vspace{0.5em}
\vspace{0.5em}
\vspace{0.5em}
TI\\
EN\\
\subsection{20.1 Contribution aux Objectifs de Développement Durable (ONU)}
• ODD 4 : Éducation de Qualité. Nous démocratisons l’accès à une formation d’élite (niveau Ivy League) sans barrière financière\\
\vspace{0.5em}
\vspace{0.5em}
\vspace{0.5em}
\vspace{0.5em}
D\\
grâce à l’ISA.\\
\vspace{0.5em}
\vspace{0.5em}
\vspace{0.5em}
\vspace{0.5em}
FI\\
• ODD 8 : Travail Décent et Croissance Économique. Nous créons des emplois à haute valeur ajoutée, exportateurs de services, et\\
N\\
rémunérés en devises fortes (via le statut local adéquat).\\
CO\\
• ODD 9 : Industrie, Innovation et Infrastructure. Nous formons les architectes de l’infrastructure financière de demain.\\
\vspace{0.5em}
\vspace{0.5em}
\subsection{20.2 Indicateurs de Performance Sociale}
Nous listons ici les indicateurs qui mesurent l’impact social réel du programme, au-delà des métriques financières.\\
\vspace{0.5em}
\vspace{0.5em}
\subsection{20.2.1 1. Inclusion des Femmes dans la Tech}
Le Web3 souffre d’un déficit de diversité criant. RBK 2.0 met en place des mesures proactives :\\
\vspace{0.5em}
\vspace{0.5em}
168\\
RBK 2.0 — Whitepaper Architecte Web3                                                     Chapitre 20. IMPACT SOCIAL \textbackslash{}& ALIGNEMENT ODD\\
\vspace{0.5em}
\vspace{0.5em}
• Bourses ”Women in Web3” : Le coût Upfront est réduit de 50\textbackslash{}% pour les candidates validant la Piscine (financé par partenaires).\\
\vspace{0.5em}
• Objectif 2026 : Atteindre 30\textbackslash{}% de femmes par cohorte (vs 5\textbackslash{}% moyenne secteur).\\
\vspace{0.5em}
\vspace{0.5em}
\subsection{20.2.2 2. Décentralisation Régionale}
Le talent est partout, les opportunités sont à la capitale.\\
\vspace{0.5em}
• Recrutement National : Roadshow dans les universités de l’intérieur (Sfax, Gabès, Gafsa).\\
\vspace{0.5em}
• Hébergement : Partenariats avec des foyers pour faciliter l’installation à Tunis durant les 4 mois intensifs.\\
\vspace{0.5em}
\vspace{0.5em}
\vspace{0.5em}
\vspace{0.5em}
EL\\
\subsection{20.2.3 3. Empreinte Carbone et Compensation}
\vspace{0.5em}
\vspace{0.5em}
\vspace{0.5em}
\vspace{0.5em}
TI\\
La blockchain est perçue comme polluante. RBK nuance et agit :\\
\vspace{0.5em}
\vspace{0.5em}
\vspace{0.5em}
\vspace{0.5em}
EN\\
• Choix Technologique : Solana est une chaîne Proof-of-Stake dont une transaction consomme moins qu’une requête Google (0.0005\\
\vspace{0.5em}
\vspace{0.5em}
\vspace{0.5em}
\vspace{0.5em}
D\\
kWh).\\
\vspace{0.5em}
\vspace{0.5em}
\vspace{0.5em}
\vspace{0.5em}
FI\\
• Compensation : Nous nous engageons à compenser 100\textbackslash{}% de l’empreinte carbone de l’école (serveurs + clim + déplacements staff)\\
\vspace{0.5em}
N\\
via l’achat de crédits carbone certifiés on-chain (ex : Toucan Protocol).\\
CO\\
\vspace{0.5em}
\vspace{0.5em}
\vspace{0.5em}
\vspace{0.5em}
\begin{confidential}CONFIDENTIEL — PROJET RBK 2.0                                    Page 169 sur 316\end{confidential}
\section{INFRASTRUCTURE SBT \textbackslash{}& CERTIFICATION}
\vspace{0.5em}
\vspace{0.5em}
La certification RBK 2.0 n’est pas un fichier PDF. C’est une preuve cryptographique On-chain, incensurable et composable, matérialisée\\
\vspace{0.5em}
\vspace{0.5em}
\vspace{0.5em}
\vspace{0.5em}
EL\\
par un Soulbound Token (SBT).\\
\vspace{0.5em}
\vspace{0.5em}
\vspace{0.5em}
\vspace{0.5em}
TI\\
EN\\
\subsection{21.1 Philosophie : ”Don’t Trust, Verify”}
Contrairement aux diplômes traditionnels (facilement falsifiables), le SBT RBK est :\\
\vspace{0.5em}
\vspace{0.5em}
\vspace{0.5em}
\vspace{0.5em}
D\\
• Permanent : Ancré sur la blockchain (Solana/Polygon).\\
\vspace{0.5em}
\vspace{0.5em}
\vspace{0.5em}
FI\\
N\\
• Révocable : En cas de triche avérée post-diplomation (mécanisme de slashing).\\
CO\\
• Riche : Contient des métadonnées prouvant les compétences (liens vers repos, hash des commits).\\
\vspace{0.5em}
\vspace{0.5em}
\subsection{21.2 Stack Technique SBT}
Nous posons ici les choix d’implémentation du SBT afin de garantir coût, performance et auditabilité.\\
\vspace{0.5em}
\vspace{0.5em}
\subsection{21.2.1 Choix du Standard}
Nous utilisons une architecture hybride optimisée pour le coût et la performance.\\
\vspace{0.5em}
\vspace{0.5em}
\vspace{0.5em}
170\\
RBK 2.0 — Whitepaper Architecte Web3                                                   Chapitre 21. INFRASTRUCTURE SBT \textbackslash{}& CERTIFICATION\\
\vspace{0.5em}
\vspace{0.5em}
• Réseau : Polygon POS (compatibilité EVM, frais nuls pour l’utilisateur) ou Solana (Compressed NFTs pour le volume).\\
\vspace{0.5em}
• Standard : ERC-5192 (Minimal Soulbound NFTs) sur EVM.\\
\vspace{0.5em}
Interface SBT (Solidity)\\
\vspace{0.5em}
\vspace{0.5em}
interface ISoulbound \textbackslash{}{\\
/// @notice Émis quand un SBT est minté. Locké par défaut.\\
event Locked(uint256 tokenId);\\
/// @notice Le transfert est impossible sauf pour le burn (révocation).\\
function locked(uint256 tokenId) external view returns (bool);\\
\vspace{0.5em}
\vspace{0.5em}
\vspace{0.5em}
\vspace{0.5em}
EL\\
/// @notice Seul l'admin (Smart Contract RBK) peut révoquer.\\
function revoke(uint256 tokenId) external onlyRole(ADMIN);\\
\vspace{0.5em}
\vspace{0.5em}
\vspace{0.5em}
\vspace{0.5em}
TI\\
\textbackslash{}}\\
\vspace{0.5em}
\vspace{0.5em}
\vspace{0.5em}
\vspace{0.5em}
D      EN\\
\subsection{21.3 Cycle de Vie de la Certification}
\vspace{0.5em}
\vspace{0.5em}
\vspace{0.5em}
FI\\
N Trigger\\
LMS (Web2)\\
Valide Gate\\
CO\\
Étudiant                                           Oracle RBK\\
\vspace{0.5em}
\vspace{0.5em}
Signe Tx\\
\vspace{0.5em}
\vspace{0.5em}
Wallet (SBT)              SC Factory\\
Mint\\
\vspace{0.5em}
FIG. 21.1 : Workflow d’Émission Automatisé\\
\vspace{0.5em}
\vspace{0.5em}
\subsection{21.4 Conformité RGPD \textbackslash{}& Privacy}
\vspace{0.5em}
\begin{confidential}CONFIDENTIEL — PROJET RBK 2.0                                   Page 171 sur 316\end{confidential}
RBK 2.0 — Whitepaper Architecte Web3                                                Chapitre 21. INFRASTRUCTURE SBT \textbackslash{}& CERTIFICATION\\
\vspace{0.5em}
\vspace{0.5em}
\vspace{0.5em}
La blockchain est publique, mais les données personnelles ne le sont pas.\\
\vspace{0.5em}
• On-Chain : Un pointeur vers le credential (hash + CID) et un indicateur d’état. Aucune donnée nominative directe n’est stockée.\\
\vspace{0.5em}
• Off-Chain (Verifiable Credentials) : Les informations identifiantes sont conservées dans une base chiffrée gérée par RBK, signées\\
par Nexus Réussite. La preuve de réussite est exposée via un schéma W3C VC.\\
\vspace{0.5em}
• Zero-Knowledge Proofs : L’étudiant présente une preuve ZK attestant qu’il détient un credential valide sans divulguer ses données\\
(utilisation de Sismo/Polygon ID).\\
\vspace{0.5em}
• Droit à l’oubli : Demande de burn + purge du credential off-chain ; la DPA RBK/Nexus documente le processus et l’horodatage des\\
\vspace{0.5em}
\vspace{0.5em}
\vspace{0.5em}
\vspace{0.5em}
EL\\
suppressions.\\
\vspace{0.5em}
\vspace{0.5em}
\vspace{0.5em}
\vspace{0.5em}
TI\\
\subsection{21.5 Cas d’Usage : Le Recrutement Instantané}
\vspace{0.5em}
\vspace{0.5em}
\vspace{0.5em}
\vspace{0.5em}
EN\\
Les partenaires B2B (Job Board) peuvent interroger le Smart Contract pour filtrer les candidats : ”Montre-moi tous les wallets qui\\
\vspace{0.5em}
\vspace{0.5em}
\vspace{0.5em}
\vspace{0.5em}
D\\
détiennent le SBT ’Solana Advanced’ ET le SBT ’Security Auditor’.” Cela réduit le temps de screening de plusieurs jours à quelques milli-\\
\vspace{0.5em}
\vspace{0.5em}
\vspace{0.5em}
\vspace{0.5em}
FI\\
secondes.\\
\vspace{0.5em}
\vspace{0.5em}
N\\
CO\\
\vspace{0.5em}
\vspace{0.5em}
\vspace{0.5em}
\vspace{0.5em}
\begin{confidential}CONFIDENTIEL — PROJET RBK 2.0                               Page 172 sur 316\end{confidential}
\section{FEUILLE DE ROUTE 120 JOURS}
\vspace{0.5em}
\vspace{0.5em}
\vspace{0.5em}
\vspace{0.5em}
EL\\
\subsection{22.1 Timeline des Opérations}
Cette feuille de route couvre l’horizon critique de J-90 (Sprint 0 confidentiel) à J+120 (Fin de la première cohorte). Elle est conçue\\
\vspace{0.5em}
\vspace{0.5em}
\vspace{0.5em}
\vspace{0.5em}
TI\\
pour sécuriser les fondations avant d’accélérer sur l’acquisition et la production publique. Nos hypothèses de départ incluent : une équipe\\
\vspace{0.5em}
\vspace{0.5em}
\vspace{0.5em}
\vspace{0.5em}
EN\\
core opérationnelle (CEO, CTO, Head of Program Nexus), la disponibilité des mentors clés à J0 et la validation du modèle juridique ISA\\
en amont.\\
\vspace{0.5em}
\vspace{0.5em}
\vspace{0.5em}
\vspace{0.5em}
D\\
FI\\
Phase 0 : J-90 → J-1 (Sprint 01 — Sécurisation des Fondations) Objectif : verrouiller juridique, finances et prototype pédagogique\\
avant toute communication externe. Livrables Must-Have2 :\\
N\\
CO\\
1. Avis juridique signé couvrant contrat ISA, statut exportateur et gouvernance RBK ↔ Nexus.\\
\vspace{0.5em}
2. Fonds de Garantie ISA doté à hauteur de 150 000 TND sur compte séquestré RBK (preuve bancaire).\\
\vspace{0.5em}
3. MVP du LMS livré : module Genesis entièrement jouable, reporting Venture Engine synchronisé.\\
\vspace{0.5em}
4. Recrutement et onboarding des trois mentors principaux (Lead Genesis, Lead Explorer, Responsable Bien-être).\\
\vspace{0.5em}
5. Cellule Discord interne opérationnelle (staff uniquement) + kit ”Mentor-in-a-Box” initial (replays, playbooks).\\
1\\
Sécurité/Infra/Ops — Anti-Money Laundering (AML) : lutte contre le blanchiment d’argent couvrant surveillance des flux, filtrage des sanctions et déclarations\\
aux autorités.\\
2\\
Projet/Qualité — Must, Should, Could, Won’t (MoSCoW) : méthode de priorisation cadrant le périmètre et arbitrant les fonctionnalités.\\
\vspace{0.5em}
\vspace{0.5em}
173\\
RBK 2.0 — Whitepaper Architecte Web3                                                               Chapitre 22. FEUILLE DE ROUTE 120 JOURS\\
\vspace{0.5em}
\vspace{0.5em}
Principe de réserve : aucune campagne marketing ou annonce publique tant que les livrables Sprint 0 ne sont pas validés en Comité\\
Exécutif Conjoint.\\
\vspace{0.5em}
Phase 1 : J1 → J60 (Production \textbackslash{}& Industrialisation)       Objectif : transformer le MVP pédagogique en usine ”Senior-by-Design” prête à\\
scaler. Livrables Must-Have :\\
\vspace{0.5em}
1. LMS (Learning Management System) durci : 100\textbackslash{}% des labs Tronc Commun (S1-S12) relus, scripts d’évaluation automatisés.\\
\vspace{0.5em}
2. Pipeline CI/CD et contrôles sécurité (lint, fuzzing) actifs sur les repos Golden Templates.\\
\vspace{0.5em}
3. Playbooks Wellbeing déployés : questionnaires hebdo, tableaux de bord santé, protocole d’escalade testé.\\
\vspace{0.5em}
\vspace{0.5em}
\vspace{0.5em}
\vspace{0.5em}
EL\\
4. Roster mentors complété (5 ”Core” + bench alumni) et contrats Nexus signés.\\
\vspace{0.5em}
\vspace{0.5em}
\vspace{0.5em}
\vspace{0.5em}
TI\\
5. Kit marketing préparé mais non diffusé : landing page, simulateur ROI, série Building in Public relue.\\
\vspace{0.5em}
\vspace{0.5em}
\vspace{0.5em}
\vspace{0.5em}
EN\\
Phase 2 : J61 → J90 (Pré-lancement \textbackslash{}& Sélection) Objectif : ouvrir la Piscine et sélectionner la cohorte Alpha en mode acquisition\\
\vspace{0.5em}
\vspace{0.5em}
\vspace{0.5em}
\vspace{0.5em}
D\\
disciplinée. Livrables Must-Have :\\
\vspace{0.5em}
\vspace{0.5em}
\vspace{0.5em}
\vspace{0.5em}
FI\\
1. Campagne Building in Public activée (Discord public, Office Hours, challenges Ship or Die).\\
\vspace{0.5em}
N\\
2. Simulateur ROI live et connecté au CRM pour scorer chaque lead.\\
CO\\
3. Piscine Rust (4 semaines) exécutée avec +100 candidats et reporting de fatigue partagé.\\
\vspace{0.5em}
4. 20 à 25 candidats ”Alpha” signés, contrats CPPS (ISA) émis et archivés.\\
\vspace{0.5em}
Phase 3 : J91 → J120 (Lancement Opérationnel)            Objectif : livrer une expérience cohorte irréprochable dès le jour 1. Livrables\\
Must-Have :\\
\vspace{0.5em}
1. Onboarding physique/virtuel orchestré (kits, accès, chartes) avec audit qualité Nexus.\\
\vspace{0.5em}
2. Sprint pédagogique Semaine 1 livré (Genesis Pool) avec NPS \textgreater{} 50 et incident log ”zéro critique”.\\
\vspace{0.5em}
\vspace{0.5em}
\begin{confidential}CONFIDENTIEL — PROJET RBK 2.0                                Page 174 sur 316\end{confidential}
RBK 2.0 — Whitepaper Architecte Web3                                                                                   Chapitre 22. FEUILLE DE ROUTE 120 JOURS\\
\vspace{0.5em}
\vspace{0.5em}
3. Rituels cadencés : Demo hebdo, Block Checks, revue Wellbeing par le Responsable Bien-être.\\
\vspace{0.5em}
\vspace{0.5em}
SPRINT 0               PRODUCTION                SÉLECTION              LANCEMENT\\
(Juridique, MVP)          (Contenu, Infra)         (Piscine, CPPS)         (Promo Alpha)\\
\vspace{0.5em}
\vspace{0.5em}
J-90                      J-60                      J0                      J60                   J120\\
\vspace{0.5em}
FIG. 22.1 : Timeline 120 jours (Vue Exécutive) — Sprint 0 inclus\\
\vspace{0.5em}
\vspace{0.5em}
\begin{center}\begin{tabular}TAB. 22.1 : Checklist Go/No-Go (Gates)\end{tabular}\end{center}
\vspace{0.5em}
\vspace{0.5em}
\vspace{0.5em}
\vspace{0.5em}
EL\\
Gate          Validation      Critères Obligatoires                                  En cas de KO\\
\vspace{0.5em}
\vspace{0.5em}
\vspace{0.5em}
\vspace{0.5em}
TI\\
Gate 0 (J-60) Sprint 0 Validé Avis juridique signé, MVP LMS, mentors clés onboardés. Pas de communication externe.\\
A (J-30)      Légal Ready     ISA conforme, Fonds bloqué (150 000 TND), assurances. Report Lancement.\\
\vspace{0.5em}
\vspace{0.5em}
\vspace{0.5em}
\vspace{0.5em}
EN\\
B (J0)        Infra Ready     Syllabus V1 figé, LMS opérationnel, Team Staffée.      Mode ”Dégradé” (Contenu JIT).\\
C (J60)       Candidats       \textgreater{} 100 Inscrits Piscine qualifiés.                      Extension période Marketing 2 sem.\\
\vspace{0.5em}
\vspace{0.5em}
\vspace{0.5em}
\vspace{0.5em}
D\\
D (J90)       Promo Ready 25 Contrats signés, 0 contentieux.                         Réduction taille promo.\\
\vspace{0.5em}
\vspace{0.5em}
\vspace{0.5em}
\vspace{0.5em}
FI\\
N\\
CO\\
\subsection{22.2 Jalons Clés \textbackslash{}& Actions}
Nous pilotons l’exécution par ”Workstreams”. Chaque action est priorisée (P0 Bloquant, P1 Critique, P2 Important).\\
\vspace{0.5em}
Workstream A : Juridique \textbackslash{}& Conformité\\
\vspace{0.5em}
Owner : COO / Legal. DoD3 : Tous les contrats sont signés et stockés sécurisés.\\
\vspace{0.5em}
• P0 : Validation modèle ISA v3 avec cabinet spécialisé.\\
3\\
Projet/Management/Qualité — Definition of Done (DoD) : liste les critères minimaux pour considérer un élément comme terminé et livrable.\\
\vspace{0.5em}
\vspace{0.5em}
\vspace{0.5em}
\vspace{0.5em}
\begin{confidential}CONFIDENTIEL — PROJET RBK 2.0                                                   Page 175 sur 316\end{confidential}
RBK 2.0 — Whitepaper Architecte Web3                                                                            Chapitre 22. FEUILLE DE ROUTE 120 JOURS\\
\vspace{0.5em}
\vspace{0.5em}
• P0 : Création structure juridique porteuse (SPV ou LLC).\\
\vspace{0.5em}
• P1 : Rédaction CGV/CGU et Politique de Confidentialité (GDPR).\\
\vspace{0.5em}
Workstream B : Produit Pédagogique\\
\vspace{0.5em}
Owner : Lead Instructor. Definition of Done : Contenu accessible sur LMS et testé par un pairs.\\
\vspace{0.5em}
• P0 : Finalisation structure Chapitres 5-8 (Syllabus détaillé).\\
\vspace{0.5em}
• P1 : Création des ”Golden Templates” (Repos de référence).\\
\vspace{0.5em}
• P1 : Banque de Quiz (300 questions) pour l’évaluation continue.\\
\vspace{0.5em}
\vspace{0.5em}
\vspace{0.5em}
\vspace{0.5em}
EL\\
TI\\
Workstream C : Stack \textbackslash{}& Ops\\
\vspace{0.5em}
\vspace{0.5em}
\vspace{0.5em}
\vspace{0.5em}
EN\\
Owner : CTO. Definition of Done : Infra stable, monitoring actif, zéro friction étudiant.\\
\vspace{0.5em}
• P0 : Configuration Workspace GitHub (Orga, Teams, Permissions).\\
\vspace{0.5em}
\vspace{0.5em}
\vspace{0.5em}
\vspace{0.5em}
D\\
FI\\
• P1 : Déploiement serveur Discord (Bots, Rôles, Channels).\\
\vspace{0.5em}
\vspace{0.5em}
N\\
• P2 : Automatisation onboarding (Zapier/Make : Typeform → Notion → Discord).\\
CO\\
Workstream D : Acquisition\\
\vspace{0.5em}
Owner : CMO. Definition of Done : Pipeline rempli à 150\textbackslash{}% des objectifs.\\
\vspace{0.5em}
• P0 : Lancement Site Web V1 (Landing, FAQ, Team).\\
\vspace{0.5em}
• P1 : Mise en ligne Simulateur ROI (Lead Magnet) et suivi NPS4 .\\
\vspace{0.5em}
• P1 : Campagne LinkedIn5 /Twitter ”Building in Public” (Daily).\\
4\\
Marketing — Net Promoter Score (NPS) : indicateur de recommandation (promoteurs - détracteurs) pour mesurer satisfaction et fidélité.\\
5\\
Organisation — Réseau professionnel utilisé pour la visibilité, l’acquisition de leads et le recrutement de partenaires ou candidats.\\
\vspace{0.5em}
\vspace{0.5em}
\begin{confidential}CONFIDENTIEL — PROJET RBK 2.0                                           Page 176 sur 316\end{confidential}
RBK 2.0 — Whitepaper Architecte Web3                                                                              Chapitre 22. FEUILLE DE ROUTE 120 JOURS\\
\vspace{0.5em}
\vspace{0.5em}
Les actions opérationnelles s’appuient sur une coordination Ops et un suivi CRM6 pour tracer l’acquisition et la montée en charge.\\
\vspace{0.5em}
\begin{center}\begin{tabular}TAB. 22.3 : Backlog Opérationnel (Extrait Top Actions)\end{tabular}\end{center}
\vspace{0.5em}
\vspace{0.5em}
ID        Action                     Prio Owner        Critère Acceptation\\
\vspace{0.5em}
LEG-01 Validation ISA Avocat          P0 Legal         Mémo juridique signé (”Safe to operate”).\\
\vspace{0.5em}
LEG-02 Setup Compte Séquestre P0 Finance               IBAN fourni, dotation 150 000 TND créditée.\\
\vspace{0.5em}
PED-01 Syllabus S1-S4 Ready           P0 Lead Inst. PDF + Markdown sur LMS.\\
\vspace{0.5em}
OPS-01 Invite Mentors Discord         P1 Ops           Tous les mentors ont le rôle ”Sensei”.\\
\vspace{0.5em}
\vspace{0.5em}
\vspace{0.5em}
\vspace{0.5em}
EL\\
ACQ-01 Page Candidature Live          P0 Mktg          Formulaire testé, data arrive dans CRM.\\
\vspace{0.5em}
\vspace{0.5em}
\vspace{0.5em}
\vspace{0.5em}
TI\\
EN\\
Risques Opérationnels Critiques à 120 Jours\\
\vspace{0.5em}
\vspace{0.5em}
\vspace{0.5em}
\vspace{0.5em}
D\\
• Retard Légal (ISA) : Blocage du modèle économique. → Mitigation : Modèle de repli ”Upfront différé”.\\
\vspace{0.5em}
\vspace{0.5em}
\vspace{0.5em}
\vspace{0.5em}
FI\\
• Déficit Candidats Qualifiés : Piscine vide. → Mitigation : Activation réseau partenaires (Bourses).\\
\vspace{0.5em}
N\\
• Churn Mentors : Départ en cours de route. → Mitigation : Roster de backup (Alumni experts).\\
CO\\
\vspace{0.5em}
\subsection{22.3 Diagramme de Gantt Macro}
Ce diagramme de Gantt illustre le chemin critique. Les tâches Juridiques et Production Pédagogique sont les goulots d’étranglement\\
initiaux.\\
6\\
Produit/Business — Customer Relationship Management (CRM) : outil ou processus de gestion de la relation client (prospects, ventes, support, suivi).\\
\vspace{0.5em}
\vspace{0.5em}
\vspace{0.5em}
\vspace{0.5em}
\begin{confidential}CONFIDENTIEL — PROJET RBK 2.0                                            Page 177 sur 316\end{confidential}
RBK 2.0 — Whitepaper Architecte Web3                                             Chapitre 22. FEUILLE DE ROUTE 120 JOURS\\
\vspace{0.5em}
\vspace{0.5em}
\vspace{0.5em}
2025\\
\vspace{0.5em}
01    02               03            04                 05           06\\
\vspace{0.5em}
\vspace{0.5em}
SPRINT 0\\
\vspace{0.5em}
Juridique ISA\\
\vspace{0.5em}
Fonds Séquestré\\
\vspace{0.5em}
MVP LMS\\
\vspace{0.5em}
\vspace{0.5em}
\vspace{0.5em}
\vspace{0.5em}
EL\\
Mentors Core\\
\vspace{0.5em}
\vspace{0.5em}
\vspace{0.5em}
\vspace{0.5em}
TI\\
PHASE 1\\
\vspace{0.5em}
\vspace{0.5em}
\vspace{0.5em}
\vspace{0.5em}
EN\\
Durcissement LMS\\
\vspace{0.5em}
\vspace{0.5em}
\vspace{0.5em}
\vspace{0.5em}
D\\
CI/CD \textbackslash{}& Sécurité\\
\vspace{0.5em}
\vspace{0.5em}
\vspace{0.5em}
\vspace{0.5em}
FI\\
Playbooks Wellbeing\\
N\\
CO\\
Roster Mentors\\
\vspace{0.5em}
PHASE 2\\
\vspace{0.5em}
Campagne BIP\\
\vspace{0.5em}
Piscine Rust\\
\vspace{0.5em}
Signature CPPS\\
\vspace{0.5em}
PHASE 3\\
\vspace{0.5em}
Onboarding Promo\\
\begin{confidential}CONFIDENTIEL — PROJET RBK 2.0                     Page 178 sur 316\end{confidential}
Genesis Week\\
\section{FEUILLE DE ROUTE : LE PLAN DE LANCEMENT (90}
JOURS)\\
\vspace{0.5em}
\vspace{0.5em}
\vspace{0.5em}
\vspace{0.5em}
EL\\
TI\\
\subsection{23.1 MOIS 1 : CADRAGE, ALLIANCE \textbackslash{}& ÉQUIPE NOYAU (J0 - J30)}
\vspace{0.5em}
\vspace{0.5em}
\vspace{0.5em}
\vspace{0.5em}
EN\\
Cette première phase fixe le socle stratégique, les partenaires et le noyau humain avant d’engager la production.\\
\vspace{0.5em}
\vspace{0.5em}
\vspace{0.5em}
\vspace{0.5em}
D\\
\subsection{23.1.1 Validation \textbackslash{}& Cadrage Stratégique}
\vspace{0.5em}
\vspace{0.5em}
\vspace{0.5em}
\vspace{0.5em}
FI\\
• Finalisation du modèle financier et des projections de cash-flow.\\
N\\
CO\\
• Signature du contrat-cadre RBK ↔ Nexus Réussite (maître d’œuvre) et contractualisation de la prestation technologique Money\\
Factory AI via Nexus.\\
\vspace{0.5em}
\vspace{0.5em}
\subsection{23.1.2 Constitution de l’Alliance Écosystémique}
• Signature des MOU avec les partenaires locaux (Universités, Incubateurs).\\
\vspace{0.5em}
• Validation du soutien de la Solana Foundation (Superteam).\\
\vspace{0.5em}
\vspace{0.5em}
\vspace{0.5em}
\vspace{0.5em}
179\\
RBK 2.0 — Whitepaper Architecte Web3                                   Chapitre 23. FEUILLE DE ROUTE : LE PLAN DE LANCEMENT (90 JOURS)\\
\vspace{0.5em}
\vspace{0.5em}
\subsection{23.1.3 Recrutement de l’Équipe Pilote}
• Recrutement du Lead Instructor / Head of Curriculum (Expert Rust/Solana).\\
\vspace{0.5em}
• Désignation d’un Responsable Administratif \textbackslash{}& Logistique dédié au Studio.\\
\vspace{0.5em}
• Identification du pool de mentors internationaux (Guest Lecturers).\\
\vspace{0.5em}
\vspace{0.5em}
\subsection{23.2 MOIS 2 : PRODUCTION DE L’ARSENAL \textbackslash{}& INFRASTRUCTURE (J31 - J60)}
L’objectif est de bâtir l’infrastructure technique et les contenus de référence.\\
\vspace{0.5em}
\vspace{0.5em}
\vspace{0.5em}
\vspace{0.5em}
EL\\
TI\\
\subsection{23.2.1 Ingénierie Pédagogique (Les « Golden Templates »)}
• Rédaction détaillée des syllabus pour le Tronc Commun et les Tracks A (Solana) \textbackslash{}& B (EVM).\\
\vspace{0.5em}
\vspace{0.5em}
\vspace{0.5em}
\vspace{0.5em}
EN\\
• Création des dépôts GitHub de référence (repos « or ») incluant les architectures de base, les tests de sécurité et les pipelines CI/CD.\\
\vspace{0.5em}
\vspace{0.5em}
\vspace{0.5em}
\vspace{0.5em}
D\\
FI\\
• Élaboration des « Incident Drills » (simulations de hacks pour les exercices du vendredi).\\
\vspace{0.5em}
\vspace{0.5em}
N\\
CO\\
\subsection{23.2.2 Mise en place du Cockpit Technique}
• Configuration des accès aux Nodes RPC premium (Helius pour Solana, Alchemy/Infura pour EVM).\\
\vspace{0.5em}
• Acquisition des licences pour les outils d’IA (Cursor, Windsurf) et de simulation.\\
\vspace{0.5em}
• Installation du LMS et du serveur Discord comme hub de communication principal.\\
\vspace{0.5em}
\vspace{0.5em}
\subsection{23.2.3 Lancement Commercial \textbackslash{}& Marketing}
• Mise en ligne du site web dédié au programme.\\
\vspace{0.5em}
\vspace{0.5em}
\vspace{0.5em}
\begin{confidential}CONFIDENTIEL — PROJET RBK 2.0                                  Page 180 sur 316\end{confidential}
RBK 2.0 — Whitepaper Architecte Web3                                   Chapitre 23. FEUILLE DE ROUTE : LE PLAN DE LANCEMENT (90 JOURS)\\
\vspace{0.5em}
\vspace{0.5em}
• Lancement de la campagne marketing « Elite Only » sur LinkedIn et Twitter (X).\\
\vspace{0.5em}
• Organisation du premier « Tunisian Web3 Builder Meetup » pour générer des leads qualifiés.\\
\vspace{0.5em}
\vspace{0.5em}
\subsection{23.3 MOIS 3 : SÉLECTION \textbackslash{}& LANCEMENT « PROMO ALPHA » (J61 - J90)}
L’objectif est de filtrer les talents et de démarrer l’immersion.\\
\vspace{0.5em}
\vspace{0.5em}
\subsection{23.3.1 Processus de Sélection d’Élite}
• Tests techniques de pré-requis (JS/TS intensif).\\
\vspace{0.5em}
\vspace{0.5em}
\vspace{0.5em}
\vspace{0.5em}
EL\\
• Entretiens de motivation pour évaluer la pensée systémique et l’autonomie.\\
\vspace{0.5em}
\vspace{0.5em}
\vspace{0.5em}
\vspace{0.5em}
TI\\
• La « Piscine » Rust : Lancement de la phase de filtrage intensif de 4 semaines sans IA pour les 20 candidats présélectionnés.\\
\vspace{0.5em}
\vspace{0.5em}
\vspace{0.5em}
\vspace{0.5em}
D   EN\\
\subsection{23.3.2 Finalisation de la Cohorte}
\vspace{0.5em}
\vspace{0.5em}
\vspace{0.5em}
\vspace{0.5em}
FI\\
• Sélection finale de la Cohorte Alpha (15 à 20 profils maximum pour garantir l’excellence).\\
\vspace{0.5em}
N\\
• Signature des contrats (incluant les clauses ISA le cas échéant).\\
CO\\
• Onboarding sur Superteam Earn pour que les étudiants voient les premières opportunités de revenus dès le début.\\
\vspace{0.5em}
\vspace{0.5em}
\subsection{23.3.3 Kick-off Opérationnel}
• Cérémonie de lancement en présence de partenaires de l’écosystème.\\
\vspace{0.5em}
• Début de la Phase 1 (Fondations \textbackslash{}& Mentalité On-chain).\\
\vspace{0.5em}
\vspace{0.5em}
\vspace{0.5em}
\vspace{0.5em}
\begin{confidential}CONFIDENTIEL — PROJET RBK 2.0                                  Page 181 sur 316\end{confidential}
RBK 2.0 — Whitepaper Architecte Web3                                   Chapitre 23. FEUILLE DE ROUTE : LE PLAN DE LANCEMENT (90 JOURS)\\
\vspace{0.5em}
\vspace{0.5em}
\vspace{0.5em}
\subsection{23.4 RÉCAPITULATIF DES JALONS CLÉS (MILESTONES)}
\vspace{0.5em}
\begin{center}\begin{tabular}TAB. 23.1 : Jalons Clés du Plan de Lancement\end{tabular}\end{center}
\vspace{0.5em}
\vspace{0.5em}
Délai Jalon                         Impact\\
J+15 MOU Solana Foundation signé Crédibilité internationale immédiate.\\
J+30 Équipe pédagogique complète Capacité de production activée.\\
J+45 Golden Templates livrés        Standard de qualité « Senior-by-Design » fixé.\\
J+60 100 leads qualifiés générés    Sécurité du taux de remplissage.\\
\vspace{0.5em}
\vspace{0.5em}
\vspace{0.5em}
\vspace{0.5em}
EL\\
J+75 Fin de la « Piscine » Rust     Cohorte d’élite validée.\\
J+90 Lancement officiel Promo Alpha Début de la transformation de RBK.\\
\vspace{0.5em}
\vspace{0.5em}
\vspace{0.5em}
\vspace{0.5em}
TI\\
EN\\
Ď Annotation Stratégique\\
\vspace{0.5em}
\vspace{0.5em}
\vspace{0.5em}
D\\
Ce plan de 90 jours est agressif mais réaliste. Il repose sur l’utilisation intensive des ressources existantes de RBK (locaux, réseau\\
\vspace{0.5em}
\vspace{0.5em}
\vspace{0.5em}
\vspace{0.5em}
FI\\
alumni) et sur l’apport d’expertise Web3 externe pour l’ingénierie de contenu.\\
\vspace{0.5em}
N\\
CO\\
\vspace{0.5em}
\vspace{0.5em}
\vspace{0.5em}
\vspace{0.5em}
\begin{confidential}CONFIDENTIEL — PROJET RBK 2.0                                 Page 182 sur 316\end{confidential}
\section{TOKEN DE RÉPUTATION \textbackslash{}& ALUMNI PROGRAM}
\vspace{0.5em}
\vspace{0.5em}
\vspace{0.5em}
\vspace{0.5em}
EL\\
\subsection{24.1 RBK Soulbound Tokens (SBTs)}
Le diplôme papier est obsolète. RBK 2.0 certifie les compétences via des Soulbound Tokens (SBTs) : des jetons numériques non-\\
\vspace{0.5em}
\vspace{0.5em}
\vspace{0.5em}
\vspace{0.5em}
TI\\
transférables, infalsifiables, et vérifiables instantanément sur la blockchain. Ce n’est pas un actif financier (pas de prix, pas de marché\\
\vspace{0.5em}
\vspace{0.5em}
\vspace{0.5em}
\vspace{0.5em}
EN\\
secondaire). C’est un CV cryptographique. Chaque SBT représente une compétence acquise (”Rust Ace”), une réalisation (”Capstone\\
Winner”) ou un rôle (”Mentor”).\\
\vspace{0.5em}
\vspace{0.5em}
\vspace{0.5em}
\vspace{0.5em}
D\\
FI\\
Architecture Technique \textbackslash{}& Privacy       Notre système respecte la confidentialité des étudiants.\\
\vspace{0.5em}
N\\
• Issuer : Un wallet Multisig (RBK Board) signe l’émission des badges.\\
CO\\
• Données : Aucune donnée personnelle (Nom/Email) n’est stockée On-chain. Le SBT contient uniquement un Hash de la preuve (ex :\\
hash du commit git ou du certificat PDF).\\
\vspace{0.5em}
• Vérification : L’employeur utilise une dApp RBK pour vérifier la possession du badge et révéler le contenu associé si l’étudiant donne\\
son accord (Signature).\\
\vspace{0.5em}
\vspace{0.5em}
\vspace{0.5em}
\vspace{0.5em}
183\\
RBK 2.0 — Whitepaper Architecte Web3                                                   Chapitre 24. TOKEN DE RÉPUTATION \textbackslash{}& ALUMNI PROGRAM\\
\vspace{0.5em}
\vspace{0.5em}
\vspace{0.5em}
Étudiant (Wallet)                 LMS (Résultats)\\
\vspace{0.5em}
\vspace{0.5em}
\vspace{0.5em}
\vspace{0.5em}
Possède                            Moteur Règles\\
\vspace{0.5em}
\vspace{0.5em}
Valid ?\\
\vspace{0.5em}
\vspace{0.5em}
Multisig Issuer\\
\vspace{0.5em}
\vspace{0.5em}
\vspace{0.5em}
\vspace{0.5em}
EL\\
Mint\\
Blockchain (SBT)\\
\vspace{0.5em}
\vspace{0.5em}
\vspace{0.5em}
\vspace{0.5em}
TI\\
EN\\
FIG. 24.1 : Architecture d’Émission SBT\\
\vspace{0.5em}
\vspace{0.5em}
\vspace{0.5em}
\vspace{0.5em}
D\\
FI\\
N\\
\begin{center}\begin{tabular}TAB. 24.1 : Catalogue des Badges SBT (Extrait)\end{tabular}\end{center}
CO\\
Badge       Niveau Critères                           Valeur Employeur\\
RS-Elite    Gold   Top 5\textbackslash{}% Piscine Rust.               Capacité cognitive, résilience.\\
Solana-Arch Silver Capstone validé avec Audit Clean. ”Production-Ready” Engineer.\\
Auditor-Jr Bronze 3 Rapports de vulnérabilité soumis. Conscience sécurité.\\
Team-Lead Silver A géré une squad de 4 devs.          Soft skills, Management.\\
\vspace{0.5em}
\vspace{0.5em}
\vspace{0.5em}
\vspace{0.5em}
\begin{confidential}CONFIDENTIEL — PROJET RBK 2.0                                    Page 184 sur 316\end{confidential}
RBK 2.0 — Whitepaper Architecte Web3                                             Chapitre 24. TOKEN DE RÉPUTATION \textbackslash{}& ALUMNI PROGRAM\\
\vspace{0.5em}
\vspace{0.5em}
Conformité \textbackslash{}& Anti-Spéculation\\
\vspace{0.5em}
Les SBT RBK sont strictement incessibles. Si un wallet est compromis, le SBT est ”brûlé” (revoked) et réémis vers une nouvelle adresse\\
après vérification d’identité (KYC). Ils n’ont aucune valeur monétaire et ne donnent droit à aucun dividende.\\
\vspace{0.5em}
\vspace{0.5em}
\vspace{0.5em}
\subsection{24.2 Usages des SBT}
Les SBT ne sont pas des objets de collection, ce sont des clés d’accès (”Token Gating”).\\
\vspace{0.5em}
1. Vérification Employeur Instantanée Plus besoin d’appeler l’école pour vérifier un diplôme. L’employeur scanne l’adresse publique\\
\vspace{0.5em}
\vspace{0.5em}
\vspace{0.5em}
\vspace{0.5em}
EL\\
du candidat et voit instantanément ses certifications.\\
\vspace{0.5em}
\vspace{0.5em}
\vspace{0.5em}
\vspace{0.5em}
TI\\
Ď Story : La Vérification en 3 secondes\\
\vspace{0.5em}
\vspace{0.5em}
\vspace{0.5em}
\vspace{0.5em}
EN\\
Avant : Un recruteur reçoit un PDF, doit appeler l’école, attendre 24h pour confirmer qu’il n’est pas falsifié. Coût : Temps +\\
Risque. Avec RBK SBT : Le recruteur colle l’adresse du candidat sur l’Explorer RBK. Le badge ”Certified Graduate” apparaît\\
\vspace{0.5em}
\vspace{0.5em}
\vspace{0.5em}
\vspace{0.5em}
D\\
instantanément avec la signature cryptographique de l’école et le lien vers le code du Capstone. Résultat : Coût 0\textbackslash{}$, 3 secondes,\\
\vspace{0.5em}
\vspace{0.5em}
\vspace{0.5em}
\vspace{0.5em}
FI\\
Confiance Absolue.\\
\vspace{0.5em}
N\\
CO\\
2. Accès au Job Board Premium Seuls les détenteurs du badge ”Ready-to-Deploy” (cursus validé) peuvent voir les offres d’emploi\\
exclusives de nos partenaires ”Gold”. Cela garantit aux recruteurs une qualité de candidature 100\textbackslash{}% filtrée.\\
\vspace{0.5em}
3. Gouvernance Alumni Le poids de vote dans la DAO Alumni est pondéré par les badges. Un ”Senior Mentor” a plus de voix qu’un\\
”New Grad” sur les décisions pédagogiques (mais pas financières).\\
\vspace{0.5em}
\vspace{0.5em}
\vspace{0.5em}
\vspace{0.5em}
\begin{confidential}CONFIDENTIEL — PROJET RBK 2.0                                Page 185 sur 316\end{confidential}
RBK 2.0 — Whitepaper Architecte Web3                                            Chapitre 24. TOKEN DE RÉPUTATION \textbackslash{}& ALUMNI PROGRAM\\
\vspace{0.5em}
\vspace{0.5em}
\begin{center}\begin{tabular}TAB. 24.3 : Usages et Bénéfices des SBT\end{tabular}\end{center}
\vspace{0.5em}
\vspace{0.5em}
Usage     Bénéfice                    SBT Requis      Mécanisme\\
Job Board Accès offres VIP            Certified Dev   Token Gating (Web3 Auth)\\
Mentoring Droit de devenir Mentor     Senior + Pedago Whitelist Manuelle\\
Bounties Accès missions audit         Auditor Level 1 Accès GitHub Repo privé\\
Events    Tickets conférence gratuits Active Member Airdrop Ticket NFT\\
\vspace{0.5em}
\vspace{0.5em}
\vspace{0.5em}
\vspace{0.5em}
\subsection{24.3 Alumni Program Structuré}
\vspace{0.5em}
\vspace{0.5em}
\vspace{0.5em}
\vspace{0.5em}
EL\\
L’Alumni Program est notre ”Moat”. C’est un réseau structuré qui continue d’apporter de la valeur des années après la sortie.\\
\vspace{0.5em}
\vspace{0.5em}
\vspace{0.5em}
\vspace{0.5em}
TI\\
EN\\
Structure en Tiers (Niveaux)     L’engagement est gamifié via des statuts qui offrent des avantages croissants.\\
\vspace{0.5em}
• Tier Bronze (New Grad) : Accès Discord Alumni, Job Board, Annuaire. Condition : Diplômé.\\
\vspace{0.5em}
\vspace{0.5em}
\vspace{0.5em}
\vspace{0.5em}
D\\
FI\\
• Tier Argent (Contributor) : Accès Bounties rémunérés, Invitations Events VIP. Condition : A parrainé 1 étudiant OU donné 10h de\\
mentorat.\\
N\\
CO\\
• Tier Or (Legend) : Accès Fonds Ventures, Siège au Conseil Pédago. Condition : A recruté un Alumni OU créé une startup RBK.\\
\vspace{0.5em}
Gouvernance Le Conseil Alumni (5 membres élus pour un mandat aligné sur un cycle du programme (48 semaines)) gère le budget\\
”Community” (financé par 1\textbackslash{}% des revenus de l’école). Ils décident des apéros, des workshops invités et des partenariats. Règle Anti-Sybil :\\
Seuls les wallets avec un SBT ”Certified” actif depuis \textgreater{} 3 mois peuvent voter.\\
\vspace{0.5em}
\vspace{0.5em}
\vspace{0.5em}
\vspace{0.5em}
\begin{confidential}CONFIDENTIEL — PROJET RBK 2.0                                Page 186 sur 316\end{confidential}
RBK 2.0 — Whitepaper Architecte Web3                                          Chapitre 24. TOKEN DE RÉPUTATION \textbackslash{}& ALUMNI PROGRAM\\
\vspace{0.5em}
\vspace{0.5em}
\vspace{0.5em}
\vspace{0.5em}
Or (Elite)\\
\vspace{0.5em}
\vspace{0.5em}
Argent (Actifs)\\
\vspace{0.5em}
\vspace{0.5em}
Bronze (Masse)\\
\vspace{0.5em}
\vspace{0.5em}
FIG. 24.2 : Pyramide des Tiers Alumni\\
\vspace{0.5em}
\vspace{0.5em}
\vspace{0.5em}
\vspace{0.5em}
EL\\
\begin{center}\begin{tabular}TAB. 24.5 : Roadmap Alumni (Année 1)\end{tabular}\end{center}
\vspace{0.5em}
\vspace{0.5em}
\vspace{0.5em}
\vspace{0.5em}
TI\\
EN\\
Trimestre Initiative             KPI                 Owner\\
Q1        Lancement Discord      100\textbackslash{}% promo inscrite Community Mgr\\
\vspace{0.5em}
\vspace{0.5em}
\vspace{0.5em}
\vspace{0.5em}
D\\
Q2        Premier Apéro Physique 30 participants     Conseil Alumni\\
\vspace{0.5em}
\vspace{0.5em}
\vspace{0.5em}
\vspace{0.5em}
FI\\
Q3        Programme Mentoring 10 binômes actifs      Lead Pédago\\
Q4\\
N\\
Annuaire On-Chain      100\textbackslash{}% profils mintés Tech Lead\\
CO\\
\vspace{0.5em}
\vspace{0.5em}
\vspace{0.5em}
\vspace{0.5em}
\begin{confidential}CONFIDENTIEL — PROJET RBK 2.0                              Page 187 sur 316\end{confidential}
\section{ÉLÉMENTS DE DIFFÉRENCIATION}
\vspace{0.5em}
\vspace{0.5em}
\vspace{0.5em}
\vspace{0.5em}
EL\\
\subsection{25.1 Le Paradigme « Senior-by-Design »}
Le terme ”Junior” est banni de notre vocabulaire. Un étudiant RBK ne sort pas pour ”apprendre le métier”, mais pour ”exécuter\\
\vspace{0.5em}
\vspace{0.5em}
\vspace{0.5em}
\vspace{0.5em}
TI\\
le métier”. L’objectif est de produire un ingénieur immédiatement opérationnel, capable de livrer du code sécurisé en production sans\\
\vspace{0.5em}
\vspace{0.5em}
\vspace{0.5em}
\vspace{0.5em}
EN\\
supervision constante.\\
\vspace{0.5em}
\vspace{0.5em}
\vspace{0.5em}
\vspace{0.5em}
D\\
Mécanisme Opérationnel\\
\vspace{0.5em}
\vspace{0.5em}
\vspace{0.5em}
\vspace{0.5em}
FI\\
• No-AI Piscine : Le filtre d’entrée se fait à la dure (Rust pur, sans Copilot) pour garantir la capacité cognitive.\\
\vspace{0.5em}
N\\
CO\\
• Standards Audit : Dès la semaine 9, tout code est soumis aux standards des cabinets d’audit (Documentation, Tests, Invariants).\\
\vspace{0.5em}
• Autonomie Radicale : Pas de ”prof” qui corrige. Peer-review et documentation technique sont les seules sources de vérité.\\
\vspace{0.5em}
\vspace{0.5em}
\vspace{0.5em}
\vspace{0.5em}
188\\
RBK 2.0 — Whitepaper Architecte Web3                                                            Chapitre 25. ÉLÉMENTS DE DIFFÉRENCIATION\\
\vspace{0.5em}
\vspace{0.5em}
\begin{center}\begin{tabular}TAB. 25.1 : Grille de Maturité Senior-by-Design\end{tabular}\end{center}
\vspace{0.5em}
\vspace{0.5em}
Axe           Niveau 0 (Junior) Niveau 4 (Senior RBK)          Preuve\\
Architecture Code monolithique Modulaire, Composabilité        Diagramme C4\\
Sécurité      ”Ça marche”       ”C’est incassable”             Threat Model\\
Tests         Manuels           CI/CD, Fuzzing, Property-Based Rapport Coverage\\
Collaboration Solo coder        Reviewer implacable            Historique PR\\
\vspace{0.5em}
\vspace{0.5em}
\vspace{0.5em}
\vspace{0.5em}
Technique          Systémique\\
\vspace{0.5em}
\vspace{0.5em}
\vspace{0.5em}
\vspace{0.5em}
EL\\
Codeur (S0)      Builder (S12)      Architecte (S28)\\
\vspace{0.5em}
\vspace{0.5em}
\vspace{0.5em}
\vspace{0.5em}
TI\\
FIG. 25.1 : Transformation Codeur → Architecte\\
\vspace{0.5em}
\vspace{0.5em}
\vspace{0.5em}
\vspace{0.5em}
EN\\
\subsection{25.2 Approche « Cyborg » : IA-Augmented Engineering}
\vspace{0.5em}
\vspace{0.5em}
\vspace{0.5em}
\vspace{0.5em}
D\\
FI\\
L’IA n’est pas une béquille, c’est un exosquelette. Chez RBK, nous formons des ”Cyborgs” : des ingénieurs qui utilisent l’IA pour multiplier\\
leur productivité par 10, tout en gardant le contrôle absolu sur la qualité et la sécurité.\\
N\\
CO\\
Protocole d’Usage\\
\vspace{0.5em}
• Autorisé : Documentation, boilerplate, génération de tests unitaires, explication d’erreurs.\\
\vspace{0.5em}
• Interdit : Copier-coller de logique métier critique sans audit ligne par ligne.\\
\vspace{0.5em}
• Traçabilité : Tout prompt générant du code prod doit être loggé (Git commit message ou comments).\\
\vspace{0.5em}
\vspace{0.5em}
\vspace{0.5em}
\vspace{0.5em}
\begin{confidential}CONFIDENTIEL — PROJET RBK 2.0                                  Page 189 sur 316\end{confidential}
RBK 2.0 — Whitepaper Architecte Web3                                                             Chapitre 25. ÉLÉMENTS DE DIFFÉRENCIATION\\
\vspace{0.5em}
\vspace{0.5em}
\begin{center}\begin{tabular}TAB. 25.3 : Checklist d’Audit Code IA\end{tabular}\end{center}
\vspace{0.5em}
\vspace{0.5em}
Point de Contrôle Risque IA                       Validation Humaine\\
Logique Invariante Hallucination de règles métier Preuve mathématique\\
Vecteurs d’Attaque Oubli de ”Reentrancy Guard” Analyse statique\\
Edge Cases         Gestion naïve des erreurs      Tests de limites\\
\vspace{0.5em}
\vspace{0.5em}
\vspace{0.5em}
\vspace{0.5em}
\subsection{25.3 Dual Track Solana/EVM : Flexibilité Stratégique}
\vspace{0.5em}
\vspace{0.5em}
\vspace{0.5em}
\vspace{0.5em}
EL\\
Pourquoi choisir ? Le marché valorise la polyvalence. Nos ingénieurs sont ”T-Shaped” : experts profonds sur une stack (ex : Solana) et\\
compétents sur l’autre (EVM). Cela garantit une employabilité maximale et une capacité à auditer des architectures cross-chain.\\
\vspace{0.5em}
\vspace{0.5em}
\vspace{0.5em}
\vspace{0.5em}
TI\\
\begin{center}\begin{tabular}TAB. 25.5 : Comparatif Technique Solana vs EVM\end{tabular}\end{center}
\vspace{0.5em}
\vspace{0.5em}
\vspace{0.5em}
\vspace{0.5em}
D   EN\\
Dimension     Solana (Track A)          EVM (Track B)\\
\vspace{0.5em}
\vspace{0.5em}
\vspace{0.5em}
\vspace{0.5em}
FI\\
Modèle Mental Stateless (Account Model) Stateful (Contract Storage)\\
Langage       Rust + Anchor             Solidity + Foundry\\
Performance\\
NParallélisme (SVM)        Séquentiel (EVM)\\
CO\\
Sécurité      Ownership checks          Reentrancy guards\\
\vspace{0.5em}
\vspace{0.5em}
\vspace{0.5em}
\vspace{0.5em}
\subsection{25.4 Intégration Superteam : Opportunités Directes}
Superteam n’est pas un partenaire, c’est notre client. RBK est conçu comme une usine à talents pour l’écosystème Superteam (Bounties,\\
Grants, Jobs).\\
\vspace{0.5em}
Processus\\
\vspace{0.5em}
\vspace{0.5em}
\vspace{0.5em}
\begin{confidential}CONFIDENTIEL — PROJET RBK 2.0                                   Page 190 sur 316\end{confidential}
RBK 2.0 — Whitepaper Architecte Web3                                                        Chapitre 25. ÉLÉMENTS DE DIFFÉRENCIATION\\
\vspace{0.5em}
\vspace{0.5em}
1. Sourcing : Les meilleurs bounties sont sélectionnés chaque lundi.\\
\vspace{0.5em}
2. Squads : Des équipes de 2-3 étudiants se forment pour attaquer les bounties complexes.\\
\vspace{0.5em}
3. Review RBK : Un mentor senior valide la soumission avant envoi (Quality Gate).\\
\vspace{0.5em}
4. Revenue : 100\textbackslash{}% des gains vont aux étudiants (preuve de concept économique).\\
\vspace{0.5em}
\vspace{0.5em}
\subsection{25.5 « On-Chain Resume » : Preuve de Travail Public}
Le CV PDF est mort. RBK délivre un ”On-Chain Resume” vérifiable cryptographiquement. Chaque compétence validée, chaque projet\\
\vspace{0.5em}
\vspace{0.5em}
\vspace{0.5em}
\vspace{0.5em}
EL\\
livré, chaque audit réalisé est ancré sur la blockchain via des SBT (Soulbound Tokens) et un historique GitHub immuable.\\
\vspace{0.5em}
\begin{center}\begin{tabular}TAB. 25.7 : Structure du On-Chain Resume\end{tabular}\end{center}
\vspace{0.5em}
\vspace{0.5em}
\vspace{0.5em}
\vspace{0.5em}
TI\\
EN\\
Composant Support       Preuve Vérifiable\\
Identité    Wallet      Signature cryptographique\\
\vspace{0.5em}
\vspace{0.5em}
\vspace{0.5em}
\vspace{0.5em}
D\\
Compétences SBT Badge   Transaction on-chain (Issuer : RBK)\\
\vspace{0.5em}
\vspace{0.5em}
\vspace{0.5em}
\vspace{0.5em}
FI\\
Projets     GitHub Repo Commit history, CI logs\\
Réputation\\
N DAO Vote    Poids de vote on-chain\\
CO\\
\vspace{0.5em}
\subsection{25.6 Ancrage Tunisie + Export : Software Factory Future}
RBK positionne la Tunisie comme la ”Base Arrière” de l’ingénierie Web3 mondiale. Moins cher que l’Europe de l’Est, plus qualifié que\\
l’Asie du Sud-Est (sur la niche Rust/Crypto), et sur le même fuseau horaire que Paris/Berlin/Lagos.\\
\vspace{0.5em}
\vspace{0.5em}
\vspace{0.5em}
\vspace{0.5em}
\begin{confidential}CONFIDENTIEL — PROJET RBK 2.0                               Page 191 sur 316\end{confidential}
RBK 2.0 — Whitepaper Architecte Web3                                                              Chapitre 25. ÉLÉMENTS DE DIFFÉRENCIATION\\
\vspace{0.5em}
\vspace{0.5em}
\begin{center}\begin{tabular}TAB. 25.9 : Risk Register Export\end{tabular}\end{center}
\vspace{0.5em}
\vspace{0.5em}
Risque          Prob.   Impact               Mitigation\\
\vspace{0.5em}
Juridique       Moyen Blocage paiements      Contrats types validés, Crypto-payments\\
\vspace{0.5em}
Fuite Talents Haut      Perte expertise locale Modèle ”Remote from Tunisia” (Salaire indexé)\\
\vspace{0.5em}
Qualité         Moyen Perte réputation       QA systématique par Senior RBK\\
\vspace{0.5em}
\vspace{0.5em}
\vspace{0.5em}
\vspace{0.5em}
EL\\
\subsection{25.6.1 Comparatif RBK 2.0 vs Bootcamps Classiques}
RBK n’est pas un bootcamp. C’est un centre d’entraînement olympique pour ingénieurs.\\
\vspace{0.5em}
\vspace{0.5em}
\vspace{0.5em}
\vspace{0.5em}
TI\\
EN\\
\begin{center}\begin{tabular}TAB. 25.11 : Matrice Comparée\end{tabular}\end{center}
\vspace{0.5em}
\vspace{0.5em}
\vspace{0.5em}
\vspace{0.5em}
D\\
Critère    RBK 2.0               Bootcamp Web2 Université\\
\vspace{0.5em}
\vspace{0.5em}
\vspace{0.5em}
\vspace{0.5em}
FI\\
Profondeur Expert (Rust/Systems) Surface (JS/React) Théorique\\
Sécurité   Obsessionnelle        Basique            Abstraite\\
Preuve\\
N\\
Audit Report          ”Projet TodoList” Diplôme Papier\\
CO\\
Modèle Éco ISA (Success fee)     Cash Upfront       Gratuit / Public\\
\vspace{0.5em}
\vspace{0.5em}
\vspace{0.5em}
\vspace{0.5em}
\begin{confidential}CONFIDENTIEL — PROJET RBK 2.0                                   Page 192 sur 316\end{confidential}
\section{CONCLUSION \textbackslash{}& FEUILLE DE ROUTE}
\vspace{0.5em}
\vspace{0.5em}
\subsection{26.1 Priorités Immédiates (Semaine 1–4)}
\vspace{0.5em}
\vspace{0.5em}
\vspace{0.5em}
\vspace{0.5em}
EL\\
Le compte à rebours est lancé. Voici le plan d’attaque pour les 30 premiers jours post-validation de ce Whitepaper.\\
\vspace{0.5em}
\vspace{0.5em}
\vspace{0.5em}
\vspace{0.5em}
TI\\
Table : Plan 4 Semaines\\
\vspace{0.5em}
\vspace{0.5em}
\vspace{0.5em}
\vspace{0.5em}
N\\
DE\\
Semaine Objectif Actions Clés                                 Owner\\
S1      Légal     Validation Contrats ISA + Setup Bancaire    CEO\\
\vspace{0.5em}
\vspace{0.5em}
\vspace{0.5em}
\vspace{0.5em}
FI\\
S2      Tech      Déploiement LMS + Setup Github Org          CTO\\
\vspace{0.5em}
\vspace{0.5em}
N\\
S3      Marketing Lancement Landing Page + Campagne ”Genesis” CMO\\
CO\\
S4      Ops       Ouverture Candidatures (Piscine Beta)       Ops\\
\vspace{0.5em}
\vspace{0.5em}
\vspace{0.5em}
\subsection{26.2 KPI de Succès}
Nous ne pilotons pas à vue. 12 indicateurs clés définissent la santé du projet.\\
\vspace{0.5em}
Table : KPI Dictionary (Extrait)\\
\vspace{0.5em}
KPI         Définition                Cible S12 Seuil Alerte\\
Selectivity \textbackslash{}% Candidats admis piscine \textless{} 10\textbackslash{}%     \textgreater{} 20\textbackslash{}%\\
\vspace{0.5em}
\vspace{0.5em}
\vspace{0.5em}
\vspace{0.5em}
193\\
RBK 2.0 — Whitepaper Architecte Web3                                                       Chapitre 26. CONCLUSION \textbackslash{}& FEUILLE DE ROUTE\\
\vspace{0.5em}
\vspace{0.5em}
\vspace{0.5em}
Attrition \textbackslash{}% Dropout durant piscine \textless{} 30\textbackslash{}%        \textgreater{} 50\textbackslash{}%\\
Job Ready \textbackslash{}% Certifiés ”Audit-Ready” \textgreater{} 80\textbackslash{}%       \textless{} 60\textbackslash{}%\\
Placement \textbackslash{}% en poste à J+90         \textgreater{} 70\textbackslash{}%       \textless{} 50\textbackslash{}%\\
\vspace{0.5em}
\vspace{0.5em}
\vspace{0.5em}
\subsection{26.3 Engagement Qualité Formel}
RBK s’engage sur une politique ”Zéro Complaisance”.\\
\vspace{0.5em}
\vspace{0.5em}
\vspace{0.5em}
\vspace{0.5em}
EL\\
• Pas de diplôme de complaisance : Si le niveau n’est pas atteint, l’étudiant double ou sort.\\
\vspace{0.5em}
\vspace{0.5em}
\vspace{0.5em}
\vspace{0.5em}
TI\\
• Code Review systématique : Aucun code ne part en prod (ou validation) sans review par un pair et un mentor.\\
\vspace{0.5em}
• Transparence totale : Les statistiques de placement et de salaire sont publiées et auditées.\\
\vspace{0.5em}
\vspace{0.5em}
\vspace{0.5em}
\vspace{0.5em}
EN\\
\subsection{26.4 Forge de l’Élite Africaine}
\vspace{0.5em}
\vspace{0.5em}
\vspace{0.5em}
\vspace{0.5em}
D\\
RBK a l’ambition de devenir le ”MIT du Web3” pour l’Afrique. Nous ne formons pas des exécutants bon marché, mais l’élite technolo-\\
\vspace{0.5em}
\vspace{0.5em}
FI\\
gique qui construira l’infrastructure financière souveraine du continent.\\
\vspace{0.5em}
[Schéma : Flywheel RBK]\\
N\\
(Sélection → Formation → Preuves → Revenus → Réputation → Sélection)\\
O\\
C\\
\vspace{0.5em}
\vspace{0.5em}
\subsection{26.5 Synthèse Valeur Stratégique}
• Pour l’Étudiant : Une carrière internationale à haute valeur ajoutée, sans dette initiale (ISA).\\
\vspace{0.5em}
• Pour l’Écosystème : Un pipeline fiable de talents ”Audit-Ready”.\\
\vspace{0.5em}
• Pour la Tunisie : Une entrée de devises forte et une montée en gamme technologique.\\
\vspace{0.5em}
\vspace{0.5em}
\vspace{0.5em}
\begin{confidential}CONFIDENTIEL — PROJET RBK 2.0                                Page 194 sur 316\end{confidential}
RBK 2.0 — Whitepaper Architecte Web3                                                      Chapitre 26. CONCLUSION \textbackslash{}& FEUILLE DE ROUTE\\
\vspace{0.5em}
\vspace{0.5em}
\vspace{0.5em}
\subsection{26.6 Synthèse Investisseur – Différenciateurs Nexus}
• Capex maîtrisé : Modèle pédagogique modulaire qui recycle 15 900 TND sur 48 semaines et active le forfait 3 300 TND uniquement\\
pour les projets Launchpad validés.\\
\vspace{0.5em}
• Cashflow récurrent : Mix Upfront + ISA (30\textbackslash{}% sur les revenus Launchpad) couvrant les coûts mentors et finançant la croissance\\
sans dilution.\\
\vspace{0.5em}
• Golden pipeline : Génération systématique d’actifs NFT Skills non transférables (Genesis Access, Build Master, Launch Proof) qui\\
sécurisent la réputation et les taux de conversion Genesis →Explorer →Bâtisseur.\\
\vspace{0.5em}
• Synergie Validator Track : Infrastructure Solana opérée en interne, monétisable auprès des partenaires (staking services, data) tout\\
\vspace{0.5em}
\vspace{0.5em}
\vspace{0.5em}
\vspace{0.5em}
EL\\
en renforçant le moat technologique.\\
\vspace{0.5em}
\vspace{0.5em}
\vspace{0.5em}
\vspace{0.5em}
TI\\
• Scalabilité continentale : Playbooks Genesis/Explorer clés en main, duplicables sur de nouveaux hubs via franchise/licence, avec\\
\vspace{0.5em}
\vspace{0.5em}
\vspace{0.5em}
\vspace{0.5em}
EN\\
Nexus en hub d’orchestration.\\
\vspace{0.5em}
\vspace{0.5em}
\vspace{0.5em}
\vspace{0.5em}
D\\
\subsection{26.7 Appel à l’Action}
\vspace{0.5em}
\vspace{0.5em}
\vspace{0.5em}
FI\\
Le marché n’attend pas. La fenêtre d’opportunité Solana/Rust est ouverte maintenant. Rejoignez la Cohorte Genesis.\\
N\\
CO\\
Next Steps\\
\vspace{0.5em}
• [J0] Validation Finale Whitepaper.\\
\vspace{0.5em}
• [J+7] Lancement Recrutement Core Team.\\
\vspace{0.5em}
• [J+30] Ouverture des Candidatures.\\
\vspace{0.5em}
\vspace{0.5em}
\vspace{0.5em}
\subsection{26.8 Message Final au CEO}
\vspace{0.5em}
\vspace{0.5em}
\vspace{0.5em}
\begin{confidential}CONFIDENTIEL — PROJET RBK 2.0                               Page 195 sur 316\end{confidential}
RBK 2.0 — Whitepaper Architecte Web3                                                      Chapitre 26. CONCLUSION \textbackslash{}& FEUILLE DE ROUTE\\
\vspace{0.5em}
\vspace{0.5em}
\vspace{0.5em}
Monsieur le CEO, Ce plan est ambitieux, risqué, mais nécessaire. Il transforme RBK d’un centre de formation classique en une Startup\\
Studio Éducative. Le modèle économique est viable (ISA + Bounties). La demande marché est validée. La technologie est mature. Il ne\\
reste qu’une variable : l’Exécution. C’est un GO.\\
\vspace{0.5em}
\vspace{0.5em}
\subsection{26.9 Profil de Sortie}
Table : Profil de Sortie Standard\\
Compétence Preuve                  Seuil\\
Rust / Solana 3 Repos GitHub Clean CI Green\\
\vspace{0.5em}
\vspace{0.5em}
\vspace{0.5em}
\vspace{0.5em}
EL\\
Sécurité      1 Rapport d’Audit    3 vulns trouvées\\
Soft Skills   Démo Vidéo           Clarté \textgreater{} 4/5\\
\vspace{0.5em}
\vspace{0.5em}
\vspace{0.5em}
\vspace{0.5em}
TI\\
D    EN\\
FI\\
N\\
CO\\
\vspace{0.5em}
\vspace{0.5em}
\vspace{0.5em}
\vspace{0.5em}
\begin{confidential}CONFIDENTIEL — PROJET RBK 2.0                               Page 196 sur 316\end{confidential}
Table des figures\\
\vspace{0.5em}
\vspace{0.5em}
\vspace{0.5em}
\vspace{0.5em}
EL\\
\subsection{1.1 La Chaîne de Valeur RBK 2.0 . . . . . . . . . . . . . . . . . . . . . . . . . . . . . . . . . . . . . . . . . . . . . . . . . .     34}
\vspace{0.5em}
\vspace{0.5em}
\vspace{0.5em}
\vspace{0.5em}
TI\\
\subsection{4.1 Corrélation talents Solana vs traction financière (2021-2026) . . . . . . . . . . . . . . . . . . . . . . . . . . . . . . . . .     57}
\vspace{0.5em}
\vspace{0.5em}
\vspace{0.5em}
\vspace{0.5em}
EN\\
\subsection{5.1 Flux décisionnel RBK → Nexus en cas d’alerte . . . . . . . . . . . . . . . . . . . . . . . . . . . . . . . . . . . . . . . . .      67}
\vspace{0.5em}
\vspace{0.5em}
\vspace{0.5em}
\vspace{0.5em}
D\\
\subsection{8.1 Timeline Macro du Cursus . . . . . . . . . . . . . . . . . . . . . . . . . . . . . . . . . . . . . . . . . . . . . . . . . . .      80}
\vspace{0.5em}
\vspace{0.5em}
\vspace{0.5em}
\vspace{0.5em}
FI\\
\subsection{9.1 Pourquoi Solana est un track d’excellence . . . . . . . . . . . . . . . . . . . . . . . . . . . . . . . . . . . . . . . . . . .     90}
\vspace{0.5em}
N\\
\subsection{9.2 Flux Anchor . . . . . . . . . . . . . . . . . . . . . . . . . . . . . . . . . . . . . . . . . . . . . . . . . . . . . . . . . . .   93}
CO\\
\subsection{10.1 Chaîne de Valeur EVM . . . . . . . . . . . . . . . . . . . . . . . . . . . . . . . . . . . . . . . . . . . . . . . . . . . . .     98}
\vspace{0.5em}
\subsection{12.1 Timeline 4 semaines — Soft Skills \textbackslash{}& Pro . . . . . . . . . . . . . . . . . . . . . . . . . . . . . . . . . . . . . . . . . . . . 107}
\vspace{0.5em}
\subsection{13.1 Pipeline Packaging . . . . . . . . . . . . . . . . . . . . . . . . . . . . . . . . . . . . . . . . . . . . . . . . . . . . . . . 114}
\vspace{0.5em}
\subsection{14.1 Boucle Guardian (SecDevOps) . . . . . . . . . . . . . . . . . . . . . . . . . . . . . . . . . . . . . . . . . . . . . . . . . 117}
\subsection{14.2 Boucle Visionnaire (Design → Simulation → Risques → Déploiement) . . . . . . . . . . . . . . . . . . . . . . . . . . . . 120}
\subsection{14.3 Boucle Builder (Produit → Release → Mesure → Itération) . . . . . . . . . . . . . . . . . . . . . . . . . . . . . . . . . . 122}
\subsection{14.4 Machine à états transactionnelle (dApp Engineer) . . . . . . . . . . . . . . . . . . . . . . . . . . . . . . . . . . . . . . . 124}
\vspace{0.5em}
\vspace{0.5em}
\vspace{0.5em}
197\\
RBK 2.0 — Whitepaper Architecte Web3                                                                                        Table des figures\\
\vspace{0.5em}
\vspace{0.5em}
\subsection{14.5 Cycle de vie tokenisé + Audit Trail (Tokenization \textbackslash{}& DePIN Architect) . . . . . . . . . . . . . . . . . . . . . . . . . . . . 126}
\vspace{0.5em}
\subsection{15.1 Mix Revenus Cible (Année 3) . . . . . . . . . . . . . . . . . . . . . . . . . . . . . . . . . . . . . . . . . . . . . . . . . . 132}
\subsection{15.2 Fonds de Garantie ISA (Cash Collateral) . . . . . . . . . . . . . . . . . . . . . . . . . . . . . . . . . . . . . . . . . . . . 140}
\subsection{15.3 Modèle de Revenus Multi-Couches . . . . . . . . . . . . . . . . . . . . . . . . . . . . . . . . . . . . . . . . . . . . . . . 141}
\vspace{0.5em}
\subsection{16.1 Maquette LinkedIn (placeholder visuel) — livrable D-21 avant lancement . . . . . . . . . . . . . . . . . . . . . . . . . . 143}
\subsection{16.2 Wireframe landing page (placeholder) — référentiel product marketing . . . . . . . . . . . . . . . . . . . . . . . . . . . 143}
\subsection{16.3 Storyboard vidéo (placeholder) — diffusion multi-canale . . . . . . . . . . . . . . . . . . . . . . . . . . . . . . . . . . . 143}
\subsection{16.4 Flywheel Building in Public . . . . . . . . . . . . . . . . . . . . . . . . . . . . . . . . . . . . . . . . . . . . . . . . . . 144}
\subsection{16.5 Funnel d’Acquisition Simplifié . . . . . . . . . . . . . . . . . . . . . . . . . . . . . . . . . . . . . . . . . . . . . . . . . 149}
\vspace{0.5em}
\vspace{0.5em}
\vspace{0.5em}
\vspace{0.5em}
EL\\
\subsection{21.1 Workflow d’Émission Automatisé      . . . . . . . . . . . . . . . . . . . . . . . . . . . . . . . . . . . . . . . . . . . . . . . 171}
\vspace{0.5em}
\vspace{0.5em}
\vspace{0.5em}
\vspace{0.5em}
TI\\
\subsection{22.1 Timeline 120 jours (Vue Exécutive) — Sprint 0 inclus . . . . . . . . . . . . . . . . . . . . . . . . . . . . . . . . . . . . . 175}
\vspace{0.5em}
\vspace{0.5em}
\vspace{0.5em}
\vspace{0.5em}
EN\\
\subsection{22.2 Gantt Macro aligné sur les phases Sprint 0 \textbackslash{}& Cohorte Alpha . . . . . . . . . . . . . . . . . . . . . . . . . . . . . . . . . . 178}
\vspace{0.5em}
\subsection{24.1 Architecture d’Émission SBT . . . . . . . . . . . . . . . . . . . . . . . . . . . . . . . . . . . . . . . . . . . . . . . . . . 184}
\vspace{0.5em}
\vspace{0.5em}
\vspace{0.5em}
\vspace{0.5em}
D\\
\subsection{24.2 Pyramide des Tiers Alumni . . . . . . . . . . . . . . . . . . . . . . . . . . . . . . . . . . . . . . . . . . . . . . . . . . . 187}
\vspace{0.5em}
\vspace{0.5em}
\vspace{0.5em}
\vspace{0.5em}
FI\\
N\\
\subsection{25.1 Transformation Codeur → Architecte . . . . . . . . . . . . . . . . . . . . . . . . . . . . . . . . . . . . . . . . . . . . . 189}
CO\\
K.1 Value Ladder et Parcours Étudiant . . . . . . . . . . . . . . . . . . . . . . . . . . . . . . . . . . . . . . . . . . . . . . . 254\\
\vspace{0.5em}
\vspace{0.5em}
\vspace{0.5em}
\vspace{0.5em}
\begin{confidential}CONFIDENTIEL — PROJET RBK 2.0                                Page 198 sur 316\end{confidential}
Liste des tableaux\\
\vspace{0.5em}
\vspace{0.5em}
\vspace{0.5em}
\vspace{0.5em}
EL\\
\subsection{1.2 Métriques de Succès RBK 2.0 . . . . . . . . . . . . . . . . . . . . . . . . . . . . . . . . . . . . . . . . . . . . . . . . . .     33}
\vspace{0.5em}
\vspace{0.5em}
\vspace{0.5em}
\vspace{0.5em}
TI\\
\subsection{1.5 Le Changement de Paradigme RBK 2.0 (Détaillé) . . . . . . . . . . . . . . . . . . . . . . . . . . . . . . . . . . . . . . .         39}
\vspace{0.5em}
\vspace{0.5em}
\vspace{0.5em}
\vspace{0.5em}
EN\\
\subsection{4.1 Segmentation des Rôles Web3 (2025) . . . . . . . . . . . . . . . . . . . . . . . . . . . . . . . . . . . . . . . . . . . . .        53}
\subsection{4.5 Grille Salariale Web3 (Remote Global) vs Local . . . . . . . . . . . . . . . . . . . . . . . . . . . . . . . . . . . . . . . .      59}
\vspace{0.5em}
\vspace{0.5em}
\vspace{0.5em}
\vspace{0.5em}
D\\
\subsection{6.1   Architecture pédagogique alignée sur le modèle 42 adapté RBK . . . . . . . . . . . . . . . . . . . . . . . . . . . . . . . .      70}
\vspace{0.5em}
\vspace{0.5em}
\vspace{0.5em}
\vspace{0.5em}
FI\\
\subsection{6.3   Mesures Anti-Burnout Intégrées dans le Cycle Hebdomadaire . . . . . . . . . . . . . . . . . . . . . . . . . . . . . . . . .       71}
\subsection{6.5}
N\\
Compétences et évaluations ciblées durant la Genesis Pool . . . . . . . . . . . . . . . . . . . . . . . . . . . . . . . . . .     72\\
CO\\
\subsection{6.7   Correspondance NFT Skills et responsabilités pédagogiques . . . . . . . . . . . . . . . . . . . . . . . . . . . . . . . . . .     73}
\vspace{0.5em}
\subsection{7.1 Modules Détaillés du Track C - Web3 Product \textbackslash{}& Ecosystem Strategy . . . . . . . . . . . . . . . . . . . . . . . . . . . . .          78}
\subsection{7.3 Timeline pédagogique détaillée (Genesis Pool →Launchpad) . . . . . . . . . . . . . . . . . . . . . . . . . . . . . . . . .          79}
\vspace{0.5em}
\subsection{8.1 Genesis Pool – Synthèse des 4 semaines . . . . . . . . . . . . . . . . . . . . . . . . . . . . . . . . . . . . . . . . . . . .      81}
\vspace{0.5em}
\subsection{9.1   Compétences Cibles vs Preuves . . . . . . . . . . . . . . . . . . . . . . . . . . . . . . . . . . . . . . . . . . . . . . . . .   90}
\subsection{9.3   Carte des Modules (Résumé Exécutif) . . . . . . . . . . . . . . . . . . . . . . . . . . . . . . . . . . . . . . . . . . . . .     91}
\subsection{9.5   Checklist Sécurité Module 1 . . . . . . . . . . . . . . . . . . . . . . . . . . . . . . . . . . . . . . . . . . . . . . . . . .   92}
\subsection{9.7   Production Readiness Review (PRR) . . . . . . . . . . . . . . . . . . . . . . . . . . . . . . . . . . . . . . . . . . . . . .     94}
\vspace{0.5em}
\vspace{0.5em}
\vspace{0.5em}
199\\
RBK 2.0 — Whitepaper Architecte Web3                                                                                          Liste des tableaux\\
\vspace{0.5em}
\vspace{0.5em}
\subsection{9.9 Stack Track A (Standard) . . . . . . . . . . . . . . . . . . . . . . . . . . . . . . . . . . . . . . . . . . . . . . . . . . . .      94}
\subsection{9.11 Itinéraire Validator (6 semaines) . . . . . . . . . . . . . . . . . . . . . . . . . . . . . . . . . . . . . . . . . . . . . . . .    95}
\subsection{9.13 Checklist Portfolio Guardian . . . . . . . . . . . . . . . . . . . . . . . . . . . . . . . . . . . . . . . . . . . . . . . . . .     96}
\vspace{0.5em}
\subsection{10.1 Carte des Modules Track B . . . . . . . . . . . . . . . . . . . . . . . . . . . . . . . . . . . . . . . . . . . . . . . . . . . 98}
\subsection{10.3 Security Checklist EVM . . . . . . . . . . . . . . . . . . . . . . . . . . . . . . . . . . . . . . . . . . . . . . . . . . . . . 100}
\subsection{10.5 Stack Track B (Foundry) . . . . . . . . . . . . . . . . . . . . . . . . . . . . . . . . . . . . . . . . . . . . . . . . . . . . 100}
\subsection{10.7 Matrice Compétences Infra EVM . . . . . . . . . . . . . . . . . . . . . . . . . . . . . . . . . . . . . . . . . . . . . . . . 101}
\vspace{0.5em}
\subsection{11.3 Matrice Compétences Product . . . . . . . . . . . . . . . . . . . . . . . . . . . . . . . . . . . . . . . . . . . . . . . . . 105}
\vspace{0.5em}
\subsection{12.1 Vue d’ensemble du module (4 semaines) . . . . . . . . . . . . . . . . . . . . . . . . . . . . . . . . . . . . . . . . . . . . 107}
\vspace{0.5em}
\vspace{0.5em}
\vspace{0.5em}
\vspace{0.5em}
EL\\
\subsection{12.3 Rubrique d’Évaluation des Soft Skills . . . . . . . . . . . . . . . . . . . . . . . . . . . . . . . . . . . . . . . . . . . . . 109}
\vspace{0.5em}
\vspace{0.5em}
\vspace{0.5em}
\vspace{0.5em}
TI\\
\subsection{13.1 Studio-Grade Checklist (Non-Négociable) . . . . . . . . . . . . . . . . . . . . . . . . . . . . . . . . . . . . . . . . . . . 110}
\vspace{0.5em}
\vspace{0.5em}
\vspace{0.5em}
\vspace{0.5em}
EN\\
\subsection{13.2 Rubric d’Évaluation Studio (100 pts) . . . . . . . . . . . . . . . . . . . . . . . . . . . . . . . . . . . . . . . . . . . . . . 114}
\vspace{0.5em}
\subsection{14.7 Revenus Indicatifs (2025)      . . . . . . . . . . . . . . . . . . . . . . . . . . . . . . . . . . . . . . . . . . . . . . . . . . . 129}
\vspace{0.5em}
\vspace{0.5em}
\vspace{0.5em}
\vspace{0.5em}
D\\
FI\\
\subsection{15.2 Projection financière — scénario réaliste (TND) . . . . . . . . . . . . . . . . . . . . . . . . . . . . . . . . . . . . . . . . 137}
\vspace{0.5em}
N\\
\subsection{15.4 Analyse de sensibilité sur le scénario réaliste . . . . . . . . . . . . . . . . . . . . . . . . . . . . . . . . . . . . . . . . . 138}
CO\\
\subsection{16.1 Calendrier Éditorial Type (Cycle 12 Semaines) . . . . . . . . . . . . . . . . . . . . . . . . . . . . . . . . . . . . . . . . . 144}
\subsection{16.3 Roadmap éditoriale Sprint 0 (90 jours détaillés) . . . . . . . . . . . . . . . . . . . . . . . . . . . . . . . . . . . . . . . . 145}
\subsection{16.5 ROI Comparatif par Option (Sortie Junior : 2 500 TND brut/mois) . . . . . . . . . . . . . . . . . . . . . . . . . . . . . . 146}
\subsection{16.7 Allocation budgétaire et objectifs CAC . . . . . . . . . . . . . . . . . . . . . . . . . . . . . . . . . . . . . . . . . . . . . 147}
\subsection{16.9 Catalogue des Incentives . . . . . . . . . . . . . . . . . . . . . . . . . . . . . . . . . . . . . . . . . . . . . . . . . . . . 150}
\vspace{0.5em}
\subsection{17.1 Matrice des Risques — Référentiel Commun (RBK – Nexus Réussite) . . . . . . . . . . . . . . . . . . . . . . . . . . . . . 153}
\subsection{17.3 Nouvelle Matrice des Risques (Version 2026) . . . . . . . . . . . . . . . . . . . . . . . . . . . . . . . . . . . . . . . . . 154}
\subsection{17.5 Déclencheurs et plans de continuité . . . . . . . . . . . . . . . . . . . . . . . . . . . . . . . . . . . . . . . . . . . . . . 155}
\subsection{17.7 Top 5 Risques et Mitigations (2026) . . . . . . . . . . . . . . . . . . . . . . . . . . . . . . . . . . . . . . . . . . . . . . 157}
\vspace{0.5em}
\vspace{0.5em}
\vspace{0.5em}
\begin{confidential}CONFIDENTIEL — PROJET RBK 2.0                                   Page 200 sur 316\end{confidential}
RBK 2.0 — Whitepaper Architecte Web3                                                                                         Liste des tableaux\\
\vspace{0.5em}
\vspace{0.5em}
\subsection{18.3 Risques Compliance \textbackslash{}& Atténuation . . . . . . . . . . . . . . . . . . . . . . . . . . . . . . . . . . . . . . . . . . . . . . . 161}
\vspace{0.5em}
\subsection{22.1 Checklist Go/No-Go (Gates) . . . . . . . . . . . . . . . . . . . . . . . . . . . . . . . . . . . . . . . . . . . . . . . . . . 175}
\subsection{22.3 Backlog Opérationnel (Extrait Top Actions) . . . . . . . . . . . . . . . . . . . . . . . . . . . . . . . . . . . . . . . . . . 177}
\vspace{0.5em}
\subsection{23.1 Jalons Clés du Plan de Lancement . . . . . . . . . . . . . . . . . . . . . . . . . . . . . . . . . . . . . . . . . . . . . . . 182}
\vspace{0.5em}
\subsection{24.1 Catalogue des Badges SBT (Extrait) . . . . . . . . . . . . . . . . . . . . . . . . . . . . . . . . . . . . . . . . . . . . . . 184}
\subsection{24.3 Usages et Bénéfices des SBT . . . . . . . . . . . . . . . . . . . . . . . . . . . . . . . . . . . . . . . . . . . . . . . . . . 186}
\subsection{24.5 Roadmap Alumni (Année 1) . . . . . . . . . . . . . . . . . . . . . . . . . . . . . . . . . . . . . . . . . . . . . . . . . . 187}
\vspace{0.5em}
\subsection{25.1 Grille de Maturité Senior-by-Design . . . . . . . . . . . . . . . . . . . . . . . . . . . . . . . . . . . . . . . . . . . . . . 189}
\vspace{0.5em}
\vspace{0.5em}
\vspace{0.5em}
\vspace{0.5em}
EL\\
\subsection{25.3 Checklist d’Audit Code IA . . . . . . . . . . . . . . . . . . . . . . . . . . . . . . . . . . . . . . . . . . . . . . . . . . . 190}
\subsection{25.5 Comparatif Technique Solana vs EVM . . . . . . . . . . . . . . . . . . . . . . . . . . . . . . . . . . . . . . . . . . . . . 190}
\vspace{0.5em}
\vspace{0.5em}
\vspace{0.5em}
\vspace{0.5em}
TI\\
\subsection{25.7 Structure du On-Chain Resume . . . . . . . . . . . . . . . . . . . . . . . . . . . . . . . . . . . . . . . . . . . . . . . . . 191}
\vspace{0.5em}
\vspace{0.5em}
\vspace{0.5em}
\vspace{0.5em}
EN\\
\subsection{25.9 Risk Register Export . . . . . . . . . . . . . . . . . . . . . . . . . . . . . . . . . . . . . . . . . . . . . . . . . . . . . . 192}
\subsection{25.11Matrice Comparée . . . . . . . . . . . . . . . . . . . . . . . . . . . . . . . . . . . . . . . . . . . . . . . . . . . . . . . 192}
\vspace{0.5em}
\vspace{0.5em}
\vspace{0.5em}
\vspace{0.5em}
D\\
C.1 Conversion locale des principaux montants . . . . . . . . . . . . . . . . . . . . . . . . . . . . . . . . . . . . . . . . . . 230\\
\vspace{0.5em}
\vspace{0.5em}
\vspace{0.5em}
\vspace{0.5em}
FI\\
C.3 Projections financières selon taille de cohorte initiale . . . . . . . . . . . . . . . . . . . . . . . . . . . . . . . . . . . . . 233\\
\vspace{0.5em}
N\\
CO\\
G.1 Scénarios de Remboursement        . . . . . . . . . . . . . . . . . . . . . . . . . . . . . . . . . . . . . . . . . . . . . . . . . 246\\
\vspace{0.5em}
H.1 Critères de sélection Piscine Rust . . . . . . . . . . . . . . . . . . . . . . . . . . . . . . . . . . . . . . . . . . . . . . . . 247\\
H.3 Barème Admission Explorer . . . . . . . . . . . . . . . . . . . . . . . . . . . . . . . . . . . . . . . . . . . . . . . . . . 248\\
\vspace{0.5em}
K.2 Structure tarifaire alignée sur le pipeline Genesis → Architecte . . . . . . . . . . . . . . . . . . . . . . . . . . . . . . . . 253\\
\vspace{0.5em}
L.1 Grille de Rémunération Mentor (Junior → Lead) . . . . . . . . . . . . . . . . . . . . . . . . . . . . . . . . . . . . . . . 257\\
\vspace{0.5em}
M.1 Tarification B2B    . . . . . . . . . . . . . . . . . . . . . . . . . . . . . . . . . . . . . . . . . . . . . . . . . . . . . . . . 259\\
\vspace{0.5em}
N.1 Matrice des Outils IA Autorisés . . . . . . . . . . . . . . . . . . . . . . . . . . . . . . . . . . . . . . . . . . . . . . . . . 262\\
\vspace{0.5em}
\vspace{0.5em}
\vspace{0.5em}
\begin{confidential}CONFIDENTIEL — PROJET RBK 2.0                                 Page 201 sur 316\end{confidential}
RBK 2.0 — Whitepaper Architecte Web3                                                                                      Liste des tableaux\\
\vspace{0.5em}
\vspace{0.5em}
O.1 Mapping Détaillé Compétences / Badges / Preuves . . . . . . . . . . . . . . . . . . . . . . . . . . . . . . . . . . . . . . 263\\
\vspace{0.5em}
U.1 Registre de risques v1 — RBK Web3 Studio (DualTrack EVM/Solana)          . . . . . . . . . . . . . . . . . . . . . . . . . . . . 276\\
\vspace{0.5em}
V.1   Registre des actions risques v1 — RBK Web3 Studio . . . . . . . . . . . . . . . . . . . . . . . . . . . . . . . . . . . . . . 282\\
V.2   KPI de pilotage v1 — RBK Web3 Studio . . . . . . . . . . . . . . . . . . . . . . . . . . . . . . . . . . . . . . . . . . . . 289\\
V.8   Journal d’incidents v1 — RBK Web3 Studio . . . . . . . . . . . . . . . . . . . . . . . . . . . . . . . . . . . . . . . . . . 300\\
V.9   Annexe X — Traceability Matrix (SSOT \textbackslash{}& Release Gate) . . . . . . . . . . . . . . . . . . . . . . . . . . . . . . . . . . . . 302\\
\vspace{0.5em}
\vspace{0.5em}
\vspace{0.5em}
\vspace{0.5em}
EL\\
TI\\
D    EN\\
FI\\
N\\
CO\\
\vspace{0.5em}
\vspace{0.5em}
\vspace{0.5em}
\vspace{0.5em}
\begin{confidential}CONFIDENTIEL — PROJET RBK 2.0                                Page 202 sur 316\end{confidential}
A \textbar{} Gabarits opérationnels \textbackslash{}& stratégiques\\
\vspace{0.5em}
\vspace{0.5em}
Cette annexe présente les matrices stratégiques validées pour le déploiement de RBK 2.0.\\
\vspace{0.5em}
\vspace{0.5em}
\vspace{0.5em}
\vspace{0.5em}
EL\\
Légende des rôles (RACI).\\
\vspace{0.5em}
\vspace{0.5em}
\vspace{0.5em}
\vspace{0.5em}
TI\\
• CEO : Direction RBK (décisions stratégiques / arbitrage).\\
\vspace{0.5em}
\vspace{0.5em}
\vspace{0.5em}
\vspace{0.5em}
EN\\
• HoP : Head of Program / Directeur pédagogique (qualité, curriculum, delivery).\\
\vspace{0.5em}
\vspace{0.5em}
\vspace{0.5em}
\vspace{0.5em}
D\\
• TechLead-EVM : Référent technique EVM (Solidity, tooling, patterns, sécurité).\\
\vspace{0.5em}
\vspace{0.5em}
\vspace{0.5em}
\vspace{0.5em}
FI\\
• TechLead-SOL : Référent technique Solana (Rust/Anchor, tooling, sécurité).\\
N\\
CO\\
• SecLead : Référent sécurité/audit (méthodologie, checklists, hardening).\\
\vspace{0.5em}
• Ops : Opérations (planning, salles, outils, LMS, comptes, supports).\\
\vspace{0.5em}
• Career : Career services (coaching, portfolio, placements, partenariats RH).\\
\vspace{0.5em}
• Mkt : Marketing \textbackslash{}& admissions (funnel, contenus, événements, conversion).\\
\vspace{0.5em}
• Legal : Conseil conformité (disclaimer, scénarios opératoires, risques).\\
\vspace{0.5em}
• Mentors : intervenants/mentors externes (reviews, office hours, panels).\\
\vspace{0.5em}
• Students : apprenants (delivery, repo, docs, demo, éthique).\\
\vspace{0.5em}
\vspace{0.5em}
203\\
RBK 2.0 — Whitepaper Architecte Web3                                                           Annexe A. Gabarits opérationnels \textbackslash{}& stratégiques\\
\vspace{0.5em}
\vspace{0.5em}
\vspace{0.5em}
A.1     Matrice SWOT (Forces, Faiblesses, Opportunités, Menaces)\\
SWOT — Analyse Stratégique RBK 2.0\\
\vspace{0.5em}
\vspace{0.5em}
Forces                                                                Faiblesses\\
• Positionnement premium : “Senior-by-Design” (engineering,      • Risque de sur-ambition : contenu trop large si non modularisé\\
sécurité, production).                                           (fatigue, dilution).\\
• DualTrack EVM/Solana : double crédibilité et employabilité • Dépendance à des experts rares (Rust/Anchor, audit), risque\\
internationale.                                                  de disponibilité.\\
• Méthodologie Studio : sprints, PR reviews, CI, tests, incident • Exigence élevée : peut réduire le volume d’inscrits si admis-\\
\vspace{0.5em}
\vspace{0.5em}
\vspace{0.5em}
\vspace{0.5em}
EL\\
drills, demo days.                                               sions trop strictes ou discours mal cadré.\\
\vspace{0.5em}
\vspace{0.5em}
\vspace{0.5em}
\vspace{0.5em}
TI\\
• Production d’un portfolio vérifiable : repos, releases, docs,  • Nécessité d’une cohérence chiffrée stricte (durées, prix, KPI)\\
rubrics, capstones.                                              pour crédibilité business/investisseurs.\\
\vspace{0.5em}
\vspace{0.5em}
\vspace{0.5em}
\vspace{0.5em}
EN\\
• Adossement possible à un réseau mentors (diaspora, builders, • Contexte local : sensibilité réglementaire autour des actifs\\
écosystèmes).                                                    numériques (communication à cadrer).\\
\vspace{0.5em}
\vspace{0.5em}
\vspace{0.5em}
\vspace{0.5em}
D\\
Opportunités                                                          Menaces\\
\vspace{0.5em}
\vspace{0.5em}
\vspace{0.5em}
\vspace{0.5em}
FI\\
• Marché remote Web3 : opportunités globales (teams distri-           • Volatilité Web3 : cycles de marché \textbackslash{}& narratives (risque sur\\
buées, rémunérations supérieures).              N                     marketing et perception).\\
CO\\
• Demande croissante pour profils sécurité / audit-readiness          • Risque réputationnel : confusion “formation dev” vs “pro-\\
(IA = plus de vulnérabilités).                                        messe crypto” si branding imprécis.\\
• Partenariats (protocols, infra, wallets, analytics) : crédibilité   • Risques sécurité : un capstone mal cadré peut exposer à des\\
+ projets réels + recrutement.                                        mauvaises pratiques (à prévenir).\\
• RBK peut devenir un hub régional (Afrique du Nord / franco-         • Concurrence MOOC/bootcamps internationaux : différencia-\\
phonie) avec cohorte pilote forte.                                    tion doit être preuves + mentoring + studio.\\
• Monétisation additionnelle : B2B (formations entreprise), stu-      • Réglementaire : incertitudes locales (paiements, communica-\\
dio services, incubation.                                             tions), besoin d’un modèle opératoire compliant.\\
\vspace{0.5em}
\vspace{0.5em}
\vspace{0.5em}
\vspace{0.5em}
\begin{confidential}CONFIDENTIEL — PROJET RBK 2.0                                  Page 204 sur 316\end{confidential}
RBK 2.0 — Whitepaper Architecte Web3                                                                Annexe A. Gabarits opérationnels \textbackslash{}& stratégiques\\
\vspace{0.5em}
\vspace{0.5em}
\vspace{0.5em}
A.2     Priorisation MoSCoW (Must, Should,\\
Could, Won’t)\\
Priorisation MoSCoW — Partie 1 (Critique)\\
\vspace{0.5em}
\vspace{0.5em}
Must have (Vital)                                                       Should have (Important)\\
• Programme DualTrack : choix EVM ou Solana avec tronc commun           • Micro-certifications (Badges) : critères vérifiables (Code + Demo).\\
+ spécialisation.                                                       • Réseau mentors : Office Hours + Reviews + Panels.\\
• Zéro réduction de contenu : restructuration sans suppression.         • Outillage standardisé : Templates PR, ADR, Runbooks.\\
\vspace{0.5em}
\vspace{0.5em}
\vspace{0.5em}
\vspace{0.5em}
EL\\
• Méthode Studio obligatoire : PR reviews, CI, tests minimaux,          • Listes automatiques : Figures, Tableaux, Acronymes.\\
conventions repo.                                                       • Sources \textbackslash{}& Hypothèses : Bibliographie, Benchmark marché, ROI.\\
\vspace{0.5em}
\vspace{0.5em}
\vspace{0.5em}
\vspace{0.5em}
TI\\
• Capstones (3) + 1 projet final : livrables, critères d’acceptation,\\
notation.\\
\vspace{0.5em}
\vspace{0.5em}
\vspace{0.5em}
\vspace{0.5em}
EN\\
• Note de cadrage remplie : SWOT/MoSCoW/RACI/risques non vides\\
dans le doc.\\
\vspace{0.5em}
\vspace{0.5em}
\vspace{0.5em}
\vspace{0.5em}
D\\
• Cohérence chiffrée : Page Factsheet unique (prix, dates, KPI).\\
\vspace{0.5em}
\vspace{0.5em}
\vspace{0.5em}
\vspace{0.5em}
FI\\
• Conformité/Éthique : Disclaimers, anti-trading, scénarios opéra-\\
toires.\\
• Employabilité : Portfolio, coaching, Demo Day. N\\
CO\\
\vspace{0.5em}
Priorisation MoSCoW — Partie 2 (Futur \textbackslash{}& Exclu)\\
\vspace{0.5em}
\vspace{0.5em}
Could have (Confort)                                                    Won’t have (Exclu)\\
\vspace{0.5em}
\vspace{0.5em}
\vspace{0.5em}
\vspace{0.5em}
\begin{confidential}CONFIDENTIEL — PROJET RBK 2.0                                   Page 205 sur 316\end{confidential}
RBK 2.0 — Whitepaper Architecte Web3                                                           Annexe A. Gabarits opérationnels \textbackslash{}& stratégiques\\
\vspace{0.5em}
\vspace{0.5em}
\vspace{0.5em}
• Simulation “Incident Drills” hebdo (War Room).                   • Contenus Trading / Spéculation.\\
• Module “Audit Readiness” avancé (Fuzzing, Invariants).           • Conseils d’évasion fiscale ou contournement légal.\\
• Offre “B2B Corporate” (Formations courtes).                      • Approche “Tout savoir” superficielle (Fluff).\\
• Incubation légère (Post-Demo Day).                               • Déploiement Mainnet sans audit préalable.\\
• Scénarios de scalabilité (Multi-promo).\\
\vspace{0.5em}
\vspace{0.5em}
\vspace{0.5em}
\vspace{0.5em}
EL\\
TI\\
D  EN\\
FI\\
N\\
CO\\
\vspace{0.5em}
\vspace{0.5em}
\vspace{0.5em}
\vspace{0.5em}
\begin{confidential}CONFIDENTIEL — PROJET RBK 2.0                                 Page 206 sur 316\end{confidential}
RBK 2.0 — Whitepaper Architecte Web3                                                    Annexe A. Gabarits opérationnels \textbackslash{}& stratégiques\\
\vspace{0.5em}
\vspace{0.5em}
\vspace{0.5em}
A.3     Matrice RACI (Responsable, Accountable,\\
Consulted, Informed)\\
Matrice RACI — Design \textbackslash{}& Build (1/2)\\
\vspace{0.5em}
\vspace{0.5em}
Activité                             R                A                C                           I\\
Définir la vision/positionnement     CEO              CEO              HoP, Mkt, Legal             Ops, Mentors, Students\\
(RBK Web3 Studio)\\
Figer la “Factsheet” (durée, co-     HoP              CEO              Ops, Career, Legal          Mkt, Students\\
\vspace{0.5em}
\vspace{0.5em}
\vspace{0.5em}
\vspace{0.5em}
EL\\
horte, prix, KPIs, tracks)\\
\vspace{0.5em}
\vspace{0.5em}
\vspace{0.5em}
\vspace{0.5em}
TI\\
Concevoir le tronc commun (objec- HoP                 HoP              TechLeads, SecLead          CEO, Ops\\
tifs, labs, DoD)\\
\vspace{0.5em}
\vspace{0.5em}
\vspace{0.5em}
\vspace{0.5em}
EN\\
Concevoir Track EVM (syllabus,       TechLead-EVM     HoP              SecLead, Mentors            CEO, Ops\\
\vspace{0.5em}
\vspace{0.5em}
\vspace{0.5em}
\vspace{0.5em}
D\\
labs, tooling)\\
\vspace{0.5em}
\vspace{0.5em}
\vspace{0.5em}
\vspace{0.5em}
FI\\
Concevoir Track Solana (syllabus,    TechLead-SOL     HoP              SecLead, Mentors            CEO, Ops\\
labs, tooling)\\
N\\
CO\\
Définir rubrics d’évaluation (PR,    SecLead          HoP              TechLeads, Mentors          CEO, Students\\
capstones, examens)\\
Définir capstones (specs, critères   HoP              HoP              TechLeads, SecLead          CEO, Mkt\\
d’acceptation, scoring)\\
Mettre en place l’infra (LMS, repos Ops               HoP              TechLeads                   CEO, Students\\
templates, CI, comptes outils)\\
\vspace{0.5em}
\vspace{0.5em}
\vspace{0.5em}
\vspace{0.5em}
\begin{confidential}CONFIDENTIEL — PROJET RBK 2.0                          Page 207 sur 316\end{confidential}
RBK 2.0 — Whitepaper Architecte Web3                                                        Annexe A. Gabarits opérationnels \textbackslash{}& stratégiques\\
\vspace{0.5em}
\vspace{0.5em}
Matrice RACI — Run \textbackslash{}& Opérations (2/2)\\
\vspace{0.5em}
\vspace{0.5em}
Activité                                                            R      A    C                        I\\
Admissions (test d’entrée, entretiens, sélection cohorte)           Mkt    CEO HoP, TechLeads, Career Ops, Students\\
Encadrement hebdo (sprints, reviews, incident drills)               HoP    HoP TechLeads, SecLead        CEO, Ops\\
Assurance qualité (cohérence doc, ToC, versioning, build LaTeX) Ops        HoP SecLead                   CEO, Mkt\\
Conformité \textbackslash{}& communication (disclaimers, scénario opératoire) Legal        CEO HoP, Ops                  Mkt, Students\\
Career services (portfolio, entretiens blancs, placements)          Career HoP Mentors, TechLeads        CEO, Mkt\\
\vspace{0.5em}
\vspace{0.5em}
\vspace{0.5em}
\vspace{0.5em}
EL\\
Organisation Demo Day (panel, format, scoring, invitations)         Career CEO HoP, Mkt, Mentors         Ops, Students\\
Suivi post-cohorte (KPIs, feedback, itérations curriculum)          HoP    CEO Ops, Career, TechLeads Mentors, Students\\
\vspace{0.5em}
\vspace{0.5em}
\vspace{0.5em}
\vspace{0.5em}
TI\\
D    EN\\
FI\\
N\\
CO\\
\vspace{0.5em}
\vspace{0.5em}
\vspace{0.5em}
\vspace{0.5em}
\begin{confidential}CONFIDENTIEL — PROJET RBK 2.0                                  Page 208 sur 316\end{confidential}
RBK 2.0 — Whitepaper Architecte Web3                                                            Annexe A. Gabarits opérationnels \textbackslash{}& stratégiques\\
\vspace{0.5em}
\vspace{0.5em}
\vspace{0.5em}
A.4     Gabarits Studio (Ops \textbackslash{}& Technique)\\
Ces modèles sont utilisés quotidiennement dans le cadre des rituels ”Studio”.\\
\vspace{0.5em}
\vspace{0.5em}
A.4.1 Template : Incident Drill Postmortem\\
\vspace{0.5em}
À remplir après chaque simulation d’attaque du Jeudi (War Room).\\
\vspace{0.5em}
Incident Postmortem \textbackslash{}#ID\\
Date : JJ/MM/AAAA                                                                                   Severity : Critical / High / Medium\\
Reporter : @StudentName                                                                                                Duration : 45min\\
\vspace{0.5em}
\vspace{0.5em}
\vspace{0.5em}
\vspace{0.5em}
EL\\
1. Le Scénario (What happened ?)\\
\vspace{0.5em}
\vspace{0.5em}
\vspace{0.5em}
\vspace{0.5em}
TI\\
Exemple : Le smart contract permettait de retirer plus que le solde disponible (underflow).\\
2. Impact (Si Mainnet)\\
\vspace{0.5em}
\vspace{0.5em}
\vspace{0.5em}
\vspace{0.5em}
EN\\
Exemple : Perte totale des fonds du pool de liquidité (TVL = \textbackslash{}$50k).\\
3. Root Cause Analysis (The ”Why”)\\
\vspace{0.5em}
\vspace{0.5em}
\vspace{0.5em}
\vspace{0.5em}
D\\
Exemple : Utilisation de unchecked block dans Solidity 0.8+ sans vérification préalable.\\
\vspace{0.5em}
\vspace{0.5em}
\vspace{0.5em}
\vspace{0.5em}
FI\\
4. Le Fix (Patch)\\
\vspace{0.5em}
N\\
Exemple : Retrait du bloc unchecked et ajout d’un revert CustomError().\\
CO\\
5. Lessons Learned (Prevention)\\
Exemple : Ajouter un test Fuzzing qui tente des retraits massifs aléatoires.\\
\vspace{0.5em}
\vspace{0.5em}
\vspace{0.5em}
A.5     Références Standards (Voir Annexe A)\\
Les gabarits obligatoires pour la structure de Repository, les ADRs et les journaux de preuves (PROOFS) sont définis exhaustivement\\
dans l’Annexe B (”Standards de Preuves \textbackslash{}& Évaluation”).\\
Tout projet Genesis, Explorer, Bâtisseur ou Architecte doit se conformer à ces standards pour être auditable.\\
\vspace{0.5em}
\vspace{0.5em}
\vspace{0.5em}
\vspace{0.5em}
\begin{confidential}CONFIDENTIEL — PROJET RBK 2.0                                   Page 209 sur 316\end{confidential}
B \textbar{} SYLLABUS TECHNIQUE\\
DÉTAILLÉ ET STANDARDS\\
\vspace{0.5em}
\vspace{0.5em}
\vspace{0.5em}
\vspace{0.5em}
EL\\
Cette annexe détaille les standards d’évaluation et le plan opérationnel semaine par semaine.\\
\vspace{0.5em}
\vspace{0.5em}
\vspace{0.5em}
\vspace{0.5em}
TI\\
EN\\
B.1 Gabarit de Dépôt Obligatoire\\
\vspace{0.5em}
\vspace{0.5em}
\vspace{0.5em}
\vspace{0.5em}
D\\
Tout projet doit respecter scrupuleusement cette arborescence standardisée. L’absence de l’un des fichiers obligatoires entraîne un rejet\\
\vspace{0.5em}
\vspace{0.5em}
\vspace{0.5em}
\vspace{0.5em}
FI\\
automatique (KO).\\
\vspace{0.5em}
Structure de Repository (Template)\\
N\\
CO\\
repo-root/\\
README.md                  \textbackslash{}# Présentation, Quickstart\\
RUNBOOK.md                 \textbackslash{}# How-to run/test/debug, Variables d'env.\\
PROOFS.md                  \textbackslash{}# Ledger hebdo des preuves (Liens PR, Démos)\\
THREAT\textbackslash{}_MODEL.md            \textbackslash{}# Analyse des risques \textbackslash{}& mitigations\\
CHANGELOG.md               \textbackslash{}# Historique des versions\\
LICENSE                    \textbackslash{}# MIT / Apache 2.0\\
OPENAPI.yaml               \textbackslash{}# Spec API (si pertinent, sinon placeholder)\\
.gitignore\\
\vspace{0.5em}
\vspace{0.5em}
\vspace{0.5em}
\vspace{0.5em}
210\\
Annexe B. SYLLABUS TECHNIQUE\\
RBK 2.0 — Whitepaper Architecte Web3                                                                       DÉTAILLÉ ET STANDARDS\\
\vspace{0.5em}
\vspace{0.5em}
\vspace{0.5em}
adr/                        \textbackslash{}# Architecture Decision Records\\
0001-template.md\\
index.md\\
\vspace{0.5em}
\vspace{0.5em}
docs/                       \textbackslash{}# Documentation additionnelle\\
architecture/             \textbackslash{}# Diagrammes (C4, Séquence)\\
security/                 \textbackslash{}# Audits, checklists\\
ops/                      \textbackslash{}# SLO, Playbooks\\
\vspace{0.5em}
\vspace{0.5em}
src/                        \textbackslash{}# Code Source\\
\vspace{0.5em}
\vspace{0.5em}
\vspace{0.5em}
\vspace{0.5em}
EL\\
tests/                      \textbackslash{}# Tests (Unit, Integ, E2E)\\
\vspace{0.5em}
\vspace{0.5em}
\vspace{0.5em}
\vspace{0.5em}
TI\\
scripts/                    \textbackslash{}# Scripts d'automatisation (OBLIGATOIRES)\\
\vspace{0.5em}
\vspace{0.5em}
\vspace{0.5em}
\vspace{0.5em}
EN\\
setup.sh                  \textbackslash{}# Installe deps, check versions\\
lint.sh                   \textbackslash{}# Format code \textbackslash{}& lint\\
\vspace{0.5em}
\vspace{0.5em}
\vspace{0.5em}
\vspace{0.5em}
D\\
test.sh                   \textbackslash{}# Run all tests\\
\vspace{0.5em}
\vspace{0.5em}
\vspace{0.5em}
\vspace{0.5em}
FI\\
build.sh                  \textbackslash{}# Compile artifacts\\
dev.sh                    \textbackslash{}# Start local dev env\\
smoke\textbackslash{}_test.sh\\
N\\
\section*{Basic healthcheck on running app}
CO\\
deploy\textbackslash{}_devnet.sh          \textbackslash{}# Deploy to devnet (si Web3)\\
\vspace{0.5em}
\vspace{0.5em}
infra/                      \textbackslash{}# Docker, CI/CD\\
\vspace{0.5em}
\vspace{0.5em}
KO-Structure (Éliminatoire)\\
L’absence de PROOFS.md, RUNBOOK.md, THREAT\textbackslash{}_MODEL.md ou du dossier adr/ est un échec automatique du livrable. Aucune évaluation\\
technique n’est réalisée si la structure n’est pas conforme.\\
\vspace{0.5em}
\vspace{0.5em}
\vspace{0.5em}
\vspace{0.5em}
\begin{confidential}CONFIDENTIEL — PROJET RBK 2.0                                 Page 211 sur 316\end{confidential}
Annexe B. SYLLABUS TECHNIQUE\\
RBK 2.0 — Whitepaper Architecte Web3                                                                      DÉTAILLÉ ET STANDARDS\\
\vspace{0.5em}
\vspace{0.5em}
B.1.1 Exigences de Reproductibilité (Scripts Obligatoires)\\
La ”Reproductibilité en une commande” est la signature de l’ingénieur RBK. Tous les scripts suivants sont obligatoires (même si\\
certains ne font qu’afficher ”N/A” avec explication) et doivent être exécutables depuis la racine du projet.\\
\vspace{0.5em}
Liste des Scripts Standardisés\\
\vspace{0.5em}
• scripts/setup.sh : Installe les dépendances système/langage.\\
\vspace{0.5em}
• scripts/lint.sh : Vérifie le style (clippy, eslint, prettier).\\
\vspace{0.5em}
• scripts/test.sh : Exécute tous les tests. Doit échouer (exit != 0) en cas d’erreur.\\
\vspace{0.5em}
\vspace{0.5em}
\vspace{0.5em}
\vspace{0.5em}
EL\\
• scripts/build.sh : Construit les binaires ou contrats.\\
\vspace{0.5em}
\vspace{0.5em}
\vspace{0.5em}
\vspace{0.5em}
TI\\
• scripts/dev.sh : Lance l’environnement de développement local.\\
\vspace{0.5em}
\vspace{0.5em}
\vspace{0.5em}
\vspace{0.5em}
EN\\
• scripts/smoke\textbackslash{}_test.sh : Vérifie que l’application démarre et répond (Healthcheck).\\
\vspace{0.5em}
\vspace{0.5em}
\vspace{0.5em}
\vspace{0.5em}
D\\
• scripts/deploy\textbackslash{}_devnet.sh : Déploie sur le réseau de test (si applicable).\\
\vspace{0.5em}
\vspace{0.5em}
\vspace{0.5em}
\vspace{0.5em}
FI\\
Règles d’Interface\\
N\\
CO\\
• Support de l’argument --help.\\
\vspace{0.5em}
• Non-interactif : Aucun prompt bloquant (sauf login explicite).\\
\vspace{0.5em}
• Codes de sortie standards (0 = Succès, !=0 = Échec).\\
\vspace{0.5em}
KO-Repro (Éliminatoire)\\
\vspace{0.5em}
Si la séquence ./scripts/setup.sh \textbackslash{}&\textbackslash{}& ./scripts/test.sh échoue sur la machine de l’évaluateur, le projet est rejeté (NO-GO).\\
\vspace{0.5em}
\vspace{0.5em}
\vspace{0.5em}
\vspace{0.5em}
\begin{confidential}CONFIDENTIEL — PROJET RBK 2.0                                 Page 212 sur 316\end{confidential}
Annexe B. SYLLABUS TECHNIQUE\\
RBK 2.0 — Whitepaper Architecte Web3                                                                               DÉTAILLÉ ET STANDARDS\\
\vspace{0.5em}
\vspace{0.5em}
\vspace{0.5em}
B.2 Rubrique d’Évaluation \textbackslash{}& Processus\\
L’évaluation chez RBK est binaire sur la forme (KO si non-conforme) et granulaire sur le fond (Grille /100). Elle s’applique à deux\\
fréquences : Hebdomadaire (Sprints) et Mensuelle/Trimestrielle (Gates).\\
\vspace{0.5em}
\vspace{0.5em}
B.2.1 Revues Hebdomadaires (Sprints)\\
Chaque vendredi, le livrable est évalué selon le Gabarit et la Grille Standardisée (voir Section B.4).\\
\vspace{0.5em}
• Validation Automatique : Exécution des scripts ‘scripts/test.sh‘ et ‘scripts/lint.sh‘.\\
\vspace{0.5em}
• Review Technique : Qualité du code, couverture de tests, et conformité de l’architecture.\\
\vspace{0.5em}
\vspace{0.5em}
\vspace{0.5em}
\vspace{0.5em}
EL\\
• Review Produit : Démo fonctionnelle et UI/UX (si applicable).\\
\vspace{0.5em}
\vspace{0.5em}
\vspace{0.5em}
\vspace{0.5em}
TI\\
EN\\
B.2.2 Format de Jury (Gate)\\
\vspace{0.5em}
\vspace{0.5em}
\vspace{0.5em}
\vspace{0.5em}
D\\
Le passage de niveau se fait devant un jury technique. Le format détaillé, les questions fixes et le modèle de verdict sont définis en\\
\vspace{0.5em}
\vspace{0.5em}
\vspace{0.5em}
\vspace{0.5em}
FI\\
Annexe B.5.\\
\vspace{0.5em}
\vspace{0.5em}
N\\
CO\\
B.3 Standards de Documentation et Templates\\
Cette section définit le contenu attendu des fichiers critiques. Le non-respect de ces formats entraîne des pénalités ou un KO.\\
\vspace{0.5em}
\vspace{0.5em}
B.3.1 Règles ADR (Architecture Decision Records)\\
Dossier : adr/\\
Convention de nommage : NNNN-slug-kebab-case.md (ex : 0002-choix-db-sql.md).\\
Fréquence : Au moins 1 ADR par module majeur ou décision structurante. Une ADR doit être immuable une fois ”Accepted” (sauf\\
changement de statut).\\
\vspace{0.5em}
\vspace{0.5em}
\vspace{0.5em}
\begin{confidential}CONFIDENTIEL — PROJET RBK 2.0                                 Page 213 sur 316\end{confidential}
Annexe B. SYLLABUS TECHNIQUE\\
RBK 2.0 — Whitepaper Architecte Web3                                                                              DÉTAILLÉ ET STANDARDS\\
\vspace{0.5em}
\vspace{0.5em}
\vspace{0.5em}
KO-ADR\\
Aucune ADR ”Accepted” alors que des choix architecturaux majeurs ont été faits = NO-GO.\\
\vspace{0.5em}
\vspace{0.5em}
B.3.2 Structure PROOFS.md (Journal des Preuves)\\
Ce fichier est le ”Ledger” de votre travail. Il doit contenir une entrée par semaine (SXX) avec les champs suivants dans l’ordre :\\
\vspace{0.5em}
1. PR(s) : Liens vers les Pull Requests mergées.\\
\vspace{0.5em}
2. Tag/Release : Version publiée (ex : v0.3.0).\\
\vspace{0.5em}
\vspace{0.5em}
\vspace{0.5em}
\vspace{0.5em}
EL\\
3. CI run(s) : Lien vers le job CI vert correspondant.\\
\vspace{0.5em}
\vspace{0.5em}
\vspace{0.5em}
\vspace{0.5em}
TI\\
4. Demo link(s) : Vidéo (Loom/Youtube) ou lien déploiement live.\\
\vspace{0.5em}
5. Reproduce commands : Commandes exactes pour vérifier le travail.\\
\vspace{0.5em}
\vspace{0.5em}
\vspace{0.5em}
\vspace{0.5em}
EN\\
6. Docs updated : Liste des fichiers de documentation mis à jour.\\
\vspace{0.5em}
\vspace{0.5em}
\vspace{0.5em}
\vspace{0.5em}
ID\\
7. ADR added/updated : Référence aux ADRs impactées.\\
\vspace{0.5em}
NF\\
8. Metrics : Benchmarks ou métriques clés (si pertinent).\\
CO\\
9. Known issues : Bugs connus et contournements.\\
\vspace{0.5em}
KO-PROOFS\\
Champs ”Reproduce commands” ou ”Demo link” manquants sans justification explicite = rejet immédiat.\\
\vspace{0.5em}
\vspace{0.5em}
B.3.3 Template : PROOFS.md\\
\section*{PROOFS --- Ledger Hebdomadaire}
\vspace{0.5em}
\vspace{0.5em}
\vspace{0.5em}
\vspace{0.5em}
\begin{confidential}CONFIDENTIEL — PROJET RBK 2.0                               Page 214 sur 316\end{confidential}
Annexe B. SYLLABUS TECHNIQUE\\
RBK 2.0 — Whitepaper Architecte Web3                                          DÉTAILLÉ ET STANDARDS\\
\vspace{0.5em}
\vspace{0.5em}
\vspace{0.5em}
\subsection*{Sxx (Semaine courante)}
- PR(s): \textless{}Lien vers GitHub\textgreater{}\\
- Tag/Release: v0.x.x\\
- CI Status: \textless{}Lien Run\textgreater{} (Green)\\
- Demo Link: \textless{}Lien Loom/YouTube\textgreater{} (30-90s)\\
- Reproduce commands: `./scripts/test\textbackslash{}_scenario\textbackslash{}_A.sh`\\
- Docs Updated: README, ADR-002\\
- ADR added/updated: ADR-005\\
- Metrics: "TPS=5000 sur machine locale"\\
- Known Issues: ...\\
\vspace{0.5em}
\vspace{0.5em}
\vspace{0.5em}
\vspace{0.5em}
EL\\
B.3.4 Template : THREAT\textbackslash{}_MODEL.md\\
\vspace{0.5em}
\vspace{0.5em}
\vspace{0.5em}
\vspace{0.5em}
TI\\
\section*{Threat Model}
\vspace{0.5em}
\vspace{0.5em}
\vspace{0.5em}
\vspace{0.5em}
EN\\
\subsection*{1. Assets to Protect}
\vspace{0.5em}
\vspace{0.5em}
\vspace{0.5em}
\vspace{0.5em}
ID\\
- User Funds (Wallet keys)\\
\vspace{0.5em}
\vspace{0.5em}
\vspace{0.5em}
NF\\
- Data Integrity (Ledger)                 CO\\
\subsection*{2. Attack Surface}
- Public RPC Endpoints\\
- Input fields (DApp)\\
\vspace{0.5em}
\vspace{0.5em}
\subsection*{3. Threats \textbackslash{}& Mitigations}
\begin{center}\begin{tabular}\textbar{} Threat \textbar{} Mitigation \textbar{} Status \textbar{}\end{tabular}\end{center}
\begin{center}\begin{tabular}\textbar{}--------\textbar{}------------\textbar{}--------\textbar{}\end{tabular}\end{center}
\begin{center}\begin{tabular}\textbar{} Replay Attack \textbar{} Nonce / Idempotency Key \textbar{} Implemented \textbar{}\end{tabular}\end{center}
\begin{center}\begin{tabular}\textbar{} XSS Injection \textbar{} Sanitize Inputs + CSP     \textbar{} Tested \textbar{}\end{tabular}\end{center}
\vspace{0.5em}
\vspace{0.5em}
\vspace{0.5em}
\vspace{0.5em}
\begin{confidential}CONFIDENTIEL — PROJET RBK 2.0                      Page 215 sur 316\end{confidential}
Annexe B. SYLLABUS TECHNIQUE\\
RBK 2.0 — Whitepaper Architecte Web3                                                  DÉTAILLÉ ET STANDARDS\\
\vspace{0.5em}
\vspace{0.5em}
\subsection*{4. Residual Risks}
- Dependency supply chain attack (audit pending).\\
\vspace{0.5em}
\vspace{0.5em}
B.3.5 Template : adr/NNNN-template.md\\
\section*{ADR NNNN --- \textless{}Titre de la Décision\textgreater{}}
\vspace{0.5em}
\vspace{0.5em}
\subsection*{Status}
Proposed \textbar{} Accepted \textbar{} Deprecated \textbar{} Superseded\\
\vspace{0.5em}
\vspace{0.5em}
\vspace{0.5em}
\vspace{0.5em}
EL\\
\subsection*{Context}
Quel est le problème ? Quelles sont les contraintes techniques ou business ?\\
\vspace{0.5em}
\vspace{0.5em}
\vspace{0.5em}
\vspace{0.5em}
TI\\
EN\\
\subsection*{Decision}
Quelle solution avons-nous choisie ? (Technologie, Pattern, Librairie)\\
\vspace{0.5em}
\vspace{0.5em}
\vspace{0.5em}
\vspace{0.5em}
D\\
FI\\
\subsection*{Consequences}
- Positive : ...\\
- Negative : ...\\
N\\
CO\\
\subsection*{Alternatives Considered}
- Option A : rejetée car...\\
- Option B : rejetée car...\\
\vspace{0.5em}
\vspace{0.5em}
\subsection*{How to Validate}
Quel test ou métrique prouve que c'est le bon choix ?\\
\vspace{0.5em}
\vspace{0.5em}
\vspace{0.5em}
B.4 Grille de Notation Standardisée (/100) et Seuils\\
\vspace{0.5em}
\vspace{0.5em}
\begin{confidential}CONFIDENTIEL — PROJET RBK 2.0                       Page 216 sur 316\end{confidential}
Annexe B. SYLLABUS TECHNIQUE\\
RBK 2.0 — Whitepaper Architecte Web3                                                                         DÉTAILLÉ ET STANDARDS\\
\vspace{0.5em}
\vspace{0.5em}
\vspace{0.5em}
Ce barème est appliqué strictement lors des revues hebdomadaires et des Gates.\\
\vspace{0.5em}
\vspace{0.5em}
B.4.1 Détail de la Notation (/100)\\
\vspace{0.5em}
\vspace{0.5em}
Critère                           Score           Niveaux de performance\\
1. Fonctionnel (Produit)          25 pts\\
• 0–10 : Feature incomplète ou cassée.\\
• 11–20 : Feature conforme mais fragile (pas de cas limites).\\
• 21–25 : Conforme + Cas limites gérés + Démo fluide.\\
\vspace{0.5em}
\vspace{0.5em}
\vspace{0.5em}
\vspace{0.5em}
EL\\
TI\\
2. Qualité Code \textbackslash{}& Tests           25 pts\\
• 0–10 : Pas de tests ou code illisible. Pas de CI.\\
\vspace{0.5em}
\vspace{0.5em}
\vspace{0.5em}
\vspace{0.5em}
EN\\
• 11–20 : Tests présents mais couverture faible. CI verte.\\
• 21–25 : Tests Unitaires/Intégration solides. Refactoring propre.\\
\vspace{0.5em}
\vspace{0.5em}
\vspace{0.5em}
\vspace{0.5em}
D\\
FI\\
3. Sécurité (Threat Model)        20 pts\\
\vspace{0.5em}
N      • 0–8 : Aucun Threat Model ou mitigations absentes.\\
CO\\
• 9–15 : TM présent mais incomplet. Risques résiduels flous.\\
• 16–20 : TM complet + Checklist appliquée + Preuves.\\
\vspace{0.5em}
4. Ops \textbackslash{}& Observabilité            15 pts\\
• 0–5 : Logs insuffisants, pas de Runbook.\\
• 6–11 : Logs présents. Runbook incomplet.\\
• 12–15 : Logs structurés, Métriques clés, Incident Playbook.\\
\vspace{0.5em}
\vspace{0.5em}
\vspace{0.5em}
\vspace{0.5em}
\begin{confidential}CONFIDENTIEL — PROJET RBK 2.0                             Page 217 sur 316\end{confidential}
Annexe B. SYLLABUS TECHNIQUE\\
RBK 2.0 — Whitepaper Architecte Web3                                                                          DÉTAILLÉ ET STANDARDS\\
\vspace{0.5em}
\vspace{0.5em}
Critère                   Score                      Niveaux de performance\\
5. Documentation \textbackslash{}& Repro. 10 pts\\
• 0–3 : Instructions floues ou absentes.\\
• 4–7 : Quickstart fonctionnel mais incomplet.\\
• 8–10 : Repro en 1 commande. Variables d’env claires.\\
\vspace{0.5em}
6. Discipline (Preuves)           5 pts\\
• 0–2 : PROOFS absent ou PR non traçables.\\
• 3–4 : PROOFS partiel.\\
• 5 : PROOFS complet + ADR + Changelog/Tags.\\
\vspace{0.5em}
\vspace{0.5em}
\vspace{0.5em}
\vspace{0.5em}
EL\\
TI\\
EN\\
B.4.2 Seuils de Décision (GO / WARN / NO-GO)\\
Le verdict dépend du score total ET des critères critiques.\\
\vspace{0.5em}
\vspace{0.5em}
\vspace{0.5em}
\vspace{0.5em}
D\\
FI\\
1. Hebdomadaire\\
\vspace{0.5em}
• ≥ 70/100 : VALIDÉ.                           N\\
CO\\
• 60–69 : VALIDÉ avec WARN (Remédiation obligatoire sous 72h).\\
\vspace{0.5em}
• \textless{} 60 : NON-VALIDÉ (NO-GO). Progression bloquée.\\
\vspace{0.5em}
2. Gates (Genesis, Explorer, Bâtisseur/Architecte)\\
\vspace{0.5em}
• ≥ 80/100 + Aucun KO : GO (Passage au niveau suivant).\\
\vspace{0.5em}
• 70–79 ou 1 Faiblesse critique : WARN (Passage conditionnel + Plan d’action).\\
\vspace{0.5em}
• \textless{} 70 ou Critère KO : NO-GO (Redoublement de phase ou sortie).\\
\vspace{0.5em}
\vspace{0.5em}
\begin{confidential}CONFIDENTIEL — PROJET RBK 2.0                                Page 218 sur 316\end{confidential}
Annexe B. SYLLABUS TECHNIQUE\\
RBK 2.0 — Whitepaper Architecte Web3                                                                         DÉTAILLÉ ET STANDARDS\\
\vspace{0.5em}
\vspace{0.5em}
\vspace{0.5em}
Critères Critiques (Échec Automatique)\\
\vspace{0.5em}
Un score insuffisant dans ces catégories entraîne un NO-GO immédiat, quel que soit le total :\\
\vspace{0.5em}
• Sécurité : Score \textless{} 8/20.\\
\vspace{0.5em}
• Reproductibilité : Score \textless{} 4/10.\\
\vspace{0.5em}
• Tests/CI : Score \textless{} 10/25.\\
\vspace{0.5em}
Exemple : Un projet avec 85/100 mais 0 tests est NO-GO.\\
\vspace{0.5em}
\vspace{0.5em}
\vspace{0.5em}
\vspace{0.5em}
EL\\
B.5 Format de Jury Gate et Procès-Verbal\\
\vspace{0.5em}
\vspace{0.5em}
\vspace{0.5em}
\vspace{0.5em}
TI\\
Cette section détaille le déroulé standard (45-60 min), les questions imposées et le modèle de verdict pour les passages de niveau\\
\vspace{0.5em}
\vspace{0.5em}
\vspace{0.5em}
\vspace{0.5em}
EN\\
(Genesis, Explorer, Bâtisseur/Architecte).\\
\vspace{0.5em}
\vspace{0.5em}
\vspace{0.5em}
\vspace{0.5em}
D\\
FI\\
B.5.1 Déroulé Standard (45–60 min)\\
\vspace{0.5em}
\vspace{0.5em}
N\\
1. Pitch (2 min) : Quoi, pour qui, pourquoi maintenant, scope exact.\\
CO\\
2. Demo Live (8–10 min) : Cas nominal (User Story principale) + 1 Cas d’erreur géré (RPC down, input invalid...).\\
\vspace{0.5em}
3. Reproduction (5 min) : Exécution en direct des scripts setup / test / run sur la machine d’un juré ou environnement vierge.\\
\vspace{0.5em}
4. Architecture (10 min) : Diagramme, composants, justification des choix (ADR).\\
\vspace{0.5em}
5. Sécurité (10 min) : Threat Model, 3 menaces majeures, mitigations prouvées, risques résiduels.\\
\vspace{0.5em}
6. Qualité / Ops (5–10 min) : Stratégie de tests, CI, Logs, Metrics, Incident Playbook.\\
\vspace{0.5em}
7. Q\textbackslash{}&A et Verdict (5 min) : Questions de creusement et décision.\\
\vspace{0.5em}
\vspace{0.5em}
\vspace{0.5em}
\begin{confidential}CONFIDENTIEL — PROJET RBK 2.0                                Page 219 sur 316\end{confidential}
Annexe B. SYLLABUS TECHNIQUE\\
RBK 2.0 — Whitepaper Architecte Web3                                                                              DÉTAILLÉ ET STANDARDS\\
\vspace{0.5em}
\vspace{0.5em}
B.5.2 Checklist de Questions Fixes\\
Le jury doit obligatoirement poser les questions suivantes pour valider la profondeur de la maîtrise.\\
\vspace{0.5em}
1. Produit / Fonctionnel\\
\vspace{0.5em}
• Quelle est la user story principale et comment la validez-vous ?\\
\vspace{0.5em}
• Montrez un cas d’erreur réel (ex : RPC down, Transaction fail) et votre gestion d’erreur.\\
\vspace{0.5em}
2. Reproductibilité\\
\vspace{0.5em}
\vspace{0.5em}
\vspace{0.5em}
\vspace{0.5em}
EL\\
• À partir d’un laptop vierge : quelles commandes exactes pour exécuter le projet ?\\
\vspace{0.5em}
\vspace{0.5em}
\vspace{0.5em}
\vspace{0.5em}
TI\\
• Quelles variables d’environnement sont indispensables et pourquoi ?\\
\vspace{0.5em}
\vspace{0.5em}
\vspace{0.5em}
\vspace{0.5em}
EN\\
3. Architecture\\
\vspace{0.5em}
\vspace{0.5em}
\vspace{0.5em}
\vspace{0.5em}
D\\
• Quelles décisions d’architecture sont irréversibles ? (Citer les ADRs)\\
\vspace{0.5em}
\vspace{0.5em}
\vspace{0.5em}
\vspace{0.5em}
FI\\
N\\
• Quelles alternatives avez-vous rejetées et pourquoi ?\\
CO\\
4. Sécurité\\
\vspace{0.5em}
• Quelles sont vos 3 menaces les plus graves (impact business/tech) ?\\
\vspace{0.5em}
• Où sont les preuves que vos mitigations fonctionnent ? (Tests, garde-fous)\\
\vspace{0.5em}
• Quels risques résiduels acceptez-vous ?\\
\vspace{0.5em}
\vspace{0.5em}
\vspace{0.5em}
\vspace{0.5em}
\begin{confidential}CONFIDENTIEL — PROJET RBK 2.0                                Page 220 sur 316\end{confidential}
Annexe B. SYLLABUS TECHNIQUE\\
RBK 2.0 — Whitepaper Architecte Web3                                                                             DÉTAILLÉ ET STANDARDS\\
\vspace{0.5em}
\vspace{0.5em}
5. Qualité \textbackslash{}& Ops\\
\vspace{0.5em}
• Montrez la stratégie de tests (Unit/Integ/E2E) et la couverture de chaque couche.\\
\vspace{0.5em}
• Si le système tombe en panne demain, que fait l’opérateur ? (Runbook, Logs).\\
\vspace{0.5em}
• Quel est votre plus gros compromis technique (Dette) ?\\
\vspace{0.5em}
6. Traçabilité\\
\vspace{0.5em}
• Où est votre ledger de preuves (PROOFS.md) ?\\
\vspace{0.5em}
\vspace{0.5em}
\vspace{0.5em}
\vspace{0.5em}
EL\\
• Quels tags/releases et changelogs permettent de suivre la progression ?\\
\vspace{0.5em}
\vspace{0.5em}
\vspace{0.5em}
\vspace{0.5em}
TI\\
B.5.3 Barème de Décision (GO / WARN / NO-GO)\\
\vspace{0.5em}
\vspace{0.5em}
\vspace{0.5em}
\vspace{0.5em}
EN\\
Le score est calculé sur la Grille /100 (Annexe B.4). Au-delà du score, des conditions éliminatoires s’appliquent.\\
\vspace{0.5em}
\vspace{0.5em}
\vspace{0.5em}
\vspace{0.5em}
D\\
Conditions KO (Éliminatoires)\\
\vspace{0.5em}
\vspace{0.5em}
\vspace{0.5em}
\vspace{0.5em}
FI\\
• KO-Repro : Impossible de lancer (setup/test/run) en conditions normales.\\
\vspace{0.5em}
N\\
CO\\
• KO-Sécurité : Absence de Threat Model ou mitigations non démontrées.\\
\vspace{0.5em}
• KO-Qualité : Pas de CI, pas de tests, aucune stratégie crédible.\\
\vspace{0.5em}
\vspace{0.5em}
Verdict\\
\vspace{0.5em}
• GO : Score ≥ 80/100, Aucun KO, Max 2 faiblesses mineures.\\
\vspace{0.5em}
• WARN : Score 70–79 OU 1 faiblesse critique (non KO). Plan de remédiation obligatoire.\\
\vspace{0.5em}
• NO-GO : Score \textless{} 70 OU Condition KO OU Incohérence majeure.\\
\vspace{0.5em}
\vspace{0.5em}
\vspace{0.5em}
\begin{confidential}CONFIDENTIEL — PROJET RBK 2.0                                 Page 221 sur 316\end{confidential}
Annexe B. SYLLABUS TECHNIQUE\\
RBK 2.0 — Whitepaper Architecte Web3                                                  DÉTAILLÉ ET STANDARDS\\
\vspace{0.5em}
\vspace{0.5em}
B.5.4 Modèle de Procès-Verbal de Jury\\
Ce template doit être rempli à l’issue de chaque Gate.\\
\vspace{0.5em}
PV de Jury (Template)\\
\vspace{0.5em}
\section*{Gate Verdict --- \textless{}Genesis/Explorer/Bâtisseur/Architecte\textgreater{} --- \textless{}Date\textgreater{}}
\vspace{0.5em}
\vspace{0.5em}
\subsection*{Decision}
GO \textbar{} WARN \textbar{} NO-GO\\
\vspace{0.5em}
\vspace{0.5em}
\subsection*{Score}
\vspace{0.5em}
\vspace{0.5em}
\vspace{0.5em}
\vspace{0.5em}
EL\\
Total: \textbackslash{}_\textbackslash{}_/100 (Voir Grille Annexe B.4)\\
\vspace{0.5em}
\vspace{0.5em}
\vspace{0.5em}
\vspace{0.5em}
TI\\
Breakdown:\\
\vspace{0.5em}
\vspace{0.5em}
\vspace{0.5em}
\vspace{0.5em}
EN\\
- Functional:        \textbackslash{}_\textbackslash{}_/25\\
- Quality/Tests:     \textbackslash{}_\textbackslash{}_/25\\
\vspace{0.5em}
\vspace{0.5em}
\vspace{0.5em}
\vspace{0.5em}
D\\
- Security:          \textbackslash{}_\textbackslash{}_/20\\
- Ops/Observability: \textbackslash{}_\textbackslash{}_/15\\
\vspace{0.5em}
\vspace{0.5em}
\vspace{0.5em}
\vspace{0.5em}
FI\\
- Docs/Repro:        \textbackslash{}_\textbackslash{}_/10\\
- Discipline/Proofs: \textbackslash{}_\textbackslash{}_/5\\
N\\
CO\\
\subsection*{Strengths}
- ...\\
- ...\\
\vspace{0.5em}
\vspace{0.5em}
\subsection*{Weaknesses / Risques}
- ...\\
- ...\\
\vspace{0.5em}
\vspace{0.5em}
\subsection*{Required Actions (if WARN/NO-GO)}
\vspace{0.5em}
\vspace{0.5em}
\vspace{0.5em}
\vspace{0.5em}
\begin{confidential}CONFIDENTIEL — PROJET RBK 2.0                               Page 222 sur 316\end{confidential}
Annexe B. SYLLABUS TECHNIQUE\\
RBK 2.0 — Whitepaper Architecte Web3                                         DÉTAILLÉ ET STANDARDS\\
\vspace{0.5em}
\vspace{0.5em}
\vspace{0.5em}
1) ...\\
2) ...\\
\vspace{0.5em}
\vspace{0.5em}
Deadline: \textless{}Date\textgreater{}\\
Proof expected: \textless{}PR link / Tag / Demo video\textgreater{}\\
\vspace{0.5em}
\vspace{0.5em}
\vspace{0.5em}
\vspace{0.5em}
B.6 Contrat de Livrable (1 page)\\
\vspace{0.5em}
\vspace{0.5em}
\vspace{0.5em}
\vspace{0.5em}
EL\\
TI\\
END\\
FI\\
N\\
CO\\
\vspace{0.5em}
\vspace{0.5em}
\vspace{0.5em}
\vspace{0.5em}
\begin{confidential}CONFIDENTIEL — PROJET RBK 2.0                      Page 223 sur 316\end{confidential}
Annexe B. SYLLABUS TECHNIQUE\\
RBK 2.0 — Whitepaper Architecte Web3                                                                                    DÉTAILLÉ ET STANDARDS\\
\vspace{0.5em}
\vspace{0.5em}
\vspace{0.5em}
\vspace{0.5em}
CONTRAT DE LIVRABLE UNIQUE (Genesis/Explorer/Bâtisseur-Architecte)\\
Ce document fixe les exigences minimales non négociables. Tout manquement à un point KO entraîne un NO-GO immédiat.\\
• 1. Structure \textbackslash{}& Discipline (KO-Structure)\\
– Fichiers obligatoires : README.md, RUNBOOK.md, PROOFS.md, THREAT\textbackslash{}_MODEL.md, CHANGELOG.md, dossier adr/, dossier scripts/, dossier\\
tests/.\\
– ADR : Naming adr/NNNN-slug.md. Min. 1/sprint ou décis. structurante. Status/Context/Decision/Validation requis. KO si absentes.\\
– PROOFS.md : Ledger hebdo. Ordre imposé : 1.PRs, 2.Tag, 3.CI, 4.Démo, 5.Commandes Repro, 6.Docs, 7.ADR, 8.Metrics, 9.Issues. KO si\\
\vspace{0.5em}
\vspace{0.5em}
\vspace{0.5em}
\vspace{0.5em}
EL\\
incomplet.\\
• 2. Reproductibilité (KO-Repro)\\
– Scripts obligatoires : setup.sh, lint.sh, test.sh, build.sh, dev.sh, smoke\textbackslash{}_test.sh.\\
\vspace{0.5em}
\vspace{0.5em}
\vspace{0.5em}
\vspace{0.5em}
TI\\
– Interface : Support --help, non-interactif, exit code standard.\\
– Règle d’Or : La séquence ./setup.sh \textbackslash{}&\textbackslash{}& ./test.sh doit réussir sur une machine vierge.\\
\vspace{0.5em}
\vspace{0.5em}
\vspace{0.5em}
\vspace{0.5em}
EN\\
• 3. Tests \textbackslash{}& Qualité (KO-Tests)\\
– Commande : ./scripts/test.sh lance TOUS les tests. CI verte obligatoire.\\
– Seuils Genesis : ≥ 10 Tests U. + 1 Intégration ou Smoke Test.\\
\vspace{0.5em}
\vspace{0.5em}
\vspace{0.5em}
\vspace{0.5em}
D\\
– Seuils Explorer : Unit + Integ + Min. 1 E2E (si UI/Flow).\\
– Seuils Bâtisseur/Architecte : E2E complets (Capstones) + Tests d’échec (Chaos/Error handling).\\
\vspace{0.5em}
\vspace{0.5em}
FI\\
– Couverture : 1 cas nominal + 1 cas d’erreur par feature majeure.\\
• 4. Démo \textbackslash{}& Vérité (KO-Démo)\\
N\\
– Hebdo : 30-90s (Vidéo/Gif) montrant 1 User Story + 1 Erreur gérée.\\
O\\
– Gate : Démo live (10 min) reproductible par le jury via le Runbook.\\
– Anti-Storytelling : Toute affirmation (”ça marche”, ”sécurisé”) doit avoir une preuve (Lien, Log, Test, Code) dans PROOFS.md.\\
C\\
\vspace{0.5em}
• 5. Sécurité \textbackslash{}& Ops (KO-Sécurité)\\
– THREAT\textbackslash{}_MODEL.md à jour. Checklist sécurité (docs/security/) cochée.\\
– Menaces majeures identifiées et mitigées (Preuves dans tests/guards).\\
Engagement : L’équipe s’engage à respecter ce contrat dès S01. Toute exception doit être justifiée par une ADR validée.\\
\vspace{0.5em}
\vspace{0.5em}
\vspace{0.5em}
\vspace{0.5em}
\begin{confidential}CONFIDENTIEL — PROJET RBK 2.0                                    Page 224 sur 316\end{confidential}
Annexe B. SYLLABUS TECHNIQUE\\
RBK 2.0 — Whitepaper Architecte Web3                                                 DÉTAILLÉ ET STANDARDS\\
\vspace{0.5em}
\vspace{0.5em}
\vspace{0.5em}
B.7 Syllabus Opérationnel (48 lignes / 48 semaines)\\
Ce tableau constitue la source de vérité pour l’exécution hebdomadaire.\\
\vspace{0.5em}
\vspace{0.5em}
\vspace{0.5em}
\vspace{0.5em}
EL\\
TI\\
D    EN\\
FI\\
N\\
CO\\
\vspace{0.5em}
\vspace{0.5em}
\vspace{0.5em}
\vspace{0.5em}
\begin{confidential}CONFIDENTIEL — PROJET RBK 2.0                              Page 225 sur 316\end{confidential}
Annexe B. SYLLABUS TECHNIQUE\\
RBK 2.0 — Whitepaper Architecte Web3                                                                                         DÉTAILLÉ ET STANDARDS\\
\vspace{0.5em}
\vspace{0.5em}
Niveau 1 : Piscine \textbackslash{}& Fondations (S01-S12)\\
\vspace{0.5em}
Sem    Cours (Objectifs)                     Labs (TP Guidés)                Livrable (Preuve)         Critères (DoD)\\
S01    Linux/CLI, Git workflow, structure Setup env, Makefile, PR template   Repo ”foundation”         CI verte, README exec\\
repo\\
S02    Rust bases : ownership, types, er- Parsing, validation, erreurs       CLI Rust v1               Zéro panic, tests ≥ 10\\
reurs\\
S03    Rust avancé : traits, generics, mo- Refactor API, crate, rustdoc      Crate réutilisable        Doc rustdoc, API clean\\
dules\\
S04    Crypto 1 : hash, encoding, serializa- SHA/Merkle, vectors tests       ”crypto-toolkit”          Vectors tests, determ.\\
tion\\
S05    Crypto 2 : signatures, keys, threats Sign/verify, key mgmt            Module signature          Keys safe, threat note\\
\vspace{0.5em}
\vspace{0.5em}
\vspace{0.5em}
\vspace{0.5em}
EL\\
S06    Blockchain 1 : tx, RPC, finality      Client RPC, parsing tx          ”explorer CLI”            Robustesse réseau\\
S07    API Design : REST, cache, limits      API endpoints, OpenAPI          Micro-API + Spec          OpenApi valid, err codes\\
\vspace{0.5em}
\vspace{0.5em}
\vspace{0.5em}
\vspace{0.5em}
TI\\
S08    Indexation : events, storage local    Indexer minimal, backfill       Indexer + Wallet API      Data cohérente\\
\vspace{0.5em}
\vspace{0.5em}
\vspace{0.5em}
\vspace{0.5em}
N\\
S09    UX Wallet : signature, pending, er- Flow connect/sign/read            Mini front/CLI UX         UX claire, erreurs gérées\\
\vspace{0.5em}
\vspace{0.5em}
\vspace{0.5em}
\vspace{0.5em}
DE\\
rors\\
S10    Tokenomics 1 : supply, incentives Simu simple, stress tests           ”tokenomics memo”         Hypothèses explicites\\
\vspace{0.5em}
\vspace{0.5em}
\vspace{0.5em}
\vspace{0.5em}
FI\\
S11    Intégration E2E : vertical slice      1 user story complète           Mini-dApp E2E             Démo fonct., tests E2E\\
S12    Block Check “Genesis”                 Nettoyage, docs, release        PACK Genesis              GO/NO-GO Jury\\
\vspace{0.5em}
\vspace{0.5em}
N\\
CO\\
\vspace{0.5em}
\vspace{0.5em}
\vspace{0.5em}
\vspace{0.5em}
\begin{confidential}CONFIDENTIEL — PROJET RBK 2.0                                     Page 226 sur 316\end{confidential}
Annexe B. SYLLABUS TECHNIQUE\\
RBK 2.0 — Whitepaper Architecte Web3                                                                                                             DÉTAILLÉ ET STANDARDS\\
\vspace{0.5em}
\vspace{0.5em}
Bâtisseur : Spécialisation (S13-S28)\\
\vspace{0.5em}
Sem    Cours (A/B/C)                                   Labs (A/B/C)                                Livrable             DoD\\
S13    A : Modèle Solana / B : Solidity bases / C : A : Prog. natif / B : ERC min / C : User Repo S13 (Tracké)          Tests + README\\
Discovery                                       Stories\\
S14    A : Sérialisation, Budget / B : Gas, Erreurs A : Constraints / B : Events / C : Proto.      PR + ADR \textbackslash{}#1          ADR clair, reviewable\\
/ C : Wireframes\\
S15    A : Sécurité Auth. / B : Access Ctrl / C : KPIs A : Auth. / B : Roles / C : Metrics Plan    ”Module pack” v1     Checklist sécu\\
S16    A : Tests (Locaux) / B : Foundry Intro / C : A : Tests état / B : Fuzz intro / C : Roadmap PRD/Tool Gate         Pipeline test OK\\
PRD Final\\
S17    A : Anchor IDL / B : Foundry Mastery / C : A : Refactor Anchor / B : Fuzz / C : Incen- Repo S17                  Tests exec, doc\\
Tokenomics                                      tives\\
S18    A : Anchor Constraints / B : ERC Stds / C : A : Seeds/PDA / B : ERCs / C : Scénarios        ”Simu pack”          Résultats repro.\\
Simu.\\
S19    A : CPI / B : Architecture / C : Risques        A : CPI / B : Timelocks / C : Mitigations   ADR \textbackslash{}#2 + PR          Menaces listées\\
\vspace{0.5em}
\vspace{0.5em}
\vspace{0.5em}
\vspace{0.5em}
EL\\
S20    A : Hardening / B : Invariants / C : Paper A : Edge cases / B : Coverage / C : Paper v1 Module 2 Clôture         Audit notes\\
S21    A : DeFi Primitives / B : DApp Intég. / C : A : Vault / B : Sign Flow / C : Tracking        Repo S21             Log/Erreurs OK\\
\vspace{0.5em}
\vspace{0.5em}
\vspace{0.5em}
\vspace{0.5em}
TI\\
Analytics\\
\vspace{0.5em}
\vspace{0.5em}
\vspace{0.5em}
\vspace{0.5em}
N\\
S22    A : Pricing / B : Oracles / C : Dashboards A : Calculs / B : Oracle mock / C : Dash v1 Dash/Indexer v1           Métriques lisibles\\
S23    A : Indexation / B : Events / C : Cohortes A : Events / B : Endpoints / C : Funnels         ”Data pack”          Schéma doc\\
\vspace{0.5em}
\vspace{0.5em}
\vspace{0.5em}
\vspace{0.5em}
DE\\
S24    A : DeFi Pack / B : L2 Scale / C : Analytics A : Stable / B : Archi. L2 / C : Final Instr. Module 3 Clôture      Démo 10 min\\
Final\\
\vspace{0.5em}
\vspace{0.5em}
\vspace{0.5em}
\vspace{0.5em}
FI\\
S25    A : Prod Perf / B : Sec Hardening / C : Gouv. A : Bench / B : Threat Model / C : Rules      Repo S25 + TM        TM présent, CI verte\\
Design\\
\vspace{0.5em}
\vspace{0.5em}
\vspace{0.5em}
\vspace{0.5em}
N\\
S26    A : Monitoring / B : Correction \textbackslash{}& Preuve / A : Metrics / B : Fix findings / C : Launch ”Ops pack”\\
CO                                                                     Runbook incidents\\
C : GTM                                         plan\\
S27    A : Wallet Intég. / B : Audit Notes / C : Play- A : Adapters / B : Audit pack / C : Metrics Release Cand.        Tag RC, Changelog\\
book\\
S28    Block Check “Build”                             Stabilisation + Docs                        PACK Bâtisseur       GO/WARN/NO-GO\\
\vspace{0.5em}
\vspace{0.5em}
\vspace{0.5em}
\vspace{0.5em}
\begin{confidential}CONFIDENTIEL — PROJET RBK 2.0                                                Page 227 sur 316\end{confidential}
Annexe B. SYLLABUS TECHNIQUE\\
RBK 2.0 — Whitepaper Architecte Web3                                                                                    DÉTAILLÉ ET STANDARDS\\
\vspace{0.5em}
\vspace{0.5em}
Architecte : Studio (S29-S48)\\
\vspace{0.5em}
Sem    Objectifs (Sprint)                  Labs (Production)            Livrable               DoD\\
S29    Capstone 1 : Reliability, ADR, TM   ADR \textbackslash{}#3, Threat Model, Spec   Spec + ADR + TM        Menaces mitigées\\
S30    Sessions, RPC Fallback              Impl. fallback, retries      Module Session         Pas de blocage\\
S31    Tx Builder, Simulation              Builder, Idempotence         Tx Pipeline v1         Scénarios test\\
S32    Observabilité, Support              Metrics, Traces, Playbook    Obs. Pack              Dashboards\\
S33    GATE CAPSTONE 1                     Tests E2E, Démo filmée       Reliability Pack       Démo stable\\
S34    Capstone 2 : Protocole On-chain     Invariants, Cas limites      Spec + Invariants      Stratégie tests\\
S35    Implémentation Core                 Core v1, Unit tests          Core v1                CI verte\\
S36    Sécurité                            Simu attaques, Mitigations   Security Checklist     Preuves fixes\\
S37    Audit Readiness                     Scripts repro, Diagrammes    Audit Pack v1          Repro 1 cmd\\
S38    GATE CAPSTONE 2                     Hardening, Perf              RC + Audit Pack        RC Tag\\
S39    Capstone 3 : Ops, SLO               SLO, Alerting                SLO + Dashboards       Alertes utiles\\
\vspace{0.5em}
\vspace{0.5em}
\vspace{0.5em}
\vspace{0.5em}
EL\\
S40    Incident Simu                       Chaos monkey, Drill          Postmortem Simulé      Actions corr.\\
S41    Deploy Prod-like                    Pipeline, Migration          Ops Playbook v1        Rollback doc\\
S42    Perf \textbackslash{}& Cout                         Profiling, Optim             Rapport Perf           Mesures\\
\vspace{0.5em}
\vspace{0.5em}
\vspace{0.5em}
\vspace{0.5em}
TI\\
S43    GATE CAPSTONE 3                     Go-live simulé               Go-Live Pack           Runbook\\
\vspace{0.5em}
\vspace{0.5em}
\vspace{0.5em}
\vspace{0.5em}
N\\
S44    Placement 1 : GitHub                Nettoyage, README            Portfolio Pack         Proof links\\
S45    Placement 2 : CV/LinkedIn           ”Project one-liners”         CV + LinkedIn          Claims vérifiables\\
\vspace{0.5em}
\vspace{0.5em}
\vspace{0.5em}
\vspace{0.5em}
DE\\
S46    Placement 3 : Demos                 Script, Vidéo, Slides        Demo Kit               Timing respecté\\
S47    Placement 4 : Interviews            Mock, System Design          Interview Notes        Réponses rigoureuses\\
\vspace{0.5em}
\vspace{0.5em}
\vspace{0.5em}
\vspace{0.5em}
FI\\
S48    Block Check “Launch Proof”          Soutenance                   PACK FINAL             GO MARCHÉ\\
\vspace{0.5em}
\vspace{0.5em}
\vspace{0.5em}
\vspace{0.5em}
N\\
CO\\
\vspace{0.5em}
\vspace{0.5em}
\vspace{0.5em}
\vspace{0.5em}
\begin{confidential}CONFIDENTIEL — PROJET RBK 2.0                                       Page 228 sur 316\end{confidential}
C \textbar{} ANNEXE B — Modèle financier \textbackslash{}& Waterfall contractuel\\
(RBK – Nexus Réussite)\\
\vspace{0.5em}
\vspace{0.5em}
\vspace{0.5em}
\vspace{0.5em}
EL\\
C.1 Référentiel financier contractuel\\
\vspace{0.5em}
\vspace{0.5em}
\vspace{0.5em}
\vspace{0.5em}
TI\\
N\\
C.1.1 Définitions comptables\\
\vspace{0.5em}
\vspace{0.5em}
\vspace{0.5em}
\vspace{0.5em}
DE\\
Pour l’application des modalités financières, les définitions suivantes sont contractuellement contraignantes :\\
\vspace{0.5em}
\vspace{0.5em}
\vspace{0.5em}
\vspace{0.5em}
FI\\
• CA\textbackslash{}_Booked : Valeur catalogue totale signée (Genesis Pool + Explorer + Bâtisseur/Architecte selon conversions validées)\\
\vspace{0.5em}
\vspace{0.5em}
N\\
CO\\
• CA\textbackslash{}_Cash : Montants effectivement encaissés (paiements directs + recouvrements ISA)\\
\vspace{0.5em}
• CA\textbackslash{}_Eligible : CA\textbackslash{}_Cash net des refunds/chargebacks/remises\\
\vspace{0.5em}
• Coûts\textbackslash{}_Directs : Mentor Pool (refacturé) + outils explicitement listés dans l’addendum technique\\
\vspace{0.5em}
• Marge\textbackslash{}_Projet : CA\textbackslash{}_Eligible – Coûts\textbackslash{}_Directs – Franchise\textbackslash{}_RBK\\
\vspace{0.5em}
\vspace{0.5em}
\vspace{0.5em}
\vspace{0.5em}
229\\
RBK 2.0 — Whitepaper Architecte Web3                                                                                     Annexe C. Annexe B — Modèle financier\\
\vspace{0.5em}
\vspace{0.5em}
C.1.2 Synthèse des montants contractuels (TND)\\
\vspace{0.5em}
\vspace{0.5em}
\begin{center}\begin{tabular}TAB. C.1 : Conversion locale des principaux montants\end{tabular}\end{center}
\vspace{0.5em}
\vspace{0.5em}
Poste                        Montant TND           Commentaires\\
Tuition annuelle cycle       15 900 TND            Inclut Genesis, Explorer, Bâtisseur et Architecte.\\
complet\\
Forfait Nexus                3 300 TND             Payable à l’entrée, complémentaire au partage 30\textbackslash{}%.\\
Launchpad\\
Management Fee               4 500 TND             Facturation fixe pour le pilotage sur 12 mois.\\
mensuel\\
\vspace{0.5em}
\vspace{0.5em}
\vspace{0.5em}
\vspace{0.5em}
EL\\
Budget acquisition           3 000 TND             Dépenses marketing approuvées et remboursables.\\
\vspace{0.5em}
\vspace{0.5em}
\vspace{0.5em}
\vspace{0.5em}
TI\\
mensuel\\
\vspace{0.5em}
\vspace{0.5em}
\vspace{0.5em}
\vspace{0.5em}
N\\
DE\\
FI\\
C.2 Cohorte 1 — Paramètres imposés + Minimum garanti\\
\vspace{0.5em}
\vspace{0.5em}
N\\
Taille minimale contractuelle de cohorte (Genesis Pool) : 15 apprenants.\\
CO\\
En cas de souscription inférieure à 15 apprenants, RBK paiera à Nexus un ”ajustement minimum” défini comme suit :\\
\vspace{0.5em}
\vspace{0.5em}
Ajustement = (15 − Genesisréel ) × 2 900 TND\\
\vspace{0.5em}
Payable au démarrage de la Genesis Pool. Objectif : assurer la viabilité du delivery et éliminer le risque ”petite cohorte”.\\
\vspace{0.5em}
\vspace{0.5em}
\vspace{0.5em}
C.3 Rémunération de Nexus — Structure à 4 composantes\\
Le modèle de compensation RBK → Nexus Réussite comprend 4 briques contractuelles, toutes obligatoires :\\
(1) Management Fee (forfait de pilotage)\\
\vspace{0.5em}
• 4 500 TND/mois, pendant 12 mois (cohorte 48 semaines)\\
\vspace{0.5em}
\vspace{0.5em}
\begin{confidential}CONFIDENTIEL — PROJET RBK 2.0                                            Page 230 sur 316\end{confidential}
RBK 2.0 — Whitepaper Architecte Web3                                                                                    Annexe C. Annexe B — Modèle financier\\
\vspace{0.5em}
\vspace{0.5em}
• Payable d’avance : le 1er de chaque mois, échéance à 15 jours max\\
\vspace{0.5em}
• Non remboursable sauf faute lourde prouvée de Nexus\\
\vspace{0.5em}
(2) Delivery Fee (part variable liée au volume)\\
\vspace{0.5em}
• DeliveryFee = 12\textbackslash{}% × CA\textbackslash{}_Booked (calculé sur la valeur catalogue signée)\\
\vspace{0.5em}
• Paiement : mensuel au prorata de l’avancement (Genesis sur mois 1–4, Explorer sur mois 5–8, Bâtisseur/Architecte sur mois 9–12)\\
\vspace{0.5em}
• Formule de proratisation : DeliveryFee\textbackslash{}_mois = DeliveryFee × (1/4) sur les 4 mois du niveau concerné\\
\vspace{0.5em}
(3) Mentor Pool (refacturation encadrée)\\
\vspace{0.5em}
• Nexus sous-traite les mentors et refacture à RBK ”au coût réel”, sur justificatifs.\\
\vspace{0.5em}
\vspace{0.5em}
\vspace{0.5em}
\vspace{0.5em}
EL\\
• CAP MentorPool : 25\textbackslash{}% × CA\textbackslash{}_Booked de la cohorte (plafond annuel).\\
\vspace{0.5em}
\vspace{0.5em}
\vspace{0.5em}
\vspace{0.5em}
TI\\
• Paiement : mensuel sur facture, délai 15 jours.\\
\vspace{0.5em}
\vspace{0.5em}
\vspace{0.5em}
\vspace{0.5em}
N\\
• Les coûts mentors sont exclus de la marge Nexus (ligne pass-through).\\
\vspace{0.5em}
\vspace{0.5em}
\vspace{0.5em}
\vspace{0.5em}
DE\\
(4) Success Fee (bonus placement / performance, simple et auditable)\\
\vspace{0.5em}
\vspace{0.5em}
\vspace{0.5em}
\vspace{0.5em}
FI\\
• 1 000 TND par apprenant placé (contrat ≥ 6 mois OU mission ≥ 3 mois), placement vérifié.\\
\vspace{0.5em}
\vspace{0.5em}
\vspace{0.5em}
N\\
• Facturable à la preuve (contrat / attestation), payable à 15 jours.\\
CO\\
C.4 Waterfall (ordre de distribution) — Clause contractuelle\\
Définir la ”Franchise\textbackslash{}_RBK” mensuelle :\\
\vspace{0.5em}
Franchise\textbackslash{}_RBK = 3 000 TND /mois + frais de paiement (si listés) + outils listés (si listés)\\
\vspace{0.5em}
Puis écrire l’ordre :\\
\vspace{0.5em}
1. Franchise\textbackslash{}_RBK\\
\vspace{0.5em}
2. Mentor Pool (coût réel, dans le cap)\\
\vspace{0.5em}
\vspace{0.5em}
\vspace{0.5em}
\begin{confidential}CONFIDENTIEL — PROJET RBK 2.0                                               Page 231 sur 316\end{confidential}
RBK 2.0 — Whitepaper Architecte Web3                           Annexe C. Annexe B — Modèle financier\\
\vspace{0.5em}
\vspace{0.5em}
3. Management Fee Nexus\\
\vspace{0.5em}
4. Delivery Fee Nexus\\
\vspace{0.5em}
5. Success Fees Nexus\\
\vspace{0.5em}
6. Résiduel RBK (marge RBK)\\
\vspace{0.5em}
\vspace{0.5em}
\vspace{0.5em}
C.5 Table de projections (OBLIGATOIRE) — 3 scénarios fixes\\
\vspace{0.5em}
\vspace{0.5em}
\vspace{0.5em}
\vspace{0.5em}
EL\\
TI\\
N\\
DE\\
FI\\
N\\
CO\\
\vspace{0.5em}
\vspace{0.5em}
\vspace{0.5em}
\vspace{0.5em}
\begin{confidential}CONFIDENTIEL — PROJET RBK 2.0               Page 232 sur 316\end{confidential}
RBK 2.0 — Whitepaper Architecte Web3                                                                            Annexe C. Annexe B — Modèle financier\\
\vspace{0.5em}
\vspace{0.5em}
\vspace{0.5em}
\vspace{0.5em}
\begin{center}\begin{tabular}TAB. C.3 : Projections financières selon taille de cohorte initiale\end{tabular}\end{center}
\vspace{0.5em}
\vspace{0.5em}
Paramètre                          Scénario A       Scénario B    Scénario C\\
Genesis Pool (taille initiale)         15               20            30\\
Explorer (projection)                   9               12            18\\
Bâtisseur/Architecte                  7,2              9,6           14,4\\
(projection)\\
CA\textbackslash{}_Booked (TND)                     160 680          214 240       321 360\\
CAP MentorPool (25\textbackslash{}%                 40 170            53 560        80 340\\
CA\textbackslash{}_Booked)\\
\vspace{0.5em}
\vspace{0.5em}
\vspace{0.5em}
\vspace{0.5em}
EL\\
Management Fee (total)                12                12            12\\
\vspace{0.5em}
\vspace{0.5em}
\vspace{0.5em}
\vspace{0.5em}
TI\\
mensualités       mensualités   mensualités\\
\vspace{0.5em}
\vspace{0.5em}
\vspace{0.5em}
\vspace{0.5em}
N\\
4 500 TND         4 500 TND     4 500 TND\\
\vspace{0.5em}
\vspace{0.5em}
\vspace{0.5em}
\vspace{0.5em}
DE\\
Delivery Fee (12\textbackslash{}%                 19 281,60         25 708,80     38 563,20\\
CA\textbackslash{}_Booked)\\
\vspace{0.5em}
\vspace{0.5em}
\vspace{0.5em}
\vspace{0.5em}
FI\\
Acquisition (total)                   12                12            12\\
\vspace{0.5em}
\vspace{0.5em}
\vspace{0.5em}
N\\
CO                  mensualités       mensualités   mensualités\\
3 000 TND         3 000 TND     3 000 TND\\
Résiduel RBK (avant autres        11 228,40         64 971,20     111 456,80\\
overhead)\\
\vspace{0.5em}
\vspace{0.5em}
\vspace{0.5em}
Données de calcul (constantes) :\\
\vspace{0.5em}
• Taux de conversion Genesis → Explorer : 60\textbackslash{}%\\
\vspace{0.5em}
• Taux de conversion Explorer → Bâtisseur/Architecte : 80\textbackslash{}%\\
\vspace{0.5em}
• Prix Genesis Pool : 2 900 TND\\
\vspace{0.5em}
• Prix Explorer : 5 900 TND\\
\vspace{0.5em}
\vspace{0.5em}
\begin{confidential}CONFIDENTIEL — PROJET RBK 2.0                                          Page 233 sur 316\end{confidential}
RBK 2.0 — Whitepaper Architecte Web3                                                                                    Annexe C. Annexe B — Modèle financier\\
\vspace{0.5em}
\vspace{0.5em}
• Prix Bâtisseur/Architecte : 8 900 TND\\
\vspace{0.5em}
• Total catalogue : 15 900 TND\\
\vspace{0.5em}
• Delivery Fee : 12\textbackslash{}% du CA\textbackslash{}_Booked\\
\vspace{0.5em}
• CAP MentorPool : 25\textbackslash{}% du CA\textbackslash{}_Booked\\
\vspace{0.5em}
• Management Fee : 4 500 TND/mois × 12 mois = 54 000 TND\\
\vspace{0.5em}
• Acquisition : 3 000 TND/mois × 12 mois = 36 000 TND\\
\vspace{0.5em}
\vspace{0.5em}
\vspace{0.5em}
C.6 Clauses comptables \textbackslash{}& fiscales (sans chiffres hasardeux)\\
\vspace{0.5em}
\vspace{0.5em}
\vspace{0.5em}
\vspace{0.5em}
EL\\
• Tous montants : HT (hors taxes) ; TVA selon régime applicable\\
\vspace{0.5em}
\vspace{0.5em}
\vspace{0.5em}
\vspace{0.5em}
TI\\
• Clause ”gross-up” : toute retenue à la source est supportée par RBK afin que Nexus reçoive le net égal au montant facturé HT\\
\vspace{0.5em}
\vspace{0.5em}
\vspace{0.5em}
\vspace{0.5em}
N\\
• Délai de paiement : 15 jours ouvrés\\
\vspace{0.5em}
\vspace{0.5em}
\vspace{0.5em}
\vspace{0.5em}
DE\\
• Pénalités de retard : 1\textbackslash{}%/mois + suspension du delivery en cas de défaut \textgreater{} 30 jours\\
\vspace{0.5em}
\vspace{0.5em}
\vspace{0.5em}
\vspace{0.5em}
FI\\
• Droit d’audit : RBK peut auditer les justificatifs mentors ; Nexus peut auditer le calcul CA\textbackslash{}_Booked/CA\textbackslash{}_Eligible\\
\vspace{0.5em}
\vspace{0.5em}
\vspace{0.5em}
N\\
CO\\
C.7 Garanties \textbackslash{}& protections contractuelles\\
• Clause de ”minimum garanti” (cf. section 2) : protection contre petite cohorte\\
\vspace{0.5em}
• CAP MentorPool : limitation du risque de surcoût personnel\\
\vspace{0.5em}
• Franchise\textbackslash{}_RBK : allocation préférentielle des charges fixes d’acquisition\\
\vspace{0.5em}
• Modalités de reporting : reporting mensuel détaillé des indicateurs financiers clés\\
\vspace{0.5em}
\vspace{0.5em}
\vspace{0.5em}
\vspace{0.5em}
\begin{confidential}CONFIDENTIEL — PROJET RBK 2.0                                            Page 234 sur 316\end{confidential}
D \textbar{} ADDENDUM JURIDIQUE — Conditions générales com-\\
plémentaires (RBK – Nexus Réussite)\\
\vspace{0.5em}
\vspace{0.5em}
\vspace{0.5em}
\vspace{0.5em}
EL\\
D.1 Périmètre contractuel RBK–Nexus\\
\vspace{0.5em}
\vspace{0.5em}
\vspace{0.5em}
\vspace{0.5em}
TI\\
N\\
Le présent addendum fait office de cadre de services entre RBK et Nexus Réussite. Il reprend les engagements opérationnels décrits au chapitre 5 et les décline\\
\vspace{0.5em}
\vspace{0.5em}
\vspace{0.5em}
\vspace{0.5em}
DE\\
sous forme d’obligations contractuelles opposables aux deux parties.\\
\vspace{0.5em}
\vspace{0.5em}
\vspace{0.5em}
\vspace{0.5em}
FI\\
D.1.1 Livrables contractuels par phase\\
\vspace{0.5em}
\vspace{0.5em}
N\\
CO\\
Livrables Nexus par phase opérationnelle\\
\vspace{0.5em}
\vspace{0.5em}
\vspace{0.5em}
Phase                 Livrables obligatoires    Format échéance\\
Sprint 0 (J -90 →     Avis juridique            Dossiers signés, dépôt Git tagué, accès LMS sandbox, enregistrements vidéo, protocole Wellbeing publié.\\
J -1)                 opérationnel,\\
curriculum Genesis\\
validé, MVP LMS\\
jouable, kit\\
Mentor-in-a-Box v1,\\
plan Wellbeing activé\\
\vspace{0.5em}
\vspace{0.5em}
\vspace{0.5em}
\vspace{0.5em}
235\\
RBK 2.0 — Whitepaper Architecte Web3                                                                                             Annexe D. Addendum juridique\\
\vspace{0.5em}
\vspace{0.5em}
\vspace{0.5em}
Phase 1 (J1 →         Dureté LMS (S1-S12),       Rapports QA hebdo, scripts CI validés, contrats mentors archivés, Notion playbooks, spec simulateur validée\\
J60)                  pipeline CI/CD \textbackslash{}&           par RBK.\\
sécurité, roster mentors\\
complet, playbooks\\
Wellbeing, simulateur\\
ROI pré-production\\
Phase 2 (J61 →        Piscine Rust exécutée,     Export CRM, dashboards fatigue, calendrier contenus, scoring candidats, CPPS en pré-signature.\\
J90)                  reporting fatigue,\\
campagne Building in\\
Public, shortlist\\
candidats Alpha,\\
dossiers CPPS prêts\\
\vspace{0.5em}
\vspace{0.5em}
\vspace{0.5em}
\vspace{0.5em}
EL\\
Phase 3 (J91 →        Onboarding promo           Procès-verbaux d’onboarding, agendas rituels, cockpit KPI, comptes rendus comité, synthèse Wellbeing.\\
J120)                 Alpha, rituels cadencés,\\
\vspace{0.5em}
\vspace{0.5em}
\vspace{0.5em}
\vspace{0.5em}
TI\\
reporting KPI hebdo,\\
\vspace{0.5em}
\vspace{0.5em}
\vspace{0.5em}
\vspace{0.5em}
N\\
comité qualité, retours\\
\vspace{0.5em}
\vspace{0.5em}
\vspace{0.5em}
\vspace{0.5em}
DE\\
Wellbeing consolidés\\
Cycle continu (post   Mise à jour contenu,       Release notes mensuelles, attestation formation mentors, rapport alumni, backlog priorisé.\\
\vspace{0.5em}
\vspace{0.5em}
\vspace{0.5em}
\vspace{0.5em}
FI\\
J120)                 train-the-trainer,\\
\vspace{0.5em}
\vspace{0.5em}
\vspace{0.5em}
N\\
support alumni,\\
amélioration continue\\
CO\\
Chaque livrable est considéré comme une obligation de moyens renforcée assortie de critères d’acceptation documentés et signés en comité exécutif conjoint.\\
\vspace{0.5em}
\vspace{0.5em}
D.1.2 Domaines transverses\\
• Pilotage pédagogique : maintien du tronc commun, animation des Block Checks et publication des DoD par module.\\
\vspace{0.5em}
• Qualité mentors : ratios mentors / apprenants respectés, audit mensuel des performances et activation du plan ”Mentor-in-a-Box”.\\
\vspace{0.5em}
• Ops apprenants \textbackslash{}& Wellbeing : traitement des tickets sous 24h, exploitation des indicateurs de fatigue, déclenchement des temps de repos via le Responsable\\
Bien-être.\\
\vspace{0.5em}
\vspace{0.5em}
\vspace{0.5em}
\vspace{0.5em}
\begin{confidential}CONFIDENTIEL — PROJET RBK 2.0                                            Page 236 sur 316\end{confidential}
RBK 2.0 — Whitepaper Architecte Web3                                                                                                 Annexe D. Addendum juridique\\
\vspace{0.5em}
\vspace{0.5em}
• Activation communauté : animation Discord, publication hebdomadaire Building in Public, coordination marketing RBK.\\
\vspace{0.5em}
• Reporting commun : alimentation du cockpit décisionnel (CABooked , CAEligible , NPS, taux de complétion) et archivage des pièces justificatives.\\
\vspace{0.5em}
• Gestion des risques : mise à jour du registre commun, déclenchement des plans de continuité, documentation post-mortem.\\
\vspace{0.5em}
• Amélioration continue : rétro-planning trimestriel des évolutions, synthèse des feedbacks apprenants, proposition des expérimentations pédagogiques.\\
\vspace{0.5em}
\vspace{0.5em}
\vspace{0.5em}
D.2 KPIs contractuels et bonus/malus\\
Indicateurs de performance Nexus\\
\vspace{0.5em}
\vspace{0.5em}
\vspace{0.5em}
\vspace{0.5em}
EL\\
KPI                           Définition \textbackslash{}& mode de mesure            Seuils contractuels\\
Qualité pédagogique           NPS cohortes, taux réussite Labs,      NPS ≥ 50, incidents critiques = 0, Labs complétés ≥ 90\textbackslash{}%.\\
\vspace{0.5em}
\vspace{0.5em}
\vspace{0.5em}
\vspace{0.5em}
TI\\
incidents critiques\\
\vspace{0.5em}
\vspace{0.5em}
\vspace{0.5em}
\vspace{0.5em}
N\\
Disponibilité mentors         Couverture slots mentoring, délai      Couverture ≥ 95\textbackslash{}%, délai ≤ 12h ouvrées.\\
\vspace{0.5em}
\vspace{0.5em}
\vspace{0.5em}
\vspace{0.5em}
DE\\
réponse Discord\\
Engagement apprenants         Taux présence rituels, complétion      Présence ≥ 85\textbackslash{}%, plans complétés ≥ 90\textbackslash{}%.\\
\vspace{0.5em}
\vspace{0.5em}
\vspace{0.5em}
\vspace{0.5em}
FI\\
plans Wellbeing\\
\vspace{0.5em}
\vspace{0.5em}
\vspace{0.5em}
N\\
Respect jalons                Respect planning livrables et SOW\\
CO                            ≥ 95\textbackslash{}% des jalons validés à date prévue.\\
Reporting financier           Mises à jour CABooked /CAEligible ,    Reporting mensuel livré J+5, écarts ≤ 2\textbackslash{}%.\\
facturation\\
\vspace{0.5em}
\vspace{0.5em}
• Bonus performance : +5\textbackslash{}% à +10\textbackslash{}% sur la facture jalon lorsque trois KPI clés dépassent le niveau ”Stretch” défini en comité.\\
\vspace{0.5em}
• Malus : –10\textbackslash{}% et plan correctif sous 15 jours si deux KPI tombent sous les seuils plancher ; possibilité de suspension de jalon.\\
\vspace{0.5em}
• Clause de revoyure : recalibrage semestriel des KPI et des niveaux de service sur décision conjointe CEC.\\
\vspace{0.5em}
\vspace{0.5em}
\vspace{0.5em}
D.3 Clauses de propriété intellectuelle \textbackslash{}& licences\\
\vspace{0.5em}
\vspace{0.5em}
\vspace{0.5em}
\begin{confidential}CONFIDENTIEL — PROJET RBK 2.0                                              Page 237 sur 316\end{confidential}
RBK 2.0 — Whitepaper Architecte Web3                                                                                                  Annexe D. Addendum juridique\\
\vspace{0.5em}
\vspace{0.5em}
\vspace{0.5em}
• Les contenus pédagogiques, playbooks et assets produits par Nexus sont co-propriété RBK–Nexus ; RBK bénéficie d’une licence perpétuelle, mondiale et libre de\\
redevance pour exploitation interne.\\
\vspace{0.5em}
• Les outils, scripts et gabarits publiés dans l’écosystème (GitHub) conservent une licence OSS compatible, avec attribution obligatoire aux deux parties.\\
\vspace{0.5em}
• Money Factory AI demeure propriétaire de Venture Engine ; Nexus garantit la réversibilité des données et l’accès continu de RBK conformément à la licence.\\
\vspace{0.5em}
• Toute réutilisation par des tiers nécessite accord écrit des deux parties et partage des revenus dérivés selon le waterfall financier (cf. C).\\
\vspace{0.5em}
\vspace{0.5em}
\vspace{0.5em}
\vspace{0.5em}
EL\\
D.4 Processus de pilotage et résolution des litiges\\
1. Détection : Ecarts détectés via cockpit KPI, audits ou feedback apprenants sont consignés dans le registre commun sous 48h.\\
\vspace{0.5em}
\vspace{0.5em}
\vspace{0.5em}
\vspace{0.5em}
TI\\
2. Remédiation : Nexus propose un plan correctif (RACI, budget, échéance) validé en CEC ; RBK confirme ou demande un arbitrage.\\
\vspace{0.5em}
\vspace{0.5em}
\vspace{0.5em}
\vspace{0.5em}
EN\\
3. Escalade : En cas de désaccord, réunion extraordinaire du Comité éthique et pédagogique sous 5 jours, avec convocation optionnelle de MFAI pour sujets\\
techniques.\\
\vspace{0.5em}
4. Arbitrage final : Si le différend persiste, médiation externe désignée conjointement ; défaut d’accord sous 30 jours déclenche clause de résiliation anticipée.\\
\vspace{0.5em}
\vspace{0.5em}
\vspace{0.5em}
\vspace{0.5em}
D\\
D.5 Référentiel financier\\
FI\\
N\\
Toutes questions relatives aux éléments financiers contractuels (tarification, calendrier de paiement, rémunération, reporting, audits) sont régies par l’annexe B\\
— Modèle financier \textbackslash{}& Waterfall contractuel (cf. C), qui constitue un élément contractuel intégral du présent contrat.\\
O\\
C\\
\vspace{0.5em}
\vspace{0.5em}
D.6 Modalités de paiement\\
• Délai de paiement : 15 jours ouvrés à compter de la date de facture\\
\vspace{0.5em}
• Les factures sont émises mensuellement au prorata de l’avancement\\
\vspace{0.5em}
• En cas de retard \textgreater{} 30 jours, suspension automatique des services\\
\vspace{0.5em}
• Pénalités de retard : 1\textbackslash{}% du montant impayé par mois de retard\\
\vspace{0.5em}
\vspace{0.5em}
\vspace{0.5em}
\begin{confidential}CONFIDENTIEL — PROJET RBK 2.0                                              Page 238 sur 316\end{confidential}
RBK 2.0 — Whitepaper Architecte Web3                                                                                           Annexe D. Addendum juridique\\
\vspace{0.5em}
\vspace{0.5em}
\vspace{0.5em}
D.7 Clause de ”Gross-Up”\\
Toutes retenues à la source, taxes ou impôts applicables sur les paiements dus à Nexus seront supportés par RBK, de sorte que Nexus reçoive in fine le montant\\
net égal au montant HT facturé.\\
\vspace{0.5em}
\vspace{0.5em}
\vspace{0.5em}
D.8 Droit d’audit\\
• RBK dispose d’un droit d’audit annuel sur les justificatifs des coûts mentors\\
\vspace{0.5em}
• Nexus dispose d’un droit d’accès aux données de CA\textbackslash{}_Booked et CA\textbackslash{}_Eligible pour vérification des calculs\\
\vspace{0.5em}
• Les audits sont effectués sur 30 jours calendriers à compter de la demande\\
\vspace{0.5em}
\vspace{0.5em}
\vspace{0.5em}
\vspace{0.5em}
EL\\
• Les coûts des audits sont supportés par la partie demandant l’audit, sauf si des écarts \textgreater{} 5\textbackslash{}% sont détectés\\
\vspace{0.5em}
\vspace{0.5em}
\vspace{0.5em}
\vspace{0.5em}
TI\\
N\\
D.9 Responsabilités contractuelles\\
\vspace{0.5em}
\vspace{0.5em}
\vspace{0.5em}
\vspace{0.5em}
DE\\
• RBK : respect des délais de paiement, communication des données CA\textbackslash{}_Booked\\
\vspace{0.5em}
\vspace{0.5em}
\vspace{0.5em}
\vspace{0.5em}
FI\\
• Nexus : fourniture conforme des services pédagogiques, reporting financier mensuel\\
\vspace{0.5em}
\vspace{0.5em}
\vspace{0.5em}
N\\
• Les deux parties : coopération pour l’audit mutuel et le reporting financier\\
CO\\
D.10 Force majeure\\
Est considéré comme un cas de force majeure tout événement échappant au contrôle raisonnable de la partie concernée, imprévisible et insurmontable malgré\\
l’application de toutes les diligences raisonnables.\\
\vspace{0.5em}
\vspace{0.5em}
\vspace{0.5em}
D.11 Résiliation anticipée\\
En cas de résiliation anticipée pour faute grave ou non-respect des engagements contractuels :\\
\vspace{0.5em}
• Arrêt immédiat des services\\
\vspace{0.5em}
\vspace{0.5em}
\begin{confidential}CONFIDENTIEL — PROJET RBK 2.0                                            Page 239 sur 316\end{confidential}
RBK 2.0 — Whitepaper Architecte Web3                                                    Annexe D. Addendum juridique\\
\vspace{0.5em}
\vspace{0.5em}
• Paiement des sommes dues à la date d’arrêt\\
\vspace{0.5em}
• Compensation pour préavis non respecté selon modalités contractuelles\\
\vspace{0.5em}
\vspace{0.5em}
\vspace{0.5em}
\vspace{0.5em}
EL\\
TI\\
N\\
DE\\
FI\\
N\\
CO\\
\vspace{0.5em}
\vspace{0.5em}
\vspace{0.5em}
\vspace{0.5em}
\begin{confidential}CONFIDENTIEL — PROJET RBK 2.0                                        Page 240 sur 316\end{confidential}
E \textbar{} TEMPLATE DE RAPPORT D’AUDIT DE SÉCURITÉ\\
\vspace{0.5em}
\vspace{0.5em}
\vspace{0.5em}
\vspace{0.5em}
EL\\
Un rapport d’audit professionnel doit être clair, complet et actionnable.\\
\vspace{0.5em}
\vspace{0.5em}
\vspace{0.5em}
\vspace{0.5em}
TI\\
E.1 Structure du Rapport\\
\vspace{0.5em}
\vspace{0.5em}
\vspace{0.5em}
\vspace{0.5em}
EN\\
1. Executive Summary : Résumé pour les décideurs (Score, Risque global).\\
\vspace{0.5em}
2. Scope : Liste des fichiers audités et Commit Hash.\\
\vspace{0.5em}
\vspace{0.5em}
\vspace{0.5em}
\vspace{0.5em}
D\\
3. Findings : Liste des vulnérabilités classées par sévérité.\\
\vspace{0.5em}
4. Recommendations : Conseils d’architecture généraux.\\
\vspace{0.5em}
FI\\
N\\
E.2 Classification des Risques\\
O\\
\vspace{0.5em}
Table : Échelle de Sévérité\\
C\\
\vspace{0.5em}
\vspace{0.5em}
Niveau   Impact                                         Exemple\\
CRITICAL Perte de fonds directe, Gel définitif          Reentrancy, Owner Key compromise\\
HIGH     Dégradation sévère du service, Perte partielle DoS, Price Oracle manipulation\\
MEDIUM   Grief mineur, Coût Gas élevé                   Griefing attack, Unbounded Loop\\
LOW/INFO Bonnes pratiques, Lisibilité                   Typo, Dead code\\
\vspace{0.5em}
\vspace{0.5em}
\vspace{0.5em}
\vspace{0.5em}
241\\
RBK 2.0 — Whitepaper Architecte Web3                                                             Annexe E. TEMPLATE DE RAPPORT D’AUDIT DE SÉCURITÉ\\
\vspace{0.5em}
\vspace{0.5em}
\vspace{0.5em}
E.3 Fiche Finding Type\\
ID-01 : Unchecked External Call (H-01)\\
\vspace{0.5em}
Sévérité : HIGH\\
Fichier : vault.rs\\
Description : L’appel CPI vers le programme Token ne vérifie pas le code retour. Impact : Un attaquant peut forcer l’échec silencieux du transfert et créditer\\
son solde interne. Recommandation : Utiliser anchor\textbackslash{}_lang::solana\textbackslash{}_program::program::invoke\textbackslash{}_signed et gérer le Result.\\
\vspace{0.5em}
\vspace{0.5em}
\vspace{0.5em}
\vspace{0.5em}
EL\\
TI\\
N\\
DE\\
FI\\
N\\
CO\\
\vspace{0.5em}
\vspace{0.5em}
\vspace{0.5em}
\vspace{0.5em}
\begin{confidential}CONFIDENTIEL — PROJET RBK 2.0                                           Page 242 sur 316\end{confidential}
F \textbar{} LE COCKPIT DE L’ARCHITECTE\\
\vspace{0.5em}
\vspace{0.5em}
Liste des outils obligatoires pour un étudiant en phase de production.\\
\vspace{0.5em}
\vspace{0.5em}
\vspace{0.5em}
\vspace{0.5em}
EL\\
F.1 Stack Outillage Minimal\\
\vspace{0.5em}
\vspace{0.5em}
\vspace{0.5em}
\vspace{0.5em}
TI\\
N\\
Table : Cockpit Tools\\
\vspace{0.5em}
\vspace{0.5em}
\vspace{0.5em}
\vspace{0.5em}
DE\\
Outil           Usage           Output Attendu\\
\vspace{0.5em}
\vspace{0.5em}
\vspace{0.5em}
\vspace{0.5em}
FI\\
Obsidian/Notion Knowledge Base Wiki du projet, Notes de recherche\\
Excalidraw      Diagramming     Schémas d’architecture C4\\
\vspace{0.5em}
\vspace{0.5em}
N\\
Linear/Jira     Task Management Tickets spécifiés et trackés\\
CO\\
Cursor/VSCode IDE               Code avec Linter et Copilot configuré\\
\vspace{0.5em}
\vspace{0.5em}
\vspace{0.5em}
F.2 Journée Type (Productivité)\\
• 09h-12h (Deep Work) : Coding (Feature complexe ou Refactoring). Pas de notifs.\\
\vspace{0.5em}
• 13h-14h (Review) : Code Review des PRs des collègues.\\
\vspace{0.5em}
• 14h-16h (Ops) : Tests, Documentation, Fixes mineurs.\\
\vspace{0.5em}
• 16h-17h (Sync) : Daily Standup, Synchro Architecte.\\
\vspace{0.5em}
\vspace{0.5em}
\vspace{0.5em}
\vspace{0.5em}
243\\
G \textbar{} Modèle ISA (Income Share Agreement)\\
\vspace{0.5em}
\vspace{0.5em}
G.1 Objet et Principes\\
\vspace{0.5em}
\vspace{0.5em}
\vspace{0.5em}
\vspace{0.5em}
EL\\
L’ISA est un mécanisme de financement sélectif destiné à aligner l’école et l’étudiant : l’étudiant ne paie que s’il dépasse un seuil de revenu, et l’école accepte un\\
\vspace{0.5em}
\vspace{0.5em}
\vspace{0.5em}
\vspace{0.5em}
TI\\
risque. L’ISA est réservé aux profils validés Top Talent.\\
\vspace{0.5em}
\vspace{0.5em}
\vspace{0.5em}
\vspace{0.5em}
N\\
Répartition des responsabilités\\
\vspace{0.5em}
\vspace{0.5em}
\vspace{0.5em}
\vspace{0.5em}
DE\\
Il est fondamental de noter que RBK porte le risque de crédit (en tant que porteur légal du contrat) tandis que Nexus Réussite porte le risque d’exécution\\
\vspace{0.5em}
\vspace{0.5em}
\vspace{0.5em}
\vspace{0.5em}
FI\\
(garant de la qualité pédagogique et de l’employabilité). Cette répartition est détaillée dans l’addendum sur la gouvernance (Chapitre 3) et l’addendum juridique\\
(Section D).\\
\vspace{0.5em}
\vspace{0.5em}
N\\
CO\\
G.2 Éligibilité (Gating)\\
• Périmètre : réservé au parcours complet Genesis Pool →Explorer →Bâtisseur →Architecte.\\
\vspace{0.5em}
• Quota : nombre de places ISA limité par cohorte (ex : 30\textbackslash{}% max).\\
\vspace{0.5em}
• Sélection : top performance + validation par comité.\\
\vspace{0.5em}
\vspace{0.5em}
\vspace{0.5em}
G.3 Paramètres ISA unifiés\\
• Seuil de déclenchement : 2500 brut/mois (revenu traçable).\\
\vspace{0.5em}
\vspace{0.5em}
\vspace{0.5em}
244\\
RBK 2.0 — Whitepaper Architecte Web3                                                                         Annexe G. Modèle ISA (Income Share Agreement)\\
\vspace{0.5em}
\vspace{0.5em}
• Taux de partage : 15\textbackslash{}% du brut mensuel (si \textgreater{} seuil).\\
\vspace{0.5em}
• Cap (plafond) : 20000 total.\\
\vspace{0.5em}
• Durée maximale : 36 mensualités payées (pas de rattrapage).\\
\vspace{0.5em}
\vspace{0.5em}
\vspace{0.5em}
G.4 Règles de Pause, Chômage, Variabilité\\
• Pause automatique : si revenu brut mensuel ≤ 2500, paiement = 0.\\
\vspace{0.5em}
• Reprise : dès que revenu brut \textgreater{} 2500.\\
\vspace{0.5em}
• Variabilité : aucun rattrapage sur les mois faibles.\\
\vspace{0.5em}
\vspace{0.5em}
\vspace{0.5em}
\vspace{0.5em}
EL\\
TI\\
G.5 Cas Limites (Edge Cases)\\
\vspace{0.5em}
\vspace{0.5em}
\vspace{0.5em}
\vspace{0.5em}
N\\
1. Revenu fluctuant : paiement déclenché uniquement les mois \textgreater{} Seuil.\\
\vspace{0.5em}
\vspace{0.5em}
\vspace{0.5em}
\vspace{0.5em}
DE\\
2. Plusieurs revenus : revenu = somme des montants bruts traçables.\\
\vspace{0.5em}
\vspace{0.5em}
\vspace{0.5em}
\vspace{0.5em}
FI\\
3. Départ à l’étranger : conversion en TND au taux mensuel.\\
\vspace{0.5em}
\vspace{0.5em}
\vspace{0.5em}
N\\
CO\\
G.6 Conformité Éthique (Musharaka)\\
Le modèle est compatible finance participative : Partage de Risque (Perte pour l’école si échec) et Partage de Profit minoritaire.\\
\vspace{0.5em}
\vspace{0.5em}
\vspace{0.5em}
G.7 Exemples Chiffrés (Seuil 2500, Taux 15\textbackslash{}%)\\
\vspace{0.5em}
\vspace{0.5em}
\vspace{0.5em}
\vspace{0.5em}
\begin{confidential}CONFIDENTIEL — PROJET RBK 2.0                                           Page 245 sur 316\end{confidential}
RBK 2.0 — Whitepaper Architecte Web3                                                     Annexe G. Modèle ISA (Income Share Agreement)\\
\vspace{0.5em}
\vspace{0.5em}
\vspace{0.5em}
\vspace{0.5em}
\begin{center}\begin{tabular}TAB. G.1 : Scénarios de Remboursement\end{tabular}\end{center}
\vspace{0.5em}
\vspace{0.5em}
Scénario         RNM Mensualité Statut Final\\
A. Junior local 3 500 525       Arrêt à 36 mois (\textless{} Cap)\\
B. Profil solide 5 000 750      Cap atteint au 27ème mois\\
C. Remote        6 000 900      Cap atteint au 23ème mois\\
D. Chômage       0     0        Drop-off à 36 mois (0 TND)\\
\vspace{0.5em}
\vspace{0.5em}
\vspace{0.5em}
\vspace{0.5em}
EL\\
TI\\
N\\
DE\\
FI\\
N\\
CO\\
\vspace{0.5em}
\vspace{0.5em}
\vspace{0.5em}
\vspace{0.5em}
\begin{confidential}CONFIDENTIEL — PROJET RBK 2.0                               Page 246 sur 316\end{confidential}
H \textbar{} GUIDE DE SÉLECTION \textbackslash{}& SCORING « PISCINE RUST »\\
\vspace{0.5em}
\vspace{0.5em}
La Piscine n’est pas un cours, c’est un filtre.\\
\vspace{0.5em}
\vspace{0.5em}
\vspace{0.5em}
\vspace{0.5em}
EL\\
H.1 Grille de Scoring\\
\vspace{0.5em}
\vspace{0.5em}
\vspace{0.5em}
\vspace{0.5em}
TI\\
N\\
Le score final (sur 100) détermine l’admission. Seuil d’admission : 75/100.\\
\vspace{0.5em}
\vspace{0.5em}
\vspace{0.5em}
\vspace{0.5em}
DE\\
\begin{center}\begin{tabular}TAB. H.1 : Critères de sélection Piscine Rust\end{tabular}\end{center}
\vspace{0.5em}
\vspace{0.5em}
\vspace{0.5em}
\vspace{0.5em}
FI\\
N\\
Critère           Poids Indicateurs\\
CO\\
Aptitude Tech     40\textbackslash{}%     Progression sur les exercices Rust, Qualité du code final.\\
Résilience        30\textbackslash{}%     Capacité à rebondir après échec, Constance de l’effort.\\
Collaboration     20\textbackslash{}%     Aide apportée aux autres (Peer-learning).\\
Communication 10\textbackslash{}%         Clarté des questions posées, Respect des mentors.\\
\vspace{0.5em}
\vspace{0.5em}
\vspace{0.5em}
\vspace{0.5em}
H.1.1 Test d’Entrée Explorer (Bypass Genesis Pool)\\
Pré-requis : Maîtrise prouvée de Rust ou C++ et des concepts Blockchain de base.\\
\vspace{0.5em}
1. Théorie (45 min) : QCM statique sur l’Account Model, le Memory Management (Stack/Heap) et la Complexité Algorithmique.\\
\vspace{0.5em}
\vspace{0.5em}
\vspace{0.5em}
247\\
RBK 2.0 — Whitepaper Architecte Web3                                                             Annexe H. GUIDE DE SÉLECTION \textbackslash{}& SCORING « PISCINE RUST »\\
\vspace{0.5em}
\vspace{0.5em}
2. Pratique (3h) : ”Mini-Piscine Express”. Implémentation d’une CLI Rust qui parse un fichier binaire et signe une payload cryptographique (Ed25519). Critère\\
Éliminatoire : Absence de tests unitaires ou usage d’IA générative détecté.\\
\vspace{0.5em}
3. Entretien (15 min) : Code review live avec le Lead Instructor. Justification des choix d’allocation mémoire.\\
\vspace{0.5em}
\vspace{0.5em}
\begin{center}\begin{tabular}TAB. H.3 : Barème Admission Explorer\end{tabular}\end{center}
\vspace{0.5em}
Critère                                        Points         Attendu                                               KO si...\\
Code Quality                                     40           Rust idiomatique,                                     ‘unwrap()‘ non géré\\
Zero CLippy warnings\\
Tests                                            30           Unit tests couvrant                                   0 tests\\
les edge cases                                        fournis\\
\vspace{0.5em}
\vspace{0.5em}
\vspace{0.5em}
\vspace{0.5em}
EL\\
Architecture                                     30           Gestion des erreurs (Result)                          Code non structuré\\
Structs propres                                       ou impropre\\
\vspace{0.5em}
\vspace{0.5em}
\vspace{0.5em}
\vspace{0.5em}
TI\\
N\\
Nous classons un dossier en KO1 lorsqu’aucun des attendus critiques n’est respecté.\\
\vspace{0.5em}
\vspace{0.5em}
\vspace{0.5em}
\vspace{0.5em}
DE\\
FI\\
H.1.2 Test d’Entrée Bâtisseur/Architecte (Bypass Explorer)\\
\vspace{0.5em}
\vspace{0.5em}
N\\
CO\\
Pré-requis : Portfolio prouvant 2+ ans d’expérience sur la stack cible (Solana ou EVM).\\
\vspace{0.5em}
1. Audit Readiness : Soumission d’un repo personnel existant. Vérification des critères ”Studio” (CI/CD, Docs, Tests E2E).\\
\vspace{0.5em}
2. Exercice de Review : L’étudiant doit auditer une PR contenant 3 vulnérabilités cachées (Reentrancy, Arithmetic Overflow, Access Control).\\
\vspace{0.5em}
\vspace{0.5em}
\vspace{0.5em}
\vspace{0.5em}
1\\
Projet/Qualité — Knock-Out (KO) : statut indiquant l’échec d’un contrôle ou d’un gate de release, nécessitant correction avant validation.\\
\vspace{0.5em}
\vspace{0.5em}
\begin{confidential}CONFIDENTIEL — PROJET RBK 2.0                                              Page 248 sur 316\end{confidential}
I \textbar{} SPÉCIFICATIONS TECHNIQUES SBT\\
\vspace{0.5em}
\vspace{0.5em}
I.1 Schéma de Métadonnées (JSON)\\
\vspace{0.5em}
\vspace{0.5em}
\vspace{0.5em}
\vspace{0.5em}
EL\\
Les SBT RBK suivent le standard Metaplex Core ou ERC-721 (Non-Transferable).\\
\vspace{0.5em}
\vspace{0.5em}
\vspace{0.5em}
\vspace{0.5em}
TI\\
Listing I.1 : Metadata SBT Standard\\
\vspace{0.5em}
\vspace{0.5em}
\vspace{0.5em}
\vspace{0.5em}
N\\
\textbackslash{}{\\
\vspace{0.5em}
\vspace{0.5em}
\vspace{0.5em}
\vspace{0.5em}
DE\\
”name” : ”RBK Guardian − Cohort 1 ” ,\\
” symbol ” : ”RBKC1−G” ,\\
\vspace{0.5em}
\vspace{0.5em}
\vspace{0.5em}
\vspace{0.5em}
FI\\
” d e s c r i p t i o n ” : ” C e r t i f i e d Solana Smart C o n t r a c t Engineer . ” ,\\
\vspace{0.5em}
\vspace{0.5em}
\vspace{0.5em}
N\\
” image ” : ” h t t p s :// arweave . n e t / . . . ” ,           CO\\
” attributes ” : [\\
\textbackslash{}{ ” t r a i t \textbackslash{}_ t y p e ” : ” Track ” , ” v a l u e ” : ” Solana ” \textbackslash{}} ,\\
\textbackslash{}{ ” t r a i t \textbackslash{}_ t y p e ” : ” L e v e l ” , ” v a l u e ” : ” Gold ” \textbackslash{}} ,\\
\textbackslash{}{ ” t r a i t \textbackslash{}_ t y p e ” : ” Cohort ” , ” v a l u e ” : ” G e n e s i s 2025” \textbackslash{}} ,\\
\textbackslash{}{ ” t r a i t \textbackslash{}_ t y p e ” : ” F i n a l G r a d e ” , ” v a l u e ” : ”92/100” \textbackslash{}}\\
],\\
” properties ” : \textbackslash{}{\\
” files” : [\\
\textbackslash{}{ ” u r i ” : ” h t t p s :// g i t h u b . com/ s t u d e n t / c a p s t o n e ” , ” type ” : ” t e x t / html ” \textbackslash{}} ,\\
\textbackslash{}{ ” u r i ” : ” h t t p s :// rbk . tn / a u d i t /S12345 ” , ” type ” : ” a p p l i c a t i o n / pdf ” \textbackslash{}}\\
]\\
\textbackslash{}}\\
\textbackslash{}}\\
\vspace{0.5em}
\vspace{0.5em}
\vspace{0.5em}
249\\
RBK 2.0 — Whitepaper Architecte Web3                                                       Annexe I. SPÉCIFICATIONS TECHNIQUES SBT\\
\vspace{0.5em}
\vspace{0.5em}
\vspace{0.5em}
I.2 Processus de Vérification\\
1. Issuer Check : Vérifier que l’adresse émettrice est bien le Multisig RBK Certifié.\\
\vspace{0.5em}
2. Owner Check : L’étudiant prouve qu’il possède le wallet (Signature message).\\
\vspace{0.5em}
3. Content Check : Le lien vers le rapport d’audit correspond au hash stocké on-chain.\\
\vspace{0.5em}
\vspace{0.5em}
\vspace{0.5em}
\vspace{0.5em}
EL\\
TI\\
N\\
DE\\
FI\\
N\\
CO\\
\vspace{0.5em}
\vspace{0.5em}
\vspace{0.5em}
\vspace{0.5em}
\begin{confidential}CONFIDENTIEL — PROJET RBK 2.0                                           Page 250 sur 316\end{confidential}
J \textbar{} DASHBOARD DE SUIVI PROMO\\
\vspace{0.5em}
\vspace{0.5em}
J.1 Indicateurs Hebdomadaires (KPI)\\
\vspace{0.5em}
\vspace{0.5em}
\vspace{0.5em}
\vspace{0.5em}
EL\\
Les KPI suivants pilotent la santé hebdomadaire de la promo.\\
\vspace{0.5em}
\vspace{0.5em}
\vspace{0.5em}
\vspace{0.5em}
TI\\
Table : Métriques de Santé Promo\\
\vspace{0.5em}
\vspace{0.5em}
\vspace{0.5em}
\vspace{0.5em}
N\\
DE\\
Catégorie   KPI             Formule                               Cible\\
Progression Velocity        Nb exercices validés / Nb total       \textgreater{} 90\textbackslash{}%\\
\vspace{0.5em}
\vspace{0.5em}
\vspace{0.5em}
\vspace{0.5em}
FI\\
Qualité     First Time Pass \textbackslash{}% Labs validés du 1er coup            \textgreater{} 50\textbackslash{}%\\
\vspace{0.5em}
\vspace{0.5em}
N\\
Engagement Attendance\\
CO        Taux présence Dailies                 \textgreater{} 95\textbackslash{}%\\
Moral       NPS Hebdo       ”Recommanderiez-vous cette semaine ?” \textgreater{} 8/10\\
\vspace{0.5em}
\vspace{0.5em}
\vspace{0.5em}
J.2 Questionnaire Bien-être Minimal\\
Envoyé chaque vendredi via Bot Discord (Anonyme).\\
\vspace{0.5em}
1. Niveau de stress (1-5) ?\\
\vspace{0.5em}
2. Charge de travail (Trop faible / OK / Trop forte) ?\\
\vspace{0.5em}
3. Sentiment de progression (Je stagne / J’apprends / Je vole) ?\\
\vspace{0.5em}
\vspace{0.5em}
\vspace{0.5em}
\vspace{0.5em}
251\\
K \textbar{} OFFRE COMMERCIALE \textbackslash{}& MODALITÉS\\
\vspace{0.5em}
\vspace{0.5em}
K.1 Le Pack RBK 2.0\\
\vspace{0.5em}
\vspace{0.5em}
\vspace{0.5em}
\vspace{0.5em}
EL\\
Ce qui est inclus pour chaque étudiant retenu.\\
\vspace{0.5em}
\vspace{0.5em}
\vspace{0.5em}
\vspace{0.5em}
TI\\
Table : Détail de l’Offre\\
\vspace{0.5em}
\vspace{0.5em}
\vspace{0.5em}
\vspace{0.5em}
N\\
DE\\
Service                                                                             Standard Inclus ?\\
Formation 48 semaines (Genesis Pool 4 + Explorer 12 + Bâtisseur 16 + Architecte 16) 1800h+ OUI\\
\vspace{0.5em}
\vspace{0.5em}
\vspace{0.5em}
\vspace{0.5em}
FI\\
Mentorat Expert (Review code hebdo)                                                 Senior    OUI\\
\vspace{0.5em}
\vspace{0.5em}
\vspace{0.5em}
N\\
Certification SBT                  CO                                               On-chain OUI\\
Accès Réseau Partenaires                                                            Superteam OUI\\
Hébergement (Piscine)                                                               Optionnel NON\\
\vspace{0.5em}
\vspace{0.5em}
\vspace{0.5em}
K.2 Pricing \textbackslash{}& Conditions (Value Ladder)\\
Notre tarification est conçue pour réduire le risque à l’entrée via un système progressif.\\
\vspace{0.5em}
\vspace{0.5em}
\vspace{0.5em}
\vspace{0.5em}
252\\
RBK 2.0 — Whitepaper Architecte Web3                                                         Annexe K. OFFRE COMMERCIALE \textbackslash{}& MODALITÉS\\
\vspace{0.5em}
\vspace{0.5em}
\begin{center}\begin{tabular}TAB. K.2 : Structure tarifaire alignée sur le pipeline Genesis → Architecte\end{tabular}\end{center}
\vspace{0.5em}
\vspace{0.5em}
Palier                                 Frais pédagogiques                 Modalités                                               Inclus\\
Genesis Pool (4 sem.)                  25\textbackslash{}% de 15 900 TND (facturé en TND) Paiement unique avant immersion                         Bootcamp\\
inten-\\
sif, Skill\\
Mirror\\
“Genesis\\
Access”,\\
support\\
mentors\\
6j/7\\
\vspace{0.5em}
\vspace{0.5em}
\vspace{0.5em}
\vspace{0.5em}
EL\\
Explorer (12 sem.)                  50\textbackslash{}% de 15 900 TND                   6 mensualités synchronisées sur les sprint reviews       Cours\\
\vspace{0.5em}
\vspace{0.5em}
\vspace{0.5em}
\vspace{0.5em}
TI\\
fonda-\\
\vspace{0.5em}
\vspace{0.5em}
\vspace{0.5em}
\vspace{0.5em}
N\\
mentaux\\
\vspace{0.5em}
\vspace{0.5em}
\vspace{0.5em}
\vspace{0.5em}
DE\\
Rust/-\\
Web3,\\
\vspace{0.5em}
\vspace{0.5em}
\vspace{0.5em}
\vspace{0.5em}
FI\\
Block\\
\vspace{0.5em}
\vspace{0.5em}
\vspace{0.5em}
N\\
CO                                                                               Check\\
“Explorer”,\\
accompa-\\
gnement\\
carrière\\
Bâtisseur \textbackslash{}& Architecte (16+16 sem.) Solde 25\textbackslash{}% de 15 900 TND             8 mensualités pendant les projets Dev DAO \textbackslash{}& Launch Proof Tracks\\
spéciali-\\
sés, Block\\
Check\\
“Build” \textbackslash{}&\\
“Launch\\
Proof”,\\
jury men-\\
tors      \textbackslash{}&\\
\begin{confidential}CONFIDENTIEL — PROJET RBK 2.0                                 Page 253 sur 316                                                    alumni\end{confidential}
Launchpad Nexus (option)            Forfait 3 300 TND + partage 30\textbackslash{}%     Forfait à l’entrée + success fee sur premiers revenus    Incubation\\
RBK 2.0 — Whitepaper Architecte Web3                                                                         Annexe K. OFFRE COMMERCIALE \textbackslash{}& MODALITÉS\\
\vspace{0.5em}
\vspace{0.5em}
K.2.1 Mécanisme d’Incitation (Upgrade)\\
• Crédit Genesis : Les frais Genesis Pool sont imputés à 100\textbackslash{}% sur la suite du parcours en cas de passage vers Explorer.\\
\vspace{0.5em}
• Fenêtre 30 jours : La décision d’upgrade doit intervenir sous 30 jours après le Block Check “Genesis” pour garantir l’étalement sur 15 900 TND.\\
\vspace{0.5em}
\vspace{0.5em}
K.2.2 Admission Directe (Passerelles)\\
Accès direct Explorer : Possible via test technique Rust/Algo/Git + soutenance Skill Mirror. Frais administratifs : 200 TND (déduits en cas d’inscription). Accès\\
\vspace{0.5em}
\vspace{0.5em}
\vspace{0.5em}
\vspace{0.5em}
EL\\
direct Bâtisseur/Architecte : Réservé aux profils expérimentés (portfolio Web3 + audit). Test + entretien. Frais : 300 TND.\\
\vspace{0.5em}
\vspace{0.5em}
\vspace{0.5em}
\vspace{0.5em}
TI\\
K.2.3 Offre ISA (Income Share Agreement)\\
Périmètre : Réservé aux ”Top Talents” qui s’engagent sur le parcours complet Genesis → Architecte ou sur le module Architecte seul. Sélection stricte (Dossier\\
\vspace{0.5em}
\vspace{0.5em}
\vspace{0.5em}
\vspace{0.5em}
EN\\
+ Technique + Social) pour l’accès ISA. Conditions Unifiées :\\
\vspace{0.5em}
• Partage : 15\textbackslash{}% du revenu brut mensuel.\\
\vspace{0.5em}
\vspace{0.5em}
\vspace{0.5em}
\vspace{0.5em}
D\\
• Déclencheur (Threshold) : Salaire \textgreater{} 2 500 TND Brut.\\
\vspace{0.5em}
\vspace{0.5em}
\vspace{0.5em}
\vspace{0.5em}
FI\\
• Durée : 36 mensualités (paiements effectifs) maximum.\\
\vspace{0.5em}
• Plafond (Cap) : 20 000 TND total remboursé (Risk Premium inclus).\\
N\\
• Clause de pause : Automatique en cas de chômage ou revenu \textless{} Seuil.\\
O\\
• Juridiction : Droit Tunisien, contrat enregistré.\\
C\\
\vspace{0.5em}
\vspace{0.5em}
Parcours Complet (15 900 TND + 3 300 TND)\\
\vspace{0.5em}
\vspace{0.5em}
Genesis Pool  Conv. 60\textbackslash{}%     Explorer     Conv. 80\textbackslash{}% Bâtisseur/Architecte                                        Launchpad\\
(25\textbackslash{}% 15 900 TND)           (50\textbackslash{}% 15 900 TND)            (25\textbackslash{}% 15 900 TND)                                           + ISA 15\textbackslash{}%\\
\vspace{0.5em}
FIG. K.1 : Value Ladder et Parcours Étudiant\\
\vspace{0.5em}
\vspace{0.5em}
\begin{confidential}CONFIDENTIEL — PROJET RBK 2.0                                           Page 254 sur 316\end{confidential}
RBK 2.0 — Whitepaper Architecte Web3                                                                           Annexe K. OFFRE COMMERCIALE \textbackslash{}& MODALITÉS\\
\vspace{0.5em}
\vspace{0.5em}
\vspace{0.5em}
K.3 Objections \textbackslash{}& Réponses\\
”C’est trop cher ?” C’est le prix d’une voiture d’occasion pour une carrière internationale. L’option ”Niveau 1” vous permet de tester pour un coût réduit.\\
”Pourquoi pas une fac publique ?” La fac offre un diplôme académique. Nous offrons une certification technique industrielle et un accès direct au réseau\\
Superteam.\\
\vspace{0.5em}
\vspace{0.5em}
\vspace{0.5em}
K.4 Politique de Remboursement et Report\\
• Satisfait ou Remboursé (Genesis Pool) : Remboursement intégral possible jusqu’à la fin de la 1ère semaine d’immersion Genesis.\\
\vspace{0.5em}
• Report de Cohorte : Possible une seule fois en cas de force majeure, sans frais, sous réserve de places disponibles.\\
\vspace{0.5em}
• Non-Garanti : RBK s’engage sur la qualité de la formation (”Obligation de Moyens”) mais ne peut garantir contractuellement une embauche ou un niveau de\\
\vspace{0.5em}
\vspace{0.5em}
\vspace{0.5em}
\vspace{0.5em}
EL\\
salaire spécifique (”Obligation de Résultats”), ceux-ci dépendant du marché et de l’effort individuel.\\
\vspace{0.5em}
\vspace{0.5em}
\vspace{0.5em}
\vspace{0.5em}
TI\\
N\\
DE\\
FI\\
N\\
CO\\
\vspace{0.5em}
\vspace{0.5em}
\vspace{0.5em}
\vspace{0.5em}
\begin{confidential}CONFIDENTIEL — PROJET RBK 2.0                                            Page 255 sur 316\end{confidential}
L \textbar{} STRATÉGIE MENTORAT \textbackslash{}& TRAIN-THE-TRAINER\\
\vspace{0.5em}
\vspace{0.5em}
La qualité de RBK 2.0 repose sur la qualité de son encadrement humain. Nous ne recrutons pas des ”profs”, mais des ”Tech Leads” capables de guider des juniors.\\
\vspace{0.5em}
\vspace{0.5em}
\vspace{0.5em}
\vspace{0.5em}
EL\\
L.1 Programme ”Mentor-in-a-Box”\\
\vspace{0.5em}
\vspace{0.5em}
\vspace{0.5em}
\vspace{0.5em}
TI\\
N\\
• Captation systématique : chaque masterclass, revue de code et Block Check est enregistrée (audio + vidéo + notes) puis indexée dans Venture Engine. Les\\
\vspace{0.5em}
\vspace{0.5em}
\vspace{0.5em}
\vspace{0.5em}
DE\\
artefacts sont tagués par module, niveau et difficultés rencontrées.\\
\vspace{0.5em}
• Knowledge base exploitable : les meilleurs segments sont transformés en micro-capsules (10 minutes) et en fiches ”anti-patterns” accessibles aux mentors et\\
\vspace{0.5em}
\vspace{0.5em}
\vspace{0.5em}
\vspace{0.5em}
FI\\
aux apprenants avancés.\\
\vspace{0.5em}
\vspace{0.5em}
N\\
CO\\
• Playbooks de remédiation : pour toute alerte burnout ou difficulté technique récurrente, un ”pack” Mentor-in-a-Box est déployé (vidéo explicative, exercices\\
guidés, scripts de coaching).\\
\vspace{0.5em}
• Contrôle qualité : le Responsable Bien-être \textbackslash{}& Résilience vérifie chaque trimestre que les sessions critiques ont bien été captées et auditées.\\
\vspace{0.5em}
\vspace{0.5em}
\vspace{0.5em}
L.2 Le Pipeline ”Train the Trainer”\\
Pour assurer la scalabilité sans perte de qualité, RBK forme ses propres mentors parmi les meilleurs Alumni.\\
\vspace{0.5em}
1. Sourcing : Top 10\textbackslash{}% des diplômés (Score Tech \textgreater{} 90/100 + Soft Skills A).\\
\vspace{0.5em}
2. Shadowing (1 cohorte) : L’aspirant-mentor suit un mentor Senior pendant 3 mois, produit des fiches ”Takeaways” après chaque sprint et contribue au dépôt\\
Mentor-in-a-Box.\\
\vspace{0.5em}
\vspace{0.5em}
256\\
RBK 2.0 — Whitepaper Architecte Web3                                                                                               Annexe L. Stratégie mentorat\\
\vspace{0.5em}
\vspace{0.5em}
3. Certification pédagogique : Formation interne de 2 semaines sur :\\
\vspace{0.5em}
• La méthode Socratique (répondre par une question).\\
\vspace{0.5em}
• La gestion de crise émotionnelle (Protocole Anti-Burnout).\\
\vspace{0.5em}
• La détection de triche par IA.\\
\vspace{0.5em}
4. Titularisation : Prise en charge d’une Squad de 15 étudiants, obligation contractuelle de produire un pack Mentor-in-a-Box par module critique.\\
\vspace{0.5em}
\vspace{0.5em}
\vspace{0.5em}
L.3 Modèle de Rémunération Incitatif\\
Nous alignons les intérêts des mentors sur la réussite des étudiants.\\
\vspace{0.5em}
\vspace{0.5em}
\vspace{0.5em}
\vspace{0.5em}
EL\\
\begin{center}\begin{tabular}TAB. L.1 : Grille de Rémunération Mentor (Junior → Lead)\end{tabular}\end{center}
\vspace{0.5em}
\vspace{0.5em}
\vspace{0.5em}
\vspace{0.5em}
TI\\
Niveau          Fixe (Mensuel) Variable (Performance)\\
\vspace{0.5em}
\vspace{0.5em}
\vspace{0.5em}
\vspace{0.5em}
N\\
Junior Mentor 2 500 TND        100 TND par étudiant validant la Genesis Pool.\\
\vspace{0.5em}
\vspace{0.5em}
\vspace{0.5em}
\vspace{0.5em}
DE\\
Senior Mentor 4 500 TND        2\textbackslash{}% du Pool ISA de sa cohorte (si placement \textgreater{} 90\textbackslash{}%).\\
Lead Instructor 7 000 TND      Part de l’EBITDA annuel (BSPCE/Tokens).\\
\vspace{0.5em}
\vspace{0.5em}
\vspace{0.5em}
\vspace{0.5em}
FI\\
N\\
CO\\
L.4 Contrat type Mentor Nexus\\
Clauses essentielles du contrat Mentor\\
\vspace{0.5em}
\vspace{0.5em}
\vspace{0.5em}
Clause                     Contenu contractuel\\
Confidentialité \textbackslash{}&          NDA renforcé, interdiction de réutiliser les artefacts pédagogiques hors Nexus/RBK, authentification MFA obligatoire.\\
sécurité\\
Propriété intellectuelle   Cession non exclusive à RBK \textbackslash{}& Nexus des supports créés, sans droit de publication publique sans accord écrit.\\
Non-sollicitation          Interdiction de recruter les apprenants ou mentors RBK pendant 12 mois post-contrat.\\
\vspace{0.5em}
\vspace{0.5em}
\vspace{0.5em}
\vspace{0.5em}
\begin{confidential}CONFIDENTIEL — PROJET RBK 2.0                                              Page 257 sur 316\end{confidential}
RBK 2.0 — Whitepaper Architecte Web3                                                                                              Annexe L. Stratégie mentorat\\
\vspace{0.5em}
\vspace{0.5em}
\vspace{0.5em}
Qualité \textbackslash{}& reporting       Mise à jour hebdomadaire de la base Mentor-in-a-Box, participation aux comités qualité et respect des process Block Check.\\
Bien-être \textbackslash{}& conformité    Acceptation du protocole Anti-Burnout : obligation de remonter toute alerte et d’appliquer les décisions du Responsable Bien-être.\\
\vspace{0.5em}
\vspace{0.5em}
\vspace{0.5em}
\vspace{0.5em}
L.5 Plan de Relève et Continuité\\
Pour éviter le ”Bus Factor” (départ d’un instructeur clé) :\\
\vspace{0.5em}
• Binômes Rotatifs : Chaque module critique (ex : Rust Advanced) est maîtrisé par au moins 2 mentors Seniors.\\
\vspace{0.5em}
• Documentation ”Playbook” : Chaque cours dispose d’un guide ”Teacher’s Notes” détaillant les points de friction habituels et les métaphores clés.\\
\vspace{0.5em}
• Guest Lecturers : Bassin de 5 experts externes (CTO partenaires) activables pour des masterclasses ponctuelles ou des remplacements d’urgence.\\
\vspace{0.5em}
\vspace{0.5em}
\vspace{0.5em}
\vspace{0.5em}
EL\\
TI\\
N\\
DE\\
FI\\
N\\
CO\\
\vspace{0.5em}
\vspace{0.5em}
\vspace{0.5em}
\vspace{0.5em}
\begin{confidential}CONFIDENTIEL — PROJET RBK 2.0                                          Page 258 sur 316\end{confidential}
M \textbar{} Offre partenariat B2B\\
\vspace{0.5em}
\vspace{0.5em}
M.1 Modèle d’Offre Corporate\\
\vspace{0.5em}
\vspace{0.5em}
\vspace{0.5em}
\vspace{0.5em}
EL\\
Ce document sert de base aux négociations avec les entreprises partenaires (ESN, Banques, Startups) souhaitant upskiller leurs équipes.\\
\vspace{0.5em}
\vspace{0.5em}
\vspace{0.5em}
\vspace{0.5em}
TI\\
N\\
M.1.1 Les Packs Entreprise\\
\vspace{0.5em}
\vspace{0.5em}
\vspace{0.5em}
\vspace{0.5em}
DE\\
FI\\
\begin{center}\begin{tabular}TAB. M.1 : Tarification B2B\end{tabular}\end{center}
\vspace{0.5em}
\vspace{0.5em}
\vspace{0.5em}
N\\
CO\\
Pack      Volume         Tarif Unitaire\\
Starter 1 à 2 sièges 18 000 TND\\
Squad     3 à 5 sièges 16 380 TND (-9\textbackslash{}%)\\
Factory 6+ sièges        15 300 TND (-15\textbackslash{}%)\\
\vspace{0.5em}
\vspace{0.5em}
\vspace{0.5em}
\vspace{0.5em}
M.2 Conditions Particulières\\
1. Engagement de Résultat : Obligation de moyens (formation). Aucun remboursement en cas d’échec aux examens.\\
\vspace{0.5em}
\vspace{0.5em}
\vspace{0.5em}
259\\
RBK 2.0 — Whitepaper Architecte Web3                                                                                              Annexe M. Offre partenariat B2B\\
\vspace{0.5em}
\vspace{0.5em}
2. Propriété Intellectuelle : Les projets réalisés par les collaborateurs sont la propriété exclusive de l’entreprise (Work for Hire).\\
\vspace{0.5em}
3. Confidentialité : NDA1 signé pour les Use-Cases métier.\\
\vspace{0.5em}
\vspace{0.5em}
\vspace{0.5em}
\vspace{0.5em}
EL\\
TI\\
N\\
DE\\
FI\\
N\\
CO\\
\vspace{0.5em}
\vspace{0.5em}
\vspace{0.5em}
\vspace{0.5em}
1\\
Finance/Juridique/Compliance — Non-Disclosure Agreement (NDA) : accord encadrant la confidentialité et l’usage des informations échangées entre parties.\\
\vspace{0.5em}
\vspace{0.5em}
\begin{confidential}CONFIDENTIEL — PROJET RBK 2.0                                            Page 260 sur 316\end{confidential}
N \textbar{} Outillage \textbackslash{}& stack technique\\
\vspace{0.5em}
\vspace{0.5em}
N.1 Stack de Développement (Cyborg-Ready)\\
\vspace{0.5em}
\vspace{0.5em}
\vspace{0.5em}
\vspace{0.5em}
EL\\
Nous listons ici la toolchain de base à installer pour reproduire l’environnement ”Cyborg” utilisé dans le cursus.\\
\vspace{0.5em}
\vspace{0.5em}
\vspace{0.5em}
\vspace{0.5em}
TI\\
N\\
N.1.1 Environnement Local\\
\vspace{0.5em}
\vspace{0.5em}
\vspace{0.5em}
\vspace{0.5em}
DE\\
• OS : Linux (Ubuntu/Pop !\textbackslash{}_OS) ou macOS. Windows via WSL2 uniquement.\\
\vspace{0.5em}
\vspace{0.5em}
\vspace{0.5em}
\vspace{0.5em}
FI\\
• IDE : VSCode (Profile RBK : Rust, Solidity, GitLens, Copilot).\\
\vspace{0.5em}
\vspace{0.5em}
\vspace{0.5em}
N\\
• Terminal : Alacritty + Tmux + Starship.\\
CO\\
N.1.2 Chain Stack\\
• Solana : Rust 1.75+, Anchor 0.29+, Solana CLI 1.18+.\\
\vspace{0.5em}
• EVM : Foundry (Forge, Cast, Anvil).\\
\vspace{0.5em}
• Indexing : The Graph (EVM), Helius/Shyft (Solana).\\
\vspace{0.5em}
\vspace{0.5em}
\vspace{0.5em}
N.2 Outils de Productivité \textbackslash{}& IA\\
\vspace{0.5em}
\vspace{0.5em}
261\\
RBK 2.0 — Whitepaper Architecte Web3                                                                   Annexe N. Outillage \textbackslash{}& stack technique\\
\vspace{0.5em}
\vspace{0.5em}
\vspace{0.5em}
\vspace{0.5em}
\begin{center}\begin{tabular}TAB. N.1 : Matrice des Outils IA Autorisés\end{tabular}\end{center}
\vspace{0.5em}
\vspace{0.5em}
Usage            Outil Validé       Politique d’Usage\\
Coding Assistant GitHub Copilot     Autorisé pour boilerplate/tests. Interdit pour algo critique.\\
Docs Search      Perplexity / Phind Recommandé pour la recherche contextuelle.\\
Review           CodeRabbit         Pré-analyse des PRs avant review humaine.\\
Diagrams         Mermaid.js         ”Diagrams as Code” obligatoire.\\
\vspace{0.5em}
\vspace{0.5em}
\vspace{0.5em}
\vspace{0.5em}
EL\\
N.3 Infrastructure CI/CD (Github Actions)\\
\vspace{0.5em}
\vspace{0.5em}
\vspace{0.5em}
\vspace{0.5em}
TI\\
Tout repo étudiant doit inclure un workflow .github/workflows/ci.yml standardisé :\\
\vspace{0.5em}
\vspace{0.5em}
\vspace{0.5em}
\vspace{0.5em}
N\\
DE\\
• Lint : cargo clippy / solhint\\
\vspace{0.5em}
\vspace{0.5em}
\vspace{0.5em}
\vspace{0.5em}
FI\\
• Test : cargo test / forge test\\
\vspace{0.5em}
\vspace{0.5em}
\vspace{0.5em}
\vspace{0.5em}
N\\
• Audit : cargo audit                          CO\\
• Format : cargo fmt --check\\
\vspace{0.5em}
\vspace{0.5em}
\vspace{0.5em}
\vspace{0.5em}
\begin{confidential}CONFIDENTIEL — PROJET RBK 2.0                                      Page 262 sur 316\end{confidential}
O \textbar{} RÉFÉRENTIEL DE COMPÉTENCES\\
\vspace{0.5em}
\vspace{0.5em}
O.1 Matrice de Compétences\\
\vspace{0.5em}
\vspace{0.5em}
\vspace{0.5em}
\vspace{0.5em}
EL\\
TI\\
\begin{center}\begin{tabular}TAB. O.1 : Mapping Détaillé Compétences / Badges / Preuves\end{tabular}\end{center}
\vspace{0.5em}
\vspace{0.5em}
\vspace{0.5em}
\vspace{0.5em}
N\\
DE\\
Module         Compétence Clé                  Badge SBT        Preuve (Livrable)      Critère de Validation\\
\vspace{0.5em}
\vspace{0.5em}
\vspace{0.5em}
\vspace{0.5em}
FI\\
Fundamentals   Algorithmique                   Algo Ninja       Profil Codewars        Rank 5kyu min.\\
\vspace{0.5em}
\vspace{0.5em}
\vspace{0.5em}
N\\
Rust Systems   Memory Safety (Borrow Checker) Rust Ace          Repo ”Rustlings”       100\textbackslash{}% Exos validés + 0 Warning clippy.\\
CO\\
Smart Contract Développement Logique           Protocol Architect Repo Capstone 2      Tests unitaires \textgreater{} 90\textbackslash{}% coverage.\\
\vspace{0.5em}
Frontend       Intégration Wallet              dApp Builder     Repo Capstone 1        Connexion \textless{} 2s, Feedback UI fluide.\\
\vspace{0.5em}
Sécurité       Audit                           Whitehat Jr      Rapport Audit Capstone Threat Model + 1 vulnérabilité (mock) trouvée.\\
\vspace{0.5em}
Delivery       CI/CD                           DevOps Ready     GitHub Actions Logs    Build Vert + Deploy auto sur Devnet.\\
\vspace{0.5em}
Soft Skills    Communication                   Squad Lead       Vidéo Demo Day         Pitch 3min clair, Slides pro.\\
\vspace{0.5em}
\vspace{0.5em}
\vspace{0.5em}
\vspace{0.5em}
263\\
P \textbar{} Charte de qualité \textbackslash{}& règles d’or\\
\vspace{0.5em}
\vspace{0.5em}
\vspace{0.5em}
\vspace{0.5em}
EL\\
RBK 2.0 repose sur un socle de valeurs non négociables. Tout manquement à ces règles entraîne une exclusion immédiate.\\
\vspace{0.5em}
\vspace{0.5em}
\vspace{0.5em}
\vspace{0.5em}
TI\\
P.1 Les 4 Commandements de l’Ingénieur RBK\\
\vspace{0.5em}
\vspace{0.5em}
\vspace{0.5em}
\vspace{0.5em}
EN\\
1. No Broken Windows : Aucun code n’est mergé sur ’main’ s’il contient des warnings de linter ou des TODOs non résolus.\\
\vspace{0.5em}
2. Don’t Trust, Verify : Chaque ligne de code générée par IA doit être auditée.\\
\vspace{0.5em}
\vspace{0.5em}
\vspace{0.5em}
\vspace{0.5em}
D\\
3. Ships or Nothing : Un projet non déployé n’existe pas.\\
\vspace{0.5em}
\vspace{0.5em}
\vspace{0.5em}
FI\\
4. Leave No One Behind : Le savoir ne vaut que s’il est partagé.\\
N\\
P.2 Matrice de Conformité (Sanctions)\\
O\\
C\\
\vspace{0.5em}
\vspace{0.5em}
Infraction          Exemple                        Sanction\\
Plagiat             Copie repo externe sans crédit Exclusion\\
Négligence Sécurité Commit de Private Key          Blâme + Reset Projet\\
Ghosting            Absence non justifiée \textgreater{} 48h Avertissement\\
\vspace{0.5em}
\vspace{0.5em}
\vspace{0.5em}
\vspace{0.5em}
264\\
RBK 2.0 — Whitepaper Architecte Web3                                               Annexe P. Charte de qualité \textbackslash{}& règles d’or\\
\vspace{0.5em}
\vspace{0.5em}
\vspace{0.5em}
P.3 Processus de Validation Qualité\\
Pipeline : Code Complete → Linter → Tests → Audit IA \textbackslash{}& Humain → Merge.\\
\vspace{0.5em}
\vspace{0.5em}
\vspace{0.5em}
\vspace{0.5em}
EL\\
TI\\
N\\
DE\\
FI\\
N\\
CO\\
\vspace{0.5em}
\vspace{0.5em}
\vspace{0.5em}
\vspace{0.5em}
\begin{confidential}CONFIDENTIEL — PROJET RBK 2.0                                   Page 265 sur 316\end{confidential}
Q \textbar{} CONTRAT DE PARTENARIAT PÉDAGOGIQUE ET DE PAR-\\
TAGE DE SUCCÈS (CPPS)\\
\vspace{0.5em}
\vspace{0.5em}
\vspace{0.5em}
\vspace{0.5em}
EL\\
TI\\
Note : Ce document est un squelette contractuel. La version définitive doit être validée par un cabinet d’avocats tunisien spécialisé en droit des affaires et fiscalité internationale.\\
\vspace{0.5em}
\vspace{0.5em}
\vspace{0.5em}
\vspace{0.5em}
EN\\
Q.1 Objet du Contrat\\
\vspace{0.5em}
\vspace{0.5em}
\vspace{0.5em}
\vspace{0.5em}
D\\
Le présent CPPS est conclu entre ReBoot Kamp Tunisie (RBK) et l’apprenant ou alumni sélectionné. RBK finance et délivre, via Nexus Réussite comme maître\\
d’œuvre pédagogique, le parcours “Architecte Web3”. En contrepartie, l’apprenant s’engage à verser à RBK une fraction de ses revenus professionnels uniquement\\
\vspace{0.5em}
\vspace{0.5em}
FI\\
lorsque les conditions de succès définies ci-après sont réunies.\\
N\\
Q.2 Périmètre des revenus éligibles\\
O\\
\vspace{0.5em}
• Domaines concernés : emplois salariés, missions freelances ou mandats de conseil dans les métiers listés dans l’Annexe 14 (Architecte Web3, Smart-Contract\\
C\\
\vspace{0.5em}
\vspace{0.5em}
Engineer, DevOps Web3, Product Strategist Web3, etc.).\\
\vspace{0.5em}
• Conditionnalité stricte : aucun paiement n’est dû si l’activité réalisée ne relève pas du périmètre Web3/blockchain/Rust/EVM décrit dans le programme et\\
n’atteint pas le seuil de revenu défini.\\
\vspace{0.5em}
• Justificatifs admissibles : contrats de travail, avenants salariaux, contrats de mission, factures clients ou attestations d’honoraires émis par un tiers vérifiable.\\
\vspace{0.5em}
\vspace{0.5em}
\vspace{0.5em}
\vspace{0.5em}
266\\
RBK 2.0 — Whitepaper Architecte Web3                        Annexe Q. CONTRAT DE PARTENARIAT PÉDAGOGIQUE ET DE PARTAGE DE SUCCÈS (CPPS)\\
\vspace{0.5em}
\vspace{0.5em}
\vspace{0.5em}
Q.3 Définitions clés\\
• Seuil de succès : 2500 brut mensuel par rapport à la devise de référence contractuelle.\\
\vspace{0.5em}
• Taux de partage : 15\textbackslash{}% du revenu brut éligible lorsque le seuil est dépassé.\\
\vspace{0.5em}
• Plafond : 20000 cumulés ou 36 mensualités payées, la première condition atteinte mettant fin aux obligations.\\
\vspace{0.5em}
• Pause automatique : suspension du versement dès que le revenu éligible repasse sous le seuil ; aucune capitalisation d’arriérés.\\
\vspace{0.5em}
\vspace{0.5em}
\vspace{0.5em}
\vspace{0.5em}
EL\\
Q.4 Obligations de RBK\\
• Délivrer le curriculum complet, les évaluations “Block Check” et l’accompagnement carrière définis dans le présent whitepaper.\\
\vspace{0.5em}
\vspace{0.5em}
\vspace{0.5em}
\vspace{0.5em}
TI\\
• Mettre à disposition du bénéficiaire un accès transparent aux relevés d’heures et de preuves d’apprentissage fournis par Nexus Réussite via Venture Engine.\\
\vspace{0.5em}
\vspace{0.5em}
\vspace{0.5em}
\vspace{0.5em}
EN\\
• Alimenter et gérer le Fonds de Garantie ISA décrit au Chapitre 15, garantissant la prise en charge des défauts éventuels.\\
\vspace{0.5em}
\vspace{0.5em}
\vspace{0.5em}
\vspace{0.5em}
D\\
Q.5 Obligations du bénéficiaire\\
\vspace{0.5em}
\vspace{0.5em}
FI\\
• Maintenir une assiduité supérieure à 95\textbackslash{}% et respecter le règlement intérieur pédagogique.\\
\vspace{0.5em}
• Déclarer à RBK tout emploi ou mission éligible dans un délai de 10 jours ouvrés, fournir les justificatifs de rémunération et signer les mandats de prélèvement\\
N\\
nécessaires.\\
O\\
• Autoriser, si nécessaire, la vérification des informations auprès de l’employeur ou du client final dans le respect du RGPD (cf. Chapitre 18).\\
C\\
\vspace{0.5em}
\vspace{0.5em}
Q.6 Collecte et validation des revenus\\
1. Transmission mensuelle par le bénéficiaire d’un bulletin de salaire, d’un contrat freelance ou d’une facture.\\
\vspace{0.5em}
2. Validation par le Comité Carrière RBK \textbackslash{}& Nexus qui cote le poste et confirme son éligibilité au périmètre Web3.\\
\vspace{0.5em}
3. En cas de doute, saisine automatique du Comité Éthique \textbackslash{}& Pédagogique décrit au Chapitre 5. Aucun prélèvement n’est effectué avant sa décision contradictoire.\\
\vspace{0.5em}
\vspace{0.5em}
\vspace{0.5em}
\vspace{0.5em}
\begin{confidential}CONFIDENTIEL — PROJET RBK 2.0                                            Page 267 sur 316\end{confidential}
RBK 2.0 — Whitepaper Architecte Web3                          Annexe Q. CONTRAT DE PARTENARIAT PÉDAGOGIQUE ET DE PARTAGE DE SUCCÈS (CPPS)\\
\vspace{0.5em}
\vspace{0.5em}
\vspace{0.5em}
Q.7 Modalités de paiement\\
• Les versements s’effectuent mensuellement via prélèvement SEPA/TND ou virement manuel selon le pays d’exercice.\\
\vspace{0.5em}
• RBK émet un relevé détaillant le montant dû, la base de calcul, le solde restant dû et la durée restante.\\
\vspace{0.5em}
• Une remise “Bonus de paiement anticipé” de 20\textbackslash{}% du solde est appliquée si l’apprenant règle l’intégralité des montants dus dans les 6 premiers mois qui\\
suivent l’accès au revenu éligible.\\
\vspace{0.5em}
\vspace{0.5em}
\vspace{0.5em}
Q.8 Gestion des litiges et cas de force majeure\\
• Toute contestation (période de chômage, maladie, changement de secteur) est portée devant le Comité Éthique \textbackslash{}& Pédagogique qui peut accorder des pauses\\
\vspace{0.5em}
\vspace{0.5em}
\vspace{0.5em}
\vspace{0.5em}
EL\\
supplémentaires, des réductions ou une annulation partielle.\\
\vspace{0.5em}
\vspace{0.5em}
\vspace{0.5em}
\vspace{0.5em}
TI\\
• Les décisions du Comité sont consignées dans Venture Engine et communiquées aux parties sous 5 jours ouvrés.\\
\vspace{0.5em}
\vspace{0.5em}
\vspace{0.5em}
\vspace{0.5em}
N\\
DE\\
Q.9 Recouvrement et protection du bénéficiaire\\
\vspace{0.5em}
\vspace{0.5em}
\vspace{0.5em}
\vspace{0.5em}
FI\\
• RBK demeure l’unique créancier du CPPS et mandate, si nécessaire, un prestataire de recouvrement certifié (ex. TCM Tunisia ou Intrum). Aucun transfert de\\
\vspace{0.5em}
\vspace{0.5em}
\vspace{0.5em}
N\\
créance n’est réalisé sans notification écrite préalable.\\
CO\\
• En cas de défaut persistant (\textgreater{}90 jours), RBK peut activer le Fonds de Garantie ISA ; le bénéficiaire reste redevable dans la limite du plafond contractuel.\\
\vspace{0.5em}
• Aucune mention publique n’est effectuée sans revue du Comité Éthique ; l’outil SBT “Launch Proof” peut être marqué “en anomalie” sans effacement de la\\
preuve de compétences.\\
\vspace{0.5em}
\vspace{0.5em}
\vspace{0.5em}
\vspace{0.5em}
\begin{confidential}CONFIDENTIEL — PROJET RBK 2.0                                            Page 268 sur 316\end{confidential}
R \textbar{} MODÈLE DE PARTENARIAT B2B (HIRING)\\
\vspace{0.5em}
\vspace{0.5em}
Cadre de collaboration pour les entreprises partenaires souhaitant recruter les talents RBK.\\
\vspace{0.5em}
\vspace{0.5em}
\vspace{0.5em}
\vspace{0.5em}
EL\\
R.1    Offre ”Hire Train Deploy”\\
\vspace{0.5em}
\vspace{0.5em}
\vspace{0.5em}
\vspace{0.5em}
TI\\
N\\
L’entreprise partenaire (le ”Client”) mandate RBK pour former sur mesure une escouade (Squad) de 3 à 5 talents sur une stack technologique spécifique.\\
\vspace{0.5em}
\vspace{0.5em}
\vspace{0.5em}
\vspace{0.5em}
DE\\
• Booking Fee : 5 000 TND HT / Talent (réservation).\\
\vspace{0.5em}
\vspace{0.5em}
\vspace{0.5em}
\vspace{0.5em}
FI\\
• Success Fee : 10\textbackslash{}% du salaire brut annuel à l’embauche.\\
\vspace{0.5em}
\vspace{0.5em}
\vspace{0.5em}
N\\
CO\\
R.2    Offre ”Corporate Upskilling”\\
Formation intensive pour les équipes tech existantes.\\
\vspace{0.5em}
• Pack Team : 25 000 TND HT (jusqu’à 10 devs).\\
\vspace{0.5em}
• Inclus : Accès LMS à vie, Certification SBT.\\
\vspace{0.5em}
\vspace{0.5em}
\vspace{0.5em}
\vspace{0.5em}
269\\
S \textbar{} Kit de survie juridique\\
\vspace{0.5em}
\vspace{0.5em}
Ce kit fournit les templates et checklists pratiques pour que l’étudiant puisse opérer professionnellement.\\
\vspace{0.5em}
\vspace{0.5em}
\vspace{0.5em}
\vspace{0.5em}
EL\\
S.1 Modèle de Contrat de Prestation Freelance (Extraits)\\
\vspace{0.5em}
\vspace{0.5em}
\vspace{0.5em}
\vspace{0.5em}
TI\\
N\\
• Objet : Description précise du Scope of Work (ex : ”Développement Vault ERC-4626”).\\
\vspace{0.5em}
\vspace{0.5em}
\vspace{0.5em}
\vspace{0.5em}
DE\\
• Paiement : Jalons clairs (30\textbackslash{}% signature, 40\textbackslash{}% tests, 30\textbackslash{}% livraison).\\
\vspace{0.5em}
\vspace{0.5em}
\vspace{0.5em}
\vspace{0.5em}
FI\\
• IP : Cession de la propriété intellectuelle au client après paiement intégral.\\
\vspace{0.5em}
\vspace{0.5em}
\vspace{0.5em}
N\\
CO\\
S.2 Checklist : Créer sa Micro-Entreprise Exportatrice\\
` Choix du Nom (INNORPI).\\
\vspace{0.5em}
` Dépôt Dossier APII.\\
\vspace{0.5em}
` Immatriculation RNE.\\
\vspace{0.5em}
` Compte Bancaire Devises.\\
\vspace{0.5em}
` Cachet Officiel.\\
\vspace{0.5em}
\vspace{0.5em}
\vspace{0.5em}
\vspace{0.5em}
270\\
RBK 2.0 — Whitepaper Architecte Web3                                                    Annexe S. Kit de survie juridique\\
\vspace{0.5em}
\vspace{0.5em}
\vspace{0.5em}
S.3 Guide Visuel : Recevoir un Salaire en Crypto\\
Stratégie de Réception :\\
\vspace{0.5em}
1. Option Simple : Passerelle (Grey/Bitwage) → Conversion Auto → Virement TND.\\
\vspace{0.5em}
2. Option Expert : Wallet Pro (USDC) → Comptabilité Devises → Cession Manuelle BCT.\\
\vspace{0.5em}
\vspace{0.5em}
\vspace{0.5em}
S.4 Red Flags (Vigilance)\\
• Refus de contrat écrit.\\
\vspace{0.5em}
• Paiement en token volatile inconnu.\\
\vspace{0.5em}
\vspace{0.5em}
\vspace{0.5em}
\vspace{0.5em}
EL\\
• Demande de ”frais d’avance”.\\
\vspace{0.5em}
\vspace{0.5em}
\vspace{0.5em}
\vspace{0.5em}
TI\\
N\\
DE\\
FI\\
N\\
CO\\
\vspace{0.5em}
\vspace{0.5em}
\vspace{0.5em}
\vspace{0.5em}
\begin{confidential}CONFIDENTIEL — PROJET RBK 2.0                                    Page 271 sur 316\end{confidential}
T \textbar{} Annexe : Modèle Piscine \textbackslash{}& Staffing\\
\vspace{0.5em}
\vspace{0.5em}
Cette annexe détaille le modèle de staffing ”piscine”, les rôles pédagogiques, les volumes d’heures contractuels et les structures de rémunération conçus pour\\
préserver la marge du projet tout en assurant une qualité d’encadrement premium.\\
\vspace{0.5em}
\vspace{0.5em}
\vspace{0.5em}
\vspace{0.5em}
EL\\
TI\\
T.1 Synthèse du Modèle Piscine et Staffing\\
\vspace{0.5em}
\vspace{0.5em}
\vspace{0.5em}
\vspace{0.5em}
N\\
DE\\
Cette section condense le modèle de staffing ”piscine”, les rôles pédagogiques, les volumes d’heures contractuels et les structures de rémunération, conçus pour\\
préserver la marge du projet tout en assurant une qualité d’encadrement premium.\\
\vspace{0.5em}
\vspace{0.5em}
\vspace{0.5em}
\vspace{0.5em}
FI\\
N\\
Synthèse du Modèle Piscine et Staffing\\
CO\\
Critère          Program Director (Alaeddine / Nexus) Lead Mentor (Senior généraliste)                        Expert(s) Fractional\\
\vspace{0.5em}
Rôle             Gouvernance pédagogique, arbitrages, Présence régulière, déblocage, revue de Interventions “chirurgicales” (revues\\
Principal        qualité, conformité ISA, placement, rela- PR, discipline d’exécution.        d’architecture, sessions audit, jurys).\\
tions écosystème.\\
\vspace{0.5em}
Cap              –                                             Min(25\textbackslash{}% du CA encaissé, Plafond TND Min(25\textbackslash{}% du CA encaissé, Plafond TND\\
Mensuel*                                                       fixe)                               fixe)\\
\vspace{0.5em}
\vspace{0.5em}
\vspace{0.5em}
\vspace{0.5em}
272\\
RBK 2.0 — Whitepaper Architecte Web3                                           Annexe T. Annexe : Modèle Piscine \textbackslash{}& Staffing\\
\vspace{0.5em}
\vspace{0.5em}
\vspace{0.5em}
\vspace{0.5em}
Barème LEAN –                          10 h/semaine (2h/jour × 5)            4 h/mois (revues milestones)\\
(8–12\\
apprenants)\\
\vspace{0.5em}
Barème         –                       15 h/semaine                          8 h/mois\\
STANDARD\\
(13–20\\
apprenants)\\
\vspace{0.5em}
Barème         –                       2 × 12 h/semaine (rotation)           12 h/mois\\
INTENSIF\\
(21–30\\
\vspace{0.5em}
\vspace{0.5em}
\vspace{0.5em}
\vspace{0.5em}
EL\\
apprenants)\\
\vspace{0.5em}
\vspace{0.5em}
\vspace{0.5em}
\vspace{0.5em}
TI\\
Rémunération –                         120–200 TND/h (ou forfait mensuel)    250–450 TND/h (ou forfait “jury + au-\\
\vspace{0.5em}
\vspace{0.5em}
\vspace{0.5em}
\vspace{0.5em}
N\\
indicative                                                                   dits”)\\
\vspace{0.5em}
\vspace{0.5em}
\vspace{0.5em}
\vspace{0.5em}
DE\\
Mode de        Forfait mensuel fixe    Forfait mensuel fixe                  Pack ”8h/mois + 1 jury” (prix fixe)\\
\vspace{0.5em}
\vspace{0.5em}
\vspace{0.5em}
\vspace{0.5em}
FI\\
Paiement\\
\vspace{0.5em}
\vspace{0.5em}
\vspace{0.5em}
N\\
CO\\
\vspace{0.5em}
\vspace{0.5em}
\vspace{0.5em}
\vspace{0.5em}
\begin{confidential}CONFIDENTIEL — PROJET RBK 2.0                 Page 273 sur 316\end{confidential}
RBK 2.0 — Whitepaper Architecte Web3                                                                    Annexe T. Annexe : Modèle Piscine \textbackslash{}& Staffing\\
\vspace{0.5em}
\vspace{0.5em}
\vspace{0.5em}
\vspace{0.5em}
Critère          TAs / Alumni Assistants\\
\vspace{0.5em}
Rôle             Support opérationnel, office-hours, suivi individuel, logistique d’évaluation.\\
Principal\\
\vspace{0.5em}
Cap              Min(25\textbackslash{}% du CA encaissé, Plafond TND fixe)\\
Mensuel*\\
\vspace{0.5em}
Barème LEAN 6 h/semaine\\
(8–12\\
apprenants)\\
\vspace{0.5em}
Barème           10 h/semaine\\
\vspace{0.5em}
\vspace{0.5em}
\vspace{0.5em}
\vspace{0.5em}
EL\\
STANDARD\\
\vspace{0.5em}
\vspace{0.5em}
\vspace{0.5em}
\vspace{0.5em}
TI\\
(13–20\\
apprenants)\\
\vspace{0.5em}
\vspace{0.5em}
\vspace{0.5em}
\vspace{0.5em}
N\\
DE\\
Barème           2 × 10 h/semaine\\
INTENSIF\\
\vspace{0.5em}
\vspace{0.5em}
\vspace{0.5em}
\vspace{0.5em}
FI\\
(21–30\\
\vspace{0.5em}
\vspace{0.5em}
\vspace{0.5em}
N\\
apprenants)            CO\\
Rémunération 40–80 TND/h\\
indicative\\
\vspace{0.5em}
Mode de            Paiement à l’heure avec enveloppe plafonnée\\
Paiement\\
* Règle d’or contractuelle : Le coût total mensuel du Mentor Pool (Lead Mentors, Experts, TAs) doit respecter min(25\textbackslash{}% du CA encaissé, Plafond TND\\
fixe). Aucune heure hors cap sans ordre de mission signé par RBK.\\
\vspace{0.5em}
\vspace{0.5em}
\vspace{0.5em}
\vspace{0.5em}
\begin{confidential}CONFIDENTIEL — PROJET RBK 2.0                                     Page 274 sur 316\end{confidential}
U \textbar{} Registre des risques opérationnels\\
\vspace{0.5em}
\vspace{0.5em}
Cette annexe présente le registre des risques opérationnels validé pour le lancement du Studio. Les impacts financiers sont évalués en TND1 et les plans d’action\\
sont pilotés par les équipes Ops.\\
\vspace{0.5em}
\vspace{0.5em}
\vspace{0.5em}
\vspace{0.5em}
EL\\
TI\\
Registre de risques v1 (version opérationnelle)\\
\vspace{0.5em}
\vspace{0.5em}
\vspace{0.5em}
\vspace{0.5em}
N\\
DE\\
Échelle (P, I).\\
• P=1 rare ; P=3 plausible ; P=5 très probable.\\
\vspace{0.5em}
\vspace{0.5em}
\vspace{0.5em}
\vspace{0.5em}
FI\\
N\\
• I=1 impact mineur (retard local, peu d’effet) ; I=3 impact significatif (dégradation qualité / retards majeurs) ; I=5 impact critique (arrêt, réputation, financier,\\
conformité).\\
CO\\
\vspace{0.5em}
\vspace{0.5em}
\vspace{0.5em}
\vspace{0.5em}
1\\
Finance/Juridique/Compliance — Dinar tunisien, devise utilisée pour exprimer les montants dans l’offre et les projections financières du programme.\\
\vspace{0.5em}
\vspace{0.5em}
275\\
RBK 2.0 — Whitepaper Architecte Web3                                                                                                           Annexe U. Registre des risques opérationnels\\
\vspace{0.5em}
\vspace{0.5em}
\vspace{0.5em}
\vspace{0.5em}
\begin{center}\begin{tabular}TAB. U.1 : Registre de risques v1 — RBK Web3 Studio (DualTrack EVM/Solana)\end{tabular}\end{center}
Déclen-          Pré-              Conti-\\
ID     Cat.     Risque          P I Sc. Niv. Own.                                                            KPI            St.\\
cheurs           vention           ngence\\
R-01   STRAT/   Sous-remplissage 3 5 15 Élevé Mkt     +Taux de conver- Funnel      admis- Réduire      taille Leads / se- Ou-\\
MKT      de la cohorte pi-             CEO      sion en baisse, sions documenté cohorte,          bas- maine, CVR vert\\
lote (inscriptions                     peu de leads qua-(ICP clair), cam- culer      format lead → en-\\
insuffisantes pour                     lifiés, abandons pagne     alumni, part-time, lancer tretien, CVR\\
lancer).                               après entretien. webinaires,       1 track d’abord, entretien\\
partenariats       renforcer offres →        admis,\\
écoles,     early- bourses/finance- coût par lead\\
bird,    preuves ment.                (CPL).\\
(capstones).\\
R-02   FIN      Dérive      budgé- 3 4 12 Élevé CEO   +Dépenses outils Budget        base Geler       options Variance       Ou-\\
taire    (experts,              Ops    \textgreater{} prévision, ho-+ réserve (10–“Could”,            rené- budget (\textbackslash{}%), vert\\
outils, infra) ré-                     noraires mentors 15\textbackslash{}%),        poli- gocier licences, burn-rate\\
duisant la marge                       imprévus, achats tique     d’achat, basculer        sur mensuel, coût\\
et la crédibilité                      hors budget.     liste       outils alternatives,       /   étudiant,\\
du business plan.                                       “Must/Should”, réduire          inter- marge brute.\\
contrats-cadres ventions          non\\
mentors.           critiques.\\
\vspace{0.5em}
\vspace{0.5em}
\vspace{0.5em}
\vspace{0.5em}
EL\\
R-03   OPS/    Retard de produc- 3 4 12 Élevé HoP     +Absence         de Plan 90 jours, ja- Décaler     lance- \textbackslash{}%    modules Ou-\\
PLANNINGtion pédagogique               Ops      “starter     kits”, lons hebdo, “gol- ment,     réduire prêts (DoD), vert\\
\vspace{0.5em}
\vspace{0.5em}
\vspace{0.5em}
\vspace{0.5em}
TI\\
(labs, templates,                       modules incom- den repos” mini- périmètre initial nb labs testés,\\
\vspace{0.5em}
\vspace{0.5em}
\vspace{0.5em}
\vspace{0.5em}
N\\
rubrics,   repos)                       plets, QA doc malistes, QA La-(1 track), inten- nb               repos\\
\vspace{0.5em}
\vspace{0.5em}
\vspace{0.5em}
\vspace{0.5em}
DE\\
avant démarrage.                        non faite à J-30. TeX/ToC, check- sifier production templates\\
lists.            en sprint, mobili-validés.\\
\vspace{0.5em}
\vspace{0.5em}
\vspace{0.5em}
\vspace{0.5em}
FI\\
ser mentors.\\
\vspace{0.5em}
\vspace{0.5em}
\vspace{0.5em}
\vspace{0.5em}
N\\
R-04   HR       Indisponibilité/ 3 4 12 Élevé CEO     +Annulations       Pool de rem- Remplacer par Heures men- Ou-\\
défection d’un ex-            HoP      répétées, latence plaçants (2 ni- mentor backup, tors livrées vs vert\\
\vspace{0.5em}
\vspace{0.5em}
\vspace{0.5em}
\vspace{0.5em}
CO\\
pert clé (Rust/An-                     réponses,     sur-veaux), sessions basculer    cer- prévues, taux\\
chor, Audit, EVM                       charge d’un seul enregistrées,     taines sessions d’annulation,\\
tooling).                              formateur.        documenta-       en asynchrone, satisfaction\\
tion   complète, réduire   scope apprenants\\
contrats    avec module avancé. (NPS        mo-\\
SLA.                              dule).\\
R-05   PED/     Hétérogénéité du 4 4 16 Cri- HoP       Beaucoup      de Test      d’entrée Séparer en deux Taux comple- Ou-\\
QUAL     niveau des appre-       tique          “rattrapage”, PR+         entretien sous-groupes,      tion labs, nb vert\\
nants (écart trop                      refusées massive- technique,         augmenter TA, PR        rework,\\
grand) → ralentis-                     ment, labs non “Bridge”        2–4 réduire features, taux d’absen-\\
sement ou décro-                       terminés.         semaines,     exi- rallonger sprint, téisme, churn\\
chage.                                                   gences TS/Git, coaching indivi- apprenant\\
regroupement duel.                  (\textbackslash{}%).\\
par binômes.\\
Suite page suivante\\
\vspace{0.5em}
\vspace{0.5em}
\vspace{0.5em}
\vspace{0.5em}
\begin{confidential}CONFIDENTIEL — PROJET RBK 2.0                                                                                       Page 276 sur 316\end{confidential}
RBK 2.0 — Whitepaper Architecte Web3                                                                                                                 Annexe U. Registre des risques opérationnels\\
\vspace{0.5em}
\vspace{0.5em}
Déclen-           Pré-              Conti-\\
ID         Cat.     Risque            P I Sc. Niv. Own.                                                             KPI           St.\\
cheurs            vention           ngence\\
R-06       PED/     Méthode “Studio” 3 5 15 Élevé HoP      CI rouge ignorée, Gates stricts : Bloquer merges, Taux            PR Ou-\\
QUAL     non      appliquée            +    Se- PR sans rubric, li- zero-merge sans “freeze” releases, conformes      vert\\
(merges       sans            cLead2 vrables non do- review + tests, séance de reme-(\textbackslash{}%), CI pass\\
review/tests) →                        cumentés.           templates   PR, diation QA, re- rate,         nb\\
perte de différen-                                         codeowners,     work obligatoire merges       di-\\
ciation et baisse                                          rubric PR.      avant capstone. rects,      dette\\
qualité.                                                                                      technique\\
(issues).\\
R-07       TECH/    Pannes/          4 3 12 Élevé Ops       +Timeouts,       er- Multi-providers, Basculer         sur Taux erreurs Ou-\\
INFRA    instabilités RPC              TL3        reurs blockhash stratégie          fall- environnement RPC, latence vert\\
(EVM/Solana)                             expiré,       rate- back,       cache, local/mocks,       moyenne,\\
pendant labs →                           limit, impossibi- mocks        locaux, déplacer     labs, incidents/\\
pertes de temps,                         lité de confirmer quotas, “incident prolonger sprint, semaine,\\
démo cassée.                             tx.                 drills” planifiés. publier runbook. temps moyen\\
de     reprise\\
(MTTR).\\
R-08       SEC      Capstone          3 5 15 Élevé SecLead Findings     cri- Checklist      sé- Bloquer    livrai- Nb findings Ou-\\
contient      une              + HoP tiques en review, curité,      threat son, exiger patch High/Critical, vert\\
vulnérabilité                          tests     d’abus model      obliga-+ re-audit, dé- couverture\\
grave    (contrat                      échouent,     lo- toire, “security mo sur testnet, tests sécurité,\\
ou logique wal-                        gique    d’authz review sprint”, documenter limi- temps            de\\
let/tx) → risque                       fragile.          tests d’abus mi- tations.             remédiation.\\
réputationnel.                                           nimaux.\\
\vspace{0.5em}
\vspace{0.5em}
\vspace{0.5em}
\vspace{0.5em}
EL\\
R-09       LEGAL/   Confusion      pu- 3 5 15 Élevé CEO +Messages/ads        Charte de com- Rectificatif pu- Nb incidents Ou-\\
\vspace{0.5em}
\vspace{0.5em}
\vspace{0.5em}
\vspace{0.5em}
TI\\
REP      blique   :    pro-              Legal +ambigus, ques- munication,         blic,     retrait comm, taux vert\\
gramme      perçu               Mkt    tions    trading, disclaimers,     contenus, FAQ de questions\\
\vspace{0.5em}
\vspace{0.5em}
\vspace{0.5em}
\vspace{0.5em}
N\\
comme “forma-                          commentaires “software engi- officielle,         for- trading, senti-\\
\vspace{0.5em}
\vspace{0.5em}
\vspace{0.5em}
\vspace{0.5em}
DE\\
tion à gagner                          négatifs, presse neering export”, mation interne ment réseau\\
de l’argent avec                       locale.           interdits expli- marketing, vali- social.\\
\vspace{0.5em}
\vspace{0.5em}
\vspace{0.5em}
\vspace{0.5em}
FI\\
crypto/trading”.                                         cites (trading). dation Legal.\\
\vspace{0.5em}
\vspace{0.5em}
\vspace{0.5em}
\vspace{0.5em}
N\\
R-10       LEGAL/   Incertitudes ré- 3 4 12 Élevé Legal +Blocage banque, Operating model Basculer         sur Nb blocages Ou-\\
\vspace{0.5em}
\vspace{0.5em}
\vspace{0.5em}
\vspace{0.5em}
CO\\
OPS      glementaires                  CEO    partenaires       par     scénarios modèle export paiement,          vert\\
locales impactant                    frileux,    refus (local/export/- +      facturation time-to-\\
paiements/com-                       sponsor,     avis hybride), paie- service, limiter contract,\\
munication/par-                      contradictoires. ments classiques, mentions, renfor- retours     ju-\\
tenariats.                                             docs juridiques. cer disclaimers. ridiques,\\
incidents\\
conformité.\\
Suite page suivante\\
\vspace{0.5em}
\vspace{0.5em}
\vspace{0.5em}
\vspace{0.5em}
2\\
Finance/Juridique/Compliance — Dinar tunisien, devise utilisée pour exprimer les montants dans l’offre et les projections financières du programme.\\
3\\
Finance/Juridique/Compliance — Dinar tunisien, devise utilisée pour exprimer les montants dans l’offre et les projections financières du programme.\\
\vspace{0.5em}
\vspace{0.5em}
\begin{confidential}CONFIDENTIEL — PROJET RBK 2.0                                                                                              Page 277 sur 316\end{confidential}
RBK 2.0 — Whitepaper Architecte Web3                                                                                                                   Annexe U. Registre des risques opérationnels\\
\vspace{0.5em}
\vspace{0.5em}
Déclen-           Pré-               Conti-\\
ID         Cat.     Risque           P I Sc. Niv. Own.                                                               KPI            St.\\
cheurs            vention            ngence\\
R-11       FIN/FX   Risque change / 3 3 9      Mo- CEO      +Hausse     coûts Hedging “soft” : Augmenter             Part coûts en Ou-\\
inflation (coûts           déré Ops      SaaS4 , factures buffer, pricing ré- prix prochaines devise      (\textbackslash{}%), vert\\
outils en devise,                        en USD5 /EUR, visable, négocia- cohortes,           ré-variance coût\\
recettes en TND)                         variation TND. tion annuelle, al- duire          outils outil,   mar-\\
impactant marge.                                          ternatives open- non        critiques, ge/étudiant.\\
source6 .           contractualiser\\
en TND si pos-\\
sible.\\
R-12       OPS/     Absence        de 3 4 12 Élevé Ops      +Chiffres contra- Single source of Gel       rédaction, Nb       inco- Ou-\\
PROC     processus docu-                HoP       dictoires,   sec- truth (Factsheet), nettoyage,    au- hérences      vert\\
mentaire     (ver-                       tions dupliquées,CHANGELOG, dit              complet, détectées,\\
sioning, sources,                        ToC erronée, in- build CI doc, réémission PDF, build LaTeX\\
changelog)     →                         puts manquants. règles d’édition. checklist qualité. green, temps\\
incohérences du                                                                                 correction.\\
livre blanc.\\
R-13       PART     Promesses       de 3 3 9   Mo- CEO      +Absence        de Pipeline parte- Remplacer par Nb partena- Ou-\\
NERS     partenariats (eco-         déré HoP      logos/letters, an- nariats    (tiers), mentors locaux/- riats signés, vert\\
system, mentors,                         nulations, délais MOU       simples, diaspora, décaler heures men-\\
sponsors)     non                        de réponse.        calendrier office annonce, recen- tors réelles,\\
concrétisées →                                              hours, transpa- trer sur preuves taux présence\\
baisse crédibilité.                                         rence “targets vs internes.          invités.\\
signed”.\\
\vspace{0.5em}
\vspace{0.5em}
\vspace{0.5em}
\vspace{0.5em}
EL\\
R-14       CAREER   Insertion     pro 3 5 15 Élevé Career     Peu       d’offres, Service carrière Prolonger      ac- Placement       Ou-\\
insuffisante                   + HoP      apprenants sans structuré : port- compagnement, S1–S24              vs vert\\
\vspace{0.5em}
\vspace{0.5em}
\vspace{0.5em}
\vspace{0.5em}
TI\\
(remote/boun-                             portfolio, faible folio      GitHub, booster “boun- S25–S48,\\
ties/placements)                          activité contribu- mock interviews, ties      sprints”, nb bounties,\\
\vspace{0.5em}
\vspace{0.5em}
\vspace{0.5em}
\vspace{0.5em}
N\\
→        promesse                         tions.              bounties   enca- mentorat     CV/- part     portfo-\\
\vspace{0.5em}
\vspace{0.5em}
\vspace{0.5em}
\vspace{0.5em}
DE\\
d’employabilité                                               drées, demo day LinkedIn, parte- lios “ready”,\\
affaiblie.                                                    panel.           nariats RH.        offres reçues.\\
\vspace{0.5em}
\vspace{0.5em}
\vspace{0.5em}
\vspace{0.5em}
FI\\
R-15       PED/     Usage irrespon- 4 4 16 Cri- HoP +PR          “copiées”, Règles IA : “IA Évaluations          Taux échec Ou-\\
\vspace{0.5em}
\vspace{0.5em}
\vspace{0.5em}
\vspace{0.5em}
N\\
ETHICS   sable de l’IA (vibe    tique SecLead incapacité expli-=        stagiaire”, orales, rewriting oral defense, vert\\
\vspace{0.5em}
\vspace{0.5em}
\vspace{0.5em}
\vspace{0.5em}
CO\\
coding) : plagiat,                   quer code, bugs explication obli- en live, sanctions taux rework\\
compréhension                        répétés, triche gatoire,         oral pédagogiques, PR,         simila-\\
faible, code non                     sur tests.         defense,     jour- rework obliga- rité         code,\\
maîtrisé.                                               naux décisions toire.                incidents de\\
(ADR).                                triche.\\
R-16       TECH/    Gestion secrets 3 5 15 Élevé SecLead Secrets   com- Politique secrets, Rotation immé- Nb fuites dé- Ou-\\
SEC      / clés (API keys,            + Ops mités,       clés scans (gitleaks), diate, suppres- tectées, temps vert\\
wallets)      mal                    partagées   sur .env.example, ro- sion historique, rotation,\\
cadrée → exposi-                     chat, env non tation clés, sand- incident report, conformité\\
tion accidentelle.                   sécurisé.       box devnet/test- formation     ex- scans repo.\\
net.              press.\\
R-17       OPS/LMS Dépendance à un 3 3 9       Mo- Ops        Panne     service, Backups, miroirs, Basculer canal Temps indis- Ou-\\
outil unique (LM-           déré           perte       accès, exports hebdo, secours,      res- ponibilité,   vert\\
S/Discord/Git)                             suppression acci- documentation taurer backup, MTTR, nb ba-\\
ou indisponibilité                         dentelle.          accès, comptes replanifier dead- ckups réussis.\\
→ perte continui-                                             admin multiples. lines.\\
té.\\
Suite page suivante\\
\vspace{0.5em}
\vspace{0.5em}
4\\
Finance/Juridique/Compliance — Dinar tunisien, devise utilisée pour exprimer les montants dans l’offre et les projections financières du programme.\\
5\\
Finance/Juridique/Compliance — Dinar tunisien, devise utilisée pour exprimer les montants dans l’offre et les projections financières du programme.\\
6\\
Finance/Juridique/Compliance — Dinar tunisien, devise utilisée pour exprimer les montants dans l’offre et les projections financières du programme.\\
\vspace{0.5em}
\vspace{0.5em}
\begin{confidential}CONFIDENTIEL — PROJET RBK 2.0                                                                                                Page 278 sur 316\end{confidential}
RBK 2.0 — Whitepaper Architecte Web3                                                                                                           Annexe U. Registre des risques opérationnels\\
\vspace{0.5em}
\vspace{0.5em}
Déclen-           Pré-             Conti-\\
ID     Cat.      Risque            P I Sc. Niv. Own.                                                             KPI            St.\\
cheurs            vention          ngence\\
R-18   SCOPE     Programme trop 4 4 16 Cri- HoP           +Backlog qui ex- Modularisation Couper “Could”, Taux comple- Ou-\\
dense (effet “tout    tique CEO           plose, labs non strict       Must/- étaler sur se- tion, heures vert\\
sur tout”) → fa-                          terminés, sprints Should/Could, maines           sup- sup,     NPS,\\
tigue, baisse qua-                        débordent.        DoD clair, ga- plémentaires,        backlog carry-\\
lité, abandon.                                              ting,     rythme proposer option over.\\
réaliste.        “Elite” séparée.\\
R-19   REP/    Document (livre 3 4 12 Élevé Ops           +ToC incorrecte, QA éditorial : Gel + refonte Nb             défauts Ou-\\
QUALITY blanc) manque                HoP            césures, tables build CI doc, mise en page,QA,               temps vert\\
de      cohérence                           mal     alignées, checklist, styles, réémission PDF, correction,\\
pro (typo, tables,                          figures     sans latex templates, validation finale retours      re-\\
figures) → im-                              légendes.         relectures croi-“Go/No-Go”.        viewers\\
pact crédibilité                                              sées.                              externes.\\
investisseurs/par-\\
tenaires.\\
R-20   FIN/      Modèle       ISA/- 2 4 8    Mo- Legal +Ambiguïtés        Option ISA “à Retirer ISA de Nb demandes Ou-\\
MODEL     bourses       mal           déré CEO   contractuelles, l’étude”, contrats l’offre     pilote, ISA,  litiges, vert\\
cadré (contrats,                       litiges,   incom- clairs,    seuils, basculer     sur taux impayés,\\
recouvrement,                          préhensions       transparence,      paiement éche- satisfaction\\
perception)     →                      candidats.        alternatives       lonné + bourses candidats.\\
risque légal/répu-                                      (échelonne-         limitées.\\
\vspace{0.5em}
\vspace{0.5em}
\vspace{0.5em}
\vspace{0.5em}
EL\\
tation.                                                  ment).\\
R-21   IP/LEGAL Propriété intellec- 3 3 9    Mo- Legal +Refus de publier Clauses          IP Retirer      code Nb incidents Ou-\\
\vspace{0.5em}
\vspace{0.5em}
\vspace{0.5em}
\vspace{0.5em}
TI\\
tuelle des caps-             déré CEO   code, désaccord simples       (open- sensible, publier IP, taux repos vert\\
tones (RBK vs étu-                      sur réutilisation, source par défaut version anony- publics,       ac-\\
diants) mal défi-                       tensions à demo pour pédagogie) misée,          accord cords signés.\\
\vspace{0.5em}
\vspace{0.5em}
\vspace{0.5em}
\vspace{0.5em}
N\\
nie → conflits.                         day.              +       exceptions, spécifique projet\\
\vspace{0.5em}
\vspace{0.5em}
\vspace{0.5em}
\vspace{0.5em}
DE\\
consentement       incubé.\\
écrit.\\
R-22   SEC/OPS Absence de PCA/- 2 4 8        Mo- Ops      +Perte repos, in- Backups,       ac- Restauration,   RPO/RTO       Ou-\\
\vspace{0.5em}
\vspace{0.5em}
\vspace{0.5em}
\vspace{0.5em}
FI\\
PRA     minimum               déré CEO      dispo outils, cor- cès MFA, PRA bascule miroirs, atteints, tests vert\\
(continuité)    en                          ruption fichiers, doc, duplication replanification, restauration,\\
cas     d’incident                          vol matériel.      dépôts,   droits communication incidents\\
\vspace{0.5em}
\vspace{0.5em}
\vspace{0.5em}
\vspace{0.5em}
N\\
majeur      (perte                                             minimaux.        apprenants.     sécurité.\\
données,     indis-                                                                            CO\\
po).\\
TL\\
R-23   TECH/     Changements     3 3 9       Mo- EVM        Breaking         Lock   versions, Patch       starter Nb incidents Ou-\\
EVM       tooling    EVM              déré           changes, erreurs images docker, kit, session “up- build, temps vert\\
(Hardhat/Foun-                             compilation,     guides     repro- grade”, bascule fix, fréquence\\
dry, forks, L2)                            docs obsolètes. ductibles, veille sur       versions updates.\\
perturbent  mo-                                             mensuelle.        stables.\\
dules.\\
TL\\
R-24   TECH/     Évolution rapide 3 3 9      Mo- SOL        Erreurs Anchor, Lock     versions, Migration    en- Nb incidents, Ou-\\
SOL       Solana/Anchor               déré           IDL    mismatch, starter kits, CI, cadrée, bascule temps patch, vert\\
(versions, APIs)                           changements      notes migration, versions, gel sur stabilité CI\\
cassant    labs/-                          CLI, docs diver-veille.            release   stable Anchor.\\
tests.                                     gentes.                            promo.\\
R-25   SEC/      Culture “security- 3 4 12 Élevé SecLead Retours         “ça Rubric      pénali- Refus validation Score sécurité Ou-\\
CULT      first” insuffisam-              + HoP marche mais…”, sante           sécurité, capstone, sprint moyen,      nb vert\\
ment      intégrée                      absence threat week             “audit de remédiation, findings, taux\\
(les étudiants op-                      model,      faibles sprint”, oral de- re-audit croisé. capstones\\
timisent la démo,                       tests abus.         fense sur risques,                  “pass security\\
pas la sûreté).                                             checklists.                          gate”.\\
\vspace{0.5em}
\vspace{0.5em}
\vspace{0.5em}
\vspace{0.5em}
\begin{confidential}CONFIDENTIEL — PROJET RBK 2.0                                                                                           Page 279 sur 316\end{confidential}
RBK 2.0 — Whitepaper Architecte Web3                                                                            Annexe U. Registre des risques opérationnels\\
\vspace{0.5em}
\vspace{0.5em}
Note d’exploitation.      Ce registre v1 est conçu pour être vivant : à chaque COPIL/COSUI (hebdo ou bi-hebdo), mettre à jour P, I, statut, actions et reporter les\\
décisions dans le CHANGELOG (et/ou un journal de pilotage).\\
\vspace{0.5em}
\vspace{0.5em}
\vspace{0.5em}
\vspace{0.5em}
EL\\
TI\\
N\\
DE\\
FI\\
N\\
CO\\
\vspace{0.5em}
\vspace{0.5em}
\vspace{0.5em}
\vspace{0.5em}
\begin{confidential}CONFIDENTIEL — PROJET RBK 2.0                                         Page 280 sur 316\end{confidential}
V \textbar{} Plan d’actions risques \textbackslash{}& tableau de bord KPI (v1)\\
\vspace{0.5em}
\vspace{0.5em}
Cette annexe détaille le plan d’action opérationnel pour mitiger les risques identifiés, ainsi que le tableau de bord de pilotage (KPI) pour le suivi hebdomadaire.\\
Le dispositif s’appuie sur un pilotage Ops et sur la traçabilité des leads dans un CRM.\\
\vspace{0.5em}
\vspace{0.5em}
\vspace{0.5em}
\vspace{0.5em}
EL\\
TI\\
Registre des actions risques v1 (transforme chaque risque en tâches exécutables)\\
\vspace{0.5em}
\vspace{0.5em}
\vspace{0.5em}
\vspace{0.5em}
N\\
Référence planning.          On note T0 la date de démarrage officiel de la cohorte. Les échéances sont exprimées relativement à T0.\\
\vspace{0.5em}
\vspace{0.5em}
\vspace{0.5em}
\vspace{0.5em}
DE\\
FI\\
N\\
CO\\
\vspace{0.5em}
\vspace{0.5em}
\vspace{0.5em}
\vspace{0.5em}
281\\
RBK 2.0 — Whitepaper Architecte Web3                                                                      Annexe V. Plan d’actions risques \textbackslash{}& tableau de bord KPI (v1)\\
\vspace{0.5em}
\vspace{0.5em}
\begin{center}\begin{tabular}TAB. V.1 : Registre des actions risques v1 — RBK Web3 Studio\end{tabular}\end{center}
Action                                         Dé-   Éché-           Livrable       Dépend- Preuve\\
Act. Risque                    Type      Prio Owner                    Statut\\
(verbe + livrable)                             but   ance            (DoD)          ances   (Trace)\\
ID\\
A- R-01     Définir ICP + mes- Préventif P1   Mkt + CEO    T0-90 T0-75 À faire   ICP validé +   Vali-      PDF/URL\\
001         sage de valeur (1                                                    FAQ publiée    dation\\
page) + FAQ anti-                                                    (site + PDF)   Legal\\
trading\\
A- R-01     Mettre en place    Préventif P1   Mkt + Ops    T0-90 T0-60 À faire   Pipeline      Outil   Capture +\\
002         funnel admissions                                                    (Lead→        CRM     procédure\\
(CRM léger + pipe-                                                   Entretien→\\
line)                                                                Admis) opéra-\\
tionnel\\
A- R-01    Lancer 3 webinaires Préventif P1   Mkt + HoP    T0-75 T0-45 À faire   3 replays +   Landing Liens replays\\
003        “Web3 careers after                                                   200 leads     + tra-  + stats\\
AI1 ” + collecte                                                      qualifiés     cking\\
leads\\
\vspace{0.5em}
\vspace{0.5em}
\vspace{0.5em}
\vspace{0.5em}
EL\\
A- R-01    Plan “early bird + Préventif P2    CEO + Ops    T0-75 T0-60 À faire   Offre tarifaire Legal     Doc offre +\\
004        alumni” (grille prix                                                  validée + CGV             CGV\\
\vspace{0.5em}
\vspace{0.5em}
\vspace{0.5em}
\vspace{0.5em}
TI\\
+ conditions)\\
\vspace{0.5em}
\vspace{0.5em}
\vspace{0.5em}
\vspace{0.5em}
N\\
A- R-01    Plan B : lancer 1     Conti-  P1   CEO + HoP    T0-30 T0-14 À faire   Règle Go/No- KPI ad- Note Go/No-\\
\vspace{0.5em}
\vspace{0.5em}
\vspace{0.5em}
\vspace{0.5em}
DE\\
005        seul track si \textless{} seuil ngence                                          Go + com-      missions Go\\
cohort                                                                munication\\
\vspace{0.5em}
\vspace{0.5em}
\vspace{0.5em}
\vspace{0.5em}
FI\\
prête\\
A- R-02     Construire budget Préventif P1    CEO + Ops    T0-90 T0-70 À faire   Budget table + Business Fichier budget\\
\vspace{0.5em}
\vspace{0.5em}
\vspace{0.5em}
\vspace{0.5em}
N\\
006         détaillé (Must/-                                                     règles d’arbi- plan\\
\vspace{0.5em}
\vspace{0.5em}
\vspace{0.5em}
A- R-02\\
Should/Could) +\\
réserve 15\textbackslash{}%\\
Contractualiser\\
CO\\
Préventif P1   CEO + HoP    T0-75 T0-50 À faire\\
trage signées\\
\vspace{0.5em}
Contrats signés Pool    Contrats\\
007         mentors (forfait +                                                   + calendrier    mentors\\
SLA + calendrier)                                                    interventions\\
A- R-02     Mettre en place    Préventif P2   Ops          T0-60 T0-30 À faire   Tableau va-     Budget Dashboard\\
008         reporting mensuel                                                    riance + alerte A-006\\
variance budget                                                      seuil 10\textbackslash{}%\\
A- R-03     Geler structure    Préventif P1   HoP          T0-90 T0-80 À faire   DoD module      –       Doc DoD\\
009         syllabus + DoD par                                                   (labs+tests+\\
module (checklist)                                                   doc) validé\\
TL\\
A- R-03     Produire “Golden Préventif P1                 T0-80 T0-45 À faire    2 repos tem-   A-009      URLs Git\\
EVM + Solana\\
010         repos” : starter kits                                                plates (CI\\
EVM + Solana                                                         verte)\\
Suite page suivante\\
\vspace{0.5em}
1\\
Finance/Juridique/Compliance — Dinar tunisien, devise utilisée pour exprimer les montants dans l’offre et les projections financières du programme.\\
\vspace{0.5em}
\vspace{0.5em}
\begin{confidential}CONFIDENTIEL — PROJET RBK 2.0                                                 Page 282 sur 316\end{confidential}
RBK 2.0 — Whitepaper Architecte Web3                                                                         Annexe V. Plan d’actions risques \textbackslash{}& tableau de bord KPI (v1)\\
\vspace{0.5em}
\vspace{0.5em}
Action                                            Dé-   Éché-           Livrable        Dépend- Preuve\\
Act. Risque                      Type    Prio Owner                       Statut\\
(verbe + livrable)                                but   ance            (DoD)           ances   (Trace)\\
ID\\
A- R-03     QA labs (reproduc- Préventif P1   TA + Ops        T0-60 T0-20 À faire   100\textbackslash{}% labs tes- A-010       Checklist QA\\
011         tibilité sur machine                                                    tés + correctifs           signée\\
neuve)\\
Avancement\\
A- R-03     Plan de réduction Conti-        P2   HoP          T0-30 T0-14 À faire   Liste modules         Note\\
QA\\
012         périmètre (si retard) ngence                                            “cut” + com\\
interne\\
A- R-04    Constituer pool       Préventif P1    HoP          T0-90 T0-60 À faire   Liste + dis-  Réseau Doc pool\\
013        de remplaçants (2                                                        ponibilité +\\
profils/compétence)                                                      tarifs\\
A- R-04    Enregistrer cours     Préventif P2    HoP + Men- T0-14 T0+     À faire   Replays +     Calen-  Liens replays\\
014        critiques + notes                     tors             14                supports ver- drier\\
complètes                                                                sionnés\\
A- R-05    Créer test d’en-      Préventif P1    HoP          T0-75 T0-60 À faire   Test + barème ICP     PDF test\\
015        trée (JS/TS2 , Git,                                                      + seuils\\
logique, anglais)\\
\vspace{0.5em}
\vspace{0.5em}
\vspace{0.5em}
\vspace{0.5em}
EL\\
A- R-05    Mettre “Bridge”       Préventif P1    HoP + TA     T0-60 T0-30 À faire   Programme      A-015       Syllabus\\
\vspace{0.5em}
\vspace{0.5em}
\vspace{0.5em}
\vspace{0.5em}
TI\\
016        2–4 semaines                                                             bridge + mini-             bridge\\
(TS+Git+tests)                                                           projet\\
\vspace{0.5em}
\vspace{0.5em}
\vspace{0.5em}
\vspace{0.5em}
N\\
A- R-05    Mettre en place       Correctif P2    HoP          T0+ T0+7 À faire      Matrice bi-    Cohorte     Planning\\
\vspace{0.5em}
\vspace{0.5em}
\vspace{0.5em}
\vspace{0.5em}
DE\\
017        binômage “strong/-                                 0                     nômes +\\
weak” + tutorat                                                          règles\\
\vspace{0.5em}
\vspace{0.5em}
\vspace{0.5em}
\vspace{0.5em}
FI\\
A- R-06    Implémenter règles Préventif P1       Ops + Tech   T0-30 T0-14 À faire   Protection     Repos       Config Git\\
\vspace{0.5em}
\vspace{0.5em}
\vspace{0.5em}
\vspace{0.5em}
N\\
018        Git : PR obligatoires                 Lead                               main + CI      A-010\\
\vspace{0.5em}
\vspace{0.5em}
CO\\
+ branch protection                                                      required\\
HoP\\
A- R-06    Créer rubric PR       Préventif P1                 T0-30 T0-14 À faire   Templates        A-018     Fichiers .gi-\\
+ SecLead\\
019        + template PR +                                                          actifs sur repos           thub\\
codeowners\\
A- R-06    Mettre “merge         Correctif P2    HoP          T0+ T0+     À faire   Procédure +     KPI CI     Procédure\\
020        freeze” si CI rouge                                0   90                déclencheurs\\
\textgreater{} 24h\\
A- R-07    Choisir 2 provi-      Préventif P1    Ops + Tech   T0-45 T0-20 À faire   Runbook         Budget     Doc runbook\\
021        ders RPC/track +                      Lead                               RPC + clés\\
fallback strategy                                                        + quotas\\
A- R-07    Ajouter mocks         Préventif P2    Tech Lead    T0-30 T0-10 À faire   Scripts mocks   A-010      Repo\\
022        locaux + modes                                                           + doc usage\\
“offline labs”\\
Suite page suivante\\
\vspace{0.5em}
\vspace{0.5em}
\vspace{0.5em}
2\\
Finance/Juridique/Compliance — Dinar tunisien, devise utilisée pour exprimer les montants dans l’offre et les projections financières du programme.\\
\vspace{0.5em}
\vspace{0.5em}
\begin{confidential}CONFIDENTIEL — PROJET RBK 2.0                                                     Page 283 sur 316\end{confidential}
RBK 2.0 — Whitepaper Architecte Web3                                                                Annexe V. Plan d’actions risques \textbackslash{}& tableau de bord KPI (v1)\\
\vspace{0.5em}
\vspace{0.5em}
Action                                          Dé- Éché-             Livrable        Dépend- Preuve\\
Act. Risque                      Type      Prio Owner                   Statut\\
(verbe + livrable)                              but ance              (DoD)           ances   (Trace)\\
ID\\
A- R-07     Tester incident      Préventif P2   Ops         T0-7 T0+7 À faire     Rapport drill + A-021      CR incident\\
023         drill “RPC down”                                                      corrections                drill\\
(scénario)\\
Sec\\
A- R-08     Exiger threat model Préventif P1                T0+ T0+ À faire       Threat model    Template Docs\\
Lead\\
024         1 page par capstone                             14    21              signé\\
+ validation Sec-\\
Lead\\
HoP\\
A- R-08     Instaurer “Security Préventif P1                T0+ T0+ À faire       Checklist sécu- KPI     Evidence\\
+ SecLead\\
025         Sprint” obligatoire                             35    49              rité complétée findings checklist\\
(S19)\\
Sec\\
A- R-08     Audit croisé inter- Préventif P2                T0+ T0+ À faire       Rapport audit A-025        PDF findings\\
Lead\\
026         équipes (2 revie-                               49    56              croisé\\
wers)\\
A- R-09     Charte marketing Préventif P1       Legal + Mkt T0-75 T0-60 À faire   Charte signée   CEO        Doc charte\\
\vspace{0.5em}
\vspace{0.5em}
\vspace{0.5em}
\vspace{0.5em}
EL\\
027         + disclaimers\\
\vspace{0.5em}
\vspace{0.5em}
\vspace{0.5em}
\vspace{0.5em}
TI\\
“no trading / no\\
speculation”\\
\vspace{0.5em}
\vspace{0.5em}
\vspace{0.5em}
\vspace{0.5em}
N\\
A- R-09     Relecture Legal      Préventif P1   Legal       T0-45 T0-20 À faire   Visa Legal      A-027      Liste URLs\\
\vspace{0.5em}
\vspace{0.5em}
\vspace{0.5em}
\vspace{0.5em}
DE\\
028         de toutes pages                                                       (liste URLs)\\
publiques + ads\\
\vspace{0.5em}
\vspace{0.5em}
\vspace{0.5em}
\vspace{0.5em}
FI\\
A- R-10     Définir modèle       Préventif P1   Legal + CEO T0-75 T0-45 À faire   Contrat + CGV –            PDFs\\
\vspace{0.5em}
\vspace{0.5em}
\vspace{0.5em}
\vspace{0.5em}
N\\
029         contractuel “soft-                                                    + mentions\\
\vspace{0.5em}
\vspace{0.5em}
CO\\
ware training ex-\\
port”\\
A- R-10     Scénarios opération- Préventif P2   CEO + Ops T0-60 T0-30 À faire     Note scénarios A-029       Note\\
030         nels (local / export                                                  + triggers\\
/ hybride)\\
A- R-11     Lister coûts en      Préventif P3   CEO + Ops T0-60 T0-30 À faire     Tableau de-     Budget     Tableau\\
031         devise + buffer +                                                     vises + clauses\\
règle révision prix\\
A- R-12     Créer “Factsheet” Préventif P1      Ops         T0-45 T0-25 À faire   Factsheet       Livre      Factsheet\\
032         (source unique                                                        versionnée +    blanc\\
chiffres) + ver-                                                      owners\\
rouillage\\
A- R-12     Mettre CI build      Préventif P2   Ops         T0-30 T0-14 À faire   Build auto +    Repo doc CI logs\\
033         LaTeX + lint doc                                                      alertes\\
(si applicable)\\
Suite page suivante\\
\vspace{0.5em}
\vspace{0.5em}
\vspace{0.5em}
\vspace{0.5em}
\begin{confidential}CONFIDENTIEL — PROJET RBK 2.0                                          Page 284 sur 316\end{confidential}
RBK 2.0 — Whitepaper Architecte Web3                                                                 Annexe V. Plan d’actions risques \textbackslash{}& tableau de bord KPI (v1)\\
\vspace{0.5em}
\vspace{0.5em}
Action                                           Dé- Éché-             Livrable        Dépend- Preuve\\
Act. Risque                      Type      Prio Owner                    Statut\\
(verbe + livrable)                               but ance              (DoD)           ances   (Trace)\\
ID\\
A- R-12     Mettre CHANGE-       Préventif P2   Ops          T0-30 T0-14 À faire   CHANGELOG Repo doc Fichiers\\
034         LOG + règles                                                           + CONTRIBU-\\
contribution (docs)                                                    TING\\
A- R-13     Pipeline partena-    Préventif P2   CEO + HoP T0-60 T0-30 À faire      Tableau pipe- Réseau Tableau\\
035         riats (Target/Con-                                                     line + owners\\
tacted/Signed)\\
A- R-13     Obtenir 2 lettres    Préventif P2   HoP          T0-45 T0-20 À faire   2 emails/LOI    A-035      PDFs\\
036         d’intention men-\\
tors/invités\\
A- R-14     Définir “Career      Préventif P1   Career + HoP T0-45 T0-14 À faire   Kit complet + –            Pack PDF\\
037         Pack” (portfolio,                                                      calendrier\\
mock interviews,\\
bounties)\\
A- R-14     Instaurer “Bounties Préventif P2    Career       T0+ T0+ À faire       1 bounty/étu- Réseau       Table bounties\\
038         Sprint” encadré                                  35    49              diant min\\
\vspace{0.5em}
\vspace{0.5em}
\vspace{0.5em}
\vspace{0.5em}
EL\\
(semaine 8+)\\
A- R-14     Mettre tableau       Correctif P2   Career       T0+ T0+ À faire       Dashboard       CRM        Dashboard\\
\vspace{0.5em}
\vspace{0.5em}
\vspace{0.5em}
\vspace{0.5em}
TI\\
039         placements (S1–                                  60    180             placements\\
\vspace{0.5em}
\vspace{0.5em}
\vspace{0.5em}
\vspace{0.5em}
N\\
S24 / S25–S48) +\\
\vspace{0.5em}
\vspace{0.5em}
\vspace{0.5em}
\vspace{0.5em}
DE\\
relances\\
A- R-15     Politique IA (ce qui Préventif P1   HoP          T0-30 T0-14 À faire   Règles IA       –          Politique\\
\vspace{0.5em}
\vspace{0.5em}
\vspace{0.5em}
\vspace{0.5em}
FI\\
040         est autorisé/inter-                                                    signées +\\
dit) + oral defense                                                    grille orale\\
\vspace{0.5em}
\vspace{0.5em}
\vspace{0.5em}
\vspace{0.5em}
N\\
A- R-15     Mettre “ADR obli- Préventif P2      Tech Lead    T0+ T0+ À faire       ADR template Repo          ADRs\\
041         gatoire” pour déci-\\
sions majeures             CO       SecLead\\
0     90              + usage vérifié\\
\vspace{0.5em}
\vspace{0.5em}
A- R-16     Mettre scan secrets Préventif P1                 T0-30 T0-14 À faire   Scan actif en CI A-018     CI logs\\
+ Ops\\
042         (gitleaks) + règles                                                    + doc\\
.env.example\\
Sec\\
A- R-16     Session “keys hy-    Préventif P2                T0+ T0+7 À faire      Support +       –          Résultats quiz\\
Lead\\
043         giene” (30 min) +                                0                     quiz\\
quiz\\
A- R-17     Mettre backups       Préventif P2   Ops          T0-30 T0-10 À faire   Backup + test Accès        Log restore\\
044         hebdo (repos/docs)                                                     restore       admin\\
+ procédure restore\\
A- R-18     Appliquer MoSCoW Préventif P1       HoP          T0+ T0+ À faire       MoSCoW          Tem-       Tickets\\
045         à chaque sprint +                                0     90              présent sur     plates\\
gating Must                                                            tickets\\
Suite page suivante\\
\vspace{0.5em}
\vspace{0.5em}
\vspace{0.5em}
\vspace{0.5em}
\begin{confidential}CONFIDENTIEL — PROJET RBK 2.0                                           Page 285 sur 316\end{confidential}
RBK 2.0 — Whitepaper Architecte Web3                                                                         Annexe V. Plan d’actions risques \textbackslash{}& tableau de bord KPI (v1)\\
\vspace{0.5em}
\vspace{0.5em}
Action                                            Dé- Éché-             Livrable       Dépend- Preuve\\
Act. Risque                        Type      Prio Owner                   Statut\\
(verbe + livrable)                                but ance              (DoD)          ances   (Trace)\\
ID\\
A- R-18     Mettre indicateur      Préventif P2   HoP + TA    T0+ T0+ À faire       Tableau charge Outils   Tableau\\
046         surcharge (heures                                 7     90              hebdo\\
sup + carry-over)\\
A- R-19     Checklist QA PDF Préventif P2         Ops         T0-30 T0-14 À faire   Checklist      Factsheet Checklist\\
047         (ToC, tables, figures,                                                  signée\\
cohérence chiffres)\\
A- R-19     Relectures croisées Préventif P2      Ops + HoP T0-21 T0-7 À faire      CR relecture + A-047    CR\\
048         (2 reviewers) +                                                         patches\\
corrections\\
A- R-20     Définir politique      Préventif P3   CEO + Legal T0-45 T0-20 À faire   Politique      –        Doc\\
049         bourses simple                                                          bourses\\
(pilot) + critères\\
A- R-21     Ajouter clause IP3 Préventif P2       Legal       T0-45 T0-20 À faire   Clause dans    A-029    Contrat\\
050         (open-source par dé-                                                    contrat\\
faut + exceptions)\\
\vspace{0.5em}
\vspace{0.5em}
\vspace{0.5em}
\vspace{0.5em}
EL\\
A- R-22     Écrire PCA/PRA mi- Préventif P2       Ops         T0-30 T0-10 À faire   Doc PCA/PRA A-044       Doc + logs\\
051         nimum (RPO/RTO,                                                         + test\\
\vspace{0.5em}
\vspace{0.5em}
\vspace{0.5em}
\vspace{0.5em}
TI\\
backups, rôles)\\
\vspace{0.5em}
\vspace{0.5em}
\vspace{0.5em}
\vspace{0.5em}
N\\
A- R-23     Verrouiller versions Préventif P2     TL (EVM)    T0-30 T0-14 À faire   Dockerfile +   A-010    Repo\\
\vspace{0.5em}
\vspace{0.5em}
\vspace{0.5em}
\vspace{0.5em}
DE\\
052         + docker image +                                                        doc\\
guide install\\
\vspace{0.5em}
\vspace{0.5em}
\vspace{0.5em}
\vspace{0.5em}
FI\\
A- R-24     Verrouiller versions Préventif P2     TL (Solana) T0-30 T0-14 À faire   Notes versions A-010    Repo\\
053         Anchor/Solana +                                                         + doc\\
\vspace{0.5em}
\vspace{0.5em}
\vspace{0.5em}
\vspace{0.5em}
N\\
guide migration\\
\vspace{0.5em}
A- R-25\\
054\\
Créer “Security\\
Gate” capstone :\\
CO\\
Préventif P1\\
Sec\\
Lead\\
T0+ T0+ À faire\\
35    56\\
Critères gate + A-025\\
checklist\\
Doc gate\\
\vspace{0.5em}
seuil findings +\\
blocage\\
SecLead\\
A- R-25     Instaurer “oral        Préventif P2               T0+ T0+ À faire       Grille + PV    A-054    PV\\
+ HoP\\
055         defense sécurité”                                 56    63\\
(10 min/équipe)\\
\vspace{0.5em}
\vspace{0.5em}
\vspace{0.5em}
\vspace{0.5em}
3\\
Finance/Juridique/Compliance — Dinar tunisien, devise utilisée pour exprimer les montants dans l’offre et les projections financières du programme.\\
\vspace{0.5em}
\vspace{0.5em}
\begin{confidential}CONFIDENTIEL — PROJET RBK 2.0                                                  Page 286 sur 316\end{confidential}
RBK 2.0 — Whitepaper Architecte Web3                                                                              Annexe V. Plan d’actions risques \textbackslash{}& tableau de bord KPI (v1)\\
\vspace{0.5em}
\vspace{0.5em}
Mode d’emploi.     À chaque COPIL/COSUI : (i) mettre à jour statuts, (ii) reporter décisions dans un journal, (iii) créer de nouvelles actions si un risque change de nature.\\
\vspace{0.5em}
\vspace{0.5em}
\vspace{0.5em}
\vspace{0.5em}
EL\\
TI\\
N\\
DE\\
FI\\
N\\
CO\\
\vspace{0.5em}
\vspace{0.5em}
\vspace{0.5em}
\vspace{0.5em}
\begin{confidential}CONFIDENTIEL — PROJET RBK 2.0                                                     Page 287 sur 316\end{confidential}
RBK 2.0 — Whitepaper Architecte Web3                           Annexe V. Plan d’actions risques \textbackslash{}& tableau de bord KPI (v1)\\
\vspace{0.5em}
\vspace{0.5em}
Tableau de pilotage KPI (v1) — Admissions, Qualité, Sécurité, Ops, Carrière, Finance\\
\vspace{0.5em}
\vspace{0.5em}
\vspace{0.5em}
\vspace{0.5em}
EL\\
TI\\
N\\
DE\\
FI\\
N\\
CO\\
\vspace{0.5em}
\vspace{0.5em}
\vspace{0.5em}
\vspace{0.5em}
\begin{confidential}CONFIDENTIEL — PROJET RBK 2.0               Page 288 sur 316\end{confidential}
RBK 2.0 — Whitepaper Architecte Web3                                                                      Annexe V. Plan d’actions risques \textbackslash{}& tableau de bord KPI (v1)\\
\vspace{0.5em}
\vspace{0.5em}
\begin{center}\begin{tabular}TAB. V.2 : KPI de pilotage v1 — RBK Web3 Studio\end{tabular}\end{center}
Seuils                Actions\\
KPI            Définition         Formule / Calcul             Cible   Cadence   Source     Owner\\
Alerte/Critique       déclenchées\\
Leads ICP      Volume de pros- \textbackslash{}#leads qualifiés / semaine ≥ 40 / Hebdo           CRM +      Mkt       Alerte \textless{} 25 ; Cri-    Intensifier webi-\\
pects répondant                            sem                    forms                tique \textless{} 15            nars, partenariats,\\
à l’ICP (profil                                                                                              retargeting\\
technique +\\
motivation)\\
CVR L→E        Efficacité qualifi- \textbackslash{}#entretiens / \textbackslash{}#leads   ≥ 25\textbackslash{}% Hebdo            CRM        Mkt +     Alerte \textless{} 18\textbackslash{}% ;        Ajuster ICP, lan-\\
cation              qualifiés                                                HoP       Critique \textless{} 12\textbackslash{}%        ding, message,\\
screening\\
CVR E→A        Efficacité sélec- \textbackslash{}#admis / \textbackslash{}#entretiens          40–     Hebdo     CRM        HoP       Alerte \textless{} 30\textbackslash{}% ;        Revoir critères,\\
tion/adéquation                                 60\textbackslash{}%                                    Critique \textless{} 20\textbackslash{}%        renforcer bridge,\\
offre                                                                                                        clarifier prérequis\\
Taille cohorte Seuil minimal de \textbackslash{}#admis confirmés               ≥ 15    Hebdo     CRM        CEO       Alerte 12–14 ; Cri-   Activer plan B :\\
lancement                                                                              tique \textless{} 12            1 track, format\\
part-time, report\\
\vspace{0.5em}
\vspace{0.5em}
\vspace{0.5em}
\vspace{0.5em}
EL\\
CPL            Coût acquisition  Budget marketing / \textbackslash{}#leads À       Hebdo         Budget + Mkt         Alerte +20\textbackslash{}% ; Cri-    Optimiser canaux,\\
\vspace{0.5em}
\vspace{0.5em}
\vspace{0.5em}
\vspace{0.5em}
TI\\
qualifiés                 définir               CRM                  tique +40\textbackslash{}%            couper ads faibles\\
Attendance     Présence aux ses- \textbackslash{}#présences / \textbackslash{}#sessions    ≥ 90\textbackslash{}% Hebdo           LMS /    HoP         Alerte \textless{} 85\textbackslash{}% ;        1 :1 coaching,\\
\vspace{0.5em}
\vspace{0.5em}
\vspace{0.5em}
\vspace{0.5em}
N\\
sions synchrones attendues                                        Discord              Critique \textless{} 75\textbackslash{}%        rattrapages, ajuster\\
\vspace{0.5em}
\vspace{0.5em}
\vspace{0.5em}
\vspace{0.5em}
DE\\
charge\\
Labs complets Livraison des       \textbackslash{}#labs validés / \textbackslash{}#labs        ≥ 85\textbackslash{}% Hebdo       Repo       HoP +     Alerte \textless{} 75\textbackslash{}% ;        Réduire scope,\\
\vspace{0.5em}
\vspace{0.5em}
\vspace{0.5em}
\vspace{0.5em}
FI\\
labs dans les       prévus                                         issues/-   TA        Critique \textless{} 60\textbackslash{}%        ajouter TA, replani-\\
\vspace{0.5em}
\vspace{0.5em}
\vspace{0.5em}
\vspace{0.5em}
N\\
délais                                                             grades                                     fier sprint\\
\vspace{0.5em}
\vspace{0.5em}
CO\\
Jalons caps-  Respect jalons      \textbackslash{}#jalons atteints / \textbackslash{}#jalons   ≥ 80\textbackslash{}% Hebdo       Board      HoP       Alerte \textless{} 70\textbackslash{}% ;        Gating Must, cou-\\
tone          capstone (S17–      planifiés                                      projet               Critique \textless{} 50\textbackslash{}%        per Could, mento-\\
S20)                                                                                                          rat intensif\\
Churn         Décrochage          \textbackslash{}#abandons / \textbackslash{}#cohorte         ≤ 5\textbackslash{}%    Mensuel   Admin      HoP       Alerte 6–10\textbackslash{}% ; Cri-   Diagnostic charge,\\
apprenants                                                                              tique \textgreater{} 10\textbackslash{}%           accompagnement,\\
adaptation rythme\\
NPS pro-       Satisfaction       \textbackslash{}%promoteurs - \textbackslash{}%détrac-       ≥ 45    Mensuel   Survey     Ops +     Alerte 20–44 ; Cri-   Plan actions quali-\\
gramme         globale (recom-    teurs                                                     HoP       tique \textless{} 20            té (top 3 irritants)\\
mandation)\\
PR pass rate   PR acceptées       \textbackslash{}#PR pass / \textbackslash{}#PR totales       ≥ 70\textbackslash{}% Hebdo       GitHub     Tech Lead Alerte \textless{} 60\textbackslash{}% ;        Coaching code\\
sans rework                                                                            Critique \textless{} 45\textbackslash{}%        review, exemples,\\
majeur                                                                                                       pairing\\
PR rework      Proportion PR      \textbackslash{}#PR rework / \textbackslash{}#PR totales ≤ 25\textbackslash{}% Hebdo           GitHub     Tech Lead Alerte \textgreater{} 35\textbackslash{}% ;        Revoir DoD, tem-\\
demandant                                                                              Critique \textgreater{} 50\textbackslash{}%        plates, formation\\
refonte                                                                                                      tests\\
Suite page suivante\\
\vspace{0.5em}
\vspace{0.5em}
\vspace{0.5em}
\vspace{0.5em}
\begin{confidential}CONFIDENTIEL — PROJET RBK 2.0                                                 Page 289 sur 316\end{confidential}
RBK 2.0 — Whitepaper Architecte Web3                                                                        Annexe V. Plan d’actions risques \textbackslash{}& tableau de bord KPI (v1)\\
\vspace{0.5em}
\vspace{0.5em}
Seuils                    Actions\\
KPI            Définition         Formule / Calcul           Cible      Cadence    Source     Owner\\
Alerte/Critique           déclenchées\\
CI pass rate   Qualité pipeline   \textbackslash{}#runs CI OK / \textbackslash{}#runs CI     ≥ 85\textbackslash{}% Hebdo           GitHub     Ops +     Alerte \textless{} 75\textbackslash{}% ;            Stop merges, fix\\
CI (lint/tests)    totaux                                           Actions    Tech Lead Critique \textless{} 60\textbackslash{}%            flakes, stabiliser\\
tests\\
MT green CI    Vitesse de répara- Temps moyen entre CI       ≤ 6ℎ       Hebdo      GitHub     Ops       Alerte \textgreater{} 12h ; Cri-       Freeze merges +\\
tion CI            rouge et CI verte                                                     tique \textgreater{} 24h               task force\\
Couverture     Couverture tests Couverture \textbackslash{}% (ou proxy :     ≥ 60\textbackslash{}% Mensuel         CI reports Tech Lead Alerte \textless{} 50\textbackslash{}% ;            Sprint tests, obliga-\\
tests          minimale (selon \textbackslash{}#tests)                       (min)                                      Critique \textless{} 35\textbackslash{}%            tions DoD\\
repos)\\
Findings       Findings sécurité \textbackslash{}#findings High+Critical     0          Hebdo      Audit      SecLead   Alerte : 1 Critical ou    Bloquer release,\\
High/Crit      majeurs sur                                   Critical ;            notes                \textgreater{}2 High ; Critique :      security sprint,\\
capstones                                     ≤ 2                                        \textgreater{}=2 Critical              re-audit\\
High\\
TM conformes Threat model         \textbackslash{}#capstones avec TM validé 100\textbackslash{}%        Hebdo      Docs repo SecLead    Alerte \textless{} 90\textbackslash{}% ;            Refus validation\\
présent et validé    / total                               (capstone)                      Critique \textless{} 75\textbackslash{}%            milestone\\
Incidents    Fuites détectées     \textbackslash{}#incidents / période      0           Hebdo      Gitleaks / SecLead   Alerte : 1 ; Critique :   Rotation immé-\\
secrets      (keys, tokens)                                                        alerts               \textgreater{}1                        diate + incident\\
\vspace{0.5em}
\vspace{0.5em}
\vspace{0.5em}
\vspace{0.5em}
EL\\
report\\
\vspace{0.5em}
\vspace{0.5em}
\vspace{0.5em}
\vspace{0.5em}
TI\\
Erreurs RPC  Erreurs RPC          \textbackslash{}#erreurs / \textbackslash{}#requêtes (ou   ≤ 2\textbackslash{}%       Hebdo      Logs app   Ops       Alerte \textgreater{} 5\textbackslash{}% ; Cri-        Basculer provider,\\
(timeouts, rate-     logs)                                                                 tique \textgreater{} 10\textbackslash{}%               activer cache/-\\
\vspace{0.5em}
\vspace{0.5em}
\vspace{0.5em}
\vspace{0.5em}
N\\
limit)                                                                                                               mocks\\
\vspace{0.5em}
\vspace{0.5em}
\vspace{0.5em}
\vspace{0.5em}
DE\\
MTTR         Temps moyen de       Moyenne temps résolution ≤ 2ℎ   Mensuel          Incident Ops         Alerte \textgreater{} 4h ; Cri-        Revoir runbooks,\\
rétablissement       incidents                                        log                  tique \textgreater{} 8h                drills, ownership\\
\vspace{0.5em}
\vspace{0.5em}
\vspace{0.5em}
\vspace{0.5em}
FI\\
Portfolios   Portfolios prêts     \textbackslash{}#apprenants ready / total ≥ 90\textbackslash{}% Hebdo            Checklist Career     Alerte \textless{} 80\textbackslash{}% ;            Bootcamp portfolio\\
\vspace{0.5em}
\vspace{0.5em}
\vspace{0.5em}
\vspace{0.5em}
N\\
prêts        (repos + docs +                                      (fin)                                 Critique \textless{} 65\textbackslash{}%            + rework\\
\vspace{0.5em}
\vspace{0.5em}
CO\\
demo)\\
Placement 3m Placement (job/-     \textbackslash{}#placés / \textbackslash{}#diplômés        ≥ 50\textbackslash{}% Mensuel         Place-     Career + Alerte \textless{} 35\textbackslash{}% ;             Intensifier réseau,\\
mission/bounty                                                        ment       CEO      Critique \textless{} 20\textbackslash{}%             mock interviews,\\
régulier) à 3                                                         sheet                                          bounties\\
mois\\
Placement    Placement post-      \textbackslash{}#placés / \textbackslash{}#diplômés        ≥ 70\textbackslash{}% Mensuel         Place-     Career + Alerte \textless{} 55\textbackslash{}% ;             Partenariats RH +\\
post         parcours                                                              ment       CEO      Critique \textless{} 40\textbackslash{}%             suivi individuel\\
sheet\\
Bounties       Bounties com-      \textbackslash{}#bounties validées / co-   ≥ 1        Mensuel    Table      Career    Alerte \textless{} 0.5 ; Cri-       Sprint bounties,\\
complétées     plétées (Super-    horte                      par étu-              bounties             tique \textless{} 0.2               mentorat ciblé\\
team/OSS etc.)                                diant\\
(min)\\
Variance       Écart budget réel (Réel - Prévu) / Prévu      ≤     Mensuel         Budget     CEO +     Alerte +10–20\textbackslash{}% ;          Couper “Could”,\\
budget         vs prévisionnel                               +10\textbackslash{}%                  sheet      Ops       Critique \textgreater{} 20\textbackslash{}%            renégocier licences\\
Encaissement   Encaissement vs Encaissement / Factura-       ≥ 95\textbackslash{}% Mensuel         Finance    CEO       Alerte \textless{} 90\textbackslash{}% ;            Relances, ajuster\\
facturé           tion                                                                   Critique \textless{} 80\textbackslash{}%            modalités paiement\\
Suite page suivante\\
\vspace{0.5em}
\vspace{0.5em}
\vspace{0.5em}
\vspace{0.5em}
\begin{confidential}CONFIDENTIEL — PROJET RBK 2.0                                                 Page 290 sur 316\end{confidential}
RBK 2.0 — Whitepaper Architecte Web3                                                                      Annexe V. Plan d’actions risques \textbackslash{}& tableau de bord KPI (v1)\\
\vspace{0.5em}
\vspace{0.5em}
Seuils                Actions\\
KPI           Définition        Formule / Calcul             Cible     Cadence   Source    Owner\\
Alerte/Critique       déclenchées\\
Marge/        Marge brute par   (Recettes - coûts directs) / À         Mensuel   Finance   CEO     Alerte baisse 15\textbackslash{}% ;   Ajuster prix, coûts,\\
étudiant      étudiant          \textbackslash{}#étudiants                   définir                               Critique baisse 30\textbackslash{}%   taille cohorte\\
\vspace{0.5em}
\vspace{0.5em}
\vspace{0.5em}
\vspace{0.5em}
EL\\
TI\\
N\\
DE\\
FI\\
N\\
CO\\
\vspace{0.5em}
\vspace{0.5em}
\vspace{0.5em}
\vspace{0.5em}
\begin{confidential}CONFIDENTIEL — PROJET RBK 2.0                                               Page 291 sur 316\end{confidential}
RBK 2.0 — Whitepaper Architecte Web3                                                                          Annexe V. Plan d’actions risques \textbackslash{}& tableau de bord KPI (v1)\\
\vspace{0.5em}
\vspace{0.5em}
Recommandation d’usage.         Ce tableau doit être revu en COPIL (hebdo pendant cohorte) avec : (i) statut des actions risques, (ii) KPI en alerte, (iii) décisions consignées dans un journal de\\
\vspace{0.5em}
pilotage.\\
\vspace{0.5em}
\vspace{0.5em}
\vspace{0.5em}
\vspace{0.5em}
EL\\
TI\\
N\\
DE\\
FI\\
N\\
CO\\
\vspace{0.5em}
\vspace{0.5em}
\vspace{0.5em}
\vspace{0.5em}
\begin{confidential}CONFIDENTIEL — PROJET RBK 2.0                                                Page 292 sur 316\end{confidential}
RBK 2.0 — Whitepaper Architecte Web3                           Annexe V. Plan d’actions risques \textbackslash{}& tableau de bord KPI (v1)\\
\vspace{0.5em}
\vspace{0.5em}
\vspace{0.5em}
KPI Heatmap — Synthèse exécutive (hebdomadaire)\\
\vspace{0.5em}
\vspace{0.5em}
\vspace{0.5em}
\vspace{0.5em}
EL\\
TI\\
N\\
DE\\
FI\\
N\\
CO\\
\vspace{0.5em}
\vspace{0.5em}
\vspace{0.5em}
\vspace{0.5em}
\begin{confidential}CONFIDENTIEL — PROJET RBK 2.0               Page 293 sur 316\end{confidential}
RBK 2.0 — Whitepaper Architecte Web3                                                                                    Annexe V. Plan d’actions risques \textbackslash{}& tableau de bord KPI (v1)\\
\vspace{0.5em}
\vspace{0.5em}
Objectif. Donner une lecture immédiate de l’état du programme par domaine (Admissions, Qualité, Sécurité, Ops, Carrière, Finance), avec un score global et des déclencheurs d’actions.\\
\vspace{0.5em}
\vspace{0.5em}
Domaine                      Score global      Tendance       Statut                   Signal / Explication (KPI en alerte + action)\\
(0–5)\\
Admissions \textbackslash{}& Marketing       3.0/5             →              Alerte                   Alerte : CVR Lead→Entretien en baisse ; cohorte à 12 confirmés. Action : A-003 webinaires + A-002 pipeline admissions (Owner :\\
Mkt/Ops).\\
Pédagogie \textbackslash{}& Engagement       4.0/5             ↗              OK                       Attendance 92\textbackslash{}%, Lab completion 88\textbackslash{}%. Action : maintenir binômage A-017 et monitoring surcharge A-046.\\
Qualité Studio (PR/CI/-      3.5/5             →              Alerte légère            CI pass rate 78\textbackslash{}% proche seuil ; MT Green 10h. Action : freeze merges si \textgreater{}24h (A-020) + stabilisation tests.\\
Tests)\\
Sécurité                     3.0/5             ↘              Alerte                   Threat models incomplets (80\textbackslash{}%). Action : A-024 threat model obligatoire + A-054 security gate capstone (Owner : SecLead).\\
Ops \textbackslash{}& Infrastructure         4.5/5             →              OK                       RPC error rate 1.5\textbackslash{}% ; drills réalisés. Action : poursuivre runbooks + tests fallback (A-021/A-023).\\
Carrière \textbackslash{}& Employabilité     3.0/5             →              Alerte                   Portfolio readiness 72\textbackslash{}% (trop tôt ou retard). Action : Career pack A-037 + sprint bounties A-038 (Owner : Career).\\
Finance                      4.0/5             →              OK                       Budget variance +6\textbackslash{}% ; revenue collected 96\textbackslash{}%. Action : maintenir reporting variance A-008.\\
\vspace{0.5em}
\vspace{0.5em}
\vspace{0.5em}
\vspace{0.5em}
EL\\
TI\\
N\\
DE\\
FI\\
N\\
CO\\
\vspace{0.5em}
\vspace{0.5em}
\vspace{0.5em}
\vspace{0.5em}
\begin{confidential}CONFIDENTIEL — PROJET RBK 2.0                                                            Page 294 sur 316\end{confidential}
RBK 2.0 — Whitepaper Architecte Web3                                                                                   Annexe V. Plan d’actions risques \textbackslash{}& tableau de bord KPI (v1)\\
\vspace{0.5em}
\vspace{0.5em}
Méthode de calcul du score (proposition). Chaque domaine est noté à partir de 3 à 6 KPI clés :\\
\vspace{0.5em}
• 5 = tous KPI au-dessus des cibles ; 4 = 1 KPI en alerte légère ; 3 = 1 KPI en alerte + plan d’action actif ;\\
\vspace{0.5em}
• 2 = plusieurs KPI en alerte ou 1 KPI critique ; 1 = plusieurs KPI critiques ; 0 = domaine non pilotable (données manquantes ou incident majeur).\\
\vspace{0.5em}
Décision COPIL requise si : un domaine ≤ 2/5 ou 2 domaines en alerte simultanée ou un KPI “critique” déclenché.\\
\vspace{0.5em}
\vspace{0.5em}
\vspace{0.5em}
\vspace{0.5em}
EL\\
TI\\
N\\
DE\\
FI\\
N\\
CO\\
\vspace{0.5em}
\vspace{0.5em}
\vspace{0.5em}
\vspace{0.5em}
\begin{confidential}CONFIDENTIEL — PROJET RBK 2.0                                                               Page 295 sur 316\end{confidential}
RBK 2.0 — Whitepaper Architecte Web3                                                                              Annexe V. Plan d’actions risques \textbackslash{}& tableau de bord KPI (v1)\\
\vspace{0.5em}
\vspace{0.5em}
\vspace{0.5em}
Modèle de Compte Rendu — COPIL / COSUI (RBK Web3 Studio)\\
Type de comité : COPIL        ou      COSUI\\
\vspace{0.5em}
Projet : RBK Web3 Studio — Dual Track (EVM / Solana)\\
\vspace{0.5em}
Date :                                    Heure :                      Lieu / visio :\\
\vspace{0.5em}
Animateur :                                              Rédacteur :\\
\vspace{0.5em}
Référence CR : RBK-W3S-CR-                        Version :\\
\vspace{0.5em}
\vspace{0.5em}
\vspace{0.5em}
\vspace{0.5em}
1. Participants\\
\vspace{0.5em}
Nom                                      Rôle                       Organisation             Présence\\
\vspace{0.5em}
\vspace{0.5em}
\vspace{0.5em}
\vspace{0.5em}
EL\\
Présent / Excusé\\
Présent / Excusé\\
\vspace{0.5em}
\vspace{0.5em}
\vspace{0.5em}
\vspace{0.5em}
TI\\
Présent / Excusé\\
\vspace{0.5em}
\vspace{0.5em}
\vspace{0.5em}
\vspace{0.5em}
N\\
DE\\
FI\\
2. Ordre du jour\\
\vspace{0.5em}
\vspace{0.5em}
\vspace{0.5em}
N\\
1. Revue KPI \textbackslash{}& Heatmap (semaine                   )\\
CO\\
2. Avancement (pédagogie, contenus, tracks, capstones)\\
\vspace{0.5em}
3. Risques / incidents / décisions sécurité\\
\vspace{0.5em}
4. Budget \textbackslash{}& ressources (mentors, outils, infra)\\
\vspace{0.5em}
5. Carrière (placements, bounties, demo day)\\
\vspace{0.5em}
6. Décisions / arbitrages / actions\\
\vspace{0.5em}
7. Prochain comité\\
\vspace{0.5em}
\vspace{0.5em}
\vspace{0.5em}
\vspace{0.5em}
\begin{confidential}CONFIDENTIEL — PROJET RBK 2.0                                                                  Page 296 sur 316\end{confidential}
RBK 2.0 — Whitepaper Architecte Web3                                                   Annexe V. Plan d’actions risques \textbackslash{}& tableau de bord KPI (v1)\\
\vspace{0.5em}
\vspace{0.5em}
3. Revue KPI (synthèse)\\
\vspace{0.5em}
Domaine      Score        État                             KPI en alerte              Décision / action\\
Admissions                OK / Alerte / Critique\\
Qualité                   OK / Alerte / Critique\\
Sécurité                  OK / Alerte / Critique\\
Ops                       OK / Alerte / Critique\\
Carrière                  OK / Alerte / Critique\\
Finance                   OK / Alerte / Critique\\
\vspace{0.5em}
\vspace{0.5em}
\vspace{0.5em}
4. Avancement \textbackslash{}& livrables\\
Points clés (fait / en cours / à risque).\\
\vspace{0.5em}
\vspace{0.5em}
\vspace{0.5em}
\vspace{0.5em}
EL\\
• Fait :\\
\vspace{0.5em}
\vspace{0.5em}
\vspace{0.5em}
\vspace{0.5em}
TI\\
• En cours :\\
\vspace{0.5em}
\vspace{0.5em}
\vspace{0.5em}
\vspace{0.5em}
N\\
• À risque / bloqué :\\
\vspace{0.5em}
\vspace{0.5em}
\vspace{0.5em}
\vspace{0.5em}
DE\\
FI\\
Décisions techniques majeures (ADR).\\
\vspace{0.5em}
\vspace{0.5em}
\vspace{0.5em}
N\\
CO\\
• ADR-             :                                                           (Owner :                           )\\
\vspace{0.5em}
• ADR-             :                                                           (Owner :                           )\\
\vspace{0.5em}
\vspace{0.5em}
\vspace{0.5em}
5. Risques \textbackslash{}& incidents\\
Risque(s) en alerte / critique (extrait registre).\\
• R-           :                                                              (Action : A-       )\\
\vspace{0.5em}
• R-           :                                                              (Action : A-       )\\
\vspace{0.5em}
\vspace{0.5em}
\vspace{0.5em}
Incidents depuis le dernier comité (résumé).\\
• INC-                 :                                                          (Sévérité :         ; MTTR :          )\\
\vspace{0.5em}
\vspace{0.5em}
\vspace{0.5em}
\begin{confidential}CONFIDENTIEL — PROJET RBK 2.0                            Page 297 sur 316\end{confidential}
RBK 2.0 — Whitepaper Architecte Web3                                                                             Annexe V. Plan d’actions risques \textbackslash{}& tableau de bord KPI (v1)\\
\vspace{0.5em}
\vspace{0.5em}
6. Décisions \textbackslash{}& arbitrages\\
\vspace{0.5em}
Déc. ID              Décision                                                        Validation\\
D-                                                                                   Approuvé / Rejeté\\
D-                                                                                   Approuvé / Rejeté\\
\vspace{0.5em}
\vspace{0.5em}
\vspace{0.5em}
7. Plan d’actions (exécutable)\\
\vspace{0.5em}
Action                                                           Éché-\\
Action ID                                                           Owner                         Statut\\
(verbe + livrable)                                               ance\\
À faire /\\
A-                                                                                                En cours /\\
Fait\\
\vspace{0.5em}
\vspace{0.5em}
\vspace{0.5em}
\vspace{0.5em}
EL\\
À faire /\\
A-                                                                                                En cours /\\
\vspace{0.5em}
\vspace{0.5em}
\vspace{0.5em}
\vspace{0.5em}
TI\\
Fait\\
\vspace{0.5em}
\vspace{0.5em}
\vspace{0.5em}
\vspace{0.5em}
N\\
À faire /\\
A-                                                                                                En cours /\\
\vspace{0.5em}
\vspace{0.5em}
\vspace{0.5em}
\vspace{0.5em}
DE\\
Fait\\
\vspace{0.5em}
\vspace{0.5em}
\vspace{0.5em}
\vspace{0.5em}
FI\\
N\\
CO\\
8. Prochain comité\\
Date/heure :                                      Type : COPIL / COSUI     Objectif :\\
\vspace{0.5em}
Annexes : (1) Heatmap KPI, (2) Extrait registre risques, (3) Journal d’incidents, (4) ADRs nouvelles.\\
\vspace{0.5em}
\vspace{0.5em}
\vspace{0.5em}
\vspace{0.5em}
\begin{confidential}CONFIDENTIEL — PROJET RBK 2.0                                                              Page 298 sur 316\end{confidential}
RBK 2.0 — Whitepaper Architecte Web3                                                                                      Annexe V. Plan d’actions risques \textbackslash{}& tableau de bord KPI (v1)\\
\vspace{0.5em}
\vspace{0.5em}
\vspace{0.5em}
Journal d’incidents (v1) — Aligné MTTR / drills / runbooks\\
But. Tracer chaque incident (réel ou drill), mesurer MTTA/MTTR, relier aux runbooks et suivre les actions correctives jusqu’à clôture.\\
\vspace{0.5em}
\vspace{0.5em}
\vspace{0.5em}
\vspace{0.5em}
EL\\
TI\\
N\\
DE\\
FI\\
N\\
CO\\
\vspace{0.5em}
\vspace{0.5em}
\vspace{0.5em}
\vspace{0.5em}
\begin{confidential}CONFIDENTIEL — PROJET RBK 2.0                                                             Page 299 sur 316\end{confidential}
RBK 2.0 — Whitepaper Architecte Web3                                                                     Annexe V. Plan d’actions risques \textbackslash{}& tableau de bord KPI (v1)\\
\vspace{0.5em}
\vspace{0.5em}
\begin{center}\begin{tabular}TAB. V.8 : Journal d’incidents v1 — RBK Web3 Studio\end{tabular}\end{center}
Sympt.                                                 Actions                  Cause\\
ID     Date     Type    Sév.    Svc                          Source    Impact                                                       Réf.\\
(rés.)                                        MTTA MTTRrest.                     rac.\\
hh :mmhh :mm\\
\vspace{0.5em}
RBKRBOPS02\\
INC- 2025-01- Réel      Sev2    CI/CD   CI rouge (tests      GitHub    Blocage       00 :15 06 :10 Freeze merges ;    Cause : test      Postmortem\\
0001 15 10 :42                          instables) : échec   Actions + merges ;                    relancer tests ;   dépendant         lien\\
intermittent sur     alerte    risque retard               isoler flaky ;     timing RPC.\\
suite e2e.           Slack     sprint.                     hotfix pipeline.   Correctifs :\\
mocks +\\
timeouts +\\
retry policy.\\
Actions : A-020,\\
A-022.\\
RBKRBOPS01\\
INC- 2025-01- Drill     Sev3    RPC/    Scénario “RPC    Exercice      Formation : 00 :05 00 :45 Fallback             Améliorer         CR drill\\
\vspace{0.5em}
\vspace{0.5em}
\vspace{0.5em}
\vspace{0.5em}
EL\\
0002 22 14 :05                  Prv     down” : timeouts planifié      démontrer                    provider B ;      runbook RPC ; lien\\
et rate-limit.   (Ops)         fallback et                  activer cache ; ajouter\\
\vspace{0.5em}
\vspace{0.5em}
\vspace{0.5em}
\vspace{0.5em}
TI\\
UX                           message UX.       observabilité\\
\vspace{0.5em}
\vspace{0.5em}
\vspace{0.5em}
\vspace{0.5em}
N\\
dégradée.                                      requêtes.\\
\vspace{0.5em}
\vspace{0.5em}
\vspace{0.5em}
\vspace{0.5em}
DE\\
Actions : A-021,\\
A-023.\\
\vspace{0.5em}
\vspace{0.5em}
\vspace{0.5em}
\vspace{0.5em}
FI\\
RBKRBSEC01\\
INC-   2025-02- Réel    Sev1    Sécurité Secret détecté   Gitleaks     Risque fuite ; 00 :10 01 :05 Révoquer clé ;    Cause :           Postmortem\\
\vspace{0.5em}
\vspace{0.5em}
\vspace{0.5em}
\vspace{0.5em}
N\\
0003   03 09 :18                         dans commit (clé CI           rotation                     purge             mauvaise          lien\\
\vspace{0.5em}
CO\\
API test).                    nécessaire.                  historique ;      hygiène .env.\\
informer équipe ; Correctifs :\\
renforcer         templates\\
pre-commit.       .env.example +\\
formation +\\
scan obligatoire.\\
Actions : A-042,\\
A-043.\\
Sev1\\
Réel    Sev2\\
INC-            Drill   Sev3\\
\vspace{0.5em}
\vspace{0.5em}
\vspace{0.5em}
\vspace{0.5em}
\begin{confidential}CONFIDENTIEL — PROJET RBK 2.0                                           Page 300 sur 316\end{confidential}
RBK 2.0 — Whitepaper Architecte Web3                                                                                         Annexe V. Plan d’actions risques \textbackslash{}& tableau de bord KPI (v1)\\
\vspace{0.5em}
\vspace{0.5em}
Règles de gouvernance incidents (recommandation).\\
• Sev1/Sev2 : postmortem obligatoire (cause, timeline, correctifs, prévention) + action(s) dans registre risques.\\
\vspace{0.5em}
• Drills : consigner scénario, temps de bascule, amélioration runbook.\\
\vspace{0.5em}
• Objectif MTTR : suivre tendance mensuelle ; si dérive, déclencher session “runbook hardening”.\\
\vspace{0.5em}
\vspace{0.5em}
\vspace{0.5em}
\vspace{0.5em}
V.1 Annexe X — Traceability Matrix (SSOT \textbackslash{}& Release Gate)\\
But. Cette matrice rend la documentation audit-ready : chaque paramètre contractuel est traçable jusqu’à sa macro SSOT4 , ses points d’exposition dans les sections critiques, et les contrôles automatisés (local\\
\vspace{0.5em}
+ CI) qui empêchent toute régression.\\
\vspace{0.5em}
SSOT. Toutes les valeurs contractuelles sont centralisées dans tex/params.tex. Toute occurrence en dur (hors SSOT) dans les fichiers critiques est un NO-GO.\\
\vspace{0.5em}
\vspace{0.5em}
\vspace{0.5em}
\vspace{0.5em}
EL\\
TI\\
N\\
DE\\
FI\\
N\\
CO\\
\vspace{0.5em}
\vspace{0.5em}
\vspace{0.5em}
\vspace{0.5em}
4\\
Finance/Juridique/Compliance — Single Source of Truth (SSOT) : garantit que les paramètres contractuels et références sensibles sont définis en un point unique.\\
\vspace{0.5em}
\vspace{0.5em}
\begin{confidential}CONFIDENTIEL — PROJET RBK 2.0                                                               Page 301 sur 316\end{confidential}
RBK 2.0 — Whitepaper Architecte Web3                                              Annexe V. Plan d’actions risques \textbackslash{}& tableau de bord KPI (v1)\\
\vspace{0.5em}
\vspace{0.5em}
\begin{center}\begin{tabular}TAB. V.9 : Annexe X — Traceability Matrix (SSOT \textbackslash{}& Release Gate)\end{tabular}\end{center}
\vspace{0.5em}
Paramètre     Macro SSOT      Fichiers critiques concer- Tests / preuves        Gate\\
nés\\
Durée pro-    \textbackslash{}ProgramWeeks   00\textbackslash{}_factsheet.tex           verify\textbackslash{}_params.sh   GO : aucun\\
gramme        \textbackslash{}ProgramDuration 01a\textbackslash{}_note\textbackslash{}_cadrage.tex      (required/forbidden)\\
hardcode,\\
Text             annexe\textbackslash{}_j\textbackslash{}_offre.tex        verify\textbackslash{}_no\textbackslash{}_         cohérence SSOT.\\
annexe\textbackslash{}_b\textbackslash{}_finance.tex       hardcode.sh        NO-GO : durées en\\
CI                 clair hors macros\\
doc-release-checks (SSOT non\\
make release-check respectée).\\
ISA — Seuil    \textbackslash{}ISAThreshold   annexe\textbackslash{}_j\textbackslash{}_offre.tex         verify\textbackslash{}_no\textbackslash{}_             GO : seuil\\
\vspace{0.5em}
\vspace{0.5em}
\vspace{0.5em}
\vspace{0.5em}
EL\\
(brut)         TND             annexe\textbackslash{}_f\textbackslash{}_isa.tex           hardcode.sh            uniquement via\\
(pattern TND)          macro.\\
\vspace{0.5em}
\vspace{0.5em}
\vspace{0.5em}
\vspace{0.5em}
TI\\
annexe\textbackslash{}_q\textbackslash{}_contrat\textbackslash{}_isa.tex\\
annexe\textbackslash{}_b\textbackslash{}_finance.tex       CI pipelines           NO-GO : seuil en\\
\vspace{0.5em}
\vspace{0.5em}
\vspace{0.5em}
\vspace{0.5em}
N\\
00\textbackslash{}_factsheet.tex           make release-check     dur ou\\
confusion\\
\vspace{0.5em}
\vspace{0.5em}
\vspace{0.5em}
\vspace{0.5em}
DE\\
brut/net.\\
ISA — Taux                                                                      GO : taux\\
\vspace{0.5em}
\vspace{0.5em}
\vspace{0.5em}
\vspace{0.5em}
FI\\
\textbackslash{}ISARate        00\textbackslash{}_factsheet.tex           verify\textbackslash{}_params.sh\\
01a\textbackslash{}_note\textbackslash{}_cadrage.tex       (SSOT)                 uniquement via\\
\vspace{0.5em}
\vspace{0.5em}
N\\
annexe\textbackslash{}_j\textbackslash{}_offre.tex         verify\textbackslash{}_no\textbackslash{}_             macro.\\
CO\\
annexe\textbackslash{}_f\textbackslash{}_isa.tex           hardcode.sh            NO-GO : 15\textbackslash{}% en\\
annexe\textbackslash{}_q\textbackslash{}_contrat\textbackslash{}_isa.tex   (15\textbackslash{}% brut/revenu)      dur en\\
annexe\textbackslash{}_b\textbackslash{}_finance.tex       CI pipelines           contexte\\
make release-check     brut/revenu.\\
ISA — Cap     \textbackslash{}ISACapTND      annexe\textbackslash{}_j\textbackslash{}_offre.tex         verify\textbackslash{}_no\textbackslash{}_             GO : cap via\\
annexe\textbackslash{}_f\textbackslash{}_isa.tex           hardcode.sh            macro.\\
annexe\textbackslash{}_q\textbackslash{}_contrat\textbackslash{}_isa.tex   (pattern cap TND)      NO-GO : cap en\\
annexe\textbackslash{}_b\textbackslash{}_finance.tex       CI pipelines           dur hors SSOT.\\
make release-check\\
Suite page suivante\\
\vspace{0.5em}
\vspace{0.5em}
\vspace{0.5em}
\vspace{0.5em}
\begin{confidential}CONFIDENTIEL — PROJET RBK 2.0                               Page 302 sur 316\end{confidential}
RBK 2.0 — Whitepaper Architecte Web3                                                  Annexe V. Plan d’actions risques \textbackslash{}& tableau de bord KPI (v1)\\
\vspace{0.5em}
\vspace{0.5em}
Paramètre       Macro SSOT       Fichiers critiques concer- Tests / preuves         Gate\\
nés\\
ISA — Max       \textbackslash{}ISAMaxPaid      00\textbackslash{}_factsheet.tex              verify\textbackslash{}_params.sh     GO : durée max\\
paiements       Months           annexe\textbackslash{}_j\textbackslash{}_offre.tex            (required)           via macro.\\
annexe\textbackslash{}_f\textbackslash{}_isa.tex              verify\textbackslash{}_no\textbackslash{}_           NO-GO : 36 en dur\\
annexe\textbackslash{}_q\textbackslash{}_contrat\textbackslash{}_isa.tex      hardcode.sh          en\\
annexe\textbackslash{}_b\textbackslash{}_finance.tex          (36 +                contexte mois /\\
mensualités/mois)    mensualités.\\
CI pipelines\\
make release-check\\
Termes inter-   (n/a) — poli-    Tout le corpus                verify\textbackslash{}_params.sh    GO : aucun interdit\\
dits (durées    tique documen-   (sections programme / KPI /   (forbidden)         détecté.\\
courtes)        taire            annexes)                      CI pipelines        NO-GO : durées\\
make release-check courtes\\
\vspace{0.5em}
\vspace{0.5em}
\vspace{0.5em}
\vspace{0.5em}
EL\\
explicites\\
\vspace{0.5em}
\vspace{0.5em}
\vspace{0.5em}
\vspace{0.5em}
TI\\
détectées.\\
\vspace{0.5em}
\vspace{0.5em}
\vspace{0.5em}
\vspace{0.5em}
N\\
Release          (n/a) — pipeline Repo (racine)                 make release-check GO : CI verte +\\
\vspace{0.5em}
\vspace{0.5em}
\vspace{0.5em}
\vspace{0.5em}
DE\\
(preuve exé-                      CI pipelines                  (local)             artefact PDF.\\
\vspace{0.5em}
\vspace{0.5em}
\vspace{0.5em}
\vspace{0.5em}
FI\\
cutable)                                                        CI                  WARN : warnings\\
doc-release-checks LaTeX non\\
\vspace{0.5em}
\vspace{0.5em}
\vspace{0.5em}
\vspace{0.5em}
N\\
CO\\
(PDF + scripts)     contractuels.\\
docs/RELEASE\textbackslash{}_DOC.md NO-GO :\\
scripts en échec ou\\
build KO.\\
\vspace{0.5em}
\vspace{0.5em}
\vspace{0.5em}
\vspace{0.5em}
\begin{confidential}CONFIDENTIEL — PROJET RBK 2.0                                  Page 303 sur 316\end{confidential}
RBK 2.0 — Whitepaper Architecte Web3                                                                                         Annexe V. Plan d’actions risques \textbackslash{}& tableau de bord KPI (v1)\\
\vspace{0.5em}
\vspace{0.5em}
Note de périmètre. Les contrôles anti-hardcode ciblent explicitement une liste de fichiers critiques. Si une nouvelle section devient contractuelle (ou si une annexe est dupliquée), mettre à jour la liste dans\\
\vspace{0.5em}
scripts/verify\textbackslash{}_no\textbackslash{}_hardcode.sh.\\
\vspace{0.5em}
\vspace{0.5em}
\vspace{0.5em}
\vspace{0.5em}
V.2 Annexe Y — Change Control (SSOT \textbackslash{}& Release Gate)\\
Objet. Cette annexe définit un mécanisme de contrôle des changements des paramètres contractuels et des sections sensibles du livre blanc. Elle vise un statut audit-ready : toute évolution est tracée, justifiée,\\
\vspace{0.5em}
et prouvée par des contrôles automatisés (local + CI).\\
\vspace{0.5em}
Périmètre. Est considéré comme changement contrôlé toute modification de :\\
\vspace{0.5em}
• la SSOT tex/\\
\vspace{0.5em}
params.tex (durée programme, ISA : seuil/taux/cap/durée max, helpers de texte) ;\\
\vspace{0.5em}
\vspace{0.5em}
\vspace{0.5em}
\vspace{0.5em}
EL\\
• les scripts de contrôles scripts/\\
\vspace{0.5em}
\vspace{0.5em}
\vspace{0.5em}
\vspace{0.5em}
TI\\
verify\textbackslash{}_params.sh et scripts/\\
\vspace{0.5em}
\vspace{0.5em}
\vspace{0.5em}
\vspace{0.5em}
N\\
verify\textbackslash{}_no\textbackslash{}_hardcode.sh ;\\
\vspace{0.5em}
\vspace{0.5em}
\vspace{0.5em}
\vspace{0.5em}
DE\\
• le pipeline CI .github/workflows/\\
\vspace{0.5em}
doc-release-checks.yml ;\\
\vspace{0.5em}
\vspace{0.5em}
\vspace{0.5em}
\vspace{0.5em}
FI\\
• les sections contractuelles / financières exposant les paramètres (factsheet, offre, ISA, contrat ISA, finance).\\
\vspace{0.5em}
\vspace{0.5em}
\vspace{0.5em}
N\\
CO\\
Règle SSOT. Toute valeur contractuelle est définie uniquement dans tex/params.tex. Toute occurrence en dur (hors SSOT) dans les fichiers critiques constitue un NO-GO.\\
\vspace{0.5em}
Procédure obligatoire (avant merge / release).\\
\vspace{0.5em}
1. Mettre à jour la SSOT (si applicable) et/ou les sections concernées uniquement via macros.\\
\vspace{0.5em}
2. Exécuter la preuve locale : make release-check.\\
\vspace{0.5em}
3. Vérifier la preuve CI : workflow doc-release-checks vert + artefact PDF.\\
\vspace{0.5em}
4. Enregistrer le changement dans la table ci-dessous (une ligne par changement).\\
\vspace{0.5em}
\vspace{0.5em}
\vspace{0.5em}
\vspace{0.5em}
\begin{confidential}CONFIDENTIEL — PROJET RBK 2.0                                                                Page 304 sur 316\end{confidential}
RBK 2.0 — Whitepaper Architecte Web3                                                                                     Annexe V. Plan d’actions risques \textbackslash{}& tableau de bord KPI (v1)\\
\vspace{0.5em}
\vspace{0.5em}
\vspace{0.5em}
Date           Élément               Macro /               Justification \textbackslash{}& impact (sections)                                                             Preuves        Gate\\
Fichier\\
AAAA-          Paramètre             \textbackslash{}ISARate /            Raison : (ex : alignement offre commerciale).                                                 make            GO /\\
MM-JJ          contractuel           tex/                  Impact : Factsheet, Offre, ISA, Contrat ISA, Finance.                                         release-check WARN /\\
params.tex                                                                                                          OK ; CI verte ; NO-GO\\
commit\\
(hash)\\
AAAA-          Contrôle     verify\textbackslash{}_no\textbackslash{}_                     Raison : (ex : ajout d’un fichier critique).                                                  make            GO /\\
MM-JJ          documentaire hardcode.sh                    Impact : périmètre anti-régression.                                                           release-check WARN /\\
OK ; CI verte ; NO-GO\\
commit\\
(hash)\\
\vspace{0.5em}
\vspace{0.5em}
\vspace{0.5em}
\vspace{0.5em}
EL\\
AAAA-          Pipeline CI                                 Raison : (ex : fiabilisation build).                                                          CI verte ;     GO /\\
\vspace{0.5em}
\vspace{0.5em}
\vspace{0.5em}
\vspace{0.5em}
TI\\
doc-\\
MM-JJ                                                      Impact : reproductibilité + artefact PDF.                                                     artefact PDF ; WARN /\\
\vspace{0.5em}
\vspace{0.5em}
\vspace{0.5em}
\vspace{0.5em}
N\\
release-\\
\vspace{0.5em}
\vspace{0.5em}
\vspace{0.5em}
\vspace{0.5em}
DE\\
checks.yml                                                                                                          commit         NO-GO\\
(hash)\\
\vspace{0.5em}
\vspace{0.5em}
\vspace{0.5em}
\vspace{0.5em}
FI\\
N\\
Définition des gates.\\
CO\\
• GO : build OK ; verify\textbackslash{}_params PASSED ; verify\textbackslash{}_no\textbackslash{}_hardcode PASSED ; CI verte + artefact PDF.\\
\vspace{0.5em}
• WARN : tout OK, mais warnings LaTeX non contractuels (typographie, overfull hbox).\\
\vspace{0.5em}
• NO-GO : build KO ; ou scripts en échec ; ou hardcode contractuel détecté ; ou incohérence brut/net.\\
\vspace{0.5em}
\vspace{0.5em}
Note. Les valeurs contractuelles elles-mêmes restent dans la SSOT. Les preuves doivent pointer vers le commit et le run CI correspondant.\\
\vspace{0.5em}
Sign-off.\\
\vspace{0.5em}
Owner (responsable documentaire) :                                                                                                                            Date :\\
\vspace{0.5em}
Reviewer (sponsor / direction) :                                                                                                                             Date :\\
\vspace{0.5em}
\vspace{0.5em}
\vspace{0.5em}
\vspace{0.5em}
\begin{confidential}CONFIDENTIEL — PROJET RBK 2.0                                                             Page 305 sur 316\end{confidential}
W \textbar{} Glossaire des termes, acronymes et jargon\\
\vspace{0.5em}
\vspace{0.5em}
Cette annexe centralise les termes techniques, acronymes et notions de gouvernance utilisés dans le livre blanc. Chaque entrée est définie à la première occurrence via note de bas de page et rappelée\\
\vspace{0.5em}
ici pour consultation rapide.\\
\vspace{0.5em}
\vspace{0.5em}
\vspace{0.5em}
\vspace{0.5em}
EL\\
Terme\\
Domaine                            Définition\\
\vspace{0.5em}
\vspace{0.5em}
\vspace{0.5em}
\vspace{0.5em}
TI\\
/Acronyme\\
\vspace{0.5em}
\vspace{0.5em}
\vspace{0.5em}
\vspace{0.5em}
N\\
SSOT                                                                         Single Source of Truth (SSOT) : garantit que les paramètres contractuels et références sen-\\
\vspace{0.5em}
\vspace{0.5em}
\vspace{0.5em}
\vspace{0.5em}
DE\\
sibles sont définis en un point unique.\\
\vspace{0.5em}
IP                                                                           Intellectual Property (IP) : droits protégeant créations, marques, contenus et logiciels, enca-\\
\vspace{0.5em}
\vspace{0.5em}
\vspace{0.5em}
\vspace{0.5em}
FI\\
drant licence, cession ou exceptions.\\
\vspace{0.5em}
\vspace{0.5em}
\vspace{0.5em}
N\\
TS                                                                           TypeScript (TS) : sur-ensemble typé de JavaScript offrant vérification statique et tooling\\
CO\\
pour des applications fiables.\\
\vspace{0.5em}
IA                                                                           Intelligence Artificielle (IA) : modèles statistiques ou LLM appliqués aux produits pour auto-\\
matiser ou augmenter des parcours et décisions.\\
\vspace{0.5em}
open-source                                                                  Modèle de diffusion de logiciel dont le code source est accessible, réutilisable et modifiable\\
sous licence ouverte.\\
\vspace{0.5em}
USD                                                                          United States Dollar (USD) : devise de référence utilisée pour exprimer des coûts, revenus\\
ou comparaisons internationales.\\
\vspace{0.5em}
SaaS                                                                         Software as a Service (SaaS) : logiciel opéré et facturé comme un service en abonnement,\\
hébergé et maintenu par le fournisseur.\\
\vspace{0.5em}
\vspace{0.5em}
\vspace{0.5em}
\vspace{0.5em}
306\\
RBK 2.0 — Whitepaper Architecte Web3                                               Annexe W. Glossaire des termes, acronymes et jargon\\
\vspace{0.5em}
\vspace{0.5em}
Terme\\
Domaine    Définition\\
/Acronyme\\
TechLead                                Référent technique (TechLead) coordonnant architecture, qualité et mentoring d’équipe\\
pour livrer un produit robuste.\\
\vspace{0.5em}
SecLead                                 Référent sécurité (SecLead) pilotant analyse de risques, contrôles et plans d’actions lors des\\
projets ou incidents.\\
\vspace{0.5em}
TND                                     Dinar tunisien, devise utilisée pour exprimer les montants dans l’offre et les projections\\
financières du programme.\\
\vspace{0.5em}
NDA                                     Non-Disclosure Agreement (NDA) : accord encadrant la confidentialité et l’usage des infor-\\
mations échangées entre parties.\\
\vspace{0.5em}
KO                                      Knock-Out (KO) : statut indiquant l’échec d’un contrôle ou d’un gate de release, nécessitant\\
\vspace{0.5em}
\vspace{0.5em}
\vspace{0.5em}
\vspace{0.5em}
EL\\
correction avant validation.\\
\vspace{0.5em}
\vspace{0.5em}
\vspace{0.5em}
\vspace{0.5em}
TI\\
DualTrack                               Organisation en deux parcours parallèles (ex. EVM et Solana) partageant un tronc commun\\
puis une spécialisation.\\
\vspace{0.5em}
\vspace{0.5em}
\vspace{0.5em}
\vspace{0.5em}
N\\
DE\\
Airdrop                                 Distribution gratuite ou ciblée de tokens à une communauté pour récompenser l’usage ou\\
stimuler l’adoption.\\
\vspace{0.5em}
\vspace{0.5em}
\vspace{0.5em}
\vspace{0.5em}
FI\\
Whitelist                               Whitelist : liste d’adresses ou d’identités autorisées à interagir avec un contrat, souvent\\
\vspace{0.5em}
\vspace{0.5em}
\vspace{0.5em}
N\\
utilisée pour les mint ou accès anticipés.\\
CO\\
CRM                                     Customer Relationship Management (CRM) : outil ou processus de gestion de la relation\\
client (prospects, ventes, support, suivi).\\
\vspace{0.5em}
LinkedIn                                Réseau professionnel utilisé pour la visibilité, l’acquisition de leads et le recrutement de\\
partenaires ou candidats.\\
\vspace{0.5em}
NPS                                     Net Promoter Score (NPS) : indicateur de recommandation (promoteurs - détracteurs) pour\\
mesurer satisfaction et fidélité.\\
\vspace{0.5em}
DoD                                     Definition of Done (DoD) : liste les critères minimaux pour considérer un élément comme\\
terminé et livrable.\\
\vspace{0.5em}
MoSCoW                                  Must, Should, Could, Won’t (MoSCoW) : méthode de priorisation cadrant le périmètre et\\
arbitrant les fonctionnalités.\\
\vspace{0.5em}
\vspace{0.5em}
\vspace{0.5em}
\begin{confidential}CONFIDENTIEL — PROJET RBK 2.0                    Page 307 sur 316\end{confidential}
RBK 2.0 — Whitepaper Architecte Web3                                               Annexe W. Glossaire des termes, acronymes et jargon\\
\vspace{0.5em}
\vspace{0.5em}
Terme\\
Domaine    Définition\\
/Acronyme\\
Sprint 0                                Phase amont de 90 jours dédiée à sécuriser juridique, finances, MVP pédagogique et staf-\\
fing avant toute communication publique du programme.\\
\vspace{0.5em}
SLA                                     Service Level Agreement contractuel décrivant engagements de service, métriques, pénali-\\
tés et modalités de mesure.\\
\vspace{0.5em}
AML                                     Anti-Money Laundering (AML) : lutte contre le blanchiment d’argent couvrant surveillance\\
des flux, filtrage des sanctions et déclarations aux autorités.\\
\vspace{0.5em}
KYC                                     Know Your Customer (KYC) : processus de connaissance client vérifiant l’identité avant\\
d’accéder à un service financier ou réglementé.\\
\vspace{0.5em}
RGPD                                    General Data Protection Regulation (GDPR / RGPD) : règlement encadrant collecte, traite-\\
\vspace{0.5em}
\vspace{0.5em}
\vspace{0.5em}
\vspace{0.5em}
EL\\
ment et droits des personnes en Europe.\\
\vspace{0.5em}
\vspace{0.5em}
\vspace{0.5em}
\vspace{0.5em}
TI\\
ROI                                     Return on Investment (ROI) : rapport entre gain généré et investissement engagé.\\
\vspace{0.5em}
\vspace{0.5em}
\vspace{0.5em}
\vspace{0.5em}
N\\
SBT                                     Soulbound Token non transférable, attaché à une identité ou un compte pour prouver un\\
\vspace{0.5em}
\vspace{0.5em}
\vspace{0.5em}
\vspace{0.5em}
DE\\
statut, une attestation ou une réussite.\\
\vspace{0.5em}
Taux de conversion                      Proportion d’utilisateurs qui passent à l’étape suivante du funnel ou réalisent l’action cible.\\
\vspace{0.5em}
\vspace{0.5em}
\vspace{0.5em}
\vspace{0.5em}
FI\\
N\\
Funnel                                  Parcours d’acquisition décomposé en étapes (awareness → consideration → conversion)\\
avec mesures à chaque étape.\\
CO\\
Rétention                               Capacité à conserver clients ou utilisateurs sur la durée, inverse du churn.\\
\vspace{0.5em}
Tokenomics                              Token economics (tokenomics) : conception économique d’un token couvrant émission,\\
distribution, utilité, incitations et mécanismes de capture de valeur.\\
\vspace{0.5em}
TVL                                     Total Value Locked (TVL) : valeur totale des actifs déposés ou immobilisés dans des proto-\\
coles DeFi, indicateur de liquidité et de traction.\\
\vspace{0.5em}
RBAC                                    Role-Based Access Control (RBAC) : modèle limitant les privilèges par rôles prédéfinis pour\\
réduire la surface d’attaque.\\
\vspace{0.5em}
API                                     Application Programming Interface (API) : interface de programmation exposant des ser-\\
vices ou données pour être consommés par d’autres applications.\\
\vspace{0.5em}
\vspace{0.5em}
\vspace{0.5em}
\vspace{0.5em}
\begin{confidential}CONFIDENTIEL — PROJET RBK 2.0                    Page 308 sur 316\end{confidential}
RBK 2.0 — Whitepaper Architecte Web3                                               Annexe W. Glossaire des termes, acronymes et jargon\\
\vspace{0.5em}
\vspace{0.5em}
Terme\\
Domaine    Définition\\
/Acronyme\\
Bridge                                  Mécanisme de transfert ou de verrouillage d’actifs entre blockchains distinctes, basé sur des\\
preuves ou des validateurs fédérés.\\
\vspace{0.5em}
Staking                                 Staking : blocage de tokens pour sécuriser un réseau Proof-of-Stake et recevoir des récom-\\
penses en échange de la participation au consensus.\\
\vspace{0.5em}
EVM                                     Ethereum Virtual Machine (EVM) : moteur d’exécution des smart contracts bytecode com-\\
patibles sur Ethereum et chaînes EVM, assurant portabilité des dApps et outils.\\
\vspace{0.5em}
Multisig                                Portefeuille multi-signature exigeant plusieurs signatures indépendantes pour autoriser une\\
transaction, réduisant le risque de compromission.\\
\vspace{0.5em}
Ops                                     « Ops » désigne les pratiques et activités d’exploitation (monitoring, incidents, runbooks,\\
\vspace{0.5em}
\vspace{0.5em}
\vspace{0.5em}
\vspace{0.5em}
EL\\
disponibilité).\\
\vspace{0.5em}
\vspace{0.5em}
\vspace{0.5em}
\vspace{0.5em}
TI\\
Token-2022                              SPL Token-2022 : standard Solana ajoutant des extensions (règles, hooks, métadonnées\\
étendues) pour des jetons programmables.\\
\vspace{0.5em}
\vspace{0.5em}
\vspace{0.5em}
\vspace{0.5em}
N\\
DE\\
On-chain                                Donnée ou logique exécutée et stockée directement sur la blockchain, vérifiée par le\\
consensus du réseau.\\
\vspace{0.5em}
\vspace{0.5em}
\vspace{0.5em}
\vspace{0.5em}
FI\\
GitHub                                  Plateforme de développement collaborative pour héberger code, pull requests, CI et revue\\
\vspace{0.5em}
\vspace{0.5em}
\vspace{0.5em}
N\\
entre équipes.\\
CO\\
Runbook                                 Guide opérationnel décrivant les étapes techniques pour exécuter une tâche récurrente ou\\
un diagnostic.\\
\vspace{0.5em}
Threat Model                            Threat Model : analyse structurée des actifs, surfaces d’attaque et scénarios adverses afin de\\
prioriser les contre-mesures de sécurité.\\
\vspace{0.5em}
README                                  Fichier d’introduction d’un dépôt décrivant le projet, son installation et son usage.\\
\vspace{0.5em}
ERC                                     Ethereum Request for Comments (ERC) : famille de standards décrivant des interfaces de\\
jetons et contrats (ex. ERC-20, ERC-721) dans l’écosystème Ethereum.\\
\vspace{0.5em}
UX                                      User Experience (UX) : qualité perçue du parcours, de la compréhension et de l’efficacité\\
pour l’utilisateur final.\\
\vspace{0.5em}
\vspace{0.5em}
\vspace{0.5em}
\vspace{0.5em}
\begin{confidential}CONFIDENTIEL — PROJET RBK 2.0                    Page 309 sur 316\end{confidential}
RBK 2.0 — Whitepaper Architecte Web3                                              Annexe W. Glossaire des termes, acronymes et jargon\\
\vspace{0.5em}
\vspace{0.5em}
Terme\\
Domaine    Définition\\
/Acronyme\\
OpenAPI                                 OpenAPI Specification (OpenAPI) : spécification standardisée décrivant de façon contrac-\\
tuelle les endpoints, schémas et flux d’une API HTTP.\\
\vspace{0.5em}
RPC                                     Remote Procedure Call (RPC) : interface exposant des points de terminaison pour lire l’état\\
de la blockchain ou soumettre des transactions.\\
\vspace{0.5em}
Gate GO/NO-GO                           Point de décision formalisé (GO/NO-GO) après vérification build, scripts et absence de\\
hardcodes ; bloque la promotion en cas d’échec.\\
\vspace{0.5em}
CLI                                     Command Line Interface (CLI) : interface en ligne de commande permettant d’exécuter des\\
actions et scripts via un terminal.\\
\vspace{0.5em}
Responsable Bien-être                   Référent Nexus habilité à ajuster charge, deadlines et protocoles de repos pour maintenir la\\
\vspace{0.5em}
\vspace{0.5em}
\vspace{0.5em}
\vspace{0.5em}
EL\\
santé mentale et physique des apprenants.\\
\vspace{0.5em}
\vspace{0.5em}
\vspace{0.5em}
\vspace{0.5em}
TI\\
Certification mentor RBK                Processus d’évaluation et de validation des mentors RBK/Nexus couvrant expertise tech-\\
nique, soft skills et conformité aux standards de mentoring Senior-by-Design.\\
\vspace{0.5em}
\vspace{0.5em}
\vspace{0.5em}
\vspace{0.5em}
N\\
DE\\
Electric Capital                        Source : Electric Capital Developer Report 2023.\\
\vspace{0.5em}
Web3.career                             Données agrégées par Web3.career \textbackslash{}& TrueUp Tech Jobs Report, Q4 2024.\\
\vspace{0.5em}
\vspace{0.5em}
\vspace{0.5em}
\vspace{0.5em}
FI\\
N\\
Stablecoin                              Stablecoin : actif numérique cherchant à maintenir une parité stable via collatéral, algo-\\
rithme ou réserves hors chaîne auditées.\\
CO\\
DAO                                     Decentralized Autonomous Organization (DAO) : organisation autonome décentralisée où\\
les détenteurs de jetons votent les décisions et les budgets via des propositions transpa-\\
rentes.\\
\vspace{0.5em}
CI/CD                                   Continuous Integration / Continuous Delivery (CI/CD) : chaîne d’intégration et de livraison\\
continues automatisant tests, builds, audits et déploiements.\\
\vspace{0.5em}
Smart contract                          Smart contract : programme déployé sur blockchain exécutant automatiquement des règles\\
vérifiables, immuables et auditées par le réseau.\\
\vspace{0.5em}
RACI                                    Responsible, Accountable, Consulted, Informed (RACI) : matrice clarifiant qui porte les\\
décisions, exécute, consulte et informe pour chaque activité.\\
\vspace{0.5em}
\vspace{0.5em}
\vspace{0.5em}
\vspace{0.5em}
\begin{confidential}CONFIDENTIEL — PROJET RBK 2.0                    Page 310 sur 316\end{confidential}
RBK 2.0 — Whitepaper Architecte Web3                                                  Annexe W. Glossaire des termes, acronymes et jargon\\
\vspace{0.5em}
\vspace{0.5em}
Terme\\
Domaine    Définition\\
/Acronyme\\
MVP                                         Minimum Viable Product (MVP) : produit minimum viable permettant de valider une hypo-\\
thèse avec l’effort minimal avant itérations.\\
\vspace{0.5em}
DeFi                                        Decentralized Finance (DeFi) : finance décentralisée construite sur des smart contracts,\\
offrant prêts, échanges et rendements sans intermédiaire bancaire traditionnel.\\
\vspace{0.5em}
ERC-20                                      Ethereum Request for Comments 20 (ERC-20) : standard définissant l’interface minimale\\
d’un jeton fongible (transfert, approbation, total supply).\\
\vspace{0.5em}
Senior-by-Design                            Principe de conception du programme visant à inculquer dès la formation les réflexes et\\
standards d’un profil senior.\\
\vspace{0.5em}
Architecte Web3                             Ingénieur capable de concevoir des systèmes décentralisés complets (Smart Contracts +\\
\vspace{0.5em}
\vspace{0.5em}
\vspace{0.5em}
\vspace{0.5em}
EL\\
Frontend + Indexing), en maîtrisant les enjeux de sécurité, de coût (Gas) et de perfor-\\
mance.\\
\vspace{0.5em}
\vspace{0.5em}
\vspace{0.5em}
\vspace{0.5em}
TI\\
Web3                                        Web 3.0 (Web3) : génération d’Internet décentralisée basée sur la propriété numérique via\\
\vspace{0.5em}
\vspace{0.5em}
\vspace{0.5em}
\vspace{0.5em}
N\\
la blockchain, redonnant le contrôle aux utilisateurs.\\
\vspace{0.5em}
\vspace{0.5em}
\vspace{0.5em}
\vspace{0.5em}
DE\\
Illustration Concrete                       Exemple concret illustratif utilisé pour faciliter la compréhension d’un concept.\\
\vspace{0.5em}
\vspace{0.5em}
\vspace{0.5em}
\vspace{0.5em}
FI\\
KPI                                         Key Performance Indicator (KPI) : indicateur clé mesurant l’atteinte d’un objectif straté-\\
\vspace{0.5em}
\vspace{0.5em}
\vspace{0.5em}
N\\
gique ou opérationnel.\\
CO\\
Contrat de Partenariat Pédago-              Accord conditionnant le remboursement à l’obtention d’un revenu dans le périmètre Web3,\\
gique                                       encadrant droits, obligations et procédures de recours pour les apprenants RBK.\\
\vspace{0.5em}
ISA                                         Income Share Agreement (ISA) : contrat où l’étudiant rembourse une part de revenu selon\\
des paramètres centralisés dans la SSOT.\\
\vspace{0.5em}
Fonds de Garantie ISA                       Compte séquestré capitalisant les contributions RBK/ISA pour couvrir défauts temporaires,\\
avances de bourses et incidents de paiement, gouverné conjointement avec Nexus.\\
\vspace{0.5em}
Maître d’œuvre pédagogique                  Partenaire responsable contractuellement de la conception, de l’exécution et de la qualité\\
pédagogique au nom de RBK, avec obligations de moyens renforcées et reporting dédié.\\
\vspace{0.5em}
\vspace{0.5em}
\vspace{0.5em}
\vspace{0.5em}
\begin{confidential}CONFIDENTIEL — PROJET RBK 2.0                        Page 311 sur 316\end{confidential}
Journal des modifications — Version 5.0\\
\vspace{0.5em}
\vspace{0.5em}
Cette version 5.0 du Livre Blanc RBK marque l’intégration formelle des outils de pilotage stratégique et la densification du volet ”Business \textbackslash{}& Operations”.\\
\vspace{0.5em}
\vspace{0.5em}
\vspace{0.5em}
\vspace{0.5em}
EL\\
TI\\
Ajouts Majeurs\\
\vspace{0.5em}
\vspace{0.5em}
\vspace{0.5em}
\vspace{0.5em}
N\\
• Direction Stratégique : Ajout ”Note de Cadrage” (Chap 1-A) avec SWOT, MoSCoW, RACI et Registre des Risques remplis.\\
\vspace{0.5em}
\vspace{0.5em}
\vspace{0.5em}
\vspace{0.5em}
DE\\
• Modèle Économique : Mise à jour Business Plan (Hypothèses, Funnel, Sources) et Factsheet unique (Chap 00).\\
\vspace{0.5em}
\vspace{0.5em}
\vspace{0.5em}
\vspace{0.5em}
FI\\
• Conformité : Ajout Guide Operating Model (Scénarios A/B/C) et Disclaimers légaux (Chap 16).\\
\vspace{0.5em}
\vspace{0.5em}
\vspace{0.5em}
N\\
• Méthodologie : Intégration du standard ”Studio” (DoD, Incident Drills) (Chap 4).\\
CO\\
• Outillage LaTeX : Création macros haute fidélité \textbackslash{}RBKSWOT etc. et harmonisation versionnelle (5.0).\\
\vspace{0.5em}
\vspace{0.5em}
\vspace{0.5em}
V5.4 — Industrialisation du ”Studio Operating System”\\
• Capstones : Spécifications industrielles (User Stories, Checklist Secu) + Rubric de notation unifiée (100pts).\\
• Compétences : Mapping précis Module → Badge → Preuve Technique (Repo) → Critère.\\
• Process : Définition du DoD Global et de la Politique de Pull Request (Template, Review).\\
• Outillage : Ajout des templates ”Incident Postmortem”, ”Repo Structure” et ”ADR” en annexe.\\
\vspace{0.5em}
\vspace{0.5em}
\vspace{0.5em}
\vspace{0.5em}
312\\
RBK 2.0 — Whitepaper Architecte Web3                                                                    Annexe W. Glossaire des termes, acronymes et jargon\\
\vspace{0.5em}
\vspace{0.5em}
\vspace{0.5em}
V5.3 — Sécurisation Réglementaire \textbackslash{}& Modèles Opérationnels\\
• Operating Models : Détaillage des 3 scénarios (Local, Export, Hybride) avec analyse d’impact et risques.\\
• Matrice Compliance : Ajout d’une vue dédiée aux risques Légal/Fiscal/Réputation.\\
• Communication Policy : Définition des ”Red Lines” (Promesses de gains) et obligations éthiques.\\
• Disclaimers : Renforcement des avertissements (”Ce document n’est pas un avis juridique”).\\
\vspace{0.5em}
\vspace{0.5em}
\vspace{0.5em}
\vspace{0.5em}
EL\\
V5.2 — Audit de Cohérence \textbackslash{}& Contenu (Factsheet)\\
• Factsheet Unique : Centralisation de toutes les métriques clés (Durée 24s, Prix, KPIs) dans le Chapitre 00.\\
\vspace{0.5em}
\vspace{0.5em}
\vspace{0.5em}
\vspace{0.5em}
TI\\
• Cohérence : Alignement des durées (16+8 sem) et des objectifs financiers (Break-even M+12) entre Note de Cadrage et Business Plan.\\
• Contenu Stratégique : Renseignement complet des matrices SWOT, MoSCoW, RACI et Risques au sein de la Note de Cadrage.\\
\vspace{0.5em}
\vspace{0.5em}
\vspace{0.5em}
\vspace{0.5em}
EN\\
• Disclaimers : Ajout systématique des avertissements légaux (Pas de conseil financier).\\
\vspace{0.5em}
\vspace{0.5em}
\vspace{0.5em}
\vspace{0.5em}
D\\
V5.1 — Édition Professionnelle (Forme \textbackslash{}& Finitions)\\
\vspace{0.5em}
\vspace{0.5em}
FI\\
• Standards LaTeX : Adoption de acro, threeparttable, bookmark.\\
• Typographie : Correction des césures agressives (Short Titles), espaces insécables, polices mathématiques (XITS).\\
N\\
• Métadonnées : Intégration complète des métadonnées PDF.\\
O\\
• Navigation : Ajouts de liste des figures, des tableaux et des acronymes automatiques.\\
C\\
\vspace{0.5em}
\vspace{0.5em}
Modifications Structurelles\\
• Conservation intégrale du contenu v4 (aucune suppression).\\
• Réorganisation de l’ordre de lecture pour prioriser le cadrage stratégique.\\
• Optimisation des largeurs de tableaux (tabularx) pour le format A4.\\
\vspace{0.5em}
\vspace{0.5em}
\vspace{0.5em}
\vspace{0.5em}
\begin{confidential}CONFIDENTIEL — PROJET RBK 2.0                                            Page 313 sur 316\end{confidential}
RBK 2.0 — Whitepaper Architecte Web3                                                    Annexe W. Glossaire des termes, acronymes et jargon\\
\vspace{0.5em}
\vspace{0.5em}
\vspace{0.5em}
Statut de Validation\\
Document prêt pour revue par le CODIR et validation finale ”Publisher-Grade”.\\
\vspace{0.5em}
\vspace{0.5em}
\vspace{0.5em}
\vspace{0.5em}
EL\\
TI\\
N\\
DE\\
FI\\
N\\
CO\\
\vspace{0.5em}
\vspace{0.5em}
\vspace{0.5em}
\vspace{0.5em}
\begin{confidential}CONFIDENTIEL — PROJET RBK 2.0                                        Page 314 sur 316\end{confidential}
REFERENCES \textbackslash{}& BIBLIOGRAPHIE\\
\vspace{0.5em}
\vspace{0.5em}
\vspace{0.5em}
\vspace{0.5em}
EL\\
TI\\
W.1 Documentation Technique\\
• Solana Docs : https://docs.solana.com\\
\vspace{0.5em}
\vspace{0.5em}
\vspace{0.5em}
\vspace{0.5em}
EN\\
• Anchor Framework : https://www.anchor-lang.com\\
• OtterSec Blog : https://osec.io/blog\\
\vspace{0.5em}
\vspace{0.5em}
\vspace{0.5em}
\vspace{0.5em}
D\\
FI\\
W.2 Rapports de l’Industrie\\
• Electric Capital Developer Report (2024)\\
N\\
• Messari State of Solana (Q4 2024)\\
O\\
• Solana Whitepaper (2017) : Proof of History.\\
C\\
\vspace{0.5em}
• The Rust Book : Steve Klabnik \textbackslash{}& Carol Nichols.\\
\vspace{0.5em}
\vspace{0.5em}
\vspace{0.5em}
W.3 Rapports de Marché\\
• HackerOne Security Report.\\
• Superteam Earn Metrics.\\
\vspace{0.5em}
\vspace{0.5em}
\vspace{0.5em}
\vspace{0.5em}
315\\
RBK 2.0 — Whitepaper Architecte Web3                                Annexe W. Glossaire des termes, acronymes et jargon\\
\vspace{0.5em}
\vspace{0.5em}
\vspace{0.5em}
W.4 Outils Cités\\
• Helius : Observabilité Solana.\\
• Trident : Solana Fuzzing Framework.\\
• Metaplex : Standard NFT.\\
\vspace{0.5em}
\vspace{0.5em}
\vspace{0.5em}
\vspace{0.5em}
EL\\
TI\\
N\\
DE\\
FI\\
N\\
CO\\
\vspace{0.5em}
\vspace{0.5em}
\vspace{0.5em}
\vspace{0.5em}
\begin{confidential}CONFIDENTIEL — PROJET RBK 2.0                    Page 316 sur 316\end{confidential}
\vspace{0.5em}
\end{document}
